\documentclass[a4paper,10pt]{article}
\usepackage[top=25mm,bottom=24mm,left=27mm,right=27mm,headheight=15pt,headsep=8mm,footskip=12mm]{geometry}
\usepackage{fontspec,amsmath,amsthm,unicode-math}
\usepackage[english,french]{babel}
\usepackage{microtype,xcolor,fancyhdr,titlesec,booktabs,longtable,array,tabularx}
\definecolor{accent}{HTML}{184C63}
\definecolor{muted}{HTML}{52616B}
\definecolor{ink}{HTML}{172A36}
\usepackage{graphicx,tikz,enumitem,float,needspace,fvextra}
\usetikzlibrary{arrows.meta,calc,positioning,decorations.pathreplacing,matrix}
\usepackage[unicode,hidelinks]{hyperref}
\usepackage[nameinlink,noabbrev]{cleveref}
\setmainfont{STIXTwoText-Regular.otf}[BoldFont=STIXTwoText-Bold.otf,ItalicFont=STIXTwoText-Italic.otf,BoldItalicFont=STIXTwoText-BoldItalic.otf]
\setmathfont{STIXTwoMath-Regular.otf}
\hypersetup{pdftitle={Relations structurelles entre les constantes physiques},pdfauthor={Oumar Haidara Fall; Rubén Ramos Balsa},pdfsubject={Action, réciprocité géométrique et mémoire des transformations}}
\newcommand{\RR}{\mathbb R}\newcommand{\ZZ}{\mathbb Z}\newcommand{\QQ}{\mathbb Q}\newcommand{\CC}{\mathbb C}\newcommand{\FF}{\mathbb F}
\newcommand{\Id}{\operatorname{id}}\newcommand{\tr}{\operatorname{tr}}\newcommand{\Tr}{\operatorname{Tr}}
\newcommand{\Spec}{\operatorname{Spec}}
\newcommand{\Span}{\operatorname{span}}\newcommand{\Cyl}{\operatorname{Cyl}}\newcommand{\dd}{\mathrm d}
\newcommand{\square}{\mdlgwhtsquare}
\newcommand{\TPK}{\mathrm{TPK}}\newcommand{\APP}{\mathrm{APP}}\newcommand{\TRIT}{\{-1,0,+1\}}\newcommand{\HMT}{\mathrm{HMT}}
\newtheorem{theorem}{Théorème}[section]
\newtheorem{lemma}[theorem]{Lemme}\newtheorem{proposition}[theorem]{Proposition}\newtheorem{corollary}[theorem]{Corollaire}
\theoremstyle{definition}\newtheorem{definition}[theorem]{Définition}\newtheorem{axiom}[theorem]{Axiome}
\theoremstyle{remark}\newtheorem{remark}[theorem]{Remarque}
\numberwithin{equation}{section}
\setlength{\parindent}{1em}\setlength{\parskip}{3pt plus 1pt minus .5pt}
\setlength{\emergencystretch}{3em}
\clubpenalty=10000\widowpenalty=10000\displaywidowpenalty=10000
\renewcommand{\normalsize}{\fontsize{10.5}{13.4}\selectfont}\normalsize
\setlist{nosep,topsep=5pt}
\titleformat{\section}{\fontsize{13}{16}\selectfont\bfseries}{\thesection}{.65em}{}
\titleformat{\subsection}{\fontsize{11.2}{14}\selectfont\bfseries}{\thesubsection}{.65em}{}
\titleformat{\subsubsection}{\normalsize\bfseries}{\thesubsubsection}{.65em}{}
\AddToHook{cmd/section/before}{\Needspace{10\baselineskip}}
\AddToHook{cmd/subsection/before}{\Needspace{7\baselineskip}}
\AddToHook{cmd/subsubsection/before}{\Needspace{6\baselineskip}}
\setcounter{tocdepth}{2}
\makeatletter
\renewcommand*\l@subsection{\@dottedtocline{2}{1.5em}{3.9em}}
\makeatother
\titlespacing*{\section}{0pt}{17pt}{9pt}\titlespacing*{\subsection}{0pt}{13pt}{6pt}
\pagestyle{fancy}\fancyhf{}
\fancyhead[L]{\fontsize{8}{10}\selectfont Relations structurelles entre les constantes physiques}
\fancyfoot[C]{\thepage}\renewcommand{\headrulewidth}{.2pt}\raggedbottom
\makeatletter
\@ifundefined{example}{\theoremstyle{definition}\newtheorem{example}[theorem]{Exemple}}{}
\@ifundefined{notatabla}{\newcommand{\notatabla}[1]{\par\smallskip{\footnotesize #1\par}}}{}
\makeatother
\providecommand{\FF}{\mathbb F}

% Apertura independiente del índice: corrección quirúrgica 20260928.
\let\HMTIndiceOriginal\tableofcontents
\renewcommand{\tableofcontents}{\clearpage\begingroup\setlength{\parskip}{0pt}\fontsize{9.2}{10.5}\selectfont\HMTIndiceOriginal\endgroup\clearpage}
\begin{document}
\thispagestyle{plain}
\begin{center}
{\fontsize{10}{12}\selectfont Manuscrit de compilation pour le transfert académique\par}
\vspace{1.5mm}
{\fontsize{17}{20.5}\selectfont\bfseries Relations structurelles\par entre les constantes physiques\par}
\vspace{1.5mm}
{\fontsize{11}{13.5}\selectfont Action, réciprocité géométrique\par et mémoire des transformations\par}
\vspace{1.5mm}
{\fontsize{11.5}{14}\selectfont Oumar Haidara Fall \textperiodcentered\ Rubén Ramos Balsa\par}
\vspace{1.5mm}
{\fontsize{8}{10}\selectfont Coautorat sans préséance de contribution\par}
\vspace{1.5mm}
{\fontsize{8}{10}\selectfont Intégration et compilation documentaires générées par ChatGPT - Astra.\par}
\end{center}
\vspace{1.5mm}
\begingroup\fontsize{9.5}{10.7}\selectfont\setlength{\parskip}{1.5pt}\renewenvironment{abstract}{\begin{center}\bfseries\abstractname\end{center}\noindent}{\par}\begin{abstract}
On établit des relations de compatibilité, de réciprocité et de reconstruction
entre les structures d’action, de masse et de réponse constitutive obtenues
par l’holographie modulaire triadique et la mécanique dimensionnelle. Le développement
part de la composition de la structure arithmétique des chemins
\emph{All Possible Paths} (APP), de la signature ternaire orientée TRIT et du
\emph{Topological Phase Kernel} (TPK), qui met à jour et transporte l’état
enrichi. Cette composition conserve orientation, retenue, mémoire et
retour. Les constantes interviennent comme évaluations des
structures préalablement construites ; leur composition permet de déterminer quelles
propriétés elles partagent et quelle information une observation conjointe récupère.
Dans la collection constitutive normalisée, on démontre que la vitesse de
propagation et la fonctionnelle de Barbero--Immirzi déterminent les valeurs propres
étiquetées des deux canaux. Les logarithmes de vitesse et d’impédance
constituent le gradient d’un potentiel spectral strictement concave.
Sa hessienne caractérise la sensibilité de l’inversion et fournit des bornes
explicites de stabilité. Le potentiel partage un moment spectral de
profondeur deux avec la construction orientée dont la valeur extrême est Catalan ;
la réponse constitutive restitue cette collection de manière unique avec sa
condition de frontière à l’origine.

La composition de quarts de tour associés à des formes elliptiques conjuguées
produit un transport relatif unimodulaire dont les dilatations réciproques
permettent de reconstruire le coefficient adimensionnel d’action, la
coordonnée de clôture et l’orientation étant conservées. Le relèvement dodécaphasique et
l’itération du transport caractérisent la mémoire de l’ordre des
opérations. L’analyse chronologique distingue la croissance de
l’information du taux entropique par position, et la réentrée de mémoire
détermine sa contribution à la réponse opératoire. Le retour inscrit et
sa composition polaire déterminent une longueur $L_\alpha$. Conserver
conjointement action et aire fixe $G=c^3L_\alpha^2/\hbar$ dans la réalisation
radiale, également sous des variations non commutatives et une réduction compatible
de mémoire.

La mémoire électronique est restituée par des invariants linéaires et
spectraux qui conservent l’ordre du registre ; le transport entre
multisections admet un critère explicite de conservation de masse.
La composition thermique réunit capacités, multiplicités et mémoire
spectrale dans une trace exacte, et construit sa réponse thermométrique
pour une lecture marquée à référence fixe. Cette réalisation relie
énergie, information, action et aire, en conservant les conditions de
normalisation et d’observation de chaque grandeur.
La composition transportée des réponses et la reconstruction
Higgs--impédance relient les canaux constitutifs au secteur
électrofaible. La substitution électronique conserve son échelle dans Fermi,
la simplifie dans les quotients angulaires et transporte la longueur inverse
par une congruence opératoire, également sans commutation.

Enfin, sous les
conditions circulaire et thermique spécifiées, une même coordonnée de masse
détermine la fréquence propre, deux longueurs réciproques de produit
invariant et la pondération thermique du cycle. Les résultats réunissent ainsi
production, compatibilité et restitution dans une description mathématique
commune, avec des démonstrations explicites et la distinction entre changement de
représentation, transport relatif et réalisation physique.


Le lecteur dodécaphasique uniforme détermine également une suite de premières racines de retour critique. On démontre que le rapport de leurs écarts converge vers la constante de Feigenbaum, en conservant la sélection initiale, l’orientation et la mémoire de la construction.

% Adición al resumen de VII; no sustituye su contenido anterior.
La même composition se prolonge en une action totale de géométrie, de jauge et de matière : elle incorpore conjointement la gravitation et les interactions forte, faible et électromagnétique, en conservant les représentations chirales et les couplages précédemment générés. On obtient son courant de spin, la source métrique totale et les termes torsionnels croisés entre espèces. La transformation de Legendre et les identités de symétrie produisent des contraintes conjointes de première classe et leur propagation. Sur leur surface régulière, on calcule la courbure de la connexion canonique totale et l’on démontre son annulation dans les directions de déformation caractéristiques. Le résultat conserve les conditions de bord et l’éventuelle holonomie globale de l’état enrichi.



\par\smallskip\noindent\textbf{Mots-clés :} action ; réciprocité géométrique ; mémoire opératoire ; constantes physiques ; holographie modulaire triadique ; mécanique dimensionnelle.
\end{abstract}


\endgroup

\clearpage
\begingroup
\setlength{\parskip}{0pt}
\fontsize{10}{11.5}\selectfont
\tableofcontents
\endgroup
\clearpage
\clearpage
\begingroup\setlength{\parskip}{2pt}\fontsize{10}{11.8}\selectfont
\section{Introduction}

\label{md:introduccion}
La lecture d’une constante comme valeur scalaire représente le terme final
d’une opération. Lorsqu’on conserve la structure qui produit cette lecture,
d’autres questions deviennent accessibles : quelles relations elle partage avec les autres
lectures, quelle transformation conserve son orientation et quelle partie de
l’histoire peut être restituée à partir d’observations conjointes. Le présent
développement situe ces questions dans la continuité entre la structure discrète
HMT et ses réalisations en mécanique dimensionnelle.

L’origine commune conserve les opérations d’APP, le régime orienté du
TRIT et la composition effective du TPK. Le retour de phase se distingue de
l’évolution de l’état enrichi. Les coordonnées régionales et leur
clôture dodécaphasique sont ensuite transmises aux lecteurs d’action,
d’incidence et de masse. Chaque transmission conserve les objets et les conditions
qu’elle utilise ; une sortie peut intervenir dans la construction suivante sans
perdre sa provenance. Le noyau initial expose cette architecture avec ses
opérations, ses prolongements et ses preuves. Les antécédents sectoriels qui
suivent matérialisent les dépendances employées par les compositions du
présent travail.

La question centrale peut être formulée en trois termes complémentaires :
\emph{production}, \emph{compatibilité} et \emph{restitution}. La
production fixe l’objet et sa généalogie. La compatibilité détermine quelles
relations ses différentes lectures doivent satisfaire. La restitution établit
quelles données de l’objet peuvent être récupérées à partir de celles-ci. Le parcours inverse
est une propriété des lectures déjà construites et conserve leurs bases,
orientations et domaines ; il ne remplace pas leur ordre générateur.

Le premier axe du manuscrit réunit l’action et la paire angulaire. Le
coefficient adimensionnel de la section d’action participe au
contre-angle ; celui-ci détermine, avec l’angle, les canaux qui alimentent la
réponse constitutive et la fonctionnelle de Barbero--Immirzi. La même section
intervient dans le quotient de retour de la lecture électronique. De cette
manière, substituer les structures dans leurs équations révèle des dépendances
partagées qu’une énumération de constantes séparées laisse implicites.
La distinction entre coefficient et grandeur dimensionnelle est essentielle :
restituer un rapport de forme et restituer une action dimensionnelle nécessitent
des données de référence différentes.

Le second axe étudie l’information conservée dans la réponse. La
vitesse et l’impédance normalisées retiennent des aspects distincts de la
répartition entre canaux. La composition permet de démontrer tant une inversion
globale dans la collection déclarée qu’une relation différentielle entre les deux
réponses. Un potentiel spectral organise cette relation, et sa hessienne
contrôle la sensibilité de l’inversion. On distingue ainsi deux propriétés :
que les données exactes déterminent l’objet de manière unique et qu’une précision
finie permette de le reconstruire avec une erreur bornée. La collection orientée dont la
valeur extrême est Catalan est incorporée par son opération spectrale,
en conservant ses arguments et sa profondeur.

La composition transportée des réponses fournit sa coordonnée
additive ; les signatures angulaires de jauge et de Higgs permettent de reconstruire
la même paire de canaux. Les moments du défaut cyclique complètent
cette lecture avec des bornes de convergence et une représentation de Mellin.
Plus loin, la section électronique est substituée dans Fermi et dans l’échelle
de Higgs : la preuve conserve quels facteurs demeurent, lesquels se simplifient
et comment se transforme l’inverse de l’opérateur de masse sans supposer de commutation.

Le troisième axe transpose la discussion à l'ordre des opérations. Le plan de quart de tour du cycle dodécaphasique et les formes conjuguées de l'ellipse permettent de construire un commutateur explicite. Le bilan des incréments de phase se ferme tandis que l'action sur l'état conserve une dilatation réciproque. Deux programmes ayant les mêmes multiplicités d'opérations peuvent produire des réponses différentes en raison de leur ordre. Cette différence admet des détecteurs et des dispositifs de restitution définis : la trace enregistre certains invariants, tandis que les réponses sur des axes orientés distinguent des sens que la trace identifie.

La lecture physique est présentée avec les conditions de chaque réalisation. Dans la section gravitationnelle circulaire, la longueur de Compton réduite et le rayon gravitationnel réduit se dilatent réciproquement et maintiennent une aire commune. Dans la réalisation thermique du cycle, la même coordonnée structurelle intervient comme exposant d'un poids d'équilibre. Dans le transport elliptique, la conservation de l'aire symplectique et celle d'une énergie quadratique commune sont examinées au moyen de leurs opérateurs respectifs. La réalisation expérimentale éventuelle du programme relatif exige de spécifier conjointement le système, la référence, les opérations disponibles et le contrôleur. Ces conditions appartiennent au sens du résultat et restent contiguës à ses preuves.

La normalisation gravitationnelle est développée avant d'utiliser ses conséquences scalaires. Le retour de Hadamard, l'inscription des incidences et la norme locale de $\alpha$ produisent une section de longueur. Son transport polaire et le même caractère de parcours produisent deux lecteurs réciproques. L'énergie conserve à son tour la longueur d'action. En réunissant $R\Lambda=L_\alpha^2I$ et $H\Lambda=\hbar cI$, la constante gravitationnelle est fixée comme coefficient commun de tous les modes positifs de la réalisation. L'horloge circulaire s'exprime ensuite à partir de la longueur et de la vitesse. Cette antériorité permet d'expliquer la normalisation sans utiliser la valeur expérimentale de $G$ ni une longueur de Planck tabulée.

Les compositions électronique et thermique précisent comment la mémoire intervient dans ces relations. Dans la première, on reconstruit le registre dodécaphasique à partir de sa somme, de ses différences et de son orientation ; le spectre distingue des ordres partageant les mêmes multiplicités. Dans la seconde, le spectre géométrique déjà construit par le transport de mémoire apporte la raison $1/10$. Les capacités chronologiques et les secteurs d'occupation déterminent une trace marquée qui se compose avec un état de Gibbs. L'énergie conditionnée, l'information marquée et la température résultante sont spécifiées séparément ; ce choix d'observations permet de démontrer le transducteur thermométrique et de suivre sa relation avec l'échelle d'action et l'aire. Les sections correspondantes conservent les données de préparation, les graduations et les hypothèses de variation utilisées.

L'exposition développe d'abord les antécédents, conserve ensuite la narration explicative du parcours et présente les compositions complètes. L'étude chronologique des vacances, l'épisode régional de réitération visible et la réentrée de mémoire permettent de suivre l'origine des questions de compatibilité. Les développements énergétiques et gaussiens sont conservés en appendices comme réalisations auxiliaires, avec leurs hypothèses et leurs preuves complètes. Les conclusions réunissent les caractérisations structurelles obtenues. Le supplément numérique préserve les sources de travail littérales, les programmes et leurs résultats ; le manuscrit contient les démonstrations employées et distingue leur lecture mathématique des contrôles finis qui les accompagnent.

Cette organisation permet de formuler précisément l'apport transversal du travail. Une constante conserve sa valeur numérique, mais sa définition structurelle spécifie également l'opération qui la produit, le domaine dans lequel elle agit et les transformations qu'elle conserve. Deux grandeurs partageant une dépendance peuvent donc se restreindre mutuellement : connaître une réponse et son orientation permet de reconstruire une répartition qu'une lecture scalaire isolée ne distingue pas. Le problème inverse est posé sur l'image déjà engendrée et conserve ses données de référence. Génération et reconstruction s'articulent ainsi sans inverser l'antériorité de la construction.

La même distinction ordonne la relation entre mémoire et information. L'état présent, le bilan accumulé et l'opérateur qui répondra à une continuation sont des objets différents. Le retour d'une phase peut coexister avec une transformation persistante de l'état ; deux suites ayant des dénombrements identiques peuvent se distinguer par l'ordre de composition. Les preuves développent cette différence au moyen d'opérateurs, d'invariants et d'observations permettant la restitution. Leur intérêt physique réside dans la précision de ce que conserve chaque description et de l'opération permettant d'accéder à l'information retenue.

Le passage aux grandeurs dimensionnelles s'effectue en maintenant les bases et les normalisations correspondantes. La reconstruction d'un coefficient d'action, la conservation d'une aire et la caractérisation d'une température relèvent de résultats liés, de types différents. L'exposition réunit ces résultats dans leurs conditions de validité et montre comment ils se composent. Les conclusions explicitent aussi bien la caractérisation obtenue pour chaque objet que la relation qu'il établit avec les autres structures.

La constante de Feigenbaum reçoit une définition structurelle précise : elle est l’invariant asymptotique de séparation des paramètres sélectionnés par les retours critiques du lecteur dodécaphasique uniforme. La construction conserve sa provenance APP--TRIT--TPK, le continu conjoint et l’ordre des premières racines. L’identification des centres sélectionnés aux centres analytiques permet d’appliquer le théorème d’échelle à cette même suite ; les démonstrations de \ref{hmtfg:escalado:15}--\ref{hmtfg:escalado:17} et \eqref{hmtfg:articulacion:limite} établissent le résultat paramétrique. La représentation de Jacobi restitue le profil à partir de sa mesure à douze nœuds, tandis que l’inversion constitutive restitue les canaux orientés et le moment dont l’extrémité est la constante de Catalan. Ces opérations manifestent les relations entre lecteurs d’un même antécédent, en conservant le domaine et l’information propres à chacun.

La composition se prolonge également en une action commune de géométrie, de jauge et de matière. La gravitation et les interactions forte, faible et électromagnétique interviennent par les réalisations du même état enrichi, avec les couplages issus de la construction précédente. La variation conjointe détermine les sources et la réponse torsionnelle ; la transformation de Legendre produit une algèbre commune de contraintes de première classe. Leur propagation et la fermeture de la connexion canonique sont démontrées sur la surface régulière des contraintes, en conservant les conditions au bord et l’holonomie globale. L’unification variationnelle et canonique apporte ainsi une composition dynamique explicite aux relations de compatibilité et de restitution qui organisent le manuscrit.

\par\endgroup
\clearpage
% Núcleo común de I conservado literalmente; registro ampliado por inclusión.
\begingroup
\makeatletter
\def\input@path{{base_articulo_I/}}
\makeatother
\section{Noyau formel de l'holographie modulaire triadique}
\label{sec:nucleo}

L'holographie modulaire triadique, abrégée HMT, commence par une structure finie de positions et d'opérations. L'ordre de construction importe : les positions reçoivent des évaluations arithmétiques, les transitions acquièrent une signature ternaire et la dynamique conserve l'information nécessaire pour les prolonger. Ce n'est qu'ensuite que les réalisations numériques sont formées. Une constante reconnue à la fin de ce processus ne remplace ni la règle qui a produit sa région ni l'histoire qui détermine son évaluation.

La dénomination \emph{All Possible Paths}, APP, exprime le rôle des chemins sur le support. Elle est adoptée comme dénomination de la construction, avec une référence conceptuelle à la formulation par chemins de la mécanique quantique ; elle ne suppose pas que sa loi arithmétique soit l'intégrale de chemins de Feynman. Le \emph{Topological Phase Kernel}, TPK,\footnote{Les auteurs dédient également les initiales TPK à Tan Pengkiang. Cette dédicace est indépendante de la définition mathématique.} désigne la composition de la sélection, du transport, de la mise à jour de la mémoire et du prolongement. Les sigles ne sont maintenus qu'après la définition des objets qu'ils représentent.

\subsection{APP : support positionnel et évaluations arithmétiques}

Considérons les neuf marques intérieures d'une droite graduée entre zéro et dix,
\[
 D_9=\{1,2,3,4,5,6,7,8,9\}.
\]
Leur disposition horizontale et verticale produit les quatre-vingt-une positions de \(D_9^2\). Les marques intérieures ne doivent pas être confondues avec les dix intervalles unitaires de la droite de référence. Cette disposition fixe un alphabet et deux directions ; elle n'utilise pas encore un développement réel pour décider d'une trajectoire.

L'identification cyclique des positions conduit au support
\[
 X_9=(\ZZ/9\ZZ)^2.
\]
La coordinatisation par représentants positifs conserve le symbole neuf pour la classe nulle. Pour tout entier positif, on définit
\begin{equation}
 \rho_9(n)=1+((n-1)\bmod9),\qquad
 q_9(n)=\frac{n-\rho_9(n)}9,
 \qquad n=\rho_9(n)+9q_9(n).
 \label{eq:residuo-cociente}
\end{equation}
La racine numérique positive est donc un choix de représentant de la classe résiduelle. Elle ne conserve pas à elle seule l'entier antérieur à la réduction ; la coordonnée \(q_9\) fournit précisément l'information manquante.

À chaque position \((i,j)\), l'accumulation de \(i\) répétée \(j\) fois produit \(ij\). Les deux évaluations entières et leurs projections sont
\[
 S(i,j)=i+j,\qquad P(i,j)=ij,\qquad
 \Sigma(i,j)=\rho_9(i+j),\quad \Pi(i,j)=\rho_9(ij).
\]
La symétrie \((i,j)\mapsto(j,i)\) échange les directions d'accumulation et conserve les deux évaluations. C'est une propriété d'un unique support, et non une identification ultérieure entre deux matrices indépendantes.

\begin{definition}[Évaluation arithmétique relevée]
L'évaluation d'APP qui conserve les deux feuillets est
\begin{equation}
 \mathcal A(i,j)=
 \bigl((R^+_{ij},q^+_{ij}),(R^\times_{ij},q^\times_{ij})\bigr)
 =\bigl((\rho_9(i+j),q_9(i+j)),(\rho_9(ij),q_9(ij))\bigr).
 \label{eq:app-levantada}
\end{equation}
Le terme \emph{feuillet} indique ici une évaluation sur le même support : additive ou multiplicative. Les deux coordonnées de quotient appartiennent à l'état arithmétique et ne sont pas éliminées lors de la réduction des résidus.
\end{definition}

\input{figures/app.tex}

\begin{proposition}[Reconstruction des évaluations]
Les deux couples de \eqref{eq:app-levantada} déterminent exactement \(i+j\) et \(ij\). Les résidus isolés ne déterminent ni ces évaluations entières ni leurs quotients.
\end{proposition}
\begin{proof}
La première affirmation est \eqref{eq:residuo-cociente} appliquée à chaque feuillet. Pour la seconde, \(10=1+9\) et \(19=1+2\cdot9\) ont le même résidu et des quotients distincts. Sur le support bidimensionnel, les positions \((5,5)\) et \((8,2)\) donnent la même somme et le même résidu du produit, mais les produits \(25=7+18\) et \(16=7+9\). La projection résiduelle identifie des données que l'évaluation relevée distingue.
\end{proof}

\subsubsection{Graphe de Cayley, topologie toroïdale et degré d'enroulement}

Les translations cardinales sont données, en coordonnées de ligne et de colonne, par
\[
 \nu_N=(-1,0),\quad\nu_E=(0,1),\quad
 \nu_S=(1,0),\quad\nu_O=(0,-1),\qquad T_d(x)=x+\nu_d.
\]
La première coordonnée augmente vers le sud et la seconde vers l'est. Le graphe
\[
 \mathscr G_9=\operatorname{Cay}(X_9,\{\nu_N,\nu_E,\nu_S,\nu_O\})
\]
est le réseau cyclique bidimensionnel : il possède quatre-vingt-un sommets et quatre arêtes orientées sortantes par sommet. C'est un graphe de Cayley pour la structure \emph{additive} du support. La multiplication des résidus possède des diviseurs de zéro et ne transforme pas toutes les lignes d'APP en un groupe multiplicatif.

Un mot cardinal \(w=d_1\cdots d_r\) et une origine \(x_0\) déterminent
\[
 x_k=x_0+\sum_{j=1}^k\nu_{d_j}\pmod9.
\]
Il existe \(81\cdot4^r\) chemins orientés de longueur \(r\), retours et retours en arrière compris. Lorsque le chemin se ferme, son déplacement entier avant réduction définit
\begin{equation}
 \operatorname{wind}(w)=\frac19
 \bigl(\#S(w)-\#N(w),\#E(w)-\#O(w)\bigr)\in\ZZ^2.
\end{equation}
L'inversion du chemin change le signe de ce vecteur. Le retour à la même position résiduelle n'impose donc pas que l'enroulement soit nul. C'est le premier exemple d'une différence qui réapparaîtra dans le retour de phase : une coordonnée peut se refermer tandis qu'une autre enregistre une évolution non triviale.

\subsubsection{Quotient de transport et identité de cocycle}

Soit \(s_9:\ZZ/9\ZZ\to D_9\) la section des représentants positifs. Le quotient de transport associé à la composition additive, également appelé \emph{retenue} en arithmétique positionnelle, est
\[
 c_9(a,b)=\frac{s_9(a)+s_9(b)-s_9(a+b)}9.
\]
Le caractère compositionnel de cette mémoire s'exprime par une identité exacte.

\begin{lemma}[Identité de cocycle]
Pour \(a,b,c\in\ZZ/9\ZZ\),
\[
 c_9(a,b)+c_9(a+b,c)=c_9(b,c)+c_9(a,b+c).
\]
\end{lemma}
\begin{proof}
Les deux sommes sont égales à
\[
 \frac{s_9(a)+s_9(b)+s_9(c)-s_9(a+b+c)}9.
\]
L'associativité de la somme des représentants relevés requiert précisément ce terme supplémentaire lors de la composition dans le quotient.
\end{proof}

Avec des représentants positifs, ce cocycle n'est pas normalisé en l'identité : \(c_9(0,a)=1\). Si \(s_0\) prend ses représentants dans \(\{0,\ldots,8\}\), sa retenue normalisée \(c_0\) satisfait
\[
 c_9(a,b)=c_0(a,b)+\mathbf1_{a=0}+\mathbf1_{b=0}-\mathbf1_{a+b=0}.
\]
La différence est le changement de section ; elle n'altère pas l'identité de cocycle. La décomposition ultérieure du calendrier en phase et fenêtre utilisera les représentants de zéro à huit.

Les totaux des deux évaluations distinguent la contribution visible et celle conservée dans le quotient :
\[
 \sum S=810,\quad\sum P=2025,\quad
 \sum\Sigma=405,\quad\sum\Pi=459,
\]
\[
 \sum q^+=45,\qquad\sum q^\times=174.
\]
On en obtient
\begin{equation}
 2025-810=(459-405)+9(174-45)=54+9\cdot129.
 \label{eq:defecto-integrado}
\end{equation}
La différence somme--produit ne se réduit pas au défaut visible \(54\). L'identité conserve également le défaut de quotient \(129\). Le supprimer modifierait l'arithmétique transportée, même si le tableau résiduel demeurait identique.

\subsubsection{Réduction ternaire et radical multiplicatif}

L'anneau \(R=\ZZ/9\ZZ\) a pour groupe des unités \(\{1,2,4,5,7,8\}\) et pour idéal maximal
\[
 \mathfrak m=(3)=\{0,3,6\},\qquad\mathfrak m^2=0.
\]
Toute ligne additive est une permutation des neuf représentants. Une ligne multiplicative l'est exactement lorsque son indice est une unité. Les lignes trois et six contiennent trois valeurs répétées ; la ligne neuf contient uniquement le représentant de la classe nulle. Le repliement multiplicatif distingue donc déjà, avant toute géométrie continue, les unités du secteur nilpotent.

L'application \(x\mapsto3x\) a un noyau et une image égaux à \(\mathfrak m\). Elle induit un isomorphisme d'espaces vectoriels sur \(\FF_3\),
\[
 R/\mathfrak m\longrightarrow\mathfrak m,\qquad [x]\longmapsto3x,
\]
mais pas un isomorphisme d'anneaux. En particulier, la suite visible \(3,6,9\) admet deux lectures différentes : sa réduction directe modulo trois est \(0,0,0\), tandis que l'extraction du facteur trois suivie de la réduction de sa coordonnée interne donne \(1,2,0\). Cette seconde opération conserve une coordonnée de l'idéal ; ce n'est pas la même projection que la réduction directe du nombre original.

\subsubsection{Réalisation latine de l'évaluation additive}

La différence entre les deux évaluations peut également s'exprimer au moyen d'
opérateurs. Soit \((e_r)_{r\in\ZZ/9\ZZ}\) la base orthonormale de \(\CC^9\).
L'affectation \(v_{ab}=e_{a+b}\) produit une base orthonormale dans chaque ligne
et chaque colonne. Les projecteurs \(P_{ab}=v_{ab}v_{ab}^{*}\) satisfont
\[
 \sum_bP_{ab}=I_9=\sum_aP_{ab},\qquad P_{ab}^2=P_{ab}.
\]
En effet, chaque translation \(b\mapsto a+b\) permute les neuf résidus. On obtient une réalisation par des projecteurs de rang un d'un carré latin, dans la terminologie des carrés latins quantiques \cite{mustovicary,delascuevas2023}. Cette construction part du feuillet additif déjà défini et explicite l'application vers sa représentation opératorielle.

L'évaluation multiplicative conserve un autre contenu : pour \(a=3\), les
vecteurs \(e_{3b}\) répètent les résidus \(0,3,6\), et
\[
 \sum_{b\in\ZZ/9\ZZ}|e_{3b}\rangle\langle e_{3b}|
 =3\bigl(|e_0\rangle\langle e_0|+|e_3\rangle\langle e_3|
                    +|e_6\rangle\langle e_6|\bigr).
\]
La multiplicité est enregistrée dans l'opérateur et caractérise le radical
multiplicatif. Le support toroïdal, la propriété latine et la
résolution opératorielle de l'identité sont donc des propriétés qui exigent
des définitions différentes. Les carrés magiques quantiques étudiés par
De las Cuevas, Drescher et Netzer exigent des entrées positives et des sommes de lignes
et de colonnes égales à l'identité ; leur théorie distingue en outre les
dilatations à entrées commutatives\cite{delascuevas2020}. Cette référence
fournit un langage précis pour comparer les réalisations d'APP.

\subsection{TRIT : orientation et régimes quadratiques}

\begin{definition}[Signature ternaire et relèvement local]
La signature TRIT utilise l'identification équilibrée
\[
 \FF_3\longrightarrow\TRIT,\qquad 0\mapsto0,\quad1\mapsto+1,\quad2\mapsto-1.
\]
Pour \(n\in\ZZ\), son relèvement conserve \(n=r+3c\), où \(r\in\TRIT\) et \(c\in\ZZ\). Dans une transition d'amplitude locale bornée, les trois états \((0,-1),(0,0),(0,+1)\) peuvent présenter le même résidu nul et des retenues distinctes.
\end{definition}

La signature ne remplace pas la trajectoire. Un zéro résiduel ne signifie pas l'absence d'information d'orientation, de phase ou de quotient. Il n'introduit pas non plus à lui seul une grandeur physique. Sa fonction est de distinguer algébriquement les classes de transport qui seront réalisées ci-après.

\subsubsection{Trois régimes quadratiques et une exponentielle commune}

Une fois la signature ternaire obtenue, sa réalisation quadratique est définie par
\begin{equation}
 \mathbb A_\tau=\RR[X]/(X^2+\tau),\qquad \tau\in\TRIT.
 \label{eq:algebras-trit}
\end{equation}
Si \(J_\tau\) est la classe de \(X\), alors \(J_\tau^2=-\tau I\). La structure distingue le type elliptique \(\tau=+1\), le type parabolique \(\tau=0\) et le type hyperbolique \(\tau=-1\). La réalisation réelle ne génère pas la signature : elle exprime les types que la réduction orientée a déjà produits.

\begin{proposition}[Exponentiation des trois types]
Dans chaque algèbre \eqref{eq:algebras-trit},
\begin{equation}
 \exp(tJ_\tau)=C_\tau(t)I+S_\tau(t)J_\tau,
\end{equation}
où
\[
 C_\tau(t)=\sum_{n\ge0}\frac{(-\tau)^nt^{2n}}{(2n)!},\qquad
 S_\tau(t)=\sum_{n\ge0}\frac{(-\tau)^nt^{2n+1}}{(2n+1)!}.
\]
Les trois réalisations sont
\[
 \begin{array}{c|c|c}
 \tau&J_\tau^2&\exp(tJ_\tau)\\\hline
 +1&-I&\cos t\,I+\sin t\,J_\tau\\
 0&0&I+tJ_\tau\\
 -1&I&\cosh t\,I+\sinh t\,J_\tau.
 \end{array}
\]
\end{proposition}
\begin{proof}
Séparer les termes pairs et impairs de la série exponentielle et substituer \(J_\tau^{2n}=(-\tau)^nI\) et \(J_\tau^{2n+1}=(-\tau)^nJ_\tau\) donne les formules. Pour \(\tau=0\), tous les termes de degré au moins deux s'annulent. Dans les deux autres cas, on retrouve les séries trigonométriques ou hyperboliques.
\end{proof}

L'exponentielle est commune aux trois régimes : elle n'est pas attribuée exclusivement à la branche parabolique. Les invariants
\[
 C_\tau(t)^2+\tau S_\tau(t)^2=1
\]
se déduisent de \(\exp(tJ_\tau)\exp(-tJ_\tau)=I\). Le même principe algébrique conserve l'unité à travers des réalisations qui diffèrent par leur type quadratique.

Dans le régime hyperbolique, \(P_\pm=(I\pm J_{-1})/2\) sont des projecteurs complémentaires et
\[
 \exp(tJ_{-1})=e^tP_++e^{-t}P_-.
\]
La croissance et la contraction apparaissent conjointement dans deux directions propres. Dans le régime elliptique, le paramètre décrit une rotation ; dans le régime parabolique, une translation nilpotente. L'interprétation temporelle adoptée par HMT associe \(+1\) au retour à mémoire, \(0\) au seuil et \(-1\) à l'ouverture bidirectionnelle. Cette interprétation exige un ordre de paramétrage ; elle n'ajoute pas un quatrième type algébrique.

\input{sections/trit_desarrollo.tex}

\input{figures/regimenes_trit.tex}

\newpage
\subsection{TPK : noyau topologique de phase}

La dynamique utilise le sélecteur de phase
\[
 M_{\rm ph}(r)=
 \begin{cases}
 +1,&r\in\{1,4,7\},\\
 -1,&r\in\{2,5,8\},\\
 0,&r\in\{3,6,9\}.
 \end{cases}
\]
Sa valeur décide quelle évaluation opère. Le transport cardinal déplace le curseur correspondant ; le registre conserve la valeur précédente, le nouveau quotient, l'orientation et la transition. La sélection et le transport sont donc des fonctions distinctes : un signe ternaire n'est pas un déplacement du tore.

Nous désignerons par \(\widehat X\) l'espace des états arithmétiques avec mémoire de trajectoire. Un élément contient les positions et les directions des deux curseurs, le mode actif, la signature TRIT, la phase, les résidus et les quotients des deux évaluations, les retenues, le préfixe construit et les données de frontière. Lorsqu'un cylindre de raffinement est introduit, il fait partie de l'état. La notation évite de remplacer l'objet par une version réduite qui ne conserve que le symbole visible.

\begin{definition}[Mise à jour du TPK]
Soient \(\mathsf{Sel}_t\) la sélection du mode au moyen de \(M_{\rm ph}\), \(\mathsf{Tra}_t\) le transport cardinal relevé et \(\mathsf{Mem}_t\) la mise à jour des registres. La transition est
\begin{equation}
 \mathcal U_t=\mathsf{Mem}_t\circ\mathsf{Tra}_t\circ\mathsf{Sel}_t.
 \label{eq:actualizacion}
\end{equation}
Chaque application conserve les champs qu'elle ne modifie pas et met à jour les autres au moyen de la division euclidienne et de la concaténation de la trajectoire.
\end{definition}

La notation \eqref{eq:actualizacion} exprime l'ordre de composition et ne remplace pas les règles effectives. La section suivante spécifie complètement sa réalisation finie à deux curseurs. Ensuite, le prolongement ne s'obtiendra pas en effaçant sa mémoire et en répétant la première émission, mais en transportant les préfixes et leurs frontières compatibles. Ainsi, APP fournit les évaluations, TRIT distingue leurs régimes et TPK compose leur évolution.

\input{sections/tpk_desarrollo_integrado.tex}

\input{sections/memoria_resolvente.tex}

% BEGIN NUCLEO_ADITIVO_20260921_CENSO

\section{Noyau topologique de phase et génération du catalogue de germes}
\label{act21:tpk:capitulo}

Les deux évaluations de l'APP acquièrent une efficacité générative lorsque l'on précise
quelle évaluation agit, en quelle position elle s'effectue, comment cette position se déplace
et quelle information chaque transition conserve. Le \emph{Topological Phase Kernel},
abrégé en TPK et dénommé en espagnol \emph{núcleo topológico de fase}, réunit
ces opérations dans une dynamique sur des états enrichis. Sa définition
comprend la sélection, le transport, la mise à jour arithmétique et la conservation du
registre. La construction finie développée dans ce chapitre fournit
les germes sur lesquels agiront les relèvements régionaux et le prolongement
coinductif.

Le résultat principal de cette étape est un catalogue généré et classé :
les conditions initiales des deux curseurs produisent des émissions entières ;
leur lecture ternaire détermine des mots orientés ; l'action diédrale organise
ces mots en orbites. Chaque étape conserve ses préimages et les relations
qui permettent de revenir à la description précédente. La chaîne de cardinaux
\[
 104\,976\longrightarrow468\longrightarrow243\subset729,
 \qquad 243\longrightarrow43,
\]
désigne sous une forme abrégée des objets différents. Chaque objet est ensuite construit et chaque dénombrement
est justifié, avant que l'un d'eux soit utilisé comme donnée de
la sélection régionale.

\subsection{Sélection de phase, transport et état enrichi}
\label{act21:tpk:operaciones}

\subsubsection{Sélecteur tritique de phase opérationnelle}

La phase nonadique prend ses valeurs en \(D_9=\{1,\ldots,9\}\). Le sélecteur  
tritique est la fonction
\begin{equation}
 M_{\mathrm{ph}}:D_9\longrightarrow\{-1,0,+1\},\qquad
 M_{\mathrm{ph}}(g)=
 \begin{cases}
 +1,&g\in\{1,4,7\},\\
 -1,&g\in\{2,5,8\},\\
 0,&g\in\{3,6,9\}.
 \end{cases}
 \label{act21:tpk:selector}
\end{equation}
Dans la réalisation émettrice, la valeur \(+1\) active l'évaluation additive,
\(-1\) active l'évaluation multiplicative et \(0\) enregistre l'événement
neutre. L'orientation cardinale du curseur est une autre coordonnée de l'état.
Cette séparation permet de fixer simultanément quelle opération est lue et dans quelle
direction est transporté son point d'évaluation.

Pour les transitions numérotées par \(t\geq1\), la phase précédant la lecture est
\begin{equation}
 g_t=1+[(t-1)]_9,\qquad \epsilon_t=M_{\mathrm{ph}}(g_t),
 \label{act21:tpk:fase}
\end{equation}
où \([n]_b\) désigne le représentant de \(n\) modulo \(b\) dans
\(\{0,\ldots,b-1\}\). Une fenêtre de neuf transitions contient la suite
\((+1,-1,0,+1,-1,0,+1,-1,0)\).

\subsubsection{Curseurs orientés et transport cardinal}

Des indices \(I_9=\{0,\ldots,8\}\) sont utilisés pour les positions et des étiquettes positives \((a,b)=(i+1,j+1)\) pour les évaluations APP. Un curseur fini est
\[
 c=(i,j;d)\in\mathcal S:=I_9^2\times\mathcal D,
 \qquad \mathcal D=\{N,E,S,O\}.
\]
Les quatre directions agissent par
\[
 v_N=(-1,0),\quad v_E=(0,1),\quad
 v_S=(1,0),\quad v_O=(0,-1).
\]
L'application cardinale correspondant à \(d\) est
\[
 T_d(i,j)=\bigl([i+(v_d)_1]_9,[j+(v_d)_2]_9\bigr).
\]
Ainsi, \(M_{\mathrm{ph}}\) et \(T_d\) ont des domaines, des codomaines et des actions différents. L'activation tritique sélectionne l'un des deux curseurs ; le transport cardinal modifie sa position dans la direction que ce curseur conserve déjà.

Le relèvement entier du curseur est
\(\widetilde c=(x,y;d)\in\mathbb Z^2\times\mathcal D\), avec position finie
\(([x]_9,[y]_9)\). Ses nombres de tours sont récupérés en
\begin{equation}
 x=9\lfloor x/9\rfloor+[x]_9,\qquad
 y=9\lfloor y/9\rfloor+[y]_9.
 \label{act21:tpk:vueltas}
\end{equation}
Ces deux divisions concernent le transport spatial. La division d'une
évaluation additive ou multiplicative constitue un registre arithmétique
distinct, qui sera conservé conjointement avec elles.

\subsubsection{Composition de la transition}

L'état de la réalisation finie comprend les deux curseurs relevés,
leurs directions, la phase, les accumulateurs de la fenêtre active, la liste
ordonnée des blocs déjà émis et le registre des transitions. Chaque registre
local contient les positions antérieures, les évaluations entières des deux
feuillets, leurs restes et quotients, le régime actif et les orientations.
Les champs de préfixe et de frontière qui apparaîtront dans un prolongement seront
incorporés au même état avec leurs règles de compatibilité.

La transition s'ordonne comme
\begin{equation}
 \mathcal U_t=\mathsf{Mem}_t\circ\mathsf{Tra}_t\circ\mathsf{Sel}_t.
 \label{act21:tpk:composicion}
\end{equation}
La sélection détermine le régime et conserve l'échantillon antérieur à tout déplacement. Le transport applique le mouvement actif et l'inversion cardinale prescrite par le calendrier. L'actualisation incorpore cet échantillon aux accumulateurs et au registre ; à la fin d'une fenêtre, elle émet son bloc. La composition utilise donc un échantillon antérieur au transport bien que son archivage définitif s'effectue à la fin de la transition.

\subsection{Conditions initiales et énumération réversible}
\label{act21:tpk:semillas}

Les feuillets additif et multiplicatif disposent de curseurs indépendants. Le
domaine des conditions initiales est le produit ordonné
\[
 \Omega_{\mathrm{seed}}=\mathcal S_+\times\mathcal S_\times,
 \qquad |\mathcal S|=9^2\cdot4=324.
\]
Par conséquent,
\begin{equation}
 |\Omega_{\mathrm{seed}}|=324^2=104\,976.
 \label{act21:tpk:censo-inicial}
\end{equation}
Le nombre \(81\) compte des positions ; \(324\) compte des positions orientées d'un curseur ; \(104\,976\) compte des paires ordonnées de curseurs. Le choix d'une position en APP conserve uniquement une partie d'une condition initiale de l'émetteur.

On fixe la numérotation
\[
 \operatorname{ind}(N)=0,\quad\operatorname{ind}(E)=1,\quad
 \operatorname{ind}(S)=2,\quad\operatorname{ind}(O)=3,
\]
et l'on définit
\begin{equation}
 \nu(i,j;d)=4(9i+j)+\operatorname{ind}(d).
 \label{act21:tpk:indice-semilla}
\end{equation}

\begin{proposition}[Énumération exacte des conditions initiales]
\label{act21:tpk:enumeracion}
L'application \(\nu\) est une bijection de \(\mathcal S\) sur
\(\{0,\ldots,323\}\). Le couple d'indices se code de manière réversible
par \(324\nu_++\nu_\times\in\{0,\ldots,104\,975\}\).
\end{proposition}
\begin{proof}
La division par quatre restaure
\[
 \operatorname{ind}(d)=[\nu]_4,\qquad
 j=[\lfloor\nu/4\rfloor]_9,\qquad i=\lfloor\nu/36\rfloor.
\]
Les bornes du domaine garantissent que ces trois quantités appartiennent à
leurs ensembles d'indices et reconstruisent \(\nu\) au moyen de
\eqref{act21:tpk:indice-semilla}. Pour le couple, la division euclidienne par
\(324\) restitue le quotient \(\nu_+\) et le reste \(\nu_\times\).
\end{proof}

L'énumération parcourt le domaine complet avant tout regroupement
régional. Les indices servent à localiser les préimages d'une
émission, en conservant explicitement la condition qui a produit chaque trajectoire.

\subsection{Évaluation des feuillets et registre arithmétique}
\label{act21:tpk:lecturas}

En une position positive \((a,b)\), les évaluations entières sont
\(A_+(a,b)=a+b\) et \(A_\times(a,b)=ab\). Pour \(n>0\), on conserve
\begin{equation}
 \rho_9(n)=1+[n-1]_9,\qquad
 q_9(n)=\lfloor(n-1)/9\rfloor,\qquad
 n=\rho_9(n)+9q_9(n).
 \label{act21:tpk:division-evaluacion}
\end{equation}
L'activation additive incorpore \(\rho_9(A_+)\) à l'accumulateur correspondant. L'activation multiplicative utilise un observable du résidu \(\rho_9(A_\times)\), défini ci-dessous.

\subsubsection{Exposant discret dans les unités modulo neuf}

Les puissances de \(2\) modulo \(9\) parcourent
\[
 2^0,2^1,\ldots,2^5\equiv1,2,4,8,7,5\pmod9,
 \qquad 2^6\equiv1\pmod9.
\]
Le logarithme discret dans ce groupe de six unités se représente par
\begin{equation}
 \ell_9(1)=0,\quad\ell_9(2)=1,\quad\ell_9(4)=2,
 \quad\ell_9(8)=3,\quad\ell_9(7)=4,\quad\ell_9(5)=5.
 \label{act21:tpk:logaritmo}
\end{equation}
La loi émettrice étend l'observable par le biais de \(\ell_9(3)=\ell_9(6)=\ell_9(9)=0\). Pour conserver son origine, l'échantillon enregistré inclut le résidu original et l'indicateur d'appartenance au groupe d'unités. Ainsi, une valeur observée nulle peut provenir de l'unité \(1\) ou de l'un des trois résidus non inversibles, et son origine demeure récupérable.

\subsubsection{Ordre de lecture et mouvement}

Les évaluations s'effectuent aux positions antérieures à la transition.
Si \(\epsilon_t=+1\), on accumule la lecture additive et on déplace son curseur.
Si \(\epsilon_t=-1\), on accumule l'exposant discret multiplicatif et on
déplace le second curseur. Si \(\epsilon_t=0\), on incrémente le compteur
neutre et on conserve les deux positions.

À la fin des étapes \(27\) et \(54\), les deux directions sont remplacées
par leurs opposées. Les fenêtres s'enchaînent en conservant leurs positions et
orientations. Au début de chaque fenêtre, seuls ses trois accumulateurs locaux sont remis à zéro ; le registre précédent reste disponible.

\begin{proposition}[Inversion du transport relevé]
\label{act21:tpk:inversion}
Soient \(a,b\in\{0,1\}\) les indicateurs d'activation et d'inversion d'un
curseur. Si \(\operatorname{opp}\) échange \(N\leftrightarrow S\) et
\(E\leftrightarrow O\), la transformation
\[
 F_{a,b}(z,d)=\bigl(z+a v_d,\operatorname{opp}^{\,b}(d)\bigr)
\]
admet pour inverse
\[
 F_{a,b}^{-1}(z',d')=
 \bigl(z'-a v_{\operatorname{opp}^{\,b}(d')},
       \operatorname{opp}^{\,b}(d')\bigr).
\]
La réduction spatiale modulo neuf entrelace cette transformation avec le
transport fini.
\end{proposition}
\begin{proof}
Appliquer deux fois l'opposition restaure la direction initiale. Une fois
cette direction récupérée, la soustraction de \(a v_d\) annule le déplacement.
Pour la projection spatiale, il suffit de l'égalité
\([z+a v_d]_9=[[z]_9+a v_d]_9\), composante par composante. L'inversion de
direction conserve cette égalité.
\end{proof}

La phase détermine les indicateurs utilisés par l'inverse. Le registre
spatial et l'arithmétique peuvent ainsi être conservés conjointement, même en
traversant le bord du représentant \(I_9^2\).

\subsection{Six fenêtres et émission décimale par blocs}
\label{act21:tpk:emision}

\subsubsection{Accumulation dans une fenêtre nonadique}

La fenêtre \(m\), avec \(1\leq m\leq6\), comprend les étapes
\(9m-8,\ldots,9m\). Soient \(S_m\) la somme de ses résidus positifs
additifs actifs, \(E_m\) la somme de ses exposants discrets actifs et
\(Z_m\) son nombre d'événements neutres.

\begin{lemma}[Multiplicités d'activation]
\label{act21:tpk:tres-eventos}
Chaque fenêtre contient trois activations additives, trois multiplicatives et
trois neutres. En particulier, \(Z_m=3\).
\end{lemma}
\begin{proof}
Les phases d'une fenêtre parcourent exactement \(1,\ldots,9\). Chacun
des trois ensembles de \eqref{act21:tpk:selector} a pour cardinal trois.
\end{proof}

Les six fenêtres forment \(54\) transitions et produisent six blocs.  
Cette longueur appartient à l'émetteur initial défini ici. La continuation  
en profondeur agit ensuite sur les mots et leurs registres compatibles.

\subsubsection{Règle de codification décimale}

La loi émettrice spécifie les chiffres
\begin{align}
 d_{m,1}&=[S_m]_{10},\notag\\
 d_{m,2}&=[S_m+E_m]_{10},\notag\\
 d_{m,3}&=[3S_m+5E_m+7Z_m]_{10},\notag\\
 U_m&=100d_{m,1}+10d_{m,2}+d_{m,3}.
 \label{act21:tpk:bloque-decimal}
\end{align}
Les poids \(100,10,1\) sont les positions des centaines, des dizaines et des unités.
L'indice \(10\) indique le reste décimal, et chaque bloc satisfait
\(0\leq U_m\leq999\). Le mot émis est la suite ordonnée
\(\mathbf U=(U_1,\ldots,U_6)\), de largeur fixe de trois chiffres par bloc.

Les coefficients \(3,5,7\) constituent les poids explicites de cette loi
émettrice. Une fois la règle fixée, chaque chiffre se calcule à partir des
accumulateurs, et les résultats de ce chapitre se déduisent de cette règle.
Le nombre \(1\) qui apparaît sous sa forme réduite a une origine
additionnelle concrète : \(7Z_m=21\equiv1\pmod{10}\).

Si \(s_m=[S_m]_{10}\) et \(e_m=[E_m]_{10}\), on obtient
\begin{equation}
 G(s_m,e_m)=100s_m+10[s_m+e_m]_{10}
                   +[3s_m+5e_m+1]_{10}=U_m.
 \label{act21:tpk:acoplamiento-firmas}
\end{equation}
La lettre \(e_m\) désigne ici une composante de la signature d'exposants discrets. Sa définition appartient à l'émetteur fini.

\begin{proposition}[Récupération des deux signatures]
\label{act21:tpk:recuperacion-firmas}
L'application \(G:\{0,\ldots,9\}^2\to\{0,\ldots,999\}\) est injective.
Pour un bloc \(U=100a+10b+c\) de son image, l'inverse est
\[
 s=a,\qquad e=[b-a]_{10},\qquad
 c=[3s+5e+1]_{10}.
\]
En conséquence, le mot \(\mathbf U\) détermine les deux signatures de six composantes \(\mathbf s\) et \(\mathbf e\).
\end{proposition}
\begin{proof}
Les deux derniers termes de \eqref{act21:tpk:acoplamiento-firmas} appartiennent à
\(\{0,\ldots,99\}\), de sorte que la centaine récupère \(s\). La dizaine
est \([s+e]_{10}\) ; en soustrayant \(s\) et en réduisant modulo dix, on récupère
\(e\). L'unité vérifie la troisième congruence. La reconstruction s'applique
indépendamment aux six positions du mot.
\end{proof}

Les totaux entiers \(S_m,E_m\) se retrouvent à partir des lectures enregistrées ;
les centaines et les dizaines retrouvent leurs restes. Par exemple, un total
\(S_m=15\) produit \(s_m=5\). Cette différence identifie exactement
quelle information chaque niveau de codage conserve.

\subsubsection{Codification positionnelle du mot de six blocs}

Les six blocs admettent une codification entière à largeur fixe,
\begin{equation}
 \mathscr N(\mathbf U)=\sum_{m=1}^{6}U_m1000^{6-m},
 \qquad0\leq\mathscr N(\mathbf U)<1000^6.
 \label{act21:tpk:entero-base1000}
\end{equation}
Le premier bloc occupe la position de plus haut poids et le sixième, celle des unités de base mille. Un bloc \(058\) occupe une position complète avec une valeur entière \(58\) ; sa largeur de trois chiffres fait partie de la présentation décimale de cette position.

\begin{proposition}[Récupération positionnelle des émissions]
\label{act21:tpk:inversa-base1000}
La codification \eqref{act21:tpk:entero-base1000} détermine le mot complet
par le biais de
\[
 U_m=\left[\left\lfloor
 \frac{\mathscr N(\mathbf U)}{1000^{6-m}}\right\rfloor\right]_{1000},
 \qquad1\leq m\leq6.
\]
\end{proposition}
\begin{proof}
En divisant par \(1000^{6-m}\), les blocs précédents produisent des multiples entiers de \(1000\), le bloc \(m\) produit \(U_m\), et la contribution des suivants appartient à \([0,1)\). La partie entière élimine cette dernière contribution et le résidu modulo \(1000\) élimine les premières.
\end{proof}

Cette codification conserve les six émissions. Son domaine possède des positions ordonnées de base mille, tandis que le domaine APP possède des positions spatiales modulo neuf. La composition maintient les deux indices : \(m\) identifie une fenêtre d'émission et \((i,j)\) identifie la cellule visitée lors d'une transition de cette fenêtre.

\subsection{Calcul complet d'une émission qui commence par 501}
\label{act21:tpk:ejemplo501}

On fixe les conditions initiales
\[
 c_{+,0}=(0,4;N),\qquad c_{\times,0}=(2,3;N),
 \qquad \nu_+=16,\quad\nu_\times=84.
\]
Les positions positives initiales sont \((1,5)\) et \((3,4)\). L'indice
chronologique précédent est \(t=0\), la première phase lue est \(g_1=1\) et le
registre est vide. Ces coordonnées sont les entrées
du calcul ; les blocs s'obtiennent en exécutant la récurrence.

\begin{table}[htbp]
\caption{Première fenêtre : lectures précédentes et accumulation active.}
\label{act21:tpk:tabla-501}
\small
\begin{tabular}{@{}rcccrrr@{}}
\toprule
\(t\)&\(\epsilon_t\)&position \(+\)&position \(\times\)&
\(\Delta S\)&\(\Delta E\)&\((S,E,Z)\)\\
\midrule
1&\(+1\)&\((1,5)\)&\((3,4)\)&6&0&\((6,0,0)\)\\
2&\(-1\)&\((9,5)\)&\((3,4)\)&0&0&\((6,0,0)\)\\
3&0&\((9,5)\)&\((2,4)\)&0&0&\((6,0,1)\)\\
4&\(+1\)&\((9,5)\)&\((2,4)\)&5&0&\((11,0,1)\)\\
5&\(-1\)&\((8,5)\)&\((2,4)\)&0&3&\((11,3,1)\)\\
6&0&\((8,5)\)&\((1,4)\)&0&0&\((11,3,2)\)\\
7&\(+1\)&\((8,5)\)&\((1,4)\)&4&0&\((15,3,2)\)\\
8&\(-1\)&\((7,5)\)&\((1,4)\)&0&2&\((15,5,2)\)\\
9&0&\((7,5)\)&\((9,4)\)&0&0&\((15,5,3)\)\\
\bottomrule
\end{tabular}
\notatabla{Les positions sont des étiquettes positives des cellules antérieures au mouvement. Les événements neutres augmentent uniquement \(Z\).}
\end{table}

Les trois évaluations additives actives sont
\[
 1+5=6,\qquad9+5=14=5+9,\qquad8+5=13=4+9,
\]
de telle sorte que \(S_1=6+5+4=15\). Les trois évaluations multiplicatives actives
sont \(3\cdot4=12\), \(2\cdot4=8\) et \(1\cdot4=4\). Leurs résidus
positifs sont \(3,8,4\), avec des observables \(0,3,2\) ; par conséquent,
\(E_1=5\). Enfin \(Z_1=3\). La règle décimale donne
\[
 d_{1,1}=[15]_{10}=5,\quad
 d_{1,2}=[15+5]_{10}=0,\quad
 d_{1,3}=[45+25+21]_{10}=1,
\]
et on obtient \(U_1=501\).

La seconde lecture additive précédente a pour quotient arithmétique \(1\),
tandis que son curseur relevé se trouve en \((-1,4;N)\), avec un nombre de tours
verticaux \(-1\). Ce sont deux données différentes du même registre. Toutes deux
interviennent dans la récupération intégrale de l'évaluation et de sa position.

\begin{table}[htbp]
\caption{Les six fenêtres de la condition initiale \((16,84)\).}
\label{act21:tpk:tabla-seis501}
\begin{tabular}{@{}rrrrrrrr@{}}
\toprule
\(m\)&pas&\(S_m\)&\(E_m\)&\(Z_m\)&\(s_m\)&\(U_m\)&\([-U_m]_3\)\\
\midrule
1&1--9&15&5&3&5&501&0\\
2&10--18&6&5&3&6&614&1\\
3&19--27&24&5&3&4&498&0\\
4&28--36&21&5&3&1&169&2\\
5&37--45&12&5&3&2&272&1\\
6&46--54&12&5&3&2&272&1\\
\bottomrule
\end{tabular}
\end{table}

Le mot résultant et ses signatures sont
\begin{align*}
 \mathbf U&=(501,614,498,169,272,272),\\
 \mathbf s&=(5,6,4,1,2,2),&
 \mathbf e&=(5,5,5,5,5,5).
\end{align*}
L'inversion cardinale après l'étape \(27\) modifie les parcours des
trois fenêtres finales. Le bloc répété \(272\) correspond à deux
fenêtres successives possédant leurs propres positions et leur propre registre chronologique.

Comme deuxième cas, la condition initiale minimale \((\nu_+,\nu_\times)=(0,0)\)
produit
\[
 \mathbf U=(252,109,252,903,850,850),\quad
 \mathbf s=(2,1,2,9,8,8),\quad
 \mathbf e=(3,9,3,1,7,7).
\]
Dans ce cas, les deux curseurs retournent à leur position et orientation initiales
après \(54\) transitions. La phase nonadique effectue six tours, le
compteur chronologique avance de \(54\) unités et le registre contient
\(54\) échantillons. La lecture positionnelle du retour et
l'histoire enregistrée décrivent des propriétés distinctes de l'exécution.

\subsection{Factorisation du recensement fini}
\label{act21:tpk:censo}

Les signatures \(f_+(\nu)=\mathbf s\) et \(f_\times(\nu)=\mathbf e\) se
calculent en exécutant le curseur correspondant pendant les six fenêtres.
L'indépendance des positions et des directions entre les deux curseurs donne
\begin{equation}
 T_{\mathrm{em}}(\nu_+,\nu_\times)
   =G^{(6)}\bigl(f_+(\nu_+),f_\times(\nu_\times)\bigr),
 \label{act21:tpk:factorizacion}
\end{equation}
où \(G^{(6)}\) applique \eqref{act21:tpk:acoplamiento-firmas} à chaque position.
La factorisation permet d'énumérer d'abord \(324+324\) trajectoires et
de combiner ensuite les signatures distinctes, en conservant la multiplicité de
leurs préimages.

\par\noindent\begin{minipage}{\linewidth}
\subsubsection{Images des deux signatures}

Les tableaux suivants comprennent toutes les signatures. Leur multiplicité est le
nombre de conditions initiales qui les produisent ; l'indice minimal localise
l'une de ces conditions. La préimage complète est définie par
\[
 f_\sigma^{-1}(a)=\{\nu\in\{0,\ldots,323\}:f_\sigma(\nu)=a\}.
\]
La récurrence des sections précédentes détermine chaque élément de cet
ensemble par un calcul entier fini.

\begin{table}[H]
\centering\fontsize{9.5}{11.4}\selectfont
\setlength{\abovecaptionskip}{8pt}\setlength{\belowcaptionskip}{6pt}
\renewcommand{\arraystretch}{1.12}
\caption{Les dix-huit signatures additives.}
\label{act21:tpk:firmas-aditivas}
\begin{tabular}{@{}lrr@{\qquad}lrr@{}}
\toprule
signature&multiplicité&\(\nu_{\min}\)&signature&multiplicité&\(\nu_{\min}\)\\
\midrule
\texttt{122564}&18&17&\texttt{564122}&18&16\\
\texttt{122898}&18&24&\texttt{645212}&18&4\\
\texttt{212645}&18&5&\texttt{654889}&18&33\\
\texttt{212988}&18&0&\texttt{889221}&18&13\\
\texttt{221456}&18&29&\texttt{889654}&18&32\\
\texttt{221889}&18&12&\texttt{898122}&18&25\\
\texttt{456221}&18&28&\texttt{898465}&18&20\\
\texttt{465898}&18&21&\texttt{988212}&18&1\\
\texttt{546988}&18&9&\texttt{988546}&18&8\\
\bottomrule
\end{tabular}

\par\vspace{12pt}

\caption{Les vingt-six signatures multiplicatives.}
\label{act21:tpk:firmas-multiplicativas}
\begin{tabular}{@{}lrr@{\qquad}lrr@{}}
\toprule
signature&multiplicité&\(\nu_{\min}\)&signature&multiplicité&\(\nu_{\min}\)\\
\midrule
\texttt{000000}&108&8&\texttt{555555}&24&84\\
\texttt{177393}&8&1&\texttt{555717}&8&14\\
\texttt{177555}&8&54&\texttt{555771}&8&18\\
\texttt{177771}&8&11&\texttt{555933}&8&24\\
\texttt{339555}&8&6&\texttt{717555}&8&12\\
\texttt{339771}&8&15&\texttt{717717}&8&21\\
\texttt{339933}&8&59&\texttt{717933}&8&19\\
\texttt{393177}&8&0&\texttt{771177}&8&9\\
\texttt{393393}&8&45&\texttt{771339}&8&13\\
\texttt{393555}&8&40&\texttt{771555}&8&16\\
\texttt{555177}&8&52&\texttt{933339}&8&57\\
\texttt{555339}&8&4&\texttt{933555}&8&26\\
\texttt{555393}&8&41&\texttt{933717}&8&17\\
\bottomrule
\end{tabular}
\end{table}
\end{minipage}\par


\begin{proposition}[Image et multiplicités de l'émetteur]
\label{act21:tpk:468}
L'image \(\mathcal U_6=\operatorname{im}T_{\mathrm{em}}\) contient
\(468\) mots distincts. Ses préimages se répartissent en \(432\)
fibres de cardinal \(144\), \(18\) de cardinal \(432\) et \(18\) de
cardinal \(1944\).
\end{proposition}
\begin{proof}
L'évaluation exhaustive de l'algorithme de la section ~\ref{act21:tpk:algoritmo} sur les indices \(0,\ldots,323\) donne les deux tableaux ci-dessus. La première partition a \(18\) classes de cardinal \(18\), et la seconde \(24\) classes de cardinal \(8\), une de cardinal \(24\) et une de cardinal \(108\). Les sommes \(18\cdot18=324\) et \(24\cdot8+24+108=324\) vérifient la masse totale des deux partitions. La proposition ~\ref{act21:tpk:recuperacion-firmas} rend injectif le couplage de ses images ; par conséquent, il y a \(18\cdot26=468\) émissions. Par \eqref{act21:tpk:factorizacion}, la préimage d'un couple de signatures est le produit de ses deux préimages. On obtient \(18\cdot24=432\) produits de cardinal \(18\cdot8=144\), \(18\) de cardinal \(18\cdot24=432\) et \(18\) de cardinal \(18\cdot108=1944\). Enfin,
\[
 432\cdot144+18\cdot432+18\cdot1944=104\,976.
\]
La partie énumérative est une vérification finie exacte de la récurrence ;
la factorisation explique pourquoi ce décompte couvre toutes les paires.
\end{proof}

L'émission de l'exemple \ref{act21:tpk:ejemplo501} a une fibre de cardinal
\(18\cdot24=432\). L'exemple d'indices \((0,0)\) a une fibre de
cardinal \(18\cdot8=144\). Dans les deux cas, le mot émis récupère
les signatures, tandis que l'identification d'une condition initiale particulière
utilise son indice ou son registre de trajectoire.

\subsection{Lecture ternaire orientée et conservation des fibres}
\label{act21:tpk:reduccion-ternaria}

La réduction orientée agit séparément sur chaque bloc :
\begin{equation}
 \Psi:\mathcal U_6\longrightarrow\mathbb F_3^6,\qquad
 \Psi(U_1,\ldots,U_6)=([-U_1]_3,\ldots,[-U_6]_3).
 \label{act21:tpk:psi}
\end{equation}
Son signe fixe une orientation du catalogue. L'identification équilibrée
\(0\mapsto0\), \(1\mapsto+1\), \(2\mapsto-1\) s'applique ensuite,
composante par composante.

Le module neuf organise les évaluations et positions d'APP ; les restes modulo dix forment chaque chiffre ; la base mille regroupe trois positions décimales ; le module trois détermine la signature ternaire de chaque bloc. Ces quatre usages ont des objets et des fonctions définis. En particulier, \(\Psi\) lit six blocs et produit six trits ; une conversion positionnelle de l'entier formé par dix-huit chiffres serait une autre opération.

Pour les deux conditions calculées, on obtient
\begin{align*}
 \Psi(501,614,498,169,272,272)&=(0,1,0,2,1,1),\\
 \Psi(252,109,252,903,850,850)&=(0,2,0,0,2,2).
\end{align*}
L'émetteur décimal est défini avant cette réduction. La fonction de
\(\Psi\) est de produire l'objet ternaire sur lequel agissent la symétrie
orbitale et les sélecteurs régionaux. Cette fonction explique sa position dans la
généalogie : elle relie une émission déjà calculée à une classification structurale.

L'image \(\mathcal W_6:=\Psi(\mathcal U_6)\) contient exactement
\(243\) mots au sein de l'ensemble ambiant de \(3^6=729\) mots.
L'énumération de l'algorithme final donne l'histogramme
\begin{equation}
\begin{array}{c|rrrrrr}
 |\Psi^{-1}(w)|&1&2&3&4&5&6\\\hline
 \text{nombre de mots }w&132&66&18&3&6&18.
\end{array}
\label{act21:tpk:fibras-ternarias}
\end{equation}
Les sommes des deux lectures sont
\[
132+66+18+3+6+18=243,
\qquad132+2\cdot66+3\cdot18+4\cdot3+5\cdot6+6\cdot18=468.
\]
La première compte les mots ternaires ; la seconde récupère le nombre d'émissions réparties entre ses fibres.

\begin{example}[Deux émissions avec la même signature ternaire]
\label{act21:tpk:colision}
Les émissions
\[
 (498,501,614,252,252,109),\qquad
 (870,810,923,252,252,109)
\]
sont distinctes et toutes deux ont pour image \((0,0,1,0,0,2)\) par \(\Psi\).
Chacune a \(144\) conditions initiales. Le mot ternaire conserve
une classe de lecture ; les émissions et leurs préimages distinguent les données
qui convergent vers cette classe.
\end{example}

La conservation des fibres s'exprime concrètement par
\[
 \mathcal F_w=\{\mathbf U\in\mathcal U_6:\Psi(\mathbf U)=w\},
 \qquad
 \Omega_{\mathbf U}=T_{\mathrm{em}}^{-1}(\mathbf U).
\]
Ainsi, la préimage complète de \(w\) dans le domaine initial est l'union disjointe des ensembles \(\Omega_{\mathbf U}\), avec \(\mathbf U\in\mathcal F_w\). Le registre enrichi peut conserver l'émission, ses signatures et la condition initiale ainsi que sa lecture ternaire. L'inverse exacte se réfère alors à la donnée enregistrée avec son type ; le mot de six trits pris isolément conserve la portée précise de \eqref{act21:tpk:psi}.

\subsection{Action diédrale et les quarante-trois orbites}
\label{act21:tpk:orbitas}

Sur six coordonnées se définissent les permutations
\begin{align}
 r(u_1,u_2,u_3,u_4,u_5,u_6)
   &=(u_3,u_1,u_2,u_5,u_6,u_4),\\
 s(u_1,u_2,u_3,u_4,u_5,u_6)
   &=(u_4,u_5,u_6,u_1,u_2,u_3).
 \label{act21:tpk:generadores-diedrales}
\end{align}
La première fait tourner les deux groupes de trois coordonnées dans des sens conjugués ;
la seconde échange les deux groupes. Leur composition satisfait
\(r^3=s^2=1\) et \(srs=r^{-1}\). Elles engendrent ainsi une action du groupe diédral
\(D_3\), à six éléments.

\begin{proposition}[Équivariance et stabilité du catalogue]
\label{act21:tpk:equivarianza}
La réduction \(\Psi\) commute avec l'action de \(D_3\). Les ensembles
\(\mathcal U_6\) et \(\mathcal W_6\) sont stables sous leurs générateurs.
\end{proposition}
\begin{proof}
L'égalité \(\Psi(g\mathbf U)=g\Psi(\mathbf U)\) se vérifie en chaque
coordonnée parce que \(g\) permute les positions et \(\Psi\) applique la même
réduction à chaque position. La stabilité du catalogue constitue une autre
vérification : en générant ses \(468\) éléments par
\eqref{act21:tpk:factorizacion}, on vérifie que \(r\mathbf U\) et
\(s\mathbf U\) appartiennent à nouveau à cet ensemble. L'algorithme de la
section~\ref{act21:tpk:algoritmo} effectue ces vérifications d'appartenance.
L'équivariance transmet ensuite la stabilité à son image ternaire.
\end{proof}

\begin{proposition}[Décomposition orbitale complète]
\label{act21:tpk:43}
L'ensemble \(\mathcal W_6\) se décompose en \(43\) orbites : \(38\)
de cardinal six et \(5\) de cardinal trois.
\end{proposition}
\begin{proof}
Pour chaque \(w\in\mathcal W_6\) on forme
\[
 \mathcal O(w)=\{w,rw,r^2w,sw,srw,sr^2w\}.
\]
On choisit comme représentant le mot lexicographiquement minimal de cet ensemble. Deux mots ont le même représentant exactement lorsqu'ils appartiennent à la même orbite. L'algorithme génère la liste des représentants du tableau~\ref{act21:tpk:representantes} ; la somme de leurs cardinaux vaut \(38\cdot6+5\cdot3=243\). Comme chaque mot du catalogue participe à cette opération et que le catalogue est stable, la liste couvre toute l'image.
\end{proof}

\begin{table}[!htbp]
\caption{Représentants minimaux et cardinaux des quarante-trois orbites.}
\label{act21:tpk:representantes}
\small
\begin{tabular}{@{}lr@{\qquad}lr@{\qquad}lr@{}}
\toprule
représentant&cardinal&représentant&cardinal&représentant&cardinal\\
\midrule
\texttt{001002}&6&\texttt{002112}&6&\texttt{012222}&6\\
\texttt{001011}&6&\texttt{002211}&6&\texttt{021022}&6\\
\texttt{001012}&6&\texttt{002220}&6&\texttt{021121}&6\\
\texttt{001021}&6&\texttt{011021}&6&\texttt{021202}&6\\
\texttt{001022}&6&\texttt{011112}&6&\texttt{022022}&3\\
\texttt{001100}&3&\texttt{011120}&6&\texttt{022112}&6\\
\texttt{001101}&6&\texttt{011201}&6&\texttt{022121}&6\\
\texttt{001110}&6&\texttt{011211}&6&\texttt{022202}&3\\
\texttt{001112}&6&\texttt{011220}&6&\texttt{022211}&6\\
\texttt{001202}&6&\texttt{011222}&6&\texttt{022222}&6\\
\texttt{001210}&6&\texttt{012112}&6&\texttt{112112}&3\\
\texttt{001211}&6&\texttt{012121}&6&\texttt{112211}&3\\
\texttt{001220}&6&\texttt{012202}&6&\texttt{112222}&6\\
\texttt{001222}&6&\texttt{012210}&6&&\\
\texttt{002012}&6&\texttt{012220}&6&&\\
\bottomrule
\end{tabular}
\end{table}

Les orbites héritent des invariants entiers de leurs fibres. Pour chaque mot,
\(n(w)=|\mathcal F_w|\) compte les émissions, tandis que
\[
 M(w)=\sum_{\mathbf U\in\mathcal F_w}|\Omega_{\mathbf U}|
\]
compte les conditions initiales. Le recensement des symboles
\(c(w)=(n_0,n_1,n_2)\) enregistre combien de fois chaque classe ternaire apparaît.
Ces données, conjointement avec l'orientation et le support des parcours,
fournissent les prédicats de la sélection régionale ultérieure. La
classification conserve ainsi plus d'information que la forme d'un mot
représentant.

\subsection{Algorithme autonome d'énumération et de contrôle}
\label{act21:tpk:algoritmo}

L'énumération utilise les entiers, les listes finies et les opérations déjà définies.  
Le procédé suivant précise son exécution ; \(\sigma\) prend les valeurs  
\(+\) et \(\times\).

\begin{enumerate}
\item Pour chaque \(\nu=0,\ldots,323\), récupérer \((i,j;d)\) par
l'inverse de \eqref{act21:tpk:indice-semilla}. Initialiser une liste vide de signatures.
\item Pour chaque fenêtre \(m=1,\ldots,6\), initialiser l'accumulateur \(A=0\).
Parcourir \(t=9m-8,\ldots,9m\), calculant \(\epsilon_t\) par
\eqref{act21:tpk:fase}.
\item Si \(\sigma=+\) et \(\epsilon_t=+1\), ajouter
\(\rho_9(i+j+2)\) à \(A\) et mettre à jour
\((i,j)\leftarrow T_d(i,j)\). Si \(\sigma=\times\) et
\(\epsilon_t=-1\), ajouter
\(\ell_9(\rho_9((i+1)(j+1)))\) et effectuer le même transport.
\item À la fin de \(t=27\) ou \(t=54\), remplacer \(d\) par son opposée.  
À la conclusion de la fenêtre, ajouter \([A]_{10}\) à la liste des signatures.  
Conserver la position et la direction pour la fenêtre suivante.
\item Grouper les \(324\) indices par égalité de leurs listes de six
composants. Ces deux groupements donnent \(f_+\), \(f_\times\) et toutes
leurs préimages.
\item Pour chaque paire de signatures distinctes, appliquer \(G^{(6)}\).
Attribuer à l'émission la multiplicité produit des deux préimages.
Vérifier la récupération des signatures dans chaque bloc et la somme des
multiplicités \(104\,976\).
\item Appliquer \(\Psi\) à chaque émission et regrouper par égalité de mots.
Compter les \(243\) valeurs et les multiplicités de
\eqref{act21:tpk:fibras-ternarias}.
\item Appliquer \(r,s\) à chaque émission et vérifier son appartenance au catalogue.
Pour chaque mot ternaire, former ses six images diédriques, éliminer les
répétitions et enregistrer son minimum lexicographique. Compter les \(43\)
orbites et les deux tailles du tableau~\ref{act21:tpk:representantes}.
\end{enumerate}

L'exécution se termine : le premier tronçon évalue \(648\) trajectoires de \(54\) transitions ; le deuxième combine \(468\) couples de signatures ; le troisième regroupe un ensemble de \(243\) mots. Une vérification directe supplémentaire peut exécuter les \(104\,976\) couples et vérifier la factorisation \eqref{act21:tpk:factorizacion} ; sa portée coïncide avec le même domaine fini.

\subsection{Conservation du registre et fonction des germes}
\label{act21:tpk:conclusion}

Soit \(\mathsf{Sample}(q_t)\) l'échantillon antérieur à la transition et \(\mathcal M_t\) le registre accumulé. La mise à jour de la mémoire est
\begin{equation}
 \mathcal M_{t+1}=\mathcal M_t\mathbin{\Vert}\mathsf{Sample}(q_t),
 \label{act21:tpk:archivo}
\end{equation}
où \(\Vert\) désigne la concaténation ordonnée. La récurrence donne
\[
 \mathcal M_n=\mathcal M_0\mathbin{\Vert}
 (\mathsf{Sample}(q_0),\ldots,\mathsf{Sample}(q_{n-1})).
\]
Les accumulateurs se reconstruisent en additionnant les contributions actives de
chaque segment de neuf enregistrements. Leurs quotients, les tours des
curseurs et les orientations restent associés aux mêmes transitions.
L'exécution détermine ce registre au moyen de la récurrence ; le fichier
préserve la provenance concrète des lectures.

À la fin de cette étape sont construits le domaine complet des conditions
initiales, l'émission décimale, ses deux signatures récupérables, l'image ternaire,
ses fibres et sa classification orbitale. Les recensements \(104\,976\), \(468\),
\(243\), \(729\) et \(43\) ont désormais un domaine et un mécanisme d'obtention.
Les sélections régionales recevront ces structures avec leurs invariants,
et les élévations agiront sur des mots orientés dont la provenance est
conservée. Cette continuité fait du catalogue fini un antécédent
constructif du développement coinductif.

% END NUCLEO_ADITIVO_20260921_CENSO
\section{Développement du noyau formel : conditions initiales, géométrie et raffinement}
\label{sec:extension}

Le support de quatre-vingt-une positions ne coïncide pas avec le domaine des conditions initiales de la dynamique. Une condition initiale doit également indiquer une orientation ; les feuillets additif et multiplicatif disposent en outre de curseurs indépendants. Cette distinction permet de reconstruire exhaustivement l’émission sans sélectionner au préalable une région numérique.

\subsection{Les états initiaux et la récurrence d’émission}

Pour l’énumération, nous utilisons les indices \(I_9=\{0,\ldots,8\}\), reliés aux étiquettes positives par \(i\mapsto i+1\). Soit
\[
 \mathcal S=I_9^2\times\{N,E,S,O\},\qquad
 \Omega_{\rm seed}=\mathcal S_+\times\mathcal S_\times.
\]
Chaque condition initiale \(\omega=(s_+,s_\times)\) détermine deux positions et deux directions. Le domaine comprend tous les couples ordonnés, de sorte que
\begin{equation}
 |\mathcal S|=324,\qquad |\Omega_{\rm seed}|=324^2=104\,976.
 \label{eq:semillas}
\end{equation}
L’annexe indépendante énumère toutes ces conditions ; le corps du texte présente la loi qui les transforme. Les lignes d’une table d’émissions sont des fibres de ce domaine de conditions initiales.

La réalisation finie utilise six fenêtres de neuf transitions. À chaque instant \(t\), la signature de phase est \(M_{\rm ph}(1+((t-1)\bmod9))\). Si elle vaut \(+1\), on lit le résidu positif \(\Sigma(i+1,j+1)=\rho_9(i+j+2)\) du premier curseur, puis on déplace ce curseur ; si elle vaut \(-1\), on lit le résidu positif \(\Pi(i+1,j+1)=\rho_9((i+1)(j+1))\) du second curseur, puis on déplace le second. La valeur zéro incrémente le compteur neutre. Au terme des transitions vingt-sept et cinquante-quatre, les orientations cardinales sont inversées. Les évaluations entières et leurs quotients sont conservés dans le registre, mais l’émetteur utilise les lectures résiduelles spécifiées ici.

Dans le secteur des unités, le lecteur multiplicatif utilise le logarithme discret de base deux dans \((\ZZ/9\ZZ)^\times\) :
\[
 \ell_9(1)=0,\quad\ell_9(2)=1,\quad\ell_9(4)=2,\quad
 \ell_9(8)=3,\quad\ell_9(7)=4,\quad\ell_9(5)=5.
\]
Pour l’émission finie, il est étendu avec la valeur zéro au secteur non unitaire. Cette extension est un observable et non une identification des éléments \(3,6,9\) à l’unité multiplicative. Le feuillet et son appartenance au secteur radical demeurent dans le registre de trajectoire.

Soient \(S_m\) la somme des trois résidus positifs additifs \(\Sigma\) de la fenêtre \(m\), \(E_m\) la somme des trois exposants discrets des résidus multiplicatifs \(\Pi\) et \(Z_m=3\) le nombre d’événements neutres. Ainsi, \(S_m\) ne somme pas les évaluations entières \(S=i+j\), dont les quotients demeurent dans le registre. Avec \([a]_{10}\in\{0,\ldots,9\}\), la règle d’émission décimale est
\begin{align}
 d_{m,1}&=[S_m]_{10},\notag\\
 d_{m,2}&=[S_m+E_m]_{10},\notag\\
 d_{m,3}&=[3S_m+5E_m+7Z_m]_{10},\notag\\
 U_m&=100d_{m,1}+10d_{m,2}+d_{m,3}.
 \label{eq:emisor-seis}
\end{align}
Les coefficients entiers et le calendrier de cette récurrence font explicitement partie de la loi d’émission déclarée. Aucune valeur de \(\pi,\varphi,e\) ni de \(\alpha\) n’est reçue. La distinction entre une règle d’émission et une constante reconnue à partir de sa sortie évite que la généalogie soit dissimulée dans une table de décimales.

Si \(s_m=[S_m]_{10}\) et \(e_m=[E_m]_{10}\), la même règle s’écrit
\begin{equation}
 U_m=100s_m+10[s_m+e_m]_{10}+[3s_m+5e_m+1]_{10}.
 \label{eq:acoplamiento}
\end{equation}
Le terme final égal à un est le résidu de \(7Z_m=21\) modulo dix. La réduction ternaire orientée du catalogue adopte la convention
\begin{equation}
 \Psi(U_1,\ldots,U_6)=([-U_1]_3,\ldots,[-U_6]_3).
 \label{eq:psi-catalogo}
\end{equation}
Le signe de cette coordinatisation est explicite : le changer modifie les mots et exige de transporter les orientations ultérieures. L'identification équilibrée de chaque composante de \(\FF_3\) au TRIT s'effectue après \eqref{eq:psi-catalogo}.

\begin{proposition}[Factorisation injective de l’émission]
Le mot \(U=(U_1,\ldots,U_6)\) détermine les deux signatures \(s=(s_1,\ldots,s_6)\) et \(e=(e_1,\ldots,e_6)\). L’image de l’émission possède \(18\cdot26=468\) éléments. La projection ternaire possède \(243\) valeurs distinctes dans l’espace ambiant de \(729\) mots.
\end{proposition}
\begin{proof}
Les centaines de \(U_m\) permettent de retrouver \(s_m\). Si \(b_m\) est son chiffre des dizaines, alors \(e_m=[b_m-s_m]_{10}\). Le chiffre des unités vérifie la troisième congruence de \eqref{eq:acoplamiento}. Le couplage est donc injectif sur le produit des signatures.

L’énumération des \(324\) conditions d’un curseur par la récurrence précédente produit dix-huit signatures additives, chacune ayant dix-huit antécédents ; et vingt-six signatures multiplicatives, dont vingt-quatre ont huit antécédents, une en a vingt-quatre et une en a cent huit. L’annexe conserve ces signatures et leurs conditions initiales, ce qui permet de répéter l’énumération finie sans catalogue numérique de référence. L’injectivité du couplage donne \(18\cdot26\) émissions et les multiplicités
\[
 \underbrace{432}_{18\cdot24}\text{ fibres de }144,
 \qquad18\text{ fibres de }432,
 \qquad18\text{ fibres de }1944.
\]
La somme est \(432\cdot144+18\cdot432+18\cdot1944=104\,976\). Appliquer \eqref{eq:psi-catalogo} aux \(468\) émissions produit exactement \(243\) mots. Ce dernier nombre est le cardinal de l’image, non celui de l’espace \(\FF_3^6\).
\end{proof}

Cette factorisation explique la différence entre exhaustivité finie et profondeur illimitée. Le domaine \eqref{eq:semillas} est complet pour le support et l’émetteur déclarés ; il ne contient pas par anticipation chaque préfixe de chaque prolongement. La profondeur appartient au système d’extensions compatibles construit sur ces conditions.

\subsection{Clôture du calendrier et progression de la mémoire}

Le calendrier de cent huit transitions se compose de douze fenêtres nonadiques. Sa structure de groupe est plus précise qu’un produit informel de deux périodes :
\begin{equation}
 0\longrightarrow C_{12}\xrightarrow{\iota}C_{108}
 \xrightarrow{p}C_9\longrightarrow0,
 \quad\iota([b])=[9b],\quad p([t])=[t]_9.
 \label{eq:extension108}
\end{equation}
Si \(t=9b+r\) avec \(0\le r<9\), la somme de deux phases produit la retenue
\[
 (b,r)+(b',r')=
 \left(b+b'+\left\lfloor\frac{r+r'}9\right\rfloor,\ [r+r']_9\right),
\]
où la première coordonnée est réduite modulo douze. La composition conserve le terme qui relie les deux coordonnées.

\begin{proposition}[Non-scindement du calendrier]
La suite \eqref{eq:extension108} est exacte et n’admet pas de section qui soit un homomorphisme de groupes.
\end{proposition}
\begin{proof}
Le noyau de la réduction modulo neuf est le sous-groupe des multiples de neuf, image de \(\iota\). Si l’extension était scindée, la commutativité donnerait \(C_{108}\cong C_{12}\times C_9\). Le premier groupe est d’exposant \(108\), tandis que le second est d’exposant \(\operatorname{mcm}(12,9)=36\), d’où une contradiction.
\end{proof}

Neuf transitions rétablissent la phase nonadique et font avancer d’une fenêtre ; cent huit rétablissent le calendrier fini. Aucun de ces faits n’impose d’effacer la trajectoire ou de rétablir le préfixe d’un prolongement. L’holonomie nonadique \(\Gamma_9\) désigne le transport d’un tour sur des états avec mémoire, de sorte que
\begin{equation}
 g(\Gamma_9x)=g(x),\qquad M(\Gamma_9x)=M(x)+1.
 \label{eq:holonomia-memoria}
\end{equation}
L’égalité de phase et l’incrément de mémoire sont compatibles et constituent deux observations du même transport.

\input{sections/temporalidad_trit.tex}

La représentation monodromique de ce retour est développée dans la
section~\ref{mono-ampl:conexion}. L’holonomie de
l’état y est distinguée de sa représentation sur la phase et le compteur. La suspension
symbolique de la section~\ref{mono-ampl:cristal} définit le modèle de
cristal temporel apériodique, et la section~\ref{mono-ampl:euler}
compose cette temporalité avec les trois réalisations d’Euler.

\subsection{Compression observable et conservation orthogonale}

Une réalisation isométrique permet d’exprimer quantitativement ce que perd une observation réduite. Soit \(U:\mathcal H\to\mathcal H\) une réalisation isométrique du transport d’un tour, \(U^*U=I\), et \(J:\mathcal H_{\rm vis}\to\mathcal H\) une inclusion isométrique. Nous définissons
\[
 T=J^*UJ,\qquad D=(I-JJ^*)UJ.
\]
Alors
\begin{equation}
 UJ=JT+D,\qquad J^*D=0,\qquad I-T^*T=D^*D.
 \label{eq:conservacion-ortogonal}
\end{equation}
La dernière égalité se déduit en prenant le produit adjoint de la première et en utilisant l’orthogonalité. La contraction de la composante visible ne constitue pas à elle seule une destruction de l’information de l’état total.

\begin{proposition}[Conservation à horizon fini]
Si \(I-T^*T=D^*D\), alors, pour tout \(n\ge1\),
\begin{equation}
 I=\sum_{r=0}^{n-1}T^{*r}D^*DT^r+T^{*n}T^n.
 \label{eq:telescopia}
\end{equation}
\end{proposition}
\begin{proof}
Multiplier \(D^*D=I-T^*T\) par \(T^{*r}\) et \(T^r\), puis sommer, produit une somme télescopique. Les termes intérieurs s’annulent et il reste \(I-T^{*n}T^n\).
\end{proof}

Le dernier terme conserve l’état terminal. Le supprimer avant de justifier une limite modifierait l’égalité. Cette distinction est le contenu opératoire de la conservation de l’unité dans la compression : la norme se répartit entre la composante visible, la mémoire et le terminal, sans être remplacée par un résidu unique.

\subsection{Développement local : commutation, signature lorentzienne et courbure discrète}

Le bloc diagonal adjacent dont les diagonales sont égales modulo neuf est unique. En effet,
\[
 k^2\equiv(k+1)^2\pmod9\quad\Longleftrightarrow\quad2k+1\equiv0\pmod9
\]
donne \(k=4\). Sa matrice visible est
\[
 B_c=\begin{pmatrix}7&2\\2&7\end{pmatrix}=7I+2R,
 \quad R=\begin{pmatrix}0&1\\1&0\end{pmatrix},\quad R^2=I.
\]
Avec \(P_\pm=(I\pm R)/2\), on obtient \(B_c=9P_++5P_-\). L’unité de la normalisation stochastique et la conservation d’une forme indéfinie sont deux réalisations de ce bloc :
\[
 \frac{B_c}{9}=P_++\frac59P_-,\qquad
 \frac{B_c}{\sqrt{45}}=\exp(\eta R),\quad\eta=\frac12\log\frac95.
\]
Pour \(Q=\operatorname{diag}(1,-1)\), la seconde matrice satisfait \(B_c^{\mathsf T}QB_c=45Q\). Il en résulte une réalisation locale dans \(SO^+(1,1)\) ; il n’est pas nécessaire d’introduire au préalable un espace-temps à quatre dimensions pour la démontrer.

La moyenne uniforme sur les classes \(C_1=\{1,4,7\}\), \(C_2=\{2,5,8\}\), \(C_0=\{3,6,9\}\) donne
\[
 M_\Pi=\begin{pmatrix}4&5&6\\5&4&6\\6&6&9\end{pmatrix},\quad
 M_\Sigma=\begin{pmatrix}5&6&4\\6&4&5\\4&5&6\end{pmatrix}.
\]
Le polynôme caractéristique de \(M_\Pi\) est
\[
 (\lambda+1)((\lambda-9)^2-72),
\]
d’où la signature \((2,1)\) de sa forme quadratique. L’écart spectral local \(9-5=4\), le coefficient \(2\), la différence agrégée \(51-45=6\) et le défaut intégré \(9\cdot6=54\) décrivent des niveaux distincts de la même comparaison arithmétique. Ce ne sont pas des indices interchangeables d’un même opérateur.

La courbure discrète peut être décrite directement avant cette réalisation spectrale. Dans les résidus, soient \(S(x,y)=x+y\), \(P(x,y)=xy\) et \(\Delta_x,\Delta_y\) les différences progressives. Les variables de liaison sont des différences de potentiel, \(A_x^T=\Delta_xT\), \(A_y^T=\Delta_yT\), et l’on fixe
\[
 F(T_x,T_y)=\Delta_x\Delta_yT_y-\Delta_y\Delta_xT_x.
\]
Puisque \(\Delta_xS=\Delta_yS=1\), \(\Delta_xP=y\) et \(\Delta_yP=x\),
\begin{equation}
 F(S,S)=F(P,P)=0,\qquad F(S,P)=+1,\quad F(P,S)=-1\pmod9.
\end{equation}
La somme et le produit, pris séparément, induisent des connexions plates ; leur mélange ordonné change le signe de la courbure de plaquette. En sommant sur un rectangle de côtés \(a,b\), les arêtes intérieures s’annulent et la somme orientée de bord vaut \(\pm ab\) modulo neuf. Une relation exacte apparaît ainsi entre l’ordre de composition et l’orientation de la frontière.

\subsection{Complétion cubique et deux notions de frontière}

L’extension du support cubique de côté neuf à celui de côté dix admet une stratification positionnelle exacte :
\begin{equation}
 10^3=9^3+3\cdot9^2+3\cdot9+1=729+243+27+1.
 \label{eq:cubo}
\end{equation}
Le premier terme correspond aux trois coordonnées intérieures ; le deuxième, à une nouvelle coordonnée ; le troisième, à deux ; et le dernier, au sommet comportant trois nouvelles coordonnées. La couronne d’extension contient \(271\) positions. En revanche, la frontière intérieure du cube discret de côté neuf en contient \(9^3-7^3=386\). La différence entre ces dénombrements est une différence de domaine, non un conflit de valeurs.

Diviser \eqref{eq:cubo} par \(1000\) donne une partition normalisée de l’unité. L’incidence du cube distingue six faces, douze arêtes et huit sommets. Le lien ultérieur avec les hexades et les octades de Steiner requiert les applications combinatoires correspondantes ; les cardinaux du cube ne s’y substituent pas. Cette complétion apporte ici la relation entre les bases \(729\) et \(1000\), tandis que le prolongement exceptionnel sera développé après la génération numérique et l’électron.

\input{sections/revision_emrh.tex}

\input{sections/geometria_modular.tex}
\input{sections/aritmetica_historias.tex}

\section{Génération coinductive de \texorpdfstring{\(\pi,\varphi,e\)}{pi, phi, e} à profondeur arbitraire}
\label{sec:generacion}

L'émission de six blocs est une étape finie de la construction, et non un développement décimal complet. Sa fonction consiste à produire des régions munies de multiplicités, d'orientations et de règles de transport. Le prolongement doit conserver ces données et déterminer une collection compatible de préfixes de longueur arbitraire. Le mot \emph{génération} désigne cette opération ; \emph{généalogie} désigne la reconstruction de ses dépendances. Les deux notions interviennent, mais remplissent des fonctions distinctes.

\subsection{Orbites, régions et caractères de clôture, de propagation et d'auto-échelle}

Sur un mot \(w=(w_1,\ldots,w_6)\), nous définissons
\[
 r(w)=(w_3,w_1,w_2,w_5,w_6,w_4),\qquad
 s(w)=(w_4,w_5,w_6,w_1,w_2,w_3).
\]
On a \(r^3=s^2=I\) et \(srs=r^{-1}\) ; ces permutations réalisent donc le groupe diédral \(D_3\). L'action conserve le catalogue et les multiplicités et commute avec \(\Psi\). Le quotient des \(243\) mots visibles possède \(43\) orbites : trente-huit de cardinal six et cinq de cardinal trois.

La classe de clôture se reconnaît par un prédicat de fibre : ses six orientations ont cinq émissions par mot, une multiplicité totale \(1008\) et une distribution \(432+4\cdot144\). Il existe une seule orbite possédant ces propriétés dans le catalogue. Pour préciser les deux autres sélections, nous écrivons \(f(w)=\#\Psi^{-1}(w)\), \(m(w)=\sum_{U\in\Psi^{-1}(w)}\operatorname{mult}(U)\) et \(\nu(w)=(n_0,n_1,n_2)\). La classe diagonale additive d'une émission est \(d(U)=[i+j]_9\) dans les indices de zéro à huit ; sa signature additive conserve cette classe et l'orientation \(ES\) ou \(NO\). Soit \(\mathcal D(\mathcal O)\) l'ensemble de ces classes pour toutes les émissions sur une orbite.

Parmi les orbites de taille six, le prédicat conjoint
\[
 f(w)=1,\qquad m(w)=144\quad(w\in\mathcal O),\qquad
 \mathcal D(\mathcal O)=\{0,3,6\}
\]
laisse exactement six orbites. Dans ce secteur, le recensement \(\nu=(2,3,1)\), égal à celui de clôture, sélectionne uniquement la propagation ; le recensement \(\nu=(2,2,2)\) sélectionne uniquement l'auto-échelle. Le support diagonal ne peut pas être omis : sans lui, il y a quatre orbites à un seul point équilibrées de multiplicité \(144\), et non une seule. Enfin, l'existence d'une émission avec \(d=6\) et une orientation \(ES\) sélectionne un représentant de chaque orbite et conduit à
\begin{equation}
 w^{\rm cl}=010211,\qquad w^{\rm pr}=201101,\qquad w^{\rm au}=121200.
 \label{eq:palabras-regionales}
\end{equation}
Ces symboles désignent d'abord des régions du catalogue. L'identification scalaire sera démontrée au moyen de leurs lecteurs ; ils ne sont pas choisis parce qu'ils coïncident avec des chiffres d'une constante connue.

Le mot de clôture \(010211\) possède les cinq préimages suivantes :
\begin{center}\small
\begin{tabular}{@{}lll r@{}}\toprule
Émission de six blocs&Signature multiplicative&Rôle&Multiplicité\\\midrule
\(501,614,498,169,272,272\)&\(555555\)&transversal&432\\
\(810,923,870,169,272,272\)&\(339555\)&orientation positive&144\\
\(810,983,810,169,272,272\)&\(393555\)&orientation positive&144\\
\(870,923,810,169,272,272\)&\(933555\)&orientation positive&144\\
\(870,923,810,418,674,521\)&\(933717\)&retour orienté&144\\\bottomrule
\end{tabular}
\end{center}
La multiplicité \(432\) et la signature constante \(555555\) caractérisent la région transversale. Elles n'identifient pas la position centrale \((5,5)\) d'APP. Confondre une signature de six fenêtres avec une position du support éliminerait la distinction entre région numérique et structure électronique.

La microfibre possède en outre une morphologie bipartite précise. Chaque émission est un produit d'une classe de curseur additif et d'une classe multiplicative. Les trois régions positives partagent la même classe additive de dix-huit conditions ; ensemble, elles contiennent trois classes multiplicatives de huit conditions chacune. La région transversale possède un produit \(18\times24\) ; les trois positives réunies, un autre \(18\times24\) ; et le retour, un produit \(18\times8\). La décomposition par émissions comporte cinq parties, tandis que la décomposition en composantes connexes du graphe biparti en comporte trois. Toutes deux conservent les mêmes \(1008\) germes sans réduire la forme à un cardinal.

\subsection{Transport fini et frontière de prolongement}

Les endomorphismes \(L_0,L_1\) de l'espace des mots ternaires agissent avant la réalisation archimédienne. Leur détermination commence par les transitions de l'état arithmétique, et non par la prescription de deux matrices. Si \(U_t\) est l'étape composée de sélection, de transport et de mise à jour, le bloc observé satisfait
\[
 b_j^\chi=\Pi_6(x_j^\chi),\qquad
 x_{j+1}^\chi=U_{9j+9}\cdots U_{9j+1}(x_j^\chi),
 \qquad \chi\in\{\mathrm{cl},\mathrm{pr},\mathrm{au}\}.
\]
On forme six transitions pour chaque régime : deux consécutives dans chacune des trois régions. Les couples de matrices d'origines et de destinations obtenus dans le registre sont
\[
\begin{array}{c|c|c|c}
X_0&Y_0&X_1&Y_1\\\hline
010211&012222&010211&002111\\
012222&010211&002111&110221\\
201101&121221&102011&012222\\
121221&102011&012222&102011\\
121200&112202&121020&010210\\
112202&121020&010210&010200
\end{array}
\]
Chaque chaîne de six chiffres représente une ligne dans \(\FF_3^6\). Puisque \(\det X_0=1\) et \(\det X_1=-1\) dans \(\FF_3\), les équations \(X_aL_a=Y_a\) déterminent de manière unique \(L_a=X_a^{-1}Y_a\). Dans cette coordinatisation par vecteurs-lignes, on obtient
\[
 L_0=\begin{pmatrix}
 2&2&2&1&2&1\\2&1&2&2&1&1\\1&0&0&1&0&2\\
 0&0&1&1&1&0\\2&2&2&0&2&1\\2&1&2&1&0&0
 \end{pmatrix},\quad
 L_1=\begin{pmatrix}
 2&1&2&0&2&2\\1&2&1&1&2&0\\1&2&1&1&1&2\\
 0&1&0&0&2&1\\2&0&2&1&0&1\\0&2&2&2&1&1
 \end{pmatrix}\quad(\bmod3).
\]
Le calendrier \((L_0,L_0,L_1,L_1)\) produit quatre nouveaux blocs de six composantes. Leurs concaténations forment
\[
 w_6\longrightarrow w_{12}\longrightarrow w_{18}
 \longrightarrow w_{24}\longrightarrow w_{30}.
\]
Les sous-indices de \(w\) indiquent la longueur accumulée. La frontière \(R_{36}\) conserve la structure terminale et ouvre la relation de prolongement ; elle n'est pas renommée comme un dernier bloc périodique qui pourrait se répéter indéfiniment.

L'identité \(L=X^{-1}Y\) démontre l'unicité une fois les transitions fixées ; elle ne remplace pas la production de ces transitions par \(U_t\). La provenance de leur registre et l'examen des seize calendriers binaires sont documentés dans le supplément technique. Le calendrier indiqué est le seul de ces seize calendriers qui reproduise conjointement les trois suites de cinq blocs.

\subsection{Cinq régions de clôture et compensation angulaire}

La région de clôture contient une composante transversale et quatre accompagnatrices. Le lecteur symétrique agit dans l'espace de leurs cinq positions au moyen de
\[
 P_5=\frac15\mathbf1\mathbf1^{\mathsf T},\qquad
 \mathbf1=(1,1,1,1,1)^{\mathsf T}.
\]
L'invariance \(P_5\lambda=\lambda\) et la conservation \(\mathbf1^{\mathsf T}\lambda=1\) imposent \(\lambda=\mathbf1/5\). Cette normalisation répartit l'unité entre les positions du lecteur ; elle ne remplace pas les multiplicités des germes \(432,144,144,144,144\) par des fréquences uniformes. La coordinatisation elliptique du lecteur réalise chaque coordonnée par \(1+\lambda_jJ\) ; sa pente est \(\lambda_j\), d'où la pente primitive \(1/5\). La règle qui fait passer d'une coordonnée normalisée à une pente est ainsi explicitée, au lieu d'identifier les deux notions sans application. Le lecteur compose les quatre positions associées selon la loi
\begin{equation}
 a\oplus b=\frac{a+b}{1-ab}.
 \label{eq:suma-pendientes}
\end{equation}
Cette loi s'obtient en multipliant \((1+aJ)(1+bJ)\) dans le régime \(J^2=-I\) et en divisant par sa composante scalaire. La composition des pentes est donc une réalisation de la structure quadratique de TRIT, et non une somme ordinaire de quatre fractions.

\begin{proposition}[Compensateur de la clôture]
La composition de quatre pentes \(1/5\) est \(120/119\). La condition de retour à une pente un détermine un unique compensateur positif de la forme \(1/q\), avec \(q>1\), et donne \(q=239\).
\end{proposition}
\begin{proof}
La première duplication donne \((1/5)\oplus(1/5)=5/12\) ; la seconde, \((5/12)\oplus(5/12)=120/119\). La condition
\[
 \frac{120/119-1/q}{1+(120/119)/q}=1
\]
équivaut à \((120-119)q=120+119\), d'où \(q=239\).
\end{proof}

Le caractère de clôture résultant a pour réalisation
\begin{equation}
 \Pi_{\rm cl}=4\left(4\arctan\frac15-\arctan\frac1{239}\right).
 \label{eq:lector-pi}
\end{equation}
L'identité angulaire de Machin reconnaît \(\Pi_{\rm cl}=\pi\). Le sens de la construction est précis : la pente et le nombre de compositions sont déclarés par le lecteur de la pentafibre ; l'équation de compensation impose \(239\). Une fois ce caractère produit, son évaluation par des séries alternées ne sélectionne pas à nouveau l'orbite des germes.

Pour \(0<x<1\), les sommes alternées
\[
 A_n(x)=\sum_{j=0}^n\frac{(-1)^jx^{2j+1}}{2j+1}
\]
encadrent \(\arctan x\) entre deux troncatures consécutives, dont la différence est \(x^{2n+3}/(2n+3)\). Appliquées aux deux pentes de \eqref{eq:lector-pi}, elles fournissent des extrémités rationnelles d'erreur explicite. La valeur numérique est ainsi calculée comme conséquence du caractère exact, sans introduire de préfixe cible.

\subsection{Lecteur de propagation et normalisation de l'unité}

La région de propagation se réalise au moyen d'une collection formelle \(P_t\in\mathcal A[[t]]\), où \(\mathcal A\) est une algèbre commutative unitaire sur \(\QQ\). Sa composition satisfait \(P_{t+s}=P_tP_s\), d'unité \(P_0=1\), et le transport élémentaire orienté fixe le générateur scalaire \(P'_0=+1\). La condition fixe son signe, et non sa seule grandeur. En écrivant
\[
 P_t=\sum_{n\ge0}a_nt^n,
\]
la normalisation détermine \(a_0=a_1=1\). La comparaison du coefficient de \(s^nt\) dans \(P_{s+t}=P_sP_t\) donne \((n+1)a_{n+1}=a_na_1\) ; de manière équivalente, \(P'_t=P_t\). Ainsi,
\begin{equation}
 (n+1)a_{n+1}=a_n,\qquad a_n=\frac1{n!}.
 \label{eq:propagacion}
\end{equation}
Le caractère en une unité de paramètre est \(\Pi_{\rm pr}=P_1\), reconnu ultérieurement comme \(e\).

La normalisation du générateur importe. La loi de semi-groupe sans cette donnée admettrait \(P_t=e^{ct}\) pour différents \(c\) ; c'est pourquoi on n'omet pas la condition qui fixe l'unité du paramètre dans le lecteur. Cette condition est antérieure à la comparaison décimale et appartient à l'exposition de la règle opératoire, et non au choix d'une valeur expérimentale.

\begin{proposition}[Bornes rationnelles de propagation]
Pour \(n\ge1\), soit \(E_n=\sum_{j=0}^n1/j!\). Alors
\[
 E_n<\Pi_{\rm pr}<E_n+\frac{n+2}{(n+1)(n+1)!}.
\]
\end{proposition}
\begin{proof}
Le premier terme omis est \(1/(n+1)!\). Chaque quotient ultérieur est au plus égal à \(1/(n+2)\). La somme géométrique majorant la queue est
\[
 \frac1{(n+1)!}\frac1{1-1/(n+2)}
 =\frac{n+2}{(n+1)(n+1)!}.
\]
L'inégalité stricte découle de la décroissance des quotients successifs.
\end{proof}

\subsection{Incidence modale et génération de l'auto-échelle}

Le sélecteur tritique produit le cycle \((+1,-1,0)\). Dans les modes \(\Sigma,\Pi\), la règle est \(+1:m\mapsto\Sigma\), \(-1:m\mapsto\Pi\), \(0:m\mapsto m\). Avec fermeture cyclique, le mode observé est \((\Sigma,\Pi,\Pi)\). Ses arêtes sont \(\Sigma\to\Pi\), \(\Pi\to\Pi\) et \(\Pi\to\Sigma\). Son incidence, dans cet ordre des modes, est donc
\begin{equation}
 F_{\rm av}=\begin{pmatrix}0&1\\1&1\end{pmatrix}.
 \label{eq:fav}
\end{equation}
Les quatre entrées s'obtiennent à partir du calendrier, et non d'une équation qui aurait préalablement reçu la valeur de \(\varphi\). La répétition du calendrier donne les dénombrements \(3F_{\rm av},9F_{\rm av},36F_{\rm av}\) aux longueurs neuf, vingt-sept et cent huit. Ces dénombrements ne sont pas les puissances de la matrice.

\begin{theorem}[Caractère positif d'auto-échelle]
L'incidence \eqref{eq:fav} possède un unique rayon propre strictement positif. En le normalisant sous la forme \((1,x)^{\mathsf T}\), sa coordonnée \(x\) satisfait \(x^2-x-1=0\), \(x>0\), et vaut \(\varphi=(1+\sqrt5)/2\).
\end{theorem}
\begin{proof}
L'équation propre donne \(x=\lambda\) et \(1+x=\lambda x\). L'élimination de \(\lambda\) produit le polynôme. Une seule de ses deux racines est positive. Comme \(F_{\rm av}^2\) a toutes ses entrées positives, le rayon positif est unique et simple. L'expression radicale reconnaît la coordonnée obtenue ; elle ne participe pas à la production de \(F_{\rm av}\).
\end{proof}

Avec \(F_0=0,F_1=1,F_{n+1}=F_n+F_{n-1}\), l'incidence induit, pour \(n\geq1\),
\[
 F_{\rm av}^n=\begin{pmatrix}F_{n-1}&F_n\\F_n&F_{n+1}\end{pmatrix}.
\]
Les quotients consécutifs \(F_{n+1}/F_n\) alternent autour de \(\varphi\). L'identité de déterminant
\[
 F_{n+1}^2-F_nF_{n+2}=(-1)^n
\]
donne une largeur exacte \(1/(F_nF_{n+1})\) entre deux quotients consécutifs. Sa croissance rend cet encadrement rationnel arbitrairement petit.

La mesure sur toutes les histoires compatibles ne coïncide pas avec la fréquence des modes du calendrier à trois événements. La fréquence chronologique est \((1/3,2/3)\). La mesure d'entropie maximale de l'enveloppe symbolique définie par \(F_{\rm av}\) se construit à partir de ses vecteurs propres et a pour rapport de poids \(\varphi^{-2}\). Cette dernière intervient dans la lecture surfacique--volumétrique ; elle ne s'obtient pas en remplaçant simplement une fréquence de comptage par une autre.

\input{sections/retorno_areal_volumetrico.tex}

Une fois les caractères générés, les deux orientations du lecteur
\[
 R(x,y)=\frac{x}{y^2}
\]
fournissent
\begin{equation}
 \rho_A=\frac\pi{\varphi^2},\qquad
 \rho_E=\frac\varphi{\pi^2},\qquad
 \varphi^3\rho_A^2\rho_E=1.
 \label{eq:areal-electron}
\end{equation}
La seconde expression apparaîtra dans l'exposant électronique. La première s'obtient également comme limite de la lecture marquée \(\pi F_{n-1}/F_{n+1}\). L'échange des arguments les relie, mais ne transforme pas un reparamétrage de l'équation électronique en une deuxième dérivation indépendante de tous ses coefficients.

\input{sections/euler_trit_articulacion.tex}

\subsection{Cylindres compatibles et profondeur arbitraire}

La réalisation d'un bloc ternaire \(b=(b_1,\ldots,b_6)\) est l'entier
\[
 \nu(b)=\sum_{j=1}^6b_j3^{6-j}\in\{0,\ldots,728\}.
\]
Une suite de blocs définit des préfixes entiers par \(N_{n+1}=729N_n+\nu(b_{n+1})\), en commençant en \(N_0=m\), où \(m\) est la partie entière déterminée par le lecteur. Les fourchettes précédentes situent les caractères de clôture, de propagation et d'auto-échelle dans \((3,4)\), \((2,3)\) et \((1,2)\), respectivement ; leurs valeurs initiales sont donc \(m=3,2,1\). Les blocs suivants raffinent la partie fractionnaire. Le cylindre correspondant est
\begin{equation}
 I_n=\left[\frac{N_n}{729^n},\frac{N_n+1}{729^n}\right],
 \qquad |I_n|=729^{-n}.
 \label{eq:cilindro}
\end{equation}
La convention semi-ouverte est utilisée pour sélectionner un représentant lorsqu'une valeur tombe sur une frontière. Pour le passage à la limite, on utilise les fermetures de ces intervalles et on identifie les deux développements de frontière équivalents. Cette précision empêche de supposer que toute intersection d'intervalles semi-ouverts emboîtés est automatiquement non vide.

\begin{theorem}[Détermination de la réalisation scalaire]
Une histoire de blocs compatible détermine une unique limite des extrémités gauches de \eqref{eq:cilindro}. Si un caractère exact dispose d'encadrements rationnels de largeur arbitrairement petite et se sépare de la frontière de chaque cylindre décimal examiné, chaque préfixe décimal de longueur finie est déterminé après un calcul fini. Les préfixes ainsi obtenus sont compatibles entre eux.
\end{theorem}
\begin{proof}
L'extension d'un bloc conserve le préfixe précédent par division euclidienne. Les cylindres fermés sont emboîtés et leur diamètre tend vers zéro ; leurs extrémités gauches forment une suite de Cauchy et possèdent une unique limite. Pour une longueur décimale fixée, une valeur séparée des frontières de la partition possède une distance positive à celles-ci. Une fourchette suffisamment étroite est contenue dans une unique cellule et détermine son préfixe. L'unicité de la cellule et l'emboîtement des partitions prouvent la compatibilité par troncature. Aux points à développement terminal, il faut ajouter la décision exacte de frontière et sa convention de représentation ; cette décision ne se déduit pas d'une largeur numérique seule.
\end{proof}

\begin{corollary}[Obtention finie et compatibilité des préfixes régionaux]
\label{gen:prefijos-regionales}
Pour chacune des réalisations \(\pi,\varphi,e\) et pour toute longueur
décimale finie, les encadrements construits déterminent le préfixe correspondant
par un calcul fini. Les préfixes obtenus à différentes profondeurs
sont compatibles par troncature.
\end{corollary}
\begin{proof}
Les caractères de clôture, de propagation et d'auto-échelle possèdent les encadrements précédents. Après leur reconnaissance, l'irrationalité de \(\pi,e,\varphi\) garantit leur séparation des frontières rationnelles de chaque partition décimale finie. On applique le théorème précédent à chaque caractère.
\end{proof}
La profondeur arbitraire exprime le quantificateur du corollaire : pour tout
horizon fini, il existe un calcul qui détermine son préfixe. L'objet
mathématique est la collection compatible complète ; chaque exécution en obtient une
projection finie.

\subsection{Involution spéculaire et synchronisation des régions}

La connexion nonadique transporte conjointement les trois caractères. Dans la coordinatisation de la neuvième phase, l'axe d'auto-échelle et la somme ternaire de clôture et de propagation restent fixés sous l'échange spéculaire. L'inversion hélicoïdale change le sens du transport ; la réflexion des deux canaux change leur orientation relative. Ce sont des involutions différentes, bien qu'elles agissent sur le même système de trajectoires.

Les coïncidences de recensement illustrent cette coordination :
\begin{center}
\begin{tabular}{@{}c rrr l@{}}\toprule
Phase&\(N_\pi\)&\(N_e\)&\(N_\varphi\)&Égalité des dénombrements\\\midrule
4&207&207&122&\(N_\pi=N_e\)\\
5&39&151&151&\(N_e=N_\varphi\)\\
6&110&110&69&\(N_\pi=N_e\)\\
7&81&29&81&\(N_\varphi=N_\pi\)\\
8&59&59&59&synchronisation triple\\
9&43&43&43&synchronisation triple\\\bottomrule
\end{tabular}
\end{center}
Ces entiers comptent des coordonnées compatibles des régions. Ils n'affirment pas une égalité entre \(\pi\), \(e\) et \(\varphi\) en tant que nombres réels. Si
\[
 \partial(a,b,c)=(a-b,b-c,c-a),
\]
son image appartient au réseau \(A_2=\{(u,v,w)\in\ZZ^3:u+v+w=0\}\). Aux phases huit et neuf, la composante relative s'annule, mais la composante diagonale passe de \((59,59,59)\) à \((43,43,43)\). La synchronisation relative coexiste avec un déplacement de \(-16\) par canal. La fermeture qui conduira à \(\alpha\) doit conserver les deux informations : coïncidence relative et évolution de la mémoire.

\subsection{Suspension hélicoïdale sturmienne}

La relation entre les bases neuf et dix de la complétion définit le paramètre adimensionnel
\begin{equation}
 \delta=\frac{\log(10/9)}{\log10}.
 \label{eq:delta-sturm}
\end{equation}
Elle ne s'obtient pas à partir d'une constante physique cible. C'est l'échelle logarithmique relative de deux bases déjà présentes dans la construction. Son emploi comme lecture symbolique du raffinement a une conséquence exacte.

\begin{lemma}[Irrationalité de la pente]
Le nombre \(\delta\) est irrationnel.
\end{lemma}
\begin{proof}
Si \(\delta=p/q\), \(q>0\), alors \((10/9)^q=10^p\) et \(9^q=10^{q-p}\). La factorisation du premier membre contient le nombre premier trois et celle du second uniquement les nombres premiers deux et cinq. L'égalité est impossible.
\end{proof}

Le mot mécanique inférieur, d'ordonnée à l'origine \(\theta\), est
\[
 z_n(\theta)=\lfloor(n+1)\delta+\theta\rfloor-\lfloor n\delta+\theta\rfloor\in\{0,1\}.
\]
On l'obtient en observant la rotation \(\theta\mapsto\theta+\delta\pmod1\) au moyen de l'intervalle \([1-\delta,1)\). La suite d'événements a pour densité \(\delta\), et la somme sur toute fenêtre de longueur \(L\) vaut \(\lfloor L\delta\rfloor\) ou \(\lceil L\delta\rceil\). Puisque \(21<1/\delta<22\), les séparations entre événements sont vingt et un ou vingt-deux. Le mot qui indique quels sont les intervalles courts est un autre mot mécanique de pente \(\sigma=22-1/\delta\) ; ses événements sont séparés par six ou sept intervalles primaires. La seconde échelle est induite par la première et n'est pas un second paramètre libre.

\input{sections/vacancias_capacidad.tex}

\begin{theorem}[Complexité des lectures de phase]
Pour \(m\ge1\), le mot mécanique a une complexité factorielle \(p(m)=m+1\). En conservant la phase \(n\bmod9\), la complexité est \(p_9(m)=9(m+1)\) ; en conservant seulement sa classe ternaire, elle est \(p_3(m)=3(m+1)\). Les trois entropies topologiques sont nulles.
\end{theorem}
\begin{proof}
Les prédécesseurs des deux extrémités de l'intervalle de codage forment, pour les facteurs de longueur \(m\), \(m+1\) points distincts du cercle : la coïncidence entre extrémités translatées élimine les doublons et l'irrationalité empêche toute coïncidence supplémentaire. Ces points déterminent \(m+1\) intervalles à itinéraires constants et distincts. Toute orbite irrationnelle visite chaque intervalle, d'où \(p(m)=m+1\).

Une phase modulo neuf étant fixée, les positions de cette phase parcourent une rotation de pas \(9\delta\), également irrationnel. Par densité, tous les facteurs apparaissent dans chaque phase, et la première coordonnée distingue leurs neuf copies. Le même argument avec le pas \(3\delta\) donne trois copies après réduction ternaire. Enfin, \(m^{-1}\log(c(m+1))\to0\) pour \(c=1,3,9\).
\end{proof}

Le produit de phases \(C_9\times\mathbb T\), translaté par \((1,\delta)\), possède des caractères de fréquence
\[
 \frac a9+k\delta\pmod1,\qquad a\in\ZZ/9\ZZ,\ k\in\ZZ.
\]
Le module spectral n'est pas \(\log10\) : c'est le module additif engendré par \(1/9\) et \(\delta\), modulo les entiers. Le facteur \(\log10\) appartient à la normalisation de \eqref{eq:delta-sturm}.

Une présentation opératorielle rend la mémoire explicite. Dans \(\ell^2(\ZZ)\), soient \(T|n\rangle=|n+1\rangle\), \(V_9|n\rangle=\zeta_9^n|n\rangle\), \(W_\delta|n\rangle=e^{2\pi i n\delta}|n\rangle\) et \(D_M|n\rangle=\lfloor n/9\rfloor|n\rangle\). Sur les vecteurs à support fini,
\begin{equation}
 V_9T=\zeta_9TV_9,\qquad
 W_\delta T=e^{2\pi i\delta}TW_\delta,\qquad
 [D_M,T^9]=T^9.
 \label{eq:weyl-relojes}
\end{equation}
Les deux premières relations sont des covariances de Weyl ; la troisième exprime qu'un tour nonadique augmente le degré de mémoire. La rotation de base est abélienne, tandis que l'algèbre des observables et des translations ne l'est pas. L'inversion \(|n\rangle\mapsto|-n\rangle\) conjugue l'avancée avec le recul. Cette réalisation fournit une suspension hélicoïdale apériodique à retour fini de phase, et non un mot périodique de longueur neuf.

\input{figures/holonomia.tex}

Pour ce dénombrement, nous fixons l'intercept \(\theta=0\) et les extrémités cumulées \((N_0,\ldots,N_4)=(0,729,972,999,1000)\), correspondant à l'ordre cubique \((729,243,27,1)\). Les dénombrements inférieurs sont \(\lfloor N_j\delta\rfloor-\lfloor N_{j-1}\delta\rfloor\) et donnent \((33,11,1,0)\). Les supérieurs sont \(\lceil N_j\delta\rceil-\lceil N_{j-1}\delta\rceil\), avec la même frontière initiale nulle, et donnent \((34,11,1,0)\). Ces valeurs dépendent de l'intercept et de l'ordre des strates ; elles ne sont pas affirmées pour une phase initiale arbitraire. Leurs totaux sont quarante-cinq et quarante-six. L'incrément correspond à une unique unité de frontière, située dans la première strate de cette partition ordonnée. Les nombres trente-quatre et onze apparaissent ainsi dans un même dénombrement antérieur à toute coordinatisation d'unités ; les identifier à des exposants physiques exige en outre la dérivation dimensionnelle correspondante.

La dénomination \emph{cristal apériodique} est employée ici pour cet ordre symbolique, dont la complexité, le spectre et le retour sont spécifiés. L'entropie topologique nulle n'équivaut pas à elle seule à une dissipation physique nulle : une réalisation thermodynamique exige une dynamique énergétique et une opération d'effacement définies. De même, cette lecture sturmienne ne remplace pas tous les opérateurs du TPK et ne démontre pas à elle seule la périodicité ou l'apériodicité de la sortie complète de \(\alpha\). Sa fonction est d'expliciter une structure d'ordre que le simple retour du calendrier ne décrit pas.

\input{sections/revision_monodromia.tex}

\input{sections/primos_retornos.tex}

% BEGIN NUCLEO_ADITIVO_20260921_SELECCION

% Ampliación demostrativa DF003--DF007. No modifica la fuente integral.
% Propietario común de los archivos Lean citados en estos comentarios:
% output/PAQUETE_ARTICULO_I_INTEGRACION_ESTADO_CAMPOS_20260919/.
% Residencia causal de la primera sección: después de los lectores regionales
% coinductivos y antes de emplear sus registros como entrada de la selección K.
% Residencia causal de la segunda sección: después de la construcción marcada
% Paley--Witt--Golay--Leech y antes de la reconstrucción de la estructura VOA.
% Los dos cortes expositivos conservan el mismo origen APP--TRIT--TPK.

\section[Génération régionale et raffinement]{Compatibilité intégrale de la génération régionale, du raffinement cylindrique et de l'incidence}
\label{act21:sec:df-lectores-cilindros}

La prolongation régionale produit des suites ordonnées de blocs, ainsi que
leurs profondeurs de lecture. La projection archimédienne et l'incidence
ternaire se calculent sur ces mêmes suites. Le lien entre ces deux
réalisations peut s'exprimer entièrement par des opérations entières :
sélection par intervalles rationnels, division euclidienne, élévation de six
trits, lecture de retenue et transformation de Witt. Cette formulation permet
de suivre chaque coefficient depuis son bloc générateur jusqu'à son emploi ultérieur.

La construction utilise les trois lecteurs régionaux déjà obtenus de l'état APP--TRIT--TPK. Nous les désignerons par les indices \(c\in\{\mathrm P,\mathrm E,\mathrm F\}\), correspondant respectivement aux régions de clôture, de propagation et d'auto-échelle. Leur réalisation réelle s'écrira \(v_c\). Les identifications ultérieures \(v_{\mathrm P}=\pi_{\mathrm{HMT}}\), \(v_{\mathrm E}=e_{\mathrm{HMT}}\) et \(v_{\mathrm F}=\varphi_{\mathrm{HMT}}\) conservent cet ordre causal. L'opération qui émet un bloc inspecte les intervalles rationnels du lecteur régional ; la valeur réelle intervient dans la démonstration de correction de cette opération.

\subsection{Émission positionnelle par séparation rationnelle de cellules}
\label{act21:subsec:df-emision-racional}

% Fuente: lean/biblioteca/RegionalPublicationComposition.lean,
% produced_brackets, eventual_publication, joint_eventual_publication,
% search_terminates, stopped_cell_correct, publications_compatible.

Pour chaque région, on dispose de deux suites rationnelles
\(\ell_c(j),u_c(j)\), construites par son lecteur, qui satisfont
\[
 \ell_c(j)\leq v_c\leq u_c(j).
\]
La propriété de séparation des cellules démontrée pour ces lecteurs affirme
que, pour chaque entier \(s>0\), il existe \(J\) tel que, pour tout \(j\geq J\),
\begin{equation}
 \lfloor s\ell_c(j)\rfloor
 =\lceil su_c(j)\rceil-1
 =\lfloor sv_c\rfloor.
 \label{act21:eq:df-separacion-celdas}
\end{equation}
L'évitement des frontières rationnelles par les valeurs régionales est l'hypothèse des théorèmes précédents qui garantit cette séparation. Ici, on emploie sa conclusion opérationnelle, avec les lecteurs et les intervalles qui la produisent.

\begin{definition}[Indice de séparation et préfixe émis]
Pour \(s>0\), soit
\[
 j_c(s)=\min\{j\in\mathbb N:
 \lfloor s\ell_c(j)\rfloor=\lceil su_c(j)\rceil-1\}.
\]
Pour une base entière \(b\geq2\) et une profondeur \(n\geq0\), on définit
\begin{equation}
 P_c(b,n)=\lfloor b^n\ell_c(j_c(b^n))\rfloor.
 \label{act21:eq:df-publicacion-prefijo}
\end{equation}
Le préfixe contient la partie entière et les premières \(n\) positions en base \(b\) ; par conséquent, \(P_c(b,0)\) est la partie entière régionale.
\end{definition}

\begin{theorem}[Terminaison et correction de l'émission]
\label{act21:thm:df-emision-correcta}
L'indice \(j_c(s)\) existe pour tout \(s>0\). Les préfixes émis satisfont
\begin{equation}
 P_c(b,n)=\lfloor b^n v_c\rfloor,
 \qquad
 P_c(b,n)\leq b^n v_c<P_c(b,n)+1.
 \label{act21:eq:df-prefijo-correcto}
\end{equation}
Pour chaque échelle \(s\) il est possible en outre de choisir un unique indice de suffisance \(J(s)\) valide simultanément pour les trois régions.
\end{theorem}

\begin{proof}
L'égalité à partir d'un certain rang de \eqref{act21:eq:df-separacion-celdas} rend non vide
l'ensemble qui définit \(j_c(s)\). Son minimum existe par le bon ordre de
\(\mathbb N\). À tout indice où l'égalité se produit, la borne
inférieure rationnelle fournit
\(\lfloor s\ell_c(j)\rfloor\leq\lfloor sv_c\rfloor\).
La borne supérieure, jointe à l'évitement d'une frontière de cellule, donne
\(\lfloor sv_c\rfloor\leq\lceil su_c(j)\rceil-1\).
Les deux extrémités entières coïncidant, l'entier intermédiaire coïncide avec
elles. Cela démontre \eqref{act21:eq:df-prefijo-correcto}. Pour l'affirmation
conjointe, il suffit de prendre le maximum des trois indices de suffisance obtenus
pour la même échelle. L'indice conjoint contrôle simultanément les trois
émissions, tandis que chaque région conserve ses propres intervalles.
\end{proof}

\begin{proposition}[Compatibilité exacte des préfixes]
\label{act21:prop:df-prefijos-compatibles}
Pour \(n,m\geq0\),
\begin{equation}
 \left\lfloor\frac{P_c(b,n+m)}{b^m}\right\rfloor=P_c(b,n).
 \label{act21:eq:df-truncamiento}
\end{equation}
En particulier, le bloc suivant
\(d_c(b,n)=P_c(b,n+1)\bmod b\) satisfait
\begin{equation}
 P_c(b,n+1)=bP_c(b,n)+d_c(b,n),\qquad 0\leq d_c(b,n)<b.
 \label{act21:eq:df-paso-posicional}
\end{equation}
\end{proposition}

\begin{proof}
Écrivons \(b^n v_c=N+\theta\), avec
\(N=P_c(b,n)\) et \(0\leq\theta<1\). Alors
\(P_c(b,n+m)=b^mN+\lfloor b^m\theta\rfloor\), et le second terme
appartient à \(\{0,\ldots,b^m-1\}\). La division euclidienne par \(b^m\)
récupère \(N\). La spécialisation \(m=1\) produit
\eqref{act21:eq:df-paso-posicional} et sa borne résiduelle.
\end{proof}

\subsection{Blocs régionaux de six trits à profondeur arbitraire}
\label{act21:subsec:df-bloques-arbitrarios}

% Fuente: deltas/integracion/JointRegionalFrontier.lean,
% publish_base729_eq_base3, generated_block_from_longer_prefix,
% generated_block_eq_finite, generated_signature_eq_finite.

L'égalité entière \(729=3^6\) détermine le regroupement de six positions ternaires en un bloc. Pour \(t\geq0\), nous définissons
\begin{equation}
 u_c(t)=P_c(729,t+1)\bmod729,
 \qquad
 B_{c,j}(t)=
 \left\lfloor\frac{u_c(t)}{3^{5-j}}\right\rfloor\bmod3,
 \quad 0\leq j<6.
 \label{act21:eq:df-bloque-y-banda}
\end{equation}
Ainsi, \(B_c(t)\) est un mot ordonné de six trits et
\begin{equation}
 u_c(t)=\sum_{j=0}^{5}B_{c,j}(t)3^{5-j}.
 \label{act21:eq:df-elevacion-banda}
\end{equation}
Le mot et son élévation entière sont deux représentations mutuellement récupérables du même bloc émis.

\begin{theorem}[Stabilité d'extraction à toute profondeur]
\label{act21:thm:df-bloque-estable}
On vérifie, pour tout \(n\),
\[
 P_c(729,n)=P_c(3,6n).
\]
Si \(t<n\), alors
\begin{equation}
 u_c(t)=
 \left\lfloor\frac{P_c(729,n)}{729^{n-t-1}}\right\rfloor\bmod729.
 \label{act21:eq:df-bloque-prefijo-largo}
\end{equation}
En conséquence, toute extension du préfixe préserve tous les
blocs précédemment émis et toutes les signatures calculées dessus.
\end{theorem}

\begin{proof}
La première égalité résulte de ce que les deux préfixes se séparent à la même
échelle, \(729^n=3^{6n}\), et du
théorème~\ref{act21:thm:df-emision-correcta}. Pour la seconde, nous appliquons
\eqref{act21:eq:df-truncamiento} aux profondeurs \(n\) et \(t+1\), puis
nous prenons le résidu modulo \(729\). L'extraction ternaire de
\eqref{act21:eq:df-bloque-y-banda} est déterministe. Toute fonction explicite
de ces dix-huit coordonnées conserve, par conséquent, le même résultat en
étendant le préfixe.
\end{proof}

La construction est quantifiée sur \(t\in\mathbb N\). Une table qui
matérialise cent blocs constitue une évaluation finie de cette
construction ; la formule \eqref{act21:eq:df-bloque-prefijo-largo} détermine
exactement sa compatibilité avec toute extension ultérieure.

\begin{example}[Premier bloc conjoint]
L'évaluation des lecteurs régionaux sélectionnés donne les mots
\[
 B_{\mathrm P}(0)=010211,\qquad
 B_{\mathrm E}(0)=201101,\qquad
 B_{\mathrm F}(0)=121200.
\]
Leur élévation par \eqref{act21:eq:df-elevacion-banda} produit
\[
 u_{\mathrm P}(0)=103,\qquad
 u_{\mathrm E}(0)=523,\qquad
 u_{\mathrm F}(0)=450.
\]
Par exemple,
\(201101_3=2\cdot243+27+9+1=523\).
Les six positions, y compris les zéros initiaux, sont conservées dans la
représentation en bande. Le zéro initial de \(010211\) possède une position
définie, bien que l'écriture décimale de \(103\) en soit dépourvue.
\end{example}

\subsection{Registres de transition et reconstruction des élévations linéaires}
\label{act21:subsec:df-registros-transicion}

% Fuentes: deltas/integracion/RegionalTransitionRecords.lean y
% GeneratedTransitionRecords.lean; lean/biblioteca/TPKLifts.lean.
% Los enlaces reconstruidos son R(0)->R(1) y R(2)->R(3).

Sur \(\mathbb F_3\), nous réunissons deux émissions consécutives de chaque région dans la matrice
\begin{equation}
 R(t)=
 \begin{pmatrix}
 B_{\mathrm P}(t)\\ B_{\mathrm P}(t+1)\\
 B_{\mathrm E}(t)\\ B_{\mathrm E}(t+1)\\
 B_{\mathrm F}(t)\\ B_{\mathrm F}(t+1)
 \end{pmatrix}.
 \label{act21:eq:df-registro-matricial}
\end{equation}
Cette définition fixe à la fois l'ordre régional et l'ordre temporel. Les
quatre premiers registres sont
\[
 X_0=R(0),\quad Y_0=R(1),\quad X_1=R(2),\quad Y_1=R(3).
\]
Ses valeurs, obtenues par extraction des préfixes régionaux, sont
\begin{align*}
 X_0&=\begin{pmatrix}
0&1&0&2&1&1\\0&1&2&2&2&2\\2&0&1&1&0&1\\
1&2&1&2&2&1\\1&2&1&2&0&0\\1&1&2&2&0&2
\end{pmatrix},&
 Y_0&=\begin{pmatrix}
0&1&2&2&2&2\\0&1&0&2&1&1\\1&2&1&2&2&1\\
1&0&2&0&1&1\\1&1&2&2&0&2\\1&2&1&0&2&0
\end{pmatrix},\\[1ex]
 X_1&=\begin{pmatrix}
0&1&0&2&1&1\\0&0&2&1&1&1\\1&0&2&0&1&1\\
0&1&2&2&2&2\\1&2&1&0&2&0\\0&1&0&2&1&0
\end{pmatrix},&
 Y_1&=\begin{pmatrix}
0&0&2&1&1&1\\1&1&0&2&2&1\\0&1&2&2&2&2\\
1&0&2&0&1&1\\0&1&0&2&1&0\\0&1&0&2&0&0
\end{pmatrix}.
\end{align*}

\begin{theorem}[Reconstruction unique des deux liens initiaux]
\label{act21:thm:df-elevaciones-reconstruidas}
Les équations \(X_0L_0=Y_0\) et \(X_1L_1=Y_1\) ont, chacune, une
unique solution dans \(M_6(\mathbb F_3)\). Ces solutions sont
\begin{align*}
 L_0&=\begin{pmatrix}
2&2&2&1&2&1\\2&1&2&2&1&1\\1&0&0&1&0&2\\
0&0&1&1&1&0\\2&2&2&0&2&1\\2&1&2&1&0&0
\end{pmatrix},&
 L_1&=\begin{pmatrix}
2&1&2&0&2&2\\1&2&1&1&2&0\\1&2&1&1&1&2\\
0&1&0&0&2&1\\2&0&2&1&0&1\\0&2&2&2&1&1
\end{pmatrix}.
\end{align*}
Les deux matrices sont inversibles.
\end{theorem}

\begin{proof}
Les matrices
\begin{align*}
 A_0&=\begin{pmatrix}
1&1&2&0&1&2\\0&0&1&0&0&1\\2&1&2&1&0&1\\
0&2&0&1&0&2\\2&1&1&0&2&2\\2&1&1&1&1&2
\end{pmatrix},&
 A_1&=\begin{pmatrix}
1&0&1&2&0&0\\0&0&1&1&2&2\\0&1&2&0&1&1\\
1&1&0&2&0&0\\1&1&2&1&1&2\\1&0&0&0&0&2
\end{pmatrix}
\end{align*}
satisfont \(A_iX_i=X_iA_i=I\), par multiplication dans \(\mathbb F_3\).
Par conséquent, les solutions sont imposées par \(L_i=A_iY_i\) ; la
multiplication produit les deux matrices montrées. Leurs inverses sont
\begin{align*}
 L_0^{-1}&=\begin{pmatrix}
1&2&0&2&0&2\\2&0&2&2&0&0\\2&1&2&0&2&0\\
1&0&0&0&2&0\\0&2&1&1&2&0\\2&2&2&2&2&2
\end{pmatrix},&
 L_1^{-1}&=\begin{pmatrix}
2&0&2&1&2&1\\2&1&0&0&2&0\\2&0&2&2&0&2\\
2&0&0&1&1&0\\1&1&1&1&1&0\\2&0&1&2&2&0
\end{pmatrix}.
\end{align*}
La vérification bilatérale \(L_iL_i^{-1}=L_i^{-1}L_i=I\) démontre la
récupération exacte de chaque mot transformé. Tous les produits
utilisent des résidus \(0,1,2\), avec leur interprétation en \(\mathbb F_3\).
\end{proof}

La lecture causale de ce résultat est précise : les émissions régionales produisent \(R(t)\) ; les registres \(R(0),\ldots,R(3)\) déterminent les deux élévations précédentes ; leurs inverses récupèrent les mots transformés.  
La continuité pour tout \(t\) appartient aux lecteurs coinductifs et au théorème ~\ref{act21:thm:df-bloque-estable}. La reconstruction matricielle démontrée ici a pour domaine les deux liens spécifiés dans son énoncé.

\subsection{Raffinement simultané dans les bases ternaire et décimale}
\label{act21:subsec:df-cilindros-mixtos}

% Fuente: deltas/integracion/RegionalCylinderTransition.lean.

Une émission ternaire de six trits et une émission décimale de trois chiffres
constituent deux lectures de bases différentes. Pour conserver leur relation, on
enregistre l'état
\[
 s=(n,k,N,D),
\]
où \(N\) est un préfixe de profondeur \(n\) en base \(729\), et \(D\) est
un préfixe de profondeur \(k\) en base \(1000\). Leurs cylindres
archimédiens sont
\begin{equation}
 I_{729}(s)=\left[\frac{N}{729^n},\frac{N+1}{729^n}\right),\qquad
 I_{1000}(s)=\left[\frac{D}{1000^k},\frac{D+1}{1000^k}\right).
 \label{act21:eq:df-cilindros}
\end{equation}
L'extrémité supérieure semi-ouverte attribue chaque valeur à une unique cellule d'une base et d'une profondeur données.

\begin{definition}[Défauts entiers de compatibilité]
On définit
\begin{align}
 A(s)&=N1000^k-D729^n,\label{act21:eq:df-defecto-inferior}\\
 B(s)&=(N+1)1000^k-D729^n=A(s)+1000^k.
 \label{act21:eq:df-defecto-superior}
\end{align}
L'état est compatible lorsque
\begin{equation}
 A(s)<729^n,\qquad B(s)>0.
 \label{act21:eq:df-compatibilidad}
\end{equation}
\end{definition}

\begin{lemma}[Critère entier d'intersection]
Les cylindres de \eqref{act21:eq:df-cilindros} ont une intersection non vide si et
seulement si \eqref{act21:eq:df-compatibilidad} est vérifiée.
\end{lemma}

\begin{proof}
Deux intervalles semi-ouverts \([a,b)\) et \([c,d)\), avec \(a<b\) et
\(c<d\), se coupent exactement lorsque \(a<d\) et \(c<b\). En appliquant cette
condition à \eqref{act21:eq:df-cilindros} et en multipliant par les dénominateurs
positifs, nous obtenons
\(N1000^k<(D+1)729^n\) et
\(D729^n<(N+1)1000^k\). Ce sont les deux inégalités de
\eqref{act21:eq:df-compatibilidad}.
\end{proof}

\begin{definition}[Mise à jour des profondeurs et des préfixes]
Pour \(0\leq u<729\) et \(0\leq d<1000\), nous définissons
\begin{align*}
 T_u(n,k,N,D)&=(n+1,k,729N+u,D),\\
 T_{u,d}(n,k,N,D)&=(n+1,k+1,729N+u,1000D+d).
\end{align*}
La première opération raffine uniquement le cylindre ternaire. La seconde  
raffine les deux lectures et conserve les deux profondeurs dans l'état.
\end{definition}

\begin{proposition}[Récurrences exactes de frontière]
\label{act21:prop:df-rec-frontera}
Les mises à jour précédentes satisfont
\begin{align}
 A(T_u s)&=729A(s)+u1000^k,\label{act21:eq:df-rec-ternaria}\\
 A(T_{u,d}s)&=729000A(s)+u1000^{k+1}-d729^{n+1}.
 \label{act21:eq:df-rec-conjunta}
\end{align}
\end{proposition}

\begin{proof}
Pour la première identité,
\[
 (729N+u)1000^k-D729^{n+1}
 =729(N1000^k-D729^n)+u1000^k.
\]
Pour la seconde,
\begin{align*}
 &(729N+u)1000^{k+1}-(1000D+d)729^{n+1}\\
 &\qquad=729000(N1000^k-D729^n)
      +u1000^{k+1}-d729^{n+1}.
\end{align*}
Les deux sont des identités entières antérieures à toute approximation réelle.
\end{proof}

\begin{theorem}[Raffinement des cylindres régionaux générés]
\label{act21:thm:df-refinamiento-generado}
Soit
\[
 s_c(n,k)=(n,k,P_c(729,n),P_c(1000,k)).
\]
Alors \(s_c(n,k)\) est compatible quels que soient \(n,k\). De plus,
\begin{align*}
 T_{u_c(n)}s_c(n,k)&=s_c(n+1,k),\\
 T_{u_c(n),d_c(1000,k)}s_c(n,k)&=s_c(n+1,k+1).
\end{align*}
La division euclidienne des préfixes raffinés récupère exactement les
préfixes précédents. Pour tout \(a,b\geq0\),
\[
 \frac{P_c(729,n+a)-P_c(729,n+a)\bmod729^a}{729^a}
 =P_c(729,n),
\]
et l'identité analogue est vérifiée en base \(1000\).
\end{theorem}

\begin{proof}
Le théorème~\ref{act21:thm:df-emision-correcta} place \(v_c\) simultanément dans
les deux intervalles \eqref{act21:eq:df-cilindros} ; le critère d'intersection donne
la compatibilité. Les identités de mise à jour sont
\eqref{act21:eq:df-paso-posicional} dans chaque base. La récupération s'obtient à partir de
\eqref{act21:eq:df-truncamiento}. Par conséquent, la récurrence du défaut et la
récupération du préfixe sont des conséquences d'une même mise à jour.
\end{proof}

\begin{example}[Premier couple de cylindres régionaux]
Pour \(n=k=1\), l'évaluation des trois régions donne
\[
\begin{array}{c|rr|rr}
 c&N&D&A&B\\\hline
 \mathrm P&2290&3141&211&1211\\
 \mathrm E&1981&2718&-422&578\\
 \mathrm F&1179&1618&-522&478
\end{array}
\]
Dans les trois lignes \(A<729\) et \(B>0\). La première colonne de préfixes
se décompose comme \(2290=3\cdot729+103\),
\(1981=2\cdot729+523\) et \(1179=729+450\) : réapparaissent exactement les
blocs de six trits déjà émis. Les signes différents de \(A\) expriment
la position relative des deux extrémités inférieures, sans altérer la
appartenance commune de la valeur régionale.
\end{example}

\subsection{Sélection du bloc suivant par le lecteur régional}

La compatibilité de deux intervalles exprime une relation entre lectures. La détermination du bloc suivant utilise en outre l'intervalle rationnel produit par la région. Cette opération peut se formuler sans éliminer l'histoire de profondeur.

\begin{definition}[Admissibilité de la prolongation régionale]
\(c,n\) étant fixés, soient \(s=729^{n+1}\) et \(j=j_c(s)\). Un résidu
\(u\in\{0,\ldots,728\}\) est admissible lorsque
\begin{equation}
 \lfloor s\ell_c(j)\rfloor
 \leq729P_c(729,n)+u
 \leq\lceil su_c(j)\rceil-1.
 \label{act21:eq:df-admisibilidad}
\end{equation}
\end{definition}

\begin{theorem}[Unicité de la prolongation compatible avec le lecteur]
\label{act21:thm:df-prolongacion-unica}
La condition \eqref{act21:eq:df-admisibilidad} est équivalente à \(u=u_c(n)\).
Par conséquent, il existe exactement une prolongation admissible de chaque région
à chaque profondeur.
\end{theorem}

\begin{proof}
À l'indice de séparation, les deux extrémités entières de
\eqref{act21:eq:df-admisibilidad} coïncident avec \(P_c(729,n+1)\).
La condition est ainsi équivalente à
\(729P_c(729,n)+u=P_c(729,n+1)\).
Par \eqref{act21:eq:df-paso-posicional}, le côté droit est
\(729P_c(729,n)+u_c(n)\). L'annulation du terme commun donne l'égalité des résidus. Réciproquement, cette égalité satisfait la condition.
\end{proof}

La dépendance effective de cette sélection reste complètement déterminée :
région sélectionnée, lecteur rationnel, profondeur, indice de séparation,
préfixe antérieur et résidu suivant. Le défaut \(A\) résume une relation
de frontière ; l'état \((n,k,N,D)\) et les intervalles régionaux conservent
les données employées par la mise à jour.

\subsection{Signature ternaire, retenue, incidence de Witt et horloge de lecture}
\label{act21:subsec:df-firma-conjunta}

% Fuentes: lean/regional/N69RegionalSignature.lean,
% deltas/integracion/JointRegionalFrontier.lean y JointCylinderSignature.lean.

Une profondeur \(t\) étant fixée, écrivons \(B_{c,j}=B_{c,j}(t)\). Pour chaque colonne \(j\), on calcule les entiers et les résidus suivants :
\begin{align}
 S_j&=B_{\mathrm P,j}+B_{\mathrm E,j}+B_{\mathrm F,j},
 &q_j&=S_j\bmod3,
 &c_j&=\left\lfloor S_j/3\right\rfloor,
 \label{act21:eq:df-firma-qc}\\
 a_j&=(B_{\mathrm P,j}+B_{\mathrm E,j}-B_{\mathrm F,j})\bmod3,
 &w_j&=\sum_c\mathbf1_{B_{c,j}\ne0},
 &\bar w_j&=w_j\bmod3.
 \label{act21:eq:df-firma-axial}
\end{align}
Le résidu de l'expression axiale est pris dans les entiers. Une implémentation
avec des naturels le calcule comme
\((B_{\mathrm P,j}+B_{\mathrm E,j}+3-B_{\mathrm F,j})\bmod3\), qui le conserve
exactement.

La matrice d'incidence ternaire utilisée par le lecteur est
\begin{equation}
 W=\begin{pmatrix}
0&1&1&1&1&1\\
1&0&1&1&2&2\\
1&1&0&2&1&2\\
1&1&2&0&2&1\\
1&2&1&2&0&1\\
1&2&2&1&1&0
\end{pmatrix}.
\label{act21:eq:df-matriz-witt}
\end{equation}
Pour chaque région, deux charges résiduelles sont conservées :
\begin{equation}
 v_c^{\mathrm{res}}=\sum_j B_{c,j}\bmod3,
 \qquad
 v_c^{\mathrm W}=\sum_j(B_cW)_j\bmod3.
 \label{act21:eq:df-cargas-residuales}
\end{equation}
La signature complète du lecteur est
\(\mathcal S(B)=(q,a,c,w,\bar w,v^{\mathrm{res}},v^{\mathrm W})\).
Sa structure conserve séparément le reste de colonne, la retenue, le
contraste axial, le poids de support et les deux lectures de charge.

\begin{theorem}[Récupération exacte et bornes de la signature]
\label{act21:thm:df-recuperacion-firma}
Pour chaque colonne et chaque région, on a
\begin{align}
 S_j&=q_j+3c_j,&0\leq q_j,a_j&<3,&0\leq c_j&\leq2,
 \label{act21:eq:df-recuperacion-total}\\
 B_{\mathrm F,j}&=2(q_j+3-a_j)\bmod3,&0\leq w_j&\leq3,
 &0\leq\bar w_j&<3.
 \label{act21:eq:df-recuperacion-eje}
\end{align}
De plus,
\begin{equation}
 v_c^{\mathrm W}
 =\left(\sum_j\bigl((B_cW)_j\bmod3\bigr)\right)\bmod3.
 \label{act21:eq:df-witt-reduccion}
\end{equation}
Les récupérations énoncées déterminent le total de chaque colonne et sa
coordonnée d'auto-échelle ; le registre régional conserve en outre
les trois bandes ordonnées.
\end{theorem}

\begin{proof}
L'identité \eqref{act21:eq:df-recuperacion-total} est la division euclidienne de
\(S_j\) par \(3\). Chaque terme de \(S_j\) est entre \(0\) et \(2\),
de telle sorte que \(0\leq S_j\leq6\) et \(c_j\leq2\). En \(\mathbb F_3\),
la soustraction des deux congruences qui définissent \(q_j\) et \(a_j\) donne
\(q_j-a_j=2B_{\mathrm F,j}\). Comme l'inverse de \(2\) est \(2\), on
obtient \eqref{act21:eq:df-recuperacion-eje} ; le terme \(3\) permet
d'écrire celle-ci avec des nombres naturels. Les bornes de \(w_j\) résultent de sommer trois
indicateurs. Enfin, réduire une somme entière modulo \(3\) équivaut
à réduire chaque terme et à réduire à nouveau la somme, ce qui prouve
\eqref{act21:eq:df-witt-reduccion}.
\end{proof}

\begin{example}[Signature du premier bloc conjoint]
Les trois bandes initiales produisent
\begin{align*}
 q&=(0,0,2,2,1,2),&a&=(1,2,0,1,1,2),\\
 c&=(1,1,0,1,0,0),&w&=(2,2,2,3,1,2),\\
 v^{\mathrm{res}}&=(2,2,0),&v^{\mathrm W}&=(2,1,1).
\end{align*}
La quatrième colonne a un total \(2+1+2=5\) ; sa signature enregistre \(q_3=2\) et \(c_3=1\), par conséquent \(q_3+3c_3=5\).
La récupération axiale donne \(2(2+3-1)\bmod3=2\), précisément la coordonnée d'auto-échelle.
Les images de Witt, réduites composante par composante, sont
\[
 (2,0,2,1,1,2),\qquad
 (0,0,0,2,0,2),\qquad
 (2,1,1,2,1,0).
\]
Sa somme résiduelle par ligne produit les trois charges \((2,1,1)\).
\end{example}

\begin{definition}[Horloge entière de lecture décimale]
L'indice de transition \(t\) et la profondeur décimale de représentation sont
reliés par
\begin{equation}
 \kappa(t)=\max\{k\in\mathbb N:1000^k\leq3^{6(t+1)}\}.
 \label{act21:eq:df-reloj}
\end{equation}
\end{definition}

\begin{proposition}[Caractérisation et monotonie de l'horloge]
\label{act21:prop:df-reloj}
L'entier \(\kappa(t)\) est déterminé de manière univoque par
\begin{equation}
 1000^{\kappa(t)}\leq729^{t+1}<1000^{\kappa(t)+1}.
 \label{act21:eq:df-reloj-cotas}
\end{equation}
La fonction \(\kappa\) est monotone et satisfait
\(\kappa(20)=\kappa(21)=20\).
\end{proposition}

\begin{proof}
Les puissances de \(1000\) sont strictement croissantes et non bornées ;
il existe donc un unique intervalle de puissances consécutives qui
contient \(729^{t+1}\). La monotonie de cette dernière expression démontre
celle de \(\kappa\). Les deux évaluations finales sont les comparaisons
entières
\[
 1000^{20}\leq729^{21}<1000^{21},\qquad
 1000^{20}\leq729^{22}<1000^{21}.
\]
Ces comparaisons s'effectuent avec des puissances entières exactes.
\end{proof}

L'égalité de profondeur décimale en deux transitions distinctes explique la nécessité de conserver les deux indices : l'horloge de lecture peut rester constante tandis que la bande ternaire avance et que sa signature est mise à jour. L'information chronologique se trouve dans le registre ordonné des transitions.

\subsection{Théorème de composition du bloc, du cylindre et de la signature}

\begin{theorem}[Compatibilité conjointe des lectures régionales]
\label{act21:thm:df-composicion-conjunta}
Pour chaque profondeur \(n\), il existe un unique triplet de blocs
\((u_{\mathrm P},u_{\mathrm E},u_{\mathrm F})\) qui satisfait
l'admissibilité régionale \eqref{act21:eq:df-admisibilidad} dans les trois régions.
Ce triplet est \((u_{\mathrm P}(n),u_{\mathrm E}(n),u_{\mathrm F}(n))\).
Simultanément :
\begin{enumerate}
\item chaque bloc met à jour son cylindre par les récurrences
\eqref{act21:eq:df-rec-ternaria} et \eqref{act21:eq:df-rec-conjunta} ;
\item ses six trits produisent la signature
\(\mathcal S(B_{\mathrm P}(n),B_{\mathrm E}(n),B_{\mathrm F}(n))\);
\item les totaux de colonne se retrouvent comme \(q+3c\), et les charges de Witt sont obtenues à partir des mêmes bandes par
\eqref{act21:eq:df-witt-reduccion} ;
\item tout préfixe de plus grande profondeur renvoie exactement ces
blocs, cylindres tronqués et signatures.
\end{enumerate}
\end{theorem}

\begin{proof}
Nous appliquons le théorème ~\ref{act21:thm:df-prolongacion-unica} à chacun des trois
indices régionaux. L'unicité coordonnée par coordonnée donne l'unicité
conjointe. Nous substituons ces mêmes résidus dans les opérations
\(T_u,T_{u,d}\), ce qui permet d'appliquer le
théorème ~\ref{act21:thm:df-refinamiento-generado}. L'application déterministe de
\eqref{act21:eq:df-bloque-y-banda} produit les bandes utilisées par
\eqref{act21:eq:df-firma-qc}--\eqref{act21:eq:df-cargas-residuales}. Les identités
de récupération sont celles du théorème ~\ref{act21:thm:df-recuperacion-firma}.
Finalement, le théorème ~\ref{act21:thm:df-bloque-estable} et le troncage de
chaque base démontrent la stabilité à plus grande profondeur. À chaque étape, on a conservé l'identité du bloc émis ; aucune des deux
lectures ne requiert de sélectionner un autre bloc.
\end{proof}

La conclusion est une compatibilité opérationnelle entre la contraction
cylindrique et la conservation de l'incidence. Les deux utilisent le même
registre généré, maintiennent les profondeurs et admettent la récupération
entière. Ce résultat est la forme explicite du lien
\[
 \text{bloc régional émis}
 \longrightarrow
 \begin{cases}
 \text{cylindre et frontière entière},\\
 \text{bande, signature, retenue et charge de Witt}.
 \end{cases}
\]
Le développement ultérieur de l'incidence et des champs emploie la deuxième lecture, conservant l'origine commune avec la première.

\subsection{Recensement des calendriers et compatibilité simultanée des biographies régionales}
Les lignes précédentes se regroupent, sans appliquer encore aucun lecteur décimal,
dans les trois biographies de blocs produites par les transitions :
\begin{align}
 B^{\rm cl}
  &=010211|012222|010211|002111|110221,\notag\\
 B^{\rm pr}
  &=201101|121221|102011|012222|102011,
 \label{act21:eq:biografias-tpk-tres-caracteres}\\
 B^{\rm au}
  &=121200|112202|121020|010210|010200.\notag
\end{align}
En effet, les paires consécutives des deux premiers tronçons sont les lignes de \(X_0,Y_0\), et celles des deux derniers sont les lignes de \(X_1,Y_1\). Ainsi, les biographies sont une autre écriture du registre de douze transitions, non une collection de préfixes conventionnels.

Pour vérifier le mot opératoire du calendrier, soit
\(c=c_1c_2c_3c_4\in\{0,1\}^4\) et, pour chaque régime
\(\chi\in\{{\rm cl},{\rm pr},{\rm au}\}\), définissons
\[
 u_0^\chi=b_0^\chi,\qquad
 u_j^\chi=u_{j-1}^\chi L_{c_j}\quad(1\le j\le4).
\]
Pour condenser le recensement, notons par
\[
 N_{\mathrm{ar}}(c):=
 \#\{(j,\chi):u_j^\chi=b_j^\chi,\ 1\le j\le4\},\qquad
 N_{\mathrm{bio}}(c):=
 \#\{\chi:(u_0^\chi,\ldots,u_4^\chi)=B^\chi\}.
\]
La multiplication exacte dans \(\FF_3\) des seize calendriers donne :
\[
\begin{array}{c|c|c}
 c&N_{\mathrm{ar}}(c)&N_{\mathrm{bio}}(c)\\ \hline
0000,0001&6&0\\
0010&9&0\\
0011&12&3\\
0100,0101,0110,0111&3&0\\
1000,1001,1010,1011&0&0\\
1100,1101,1110,1111&0&0
\end{array}
\tag{D}
\label{act21:eq:censo-dieciseis-calendarios}
\]
Le tableau épuise les \(2^4=16\) possibilités et chacune de ses entrées s'obtient
par quatre multiplications d'une ligne par l'une des deux matrices reproduites.
Par conséquent, \(0011\) est l'unique calendrier qui reconstruit
simultanément les trois biographies complètes. En notation opératorielle,
\begin{equation}
 (L_0,L_0,L_1,L_1)
 \label{act21:eq:calendario-0011-interno}
\end{equation}
ce n'est pas une prescription indépendante : c'est la seule lecture binaire compatible
avec les douze arêtes déjà produites par \(U_t\).

Le calendrier \eqref{act21:eq:calendario-0011-interno} produit quatre nouveaux
blocs de six composants. Leurs concaténations forment
\[
 w_6\longrightarrow w_{12}\longrightarrow w_{18}
 \longrightarrow w_{24}\longrightarrow w_{30}.
\]
Les sous-indices de \(w\) indiquent la longueur accumulée. La frontière \(R_{36}\) conserve la structure terminale et ouvre la relation de prolongement ; elle n'est pas renommée comme un dernier bloc périodique qui pourrait se répéter indéfiniment.

L'identité \(L=X^{-1}Y\) démontre l'unicité une fois les
transitions fixées ; elle ne remplace pas leur production par \(U_t\). Réciproquement,
\eqref{act21:eq:censo-dieciseis-calendarios} prouve dans l'article l'unicité
de l'ordre des quatre transports. La chaîne causale est ainsi close :
le recensement APP produit les régions initiales, le TRIT fixe le régime et
l'orientation, le TPK produit le registre des transitions, et ce registre
détermine \(L_0,L_1\) et \(0011\) avant toute réalisation archimédienne.



% Desarrollo reunido K31-07--K31-15. Inserción anterior a registro_k y alpha.
% Mantiene explícita la interfaz del repertorio terminal de ocho bloques.
% Propietarios y localizadores: metadata/K_ALPHA_PROPIETARIOS.md.
\section{Détermination régionale du descripteur dodécaphasique et sélection du registre de couplage}
\label{act21:chap:despliegue-selector-k}

La détermination du registre de couplage articule deux constructions
qu'il convient de déployer dans leur ordre effectif. La première produit les
antécédents du sélecteur à partir des représentations régionales du TPK :
bandes de six trits, signatures de colonne, calendrier de conversion,
descripteur orienté et tableau fonctionnel. La seconde opère sur ces
antécédents et sur le répertoire terminal de huit blocs, énumère leurs
ordonnancements admissibles et sélectionne un unique registre au moyen de deux
conditions simultanées : incidence exceptionnelle et hétérotypie de lecture.
La valeur du registre apparaît au terme de cette composition.

Le présent développement présuppose les représentations régionales obtenues
dans les chapitres précédents à partir d'APP--TRIT--TPK. Son propos est de rendre
explicite chaque opération intermédiaire qui conduit de ces représentations
au sélecteur fini. Sont ensuite incorporées la réalisation du même
registre par incidences signées, l'inversion de Hadamard et la
représentation de la constante de structure fine. Les trois constructions
conservent leurs domaines respectifs : sélection du registre, récupération
à partir de ses canaux et évaluation arithmétique de ses représentations.

\subsection{Bandes régionales, codification et signatures entières}

\subsubsection{Domaines et codification des bandes régionales}

On fixe l'ordre des canaux
\[
 \chi\in\{\pi,e,\varphi\},
 \qquad B_\chi(t)\in\mathbb F_3^6,
 \qquad t\in\mathbb N.
\]
Les symboles de canal identifient les régions déjà engendrées. Chaque
$B_\chi(t)$ conserve la place qu'il occupe dans sa suite de représentations ;
l'indice $t$ compte les blocs ternaires. L'évaluation naturelle d'une
bande $b=(b_0,\ldots,b_5)$ est
\begin{equation}
 \operatorname{enc}_6(b)=
 243b_0+81b_1+27b_2+9b_3+3b_4+b_5,
 \qquad 0\leq\operatorname{enc}_6(b)<729.
 \label{act21:eq:kdes-enc6}
\end{equation}
Son inverse, avec une largeur fixée, est
\[
 \operatorname{dig}_j(n)=
 \left\lfloor\frac{n}{3^{5-j}}\right\rfloor\bmod3,
 \qquad 0\leq j<6.
\]
Ainsi sont retenus les zéros initiaux : les mots $002001$ et $2001$
pourraient partager une écriture entière abrégée, mais le type du
premier contient six positions. Toutes les concaténations qui suivent
utilisent cette largeur.

Dans la fenêtre de cent blocs, la représentation de six cents trits du
canal $\chi$ détermine un entier $P_\chi^{(600)}$. L'extraction s'effectue
par
\begin{equation}
 b_\chi(t)=
 \left\lfloor\frac{P_\chi^{(600)}}{729^{99-t}}\right\rfloor\bmod729,
 \qquad 0\leq t<100,
 \qquad B_\chi(t)=\operatorname{enc}_6^{-1}(b_\chi(t)).
 \label{act21:eq:kdes-extraccion}
\end{equation}
L'entier de six cents trits est une représentation du producteur régional ;
la table de cent lignes est l'évaluation de
\eqref{act21:eq:kdes-extraccion}. Cette séparation permet de remplacer une table
historique par son producteur sans modifier le sélecteur qui l'utilise.

\subsubsection{Résidus, quotients, occupation et orientation}

Pour chaque colonne $j$ et temps $t$, on introduit
\begin{align}
 T_j(t)&=B_\pi(t)_j+B_e(t)_j+B_\varphi(t)_j,\\
 q_j(t)&=T_j(t)\bmod3,
 &c_j(t)&=\left\lfloor T_j(t)/3\right\rfloor,\\
 a_j(t)&=\bigl(B_\pi(t)_j+B_e(t)_j-B_\varphi(t)_j\bigr)\bmod3,\\
 w_j(t)&=\mathbf1_{B_\pi(t)_j\ne0}
       +\mathbf1_{B_e(t)_j\ne0}
       +\mathbf1_{B_\varphi(t)_j\ne0},
 &\bar w_j(t)&=w_j(t)\bmod3.
 \label{act21:eq:kdes-firmas}
\end{align}
La soustraction de la troisième expression s'effectue sur les entiers et se réduit
ensuite. Dans une implémentation sur les naturels, l'expression équivalente
est $(B_\pi+B_e+3-B_\varphi)\bmod3$ ; la soustraction tronquée altérerait le lecteur.

\begin{proposition}[Reconstruction de la somme de colonne]
Pour toutes les bandes admissibles,
\[
 T_j(t)=q_j(t)+3c_j(t),\qquad
 0\leq q_j(t)<3,\quad 0\leq c_j(t)\leq2,
 \quad0\leq w_j(t)\leq3.
\]
\end{proposition}
\begin{proof}
Chaque colonne additionne trois représentants de $\{0,1,2\}$, de sorte que
$0\leq T_j(t)\leq6$. La division euclidienne par trois fournit le reste
et le quotient indiqués. L'occupation compte trois conditions binaires ;
sa borne se déduit directement de cette définition.
\end{proof}

Cette identité conserve le quotient de la somme locale. La reconstruction
des trois bandes individuelles requiert les bandes elles-mêmes ou un lecteur
supplémentaire qui les distingue : une somme de colonne a moins de coordonnées
que le triplet qui la produit. Par conséquent, la ligne temporelle utilisée plus
tard conserve conjointement les trois bandes et les quatre champs
$q,a,c,\bar w$.

\subsection{Horloge de conversion et premier avancement nul}

\subsubsection{Capacité ternaire et profondeur décimale}

Une représentation de $t+1$ bandes contient $6(t+1)$ trits. Le calendrier
de conversion associé est défini par l'entier
\begin{equation}
 d_t=\max\{k\in\mathbb N:1000^k\leq3^{6(t+1)}\}
     =\max\{k\in\mathbb N:1000^k\leq729^{t+1}\}.
 \label{act21:eq:kdes-reloj}
\end{equation}
La définition compare des capacités entières. Le compteur $t$ et la profondeur $d_t$ enregistrent des grandeurs distinctes : une transition ajoute une bande de six trits ; l'horloge compte les puissances complètes de la base décimale mille contenues dans la capacité ternaire accumulée.

\begin{lemma}[Caractérisation entière de l'horloge]
Pour $t,k\in\mathbb N$ on a
\[
 d_t=k\quad\Longleftrightarrow\quad
 1000^k\leq729^{t+1}<1000^{k+1}.
\]
De plus, $d_t\leq d_{t+1}$.
\end{lemma}
\begin{proof}
Les puissances de mille forment une suite strictement croissante et non bornée. La puissance $729^{t+1}$ se situe entre deux termes consécutifs ; l'exposant inférieur est exactement le maximum de \eqref{act21:eq:kdes-reloj}. L'inégalité $729^{t+1}<729^{t+2}$ implique la monotonie du maximum.
\end{proof}

\begin{proposition}[Premier avancement nul]
Le plus petit temps qui satisfait $d_{t+1}=d_t$ est
\[
 t_*=20.
\]
En particulier,
\[
 d_0=0,\quad d_{19}=19,\quad d_{20}=20,
 \quad d_{21}=20,\quad d_{22}=21.
\]
\end{proposition}
\begin{proof}
Les comparaisons entières exactes
\[
 1000^{20}\leq729^{21},\qquad
 729^{22}<1000^{21},\qquad
 1000^{21}\leq729^{23}<1000^{22}
\]
donnent les valeurs aux temps vingt, vingt-et-un et vingt-deux. Pour
$0\leq t\leq20$, l'expression
$729^{t+1}/1000^t=729(729/1000)^t$ est décroissante et sa valeur en vingt
est supérieure ou égale à un. Par conséquent,
$1000^t\leq729^{t+1}<1000^{t+1}$ et $d_t=t$. Aucun temps inférieur à
vingt ne peut satisfaire $d_{t+1}=d_t$ ; en vingt, c'est le cas.
Toutes les comparaisons utilisées sont de puissances entières.
\end{proof}

Le nombre vingt est donc une sortie du critère minimal. La règle qui sélectionne l'instant est ``première avance nulle de l'horloge'' ; l'algorithme et sa preuve peuvent utiliser cette règle sans recevoir vingt comme paramètre externe.

\subsubsection{Retour nonadique et antécédent temporel}

La longueur du retour de phase détermine
\begin{equation}
 r=t_*-9=11.
 \label{act21:eq:kdes-retorno}
\end{equation}
Les requêtes du descripteur ont lieu en $t_*+1$, $r+1$ et $r$.
Cette distribution conserve la relation entre l'instant de pause, la
translation nonadique et la lecture de phase immédiatement associée. Les
temps onze, douze et vingt-et-un désignent des positions du calendrier
régional ; leur différence temporelle reste enregistrée même lorsque deux
temps produisent la même profondeur décimale.

\subsection{Construction orientée du descripteur de vingt-quatre trits}

\subsubsection{Défauts opposés de semi-bande}

On considère les défauts orientés de la coordinatisation régionale
\[
 \delta_+=(0,0,0,1,1,1),\qquad
 \delta_-=(0,0,0,2,2,2)\in\mathbb F_3^6.
\]
Son opposition se vérifie composante par composante :
\[
 \delta_++\delta_-=0\quad\text{en }\mathbb F_3^6.
\]
Le représentant $2$ réalise $-1$ dans la coordinatisation ternaire. La concaténation
des queues de trois positions, dans l'ordre négatif--positif, produit
le mot $222111$. La somme des défauts et la concaténation de leurs
queues sont des opérations distinctes : la première a pour codomaine
$\mathbb F_3^6$, tandis que la seconde conserve l'ordre de deux
segments de longueur trois.

\subsubsection{Charge duale et sélection de phase}

Dans la coordinatisation régionale, la transformation de six coordonnées utilise
\[
 A_{\rm reg}=
 \begin{pmatrix}
 0&1&1&1&1&1\\
 1&0&1&1&2&2\\
 1&1&0&2&1&2\\
 1&1&2&0&2&1\\
 1&2&1&2&0&1\\
 1&2&2&1&1&0
 \end{pmatrix}.
\]
Avec des vecteurs ligne, la charge duale de $b\in\mathbb F_3^6$ est la somme
des coordonnées de $bA_{\rm reg}$. Les sommes entières des lignes de
$A_{\rm reg}$ sont $(5,7,7,7,7,7)$ ; d'où résulte l'expression locale
\begin{equation}
 \eta^\vee(b)=2b_0+b_1+b_2+b_3+b_4+b_5\pmod3.
 \label{act21:eq:kdes-cargadual}
\end{equation}
La phase de lecture à l'instant $r$ est
$\theta_r=(r-1)\bmod3$. Pour chaque canal, on définit
\[
 \varepsilon_\chi(r)=
 \begin{cases}
 1,&\eta^\vee(B_\chi(r))=\theta_r,\\
 2,&\eta^\vee(B_\chi(r))\ne\theta_r.
 \end{cases}
\]
La représentation phase--charge est
\begin{equation}
 D(r)=\varepsilon_\pi(r)B_\pi(r)
      +\varepsilon_e(r)B_e(r)
      +\varepsilon_\varphi(r)B_\varphi(r)
      \quad\text{en }\mathbb F_3^6.
 \label{act21:eq:kdes-fasecarga}
\end{equation}
Le coefficient se détermine par une égalité de charges, avant d'évaluer le mot qui résultera de la combinaison.

\subsubsection{Composition du descripteur}

On désigne par $\Vert$ la concaténation de mots de largeur fixée.
La construction complète est
\begin{equation}
 \begin{split}
 \Omega_{24}={}&
 (B_e(t_*+1)+\delta_-)
 \Vert (\delta_-|_{3,4,5}\Vert\delta_+|_{3,4,5})\\
 &\Vert(B_\pi(r+1)+B_e(r+1))\Vert D(r).
 \end{split}
 \label{act21:eq:kdes-omega}
\end{equation}
Chaque terme de longueur six se calcule en $\mathbb F_3^6$ ; la
concaténation finale a une longueur $6+6+6+6=24$. Cette construction
sépare le descripteur $\Omega_{24}$ des étapes régionales
$w_{24}^{\chi}$ : les deux objets partagent la même longueur, mais leurs producteurs
et leurs places dans la composition sont différents.

\begin{proposition}[Évaluation régionale du descripteur]
Les représentations régionales et le critère de la première avancée nulle
déterminent
\[
 \Omega_{24}=020022\,222111\,211201\,021101,
 \qquad [\Omega_{24}]_3=66241910521.
\]
\end{proposition}
\begin{proof}
Les bandes nécessaires, conservant l'ordre $(\pi,e,\varphi)$, sont
\[
 \begin{array}{c|ccc}
 t&B_\pi(t)&B_e(t)&B_\varphi(t)\\\hline
 11&022212&110001&100021\\
 12&012012&202222&000120\\
 21&221022&020100&211021
 \end{array}
\]
et proviennent de l'extraction \eqref{act21:eq:kdes-extraccion}. Au temps
vingt-et-un,
$020100+000222=020022$ en $\mathbb F_3^6$. Les queues des défauts
apportent $222111$. Au temps douze,
$012012+202222=211201$.

Pour le temps onze, les charges duales sont respectivement $0,1,2$ ; la phase vaut $(11-1)\bmod3=1$. Les coefficients sont $(2,1,2)$ et
\[
 2(022212)+(110001)+2(100021)=021101
 \quad\text{en }\mathbb F_3^6.
\]
La concaténation de ces quatre résultats donne le mot annoncé.
Son évaluation entière s'obtient par Horner en base trois, ou par trois
concaténations successives en base $729$.
\end{proof}

Le tableau précédent joue un rôle vérificateur : chacune de ses entrées provient du producteur régional. La règle de construction \eqref{act21:eq:kdes-omega} est celle qui permet de reproduire le descripteur sans introduire en tant que donnée la chaîne finale de vingt-quatre trits.

\subsection{Panneau fonctionnel et répertoire terminal}

\subsubsection{Neuf positions régionales}

Le panneau fonctionnel conserve les trois premières bandes de chacun des trois canaux, dans l'ordre régional déclaré :
\begin{equation}
 \mathcal P_9=
 \bigl(b_\pi(0),b_\pi(1),b_\pi(2),
       b_e(0),b_e(1),b_e(2),
       b_\varphi(0),b_\varphi(1),b_\varphi(2)\bigr).
 \label{act21:eq:kdes-panel}
\end{equation}
Son évaluation est
\[
 \mathcal P_9=(103,161,103,523,457,301,450,398,438).
\]
Les neuf positions conservent l'origine régionale, bien que la valeur
103 apparaisse deux fois. Le panneau comporte neuf incidences ordonnées et
huit valeurs distinctes. L'énumération ultérieure utilise les incidences ;
la déduplication des candidats s'effectue à une étape postérieure.

\subsubsection{Répertoire non ordonné de huit blocs}

Le répertoire terminal qui intervient dans cette construction est
\begin{equation}
 \mathcal S_{\rm term}=
 \{002001,020111,200110,121012,
   010122,001001,221110,222110\}\subset\mathbb F_3^6.
 \label{act21:eq:kdes-sterm}
\end{equation}
Ses huit éléments sont distincts. La codification entière, ordonnée par
valeur uniquement pour fixer une représentation informatique, est
\[
 \{28,55,98,175,437,498,687,714\}.
\]
La provenance documentaire de ce répertoire est la segmentation terminale du mot de soixante-douze trits et la récupération par positions au moyen des incidences des lecteurs phase--charge et affine. Le sélecteur développé ci-après reçoit le répertoire sans son ordonnancement historique et met à l'épreuve ses $8!$ ordonnancements.

La formulation exacte de l'implémentation réunie maintient ici une
interface explicite : les producteurs régionaux remplacent matériellement
le descripteur, le panneau et les lignes temporelles ; le répertoire
\eqref{act21:eq:kdes-sterm} conserve son producteur documentaire antérieur. Cette
distinction identifie où intervient chaque antécédent. L'appartenance de
huit mots à un répertoire temporel de lecteurs, la sélection de ces
huit mots et la détermination de leur ordre terminal sont trois
énoncés différents. La preuve de sélection terminale démontre le
troisième pour le répertoire déclaré.

La récupération documentaire par positions peut s'exprimer de manière
exacte. Soit $\mathcal I$ le registre des incidences phase--charge et soit
$\mathcal J$ le registre de l'extension affine. Chaque incidence conserve
la position $\ell$ du mot de soixante-douze trits et le mot
émis $u$. Pour $5\leq\ell\leq12$, on forme
\[
 \mathcal S_\ell=
 \{u:(\ell,u)\in\mathcal I\cup\mathcal J\}.
\]
Le récupérateur vérifie que les huit positions apparaissent et que chaque
$\mathcal S_\ell$ est un singleton. Si $s_\ell$ est son élément unique,
la sortie livrée au sélecteur est
\[
 \mathcal S_{\rm term}=\{s_5,s_6,\ldots,s_{12}\}.
\]
La preuve de singleton par position préserve la cohérence des témoins historiques ; l'exportation ultérieure préserve les huit mots et supprime uniquement la mise en ordre historique. Le sélecteur reconstruit l'ordre par l'énumération complète et les conditions terminales, au lieu de le recevoir du registre d'origine.

% RESULTADO_RECUPERADO: K31-10/P10 del inventario REV05.
% Fuente de las incidencias:
% 16_CIERRE_GLOBAL_HMT_MD_2026-07-22/02_LECTURA_DODECAFASICA/continuacion_fases/
% verificar_obstruccion_y_dato_minimo_w72.py, lineas 45--60 y 141--233.
% RESULTADO_OBSTRUCCION_Y_DATO_MINIMO_W72.json, phase_charge_reader.
% Fuente de las bandas: N69_signatures_t0_t99.csv, tiempos impresos abajo.
% Consumidor: PAQUETE_CONTINUIDAD_K_UNIDAD_20260918_BASE_COMPARTIDA/python/
% selector_orbital_original.py, verify_historical_inputs, lineas 159--193.
% No se transfiere la restriccion del lector comparativo al generador completo.
\subsubsection{Incidences temporelles du répertoire terminal et récupération par position}
\label{act21:sec:repertorio-incidencias-completas}

La représentation du répertoire terminal utilise les mêmes blocs régionaux,
les mêmes retenues et le même calendrier qui produisent le descripteur initial
de vingt-quatre trits. Il convient d'expliciter séparément deux opérations de
cette composition. La première récupère huit mots à partir d'incidences
temporelles enregistrées. La seconde présente ces huit mots au sélecteur
dodécaphasique comme un ensemble non ordonné. Le sélecteur rétablit
l'ordre par ses conditions d'incidence, de cylindre et de clôture.
Cette distinction permet de conserver les témoins de provenance sans transformer
l'ordre d'une segmentation historique en donnée d'entrée du sélecteur.

Soit $B_t=(b_{t,0},b_{t,1},b_{t,2})\in(\mathbb F_3^6)^3$ la frontière
régionale déjà engendrée au temps $t$, avec $0\leq t<100$. Dans la coordinatisation
utilisée par le registre historique, la matrice d'incidence est
\[
 A_{\mathrm{hist}}=
 \begin{pmatrix}
 0&1&1&1&1&1\\
 1&0&1&1&2&2\\
 1&1&0&2&1&2\\
 1&1&2&0&2&1\\
 1&2&1&2&0&1\\
 1&2&2&1&1&0
 \end{pmatrix}.
\]
Les mots sont considérés comme des vecteurs ligne. Leurs deux charges et la phase
chronologique sont
\[
 r_{t,i}=\sum_{j=1}^{6}b_{t,i,j},\qquad
 d_{t,i}=\sum_{j=1}^{6}(b_{t,i}A_{\mathrm{hist}})_j,
 \qquad \vartheta_t=t-1\pmod 3.
\]
Toutes les opérations de ces trois expressions ont lieu dans
$\mathbb F_3$. L'indice historique identifie une coordinatisation précise de
l'incidence de Witt ; il maintient explicite la transformation de coordonnées
lorsqu'une autre présentation de cette même incidence est utilisée.

Pour $c\in\{r_t,d_t\}$, $\delta\in\mathbb F_3$ et
$\epsilon\in\{1,-1\}$, nous définissons
\[
 \eta_i(c,\delta,\epsilon)=
 \begin{cases}
  \epsilon,&c_i=\vartheta_t+\delta,\\
  -\epsilon,&c_i\ne\vartheta_t+\delta,
 \end{cases}
 \qquad
 L_{t,c,\delta,\epsilon}
   =\sum_{i=0}^{2}\eta_i(c,\delta,\epsilon)b_{t,i}.
\]
L'orientation $-1$ s'écrit $2$ lors de l'impression des coefficients
ternaires. Outre ces douze lecteurs comparatifs par frontière, le
registre dispose du lecteur affine
\[
 L^{\mathrm{af}}_t
    =\sum_{i=0}^{2}(d_{t,i}+\vartheta_t)b_{t,i}.
\]
L'addition de la phase dans le lecteur affine et la comparaison des charges dans
le lecteur précédent sont des opérations distinctes, toutes deux sur des données déjà
produites par la frontière TPK.

\begin{table}[htbp]
\centering\small
\begin{tabular}{rccc cc}
\toprule
$t$&$b_{t,0}$&$b_{t,1}$&$b_{t,2}$&$r_t$&$d_t$\\
\midrule
8 &200121&212212&110110&011&202\\
11&022212&110001&100021&001&012\\
30&210112&110202&211110&100&012\\
34&121022&000011&112110&220&021\\
37&120121&011202&112221&100&201\\
49&110212&200212&100122&110&201\\
58&111121&022121&122201&122&220\\
67&010012&001220&011120&122&122\\
75&220000&020111&100002&120&021\\
81&122201&020202&201122&202&001\\
90&022112&022110&121010&202&200\\
\bottomrule
\end{tabular}
\caption{Frontières temporelles qui réalisent les incidences imprimées du
répertoire terminal. Chaque entrée de six chiffres conserve l'ordre des
six coordonnées de la frontière régionale.}
\label{act21:tab:terminal-bandas-testigo}
\end{table}

Le tableau suivant déploie toutes les incidences comparatives du
registre sur les douze blocs de la segmentation considérée. Les
lettres $r$ et $d$ désignent respectivement la charge visible et la charge
duale. La colonne $j$ conserve l'indice de position du témoin historique ;
l'opération de livraison au sélecteur oublie cet indice après avoir formé
l'ensemble de mots.

\begin{table}[htbp]
\centering\small
\begin{tabular}{rrcrrc c}
\toprule
$j$&$t$&charge&$\delta$&$\epsilon$&$(\eta_0\eta_1\eta_2)$&$L_{t,c,\delta,\epsilon}$\\
\midrule
4&11&$d$&0& 1&212&021101\\
5&67&$d$&1&$-1$&211&002001\\
5&67&$d$&2& 1&211&002001\\
5&67&$r$&1&$-1$&211&002001\\
5&67&$r$&2& 1&211&002001\\
6&81&$d$&0& 1&222&020111\\
6&81&$r$&2& 1&222&020111\\
7&34&$r$&1&$-1$&111&200110\\
7&75&$d$&1&$-1$&211&200110\\
7&75&$r$&2&$-1$&211&200110\\
8&90&$r$&0& 1&121&121012\\
8&90&$r$&1&$-1$&121&121012\\
9&49&$d$&0&$-1$&121&010122\\
11&11&$d$&2&$-1$&211&221110\\
11&37&$d$&0&$-1$&121&221110\\
12&8&$d$&0&$-1$&111&222110\\
12&8&$r$&1&$-1$&111&222110\\
12&58&$d$&1&$-1$&111&222110\\
12&58&$r$&0&$-1$&111&222110\\
\bottomrule
\end{tabular}
\caption{Incidences phase--charge complètes sur la segmentation terminale  
aux cent frontières du registre. Les coïncidences de mots entre  
lignes conservent leurs temps, orientations et types de charge.}
\label{act21:tab:terminal-incidencias-completas}
\end{table}

Le dixième bloc dispose d'un témoin affine. Dans $t=30$ on a
$d_{30}=(0,1,2)$ et $\vartheta_{30}=2$, d'où
\[
 (d_{30,0}+\vartheta_{30},d_{30,1}+\vartheta_{30},
 d_{30,2}+\vartheta_{30})=(2,0,1)
\]
et, en utilisant les bandes du tableau ~\ref{act21:tab:terminal-bandas-testigo},
\[
 L^{\mathrm{af}}_{30}
 =2(210112)+0(110202)+(211110)
 =(001001)\quad\text{en }\mathbb F_3^6.
\]
L'égalité est coordonnée par coordonnée : avant de réduire modulo trois,  
le vecteur de la somme est $(6,3,1,3,3,4)$.

\begin{proposition}[Récupération intégrale du répertoire par incidences]
\label{act21:prop:repertorio-terminal-ocho-incidencias}
Considérons le registre des incidences comparatives du
tableau~\ref{act21:tab:terminal-incidencias-completas}, augmenté du témoin
affine $(j,t,L)=(10,30,001001)$. Pour chaque position $5\leq j\leq12$,
l'ensemble des mots de sortie de ses incidences est un singleton. L'
extraction de ces huit mots et l'oubli ultérieur de leur ordre
produisent exactement
\[
 \mathcal S_{\mathrm{term}}=
 \{002001,020111,200110,121012,010122,001001,221110,222110\}.
\]
Sa codification en tant que nombres entiers en base trois est
\[
 \{28,55,98,175,437,498,687,714\}.
\]
\end{proposition}
\begin{proof}
La multiplicité d'incidences par position $j=5,6,7,8,9,10,11,12$
est, respectivement,
\[
 (4,2,3,2,1,1,2,4).
\]
À chaque position, toutes les lignes impriment le même mot. Chaque ligne se vérifie en substituant ses trois bandes et les trois coefficients dans la formule de $L$ ; le dixième cas a été calculé explicitement au moyen de $L^{\mathrm{af}}_{30}$. Par exemple, en $t=67$ et avec les coefficients $(2,1,1)$, on obtient
\[
 2(010012)+(001220)+(011120)=(002001)\pmod3.
\]
Les huit résultats sont distincts en tant que mots de longueur six. Leur codification $\sum_{k=1}^{6}w_k3^{6-k}$ donne, dans l'ordre des positions, $(55,175,498,437,98,28,714,687)$. En oubliant l'ordre, on obtient l'ensemble des entiers énoncé. Ainsi, l'extraction préserve chaque mot complet, y compris ses positions nulles, et fournit huit éléments distincts au sélecteur.
\end{proof}

La preuve précédente est une récupération exacte à partir d'incidences
identifiées. Sa donnée d'entrée comprend le registre d'incidences
avec leurs positions et leurs témoins. La sélection orbitale ultérieure reçoit
uniquement $\mathcal S_{\mathrm{term}}$ et le descripteur régional initial ;
elle parcourt les $8!$ ordonnancements et détermine leur compatibilité avec les
frontières fonctionnelles et les incidences orientées de la roue
dodécaphasique. Les deux opérations sont ainsi contiguës, avec des domaines
explicites : provenance temporelle des huit mots et détermination
constructive de leur ordre.

Le tableau comparatif contient les dix-neuf incidences que sa formule produit sur la segmentation de douze blocs aux cent temps déclarés. Cette exhaustivité finie se vérifie en parcourant $100\cdot2\cdot3\cdot2=1200$ évaluations. Le lecteur affine étend la couverture des positions terminales de sept à huit. Ce décompte caractérise le répertoire local de lecteurs qui vient d'être défini ; les autres opérations du TPK conservent les domaines et les fonctions constructives établis dans leurs dérivations respectives.

\subsection{Lecteurs comparatifs, lecteur affine et incidence temporelle}

\subsubsection{Transformation de Witt du sélecteur}

Le sélecteur utilise la coordinatisation à six coordonnées
\[
 A_{\rm ter}=
 \begin{pmatrix}
 0&1&1&1&1&1\\
 1&0&1&2&2&1\\
 1&1&0&1&2&2\\
 1&2&1&0&1&2\\
 1&2&2&1&0&1\\
 1&1&2&2&1&0
 \end{pmatrix}.
\]
On définit $W(b)=bA_{\rm ter}$ dans $\mathbb F_3^6$, ainsi que les charges
\[
 \eta(b)=\sum_{j=0}^5 b_j\pmod3,
 \qquad\eta_{\rm ter}^\vee(b)=\eta(W(b)).
\]
La coordinatisation du descripteur et celle du sélecteur sont reliées par la
permutation $\sigma=(0,1,2,4,5,3)$, écrite comme liste d'images :
\begin{equation}
 (A_{\rm ter})_{ij}=(A_{\rm reg})_{\sigma(i),\sigma(j)}.
 \label{act21:eq:kdes-cambio-carta}
\end{equation}

\begin{lemma}[Compatibilité de la fonctionnelle de charge]
Pour toute bande $b$, on a
\[
 \eta_{\rm ter}^\vee(b)=\eta^\vee(b).
\]
\end{lemma}
\begin{proof}
La somme entière de chaque ligne correspondante des deux matrices coïncide :
elle vaut cinq pour la première ligne et sept pour les cinq autres. Par
distributivité,
\[
 \sum_j\sum_i b_i(A_{\rm ter})_{ij}
 =\sum_i b_i\sum_j(A_{\rm ter})_{ij}
 =\sum_i b_i\sum_j(A_{\rm reg})_{ij}.
\]
La réduction modulo trois donne l'égalité des charges. L'identité des matrices \eqref{act21:eq:kdes-cambio-carta} se vérifie par substitution de ses trente-six entrées et conserve, en outre, le changement de coordonnées.
\end{proof}

L'égalité de la fonctionnelle totale de charge permet de comparer ses
représentations. Les transformations vectorielles restent identifiées
par leurs matrices et par la permutation explicite ; la preuve précédente
conserve cette information de coordonnées.

\subsubsection{Douze lecteurs comparatifs par instant}

On définit la phase temporelle
\[
 \theta(t)=(t+2)\bmod3.
\]
Cette écriture réalise également $(t-1)\bmod3$ pour $t=0$, où la phase
vaut deux. Pour chaque type de charge
$\nu\in\{\eta,\eta_{\rm ter}^\vee\}$, déplacement
$o\in\{0,1,2\}$ et orientation $s\in\{1,2\}$, on prend
\[
 \epsilon_\chi^{\nu,o,s}(t)=s
 \begin{cases}
 1,&\nu(B_\chi(t))=\theta(t)+o\pmod3,\\
 2,&\nu(B_\chi(t))\ne\theta(t)+o\pmod3.
 \end{cases}
\]
Le mot comparatif correspondant est
\begin{equation}
 C_{\nu,o,s}(t)=
 \sum_{\chi\in\{\pi,e,\varphi\}}
 \epsilon_\chi^{\nu,o,s}(t)B_\chi(t)
 \quad\text{en }\mathbb F_3^6.
 \label{act21:eq:kdes-comparativo}
\end{equation}
Les paramètres engendrent douze représentations étiquetées par ligne. Deux
étiquettes peuvent produire le même mot ; les témoins temporels et
les étiquettes restent disponibles dans le calendrier complet.

\subsubsection{Lecteur affine de phase et de charge}

Le lecteur affine emploie des coefficients
\[
 a_\chi^{\rm af}(t)=\eta_{\rm ter}^\vee(B_\chi(t))+\theta(t)
 \pmod3,
\]
et produit
\begin{equation}
 F(t)=\sum_{\chi\in\{\pi,e,\varphi\}}
 a_\chi^{\rm af}(t)B_\chi(t)
 \quad\text{en }\mathbb F_3^6.
 \label{act21:eq:kdes-afin}
\end{equation}
La règle comparative distingue la coïncidence et la différence de charges ; la règle affine additionne charge et phase avant de combiner les bandes. C'est cette différence opératoire qu'utilise la sélection hétérotypique.

\subsubsection{Incidence mot--résidu}

Soit $R$ la rotation cyclique des six positions d'une bande. À la
ligne temporelle $t$, on forme le répertoire résiduel
\[
 \mathcal R(t)=\{R^jv:
      v\in\{q(t),a(t),c(t),\bar w(t)\},\ 0\leq j<6\}.
\]
Les relations d'incidence temporelle sont
\begin{align}
 \mathcal A_{\rm cmp}
 &=\{(C_{\nu,o,s}(t),u):
      0\leq t<100,\ u\in\mathcal R(t),\ \nu,o,s\text{ admissibles}\},\\
 \mathcal A_{\rm af}
 &=\{(F(t),u):0\leq t<100,\ u\in\mathcal R(t)\}.
 \label{act21:eq:kdes-atlas}
\end{align}
Chaque couple relie un mot lu à un résidu de la même ligne.
Le temps se quantifie de manière commune dans les deux coordonnées du couple.
Dans le calendrier employé, les temps $0,\ldots,99$ sont distincts ;
identifier une même ligne et identifier un même temps sont, par conséquent,
des opérations équivalentes.

\begin{proposition}[Compression existentielle du répertoire temporel]
La fonction
\[
 M(v,u)=\mathbf1_{(v,u)\in\mathcal A_{\rm cmp}}
        +2\mathbf1_{(v,u)\in\mathcal A_{\rm af}}
        \in\{0,1,2,3\}
\]
conserve exactement l'existence de chaque type de lecture pour le couple  
$(v,u)$.
\end{proposition}
\begin{proof}
Le premier bit vaut un exactement lorsqu'il existe un temps, un champ
résiduel, une rotation et une étiquette comparative qui produisent le couple.
Le second bit exprime la condition correspondant au lecteur affine.
L'union par disjonction des bits conserve chaque condition d'existence
lors du parcours de toutes les lignes. Un couple peut admettre les deux types et, dans ce
cas, le masque vaut trois.
\end{proof}

La fonction $M$ est un lecteur d'existence. Les temps témoins, leurs multiplicités et les étiquettes individuelles demeurent dans \eqref{act21:eq:kdes-atlas} et dans le calendrier qui les produit. L'implémentation finie emploie $M$ pour décider de l'hétérotypie, tandis que l'archive généalogique conserve les données de provenance plus riches.

\subsection{Énumération de cylindres et candidats décimaux}

\subsubsection{Concaténation de soixante-douze et soixante-dix-huit trits}

Pour un ordonnancement $\tau=(s_1,\ldots,s_8)$ de
$\mathcal S_{\rm term}$, on définit
\[
 P_{72}(\tau)=[\Omega_{24}\Vert s_1\Vert\cdots\Vert s_8]_3.
\]
De manière récursive, on part de $p_0=[\Omega_{24}]_3$ et on calcule
$p_i=729p_{i-1}+[s_i]_3$ pour $1\leq i\leq8$. Cette récurrence
conserve les zéros initiaux de chaque bloc par la multiplication
explicite par $729$. Pour chaque position $b$ du panneau fonctionnel,
\[
 P_{78}(\tau,b)=729P_{72}(\tau)+b.
\]
Les longueurs sont respectivement $24+8\cdot6=72$ et $72+6=78$.
Le cylindre ternaire correspondant est
\begin{equation}
 I_{\tau,b}=
 \left[\frac{P_{78}(\tau,b)}{3^{78}},
       \frac{P_{78}(\tau,b)+1}{3^{78}}\right).
 \label{act21:eq:kdes-cilindro}
\end{equation}

\subsubsection{Intersection exacte avec la grille de douze triades}

La lecture décimale utilise douze triades, c'est-à-dire la grille $N/10^{36}$. On utilise les abréviations $D=10^{36}$ et $Q=3^{78}$.

\begin{lemma}[Extrémités entières du cylindre]
Pour un préfixe entier $p$ de longueur soixante-dix-huit, les numérateurs
des points de la grille qui appartiennent à son cylindre sont
\begin{equation}
 \left\{N\in\mathbb Z:
 \left\lceil\frac{pD}{Q}\right\rceil
 \leq N<
 \left\lceil\frac{(p+1)D}{Q}\right\rceil\right\}.
 \label{act21:eq:kdes-rejilla}
\end{equation}
Chaque cylindre contient au plus un point de cette grille.
\end{lemma}
\begin{proof}
La condition $N/D\in[p/Q,(p+1)/Q)$ est équivalente à
$pD/Q\leq N<(p+1)D/Q$. Pour $N$ entier, la première inégalité
est équivalente à $\lceil pD/Q\rceil\leq N$ ; la seconde est équivalente à
$N<\lceil(p+1)D/Q\rceil$. La longueur de l'intervalle en coordonnée
$N$ est $D/Q<1$, car $10^{36}<3^{78}$. Par conséquent, il peut contenir,
au maximum, un entier.
\end{proof}

Les plafonds sont évalués par division entière :
$\lceil a/Q\rceil=\lfloor(a+Q-1)/Q\rfloor$ pour $a\geq0$.
Toute l'énumération utilise des entiers exacts ; le choix des points de frontière est fixé par l'extrémité droite semi-ouverte.

On conserve d'abord la liste indexée d'incidences
\[
 \mathcal C_{\rm ind}=
 \{(\tau,j,N):N/D\in I_{\tau,\mathcal P_9(j)}\},
 \qquad1\leq j\leq9,
\]
et l'on obtient ensuite l'ensemble des numérateurs
\[
 \mathcal C=\{N:\exists\tau,j\ (\tau,j,N)\in\mathcal C_{\rm ind}\}.
\]
Le passage à $\mathcal C$ élimine des répétitions d'un numérateur, conservées
dans $\mathcal C_{\rm ind}$ avec leur origine.

\begin{lemma}[Récupération de l'ordre à partir du point du cylindre]
Si $(\tau,j,N)\in\mathcal C_{\rm ind}$, le bloc numéro $i$ du
mot de soixante-dix-huit trits se récupère par
\[
 \left\lfloor\frac{N\,729^i}{D}\right\rfloor\bmod729,
 \qquad1\leq i\leq13.
\]
En particulier, les blocs cinquième à douzième récupèrent $\tau$.
\end{lemma}
\begin{proof}
L'appartenance au cylindre signifie
$p\leq QN/D<p+1$, où $p=P_{78}(\tau,\mathcal P_9(j))$.
Par conséquent, $\lfloor QN/D\rfloor=p$. L'extraction des blocs
précédents coïncide avec la division successive du préfixe $p$ par
des puissances de $729$ : tronquer une expansion positionnelle avant d'extraire
un bloc conserve ce bloc. La formule écrite réalise
exactement cette extraction. Les quatre premiers blocs forment
$\Omega_{24}$ et les huit suivants sont l'ordonnancement $\tau$.
\end{proof}

\begin{proposition}[Recensement exact des candidats]
Pour $\Omega_{24}$, $\mathcal S_{\rm term}$ et $\mathcal P_9$
définis précédemment,
\[
 \#\operatorname{Ord}(\mathcal S_{\rm term})=40320,
 \qquad\#\mathcal C_{\rm ind}=21902,
 \qquad\#\mathcal C=19446.
\]
\end{proposition}
\begin{proof}
Les huit blocs sont distincts, donc leurs ordonnements sont les
$8!=40320$ permutations. Pour chaque permutation et chacune des
neuf positions du panneau, on calcule $p=P_{78}(\tau,b)$, on applique
les extrêmes \eqref{act21:eq:kdes-rejilla} et on ajoute son unique entier lorsque
l'intervalle entier n'est pas vide. La somme des cardinaux de ces
intervalles est $21902$. La réunion de leurs entiers contient $19446$
éléments distincts. L'évaluation exhaustive de ces opérations
est certifiée par les théorèmes
\texttt{indexed\_cylinder\_count} et \texttt{candidate\_count} du
sélecteur fini. L'algorithme réalise exactement l'énumération
décrite : permutations, neuf incidences du panneau, limites par
division entière et élimination des doublons.
\end{proof}

Les nombres $21902$ et $19446$ comptent des objets différents. Le premier
compte des incidences indexées qui contiennent un point décimal ; le second
compte des numérateurs distincts. Le nombre total de requêtes de cylindre
est $40320\cdot9=362880$, y compris les cylindres dont l'intersection avec
la grille est vide.

\subsection{Incidence exceptionnelle et condition d'hétérotypie}

\subsubsection{Face exceptionnelle produite par le panneau}

On élève une bande à douze coordonnées par le biais de
\[
 \Gamma(b)=(b,W(b))\in\mathbb F_3^{12}.
\]
Le support d'un mot est l'ensemble des positions de ses coordonnées non nulles, avec les indices humains $1,\ldots,12$. Les premières bandes régionales du panneau produisent
\begin{equation}
 \mathcal B=\operatorname{supp}\Gamma(B_\pi(0))
          \cap\operatorname{supp}\Gamma(B_e(0))
          =\{4,6,9,10\}.
 \label{act21:eq:kdes-cara}
\end{equation}
L'opération utilise les mots régionaux $103$ et $523$ et la matrice de Witt du sélecteur. La face se calcule avant de consulter les coordonnées des candidats.

Pour $N\in\mathcal C$, ses douze triades sont
\begin{equation}
 K_i(N)=\left\lfloor\frac{N}{1000^{12-i}}\right\rfloor\bmod1000,
 \qquad1\leq i\leq12.
 \label{act21:eq:kdes-triadas}
\end{equation}
Le support de couronne est
\[
 \mathcal H(N)=\{i:K_i(N)\geq729\},
 \qquad
 \mathcal C_{\mathcal B}=\{N\in\mathcal C:\mathcal H(N)=\mathcal B\}.
\]
La frontière $729=3^6$ relie la bande de six trits avec la triade décimale. L'évaluation exhaustive de l'égalité des supports donne
\begin{equation}
 \#\mathcal C_{\mathcal B}=79.
 \label{act21:eq:kdes-79}
\end{equation}
Le support conserve les positions, tandis que le registre complet
conserve aussi les douze grandeurs. L'étape suivante utilise cette
information positionnelle conjointement avec des relations entre grandeurs.

\subsubsection{Mots ternaires et différences orientées du candidat}

La $i$-ième bande ternaire du point $N/D$ s'obtient par
\begin{equation}
 T_i(N)=\left\lfloor\frac{N\,729^i}{D}\right\rfloor\bmod729,
 \qquad1\leq i\leq13.
 \label{act21:eq:kdes-tercandidato}
\end{equation}
L'arête orientée $a\to b$ a un résidu
\begin{equation}
 R_{a,b}(N)=(K_a(N)-K_b(N))\bmod729.
 \label{act21:eq:kdes-resarista}
\end{equation}
La soustraction s'effectue en $\mathbb Z$ avant la réduction. Pour cette
arête, on consulte dans les relations \eqref{act21:eq:kdes-atlas} le couple
\[
 \bigl(\operatorname{enc}_6^{-1}(T_b(N)),
  \operatorname{enc}_6^{-1}(R_{a,b}(N))\bigr).
\]
C'est donc une consultation conjointe sur le mot cible et sur la différence orientée des triades.

On écrit $\mathrm{Cmp}(a,b;N)$ lorsque le couple appartient à $\mathcal A_{\rm cmp}$ et $\mathrm{Af}(a,b;N)$ lorsqu'il appartient à $\mathcal A_{\rm af}$.

\subsubsection{Détermination de l'orbite transversale}

Le descripteur de vingt-quatre trits détermine les trois premières triades décimales avant les filtrations terminales. En effet, son cylindre satisfait l'inclusion exacte
\[
 \left[\frac{66241910521}{3^{24}},
       \frac{66241910522}{3^{24}}\right)
 \subseteq
 \left[\frac{234543140}{10^9},
       \frac{234543141}{10^9}\right).
\]
Les deux inégalités aux extrémités se vérifient par des produits croisés
entiers. Tout candidat partage donc les trois positions de
référence $1,2,3$ et leurs représentations $(234,543,140)$.

La translation de pas quatre sur $\mathbb Z/12\mathbb Z$ possède les
quatre orbites
\[
 (1,5,9),\qquad(2,6,10),\qquad(3,7,11),\qquad(4,8,12).
\]
Les trois premières contiennent respectivement l'une des positions de référence ; la quatrième est la seule disjointe de $\{1,2,3\}$. La forêt orientée des pas trois et quatre est fixée au moyen des arêtes
\[
 \begin{array}{lll}
 1\to4,&1\to5,&2\to6,\\
 3\to7,&4\to8,&5\to9,\\
 6\to10,&7\to11,&8\to12.
 \end{array}
\]
Sa première arête a un pas de trois et les autres ont un pas de quatre. Les deux
arêtes internes de l'orbite distinguée sont précisément
$4\to8$ et $8\to12$. Ainsi est fixé le lieu de la condition de
lecture par la géométrie des positions et par le préfixe initial,
avant d'examiner les candidats de la face exceptionnelle.

\subsubsection{Condition terminale d'hétérotypie}

La condition terminale est
\begin{equation}
 \begin{split}
 \mathcal T(N)={}&
 [\mathrm{Cmp}(4,8;N)\wedge\mathrm{Af}(8,12;N)]\\
 &\vee[\mathrm{Af}(4,8;N)\wedge\mathrm{Cmp}(8,12;N)].
 \end{split}
 \label{act21:eq:kdes-heterotipia}
\end{equation}
Les deux affectations de types aux deux arêtes sont admises. La
sélection exige qu'il existe une réalisation de types distincts ; les
étiquettes individuelles et leurs temps restent disponibles dans les
témoins d'incidence.

\begin{theorem}[Sélection terminale du registre de couplage]
Sur le domaine de candidats produit par
$\Omega_{24}$, $\mathcal S_{\rm term}$, $\mathcal P_9$ et le calendrier
régional de cent lignes, existe un unique numérateur qui satisfait
simultanément l'incidence exceptionnelle et l'hétérotypie :
\[
 \{N\in\mathcal C_{\mathcal B}:\mathcal T(N)\}
 =\{234543140729659824621058914794146601\}.
\]
Son registre de douze triades est
\begin{equation}
 K=(234,543,140,729,659,824,621,058,914,794,146,601).
 \label{act21:eq:kdes-kfinal}
\end{equation}
\end{theorem}
\begin{proof}
Le recensement précédent produit les $79$ éléments de
$\mathcal C_{\mathcal B}$. Pour chacun, on extrait ses douze
coordonnées par \eqref{act21:eq:kdes-triadas} ; on calcule les bandes de
destination $T_8,T_{12}$ et les résidus $R_{4,8},R_{8,12}$ ; on consulte
les deux masques du répertoire temporel et on applique la disjonction
\eqref{act21:eq:kdes-heterotipia}. La liste résultante contient exactement
l'entier indiqué, selon l'évaluation exhaustive certifiée par
\texttt{terminal\_selection}. L'égalité de liste avec un singleton
implique existence et unicité dans le domaine déclaré. L'extraction
base mille de cet entier donne \eqref{act21:eq:kdes-kfinal}.

La procédure de sélection reçoit les producteurs et les règles énumérées. L'entier final intervient dans l'énoncé de l'égalité vérifiée, après avoir exécuté la sélection. Cette position logique permet de comparer le résultat sans l'utiliser comme critère de choix.
\end{proof}

\subsubsection{Témoin arithmétique d'appartenance à un cylindre}

Une des incidences du numérateur sélectionné utilise la mise en ordre
\[
 (002001,020111,200110,121012,010122,001001,221110,222110)
\]
et la frontière $438=121020_3$. Pour elle,
\[
 P_{72}=5283881584882673136514102926090957,
\]
\[
 P_{78}=3851949675379468716518781033120308091.
\]
La formule des extrémités entières produit
\begin{align*}
 \left\lceil\frac{P_{78}10^{36}}{3^{78}}\right\rceil
 &=234543140729659824621058914794146601,\\
 \left\lceil\frac{(P_{78}+1)10^{36}}{3^{78}}\right\rceil
 &=234543140729659824621058914794146602.
\end{align*}
L'unique entier de l'intervalle est donc $N_K$. Les treize bandes ternaires de son point décimal sont
\[
 \begin{array}{rrrr}
 020022&222111&211201&021101\\
 002001&020111&200110&121012\\
 010122&001001&221110&222110\\
 121020&&&
 \end{array}
\]
Ce calcul exhibe un témoin concret d'appartenance au domaine énuméré. L'unicité du registre procède de l'examen exhaustif de tous les candidats, outre ce témoin particulier.

\subsection{Composition régionale et continuité des réalisations}

\subsubsection{Composition des générateurs régionaux}

Soit $\operatorname{Sel}(\Omega,\mathcal S,\mathcal P,\mathcal N)$ le
sélecteur défini par l'énumération et les deux filtres précédents.
Les producteurs régionaux satisfont les égalités
\[
 \Omega_{\rm reg}=\Omega_{24},\qquad
 \mathcal P_{\rm reg}=\mathcal P_9,\qquad
 \mathcal N_{\rm reg}=\mathcal N_{100},
\]
où $\mathcal N_{\rm reg}$ se calcule par extraction de bandes et
lecture de signatures. Par congruence de fonctions,
\begin{equation}
 \begin{split}
 &\operatorname{Sel}(\Omega_{\rm reg},\mathcal S_{\rm term},
                    \mathcal P_{\rm reg},\mathcal N_{\rm reg})\\
 &\hspace{10mm}=
 \operatorname{Sel}(\Omega_{24},\mathcal S_{\rm term},
                    \mathcal P_9,\mathcal N_{100})=\{N_K\}.
 \end{split}
 \label{act21:eq:kdes-composicion}
\end{equation}
L'égalité remplace trois interfaces documentaires par leur construction
régionale effective. Le répertoire terminal conserve la provenance
spécifiée dans \eqref{act21:eq:kdes-sterm}.

Le registre se forme alors par les douze divisions euclidiennes de
\eqref{act21:eq:kdes-triadas}. Chaque coordonnée appartient à
$\{0,\ldots,999\}$ et la longueur est douze par construction. L'
implémentation qui prend le premier élément de la liste sélectionnée
est justifiée par l'égalité avec le singleton : la valeur par défaut
d'une opération informatique sur des listes vides n'intervient jamais dans
ce résultat.

\subsubsection{Relation avec les incidences signées et l'inversion}

La construction présente fournit une sélection prospective dans un
domaine fini explicite. Le développement ultérieur du registre
dodécaphasique produit ses canaux orientés à partir de l'état
terminal enrichi et démontre une inversion intégrale par
Hadamard. Cette inversion récupère un registre à partir de canaux
compatibles ; ici, on sélectionne un registre à partir de ses
antécédents régionaux et incidentiels. L'égalité de ses résultats
permet de composer les deux réalisations en maintenant l'origine de chaque
opération.

L'unicité de sélection et l'unicité d'inversion sont des propriétés distinctes. La première quantifie les candidats construits dans ce chapitre. La seconde quantifie les registres qui partagent une signature compatible. Leur conjonction justifie que le registre sélectionné puisse être transporté vers la réalisation signée et être récupéré à partir d'elle, sans convertir une récupération rétrospective en producteur du registre.

\subsubsection{Relation entre les recensements terminaux}

La suite
\[
 40320\ \longrightarrow\ 21902\ \longrightarrow\ 19446
 \ \longrightarrow\ 79\ \longrightarrow\ 1
\]
enregistre successivement des ordonnancements, des incidences indexées avec la grille, des numérateurs distincts, des candidats de la face exceptionnelle et des survivants hétérotypiques. La notation à flèches indique l'ordre des opérations ; le deuxième nombre compte une classe d'objets différente de la première.

Le recensement des catalogues régionaux
$196\cdot197\cdot196\longrightarrow331\longrightarrow2
\longrightarrow1$, développé à l'endroit correspondant, part d'une
autre population : blocs régionaux de trois catalogues et conditions
terminales de Gram, d'occupation, de distance et d'orientation. Les deux recensements
conservent leurs domaines. Leur convergence vers une même représentation
s'exprime par l'identification de leurs lecteurs et de leurs sorties ;
les chiffres intermédiaires décrivent des opérations différentes et restent
documentés séparément.

\subsubsection{Composition régionale de la constante de structure fine}

La sortie $K$ intervient ensuite dans la composition régionale
ordonnée et dans la retenue en base mille. Son emploi ultérieur conserve les
douze positions et la frontière de lecture. Le couplage avec la racine
analytique et les représentations à profondeur arbitraire utilise cette
même sortie, avec les coordonnées régionales déjà engendrées et les
règles de la coordinatisation analytique. Cette dépendance prépare le chapitre
consacré aux deux réalisations corrélées de la constante de
structure fine.

La portée de ce chapitre est précise : production du descripteur et de ses antécédents régionaux, détermination du domaine fini, sélection unique sur ce domaine et composition matérielle avec les producteurs. La récurrence analytique de tous les degrés et la compatibilité des retenues infinies sont développées ensuite avec leurs démonstrations propres.

\subsection{Certification exécutable et traçabilité des opérations}

Le développement formel conserve la correspondance entre les
définitions mathématiques et leurs réalisations exécutables. Les preuves
des signatures et de l'horloge se trouvent dans
\texttt{N69RegionalSignature.lean} ; l'extraction des bandes et la
reproduction des cent lignes, dans
\texttt{GeneratedN69Rows.lean} ; le critère minimal et la construction
de $\Omega_{24}$, dans \texttt{RegionalW24.lean}. La compatibilité des
deux matrices de Witt est réunie dans
\texttt{WittReaderCompatibility.lean}.

Les lecteurs temporels, les rotations, les résidus orientés et la
condition hétérotypique sont formalisés dans
\texttt{TerminalReaders.lean}. L'énumération et les recensements sont
formalisés dans \texttt{Terminal}\allowbreak\texttt{Orbital}\allowbreak\texttt{Selection.lean}. Le panneau
régional se relie par \texttt{Regional}\allowbreak\texttt{Panel}\allowbreak\texttt{Bridge.lean}, et la
composition \eqref{act21:eq:kdes-composicion} s'exprime en
\texttt{Selected}\allowbreak\texttt{From}\allowbreak%
\texttt{Regional}\allowbreak\texttt{Inputs.lean}.

Les vérifications finies d'énumération utilisent l'évaluation native de Lean ; leurs rapports enregistrent \texttt{Lean.ofReduceBool}, ainsi que les axiomes logiques habituels de la bibliothèque. La reconstruction régionale du descripteur utilise des preuves du noyau et une réduction ordinaire. Cette distinction décrit les mécanismes de vérification réellement employés. Les reçus de compilation certifient ses modules et dépendances déclarés ; les preuves mathématiques conservent les domaines, producteurs et hypothèses qui ont été explicités dans le corps du chapitre.

\subsection{Recensement complet des cylindres et sélection de la frontière terminale}
\label{act21:censo:terminal-completo}

La lecture terminale utilise conjointement deux informations : la compatibilité des cylindres régionaux et un profil d'incidence avec des canaux et des positions étiquetés. Dans cette section, on construit les catalogues complets, on énumère leurs triplets et on démontre l'effet de chaque condition du lecteur. Le point de départ est constitué des sorties régionales déjà produites dans la section~\ref{sec:generacion}, et non de valeurs introduites pour définir le transport APP--TRIT--TPK. La réduction à une matrice appartient à une lecture ultérieure de ces sorties et conserve l'ordre des trois canaux.

Le résultat est un théorème de recensement à structure de départ explicite.
Le profil de Gram, des poids, des distances et des marges est déclaré comme donnée du
lecteur terminal. L'exhaustivité de l'énumération démontre ce que sélectionne
ce profil sur les catalogues fixés ; elle ne démontre pas, à elle seule, que le
profil soit une loi universelle ni qu'il puisse être reconstruit à partir du seul triplet
qui le satisfait finalement. Cette séparation permet d'intégrer la preuve
finie complète sans inverser sa généalogie.

\subsubsection{Des sorties régionales aux catalogues ternaires}

Soient \(x_\pi,x_e,x_\varphi\in[0,1)\) les parties fractionnaires des
trois coordonnées régionales déjà construites. On utilise leurs cylindres
exprimés à mille positions décimales,
\[
 J_\chi^{(1000)}=
 \left[\frac{d_\chi}{10^{1000}},
       \frac{d_\chi+1}{10^{1000}}\right),
 \qquad \chi\in\{\pi,e,\varphi\}.
\]
Les entiers \(d_\chi\) sont des lectures finies de sortie. La longueur mille
est suffisante pour les entiers qui suivent, comme on le vérifie par les
deux bornes de chaque intervalle ; ce n'est pas une profondeur supposée infinie.
Définissons
\begin{equation}
 p=30,\quad r=1050,\quad L=p+r=1080,\quad k=171,
 \qquad D=3^L,\quad Q=1000^k,
 \label{act21:censo:parametros}
\end{equation}
et
\[
 P_\chi=\lfloor3^{30}x_\chi\rfloor,\qquad
 T_\chi=\lfloor Qx_\chi\rfloor.
\]
Ici, \(k\) compte des triades décimales et ne désigne pas le registre dodécaphasique \(K\). On a \(1000^{171}\leq3^{1080}<1000^{172}\).

\begin{lemma}[Extraction stable d'un entier]
\label{act21:censo:estabilidad}
Pour \(d,m,h\in\mathbb Z\), avec \(m,h>0\), la valeur de
\(\lfloor mx\rfloor\) est constante en \([d/h,(d+1)/h)\) si et seulement si
\begin{equation}
 \left\lfloor\frac{md}{h}\right\rfloor
 =\left\lfloor\frac{m(d+1)-1}{h}\right\rfloor.
 \label{act21:censo:prueba-estabilidad}
\end{equation}
Lorsqu'elle est satisfaite, chacun des deux membres donne cette valeur.
\end{lemma}
\begin{proof}
Le minimum de \(mx\) est atteint à l'extrémité gauche. Son supremum
est \(m(d+1)/h\) et n'est pas atteint. Le plus grand entier que peut prendre
\(\lfloor mx\rfloor\) est
\(\lceil m(d+1)/h\rceil-1\), égal au second membre car
le numérateur et le dénominateur sont des entiers. L'égalité des extrémités
de l'intervalle de valeurs prouve l'affirmation.
\end{proof}

La preuve de stabilité se vérifie sur les trois canaux pour
\(m=3^{30},Q,3^{1080}\). Le dernier choix sera utilisé uniquement
pour une comparaison ultérieure, pas pour sélectionner le profil. Les
premières trente positions récupérées sont précisément les cinq
étapes de six trits du transport régional :
\begin{center}\small
\begin{tabular}{@{}c l r@{}}\toprule
Canal&Préfixe de trente trits&\(P_\chi\)\\\midrule
\(\pi\)&\texttt{010211012222010211002111110221}&29152671743887\\
\(e\)&\texttt{201101121221102011012222102011}&147887858824447\\
\(\varphi\)&\texttt{121200112202121020010210010200}&127247717616687\\\bottomrule
\end{tabular}
\end{center}
La coïncidence conserve la prolongation déjà construite ; la lecture décimale
n'est pas utilisée pour remplacer rétroactivement ses matrices de transport.

Un canal étant fixé, une queue de \(r\) trits se représente par
\(0\leq R<3^r\). Son cylindre et le cylindre des \(k\) premières
triades sont
\begin{equation}
 I_{\chi,R}=\left[\frac{P_\chi3^r+R}{D},
                       \frac{P_\chi3^r+R+1}{D}\right),\qquad
 J_{\chi,T}=\left[\frac{T_\chi}{Q},\frac{T_\chi+1}{Q}\right).
 \label{act21:censo:dos-cilindros}
\end{equation}
Le catalogue inclut toutes les queues pour lesquelles ces deux cylindres
ont une intersection non vide.

\begin{proposition}[Catalogue par intersection exacte]
\label{act21:censo:catalogos}
Les queues compatibles sont les entiers de l'intervalle \([a_\chi,b_\chi]\),
avec
\begin{align}
 a_\chi&=\max\left(0,
       \left\lfloor\frac{T_\chi D}{Q}\right\rfloor-P_\chi3^r\right),
       \label{act21:censo:extremo-a}\\
 b_\chi&=\min\left(3^r-1,
       \left\lfloor\frac{(T_\chi+1)D-1}{Q}\right\rfloor-P_\chi3^r\right).
       \label{act21:censo:extremo-b}
\end{align}
Pour les cylindres régionaux fixés, les troncatures n'interviennent pas et l'on obtient
\begin{equation}
\begin{array}{c|r|r|r}
 \chi&a_\chi\bmod729&b_\chi\bmod729&b_\chi-a_\chi+1\\\hline
 \pi&67&262&196\\
 e&584&51&197\\
 \varphi&117&312&196
\end{array}
\label{act21:censo:tabla-catalogos}
\end{equation}
\end{proposition}
\begin{proof}
Écrivons \(N=P_\chi3^r+R\). Deux intervalles semi-ouverts de
\eqref{act21:censo:dos-cilindros} se coupent exactement lorsque
\[
 NQ<(T_\chi+1)D,\qquad (N+1)Q>T_\chi D.
\]
La deuxième inégalité est équivalente à  
\(N\geq\lfloor T_\chi D/Q\rfloor\) ; la première, à  
\(N\leq\lfloor((T_\chi+1)D-1)/Q\rfloor\). Soustraire  
\(P_\chi3^r\) et conserver \(0\leq R<3^r\) prouve les formules.  
L'évaluation par divisions entières des \(T_\chi\) extraits par  
\eqref{act21:censo:prueba-estabilidad} donne le tableau. Comme \(r\geq6\),  
\(P_\chi3^r\equiv0\pmod{729}\), de sorte que le reste de la première  
extrémité s'obtient également directement de  
\(\lfloor T_\chi D/Q\rfloor\bmod729\).
\end{proof}

Les six derniers trits d'une queue forment le mot
\(\nu_6(R\bmod729)\), où \(\nu_6\) est l'écriture ternaire
de longueur six, y compris ses zéros initiaux. D'après le tableau, chaque
catalogue possède moins de \(729\) éléments consécutifs, de sorte que
cette réduction y est injective. Ainsi, les trois catalogues de frontière
sont entièrement décrits, sans choisir les triplets finaux :
\begin{equation}
 \begin{split}
 \mathcal B_\pi&=\{\nu_6(67+i):0\leq i<196\},\\
 \mathcal B_e&=\{\nu_6((584+j)\bmod729):0\leq j<197\},\\
 \mathcal B_\varphi&=\{\nu_6(117+\ell):0\leq\ell<196\}.
 \end{split}
 \label{act21:censo:fronteras-explicitas}
\end{equation}
Exiger que l'extrémité gauche de \(I_{\chi,R}\) appartienne à
\(J_{\chi,T}\) serait une autre condition : elle perdrait le premier cylindre
qui coupe la frontière et donnerait \(195,196,195\). Le recensement utilise l'intersection des cylindres complets, conformément à sa définition.

\subsubsection{Profil déclaré et indicateurs de compatibilité}

Pour \(u,v\in\mathbb F_3^6\), soient
\(g(u,v)=\sum u_iv_i\in\mathbb F_3\),
\(n(u)=g(u,u)\), \(w(u)=\#\{i:u_i\ne0\}\) et
\(h(u,v)=\#\{i:u_i\ne v_i\}\). Les deux dernières fonctions prennent
des valeurs entières ; les deux premières, des valeurs dans \(\mathbb F_3\).
Le profil local du lecteur, en ordre \((\pi,e,\varphi)\), est
\begin{equation}
 \begin{split}
 (g_{\pi e},g_{\pi\varphi},g_{e\varphi})&=(2,2,0),\qquad
 (n_\pi,n_e,n_\varphi)=(0,0,0),\\
 (w_\pi,w_e,w_\varphi)&=(3,3,3),\qquad
 (h_{\pi e},h_{\pi\varphi},h_{e\varphi})=(2,5,3).
 \end{split}
 \label{act21:censo:perfil-local}
\end{equation}
Les mots sont ordonnés comme dans \eqref{act21:censo:fronteras-explicitas}.
Nous utiliserons \([E]\in\{0,1\}\) pour l'indicateur d'une condition \(E\).
Pour chaque paire de canaux \(\chi,\psi\), la matrice rectangulaire
\(A_{\chi\psi}\) a l'entrée \([g(u,v)=g_{\chi\psi}]\).
La matrice \(A^{h}_{\chi\psi}\) multiplie cette entrée par
\([h(u,v)=h_{\chi\psi}]\). Soient \(D_\chi^n\) la matrice diagonale
de \([n(u)=0]\), et \(D_\chi^{nw}\) celle de
\([n(u)=0][w(u)=3]\).

\begin{theorem}[Énumération exhaustive par sommes finies]
\label{act21:censo:teorema-conteo}
Le produit des trois catalogues contient \(7\,567\,952\) triplets.
Les conditions successives de \eqref{act21:censo:perfil-local} laissent
\begin{equation}
 7\,567\,952\longrightarrow275\,794\longrightarrow6235
 \longrightarrow3127\longrightarrow331.
 \label{act21:censo:cadena-local}
\end{equation}
Ces quatre comptages sont, respectivement,
\begin{align}
 N_g&=\operatorname{tr}(A_{\pi e}A_{e\varphi}A_{\varphi\pi}),
 \nonumber\\
 N_{gn}&=\operatorname{tr}(D_\pi^n A_{\pi e}
             D_e^n A_{e\varphi}D_\varphi^n A_{\varphi\pi}),
 \nonumber\\
 N_{gnw}&=\operatorname{tr}(D_\pi^{nw} A_{\pi e}
             D_e^{nw} A_{e\varphi}D_\varphi^{nw} A_{\varphi\pi}),
 \nonumber\\
 N_{gnwh}&=\operatorname{tr}(D_\pi^{nw} A^h_{\pi e}
             D_e^{nw} A^h_{e\varphi}D_\varphi^{nw} A^h_{\varphi\pi}).
 \label{act21:censo:trazas}
\end{align}
Tous les produits et traces de cette formule sont calculés dans les entiers,  
non dans le corps de trois éléments.
\end{theorem}
\begin{proof}
La taille du produit est \(196\cdot197\cdot196\). En développant
une trace de \eqref{act21:censo:trazas}, ses indices parcourent exactement un
triplet de mots. Le terme vaut un s'il satisfait toutes les conditions
représentées et zéro sinon. Chaque triplet est compté une fois.
Les matrices diagonales ajoutent les conditions individuelles ; les facteurs
rectangulaires ajoutent les conditions par paires. Cela démontre que les
traces sont précisément les cardinaux annoncés.

Pour expliciter son évaluation, une fois fixés les indices \(i,j\) des deux premiers canaux, on forme les ensembles d'indices \(\ell\) du troisième qui satisfont chaque condition par paires. On intersecte leurs indicateurs, on ajoute les indicateurs individuels pertinents et on additionne leur cardinal. La somme finale parcourt \(0\leq i<196\), \(0\leq j<197\). La substitution des mots de \eqref{act21:censo:fronteras-explicitas} et du profil \eqref{act21:censo:perfil-local} donne, dans cet ordre, \(275794,6235,3127,331\). Le supplément enregistre les quatre contributions de chacun des \(196\) indices \(i\) et les \(331\) triplets complets ; son programme évalue ces sommes avec des indicateurs binaires exacts et vérifie scalairement chaque survivant. Aucun échantillon ni aucune tolérance n'intervient.
\end{proof}

\subsubsection{Incidence étiquetée et charge d'orientation}

Pour un triplet survivant, soit \(B\in\mathbb F_3^{3\times6}\) sa
matrice de lignes ordonnées. Le second profil exige
\begin{equation}
 \operatorname{rank}_{\mathbb F_3}B=3,\qquad
 w_{\mathrm{col}}(B)=(1,1,2,2,3,0),\qquad
 q_{\partial}(B)=(1,2,2,1,2,0),
 \label{act21:censo:perfil-columnas}
\end{equation}
où \(w_{\mathrm{col}}\) compte les éléments non nuls de chaque
colonne et \(q_{\partial}\) additionne ses entrées modulo trois. Il s'agit d'un profil
de colonnes étiquetées ; les permuter sans transporter le profil définit un autre
lecteur. Notons \(\mathcal F\) l'ensemble des \(331\) triplets
du théorème précédent.

\begin{proposition}[Réduction du recensement au couple terminal]
\label{act21:censo:dos-matrices}
Les sommes d'indicateurs sur \(\mathcal F\), en exigeant successivement
le rang, les poids de colonne et les sommes de colonne, sont
\begin{equation}
 \sum_{B\in\mathcal F}[\operatorname{rank}B=3]=331,
 \qquad N_{\mathrm{rango,pesos}}=18,
 \qquad N_{\mathrm{rango,pesos,sumas}}=2.
 \label{act21:censo:conteo-columnas}
\end{equation}
Avec des indices de catalogue commençant à un, les deux triplets finaux sont
\((120,197,157)\) et \((129,197,148)\), et leurs matrices sont
\[
 B_0=\begin{pmatrix}0&2&0&2&2&0\\0&0&1&2&2&0\\1&0&1&0&1&0\end{pmatrix},
 \qquad
 B_1=\begin{pmatrix}0&2&1&0&2&0\\0&0&1&2&2&0\\1&0&0&2&1&0\end{pmatrix}.
\]
\end{proposition}
\begin{proof}
Dans chaque triplet de l'ensemble déjà énuméré, on calcule le rang par élimination en \(\mathbb F_3\), on compte les éléments non nuls des six colonnes et on les additionne modulo trois. Les produits des indicateurs de \eqref{act21:censo:perfil-columnas} donnent \eqref{act21:censo:conteo-columnas}.
La substitution des deux indices finaux en
\eqref{act21:censo:fronteras-explicitas} donne les matrices exposées. La liste
complète des survivants et les filtres exacts du supplément permettent
de reproduire l'élimination sans supposer ces deux matrices comme domaine.
\end{proof}

L'orientation \(\sigma=(1,1,-1)=(1,1,2)\in\mathbb F_3^3\)
sélectionne la droite
\(\langle\sigma\rangle=\{(0,0,0),(1,1,2),(2,2,1)\}\).
Les charges des lignes de \(B_0,B_1\) sont \((0,2,0)\) et \((2,2,1)\) ;
seule la seconde appartient à cette droite. On retrouve ainsi la dernière étape
du sélecteur de la section~\ref{act21:censo:dos-matrices}, désormais précédée de
la construction et de l'énumération de son domaine complet. La lecture stable
de \(\lfloor3^{1080}x_\chi\rfloor\), effectuée après les filtres,
donne les frontières \texttt{021020}, \texttt{001220}, \texttt{100210},
avec les mêmes rangs \((129,197,148)\). Cette comparaison ultérieure ne définit
pas le profil et ne participe pas aux sommes d'indicatrices.

\subsubsection{Information conservée et portée du résultat}

Le procédé conserve le préfixe de trente trits, la queue compatible, l'ordre régional, la position de chaque trit et l'orientation. La réduction modulo \(729\) conserve ici l'identité de chaque candidat car l'intervalle des queues a une longueur inférieure à \(729\) ; cette injectivité n'est pas affirmée pour une histoire arbitraire. Le recensement de \(331\) démontre également que le profil local isolé ne sélectionne pas une seule matrice : les marges étiquetées et la charge apportent des conditions effectives supplémentaires.

On a donc récupéré une énumération finie intégrale et son sélecteur sur des cylindres de sortie, avec un profil de lecture explicite. La construction du registre signé utilise en outre son lecteur de mémoire et ses canaux, comme cela est spécifié ci-dessous ; la matrice visible ne se confond pas avec ce registre par égalité du nombre de composants. Ce compte ne démontre pas non plus l'universalité du profil, ne génère pas à nouveau les coordonnées régionales et ne décide pas d'une affirmation sur alpha ou sur le continu complet.


% END NUCLEO_ADITIVO_20260921_SELECCION
\section{Registre dodécaphasique et constante rationnelle de phase}
\label{sec:k-registro}

La génération des coordonnées régionales n'épuise pas l'information de l'état APP--TRIT--TPK. Une lecture archimédienne détermine une valeur au moyen de cylindres compatibles ; une autre lecture conserve l'opposition des feuillets, l'orientation de la route et le registre des incidences traversées. Le registre dodécaphasique appartient à cette seconde lecture avant d'intervenir dans la clôture d'alpha. Il n'est donc pas défini en soustrayant une constante physique à une combinaison de valeurs préalablement connues. Il est reconstruit à partir d'une coordonnée signée du même état arithmétique avec mémoire de trajectoire qui produit les coordonnées régionales.

La distinction est nécessaire même lorsque deux objets possèdent douze
coordonnées. Un mot de douze blocs du catalogue, une liste de douze
événements orientés, un vecteur entier signé et un registre dodécaphasique en base mille sont des
objets distincts. Leur comparaison requiert les applications qui les
relient. L'égalité de longueur ne fournit pas ces applications, et la
perte des signes n'est pas une simple simplification de notation.

Ce chapitre maintient la délimitation causale fixée dans les chapitres précédents :
les neuf positions APP des axes horizontal et vertical produisent les
feuillets arithmétiques ; TRIT type le régime et l'orientation ; TPK sélectionne,
transporte, met à jour et prolonge, en conservant le résidu, le quotient, la retenue,
le feuillet, la route, la frontière et la mémoire. À partir de la structure discrète conjointe du
continu s'obtiennent les coordonnées utilisées ici. La chaîne
\(w_6\to w_{12}\to w_{18}\to w_{24}\to w_{30}\to R_{36}\to G_9\)
n'est pas remplacée par un calendrier fini. L'ordre
\(U_t\to\Gamma_9\to K_{\mathrm{ph}}\) n'est pas non plus inversé : la dernière application est une
lecture à mémoire réduite, et non le générateur du registre signé.

\input{sections/revision_k.tex}

\subsection{Extension cyclique et graduation de mémoire}
\label{subsec:k-calendario}

Le calendrier observable contient cent huit positions, organisées en douze fenêtres nonadiques. Avec des indices positifs, soient
\[
 E_m=\{9(m-1)+1,\ldots,9m\},\qquad 1\leq m\leq12.
\]
Le support ordonné \(E_1\sqcup\cdots\sqcup E_{12}\) conserve la position et l'
appartenance à une fenêtre. Le groupe \(C_{108}=\mathbb Z/108\mathbb Z\)
conserve, en revanche, la loi de translation cyclique. Pour le représentant
\(0\leq\widetilde\theta<108\), les coordonnées
\[
 \beta(\theta)=\left[\left\lfloor\frac{\widetilde\theta}{9}\right\rfloor\right]_{12},
 \qquad r(\theta)=1+(\widetilde\theta\bmod9)
\]
séparent le secteur dodécaphasique et la position nonadique intérieure. Aucune d'elles ne conserve à elle seule le nombre total de tours.

\begin{proposition}[La retenue du calendrier]
\label{prop:k-extension}
Les applications
\[
 \iota:C_{12}\longrightarrow C_{108},\quad[b]\longmapsto[9b],
 \qquad p:C_{108}\longrightarrow C_9,\quad[\theta]\longmapsto[\theta]_9
\]
forment une suite exacte non scindée
\[
 0\longrightarrow C_{12}\xrightarrow{\iota}C_{108}
 \xrightarrow{p}C_9\longrightarrow0.
\]
Pour la section ensembliste qui choisit les représentants \(0,\ldots,8\), le défaut d'additivité est la retenue
\[
 c(r,r')=\left[\left\lfloor\frac{\widetilde r+\widetilde r'}9\right\rfloor\right]_{12}.
\]
\end{proposition}

\begin{proof}
Le noyau de \(p\) est formé des multiples de neuf et coïncide avec
l'image de \(\iota\). La réduction est surjective et la multiplication
par neuf est injective sur le domaine indiqué. Si la suite était
scindée, le groupe central serait isomorphe à \(C_{12}\times C_9\), dont l'
exposant est \(\operatorname{mcm}(12,9)=36\) ; l'exposant de
\(C_{108}\) est 108. La contradiction prouve la non-scission. Enfin,
la division euclidienne de \(\widetilde r+\widetilde r'\) par neuf donne la
retenue écrite et vérifie l'identité de la section.
\end{proof}

Le résultat explique une différence qui réapparaîtra dans alpha : revenir à la
phase n'équivaut pas à revenir à l'état. Si \(q_{\mathrm{mem}}\) compte les
tours nonadiques, le calendrier ne conserve que
\(\beta=[q_{\mathrm{mem}}]_{12}\). Après cent huit pas,
\(\beta\) revient, mais \(q_{\mathrm{mem}}\) augmente de douze. Le
registre conserve en outre le cylindre, le livre d'incidences et la mémoire
verticale. La clôture d'une coordinatisation finie ne transforme pas cet historique en un
cycle sans mémoire.

\subsection{Le registre orienté et la coordinatisation intégrale \texorpdfstring{\(1+5+6\)}{1+5+6}}
\label{subsec:k-registro-integral}

Pour une histoire produite par TPK, soient \(\tau_m\) la sous-route sur
\(E_m\), \(\sigma_m\) son orientation, \(\kappa_m\) la retenue de
frontière et \(\Lambda_m\) le livre local des incidences. Le registre est
\[
 \mathcal R_{12}(x)=
 \bigl((E_m,\tau_m,\sigma_m,\kappa_m,\Lambda_m)\bigr)_{m=1}^{12}.
\]
Son domaine est l'espace des histoires arithmétiques avec mémoire de trajectoire, et non l'ensemble des
étiquettes de phase. La projection sur le calendrier oublie précisément certaines
des données dont le lecteur signé a besoin.

Une coordinatisation initiale du registre peut s’écrire au moyen d’événements
orientés. La construction de référence~(sous-section~\ref{subsec:k-registro-integral}) fixe l’ensemble des codes d’événement
\(\{2,4,6,8\}\) et l’orientation sectorielle
\[
 \Sigma_{12}=(+,+,+,-,-,-,+,+,+,-,-,-).
\]
Si \(d_t\) est le code de direction au pas \(t\), on obtient
\[
 e_i=\Sigma_{12}(i+1)
 \#\{t\in E_{i+1}:d_t\in\{2,4,6,8\}\},
 \qquad0\leq i<12.
\]
La formule maintient distincts le recensement non négatif des événements et
l'orientation. Elle n'identifie pas une direction cardinale à un signe TRIT. Le
vecteur \(e\), dont les composantes sont des recensements locaux signés, n'est pas non plus
encore le vecteur \(U\) qui apparaîtra plus bas.

Les positions opposées déterminent, pour \(0\leq i<6\),
\[
 p_i=e_i+e_{i+6},\qquad a_i=e_i-e_{i+6}.
\]
Avec
\[
 s_C=\sum_{i=0}^{5}p_i,\qquad M_i=p_i-p_5\quad(0\leq i<5),
\]
on définit la coordinatisation intégrale
\[
 \mathcal L_{12}(e)
 =(s_C,M_0,\ldots,M_4,a_0,\ldots,a_5).
\]
La décomposition \(1+5+6\) n'est pas une répartition arbitraire des coordonnées :
une coordonnée conserve la charge, cinq comparent les paires par rapport à une
paire de référence et six conservent l'opposition des demi-couronnes.

\begin{lemma}[Inversion de la coordinatisation intégrale]
\label{lem:k-carta-integral}
Un tuple entier \((s_C,M_0,\ldots,M_4,a_0,\ldots,a_5)\) appartient à l'image de \(\mathcal L_{12}\) si et seulement si
\[
 s_C-\sum_{i=0}^{4}M_i\equiv0\pmod6,
 \qquad p_i\equiv a_i\pmod2\quad(0\leq i<6),
\]
où
\[
 p_5=\frac{s_C-\sum_{i=0}^{4}M_i}{6},\qquad p_i=p_5+M_i\quad(i<5).
\]
Dans ce cas, sa préimage est unique et est donnée par
\[
 e_i=\frac{p_i+a_i}{2},\qquad e_{i+6}=\frac{p_i-a_i}{2}.
\]
\end{lemma}

\begin{proof}
La somme des cinq marges est \(s_C-6p_5\), d'où l'on obtient la
division par six. Une fois les \(p_i\) reconstruits, chaque paire d'équations
\(p_i=e_i+e_{i+6}\), \(a_i=e_i-e_{i+6}\) possède la solution indiquée.
Cette solution est entière exactement sous la congruence de parité. Toutes
les opérations s'inversent, de sorte qu'aucune liberté résiduelle ne subsiste.
\end{proof}

Cette réversibilité montre quelle information est conservée ; elle n'autorise pas à
supprimer le reste du livre. Le registre complet possède un domaine plus vaste
que cette coordinatisation des recensements. L'extraction des canaux qui alimentent
\(U\) utilise ce registre complet et ne se déduit pas de la liste
\((e_i)\) par une identification de noms.

% Propuesta separada; no modifica el bloque común del artículo I.
% Procedencia: artículo II, revisión 05, registro_incidencias.tex y
% registro_imagen_integral.tex; aclaraciones de índices de revisiones 06/07.
% Inserción prevista inmediatamente antes de subsec:k-terminal.
\subsection{Évaluation des incidences et détermination des coordonnées transversales}
\label{iii:r25:incidencias}

Le registre dodécaphasique conserve une histoire de visites et de traces
orientées. Ses coordonnées transversales s’obtiennent au moyen d’une
composition d’opérations : évaluation arithmétique des visites,
dénombrement par fenêtres, normalisation décimale et différences cycliques.
La reconstruction du registre à partir de ces différences constitue
une autre opération, dont l’existence et l’unicité seront démontrées ensuite.
Cette séparation permet de préciser tant les données utilisées par
l’extracteur que l’information qu’il transmet à la lecture harmonique.

Soit $\mathscr L$ un registre d’incidences produit par la sélection,
le transport et la mise à jour d’une histoire APP--TRIT--TPK.
Chaque visite conserve son chemin, sa position APP, son feuillet, son instant,
sa multiplicité et sa classe de frontière ; chaque trace conserve ses membres,
son orientation, sa multiplicité et sa longueur dans l’unité déclarée par son
lecteur. Les fenêtres sont
\[
 E_m=\{9(m-1)+1,\ldots,9m\},\qquad 1\leq m\leq12.
\]
L’évaluation du feuillet de somme ou de produit détermine l’entier
$n(v)$ de chaque visite et sa division positive
\[
 n(v)=r(v)+9q(v),\qquad 1\leq r(v)\leq9.
\]
Le résidu se calcule sur cette évaluation. La classe de frontière d’une
visite de résidu $9$ est conservée au moyen de
$\epsilon_\partial(v)=1$ dans la couronne et de $-1$ à l’intérieur.

La présentation d’incidence du registre utilise les trois dénombrements
\begin{align}
 A_m&=\sum_{v\in E_m}w(v)
       \mathbf1_{\{1,3,5,7\}}(r(v)),\label{iii:r25:recuento-a}\\
 C_m&=\sum_{j\geq0}
       \bigl(\nu^-_{120,m}(j)-\nu^+_{120,m}(j)\bigr),
       \label{iii:r25:recuento-c}\\
 V_m&=\sum_{v\in E_m}w(v)
       \mathbf1_{\{9\}}(r(v))\epsilon_\partial(v).
       \label{iii:r25:recuento-v}
\end{align}
Ici, $w(v)$ est la multiplicité d’incidence et
$\nu^\pm_{120,m}(j)$ compte les traces orientées du canal indiqué
de longueur réduite $120(1+3j)$, selon l’énumération du registre.
La famille de traces et leur multiplicité sont des données de l’évaluation,
avec la règle de production qui leur correspond. La somme entière nécessite
un support fini ou une construction qui détermine son sens.
La longueur réduite d’une trace et le nombre d’instants de la
fenêtre sont des grandeurs typées différentes ; l’unité du lecteur
spécifie leur relation. Les symboles $A_m,C_m,V_m$ désignent ici
des dénombrements entiers, distincts de la paire angulaire et de l’opérateur constitutif.

Soit $[x]_{10}\in\{0,\ldots,9\}$ le représentant euclidien.
La normalisation conserve également le quotient dans
\[
 x=[x]_{10}+10\lfloor x/10\rfloor.
\]
On forme le bloc
\begin{equation}
 z_m=100[A_m]_{10}+10[C_m]_{10}+[V_m]_{10},
 \qquad 0\leq z_m\leq999.
 \label{iii:r25:bloque}
\end{equation}
Avec des indices cycliques modulo $12$, soient
\[
 (D_a z)_m=z_m-z_{m+a},\qquad a\in\{3,4\},
 \qquad Q(z)=\sum_{m=1}^{12}z_m.
\]
La composition d’incidence s’exprime par
\begin{equation}
 \mathcal E_{\rm cnt}(\mathscr L)
 =\mathcal D(z):=(D_3z,D_4z,Q(z))\in\mathbb Z^{25}.
 \label{iii:r25:composicion}
\end{equation}
Ses arguments sont les visites et les traces avec leurs règles de dénombrement.
L’évaluation ne reçoit ni $K$, ni $U$, ni $\alpha$, ni les coordonnées
transversales attendues. Pour l’identifier à
$\mathcal E_{108}^{90,120}$, on conserve en outre la composition avec
la famille et les incidences effectives de l’état terminal.
L’inversibilité de $\mathcal D$ détermine une reconstruction ultérieure ;
à elle seule, elle ne sélectionne pas cette famille d’incidences.

\begin{proposition}[Information résiduelle déterminée par l’extracteur]
\label{iii:r25:residuos}
Soient $N,N'\in\mathbb Z^{12\times3}$ deux registres de dénombrements,
ordonnés par $(A,C,V)$. Pour les opérations
\eqref{iii:r25:bloque}--\eqref{iii:r25:composicion},
\[
 \mathcal E_{\rm cnt}(N)=\mathcal E_{\rm cnt}(N')
 \quad\Longleftrightarrow\quad N-N'\in10\mathbb Z^{36}.
\]
Les $36$ résidus sectoriels déterminent exactement la sortie de
$25$ coordonnées. Les quotients et les provenances constituent
une information supplémentaire du registre.
\end{proposition}
\begin{proof}
L’égalité modulo $10$ produit les mêmes blocs et leurs mêmes
coordonnées transversales. Réciproquement, l’égalité des
coordonnées donne $D_3(z-z')=D_4(z-z')=0$. Les pas $3$ et $4$
engendrent le cycle de longueur $12$, donc $z-z'$ est constant.
L’égalité de $Q$ annule cette constante. Le développement de chaque entier
de $0$ à $999$ en trois chiffres décimaux est unique ; on retrouve ainsi
les trois résidus de chaque fenêtre.
\end{proof}

La proposition est valable sur tout le domaine entier indiqué. Elle ne
postule pas que les $36$ classes soient indépendantes sur le sous-ensemble
des registres admissibles produit par TPK.

\subsection{Image entière et inversion des coordonnées transversales}
\label{iii:r25:imagen-seccion}

Pour exprimer l’inversion, on réindexe les blocs par
$i\in\mathbb Z/12\mathbb Z$, de $0$ à $11$. Soient
$(Sz)_i=z_{i+1}$, $D_3=I-S^3$ et $D_4=I-S^4$.
Étant donné un triplet entier $(b,c,q)$, définissons
\begin{equation}
 w_i=c_i-b_{i+1},\qquad
 P_0=0,\quad P_i=\sum_{j=0}^{i-1}w_j\ (1\leq i\leq12),
 \qquad T(w)=\sum_{i=0}^{11}P_i.
 \label{iii:r25:incrementos}
\end{equation}

\begin{theorem}[Caractérisation de l’image entière]
\label{iii:r25:imagen}
Un triplet $(b,c,q)\in\mathbb Z^{25}$ appartient à
$\operatorname{im}\mathcal D$ si et seulement si
\begin{align}
 b_i&=w_i+w_{i+1}+w_{i+2},\qquad 0\leq i<12,
       \label{iii:r25:relaciones}\\
 \sum_{i=0}^{11}w_i&=0,\label{iii:r25:cierre}\\
 q+T(w)&\equiv0\pmod {12}.\label{iii:r25:congruencia}
\end{align}
Son unique préimage entière est
\begin{equation}
 z_i=\frac{q+T(w)}{12}-P_i,\qquad0\leq i<12.
 \label{iii:r25:inversa}
\end{equation}
Les douze relations \eqref{iii:r25:relaciones} et la relation
\eqref{iii:r25:cierre} sont linéairement indépendantes sur
$\mathbb Q$.
\end{theorem}
\begin{proof}
Pour une sortie de $\mathcal D$,
\[
 c_i-b_{i+1}
 =(z_i-z_{i+4})-(z_{i+1}-z_{i+4})=z_i-z_{i+1}.
\]
Trois accroissements consécutifs ont pour somme $b_i$ et leur somme cyclique
totale est nulle. De plus, $P_i=z_0-z_i$, d’où
$q=12z_0-T(w)$. On obtient les conditions et la formule inverse.
Réciproquement, la congruence rend entières les coordonnées de
\eqref{iii:r25:inversa} et la clôture permet de les prolonger
cycliquement. Leurs différences adjacentes sont $w_i$, donc $D_3z=b$
et
\[
 (D_4z)_i=w_i+w_{i+1}+w_{i+2}+w_{i+3}
 =w_i+b_{i+1}=c_i.
\]
La somme donne $Q(z)=q$. La formule détermine également l’unicité.
Enfin, le changement entier inversible
$(b,c)\mapsto(b,w=c-Sb)$ transforme les douze premières relations
en $b=(I+S+S^2)w$ : chacune isole une coordonnée différente de
$b$. La clôture $\sum_iw_i=0$ ajoute une relation indépendante.
\end{proof}

L’image rationnelle a pour dimension $12$. La condition entière ajoute
la congruence d’ordre $12$ de \eqref{iii:r25:congruencia}.
Les onze accroissements $w_0,\ldots,w_{10}$ et la charge $q$ sont
des coordonnées suffisantes, avec $w_{11}=-\sum_{i=0}^{10}w_i$.
Les $25$ coordonnées conservent ces $12$ degrés de liberté
au moyen de $13$ relations linéaires. Si le triplet provient en outre de
\eqref{iii:r25:bloque}, l’inverse appartient à
$\{0,\ldots,999\}^{12}$ et retrouve ses trois résidus par fenêtre.

\subsubsection{Concordance avec la transformation harmonique}

La transformation $\mathcal H_{12}$ agit au moyen de $H_4$ sur les
orbites $(r,r+3,r+6,r+9)$, $r=0,1,2$, où
\[
 H_4=\begin{pmatrix}
 1&1&1&1\\1&1&-1&-1\\1&-1&1&-1\\1&-1&-1&1
 \end{pmatrix},\qquad H_4^2=4I_4.
\]
Pour un triplet compatible, soient
\[
 B_r=\sum_{j=0}^3c_{r+3j},\qquad
 a_0=\frac{q+2B_0+B_1}{3},\quad
 a_1=a_0-B_0,\quad a_2=a_1-B_1.
\]
Le lecteur harmonique a pour blocs
\[
 (\Pi_H(b,c,q))^{(r)}
 =(a_r,b_{r+3}-b_{r+9},b_r+b_{r+6},b_r-b_{r+6}).
\]

\begin{proposition}[Concordance des reconstructions]
\label{iii:r25:hadamard}
Sur $\operatorname{im}\mathcal D$, on a
\[
 \Pi_H\mathcal D=\mathcal H_{12},\qquad
 \mathcal D^{-1}=\tfrac14\mathcal H_{12}\Pi_H.
\]
L’image de $\Pi_H$ restreinte à ce domaine est
\[
 \mathcal L_H=
 \{u\in\mathbb Z^{12}:\mathcal H_{12}u\equiv0\pmod4\}.
\]
\end{proposition}
\begin{proof}
Pour $z=\mathcal D^{-1}(b,c,q)$, soient
$h_r=\sum_{j=0}^3z_{r+3j}$. Le décalage de quatre positions
envoie chaque orbite sur la suivante ; par conséquent, $B_r=h_r-h_{r+1}$.
L’identité $q=h_0+h_1+h_2$ donne $a_r=h_r$.
Sur une orbite, écrivons $(x_0,x_1,x_2,x_3)$ pour ses coordonnées
et $d_j=x_j-x_{j+1}$. On vérifie
\[
 \begin{aligned}
 d_1-d_3&=x_0+x_1-x_2-x_3,\\
 d_0+d_2&=x_0-x_1+x_2-x_3,\\
 d_0-d_2&=x_0-x_1-x_2+x_3.
 \end{aligned}
\]
Ce sont les trois lignes restantes de $H_4x$. L’inversion et la
caractérisation entière découlent de $\mathcal H_{12}^2=4I$.
\end{proof}

La condition de congruence harmonique contrôle la sortie de $\Pi_H$ ;
l’appartenance des canaux à $\operatorname{im}\mathcal D$
nécessite en outre \eqref{iii:r25:relaciones}--\eqref{iii:r25:congruencia}.
Un contrôle négatif exact consiste à prendre $b=0$, $q=0$ et
$c=e_0-e_3$. Les trois sommes orbitales de $c$ s’annulent, de sorte
que $\Pi_H(b,c,q)=0$, mais les relations d’image ne sont pas satisfaites.
La lecture harmonique ne remplace pas le contrôle de compatibilité des
canaux.

\subsection{Conjugaison, covariance et mémoire du registre d’incidence}
\label{iii:r25:transporte}

Soit $\rho(m)=13-m$. Si la conjugaison des incidences inverse
l’ordre des fenêtres, conserve les classes arithmétique et de frontière
des visites et échange les canaux des traces, alors
\[
 (A_m^\dagger,C_m^\dagger,V_m^\dagger)
 =(A_{\rho(m)},-C_{\rho(m)},V_{\rho(m)}).
\]
Avec $c_m=[C_m]_{10}$, la normalisation décimale donne
\begin{equation}
 z_m^\dagger=z_{\rho(m)}
 +100\mathbf1_{\{c_{\rho(m)}\ne0\}}-20c_{\rho(m)}.
 \label{iii:r25:conjugacion}
\end{equation}
En effet, $[-C]_{10}=0$ si $[C]_{10}=0$ et vaut
$10-[C]_{10}$ dans les autres cas. La négation du canal
s’applique avant de reconstruire le bloc décimal.

L’application à une involution concrète conserve également les
extrémités des fenêtres. Pour un chemin
$\gamma=(x_0,\ldots,x_n)$ lu avant l’avancée,
\begin{equation}
 \sum_{k=0}^{n-1}f(x_{n-k})
 =\sum_{k=0}^{n-1}f(x_k)+f(x_n)-f(x_0).
 \label{iii:r25:extremos}
\end{equation}
Les deux sommes contiennent les mêmes termes intérieurs ; la différence
est l’extrémité finale moins l’initiale. Sur les chemins fermés, cette
correction s’annule. Dans les fenêtres ouvertes, elle est incorporée avant de
réduire modulo $10$.

\begin{proposition}[Conséquence d’une réduction périodique]
\label{iii:r25:periodo}
Si l’observable d’une famille fixe de chemins satisfait
$o(t+54)=o(t)$, si ses multiplicités sont conservées sous ce retour
et si le même lecteur local agit dans des fenêtres de neuf instants,
alors
\[
 (A,C,V)_{m+6}=(A,C,V)_m,\qquad z_{m+6}=z_m.
\]
\end{proposition}
\begin{proof}
La translation $t\mapsto t+54$ envoie chaque fenêtre bijectivement
sur celle située six positions plus loin et conserve chaque incidence
comptée et sa multiplicité. L’égalité se transmet aux dénombrements
et à leurs normalisations décimales.
\end{proof}

Une présentation à douze blocs distincts conserve donc
une information qui ne se factorise pas par cette observable de période $54$ :
sélection des incidences, orientation, extrémités, multiplicité ou
mémoire. La proposition délimite l’effet de cette réduction ; elle
n’attribue pas une période $54$ à l’histoire complète.

\subsubsection{Détermination sur les orbites et accumulation d’événements}

Fixer uniquement $Q(z)=q$ laisse une classe affine du réseau
\[
 \mathcal Z_0=\{v\in\mathbb Z^{12}:\sum_i v_i=0\}
 =\bigoplus_{i=0}^{10}\mathbb Z(e_i-e_{11}).
\]
Par exemple, $100\mathbf1+e_0$ et $100\mathbf1+e_1$ ont pour
charge $1201$ et satisfont les contrôles d’image au moyen de leurs
canaux respectifs. Ce sont des témoins des équations de compatibilité,
sans attribution d’appartenance aux histoires terminales HMT.

\begin{proposition}[Covariance cyclique]
\label{iii:r25:covariancia}
Soit $\rho$ la translation des fenêtres d’un registre $R$, de
période minimale $p\mid12$, et soit $S$ la translation de ses
coordonnées. Un lecteur covariant sur cette orbite est déterminé
par $z\in\operatorname{Fix}(S^p)$. La condition $Q(z)=q$
admet des solutions entières si et seulement si $12/p$ divise $q$ ;
lorsqu’elles existent, elles forment une classe affine de rang $p-1$.
\end{proposition}
\begin{proof}
La règle $F(\rho^jR)=S^jz$ est bien définie exactement si
$S^pz=z$. Ces vecteurs répètent $12/p$ fois un bloc de
longueur $p$, et leur charge est $\frac{12}{p}\sum_{i=0}^{p-1}z_i$.
La formule prouve la divisibilité et le rang. La covariance des
canaux découle de $D_aS=SD_a$ et $QS=Q$.
\end{proof}

Dans la coordinatisation harmonique, l’action correspondante est
$\widehat S_H=\mathcal H_{12}S\mathcal H_{12}/4$.
Si l’origine est marquée, l’action transporte également cette marque.
La période de l’orbite indiquée ici ne s’identifie pas à celle de
l’histoire enrichie. Pour une lecture initiale $F_{108}$ déjà
définie, la naturalité sous troncature s’exprime au moyen de
$F_N=F_{108}\circ\tau_{N,108}$ ; cette identité conserve
l’évaluation initiale sans sélectionner sa règle de production.

\begin{proposition}[Composition additive des contributions d’incidence]
\label{iii:r25:eventos}
Supposons que l’état initial et chaque événement $a_t$ produisent,
au moyen de règles préalablement fixées sur le registre, des contributions
$(\delta w_t,\delta q_t)$ telles que
\[
 \sum_i\delta w_{t,i}=0,
 \qquad \delta q_t+T(\delta w_t)\equiv0\pmod {12}.
\]
Avec des conditions initiales compatibles $(w^{(0)},q^{(0)})$,
la mise à jour par somme détermine un unique registre à chaque étape.
Son accroissement est
\[
 \delta z_{t,i}=
 \frac{\delta q_t+T(\delta w_t)}{12}-P_i(\delta w_t).
\]
La construction respecte la concaténation des segments et conserve
la lecture d’un préfixe sous les prolongements ultérieurs.
\end{proposition}
\begin{proof}
Les conditions de compatibilité se conservent par somme et
l’inverse \eqref{iii:r25:inversa} est additif en $(w,q)$.
La récurrence fixe chaque état. La somme sur des segments concaténés
est la somme de leurs contributions ; un préfixe reçoit les mêmes
termes avant et après le prolongement.
\end{proof}

Cette proposition donne un critère suffisant pour l’accumulation
additive, non une exigence de linéarité de tout lecteur TPK.
Lorsque les événements se transportent par des actions non triviales, leurs
contributions s’expriment dans la coordinatisation terminale au moyen du cocycle
affine correspondant. Les règles d’événement et l’état initial
conservent dans les deux cas leur provenance d’incidence.

Le développement précédent détermine le passage
\[
 \mathscr L\longmapsto(A,C,V)\longmapsto z
 \longmapsto(D_3z,D_4z,Q(z))\longmapsto\mathcal H_{12}z
\]
avec les conditions de chaque application. Son domaine est distinct du
registre régional de l’appendice~\ref{iii:nucleo-finito} : les
équations $X_aL_a=Y_a$ y déterminent les relèvements à partir de
transitions ternaires. Ici, les incidences des visites et des traces
déterminent des blocs entiers et leurs canaux. L’unicité d’une
reconstruction ne remplace pas la production du registre qu’utilise
l’autre. Les deux conservent leurs primitives, opérations et preuves dans
la chaîne APP--TRIT--TPK.


\subsection{Deux lectures de l'état terminal}
\label{subsec:k-terminal}

Désignons par \(x_{\mathrm{term}}\) l'état produit par la composition de sélection, de transport et de mise à jour jusqu'à la coupe terminale. L'écriture opératorielle est
\[
 x_{\mathrm{term}}=
 (\mathcal U_{108}\circ\cdots\circ\mathcal U_1)(x_{\mathrm{seed}}),
 \qquad \mathcal U_t=\mathsf{Upd}_t\circ\mathsf{Tra}_t\circ\mathsf{Sel}_t.
\]
Ses deux lectures pertinentes sont
\[
 x_{\mathrm{term}}
 \xrightarrow{(\pi_{\mathrm{term}},R_{\mathrm{sgn}})}
 (B_{\mathrm{term}},U).
\]
La matrice \(B_{\mathrm{term}}\) retient la lecture ternaire visible ; le vecteur \(U\) conserve une autre coordonnée de la même histoire. On n'insère pas entre eux une flèche \(B_{\mathrm{term}}\to U\).

Le sélecteur terminal de la construction de référence~(sous-section~\ref{subsec:k-terminal}) opère sur des catalogues
de tailles 196, 197 et 196. Les signatures locales réduisent leurs
\(196\cdot197\cdot196=7\,567\,952\) triplets à 331, et la marge des
colonnes \(q_{\partial}=(1,2,2,1,2,0)\) laisse deux matrices :
\[
 B_0=\begin{pmatrix}0&2&0&2&2&0\\0&0&1&2&2&0\\1&0&1&0&1&0\end{pmatrix},
 \qquad
 B_1=\begin{pmatrix}0&2&1&0&2&0\\0&0&1&2&2&0\\1&0&0&2&1&0\end{pmatrix}.
\]
L'orientation des lignes est \(\sigma=(1,1,-1)=(1,1,2)\) dans
\(\mathbb F_3^3\). Les charges des lignes des candidates sont
\((0,2,0)\) et \((2,2,1)\). Seule la seconde appartient à
\(\langle\sigma\rangle=\{(0,0,0),(1,1,2),(2,2,1)\}\) ; par conséquent,
\(B_{\mathrm{term}}=B_1\). Cette vérification prouve la dernière étape du
sélecteur à partir du recensement déclaré. Elle ne transforme pas la marge des colonnes
en une conséquence de la charge des lignes, et ne remplace pas l'énumération
précédente par l'inspection de ces deux matrices.

Le lecteur signé possède la composition typée
\begin{equation}
 R_{\mathrm{sgn}}
 =\Pi_H\circ\mathcal E_{108}^{90,120}\circ\mathcal R_{12}:
 \mathcal X_{\mathrm{TPK}}^{\mathrm{term,mem}}
 \longrightarrow\mathbb Z^{12}.
 \label{eq:k-lector-firmado}
\end{equation}
Le registre des canaux détermine
\[
 \mathcal E_{108}^{90,120}(\mathcal R_{12}(x_{\mathrm{term}}))
 =(b^{90},b^{120},Q_{\mathrm{TPK}}),
\]
avec les coordonnées
\begin{align*}
 b^{90}={}&(-495,-116,-684,108,601,-90,\\
            &\hspace{19mm}-173,-88,313,560,-397,461),\\
 b^{120}={}&(-425,-281,-481,671,-255,30,\\
             &\hspace{19mm}475,-543,680,251,6,-128),\\
 Q_{\mathrm{TPK}}={}&6263.
\end{align*}
Ces quantités sont utilisées comme coordonnées du registre préalable, non
comme paramètres ajustés au moyen d’alpha. La reconstruction démontrée
dans les paragraphes suivants commence exactement à cette sortie typée.
La trace matérielle de l’extracteur antérieur et la portée de sa transcription
autonome sont conservées dans la note technique de provenance ; l’arithmétique
ultérieure n’est pas présentée comme un substitut à cette trace.

\subsection{Du registre des canaux au vecteur signé}
\label{subsec:k-pi-h}

Le lecteur harmonique \(\Pi_H\) peut s'écrire sans utiliser ni \(K\) ni
alpha. Pour éviter de confondre ses sommes orbitales avec les blocs d'alpha,
nous les noterons \(s_1,s_2,s_3\). Soient
\[
 B_r=\sum_{j=0}^{3}b^{120}_{r+3j},\qquad
 d_{r,j}=b^{90}_{r+3j},\qquad1\leq r\leq3,
\]
et
\begin{equation}
 s_1=\frac{Q_{\mathrm{TPK}}+2B_1+B_2}{3},\qquad
 s_2=s_1-B_1,\qquad s_3=s_2-B_2.
 \label{eq:k-sumas-armonicas}
\end{equation}
La division par trois appartient au lecteur des trois orbites. Son caractère entier est vérifié sur les sorties compatibles du registre ; il n'est pas supposé pour un élément arbitraire de \(\mathbb Z^{25}\).

Les trois blocs signés sont
\begin{equation}
 U^{(r)}=(s_r,d_{r,1}-d_{r,3},d_{r,0}+d_{r,2},d_{r,0}-d_{r,2}).
 \label{eq:k-bloques-firmados}
\end{equation}
La substitution entière donne
\[
 (B_1,B_2,B_3)=(972,-1073,101),\qquad
 (s_1,s_2,s_3)=(2378,1406,2479),
\]
et, par conséquent,
\begin{align*}
 U^{(1)}&=(2378,-452,-668,-322),\\
 U^{(2)}&=(1406,998,-204,-28),\\
 U^{(3)}&=(2479,-551,-371,-997).
\end{align*}
En entrelaçant les colonnes, on obtient
\begin{equation}
 U=(2378,1406,2479,-452,998,-551,
       -668,-204,-371,-322,-28,-997).
 \label{eq:k-u-firmado}
\end{equation}
La présence de valeurs négatives et de composantes supérieures à 999
montre que \(U\) n'est pas un mot de chiffres en base mille. Ce n'est pas non plus
\(U_6\Vert U_6\), concaténation du catalogue de période six. La différence
concerne le domaine et l'information conservée, et non seulement les valeurs.

L'ordre constructif jusqu'ici est
\[
 x_{\mathrm{term}}\longrightarrow\mathcal R_{12}
 \longrightarrow(b^{90},b^{120},Q_{\mathrm{TPK}})
 \longrightarrow U.
\]
L'inversion qui suit produit le registre dodécaphasique. Le fait que nous puissions ensuite récupérer les canaux à partir du registre dodécaphasique constitue un contrôle de réversibilité ; il n'inverse pas cet ordre causal.

\subsection{Hadamard par orbites et reconstruction entière}
\label{subsec:k-hadamard}

Dans le cycle à douze positions, le pas trois détermine trois orbites :
\[
 (1,4,7,10),\qquad(2,5,8,11),\qquad(3,6,9,12).
\]
La séparation de quatre secteurs de chaque orbite est celle qui, une fois
l'origine et l'orientation fixées, admet la lecture angulaire de quatre-vingt-dix degrés. La
transformation n'agit pas sur trois quadruplets contigus de l'ordre naturel.

Soit
\[
 H_4=\begin{pmatrix}
 1&1&1&1\\1&1&-1&-1\\1&-1&1&-1\\1&-1&-1&1
 \end{pmatrix}.
\]
Si \(P_{\mathrm{orb}}\) réordonne selon les trois orbites indiquées, nous définissons
\[
 \mathcal H_{12}=P_{\mathrm{orb}}^{-1}
 (H_4\oplus H_4\oplus H_4)P_{\mathrm{orb}}.
\]

\begin{lemma}[Inversion de Hadamard et intégralité]
\label{lem:k-hadamard}
On a
\[
 H_4^2=4I_4,\qquad H_4^{-1}=\frac14H_4,\qquad\det H_4=-16.
\]
Pour \(u\in\mathbb Z^4\), la préimage \(H_4^{-1}u\) est entière si et
seulement si \(H_4u\equiv0\pmod4\).
\end{lemma}

\begin{proof}
Les lignes sont orthogonales et le carré de leur norme vaut quatre ; la matrice est symétrique. Cela donne \(H_4^2=4I_4\) et l'inverse. Une élimination entière, ou le développement du déterminant, donne le signe négatif et la valeur \(-16\). La condition de caractère entier est exactement la divisibilité par quatre des quatre coordonnées de \(H_4u\).
\end{proof}

\begin{definition}[Transformation entière dodécaphasique]
Le sous-réseau des registres signés et sa coordinatisation entière sont
\[
 \mathcal L_H=\{u\in\ZZ^{12}:\mathcal H_{12}u\equiv0\pmod4\},
 \qquad \mathscr K(u)=\tfrac14\mathcal H_{12}u.
\]
L'application \(\mathscr K:\mathcal L_H\to\ZZ^{12}\) est un isomorphisme
de groupes abéliens, d'inverse \(k\mapsto\mathcal H_{12}k\).
Le registre canonique \(K\) est l'image du registre signé généré
par les opérations précédentes, avec une origine de phase et une orientation fixées.
\end{definition}

\begin{proof}
La congruence définit exactement l'intégralité de \(\mathscr K(u)\).
L'identité \(\mathcal H_{12}^2=4I\) démontre les deux compositions
inverses. En particulier, tout \(k\in\ZZ^{12}\) a pour préimage
\(\mathcal H_{12}k\in\mathcal L_H\).
\end{proof}

\begin{proposition}[Registre entier canonique]
\label{prop:k-sello}
L'inversion \(K^{(r)}=H_4U^{(r)}/4\) des blocs de
\eqref{eq:k-bloques-firmados} produit
\begin{align*}
 K^{(1)}&=(234,729,621,794),\\
 K^{(2)}&=(543,659,58,146),\\
 K^{(3)}&=(140,824,914,601).
\end{align*}
Dans l'ordre naturel, le registre dodécaphasique est
\begin{equation}
 K=(234,543,140,729,659,824,621,058,914,794,146,601).
 \label{eq:k-vector}
\end{equation}
C'est l'unique préimage rationnelle de \(U\), et elle appartient à \(\{0,\ldots,999\}^{12}\).
\end{proposition}

\begin{proof}
Les trois produits avant division sont
\begin{align*}
 H_4U^{(1)}&=(936,2916,2484,3176),\\
 H_4U^{(2)}&=(2172,2636,232,584),\\
 H_4U^{(3)}&=(560,3296,3656,2404).
\end{align*}
Toutes les coordonnées sont divisibles par quatre. Diviser et défaire la
permutation des orbites donne \eqref{eq:k-vector}. L'unicité procède de l'
inversibilité rationnelle, et non d'une recherche parmi des candidats proches d'
une valeur physique.
\end{proof}

Le caractère entier peut également se formuler en termes de réseaux. Les diviseurs déterminantaux de \(H_4\) sont \(1,2,4,16\) : les entrées ont pour plus grand commun diviseur un ; les mineurs d'ordre deux sont pairs et l'un d'eux est de module deux ; ceux d'ordre trois sont des multiples de quatre et l'un d'eux est de module quatre. On en déduit
\[
 \operatorname{SNF}(H_4)=\operatorname{diag}(1,2,2,4),
\]
et pour la transformation globale,
\[
 \operatorname{coker}\mathcal H_{12}
 \cong(\mathbb Z/2\mathbb Z)^6\oplus(\mathbb Z/4\mathbb Z)^3.
\]
L'indice de l'image est \(16^3=4096\). Ainsi, tout vecteur signé
entier n'admet pas nécessairement une préimage entière ; le vecteur obtenu satisfait les
congruences précises exigées par l'inversion. Le facteur \(1/4\) est
imposé par la matrice et n'a pas le statut d'un coefficient libre.

\subsection{Différences cycliques et composante uniforme}
\label{subsec:k-transversal}

Avec des indices cycliques modulo douze, soient
\[
 (D_3z)_m=z_m-z_{m+3},\qquad(D_4z)_m=z_m-z_{m+4},
 \qquad Q(z)=\sum_{m=1}^{12}z_m.
\]
Le premier opérateur compare des secteurs séparés par trois positions ; le
second, par quatre. Les noms de canal 90 et 120 se réfèrent ici à
ces séparations dans le cycle orienté, et non à une identification avec un
opérateur physique d'un autre domaine.

\begin{proposition}[Complétude du système transversal]
\label{prop:k-transversal}
L'opérateur
\[
 \mathcal D=(D_3,D_4,Q):\mathbb Z^{12}\longrightarrow\mathbb Z^{25}
\]
a un rang égal à douze et des facteurs invariants non nuls
\[
 \operatorname{diag}(1,\ldots,1,12),
\]
avec onze unités. En particulier, leurs sorties déterminent un unique registre dodécaphasique.
\end{proposition}

\begin{proof}
Si \(D_3z=D_4z=0\), le vecteur est constant le long des deux pas.
Comme \(3\) et \(4\) engendrent le cycle modulo douze,
\(\ker D_3\cap\ker D_4=\langle\mathbf1\rangle\). La charge totale
élimine cette dernière liberté parce que \(Q(\mathbf1)=12\).

Pour l'affirmation entière, les lignes de \((D_3,D_4)\) sont des incidences orientées d'un graphe connexe. Un arbre couvrant fournit onze facteurs invariants unitaires. Tout mineur d'ordre douze doit contenir la ligne \(Q\). En ajoutant les onze premières colonnes à la dernière, celle-ci s'annule dans les lignes d'incidence et vaut douze dans la ligne \(Q\). Tous les mineurs maximaux sont divisibles par douze et l'arbre en produit un de module exactement douze. Le dernier facteur est donc douze.
\end{proof}

La vérification sur \eqref{eq:k-vector} donne
\[
 D_3K=b^{90},\qquad D_4K=b^{120},\qquad Q(K)=6263.
\]
Les formules \eqref{eq:k-sumas-armonicas} et
\eqref{eq:k-bloques-firmados} reconstruisent à leur tour \(U\) à partir de ces
différences. On ferme ainsi un triangle de représentations exactes du
même objet : registre transversal, vecteur signé et registre dodécaphasique. L'égalité
des résultats ne fait pas de leurs codomaines un seul et même domaine ; elle démontre que
les applications déclarées restituent l'information pertinente.

\subsection{Deux lectures rationnelles : \texorpdfstring{\(\kappa_{36}\)}{kappa36} et \texorpdfstring{\(\kappa_{\mathrm{per}}\)}{kappa périodique}}
\label{subsec:k-kappas}

Une fois \(K\) construit, la coordinatisation en base mille permet deux opérations
distinctes. La première lit un seul tour ; la seconde prolonge périodiquement le registre
dodécaphasique. Soit \(B=1000\) et définissons son entier de concaténation
\begin{equation}
 N_K=\sum_{m=1}^{12}K_mB^{12-m}
 =234543140729659824621058914794146601.
 \label{eq:k-numerador}
\end{equation}
Les zéros initiaux d'un bloc, comme celui de 058, font partie de son écriture
à trois chiffres. Ils ne changent pas sa valeur entière, mais permettent de concaténer
sans ambiguïté les douze blocs.

\begin{definition}[Lectures finie et périodique du registre dodécaphasique]
\label{def:k-kappas}
La lecture d'un tour est
\[
 \kappa_{36}=\sum_{m=1}^{12}K_mB^{-m}=\frac{N_K}{B^{12}}.
\]
Le prolongement périodique est
\[
 K_j^{\mathrm{per}}=K_{1+((j-1)\bmod12)},\qquad
 \kappa_{\mathrm{per}}=\sum_{j\geq1}K_j^{\mathrm{per}}B^{-j}.
\]
\end{definition}

\begin{proposition}[Valeur et période exactes du registre dodécaphasique prolongé]
\label{prop:k-periodo}
On a
\begin{equation}
 \kappa_{\mathrm{per}}=\frac{N_K}{B^{12}-1},\qquad
 \kappa_{\mathrm{per}}-\kappa_{36}
 =\frac{N_K}{B^{12}(B^{12}-1)}
 =B^{-12}\kappa_{\mathrm{per}}.
 \label{eq:k-diferencia-kappas}
\end{equation}
Le mot infini du registre dodécaphasique a pour période minimale douze en base mille et pour période minimale trente-six en base dix.
\end{proposition}

\begin{proof}
La somme du premier tour est \(N_KB^{-12}\) ; chaque nouveau tour
multiplie sa contribution par \(B^{-12}\). La série géométrique produit
\(N_K/(B^{12}-1)\), et la soustraction donne la seconde identité.

Les douze composantes de \(K\) sont distinctes. Une période propre du cycle
identifierait deux d'entre elles ; la période minimale des blocs est donc
douze. Le développement décimal est canonique, car le registre dodécaphasique n'a pas de queue
de blocs 999. Si sa période décimale minimale était \(d\), elle diviserait
36. Un regroupement par trois produirait une période de blocs
\(d/\gcd(d,3)\). Aucun diviseur propre de 36 ne donne douze par cette
opération. La période décimale minimale est donc 36.
\end{proof}

Les deux lectures sont rationnelles, mais ne sont pas le même rationnel. La première se termine après douze blocs ; la seconde répète le tour. En outre, \(0<\kappa_{\mathrm{per}}-\kappa_{36}<B^{-12}\). Leur proximité ne permet pas de substituer l'une à l'autre dans une identité exacte. En particulier, la clôture finie d'alpha utilise \(\kappa_{36}\), tandis que la somme de la section prolongée utilise \(\kappa_{\mathrm{per}}\).

\input{sections/k_reversibilidad.tex}

\subsection{Périodicité et conoyau de l'opérateur de transport}
\label{subsec:k-clase-acarreo}

La périodicité que nous venons de démontrer appartient au registre dodécaphasique. Elle n'implique
pas la périodicité des coordonnées régionales, des retenues de la clôture
ni d'alpha. Il est utile de préciser cette séparation au moyen d'un second
quotient, purement entier.

Soit \(S\) le décalage cyclique de douze coordonnées, d'orientation fixée, et soit \(b\geq2\) une base entière. L'opérateur de retenue périodique est \(\partial_b=S-bI\). Si une identité périodique s'écrit
\[
 A=P+E-\Phi-K+\partial_bC,
\]
alors les transformations
\[
 K\longmapsto K+\partial_bh,\qquad C\longmapsto C+h
\]
conservent son membre de droite. Cette invariance compare les représentants d'une
classe de retenue ; elle n'autorise pas à modifier le registre dodécaphasique déjà
sélectionné par le registre.

\begin{proposition}[Quotient de la retenue périodique]
\label{prop:k-cociente-acarreo}
Avec l'orientation précédente,
\[
 \operatorname{coker}(S-bI)\cong\mathbb Z/(b^{12}-1)\mathbb Z.
\]
\end{proposition}

\begin{proof}
Le module cyclique se présente sous la forme
\(\mathbb Z[t]/(t^{12}-1)\). Imposer \(S=bI\) équivaut à imposer
\(t=b\) et laisse la relation \(b^{12}-1=0\). L'évaluation des
coefficients en \(b\), avec l'ordre de concaténation choisi, représente
la classe résultante.
\end{proof}

Le dénominateur \(B^{12}-1\) de \(\kappa_{\mathrm{per}}\) et ce quotient proviennent de la même longueur cyclique, mais leurs objets restent distincts : l'un est une évaluation rationnelle et l'autre un quotient d'un module de retenues. La clôture finie d'alpha sera une chaîne orientée à frontière droite. Son prolongement transportera une frontière distante compatible, et non une retenue identifiée par décret entre la première et la treizième position.

Le chapitre suivant est ainsi préparé. Le registre dodécaphasique a été fixé avant
de résoudre la récurrence d'alpha ; sa reconstruction par Hadamard est
entière et unique ; ses différences transversales restituent le registre ; et
ses deux évaluations rationnelles possèdent des domaines et des usages distincts. Aucun
de ces faits n'utilise alpha comme entrée, et aucun ne décide encore du
type arithmétique de la sortie de la clôture complète.

\section{Dérivation de la constante de structure fine}
\label{sec:alpha}

La constante de structure fine reçoit dans ce développement une détermination arithmétique qui réunit les trois coordonnées régionales et le registre intégral dodécaphasique. L'opération combine ses blocs avec orientation fixée et transporte les quotients de la normalisation positionnelle. Le paramètre résultant exprime la compatibilité de cette différence orientée avec la prolongation des cylindres. Sa valeur et les données d'incidence utilisées dans sa construction resteront liées pendant le développement exceptionnel.

L'ordre interne est important. Les coordonnées \(\pi\), \(e\) et \(\varphi\), déjà construites dans la section précédente, sont des coordonnées
précédentes du même état arithmétique avec mémoire de trajectoire, et peuvent agir légitimement comme
entrées d'un lecteur ultérieur. Elles ne sont pas des constantes conventionnelles
injectées pour sélectionner alpha. La génération corrélée des quatre coordonnées partage une origine et
une compatibilité coinductive, sans que cela impose d'évaluer ses quatre
composantes en un seul pas.

Les deux voies expriment des fonctions mathématiques différentes du même paramètre.
La première réalise une clôture entière par division euclidienne en base
\(1000\). La seconde étudie l’équilibre analytique des vacances au moyen
de la résolvante \(E_{21/22}\) et de ses prolongements gradués. Une voie
détermine des blocs et des quotients ; l’autre détermine une racine par une
équation et un domaine d’unicité. Leur identification est formulée ensuite
dans les cylindres communs à toute profondeur. L’exposé distingue
ces trois résultats : construction positionnelle, détermination analytique
et compatibilité des deux.

\subsection{Première voie : détermination positionnelle par clôture entière}
\subsubsection{Domaine et orientation de la clôture entière}
\label{subsec:alpha-dominio}

Soient \(P_j,E_j,\Phi_j\in\{0,\ldots,999\}\) les blocs canoniques des
trois coordonnées régionales, et prolongeons le registre dodécaphasique par
\[
 K_j=K_{1+((j-1)\bmod12)}.
\]
L'orientation du canal est \((1,1,-1)\). La contribution du registre dodécaphasique se soustrait selon la clôture dodécaphasique déjà typée. Pour \(B=1000\), l'opération locale distingue le bilan signé
\[
 Z_j:=P_j+E_j-\Phi_j-K_j
\]
de la somme qui incorpore le quotient entrant. La normalisation est
\begin{equation}
 s_j=Z_j+c_{j+1},\qquad
 A_j=s_j-B\left\lfloor\frac{s_j}{B}\right\rfloor,
 \qquad c_j=\left\lfloor\frac{s_j}{B}\right\rfloor.
 \label{eq:alpha-recurrencia}
\end{equation}
Ainsi, \(A_j\in\{0,\ldots,B-1\}\), même lorsque \(s_j\) est
négatif. La retenue qui entre depuis la droite appartient au domaine de
l'opération ; elle ne peut pas être omise car le tableau s'imprime de gauche
à droite.

\begin{definition}[Clôture de profondeur finie]
\label{def:alpha-finita}
Pour une profondeur \(N\) et une frontière entière \(t\), l'application
\(\mathcal C_N\) reçoit
\((P_{1:N},E_{1:N},\Phi_{1:N},K_{1:N},c_{N+1}=t)\) et applique
\eqref{eq:alpha-recurrencia} dans l'ordre \(N,N-1,\ldots,1\). Sa sortie
est la paire de mots
\[
 (A_{1:N},c_{1:N+1}).
\]
\end{definition}

La division euclidienne ne possède pas un choix de signe résiduel : le reste se
fixe en \(\{0,\ldots,999\}\) et le quotient reste déterminé. Par
exemple, \(-431=569-1000\) et \(-1200=800-2\cdot1000\). Les
retenues négatives de la clôture ne sont pas des anomalies ; elles sont l'information
entière nécessaire pour maintenir l'égalité orientée.

\subsubsection{Fenêtre dodécaphasique et conditions de retenue}
\label{subsec:alpha-ventana}

Les douze premières triades obtenues sont
\begin{align*}
 P={}&(141,592,653,589,793,238,462,643,383,279,502,884),\\
 E={}&(718,281,828,459,045,235,360,287,471,352,662,497),\\
 \Phi={}&(618,033,988,749,894,848,204,586,834,365,638,117).
\end{align*}
Dans le canal de fermeture, les cinq régions qui déterminent la coordonnée \(P\) conservent leurs antécédents distincts. Partager cette lecture ne les identifie pas en tant que régions ou en tant qu'histoires.

Avec frontière \(c_{13}=0\), la récurrence produit le tableau suivant.
Chaque entrée est montrée, la retenue entrante, le total avant normalisation,
le bloc obtenu et la retenue sortante. Le tableau est un calcul entier
reproductible, non une liste de blocs sélectionnés par leur ressemblance avec
un objectif.

\begin{table}[htbp]
\centering
\small
\setlength{\tabcolsep}{3.5pt}
\begin{tabular}{r|rrrr|rr|rr}
\hline
\(j\)&\(P_j\)&\(E_j\)&\(\Phi_j\)&\(K_j\)&\(c_{j+1}\)&\(s_j\)&\(A_j\)&\(c_j\)\\
\hline
1&141&718&618&234&0&7&007&0\\
2&592&281&033&543&0&297&297&0\\
3&653&828&988&140&\(-1\)&352&352&0\\
4&589&459&749&729&\(-1\)&\(-431\)&569&\(-1\)\\
5&793&045&894&659&\(-2\)&\(-717\)&283&\(-1\)\\
6&238&235&848&824&\(-1\)&\(-1200\)&800&\(-2\)\\
7&462&360&204&621&0&\(-3\)&997&\(-1\)\\
8&643&287&586&058&\(-1\)&285&285&0\\
9&383&471&834&914&\(-1\)&\(-895\)&105&\(-1\)\\
10&279&352&365&794&0&\(-528\)&472&\(-1\)\\
11&502&662&638&146&0&380&380&0\\
12&884&497&117&601&0&663&663&0\\
\hline
\end{tabular}
\caption{Clôture initiale avec frontière droite nulle. La dépendance se propage de la dernière ligne à la première.}
\label{tab:alpha-ventana}
\end{table}

Le mot des retenues est
\[
 (c_1,\ldots,c_{13})=(0,0,0,-1,-1,-2,-1,0,-1,-1,0,0,0),
\]
et la lecture de la fenêtre est le rationnel
\begin{equation}
 \alpha_{12}^{(0)}
 =\frac{7297352569283800997285105472380663}{10^{36}}
 =0.007297352569283800997285105472380663.
 \label{eq:alpha-ventana-cero}
\end{equation}
L'exposant rappelle la frontière choisie. Cette fenêtre est exacte dans son
domaine fini ; on ne l'identifie pas d'emblée à la troncature de la
section infinie. La différence entre les deux lectures deviendra visible lors du
transport de la retenue de la queue lointaine.

\subsubsection{Télescopage et réversibilité dans une fibre de frontière}
\label{subsec:alpha-telescopia}

La forme locale de l'égalité est
\begin{equation}
 P_j+E_j-\Phi_j-K_j-A_j=Bc_j-c_{j+1}.
 \label{eq:alpha-identidad-local}
\end{equation}
Le terme de retenue est un lien entre des positions contiguës. Si on le supprime coordonnée par coordonnée, on obtient une égalité fausse. Seule une évaluation qui conserve ses extrémités permet de le télescoper.

\begin{proposition}[Identité entière de profondeur arbitraire]
\label{prop:alpha-telescopia}
Pour \(\iota_N(z)=\sum_{j=1}^{N}z_jB^{N-j}\), toute exécution de
\(\mathcal C_N\) satisfait
\begin{equation}
 \iota_N(P)+\iota_N(E)-\iota_N(\Phi)-\iota_N(K)-\iota_N(A)
 =B^Nc_1-c_{N+1}.
 \label{eq:alpha-telescopia}
\end{equation}
\end{proposition}

\begin{proof}
Multiplier \eqref{eq:alpha-identidad-local} par \(B^{N-j}\) et ajouter  
produit du côté droit
\[
 \sum_{j=1}^{N}(Bc_j-c_{j+1})B^{N-j}=B^Nc_1-c_{N+1}.
\]
Chaque retenue intérieure apparaît deux fois avec un signe opposé ; seuls les deux extrêmes survivent.
\end{proof}

Pour le tableau ~\ref{tab:alpha-ventana}, les deux extrémités sont nulles. Si
\(p_{12},e_{12},f_{12}\) sont les trois sommes fractionnaires tronquées,
\[
 \alpha_{12}^{(0)}=p_{12}+e_{12}-f_{12}-\kappa_{36}.
\]
La constante du registre dodécaphasique qui apparaît est \(\kappa_{36}\), non
\(\kappa_{\mathrm{per}}\). Le domaine de cette identité contient
des troncatures et une frontière finie ; sa forme n'autorise pas à effacer les deux
conditions en passant à la section complète.

L'inversion locale est aussi exacte :
\[
 K_j=P_j+E_j-\Phi_j+c_{j+1}-A_j-Bc_j.
\]
Étant donnés \(P,E,\Phi\), les blocs \(A\) et les retenues complètes,  
on restaure le registre dodécaphasique. La concaténation restaure en outre
\[
 \iota_N(K)=\iota_N(P)+\iota_N(E)-\iota_N(\Phi)-\iota_N(A)
             -B^Nc_1+c_{N+1}.
\]
Cette réversibilité est une propriété de la fibre avec données et frontières fixes. Elle ne doit pas devenir une nouvelle recette causale qui commence par alpha et fabrique ensuite le registre. La direction de construction reste encore registre signé, registre dodécaphasique, clôture et sortie.

\subsubsection{Prolongation et normalisation canonique de la section complète}
\label{subsec:alpha-completa}

La prolongation conserve le registre dodécaphasique périodique et reçoit les suites
régionales complètes. Soient
\[
 p=\sum_{j\geq1}P_jB^{-j},\qquad
 e_0=\sum_{j\geq1}E_jB^{-j},\qquad
 f=\sum_{j\geq1}\Phi_jB^{-j}.
\]
Chaque série est une réalisation du caractère interne correspondant.  
Les parties entières déjà obtenues sont 3, 2 et 1, de sorte que  
\(p=\pi-3\), \(e_0=e-2\) et  
\(f=\varphi-1\).

\begin{proposition}[Existence et unicité de la normalisation complète]
\label{prop:alpha-normalizacion}
La série
\begin{equation}
 y=p+e_0-f-\kappa_{\mathrm{per}}
   =\sum_{j\geq1}(P_j+E_j-\Phi_j-K_j)B^{-j}
 \label{eq:alpha-serie}
\end{equation}
converge absolument. Dans le domaine des blocs obtenus,
\(0<y<10^{-2}\). Il existe une unique expansion canonique \((A_j)\) de
\(y\), sans queue finalement constante 999, et une unique suite entière
bornée \((c_j)\) qui satisfait \eqref{eq:alpha-identidad-local} avec
\(c_1=0\). Ces retenues appartiennent à l'ensemble invariant
\(\mathcal C=\{-2,-1,0,1,2\}\).
\end{proposition}

\begin{proof}
Les coefficients de la série sont bornés par \(2(B-1)\) en module,  
de sorte qu'elle est dominée par une série géométrique. Si  
\(Z_j=P_j+E_j-\Phi_j-K_j\) et  
\(R_j(Z)=\sum_{r\geq0}Z_{j+r}B^{-r-1}\), chaque queue est la différence  
de deux sommes de deux queues canoniques et appartient à \((-2,2)\).  
Le premier coefficient est  
\(Z_1=141+718-618-234=7\) ; par conséquent  
\(y=(7+R_2(Z))/1000\in(0.005,0.009)\).

Soit \((A_j)\) l'expansion canonique de \(y\). Définissons
\[
 c_j=R_j(Z)-R_j(A).
\]
En multipliant l'égalité des deux séries par \(B^{j-1}\) et en séparant
ses premiers \(j-1\) termes on obtient
\[
 c_j=\sum_{r=1}^{j-1}(A_r-Z_r)B^{j-1-r}\in\mathbb Z,
\]
et \(c_1=0\). De plus, \(R_j(A)\in[0,1)\), donc
\(-3<c_j<2\), ce qui inclut toutes les retenues dans \(\mathcal C\).
Les identités de déplacement des queues donnent
\(Bc_j=Z_j-A_j+c_{j+1}\), qui est la récurrence requise.

Réciproquement, toute solution bornée à sortie canonique et
\(c_1=0\) se télescope à la limite en la même série \(y\). Le développement
canonique est unique et, une fois fixé, la formule des queues détermine
toutes les retenues. Enfin, si l'on applique une exécution finie avec
\(c_{j+1}\in\mathcal C\), on a \(-2000\leq s_j\leq2000\),
de sorte que son quotient euclidien appartient de nouveau à \(\mathcal C\).
\end{proof}

La proposition explicite la normalisation du lecteur ultérieur : elle ne remplace pas la génération des trois suites par une série conventionnelle introduite comme donnée. Une fois ces suites produites par l'état coinductif, leur combinaison et la frontière de la normalisation sont déterminées sans consulter alpha.

La procédure finie conservée dans le corpus propage les cinq
frontières de \(\mathcal C\) et conserve uniquement le préfixe commun à
toutes. Si deux profondeurs ont coalescé avant une frontière,
la récurrence déterministe de droite à gauche donne les mêmes blocs
à sa gauche. C'est la compatibilité avec la troncature. On ne
choisit pas une frontière parce qu'elle donne des chiffres favorables, et on ne déduit pas la
compatibilité à toute profondeur de l'observation d'une longue exécution.

La comparaison initiale montre pourquoi cette précaution importe. Dans
la prolongation stabilisée on obtient \(c_{13}^{\infty}=-1\), tandis
que la fenêtre isolée imposait \(c_{13}=0\). Ainsi,
\[
 884+497-117-601-1=662
\]
dans le dernier bloc du premier tour. Les onze premières triades ne
changent pas, mais la prolongation commence
\[
 (007,297,352,569,283,800,997,285,105,472,380,662,999,\ldots).
\]
Une apparition du bloc 999 ne viole pas la convention canonique : ce
qui est exclu est une queue qui soit finalement constante 999. La fenêtre
à frontière nulle se termine en 663 ; le troncage de la section complète
se termine en 662. Les deux affirmations décrivent des opérations différentes
et exactes. Elles ne doivent pas être fusionnées en une seule notation de projection.

\subsubsection{Définition opérationnelle et équation additive complète}
\label{subsec:alpha-definicion}

\begin{definition}[Coordonnée alpha de la clôture]
\label{def:alpha-operativa}
La voie entière détermine la coordonnée
\begin{equation}
 \alpha_A=\sum_{j\geq1}A_jB^{-j}
 =p+e_0-f-\kappa_{\mathrm{per}},
 \label{eq:alpha-via-a}
\end{equation}
où \((A_j,c_j)\) est la normalisation canonique compatible de
\eqref{eq:alpha-recurrencia}. Dans la réalisation scalaire des
coordonnées HMT,
\begin{equation}
 \boxed{\alpha_A=
 \pi+e-\varphi
 -4-\kappa_{\mathrm{per}}.}
 \label{eq:alpha-identidad-completa}
\end{equation}
\end{definition}

L'entier quatre n'est pas calibré : il est
\(3+2-1\), combinaison des parties entières des trois coordonnées.
La base mille provient de la coordinatisation cubique et la période douze du registre dodécaphasique. Les
signes sont ceux du canal orienté et de la clôture, et la frontière distante est
celle qu'impose la section compatible. Telles sont les origines effectives des
coefficients de l'équation.

À ce niveau, on peut décrire alpha comme la coordonnée scalaire du défaut de fermeture entre les trois réalisations régionales et le retour rationnel dodécaphasique, calculé avec une retenue compatible. Le terme \emph{défaut} désigne une différence algébrique déterminée, et non une perte arbitraire ni une imperfection du système. Cette description ne supprime pas la généalogie : sans les régions, le registre signé, l'orientation et la loi de retenue, l'égalité scalaire ne serait pas une exposition suffisante de la construction.

Il faut également maintenir séparés trois objets que la notation d'alpha peut accompagner dans des lecteurs ultérieurs : \(A\) est un mot de blocs ; le support négatif de pré-retenue s'obtient à partir du signe des totaux avant normalisation ; et un vecteur radial exceptionnel appartient à un codomaine d'incidence. Aucun d'eux n'est interchangeable avec la valeur réelle \(\alpha_A\).

\input{sections/revision_alpha_notacion.tex}
\input{sections/revision_alpha_frontera.tex}

\input{sections/alpha_dependencias.tex}

\subsection{Seconde voie : détermination analytique au moyen des vacances}
\subsubsection{Invariants du transport et résolvante graduée}
\label{subsec:alpha-via-b}

La voie analytique reçoit des invariants de l'état orienté, et non la sortie
\(\alpha_A\). Sa troncature d'ordre six est
\begin{equation}
 P_6(x)=2\pi x-\frac74x^2
 +\frac{x^3}{2\pi}
 +\frac{x^4}{20}-\frac{2x^5}{21}-\frac{x^6}{46},
 \label{eq:alpha-p6}
\end{equation}
et la condition de vacance utilise
\[
 \delta=\log_{10}\frac{10}{9}.
\]
Le rapport \(10/9\) compare la coordinatisation décimale aux neuf positions
APP de la construction nonadique. Son logarithme appartient à cette lecture
ultérieure de compactification ; ce n'est pas une valeur cible d'alpha.

La généalogie des coefficients doit être conservée complète, mais avec sa
force logique exacte. Les développements récents réunissent une signature marquée
de l'état. La matrice de synchronisation et le recensement associé sont
\[
 D_{\mathrm{syn}}=\begin{pmatrix}85&112\\41&52\end{pmatrix},
 \qquad N_9=43.
\]
De son déterminant on extrait
\[
 a=-\frac{\det D_{\mathrm{syn}}}{N_9}
   =\frac{172}{43}=4.
\]
L'horloge induite sur les trous courts a un spectre secondaire \(\{6,7\}\) ; son extrémité supérieure donne \(b=7\). La cardinalité tritique \(m=3\) détermine
\[
 t_*=mb=21,\qquad k_*=mb-1=20.
\]
La lecture de vacance dans la coordinatisation en base mille donne
\[
 v_-=\lfloor1000\delta\rfloor=45,\qquad
 v_+=\lceil1000\delta\rceil=46.
\]
Ces entiers sont des antécédents typés de l'évaluation de la troncature. Ils ne sont obtenus ni à partir de sa racine ni à partir d'une approximation métrologique de celle-ci.

La branche inférieure coïncide en outre avec le déterminant du bloc APP
\[
 B_c=\begin{pmatrix}7&2\\2&7\end{pmatrix},\qquad\det B_c=45.
\]
Le rang transversal du registre du chapitre précédent donne
\(r_{\mathrm{dod}}=11\), et la demi-couronne de phase fournit
\(f_{\mathrm{ph}}=8\binom62=120\). En particulier,
\(45r_{\mathrm{dod}}=495\).

Avec \(\omega=2\pi\), la signature
\[
 \mathcal I_9=(\omega,m,a,b,k_*,t_*,v_-,v_+,f_{\mathrm{ph}},r_{\mathrm{dod}})
\]
alimente le lecteur gradué
\begin{align}
 \mathfrak E_9(\mathcal I_9)(x)
 ={}&\omega x-\frac ba x^2+\omega^{-1}x^3
 +\frac{x^4}{k_*}-\frac{2x^5}{t_*}-\frac{x^6}{v_+}\notag\\
 &-\frac{x^7}{f_{\mathrm{ph}}}+\frac{x^8}{v_-}
 +\frac{2x^9}{v_-r_{\mathrm{dod}}}.
 \label{eq:alpha-firma-coeficientes}
\end{align}
La substitution restitue les coefficients de \(P_6\) et de sa
prolongation d'ordre neuf. Le couple \(45/46\) traverse la séparation
entre les deux : la branche supérieure occupe le degré six et la branche inférieure réapparaît
aux degrés huit et neuf.

La factorisation de \eqref{eq:alpha-firma-coeficientes} se formule  
pour un état qui conserve l'origine, l'ordre des canaux et l'orientation.  
Dans ce domaine marqué, chaque coefficient reste fixé par  
\((\omega,m,a,b,k_*,t_*,v_-,v_+,f_{\mathrm{ph}},r_{\mathrm{dod}})\).  
Deux lecteurs qui préservent la signature et appliquent la même évaluation  
graduée coïncident terme à terme jusqu'au degré neuf. Cette  
unicité ne s'étend pas aux lecteurs qui oublient ces données.

\subsubsection{Résolvante nilpotente aux degrés \(7\), \(8\) et \(9\)}
\label{subsec:alpha-resolvente}

La première prolongation admet une construction opératoire finie
spécialement explicite. Sur l'espace à base \(e_7,e_8,e_9\),
soit
\[
 Ne_7=-\frac83e_8,\qquad Ne_8=\frac2{11}e_9,\qquad Ne_9=0,
 \qquad v_7=-\frac1{120}e_7.
\]
Le symbole \(N\) désigne ici un opérateur nilpotent, et non une profondeur de troncature. Son domaine et sa graduation sont fixés avant d'évaluer le polynôme en un nombre.

\begin{lemma}[Prolongement nilpotent exact]
\label{lem:alpha-resolvente}
On vérifie \(N^3=0\) et
\[
 (I-N)^{-1}v_7
 =-\frac{e_7}{120}+\frac{e_8}{45}+\frac{2e_9}{495}.
\]
L'évaluation \(e_n\mapsto x^n\) produit
\begin{equation}
 T_{7:9}(x)=-\frac{x^7}{120}+\frac{x^8}{45}+\frac{2x^9}{495}.
 \label{eq:alpha-t79}
\end{equation}
\end{lemma}

\begin{proof}
L'opérateur augmente le degré et s'annule sur le dernier vecteur de la base,
de sorte que son cube est nul. L'inverse est la résolvante finie
\(I+N+N^2\). De plus,
\[
 Nv_7=\frac1{45}e_8,\qquad N^2v_7=\frac2{495}e_9.
\]
Additionner les trois termes et évaluer donne la formule.
\end{proof}

Ici, les dénominateurs 120, 45 et 495 ne sont pas des ajustements décimaux : ils proviennent de la demi-couronne, du déterminant APP et de son produit par le rang dodécaphasique. Les transferts \(-8/3\) et \(2/11\), ainsi que le signe et le degré du germe, appartiennent au lecteur typé. Conserver seulement la symétrie du cycle et remplacer cet opérateur par un autre opérateur symétrique ne conserve pas nécessairement les coefficients.

\subsubsection{Racine de multiplicité un de la troncature d'ordre neuf}
\label{subsec:alpha-jet}

Définissons
\[
 P_9=P_6+T_{7:9},\qquad E_{21/22}^{(9)}(x)=P_9(x)-\delta,
 \qquad I_\alpha=(0,10^{-2}).
\]
Cet opérateur fini a un résultat propre d'existence et d'unicité.

Comme vérification ultérieure des caractères déjà construits, l'isolement rationnel peut se reproduire sans une table décimale. Soient
\[
 A_n(t)=\sum_{j=0}^{n}\frac{(-1)^jt^{2j+1}}{2j+1},\qquad
 L_n(t)=2\sum_{j=0}^{n}\frac{t^{2j+1}}{2j+1}\quad(0<t<1).
\]
Le critère alternant et la somme géométrique de la queue donnent
\begin{align}
 A_{2n+1}(t)&<\arctan t<A_{2n}(t),
 \label{eq:alpha-cotas-arctan}\\
 L_n(t)&<\log\frac{1+t}{1-t}
 <L_n(t)+\frac{2t^{2n+3}}{(2n+3)(1-t^2)}.
 \label{eq:alpha-cotas-log}
\end{align}
Nous appliquons \eqref{eq:alpha-cotas-arctan} à l'identité de Machin
\(\pi=16\arctan(1/5)-4\arctan(1/239)\), et
\eqref{eq:alpha-cotas-log} à
\[
 \log(10/9)=\log\frac{1+1/19}{1-1/19},\qquad
 \log 10=\log(10/9)+2\log\frac{1+1/2}{1-1/2}.
\]
Écrivons
\[
 \delta:=\frac{\log(10/9)}{\log 10}.
\]
Avec les ordres de troncature \(80,20,40,120\), appliqués
respectivement aux quatre séries précédentes, on obtient des intervalles
rationnels pour \(\pi\) et \(\delta\), tous deux de largeur inférieure à
\(10^{-70}\). Ces bornes vérifient une réalisation ultérieure ; elles ne
sélectionnent pas l'état APP--TRIT--TPK ni ne reçoivent une valeur d'alpha.

\begin{proposition}[Racine de multiplicité un de la troncature]
\label{prop:alpha-raiz-jet}
L'opérateur \(E_{21/22}^{(9)}\) possède une unique racine
\(\xi_9\) dans \(I_\alpha\), et cette racine est simple. Elle possède en outre l'isolement rationnel
\begin{equation}
 \frac{729735256928380099}{10^{20}}
 <\xi_9<\frac{729735256928380100}{10^{20}}.
 \label{eq:alpha-intervalo-jet}
\end{equation}
\end{proposition}

\begin{proof}
Les bornes des coordonnées internes
\(3.14<\pi<3.15\) et
\(0.045<\delta<0.046\) sont suffisantes pour l'existence. En zéro,
\(E_{21/22}^{(9)}(0)=-\delta<0\). En \(10^{-2}\), la contribution
linéaire moins les termes négatifs est supérieure à \(0.0626\), donc la
valeur de l'opérateur est positive.

Pour \(0\leq x\leq10^{-2}\), en éliminant les termes positifs de
la dérivée on obtient
\[
 (E_{21/22}^{(9)})'(x)
 \geq2\pi-\frac72\cdot10^{-2}
 -\frac{10}{21}\cdot10^{-8}-\frac6{46}\cdot10^{-10}
 -\frac7{120}\cdot10^{-12}>6.24.
\]
La continuité donne l'existence ; la borne sur la dérivée donne l'unicité et la simplicité. Pour l'extrémité gauche \(l\) et l'extrémité droite \(u\) de \eqref{eq:alpha-intervalo-jet}, la substitution rationnelle des bornes \eqref{eq:alpha-cotas-arctan}--\eqref{eq:alpha-cotas-log} dans le polynôme donne, avec arrondi dirigé,
\[
 -4.6\,10^{-20}<E_{21/22}^{(9)}(l)<-4.5\,10^{-20}<0,
 \qquad
 0<1.6\,10^{-20}<E_{21/22}^{(9)}(u)<1.8\,10^{-20}.
\]
La monotonie déjà prouvée place l'unique racine entre les deux extrémités.
L'intervalle se vérifie après avoir construit l'opérateur et n'intervient pas dans la sélection de ses coefficients.
\end{proof}

Le résultat concerne \(\xi_9\). Il n'autorise pas à écrire \(\xi_9=\alpha_A\) comme une égalité exacte de nombres complets. La racine d'une troncature peut partager un préfixe avec la racine de l'opérateur prolongé et différer à une position ultérieure. La simplicité permet de contrôler le déplacement de la racine sous perturbations ; elle n'annule pas ces perturbations.

\subsection{Critère d'identification des lectures compatibles}
\label{subsec:alpha-limite-dos-vias}

L'identification de deux lectures complètes exige de conserver leurs cylindres
et les applications de restriction. Pour une collection graduée compatible,
\[
 E_{21/22}^{(6)}\longleftarrow E_{21/22}^{(9)}
 \longleftarrow E_{21/22}^{(12)}\longleftarrow\cdots,
 \qquad
 E_{21/22}=\varprojlim_m E_{21/22}^{(6+3m)},
\]
la limite désigne le lecteur avec ses transports, et non une queue sélectionnée
librement. Le critère suivant sépare la conséquence de la
compatibilité de la construction effective qui l'établit ensuite.

\begin{proposition}[Identification des limites compatibles]
\label{prop:alpha-dos-limites}
Si les valeurs des voies entière et analytique appartiennent aux mêmes cylindres \(C_N^\alpha\) pour tout \(N\), et
\(\operatorname{diam}C_N^\alpha\to0\), alors
\[
 \alpha_A=\alpha_B=\alpha.
\]
\end{proposition}

\begin{proof}
Pour tout \(N\), la distance entre les deux valeurs est bornée
par \(\operatorname{diam}C_N^\alpha\). Faire tendre \(N\) vers l'infini
donne une distance nulle. L'existence des sections compatibles assure
qu'il ne s'agit pas de l'intersection vide d'un système d'approximations.
\end{proof}

La différence peut s'exprimer algébriquement. Soit \(F=\mathbb Q(\pi,\delta)\) et soit \(j_9:F[[x]]\to F[x]/(x^{10})\) la troncature. Alors
\[
 \ker j_9=x^{10}F[[x]].
\]
Tous les opérateurs
\[
 E_h(x)=E_{21/22}^{(9)}(x)+x^{10}h(x)
\]
ont la même troncature d'ordre neuf. Cependant, en prenant
\(h=0\) et \(h=1/729\),
\[
 E_{1/729}(\xi_9)=\frac{\xi_9^{10}}{729}>0,
\]
de sorte qu'ils ne partagent pas cette racine. L'observation ne permet pas de choisir une queue étrangère au transport HMT ; elle démontre précisément pourquoi la queue générée et sa compatibilité doivent être conservées. Une série formelle arbitraire ne suffit pas non plus pour parler de racines analytiques sans son domaine de convergence ou sans la réalisation du système de cylindres.

\input{sections/alpha_publicacion_completa.tex}

\section{Prolongement exceptionnel : de l'incidence dodécaphasique à Moonshine}
\label{exc:seccion}

L'expression numérique n'épuise pas l'état qui l'a générée. Dans la construction précédente, APP conserve les deux feuillets et la division avec reste et quotient ; le TRIT détermine le régime local et son orientation ; le TPK compose sélection, transport, retenue, mémoire et retour. La lecture archimédienne contracte des cylindres compatibles. La lecture que nous étudions maintenant conserve en revanche l'incidence des positions et leur transport entre repères. Toutes deux partent du même état arithmétique avec mémoire de trajectoire. On ne reconstruit pas une géométrie exceptionnelle à partir des décimales finales de $\pi$, $\varphi$, $e$ ou $\alpha$.

Ce point permet de préciser le sens du prolongement
$\alpha/K\longrightarrow\text{incidence exceptionnelle}$. La flèche
n'affirme pas qu'un nombre réel isolé détermine un code, un réseau ou un
groupe. Elle transporte le registre dodécaphasique et les données signées qui sont
intervenues dans la détermination de $\alpha$, avec les mots et les
cadres issus d'APP--TRIT--TPK. En particulier, le registre ordonné
$K$ peut être retrouvé à partir de sa constante rationnelle $\kappa$ avec les
conventions de base, de longueur et d'origine déjà fixées. La projection sur un
support ne conserve en revanche que les incidences sélectionnées.
Pour sa part, $K_{\rm ph}$
désigne une autre lecture, à mémoire réduite, postérieure à l'holonomie
$\Gamma_9$.

L'exposé sépare trois opérations mathématiques : construire les objets
finis et leurs applications ; les prolonger en un réseau et en une algèbre d'opérateurs
de vertex ; reconnaître ensuite les objets obtenus au moyen des théorèmes
classiques pertinents. Les noms Paley, Witt, Mathieu, Golay, Leech et
Monstre désignent ces reconnaissances et les réalisations explicites
qui les rendent vérifiables. Ce ne sont pas des sélecteurs rétrospectifs des
germes ni des coefficients de la génération.

\input{sections/revision_incidencia_argumento.tex}



\subsection{Compatibilité des réalisations archimédienne et d'incidence}
\label{exc:doble-lectura}

Soit $X_n^{\rm cal}$ l'espace des préfixes admissibles de longueur $n$ de l'état arithmétique avec mémoire de trajectoire, avec les troncatures $p_{n,m}:X_n^{\rm cal}\to X_m^{\rm cal}$ pour $m\leq n$. L'exposant rappelle que le repère initial déclaré et son transport sont conservés, et non qu'un paramètre est ajusté à une valeur expérimentale. Nous écrivons $(w_j,f_j)$ pour le mot visible de six trits et le repère au pas $j$. Les repères satisfont une loi causale
\[
 f_{j+1}=g_j(x)f_j,
\]
où $g_j(x)$ dépend du préfixe déjà construit. En conséquence, une
prolongation n'altère pas les cadres antérieurs.

La chaîne
\[
 w_6\longrightarrow w_{12}\longrightarrow w_{18}
 \longrightarrow w_{24}\longrightarrow w_{30}
 \longrightarrow R_{36}\longrightarrow G_9
\]
est conservé dans cette lecture. Le retour de la phase ne réinitialise ni le
cadre ni la mémoire de l'état. De même, on maintient distincts le
recensement de $468$ émissions décimales, les $243$ mots visibles de six
trits et l'espace ambiant $\mathbb F_3^6$, de $729$ mots. L'espace ambiant sera
utilisé pour construire un code linéaire ; il ne se substitue pas au domaine
admissible des émissions.

La lecture archimédienne prend, pour un bloc $b$ de six trits, son
adresse entière $\nu(b)\in\{0,\ldots,728\}$. Un préfixe détermine
\[
 q_n=m+\sum_{j=0}^{n-1}\frac{\nu(b_j)}{729^{j+1}},
 \qquad I_n=[q_n,q_n+729^{-n}].
\]
La compatibilité des préfixes donne des cylindres emboîtés dont le diamètre tend
vers zéro. Leur classe de Cauchy appartient d'abord à la complétion interne
$\mathbb R_{\rm HMT}$ ; son identification ultérieure à la droite réelle
usuelle n'intervient pas dans le choix des blocs.

Pour la lecture d'incidence, nous introduisons, sur le même préfixe, l'application de graphe
\begin{equation}
 \gamma_A:\mathbb F_3^6\longrightarrow\mathbb F_3^{12},
 \qquad w\longmapsto(w,wA).
 \label{exc:gamma}
\end{equation}
Les mots s'écrivent comme des vecteurs lignes. Si $A^f=fAf^{-1}$, le changement de repère satisfait l'identité concrète
\begin{equation}
 \gamma_{A^f}(wf^{-1})
   =\gamma_A(w)(f^{-1}\oplus f^{-1}).
 \label{exc:covariancia}
\end{equation}
Le codomaine calibré conserve $f$ avec le mot codé.
Nous définissons alors
\begin{equation}
 Q_n(x)=\bigl((f_j,\gamma_{f_jA_Wf_j^{-1}}(w_j))\bigr)_{0\leq j<n}.
 \label{exc:qn}
\end{equation}

\begin{proposition}[Compatibilité et prolongement de la lecture d'incidence]
\label{exc:limite-inverso}
Si $r_{n,m}$ tronque les suites codées, on a
\[
 r_{n,m}Q_n=Q_m p_{n,m}.
\]
Les $Q_n$ déterminent donc une unique application $Q_\infty$ entre les limites inverses respectives. Dans chaque repère, l'application conserve exactement le mot visible, puisque $\operatorname{pr}_1\gamma_A=\operatorname{id}$.
\end{proposition}

\begin{proof}
Deux préfixes compatibles possèdent les mêmes mots et les mêmes
opérateurs de transport jusqu'à l'indice $m-1$. Ils possèdent, par récurrence,
les mêmes cadres à ces indices. Les $m$ premières composantes de
\eqref{exc:qn} coïncident, ce qui prouve le carré. La propriété
universelle de la limite inverse donne $Q_\infty$. L'affirmation de récupération
est la projection sur les six premières coordonnées de chaque graphe.
\end{proof}

Cette injectivité locale a un domaine précis. Elle ne démontre pas que
$Q_\infty$ récupère toute la mémoire profonde de l'état TPK lorsque seule
la suite des mots visibles et des repères a été retenue. Elle démontre encore moins
l'injectivité après le passage des mots aux supports.
En outre, un changement général de base dans $\operatorname{GL}_6(\mathbb F_3)$
ne conserve pas à lui seul le poids de Hamming ordinaire. Dans une coordinatisation
générale, on transporte aussi l'incidence et la métrique de la coordinatisation
initiale. Les calculs de supports qui suivent sont effectués dans la coordinatisation
de Paley déclarée ; les réétiquetages monomiaux conservent directement leur
interprétation en coordonnées.

\subsection{La réalisation finie et le code ternaire}
\label{exc:paley-codigo}

L'opérateur elliptique du TRIT admet la réalisation finie suivante dans
la coordinatisation à six positions. Celles-ci coordinatisent les six unités
$\{1,2,4,5,7,8\}$ modulo neuf, conservées par le transport modulaire.
Sur $\mathbb P^1(\mathbb F_5)$, une fois
l'ordre $\infty,0,1,2,3,4$ déclaré, soit $C$ la matrice de diagonale
nulle, avec $C_{\infty,a}=C_{a,\infty}=1$ et
$C_{a,b}=\chi(b-a)$ aux cinq positions finies. Ici,
$\chi(1)=\chi(4)=1$, $\chi(2)=\chi(3)=-1$ et $\chi(0)=0$.
Il s'agit d'une réalisation en coordonnées de la structure quadratique déjà
transportée ; l'identification des six positions à la droite
projective est une donnée de coordinatisation, et non une nouvelle entrée numérique dans le
générateur.

Les identités de la somme quadratique donnent $C^2=5I_6$. En réduisant modulo
trois, nous obtenons
\begin{equation}
 A_W=
 \begin{pmatrix}
 0&1&1&1&1&1\\
 1&0&1&2&2&1\\
 1&1&0&1&2&2\\
 1&2&1&0&1&2\\
 1&2&2&1&0&1\\
 1&1&2&2&1&0
 \end{pmatrix},
 \qquad A_W^2=2I_6=-I_6.
 \label{exc:aw}
\end{equation}
L'égalité quadratique réalise le régime elliptique ; elle n'identifie pas le
corps $\mathbb F_3$ à une réalisation archimédienne de l'unité
imaginaire. Ce sont des réalisations distinctes d'une relation opératoire
commune, avec leurs domaines respectifs.

\Needspace{8\baselineskip}
\begin{definition}[Code de la coordinatisation dodécaphasique]
\label{exc:codigo}
Le code ternaire associé à \eqref{exc:aw} est
\[
 C_W=\{(w,wA_W):w\in\mathbb F_3^6\}
       \subset\mathbb F_3^{12}.
\]
Son ensemble d'hexades est
\[
 \mathcal H=
 \{\operatorname{supp}(c):c\in C_W,
                           \operatorname{wt}(c)=6\}.
\]
\end{definition}

\begin{proposition}[Autodualité, distance et système de Witt]
\label{exc:codigo-witt}
Le code $C_W$ est autodual de paramètres $[12,6,6]_3$. Son énumérateur
des poids est
\begin{equation}
 \sum_{c\in C_W}z^{\operatorname{wt}(c)}
   =1+264z^6+440z^9+24z^{12}.
 \label{exc:enumerador}
\end{equation}
L'ensemble $\mathcal H$ possède $132$ éléments et constitue un système
$S(5,6,12)$ : chaque ensemble de cinq positions est contenu dans une
unique hexade.
\end{proposition}

\begin{proof}
La matrice génératrice $[I_6\mid A_W]$ est de rang six et, par symétrie de $A_W$, satisfait
\[
 [I_6\mid A_W][I_6\mid A_W]^{\mathsf T}
   =I_6+A_W^2=0.
\]
Le code est auto-orthogonal, de dimension égale à la moitié de sa longueur, et donc
autodual. L'évaluation finie des $3^6$ mots donne
\eqref{exc:enumerador} ; il s'agit d'une énumération sur la matrice explicite,
sans valeurs réelles ni cibles externes. En particulier, la distance
minimale est six.

Si deux mots de poids six ont le même support, parmi leurs six quotients de coordonnées non nulles, l'un des signes $1,-1$ apparaît au moins trois fois. En soustrayant le multiple correspondant, on obtiendrait un mot non nul de poids au plus trois, ce qui est impossible. Un support correspond donc exactement à la paire $\{c,-c\}$ et il y a $264/2=132$ supports. Si deux supports distincts partageaient cinq positions, le même argument annulerait au moins trois coordonnées d'une union de taille sept, produisant un poids au plus quatre. Un quintuplet appartient donc à au plus une hexade. Enfin,
\[
 132\binom65=792=\binom{12}5,
\]
de sorte qu'il appartient à une et une seule.
\end{proof}

La reconnaissance ultérieure est le code de Golay ternaire étendu et
le petit système de Witt \cite{conwaysloane}. La preuve précédente explique quel objet est
reconnu. Elle ne remplace pas la généalogie APP--TRIT--TPK par une recherche
parmi les codes connus.

\subsection{Drapeau dodécaphasique associé à \texorpdfstring{$K$ et $\alpha$}{K et alpha}}
\label{exc:bandera}

\input{sections/incidencia_registro.tex}

Dans le cadre hérité du registre dodécaphasique, les données sont
\begin{align}
 K={}&(234,543,140,729,659,824,621,\mathtt{058},914,794,146,601),
 \label{exc:k}\\
 w_\pi={}&010211, &w_e={}&201101,\nonumber\\
 w_\varphi={}&121200, &w_{\rm APP}={}&022020.\nonumber
\end{align}
En appliquant $\gamma_{A_W}$, on obtient les incidences suivantes :
\begin{align*}
 P_\pi&=\{2,4,5,6,7,8,9,10,11\},\\
 H_e&=\{1,3,4,6,9,10\},\\
 H_\varphi&=\{1,2,3,4,7,9\},\\
 H_{\rm APP}&=\{2,3,5,10,11,12\}.
\end{align*}
Le premier est le support d'un mot de poids neuf ; les trois autres
sont des hexades. Le registre fournit, par sa lecture de bloc haut,
\begin{equation}
 B_K=\{i:K_i\geq729\}=\{4,6,9,10\}=H_e\cap P_\pi.
 \label{exc:bk}
\end{equation}
L'entier $729$ indique ici une frontière du format de blocs
déclaré. Il n'est pas identifié par son seul cardinal à l'espace ambiant
ternaire utilisé dans \eqref{exc:codigo}. Dans ce registre précis, tous
les seuils entiers entre $660$ et $729$ produisent le même support :
le support ne restitue ni le seuil ni les douze blocs.

La voie arithmétique de $\alpha$ conserve, avant la retenue, le vecteur
signé
\begin{equation}
 Z_0=(7,297,353,-430,-715,-1199,-3,286,-894,-528,380,663).
 \label{exc:precarry}
\end{equation}
Son support négatif est
\begin{equation}
 H^-_\alpha=\{4,5,6,7,9,10\}.
 \label{exc:halpha}
\end{equation}
Ces signes ne sont pas extraits du développement décimal du scalaire
$\alpha$ ; ils sont conservés dans l'état qui précède son expression numérique.

\input{sections/incidencia_normalizacion.tex}

\begin{proposition}[Coïncidence des deux sélecteurs du drapeau]
\label{exc:selector-bandera}
Dans $\mathcal H$, le support \eqref{exc:halpha} est l'unique hexade $H$ qui satisfait
\[
 B_K\subset H\subset P_\pi,
 \qquad
 (|H\cap H_e|,|H\cap H_\varphi|,|H\cap H_{\rm APP}|)=(4,3,2).
\]
Ainsi, la lecture signée et la lecture d'incidence sélectionnent le même
drapeau $B_K\subset H^-_\alpha$.
\end{proposition}

\begin{proof}
Les quatre hexades qui contiennent $B_K$ s'obtiennent en ajoutant respectivement les paires
\[
 \{1,3\},\qquad\{2,11\},\qquad\{5,7\},\qquad\{8,12\}.
\]
Seules les deuxième et troisième paires sont contenues dans $P_\pi$. Pour la
deuxième, l'intersection avec $H_\varphi$ a une taille de trois et avec
$H_{\rm APP}$ une taille de trois ; elle ne satisfait pas la dernière condition. Pour la
troisième, les trois intersections ont les tailles $4,3,2$. Cette dernière
produit précisément le support négatif de \eqref{exc:precarry}.
\end{proof}

Il convient de conserver tant le drapeau que son cadre marqué :
\[
 \mathcal I_*=(B_K\subset H^-_\alpha),\qquad
 \widetilde{\mathcal I}_*
   =(\mathcal I_*;P_\pi,H_e,H_\varphi,H_{\rm APP}).
\]
Les opérations d'union, d'intersection et de différence sont équivariantes sous un réétiquetage simultané du plan combinatoire. C'est la classe marquée, et non les numéros des positions pris isolément, qui est l'objet transporté.

\begin{lemma}[Origine sélectionnée par l'incidence marquée]
\label{exc:origen}
Pour le cadre précédent,
\[
 o_*=(H_e\cap H_\varphi)
           \setminus(H^-_\alpha\cup H_{\rm APP})=\{1\}.
\]
La règle est équivariante sous tout automorphisme appliqué simultanément
aux données de $\widetilde{\mathcal I}_*$.
\end{lemma}

\begin{proof}
La première intersection est $\{1,3,4,9\}$. La réunion que l'on soustrait
contient $3,4,9$ et ne contient pas $1$. L'équivariance découle de ce qu'une
bijection commute avec les opérations ensemblistes.
\end{proof}

Cette origine sélectionnera une direction du réseau voisin lorsque la métrique $A_2$ et la chambre positive seront déclarées. L'incidence seule ne contient pas cette métrique. Le passage au réseau conservera explicitement les deux données.

\input{sections/incidencia_bandera.tex}

\subsection{Entrelacement isométrique des représentations d'incidence}
\label{exc:entrelazador}

L'incidence définit une application linéaire dont l'action peut être vérifiée
position par position :
\begin{equation}
 I:\mathbb R^{12}\longrightarrow\mathbb R^{\mathcal H},
 \qquad (Iv)(H)=\sum_{p\in H}v_p.
 \label{exc:incidencia}
\end{equation}
Soit $V_{11}=\mathbf1^\perp\subset\mathbb R^{12}$. Les deux réalisations
euclidiennes sont ici des langages ultérieurs pour le même plan fini.

\begin{proposition}[Isométrie de l'écran dodécaphasique]
\label{exc:isometria}
On a
\[
 I^{\mathsf T}I=36I_{12}+30J_{12},
\]
où $J_{12}$ est la matrice dont les entrées sont un. En conséquence,
\[
 U_W=\frac16 I\big|_{V_{11}}:
 V_{11}\xrightarrow{\ \sim\ }E_{\rm inc}:=I(V_{11})
\]
est une isométrie équivariante pour les actions par permutation des
automorphismes du système.
\end{proposition}

\begin{proof}
Chaque point appartient à
$r=\binom{11}{4}/\binom54=66$ hexades, et chaque paire à
$\lambda_2=\binom{10}{3}/\binom43=30$. Ce sont, respectivement, les
entrées diagonales et non diagonales de $I^{\mathsf T}I$. Sur
$\mathbf1^\perp$, le terme $J_{12}$ s'annule et il reste $36I$.
Enfin, si $g$ est un automorphisme, $p\in H$ équivaut à
$g(p)\in g(H)$ ; d'où $I\rho_{12}(g)=\rho_{\mathcal H}(g)I$.
\end{proof}

L'information d'incidence possède également une lecture spectrale concrète. Définissons
\[
 G_{H,H'}=\binom{|H\cap H'|}{4}.
\]

\Needspace{8\baselineskip}
\begin{lemma}[Action du noyau d'intersection]
\label{exc:espectro}
Si $J_{132,12}$ est la matrice rectangulaire dont tous les coefficients valent un,
\[
 GI=30I+15J_{132,12}.
\]
Ainsi, $G$ agit par le scalaire $30$ sur $E_{\rm inc}$.
\end{lemma}

\begin{proof}
L'entrée $(H,p)$ de $GI$ compte les couples $(T,H')$ tels que
$T\subset H\cap H'$, $|T|=4$ et $p\in H'$. Si $p\in H$, il y a dix
quadruplets $T$ contenant $p$, chacun contenu dans quatre hexades,
et cinq qui ne le contiennent pas, chacun ayant une unique extension du
quintuplet $T\cup\{p\}$. Le total est $10\cdot4+5=45$.
Si $p\notin H$, les quinze quadruplets de $H$ donnent chacun une
extension : le total est $15$. D'où l'identité. Le terme rectangulaire
s'annule sur $V_{11}$.
\end{proof}

Il existe bien ici un opérateur d'entrelacement matériel : son domaine, son codomaine, ses entrées et sa normalisation sont déterminés. Par exemple, tout projecteur orthogonal $P$ déjà construit sur $V_{11}$ se transporte en $Q=U_WPU_W^*$ sur $E_{\rm inc}$, avec $U_WP=QU_W$. Cette affirmation n'identifie pas par leur seule trace tous les opérateurs qui agissent dans ces espaces.

\subsection{L'étoile de Witt et la génération de \texorpdfstring{$M_{12}$}{M12}}
\label{exc:estrella}

\begin{definition}[Étoile sur un quadruplet]
\label{exc:def-estrella}
Pour $B\in\binom{[12]}4$, l'étoile de Witt est la collection des
quatre hexades qui contiennent $B$. Ses pétales sont les paires $H\setminus B$.
La permutation $\sigma_B$ fixe $B$ et échange les deux points de chaque
pétale.
\end{definition}

L'existence de quatre hexades se déduit de
\[
 \lambda_4=\frac{\binom81}{\binom21}=4.
\]
Deux pétales ne peuvent pas partager un point, car deux hexades contiendraient alors le même quintuplet. Ce sont donc quatre paires disjointes qui partitionnent $[12]\setminus B$. Cette description n'est ni une étoile à cinq branches ni une métaphore géométrique : c'est une incidence quatre-à-un du plan combinatoire.

\begin{figure}[htbp]
\centering
\begin{tikzpicture}[x=1cm,y=1cm, every node/.style={font=\small}]
 \node[draw,rounded corners,minimum width=3.8cm,minimum height=.7cm]
  (b) at (0,0) {$B_K=\{4,6,9,10\}$};
 \node[draw,minimum width=2.7cm,minimum height=.65cm]
  (p1) at (-4.6,-1.7) {$\{1,3\}$};
 \node[draw,minimum width=2.7cm,minimum height=.65cm]
  (p2) at (-1.55,-1.7) {$\{2,11\}$};
 \node[draw,very thick,minimum width=2.7cm,minimum height=.65cm]
  (p3) at (1.55,-1.7) {$\{5,7\}$};
 \node[draw,minimum width=2.7cm,minimum height=.65cm]
  (p4) at (4.6,-1.7) {$\{8,12\}$};
 \draw (b) -- (p1); \draw (b) -- (p2);
 \draw[very thick] (b) -- (p3); \draw (b) -- (p4);
 \node at (-4.6,-2.35) {$H_e$};
 \node at (-1.55,-2.35) {$B_K\cup\{2,11\}$};
 \node at (1.55,-2.35) {$H^-_\alpha$};
 \node at (4.6,-2.35) {$B_K\cup\{8,12\}$};
 \node at (0,-3.05)
  {$\sigma_{B_K}=(1\ 3)(2\ 11)(5\ 7)(8\ 12)$};
\end{tikzpicture}
\caption{L'étoile concrète du drapeau dodécaphasique. Chaque ligne
indique l'union du quadruplet central avec la paire à son extrémité ; la
ligne épaisse signale l'hexade également sélectionnée par les signes
antérieurs à la retenue de $\alpha$. Les quatre paires épuisent les huit
positions extérieures.}
\label{exc:fig-estrella}
\end{figure}

\begin{proposition}[La collection d'étoiles]
\label{exc:m12}
Les $\binom{12}4=495$ permutations $\sigma_B$ préservent $\mathcal H$.
Le groupe qu'elles engendrent est d'ordre $95040$ et est le groupe de Mathieu
$M_{12}$ dans son action sur le petit système de Witt.
\end{proposition}

\begin{proof}
L’affirmation de préservation est une vérification finie sur les
$132$ hexades de \eqref{exc:codigo} : pour chaque quadruplet, on construit
les quatre pétales et l’on vérifie
$\{\sigma_B(H):H\in\mathcal H\}=\mathcal H$.
La clôture par composition de ces permutations produit $95040$
éléments. Dans l’énumération lexicographique, il suffit d’ajouter successivement
six étoiles qui ne sont pas encore contenues dans le groupe ; les ordres intermédiaires
sont $2,6,18,360,7920,95040$, et les $495$ étoiles appartiennent au groupe
final. Le procédé et les six générateurs sont explicités dans
la note technique de provenance. Enfin, on applique la reconnaissance
du groupe d’automorphismes de $S(5,6,12)$ comme $M_{12}$, dont l’ordre est
$12\cdot11\cdot10\cdot9\cdot8=95040$.
\end{proof}

Une étoile individuelle engendre un groupe cyclique d'ordre deux. Elle
n'engendre pas à elle seule $M_{12}$, et encore moins le Monstre. L'
existence de $495$ générateurs d'un groupe de permutations ne prouve pas non plus que
chacun soit le transport d'un lacet spécifique du TPK. Cette affirmation
requiert l'application de lacets correspondante ; le résultat établi ici
est l'action d'incidence finie.

\begin{remark}[L'égalité des traces n'identifie pas les opérateurs]
Sur $V_{11}$, une étoile a pour spectre $+1$ avec multiplicité sept
et $-1$ avec multiplicité quatre. Sa trace normalisée est $3/11$.
Un projecteur de rang trois sur le même espace possède également une trace
normalisée $3/11$, mais un spectre $1^3,0^8$. Ils ne sont pas conjugués : leurs
polynômes minimaux et leurs spectres sont différents. L'opérateur d'entrelacement
\eqref{exc:incidencia} transporte des opérateurs définis ; il ne transforme pas
cette coïncidence scalaire en une équivalence d'opérateurs.
\end{remark}

\subsection{Le résidu binaire et les cinq régions de \texorpdfstring{$\pi$}{pi}}
\label{exc:cinco-regiones}

La reconnaissance suivante utilise le code de Golay binaire
étendu $G_{24}$, de paramètres $[24,12,8]_2$. Ses poids sont
$0,8,12,16,24$, et les supports de poids huit forment le plan
$S(5,8,24)$. Nous fixons une dodécade $D$, support d'un mot de poids
douze. Cette donnée et l'isomorphisme d'incidence qui sera déclaré
ci-après appartiennent à la coordinatisation de réalisation ; ils ne sont pas utilisés pour
engendrer ni sélectionner les décimales de $\pi$.

\begin{proposition}[Résidu d'une dodécade et élévation d'hexades]
\label{exc:residuo-binario}
L'ensemble
\[
 \mathcal H(D)=
 \{O\cap D: O\text{ est une octade},\ |O\cap D|=6\}
\]
est un système $S(5,6,12)$ sur $D$. Chacune de ses hexades a une unique élévation en une octade. Pour tout quadruplet $B\subset D$, les cinq octades qui le contiennent se répartissent en quatre élévations résiduelles et une unique octade $O_B^\perp$ avec $O_B^\perp\cap D=B$.
\end{proposition}

\begin{proof}
Un quintuplet $A\subset D$ appartient à une unique octade $O$. Si
$j=|O\cap D|$, le mot binaire somme de $O$ et $D$ a pour poids
$20-2j$. Parmi les poids permis, et avec $0\leq j\leq8$, seuls
$j=2,4,6$ sont possibles. Comme $A\subset O\cap D$, on a $j\geq5$
et donc $j=6$. Cela donne l'existence et l'unicité de l'hexade
résiduelle, ainsi que l'unicité de son élévation en choisissant cinq de ses
points. D'autre part,
\[
 \lambda_4(S(5,8,24))=\frac{\binom{20}1}{\binom41}=5,
 \qquad
 \lambda_4(S(5,6,12))=4.
\]
Les quatre hexades résiduelles sur $B$ donnent quatre octades distinctes. La cinquième ne peut pas couper $D$ en six points, car elle serait une autre hexade résiduelle ; puisqu'elle contient $B$, son intersection est de taille quatre et coïncide avec $B$.
\end{proof}

Une coordinatisation $\ell:[12]\to D$ qui préserve les hexades identifie
$\mathcal H$ à $\mathcal H(D)$. Son existence est la reconnaissance
du plan déjà construit comme le petit plan de Witt ; un choix
précis de $\ell$ fixe le réétiquetage. Le relèvement
\[
 \iota_D:\mathcal H\longrightarrow
       \{\text{octades de }G_{24}\},
 \qquad H\longmapsto O\text{ tel que }O\cap D=\ell(H),
\]
est injective parce que l'intersection avec $D$ redonne $\ell(H)$. Cette injectivité a pour domaine les hexades, et non le registre $K$.

La relation avec les cinq régions de $\pi$ conserve des données plus riches
que le nombre cinq. Les émissions et leurs signatures sont :
\begin{center}
\small
\begin{tabular}{@{}crll@{}}
\toprule
Région & Multiplicité & Émission $U_6$ & Signature \\
\midrule
$\perp$ & $432$ & $(501,614,498,169,272,272)$ & $555555$ \\
$++_1$ & $144$ & $(810,923,870,169,272,272)$ & $339555$ \\
$++_2$ & $144$ & $(810,983,810,169,272,272)$ & $393555$ \\
$++_3$ & $144$ & $(870,923,810,169,272,272)$ & $933555$ \\
$--$ & $144$ & $(870,923,810,418,674,521)$ & $933717$ \\
\bottomrule
\end{tabular}
\end{center}
Dans chaque bloc décimal de trois chiffres, la signature utilise la différence
modulaire des deux premiers chiffres. Les cinq émissions ont le
même mot tritique visible $w_\pi=010211$, mais pas la même
généalogie régionale. Les trois régions positives conservent la queue
$169|272|272$ ; la région négative la remplace par $418|674|521$ ;
la région transversale a une signature constante et une multiplicité supérieure.

L'extension marquée de la coordinatisation régionale envoie la région transversale
sur $O_{\ell(B_K)}^\perp$, la région négative sur $\iota_D(H^-_\alpha)$ et les
trois régions positives sur les trois autres relèvements de l'étoile. Ainsi se
réalise la partition $4+1$ des cinq octades sur $\ell(B_K)$.
La correspondance individuelle entre les trois noms positifs et
les trois pétales restants exige un ordre de coordinatisation ; l'incidence ne la détermine
pas sans cette donnée. L'énoncé précis est donc une
correspondance marquée des régions, compatible avec leurs rôles et avec
le drapeau, et non une bijection canonique déduite du cardinal cinq.

L'identité $432+4\cdot144=1008$ conserve les multiplicités
régionales. On a aussi $1008=36\cdot28$, et une octade possède
$\binom82=28$ paires. Ces égalités sont des contrôles arithmétiques du
catalogue, et non des substituts à l'application $\iota_D$ ni à la coordinatisation régionale.

\input{sections/estrella_pentadica_recuperada.tex}

\subsection{Du code ternaire au réseau de Niemeier}
\label{exc:niemeier}

Le passage à la dimension vingt-quatre ne nécessite pas de parcourir d'abord le
code binaire. Il s'effectue directement à partir de $C_W$ par recollement des
discriminants. On distingue ainsi deux prolongations compatibles :
l'une conserve le résidu $W_{12}\subset W_{24}$ ; l'autre construit le
réseau sur lequel sera réalisée l'algèbre d'opérateurs de vertex.

Soit $A_2=\mathbb Z a_1\oplus\mathbb Z a_2$ de matrice de Gram
\[
 \begin{pmatrix}2&-1\\-1&2\end{pmatrix},
 \qquad \rho=a_1+a_2,\qquad \rho^2=2.
\]
Nous représentons les trois classes de $A_2^*/A_2$ par
\[
 \varpi_0=0,\qquad
 \varpi_1=\frac{2a_1+a_2}{3},\qquad
 \varpi_2=\frac{a_1+2a_2}{3}.
\]
L'application $j\mapsto\varpi_j+A_2$ identifie additivement
$\mathbb F_3$ à ce discriminant. Pour $c\in C_W$, nous écrivons
$\phi(c)=(\varpi_{c_1},\ldots,\varpi_{c_{12}})$ et définissons
\begin{equation}
 N(C_W)=\bigcup_{c\in C_W}\bigl(A_2^{12}+\phi(c)\bigr).
 \label{exc:niemeier-def}
\end{equation}

\begin{theorem}[Pegado ternario]
\label{exc:niemeier-teorema}
L'objet \eqref{exc:niemeier-def} est un réseau positif, pair et
unimodulaire de rang $24$. Son système de racines est exactement
$A_2^{12}$ ; il est donc reconnu comme le réseau de Niemeier ayant
ce système de racines.
\end{theorem}

\begin{proof}
La linéarité du code rend l'union de classes stable par addition. L'auto-orthogonalité annule l'accouplement discriminant : celui-ci est un multiple de $c\cdot d/3$ modulo $\mathbb Z$. Le réseau est donc entier. La norme d'une classe de recollement est $2\operatorname{wt}(c)/3$ modulo $2\mathbb Z$. Les poids de \eqref{exc:enumerador} sont des multiples de trois, de sorte qu'il est pair.

L'indice de $A_2^{12}$ dans $N(C_W)$ est $|C_W|=3^6$. En conséquence,
\[
 \det N(C_W)=\frac{3^{12}}{(3^6)^2}=1.
\]
La positivité procède de la somme orthogonale des douze métriques $A_2$.
Dans une classe non nulle de $A_2^*/A_2$, la norme minimale est $2/3$ ;
tout mot non nul possède au moins six positions non nulles.
Par conséquent, un vecteur d'une classe de recollement non triviale a une norme
au moins égale à $6\cdot2/3=4$. Aucune racine de norme deux n'apparaît hors de
$A_2^{12}$. À l'intérieur de celui-ci, les racines sont précisément celles de ses
douze composantes.
\end{proof}

\subsection{Sélection par incidence du voisin unimodulaire sans racines}
\label{exc:leech}

\input{sections/incidencia_reticulo.tex}

L'étiquette $o_*=1$ du drapeau se réalise maintenant comme l'une des
douze composantes $A_2$. Cette opération déclare la métrique précédente et
la chambre radiale positive
\[
 v(a)=\sum_{i=1}^{12}a_i\rho_i,
 \qquad a_i=1+3b_i,\quad b_i\in\mathbb Z_{\geq0}.
\]
L'exigence de parité du voisin d'indice trois est $v(a)^2\equiv0\pmod{18}$. Puisque
\[
 v(a)^2=2\sum_i(1+3b_i)^2,
\]
cette congruence équivaut à $\sum_i b_i\equiv1\pmod3$. La plus petite
valeur de la norme dans la chambre s'obtient avec un unique $b_i=1$ et
tous les autres nuls ; elle vaut $54$. L'origine marquée sélectionne
\begin{equation}
 v_\alpha=4\rho_1+\rho_2+\cdots+\rho_{12},
 \qquad v_\alpha^2=54.
 \label{exc:vector}
\end{equation}
Le mot ``plus petit'' se rapporte à cette chambre positive. Il n'affirme pas la minimalité dans toute la classe résiduelle : le vecteur $(a_1-2a_2,\rho,\ldots,\rho)$ représente la même classe modulo trois et a pour norme $14+11\cdot2=36$. La restriction de chambre fait partie des données de la construction, et n'est pas une conclusion que l'on puisse omettre.

Soit $N=N(C_W)$. Nous définissons
\begin{equation}
 N_{\alpha,0}=\{x\in N:\langle x,v_\alpha\rangle\equiv0\pmod3\},
 \qquad
 N_\alpha=N_{\alpha,0}+\mathbb Z\frac{v_\alpha}{3}.
 \label{exc:vecino}
\end{equation}

\begin{theorem}[Voisin de Leech déterminé par le drapeau marqué]
\label{exc:teorema-leech}
Avec la métrique et la chambre déclarées, $N_\alpha$ est un réseau
positif, pair et unimodulaire de rang $24$, sans vecteurs de norme deux.
D'après le théorème de classification correspondant, il est isométrique au
réseau de Leech $\Lambda_{24}$.
\end{theorem}

\begin{proof}
Le caractère $x\mapsto\langle x,v_\alpha\rangle\bmod3$ est non nul :
une racine simple d'une composante s'apparie avec $a_i\rho_i$ en
$a_i\equiv1\pmod3$. Par conséquent, $[N:N_{\alpha,0}]=3$ et
$\det N_{\alpha,0}=9$. De plus,
$v_\alpha/3\in N_{\alpha,0}^*$ et $v_\alpha\in N_{\alpha,0}$.
Si $v_\alpha/3$ appartenait à $N$, l'intégralité de $N$ rendrait
tous les appariements avec $v_\alpha$ divisibles par trois,
ce qui contredirait le caractère non nul. La classe de $v_\alpha/3$ est
donc d'ordre trois. On en déduit
\[
 [N_\alpha:N_{\alpha,0}]=3,
 \qquad \det N_\alpha=9/3^2=1.
\]
L'extension est entière, et pour $x\in N_{\alpha,0}$,
\[
 \left(x+m\frac{v_\alpha}{3}\right)^2
   =x^2+2m\frac{\langle x,v_\alpha\rangle}{3}+6m^2
   \in2\mathbb Z.
\]
Elle est donc paire.

Il reste à exclure les racines, ce qui est l'étape décisive. Les racines de $N$ appartiennent à $A_2^{12}$. Leurs accouplements avec $v_\alpha$ sont congrus à $\pm1$ ou à $\pm2$ modulo trois ; aucune n'appartient à $N_{\alpha,0}$. Pour les deux autres classes du voisin, chaque composante appartient à l'une des classes locales
\[
 A_2+\varpi_j\pm\rho/3,\qquad j=0,1,2.
\]
Comme $4\rho/3\equiv\rho/3\pmod{A_2}$, cette description vaut
également dans la composante distinguée. En écrivant un vecteur local
sous la forme $(u a_1+v a_2)/3$, sa norme est
$2(u^2-uv+v^2)/9$. Les six classes précédentes ont des résidus non
nuls et un minimum $2/9$ : les représentants de forme quadratique minimale
sont $(\pm1,0),(0,\pm1),(1,1),(-1,-1)$ dans les résidus appropriés.
Un vecteur d'une classe nouvelle a donc une norme au moins égale à
$12\cdot2/9=8/3>2$. Aucune des trois classes ne contient de racines.
La classification des réseaux pairs unimodulaires positifs de
rang vingt-quatre identifie l'unique cas sans racines à
$\Lambda_{24}$.
\end{proof}

L'argument ne reconnaît pas Leech par une coïncidence de dimension.
Il construit le recollement, prouve l'intégralité, la parité et le déterminant, et
élimine les racines au moyen d'une borne locale. Ce n'est qu'ensuite qu'il applique le
théorème de classification. La dépendance à l'égard de $\alpha$ est conservée dans
le drapeau qui sélectionne $o_*$ et dans le vecteur \eqref{exc:vector} ;
on ne transforme pas la valeur scalaire de la constante en une longueur du
réseau.

\subsection{L'algèbre d'opérateurs de vertex et le Monstre}
\label{exc:voa}

Soit $\Lambda=N_\alpha$, déjà reconnu comme le réseau de Leech. L'étape suivante utilise son réseau entier, et non un développement décimal. La construction de l'algèbre d'opérateurs de vertex de réseau s'écrit
\[
 V_\Lambda=M(1)\otimes\mathbb C_\varepsilon[\Lambda].
\]
Ici, $M(1)$ est l'espace des oscillateurs de Heisenberg en rang vingt-quatre ; $\mathbb C_\varepsilon[\Lambda]$ est l'algèbre de groupe tordue par le cocycle de l'extension centrale du réseau. Le cocycle et le relèvement des isométries font partie de la construction : une somme vectorielle nue ne détermine pas les produits de vertex.

La négation $\lambda\mapsto-\lambda$ admet le relèvement standard
$\theta$. Avec son module tordu et le choix de parité compatible,
la construction de Frenkel--Lepowsky--Meurman produit
\begin{equation}
 V^\natural_{\rm HMT}
   =V_\Lambda^+\oplus(V_\Lambda^T)^+.
 \label{exc:flm}
\end{equation}
Le signe $+$ désigne la partie paire de l'action relevée. L'expression
inclut le produit de vertex de l'orbifold ; elle ne se limite pas à une somme
d'espaces gradués.

\begin{theorem}[Prolongement jusqu'au module Moonshine]
\label{exc:moonshine}
L'algèbre \eqref{exc:flm}, construite sur le voisin
\eqref{exc:vecino}, est une réalisation de l'algèbre Moonshine
$V^\natural$. Elle a une charge centrale $24$, une composante de poids un nulle
et un groupe d'automorphismes isomorphe au Monstre $\mathbb M$. Son caractère
gradué est
\begin{equation}
 \operatorname{Tr}_{V^\natural_{\rm HMT}}q^{L_0-1}
   =J(\tau)=j(\tau)-744
   =q^{-1}+196884q+21493760q^2+\cdots,
 \qquad q=e^{2\pi i\tau}.
 \label{exc:j}
\end{equation}
\end{theorem}

\begin{proof}
Le théorème \ref{exc:teorema-leech} vérifie les hypothèses sur le réseau :
positivité, parité, unimodularité, rang vingt-quatre et absence de
racines. On applique alors la construction et le théorème de
Frenkel--Lepowsky--Meurman pour le réseau de Leech. Ceux-ci identifient
l'orbifold à $V^\natural$, établissent son caractère et son groupe d'
automorphismes. La composition du drapeau, du voisin et du foncteur
de construction de VOA donne la réalisation indiquée ici. Le groupe ne se déduit
pas du coefficient $196884$.
\end{proof}

Le théorème utilisé est une reconnaissance classique ultérieure dont les hypothèses sont vérifiées, et non une nouvelle démonstration dans cet article de la construction FLM \cite{flm}. En outre, la notation $q=e^{2\pi i\tau}$ est la variable conventionnelle de Fourier employée pour reconnaître le caractère. Elle n'introduit ni $\pi$ ni $e$ comme entrées du générateur coinductif qui les a déjà déterminés.

Une partie du mécanisme de trace peut être observée avant l'orbifold.
Seul le vecteur nul du réseau est fixé par la négation ; chaque
oscillateur contribue avec un signe alterné. D'où
\[
 \operatorname{Tr}_{V_\Lambda}\theta q^{L_0-1}
   =t_2(\tau)
   :=q^{-1}\prod_{n\geq1}(1+q^n)^{-24}.
\]
Cette trace dans l'algèbre de réseau ne doit pas être confondue avec le
caractère sans insertion de $V^\natural$. Elle n'identifie pas non plus à elle
seule une classe de conjugaison du Monstre : l'espace, l'insertion
et le relèvement doivent être spécifiés avant de comparer les fonctions.

\input{sections/incidencia_automorfismos.tex}

\subsection{Traces graduées de Moonshine et représentation du Monstre}
\label{exc:alcance}

Pour $g\in\operatorname{Aut}(V^\natural_{\rm HMT})$, la trace avec insertion est
\[
 T_g(\tau)=
 \operatorname{Tr}_{V^\natural_{\rm HMT}}gq^{L_0-1}.
\]
Le théorème du Moonshine monstrueux identifie ces séries aux
Hauptmoduln des groupes de genre zéro correspondants
\cite[théorème~1.1]{borcherds}.
Le caractère $J$, le groupe $\mathbb M$ et la collection $g\mapsto T_g$
sont des réalisations de l'algèbre déjà construite. Ils ne forment pas une suite
causale dans laquelle on devinerait d'abord $J$ pour en déduire ensuite le Monstre.

L'unité possède ici une localisation algébrique précise. Le sous-espace de poids zéro est la droite du vide, \(V^\natural_0=\CC\mathbf1\), et le coefficient de \(q^{-1}\) dans \(J\) est donc \(1\). En poids un, la dimension est zéro, d'où l'absence du terme constant dans \(J=j-744\). Un morphisme d'algèbres de vertex conserve le vecteur du vide. Cette propriété fournit le sens technique de la conservation de l'unité à ce niveau ; l'unité des observables cylindriques est conservée par la précomposition démontrée précédemment. La relation entre les deux niveaux doit suivre les applications de construction avec leurs unités.

En poids deux apparaît l'égalité
\[
 196884=1+196883=196560+324=54(5\cdot729+1).
\]
La première séparation correspond à la droite conforme triviale et à la
composante irréductible de dimension $196883$ de la représentation
du Monstre. Les autres expressions ont des lectures en termes de réseaux ou
d'arithmétique. Leur égalité n'identifie pas automatiquement leurs termes
à des régions, des mots ou des routes du TPK. Pour ce faire, il faudrait
présenter les applications entre ces espaces et prouver qu'elles respectent les
actions pertinentes.


En particulier, une isométrie $N_\alpha\simeq\Lambda_{24}$ induit
une réalisation de \eqref{exc:flm}, et un choix d'isomorphisme
avec $V^\natural$ transporte son groupe d'automorphismes. Cela n'a pas
construit une action du Monstre sur les chiffres de
$\pi$, sur les blocs de $K$ ni sur toutes les histoires TPK.
Une telle affirmation supplémentaire requerrait une représentation
$\rho_{\rm TPK}$ et une application $F$ avec un carré explicite
\[
 F\rho_{\rm TPK}(g)=\rho_{V^\natural}(g)F,
\]
en incluant le domaine, le codomaine et l'information conservée. L'absence
de cette étape dans l'affirmation présente ne retire ni l'opérateur d'entrelacement
d'incidence $U_W$, ni le drapeau ni le voisin déjà construits ; elle
délimite seulement une promotion différente, celle de l'action sur l'état source.

La chaîne établie dans cette section peut se résumer, en gardant visibles ses données de transport, comme suit
\begin{align*}
 &\text{APP--TRIT--TPK et état arithmétique avec mémoire de trajectoire}
   \longrightarrow (w_j,f_j;K,Z_0)\\
 &\longrightarrow (C_W,\mathcal H;B_K\subset H^-_\alpha,
                   P_\pi,H_e,H_\varphi,H_{\rm APP})\\
 &\longrightarrow \bigl(N(C_W),o_*,\text{métrique et chambre}\bigr)
   \longrightarrow N_\alpha
   \longrightarrow V^\natural_{\rm HMT}.
\end{align*}
Au niveau de l'incidence, les étoiles engendrent $M_{12}$ et la
coordinatisation résiduelle réalise la structure régionale $4+1$ dans $W_{24}$.
Au niveau des réseaux et des VOA, la construction donne Leech et le module
Moonshine. Ce sont des prolongements d'une même source marquée, avec
des objets et des applications différents. La continuité causale n'oblige pas à
identifier leurs codomaines ni à effacer les choix explicites
de repère, de chambre et de relèvement.

Le résultat de la section est donc une dérivation réunie :
les preuves finies et réticulaires rendent effectif le passage depuis
l’incidence dodécaphasique jusqu’aux hypothèses des théorèmes de
reconnaissance. Ce résultat se rapporte aux constructions et
aux hypothèses spécifiées ; il n’ajoute pas de validation expérimentale
à la comparaison des constantes.

\endgroup

\section{Constantes de Planck et structure déterminantale de l'action}
\label{sec:action}
\subsection{Échelle déterminantale d'action et valuation décimale}

L'échelle absolue et le correcteur relatif proviennent d'objets
distincts. La première utilise le bloc local du registre dodécaphasique
de pas trois,
\begin{equation}
 H_4=\begin{pmatrix}
 1&1&1&1\\1&1&-1&-1\\1&-1&1&-1\\1&-1&-1&1
 \end{pmatrix},\qquad
 G_H=\mathbb Z^4/H_4\mathbb Z^4.
 \label{eq:el-h4}
\end{equation}
L'opération utilise un quotient local. Les trois blocs du registre
complet ne sont pas remplacés par une tensorisation qui multiplierait à
nouveau le nombre de feuillets de ce lecteur.

\begin{lemma}[Nombre de feuillets du quotient local]
La forme normale de Smith de \(H_4\) est
\(\operatorname{diag}(1,2,2,4)\). Par conséquent,
\[
 G_H\simeq\mathbb Z/2\mathbb Z\oplus\mathbb Z/2\mathbb Z
 \oplus\mathbb Z/4\mathbb Z,\qquad |G_H|=16.
\]
\end{lemma}
\begin{proof}
Le plus grand commun diviseur des entrées vaut un. Les mineurs d'ordre deux
sont pairs et l'un d'eux vaut \(-2\), de sorte que leur plus grand commun diviseur
vaut deux. Comme \(H_4H_4^{\mathsf T}=4I\) et \(\det H_4=-16\), les
cofacteurs valent \(\pm4\) ; le troisième diviseur déterminantal vaut quatre.
Le quatrième vaut \(16\). Les quotients successifs des diviseurs
\(1,2,4,16\) sont \(1,2,2,4\). Leur produit donne l'ordre du quotient.
\end{proof}

\begin{definition}[Lecteur déterminantal d'action]
La norme multiplicative de la section scalaire \(\alpha\), constante
sur les feuillets de \(G_H\), et une application de l'auto-échelle
\(\varphi\) produisent
\begin{equation}
 \Sigma_H=\varphi\,\mathrm N_{G_H}(\alpha),\qquad
 \mathrm N_{G_H}(\alpha)=\prod_{g\in G_H}\alpha=\alpha^{16}.
 \label{eq:el-lector-determinantal}
\end{equation}
La coordinatisation décimale associée détermine
\(\operatorname{ord}_{10}(\Sigma_H)=\lfloor\log_{10}\Sigma_H\rfloor\).
\end{definition}

La puissance seize provient ainsi du nombre de feuillets du quotient, une
fois ce lecteur local fixé. Le facteur \(\varphi\) agit une fois sur
la base de la droite d'action. Le calcul de Smith, à lui seul, ne
sélectionnerait pas n'importe quelle fonctionnelle possible sur le quotient : la norme
multiplicative et l'auto-échelle font partie de l'opération déclarée.

\begin{theorem}[Décade de la section d'action]
Le lecteur \eqref{eq:el-lector-determinantal}, appliqué aux coordonnées
HMT déjà engendrées, satisfait
\begin{equation}
 10^{-34}<\Sigma_H<10^{-33},\qquad
 \operatorname{ord}_{10}(\Sigma_H)=-34.
 \label{eq:el-decada}
\end{equation}
\end{theorem}
\begin{proof}
\label{md:accion:intervalos}Il suffit des intervalles rationnels, beaucoup plus grossiers que les cylindres
disponibles,
\[
 \frac{1618}{1000}<\varphi<\frac{1619}{1000},\qquad
 \frac{7297}{10^6}<\alpha<\frac{7298}{10^6}.
\]
La fonction \((x,a)\mapsto xa^{16}\) est croissante en ses deux arguments
positifs. Les deux inégalités entières
\[
 1618\cdot7297^{16}>10^{65},\qquad
 1619\cdot7298^{16}<10^{66}
\]
impliquent
\[
 10^{-34}
 <\frac{1618}{1000}\left(\frac{7297}{10^6}\right)^{16}
 <\Sigma_H
 <\frac{1619}{1000}\left(\frac{7298}{10^6}\right)^{16}
 <10^{-33}.
\]
La définition de l'ordre décimal conclut la preuve. Ni une
unité SI ni une masse de référence n'intervient dans ces inégalités.
\end{proof}

La section normalisée commence par
\(\Sigma_H=1.0462438381\ldots\times10^{-34}\). Sa mantisse ne
s'identifie pas à la coordonnée d'action avant ou après le retour.
Le lecteur déterminantal fixe l'ordre décimal ; les deux coordonnées de
retour s'obtiennent au moyen des lecteurs suivants.


\subsection{Sections d'action et quotient adimensionnel de retour}

Soit \(\mathcal L_S\) une droite orientée d'action et \(\mathcal U_S\)
sa base positive. Le caractère d'ordre cinq et le terme régional de
retour ont pour expressions
\begin{align}
 H_5={}&\frac12\sqrt{\frac{1000\alpha}{\varphi}}-\alpha
   +\frac9{16}\alpha^2-\frac59\alpha^3
   +\frac7{48}\alpha^4-\frac1{54}\alpha^5,
   \label{eq:el-h5}\\
 D_A={}&\exp\!\left(-\frac{100\pi\alpha}{9}\right),
   \label{eq:el-da}\\
 \mathcal C_\pi(\alpha)={}&169\alpha^6+
       \frac{D_A\alpha^7}{1-D_A\alpha},
   \label{eq:el-cpi}\\
 \eta_{\mathrm{ret}}={}&H_5-\frac{90}{\pi}\mathcal C_\pi(\alpha).
   \label{eq:el-eta}
\end{align}
Ce sont des lectures ultérieures de la fermeture dodécaphasique et de son
incidence régionale. L'entier \(169\) est conservé comme le sélecteur
régional déjà construit, et les coefficients de \(H_5\) comme ceux de son
caractère local ; ils ne sont pas remplacés par une interpolation de la constante de
Planck. Les formules explicitent exactement les lectures nécessaires
à l'étape électronique. Il n'est pas nécessaire d'introduire ici d'autres
fonctionnelles de la théorie des masses.

Pour \(0<\alpha<1\), on a \(0<D_A<1\), de sorte que
\(D_A\alpha<1\). Le terme rationnel de \eqref{eq:el-cpi} conserve la
somme complète de la queue géométrique
\begin{equation}
 \frac{D_A\alpha^7}{1-D_A\alpha}
 =\sum_{n=0}^{\infty}D_A^{n+1}\alpha^{n+7}.
 \label{eq:el-cola-regional}
\end{equation}
Il ne s'agit pas d'une troncature que l'on pourrait modifier silencieusement en changeant la
précision. Après le terme d'indice \(N\), la queue restante est
\(D_A^{N+2}\alpha^{N+8}/(1-D_A\alpha)\), ce qui permet de contrôler
l'évaluation sans recourir à une grandeur cible.

\begin{definition}[Sections d'action et lecteur de retour]
Les sections avant et après le retour sont
\begin{equation}
 \hbar_{\mathrm{pre}}=H_5\,10^{-34}\mathcal U_S,
 \qquad
 \hbar_{\mathrm{ret}}=\eta_{\mathrm{ret}}\,10^{-34}\mathcal U_S.
 \label{eq:el-secciones-accion}
\end{equation}
Lorsque \(\eta_{\mathrm{ret}}>0\), on définit
\begin{equation}
 \delta_{\mathrm{ret}}=H_5-\eta_{\mathrm{ret}},\qquad
 R_{\mathrm{act}}=\frac{\hbar_{\mathrm{pre}}}{\hbar_{\mathrm{ret}}}
 =\frac{H_5}{\eta_{\mathrm{ret}}}.
 \label{eq:el-ratio}
\end{equation}
\end{definition}

\begin{proposition}[Positivité et invariance du rapport]\label{md:accion:positividad}
Dans les intervalles utilisés dans la preuve de \eqref{eq:el-decada},
\(0<\eta_{\mathrm{ret}}<H_5\), et donc
\(R_{\mathrm{act}}>1\). Le rapport ne change pas lorsqu’on remplace la base
\(\mathcal U_S\) par toute autre base positive commune.
\end{proposition}
\begin{proof}
Les termes de \(\mathcal C_\pi\) sont positifs. Pour obtenir une borne
élémentaire, il suffit d'utiliser \(\alpha<0.008\), \(\varphi<1.619\),
\(\alpha>0.007\), \(D_A<1\) et \(\pi>3\). La racine dans
\eqref{eq:el-h5} est supérieure à un, tandis que la somme des termes
négatifs est inférieure à \(0.008001\) ; en particulier, \(H_5>0.99\).
De même,
\[
 0<\frac{90}{\pi}\mathcal C_\pi
 <30\left(169(0.008)^6+
           \frac{(0.008)^7}{1-0.008}\right)<10^{-6}.
\]
Ainsi, \(\eta_{\mathrm{ret}}>0.989999>0\) et
\(\eta_{\mathrm{ret}}<H_5\). Enfin, le changement de base divise
les deux coordonnées d'action par le même facteur positif, qui se simplifie dans
leur quotient.
\end{proof}

L'évaluation ultérieure fournit
\begin{align*}
 H_5&=1.0545718176461544696\ldots,\\
 \eta_{\mathrm{ret}}&=1.0545718169150390611\ldots,\\
 R_{\mathrm{act}}&=1.0000000006932817630\ldots.
\end{align*}
La différence et le quotient sont deux lectures des mêmes sections,
l'une additive et l'autre multiplicative. Ils ne représentent pas deux constantes de
Planck indépendantes. Le passage de l'action angulaire à l'action par tour
complet s'exprime, dans chaque section, par
\(h=2\pi\hbar\). La coordinatisation ultérieure
\(\mathcal U_S\mapsto\mathrm{J\,s}\) assigne une unité, mais ne
rétroagit pas sur l'ordre décimal déjà démontré.

Le retour régional peut aussi s'écrire dans une coordinatisation angulaire.
Cette reparamétrisation n'est pas une dépendance supplémentaire indispensable :
\eqref{eq:el-h5}--\eqref{eq:el-eta} fournissent directement
\(H_5\), \(\eta_{\mathrm{ret}}\) et \(R_{\mathrm{act}}\). Si l'on
introduit l'angle et le contre-angle, tous deux doivent être transportés dans la même
coordinatisation ; on ne mélange pas degrés et radians au sein d'une différence.

\subsubsection{Caractère d'action et provenance de ses coefficients}
\label{sec:revision-planck}

La constante de Planck réduite, \(\hbar\), est l'objet physique qui
motive la réalisation de la section d'action angulaire. Sa position dans
la dépendance électronique doit être spécifiée à deux niveaux. La norme
déterminantale \(\Sigma_H\) détermine la décade de l'échelle. Le caractère
\(H_5\) et le retour régional déterminent ensuite les coordonnées des
sections d'action. Cette séparation permet d'exposer le résultat
nécessaire à l'électron sans anticiper les développements complets
d'autres secteurs de la Mécanique Dimensionnelle.

La provenance du caractère d'ordre cinq peut être explicitée au moyen
d'invariants déjà fixés dans le support. Dans la coordinatisation centrale, on considère
\[
 B_c=\begin{pmatrix}7&2\\2&7\end{pmatrix},
 \qquad \lambda_+=9,\quad\lambda_-=5.
\]
Le calendrier est de longueur \(12\), l'orbite du pas \(3\) est de
longueur \(4\), et le quotient de dénombrement pertinent est
\(104976/1944=54\). Les coefficients du caractère satisfont
\begin{equation}
 \frac9{16}=\frac{\lambda_+}{4^2},\quad
 \frac59=\frac{\lambda_-}{\lambda_+},\quad
 \frac7{48}=\frac{\operatorname{tr}(B_c)/2}{4\cdot12},\quad
 \frac1{54}=\frac{1944}{104976}.
 \label{eq:revision-h5-coeficientes}
\end{equation}
L'attribution de ces invariants aux degrés \(2,3,4,5\), avec
l'alternance des signes de \eqref{eq:el-h5}, appartient à la définition du
caractère analytique utilisé. L'expliciter est essentiel : l'ensemble des
cardinaux, considéré sans le caractère ni l'ordre de ses degrés,
admettrait d'autres expressions. Le contenu constructif réside dans la
composition du support, de l'orientation et de la règle analytique déclarée.

\subsubsection{Continuation régionale et équation de renouvellement}

Le sélecteur \(169\) détermine l'impulsion de degré six. Le prolongement
ultérieur s'organise selon le transfert \(D_A\) et une normalisation
unitaire de la droite déterminantale. Ces règles admettent une formulation
en séries formelles avant la substitution \(z=\alpha\), avec \(D_A\) déjà
construit. Écrivons
\[
 \mathscr C(z)=169z^6+R(z),\qquad R(z)\in z^7\mathbb R[[z]].
\]
La concaténation d'un prolongement augmente le degré d'une unité et
multiplie son poids par \(D_A\). Le premier prolongement strict a pour
poids \(D_A\). Son équation de renouvellement est donc
\begin{equation}
 R(z)=D_Az^7+D_AzR(z).
 \label{eq:revision-renovacion}
\end{equation}

\begin{proposition}[Unicité formelle de la continuation normalisée]
L'équation \eqref{eq:revision-renovacion} détermine une unique série et
\[
 \mathscr C(z)=169z^6+\frac{D_Az^7}{1-D_Az}
             =169z^6+\sum_{n=7}^{\infty}D_A^{n-6}z^n.
\]
La réalisation est analytique pour \(|D_Az|<1\).
\end{proposition}
\begin{proof}
L'élément \(1-D_Az\) a pour terme constant un et est inversible
dans l'anneau des séries formelles. Résoudre l'équation produit la série
affichée. De manière équivalente, le coefficient de degré sept vaut \(D_A\)
et chaque coefficient ultérieur vaut \(D_A\) fois le précédent. La série
géométrique converge absolument dans le disque indiqué.
\end{proof}

La normalisation du premier prolongement est une donnée opératoire
indispensable. Si l'on ne conservait que l'impulsion de degré six et
le rapport de concaténation, les séries
\[
 169z^6+\lambda\frac{D_Az^7}{1-D_Az}
\]
partageraient encore ces deux données pour tout \(\lambda\).
La règle unitaire fixe \(\lambda=1\) avant l'évaluation.
La contribution de chaque règle à l'unicité de l'opérateur de lecture
est ainsi localisée, et l'expression rationnelle conserve la continuation complète.

\subsubsection{Constante de Planck réduite et retour de la section d'action}

Dans la coordinatisation dimensionnelle de \eqref{eq:el-secciones-accion}, la section
antérieure au retour est la construction du modèle destinée à la comparaison
avec la constante de Planck réduite :
\[
 \hbar_{\mathrm{pre}}=H_5\,10^{-34}\mathcal U_S.
\]
La section postérieure conserve également la compensation régionale
\(90\mathscr C(\alpha)/\pi\). Toutes deux appartiennent à la même droite
d'action et leur quotient \(R_{\mathrm{act}}\) mesure l'effet multiplicatif
du retour. Le passage à \(h\) par \(2\pi\) utilise la relation
entre la paramétrisation angulaire et un tour complet.

Cette construction intervient effectivement dans
\(R_{\mathrm{act}}^{80-54\alpha+6\Delta_4}\). L'électron conserve
sa coordonnée centrale et sa trajectoire ; l'échelle chronologique ne
remplace pas cette structure par un quantum indifférencié. On ne
multiplie pas non plus à nouveau par \(10^{-34}\) une masse déjà exprimée en MeV :
la décade appartient à la section d'action, se simplifie dans son quotient
sans dimension et ne réapparaît que là où agit la réalisation
dimensionnelle correspondante.

La comparaison numérique de \(H_5\), de la section retournée et de
\(h/(2\pi)\) est présentée dans la section~\ref{sec:revision-planck-contraste}.
On distingue ainsi la dénomination de l'objet physique, l'équation que le
modèle lui attribue et la précision avec laquelle cette équation le reproduit.

\subsection{Confrontation de la constante de Planck réduite}
\label{sec:revision-planck-contraste}

Le Système international en vigueur depuis le 20 mai 2019 fixe
exactement \(h=6.62607015\times10^{-34}\,\mathrm{J\,s}\)
\cite{bipm2018}. Sa constante réduite est donc
\[
 \hbar_{\mathrm{SI}}=\frac{6.62607015}{2\pi}\,10^{-34}\,
 \mathrm{J\,s}
 =1.0545718176461563912\ldots\times10^{-34}\,\mathrm{J\,s}.
\]
Les points de suspension correspondent au développement d'une expression
exacte, et non à une incertitude expérimentale de \(h\) dans le SI actuel.
Cette valeur de référence est utilisée ici après l'évaluation des
lecteurs de HMT ; sa fonction est comparative.

Les deux coordonnées du retour de l'article sont
\begin{align*}
 H_5&=1.0545718176461544696\ldots,\\
 \eta_{\mathrm{ret}}&=1.0545718169150390611\ldots.
\end{align*}
Dans l'identification de \(\mathcal U_S\) à \(\mathrm{J\,s}\), la
différence relative de la section antérieure par rapport à
\(\hbar_{\mathrm{SI}}\) est d'environ \(-1.82\times10^{-15}\).
Celle de la section postérieure est d'environ \(-6.93\times10^{-10}\).
La seconde différence contient l'effet du retour utilisé par le
correcteur électronique ; elle n'équivaut pas à l'erreur d'évaluation de la
première section.

La précision de ces comparaisons permet de reconnaître la proximité
quantitative de la construction de \(H_5\) et, simultanément, de distinguer
proximité et identité exacte. Dans le présent manuscrit, les expressions
présentées ne fournissent pas une égalité exacte avec \(h/(2\pi)\).
Le résultat démontré sur \(\Sigma_H\) est sa décade \(-34\) ;
le résultat analytique sur \(H_5\) est l'évaluation du caractère
déclaré. L'identification physique est examinée dans cette coordinatisation et avec les
différences explicites que nous venons de donner. La valeur SI n'intervient
ni dans la définition de la norme déterminantale ni dans la récurrence régionale.

\subsection{Coordonnées angulaires, normalisation du tour et unités d'action}
\label{sec:ii-angulos}

La grandeur d’action et la coordonnée angulaire appartiennent à des espaces
distincts. La section \(\hbar_{\rm ret}\) occupe la droite d’action
\(\mathcal L_S\) ; \(\eta_{\rm ret}\) est son coefficient dans la base
déclarée. La représentation angulaire agit sur la paire adimensionnelle
\((\alpha,\eta_{\rm ret})\), après sa construction :
\begin{equation}
 \mathcal P(\alpha,\eta_{\rm ret})
 =\bigl([1000\alpha]_{360},[2(\eta_{\rm ret}+\alpha)]_{360}\bigr).
 \label{eq:ii-par-angular}
\end{equation}
La complétion nonadique détermine le facteur 1000 ; la clôture
bidirectionnelle détermine le facteur deux. La coordinatisation circulaire de 360
unités rend visibles les partitions
\(12\cdot30=4\cdot90=3\cdot120\). Sur la branche positive de
l’article, les représentants sont
\[
 A_{\deg}=1000\alpha,\qquad
 C^*_{\deg}=2(\eta_{\rm ret}+\alpha).
\]

\begin{proposition}[Décalage de la représentation en base mille]
\label{prop:ii-desplazamiento-angular}
Soit \(\alpha=\sum_{n\geq1}b_n1000^{-n}\), avec
\(b_n\in\{0,\ldots,999\}\), le développement compatible déjà engendré.
La coordonnée relevée \(A_{\deg}=1000\alpha\) satisfait
\[
 A_{\deg}=b_1+\sum_{n\geq1}b_{n+1}1000^{-n}.
\]
La représentation conserve la suite des blocs à partir du deuxième ;
le bloc initial passe à la partie entière.
\end{proposition}
\begin{proof}
La série converge absolument. Multiplier chaque terme par 1000
et séparer le terme initial fournit l’identité. L’opération
agit sur le développement construit, avec ses blocs et sa mémoire ;
la représentation circulaire ultérieure prend sa classe modulo 360.
\end{proof}

\begin{definition}[Coordinatisations de la coordonnée angulaire]
La coordinatisation en degrés exprime un tour complet au moyen de 360 unités.
La coordinatisation en radians l’exprime au moyen de \(2\pi\). Le changement de
coordonnées et ses inverses sont
\[
 x=\frac{\pi}{180}A_{\deg},\qquad
 y=\frac{\pi}{180}C^*_{\deg},\qquad
 A_{\deg}=\frac{180}{\pi}x,\qquad
 C^*_{\deg}=\frac{180}{\pi}y.
\]
Sur un parcours dont le nombre de tours est conservé, on utilise les
coordonnées relevées ; le quotient circulaire exprime sa classe modulo
le tour correspondant.
\end{definition}

La définition caractérise les degrés et les radians par leur relation exacte
avec le tour. Le registre retient en outre l’orientation, les tours et la
mémoire de transport. La représentation de sa classe angulaire conserve
la phase, tandis que le relèvement distingue des parcours ayant la même
phase finale. La lecture physique d’action nécessite de préciser laquelle de
ces coordonnées intervient dans l’intégrale.

\begin{proposition}[Covariance du lecteur d'action]
Pour \(h_a=2\pi\hbar_a\), l'action d'une coordonnée relevée
satisfait
\begin{equation}
 \mathcal S_a(\theta)
 =\hbar_a\theta_{\rm rad}
 =h_a\frac{\theta_{\deg}}{360}.
 \label{eq:ii-accion-cartas}
\end{equation}
Un tour positif exprime \(h_a\).
\end{proposition}
\begin{proof}
Substituer \(\theta_{\rm rad}=\pi\theta_{\deg}/180\) donne
\(\hbar_a\pi\theta_{\deg}/180
=(2\pi\hbar_a)\theta_{\deg}/360\). Pour un tour,
\(\theta_{\deg}=360\), d'où \(h_a\).
\end{proof}

La périodicité d'une représentation est déclarée séparément. Si la
représentation transportée satisfait \(U(2\pi)=-I\), sa restitution
opératoire se produit en \(4\pi\). Cette propriété du relèvement
coexiste avec la normalisation par tour de
\eqref{eq:ii-accion-cartas}. Chacune conserve son objet : grandeur
intégrée, phase circulaire et état représenté.

\subsection{Covariance des canaux et de la fonctionnelle d'incidence}

Les canaux s'écrivent de manière équivalente
\begin{equation}
 q_\pm=\exp[-(x\pm y)]
 =\exp\!\left[-\frac{\pi}{180}(A_{\deg}\pm C^*_{\deg})\right].
\label{eq:ii-canales}
\end{equation}
\begin{proposition}[Récupération des coordonnées paire et impaire]
\label{prop:ii-inversion-canales-angulares}
Dans la coordinatisation réelle relevée, les canaux positifs déterminent
de manière univoque les deux coordonnées :
\[
 A_{\deg}=-\frac{90}{\pi}\log(q_+q_-),\qquad
 C^*_{\deg}=\frac{90}{\pi}\log\frac{q_-}{q_+}.
\]
L'échange \((q_+,q_-)\mapsto(q_-,q_+)\) conserve
\(A_{\deg}\) et inverse \(C^*_{\deg}\).
\end{proposition}
\begin{proof}
La somme et la différence des logarithmes de
\eqref{eq:ii-canales} sont respectivement
\(-\pi A_{\deg}/90\) et \(-\pi C^*_{\deg}/90\).
Le logarithme réel est univoque sur les canaux positifs.
Le produit est symétrique et le quotient s'inverse lors de leur échange.
\end{proof}
Dans le registre à douze positions, l'involution antipodale
\(r\mapsto r+6\pmod {12}\) forme six paires. La distribution
uniforme de la coordonnée impaire sur ces paires est
\(C^*_{\deg}/6\) : sommer ses six contributions restitue
\(C^*_{\deg}\). Cette coordonnée sectorielle se distingue de
la constante de Planck \(h\), qui appartient à la droite d'action.

La conversion est appliquée une seule fois à chaque argument. La grandeur
\(\alpha\) demeure adimensionnelle et la grandeur \(\hbar\)
appartient toujours à la droite d'action ; leurs coordonnées antérieures
interviennent dans l'application \(\mathcal P\).

La fonctionnelle de Barbero--Immirzi contient, outre les canaux, une
composante bilinéaire dans les représentants angulaires. Son transport est
\begin{equation}
 \frac{\pi}{180}A_{\deg}C^*_{\deg}
 =\frac{180}{\pi}xy.
 \label{eq:ii-barbero-cartas}
\end{equation}
L'identité s'obtient en remplaçant les deux représentants en degrés
par leurs expressions en radians. La formule complète en
coordonnées analytiques est donc
\[
 \gamma_{\HMT}
 =\frac{180}{\pi}xy
 +12\bigl[S_{90}(e^{-x+y})-S_{90}(e^{-x-y})\bigr]
 -\bigl[S_{120}(e^{-x+y})-S_{120}(e^{-x-y})\bigr].
\]
L’équivalence des deux coordinatisations transporte la composante bilinéaire
et les fonctions exponentielles simultanément. Cette transformation
explicite le rôle précis du degré dans la construction et permet
de reproduire l’évaluation sans ambiguïté d’unités angulaires.

\section{Ellipse d’action et réciprocité modulaire}
\label{sec:ellipse}
Les sorties $A,C^*,\hbar>0$ étant fixées dans la coordinatisation
$|C^*/A|<1$, nous définissons
\[
\eta_{\rm el}=\operatorname{artanh}(C^*/A),\quad
a_S=\sqrt{2\hbar}\,e^{\eta_{\rm el}/2},\quad
b_S=\sqrt{2\hbar}\,e^{-\eta_{\rm el}/2}.
\]
Les coordonnées canoniques équilibrées ont pour unité la racine d’action ;
on n’additionne pas ici la position et la quantité de mouvement sans coordinatisation dimensionnelle.
On a
\[
a_Sb_S=2\hbar,\qquad
\tau_{\rm el}=i\,\frac{b_S}{a_S},\qquad
R_{\rm el}=\sqrt{\frac{b_S}{a_S}}
=\left(\frac{A-C^*}{A+C^*}\right)^{1/4}.
\]
La racine positive détermine la composition $\tau_{\rm el}=iR_{\rm el}^2$,
et permet de retrouver le rapport angulaire :
\[
\frac{C^*}{A}=\frac{1-R_{\rm el}^4}{1+R_{\rm el}^4}.
\]
L’inversion de forme envoie $R_{\rm el}$ sur $R_{\rm el}^{-1}$ et
$\tau_{\rm el}$ sur $-1/\tau_{\rm el}$.

\begin{proposition}[Énergie coïncidente, action discernable]
Soit
\[
D(\eta)=\begin{pmatrix}e^\eta&0\\0&e^{-\eta}\end{pmatrix},
\quad
S_2=\begin{pmatrix}0&1\\1&0\end{pmatrix},
\quad
J_2=\begin{pmatrix}0&-1\\1&0\end{pmatrix}.
\]
Les deux applications satisfont $T^{\mathsf T}D(-\eta)T=D(\eta)$.
Cependant, pour $\omega=\dd Q\wedge\dd P$,
\[
S_2^*\omega=-\omega,\qquad J_2^*\omega=\omega.
\]
Si $\mathcal I(\gamma)=\oint_\gamma P\,\dd Q$ sur une courbe fermée,
\[
\mathcal I(S_2\gamma)=-\mathcal I(\gamma),\qquad
\mathcal I(J_2\gamma)=\mathcal I(\gamma).
\]
\end{proposition}
\begin{proof}
La conjugaison échange les entrées diagonales de $D$.
Les déterminants de $S_2,J_2$ sont $-1,1$, respectivement, ce qui
prouve la transformation de $\omega$. Pour $\lambda=P\,\dd Q$,
\[
S_2^*\lambda=\dd(PQ)-\lambda,\qquad
J_2^*\lambda=\lambda-\dd(PQ).
\]
Intégrer la différentielle exacte sur la courbe fermée complète la preuve.
\end{proof}

L’énergie quadratique et le module ne suffisent pas à distinguer les deux
transports. L’orientation de l’action les distingue, même
dans la forme circulaire. Le marquage géométrique fait ainsi partie de la lecture
récupérable, ce n’est pas un ornement ajouté à la valeur.

% Adición separada: no sustituye 07_elipse.tex.
% Procedencia: integral 2249, 52.4.1--2, 52.6.2--3, 53.28 y 90.4.6.
% Recuperación de resultados y explicitación geométrica; certificado local
% verificar_elipse_II.py. No usa un radio físico como entrada generadora.

\subsection{Demi-axes, foyers et rayons intérieurs}
\label{sec:ii-elipse-geometria}

La figure est obtenue à partir des sorties angulaires et d’action des sections
précédentes. APP conserve les feuillets et leurs résidus et quotients ; TRIT fixe
l’orientation ; TPK transporte frontière, retenue et mémoire jusqu’aux
représentations du même état enrichi et de la structure discrète
conjointe du continu. Ici sont reçues la paire \((A,C^*)\) et la section
positive \(\hbar\) après cette construction. La géométrie explicite
quelle information leurs lecteurs conservent, sans utiliser une longueur observée
pour sélectionner l’état source.

Soit \(\xi=C^*/A\in(-1,1)\), avec \(A\ne0\) et les deux angles dans la
même unité. Écrivons \(\eta=\operatorname{artanh}\xi\) et
\[
 a_S=\sqrt{2\hbar}\,e^{\eta/2},\qquad
 b_S=\sqrt{2\hbar}\,e^{-\eta/2},\qquad
 \gamma(\theta)=(a_S\cos\theta,b_S\sin\theta).
\]
Les demi-axes sont étiquetés par leurs coordonnées : si \(\eta<0\),
l’axe vertical est le plus grand. Tous deux appartiennent à la droite de racine
d’action de la coordinatisation équilibrée, non à la droite des longueurs spatiales.

\begin{proposition}[Géométrie focale et intérieur radial]
\label{prop:ii-elipse-focos}
Soient \(a=\max(a_S,b_S)\), \(b=\min(a_S,b_S)\). La demi-distance
focale et l'excentricité sont
\begin{equation}
 f_S=\sqrt{a^2-b^2}
 =\sqrt{4\hbar\sinh|\eta|},\qquad
 e_f=\frac{f_S}{a}=\sqrt{1-e^{-2|\eta|}}.
 \label{eq:ii-elipse-focal}
\end{equation}
Pour \(\eta>0\), les foyers sont \((\pm f_S,0)\) ; pour
\(\eta<0\), ils sont \((0,\pm f_S)\). En \(\eta=0\), ils coïncident avec
le centre. Le rayon de frontière dans la direction polaire \(\phi\) est
\begin{equation}
 \rho_{\partial}(\phi)
 =\left(\frac{\cos^2\phi}{a_S^2}
        +\frac{\sin^2\phi}{b_S^2}\right)^{-1/2}.
 \label{eq:ii-elipse-radio-polar}
\end{equation}
Le segment intérieur dans cette direction est constitué des points
\(r(\cos\phi,\sin\phi)\), avec \(0\le r<\rho_{\partial}(\phi)\).
En particulier, \(b\le\rho_{\partial}(\phi)\le a\).
\end{proposition}

\begin{proof}
Les carrés du grand et du petit demi-axe sont
\(2\hbar e^{|\eta|}\) et \(2\hbar e^{-|\eta|}\). Leur différence
donne \eqref{eq:ii-elipse-focal}. Pour vérifier le foyer, supposons
\(\eta>0\). En \(X=\gamma(\theta)\), les distances à
\(F_+=(f_S,0)\) et \(F_-=(-f_S,0)\) satisfont
\[
 |X-F_+|^2=(a_S-f_S\cos\theta)^2,\qquad
 |X-F_-|^2=(a_S+f_S\cos\theta)^2.
\]
Comme \(a_S>f_S\), leurs racines positives ont pour somme \(2a_S\).
Pour \(\eta<0\), échanger les coordonnées donne la même preuve.
Dans le cas circulaire, \(f_S=0\). Enfin, substituer
\((Q,P)=r(\cos\phi,\sin\phi)\) dans
\(Q^2/a_S^2+P^2/b_S^2=1\) donne
\eqref{eq:ii-elipse-radio-polar}. Le coefficient de \(r^2\) est
compris entre \(a^{-2}\) et \(b^{-2}\), ce qui démontre les bornes et la
caractérisation de l'intérieur.
\end{proof}

\subsection{Phase paramétrique, angle polaire et aire d'action}
\label{sec:ii-elipse-fase-area}

\begin{proposition}[Loi aréale et deux coordonnées angulaires]
\label{prop:ii-elipse-area}
La phase \(\theta\) de \(\gamma\) et son angle polaire continu
\(\phi\), relevés à partir de zéro, satisfont
\begin{equation}
 \phi(\theta)=\arg(a_S\cos\theta+i b_S\sin\theta),\qquad
 \frac{d\phi}{d\theta}
 =\frac{a_Sb_S}{a_S^2\cos^2\theta+b_S^2\sin^2\theta}>0.
 \label{eq:ii-elipse-fases}
\end{equation}
Dans une coordinatisation où les deux tangentes sont définies,
\(\tan\phi=(b_S/a_S)\tan\theta\). L’aire orientée balayée est
\begin{equation}
 \mathcal A(\theta)=\frac12\int_0^\theta(QP'-PQ')\,dt
 =\hbar\theta,\qquad
 \mathcal A(1)=\hbar,\qquad \mathcal A(2\pi)=h.
 \label{eq:ii-elipse-area-fase}
\end{equation}
\end{proposition}

\begin{proof}
La dérivée de l'argument d'un vecteur non nul est
\((QP'-PQ')/(Q^2+P^2)\). Ici, le numérateur est
\(a_Sb_S=2\hbar\), ce qui démontre la dérivée et la monotonie.
Le quotient \(P/Q\) donne la relation des tangentes ; le relèvement
continu fixe les quadrants et le nombre de tours. L'intégrale aréale
vaut \(a_Sb_S\theta/2=\hbar\theta\), et
\(h=2\pi\hbar\) donne la dernière identité.
\end{proof}

La phase paramétrique ne coïncide pas en général avec l’angle polaire. Un radian
de phase balaie \(\hbar\), mais ne s’identifie pas pour autant à un arc
circulaire ni à un secteur d’angle polaire unitaire. En coordonnées polaires,
la même loi s’écrit
\(d\mathcal A=\rho_{\partial}(\phi)^2d\phi/2\).
Avec l’orientation de \(\gamma\) utilisée ici,
\(\oint P\,dQ=-h\) : l’aire positive est \(h\), tandis que
l’intégrale de cette forme de degré un conserve son signe. Cela concorde avec la
distinction entre les transports symplectique et antisymplectique de la
section précédente.

\begin{figure}[htbp]
 \centering
 % Geometría calculada, no coordenadas focales elegidas a ojo.
% Caso algebraico de ilustración xi=1/3; no evaluación de las constantes HMT.
\begingroup
\pgfmathsetmacro{\elXi}{1/3}
\pgfmathsetmacro{\elEta}{0.5*ln((1+\elXi)/(1-\elXi))}
\pgfmathsetmacro{\elA}{exp(\elEta/2)}
\pgfmathsetmacro{\elB}{exp(-\elEta/2)}
\pgfmathsetmacro{\elF}{sqrt(\elA*\elA-\elB*\elB)}
\pgfmathsetmacro{\elTheta}{30}
\pgfmathsetmacro{\elX}{\elA*cos(\elTheta)}
\pgfmathsetmacro{\elY}{\elB*sin(\elTheta)}
\pgfmathsetmacro{\elPolar}{atan(\elY/\elX)}
\begin{tikzpicture}[x=2.18cm,y=2.18cm,font=\small,
 every node/.style={inner sep=2pt},>=Stealth]
 \begin{scope}
  \node[font=\bfseries] at (0,1.38) {Forme et orientation};
  \fill[accent!12] (0,0)--(\elA,0)
    plot[domain=0:30,samples=50] ({\elA*cos(\x)},{\elB*sin(\x)})--cycle;
  \draw[ink,thick] (0,0) ellipse[x radius=\elA,y radius=\elB];
  \draw[->,black!55] (-1.42,0)--(1.48,0) node[right]{\(u\)};
  \draw[->,black!55] (0,-1.05)--(0,1.08) node[above]{\(v\)};
  \draw[accent,thick,-{Stealth}] (\elA,0)
    plot[domain=0:30,samples=50] ({\elA*cos(\x)},{\elB*sin(\x)});
  \draw[accent,thick] (0,0)--(\elX,\elY);
  \node[below left,font=\scriptsize] at (0,0) {\(O\)};
  \draw[accent] (.35,0) arc[start angle=0,end angle=\elPolar,radius=.35];
  \node[accent] at (.47,.065) {\(\phi\)};
  \fill[accent] (\elX,\elY) circle[radius=1.4pt];
  \node[above right] at (\elX,\elY) {\(X(\theta)\)};
  \node[accent,align=center,font=\scriptsize] at (1.26,.18) {\(\theta=\pi/6\)};
  \draw[dashed,black!45] (-\elF,0)--(\elX,\elY);
  \draw[dashed,black!45] (\elF,0)--(\elX,\elY);
  \fill[ink] (-\elF,0) circle[radius=1.5pt] node[below]{\(F_-\)};
  \fill[ink] (\elF,0) circle[radius=1.5pt] node[below]{\(F_+\)};
  \draw[<->,black!60] (-\elA,-1.04)--(0,-1.04)
    node[midway,below,font=\scriptsize]{\(a_S/\sqrt{2\hbar}\)};
  \draw[<->,black!60] (-1.32,0)--(-1.32,\elB)
    node[midway,left,font=\scriptsize]{\(\frac{b_S}{\sqrt{2\hbar}}\)};
  \node[align=center,font=\scriptsize] at (0,-1.55)
    {\(\rho_{\partial}(\phi)/\sqrt{2\hbar}=|OX|\)\\
     \(f_S^2=a_S^2-b_S^2\), \quad aire totale \(h\)};
 \end{scope}
 \begin{scope}[shift={(3.65,0)}]
  \node[font=\bfseries] at (0,1.38) {Forme réciproque};
  \draw[densely dashed,black!25] (0,0) ellipse[x radius=\elA,y radius=\elB];
  \draw[ink,thick] (0,0) ellipse[x radius=\elB,y radius=\elA];
  \draw[->,black!55] (-1.28,0)--(1.38,0) node[right]{\(u\)};
  \draw[->,black!55] (0,-1.30)--(0,1.26) node[right]{\(v\)};
  \fill[ink] (0,\elF) circle[radius=1.5pt] node[right]{\(F'_+\)};
  \fill[ink] (0,-\elF) circle[radius=1.5pt] node[right]{\(F'_-\)};
  \node[align=center,font=\scriptsize] at (0,-1.55)
    {\(\eta\mapsto-\eta\), \quad \(a_S\leftrightarrow b_S\)\\
     \(R_{\rm el}\mapsto R_{\rm el}^{-1}\), \quad \(a_Sb_S=2\hbar\)};
 \end{scope}
 \draw[->,accent,thick] (1.60,.80)--(2.17,.80)
   node[midway,above,font=\scriptsize]{échange};
\end{tikzpicture}
\endgroup

 \caption{Demi-axes, foyers et réciprocité de l'ellipse d'action dans les
 coordonnées normalisées \(u=Q/\sqrt{2\hbar}\),
 \(v=P/\sqrt{2\hbar}\). Le secteur marqué possède la phase paramétrique
 \(\theta=\pi/6\), distincte de son angle polaire \(\phi\) ; son aire
 d'action est \(\hbar\pi/6\). La figure utilise le cas algébrique
 \(\xi=1/3\) pour rendre les relations lisibles : elle n'attribue pas cette valeur
 à la sortie angulaire HMT. Les foyers sont calculés à partir des demi-axes.
 À droite, l'échange des axes conserve l'aire totale \(h\)
 et inverse le rayon modulaire adimensionnel.}
 \label{fig:ii-elipse-calculada}
\end{figure}

Pour le cas de contrôle de la figure, \(\xi=1/3\), les coordonnées
normalisées sont exactement
\[
 \eta=\tfrac12\log2,\quad
 \frac{a_S}{\sqrt{2\hbar}}=2^{1/4}=1.1892071150\ldots,\quad
 \frac{b_S}{\sqrt{2\hbar}}=
 \frac{f_S}{\sqrt{2\hbar}}=R_{\rm el}=2^{-1/4}=0.8408964153\ldots.
\]
Ce sont des nombres de contrôle de non-régression du dessin, non des valeurs attribuées au générateur
angulaire. Le certificat local vérifie également la forme circulaire et les
orientations \(\xi=\pm1/3,\pm3/5\).

\subsection{Réciprocité de forme et distance intérieure}

La représentation rectangulaire retrouve
\begin{equation}
 R_{\rm el}=\sqrt{b_S/a_S}=e^{-\eta/2},\qquad
 R_{\rm el}^4=\frac{1-\xi}{1+\xi},\qquad
 \xi=\frac{1-R_{\rm el}^4}{1+R_{\rm el}^4}.
 \label{eq:ii-elipse-reciproca}
\end{equation}
La preuve consiste à substituer
\(2\eta=\log((1+\xi)/(1-\xi))\). L'échange des axes
envoie \(\eta\) sur \(-\eta\), \(R_{\rm el}\) sur
\(R_{\rm el}^{-1}\) et \(\tau=iR_{\rm el}^2\) sur \(-1/\tau\).
Le produit \(a_Sb_S\) est conservé. Le rayon \(R_{\rm el}\) est adimensionnel ;
le rayon polaire \(\rho_{\partial}\) a pour unité la racine carrée d'action.

L'intérieur admet également la distance projective de Hilbert. Si
\(z_-,x,y,z_+\) sont ordonnés sur une corde de l'ellipse, on définit
\[
 d_H(x,y)=\frac12\log
 \frac{|y-z_-|\,|x-z_+|}{|x-z_-|\,|y-z_+|}.
\]
C'est une distance adimensionnelle : les unités se simplifient dans le birapport.
La normalisation \((Q,P)\mapsto(Q/a_S,P/b_S)\) envoie l'ellipse
sur le disque unité et préserve ce rapport, car les longueurs sur une
même droite sont multipliées par un facteur commun. Elle transporte ainsi la métrique
vers le modèle de Beltrami--Klein. Ni cette distance intérieure ni la distance
entre les foyers ne s'identifie à une longueur atomique.

\subsection{Rayon de Bohr : conséquence de l'échelle d'action}
\label{sec:ii-elipse-bohr}

Dans la réalisation électrodynamique, l'échelle de Bohr s'exprime par
\(a_{\rm B}=\hbar/(m_ec\alpha)\) \cite{codata2022}. Dans la chaîne
du présent article, cette lecture succède aux constructions de
l'action, du vide, de la structure fine et de la masse électronique. On reçoit ici
\(\hbar,m_e,c,\alpha>0\) dans une même réalisation dimensionnelle.
\begin{proposition}[Expression aréale de l'échelle de Bohr]
\label{prop:ii-elipse-bohr}
Dans cette réalisation,
\begin{equation}
 a_{\rm B}=\frac{a_Sb_S}{2m_ec\alpha}
 =\frac{\mathcal A(2\pi)}{2\pi m_ec\alpha}
 =\frac{\bar\lambda_C(m_e)}{\alpha},\qquad
 \bar\lambda_C(m_e)=\frac{\hbar}{m_ec}.
 \label{eq:ii-elipse-bohr}
\end{equation}
\end{proposition}
\begin{proof}
Les identités \(a_Sb_S=2\hbar\) et
\(\mathcal A(2\pi)=2\pi\hbar\), substituées dans l'expression
de Bohr, donnent les deux premiers membres. La définition de
\(\bar\lambda_C(m_e)\) donne le troisième.
\end{proof}

Le numérateur a la dimension d'une action et \(m_ec\) celle d'une quantité de mouvement ; le
quotient est une longueur. Cette composition utilise le produit des
demi-axes et les trois autres grandeurs : elle n'affirme pas que Bohr soit un demi-axe,
un foyer, une distance de Hilbert ou le rayon modulaire. Elle n'assimile pas non plus
la masse électronique au quantum circulaire d'ordre 108. La notation
\(a_{\rm B}\) évite de confondre cette longueur avec la normalisation
\(a_0\) de l'opérateur d'aire.

\subsection{Réalisation inductive--capacitive et continuité constitutive}
\label{sec:ii-elipse-lc}

Le pont électrodynamique de l’ellipse conserve une coordinatisation supplémentaire :
une horloge \(\omega>0\) et une impédance de référence \(Z_*>0\).
La construction retrouvée dans le corpus définit
\begin{equation}
 L_\eta=\frac{Z_*}{\omega}e^{-\eta},\qquad
 C_\eta=\frac{1}{Z_*\omega}e^\eta,
 \qquad Z_\eta=\sqrt{L_\eta/C_\eta}.
 \label{eq:ii-elipse-lc}
\end{equation}
Ici, \(C_\eta\) est la capacité électrique, distincte du contre-angle \(C^*\).
\begin{proposition}[Produit dynamique et quotient constitutif]
\label{prop:ii-elipse-lc}
Dans la coordinatisation précédente,
\begin{equation}
 L_\eta C_\eta=\omega^{-2},\qquad
 \frac{Z_\eta}{Z_*}=e^{-\eta}
 =\frac{b_S}{a_S}=R_{\rm el}^{2}.
 \label{eq:ii-elipse-impedancia}
\end{equation}
Avec \(Q=\sqrt{Z_*}\,q\), \(P=\Phi/\sqrt{Z_*}\), l'énergie
\(H_{LC}=q^2/(2C_\eta)+\Phi^2/(2L_\eta)\) prend la forme
\[
 H_{LC}=\frac\omega2(e^{-\eta}Q^2+e^\eta P^2).
\]
Sur \(\gamma(\theta)\), ses composantes électrique et magnétique sont
\(U_E=\hbar\omega\cos^2\theta\) et
\(U_B=\hbar\omega\sin^2\theta\) ; leur somme vaut \(\hbar\omega\).
\end{proposition}
\begin{proof}
Dans le produit de \eqref{eq:ii-elipse-lc}, les exponentielles et
\(Z_*\) se simplifient ; son quotient vaut \(Z_*^2e^{-2\eta}\). Prendre la racine
positive et utiliser \eqref{eq:ii-elipse-reciproca} donne
\eqref{eq:ii-elipse-impedancia}. Substituer \(q=Q/\sqrt{Z_*}\)
et \(\Phi=\sqrt{Z_*}P\) dans \(H_{LC}\) donne sa forme équilibrée.
Enfin, les demi-axes de \(\gamma\) simplifient les facteurs
\(e^{\mp\eta}\), et \(\cos^2\theta+\sin^2\theta=1\).
\end{proof}

Ce pont distingue la forme, l'horloge et le transducteur. Le choix
du sens dynamique doit s'accorder avec la convention symplectique : pour
les équations \(\dot Q=\partial_PH\),
\(\dot P=-\partial_QH\), la paramétrisation précédente satisfait
\(\dot\theta=-\omega\). Les identités énergétiques ne dépendent pas
de l'inversion de ce parcours.

Le lecteur résolvant du vide agit sur les canaux déjà produits
\(q_\pm=e^{-(A_{\rm rad}\pm C^*_{\rm rad})}\) et l’incidence
\(90/120\). Ses réponses
\(r_\pm=(1-q_\pm^{90})/(1-q_\pm^{120})\) expriment
\(\widehat Z=r_-/r_+\) et \(\widehat c=(r_+r_-)^{-1}\).
Ce sont des réponses spectrales, non des rayons géométriques. L’identité
\eqref{eq:ii-elipse-impedancia} fixe ici la réalisation LC ;
l’identifier à l’impédance constitutive du vide conserve l’application
entre les deux lecteurs et leurs coordinatisations. Leurs valeurs ne sont pas égalées du fait qu’elles partagent
l’antécédent angulaire.

\section[Retour de feuillet, incidence aréale--volumétrique et orientation constitutive]{Retour de feuillet, incidence aréale--volumétrique\\et orientation constitutive}
\label{iii:sec:hojas}

La structure des feuillets intervient avant l’évaluation des grandeurs
électromagnétiques. Sa fonction consiste à distinguer la fermeture d’une projection,
la restitution d’une orientation et l’évolution du registre transporté.
Ces opérations agissent sur des objets liés, mais de types différents : une
trajectoire projetée peut s’être refermée tandis que son relèvement reste
sur le feuillet opposé ; l’orientation peut avoir été restituée tandis que la mémoire
conserve les deux tours effectués. La construction constitutive doit recevoir
cette information avec la structure qui la détermine.

\subsection{Retour projeté et restitution du revêtement orienté}

Soit $\ell$ une boucle de retour d’une extension survivante et soit
$\varepsilon:\langle\ell\rangle\to\{+1,-1\}$ le caractère de son
relèvement, avec $\varepsilon(\ell)=-1$. Le fibré d’intervalles associé
a pour présentation
\begin{equation}
 \mathcal M=([0,1]\times[-1,1])/{(0,u)\sim(1,-u)}.
 \label{iii:eq:mobius}
\end{equation}
La donnée décisive est le signe du transport sur la boucle spécifiée.
L’identification des extrémités réalise géométriquement ce caractère.

\begin{proposition}[Deux ordres de clôture]
Un tour de la boucle ferme la base et inverse l’orientation de la fibre. Deux
tours restituent l’orientation. Dans la paramétrisation angulaire du lecteur de
bord, ces retours correspondent respectivement à la clôture aréale $2\pi$ et
à la clôture volumétrique $4\pi$ du revêtement orienté.
\end{proposition}
\begin{proof}
La relation d’équivalence de \eqref{iii:eq:mobius} identifie la fibre terminale
à la fibre initiale au moyen de $u\mapsto-u$. La composition de deux transports est
$u\mapsto-(-u)=u$. Paramétrer un tour de la base par $2\pi$ attribue $4\pi$
à son carré. L’assertion concerne ce relèvement et cette lecture.
\end{proof}

La mémoire du TPK conserve en outre la trajectoire parcourue et les mises à jour du registre. Le carré du caractère de signe est l'identité, mais n'impose pas l'annulation de ces mises à jour. La représentation spinorielle, lorsqu'elle sera utilisée, devra transporter cette donnée au moyen de sa propre application.

\subsection{Descente de l’incidence depuis le calendrier tritique}

Dans le bloc primitif $(+1,-1,0)$, le TPK met à jour le mode arithmétique au moyen de
$+1:m\mapsto\Sigma$, $-1:m\mapsto\Pi$ et $0:m\mapsto m$. Son orbite cyclique
est $(\Sigma,\Pi,\Pi)$. Le lecteur de feuillets attribue $\Sigma$ au feuillet aréal et
$\Pi$ au feuillet volumétrique. Dans l’ordre des bases $(e_{\rm ar},e_{\rm vol})$,
l’incidence des arêtes est représentée par
\begin{equation}
 F_{\rm av}=\begin{pmatrix}0&1\\1&1\end{pmatrix}.
 \label{iii:eq:fav}
\end{equation}
Les trois paires consécutives sont $\Sigma\to\Pi$, $\Pi\to\Pi$ et $\Pi\to\Sigma$, chacune de multiplicité un. Les calendriers de longueurs $9$, $27$ et $108$ contiennent $3$, $9$ et $36$ copies du bloc primitif. Leurs matrices d'incidence sont donc ces multiples de $F_{\rm av}$.

Cette construction conserve la différence entre fréquence chronologique et mesure
de l’espace des continuations. La première assigne des poids $(1/3,2/3)$ au cycle
des modes. La seconde se définit sur l’enveloppe symbolique minimale de mémoire
un qui admet exactement les couples précédents. Sa matrice d’adjacence est
$F_{\rm av}$ ; les dénombrements de trajectoires appartiennent à cette enveloppe.

\Needspace{10\baselineskip}
\begin{proposition}[Poids de retour des deux feuillets]
Soient $F_0=0$, $F_1=1$ et $F_{n+1}=F_n+F_{n-1}$. Pour $n\geq1$,
\begin{equation}
 F_{\rm av}^{n}=
 \begin{pmatrix}F_{n-1}&F_n\\F_n&F_{n+1}\end{pmatrix}.
 \label{iii:eq:fav-powers}
\end{equation}
La mesure de Parry de cette enveloppe attribue aux feuillets des poids proportionnels à $(1,\varphi^2)$, et le rapport des décomptes diagonaux converge vers $\varphi^{-2}$.
\end{proposition}
\begin{proof}
La formule des puissances se prouve par induction en utilisant la récurrence.
Le carré de $F_{\rm av}$ est strictement positif. Son vecteur positif de
Perron est $(1,\varphi)$, de valeur propre $\varphi$, car
$\varphi^2=\varphi+1$. La matrice étant symétrique, les poids de Parry sont
proportionnels aux carrés des composantes de ce vecteur. L’autre racine
du polynôme caractéristique est $-\varphi^{-1}$ ; son module est strictement
inférieur à $\varphi$, ce qui donne la limite des quotients diagonaux.
\end{proof}

L’identification spectrale de cette incidence est postérieure à sa construction
à partir du calendrier. La coordonnée régionale $\varphi$ conserve sa généalogie
coinductive antérieure : la reconnaissance de la valeur propre ne sélectionne pas le calendrier.

\subsection{Marquage de la frontière et rapport surfacique--volumétrique}

Soit $\Pi_{\rm HMT}(X_\infty)$ la valeur du lecteur régional de clôture déjà construit. Avec les projecteurs diagonaux de feuillet, on définit
\begin{equation}
 D_\pi=\Pi_{\rm HMT}(X_\infty)P_{\rm ar}+P_{\rm vol},\qquad
 \Theta_n=
 \frac{\langle e_{\rm ar},D_\pi F_{\rm av}^{n}e_{\rm ar}\rangle}
 {\langle e_{\rm vol},F_{\rm av}^{n}e_{\rm vol}\rangle}.
 \label{iii:eq:marked-return}
\end{equation}
La normalisation volumétrique est unitaire ; la clôture régionale marque une seule fois
le retour aréal. En appliquant \eqref{iii:eq:fav-powers},
\begin{equation}
 \Theta_n=\Pi_{\rm HMT}(X_\infty)\frac{F_{n-1}}{F_{n+1}}
 \longrightarrow\frac{\pi}{\varphi^2}.
 \label{iii:eq:av-limit}
\end{equation}
Le rapport de retour est déterminé par l'incidence, la mesure et le marquage. Le quotient $2\pi/4\pi$ décrit, en revanche, les deux ordres de clôture du revêtement. Les deux relations participent à la lecture des feuillets sans se substituer l'une à l'autre.

\subsection{Transport de l’orientation aux canaux constitutifs}

Dans une représentation à deux canaux, soit $\mathcal R$ une involution
autoadjointe et $\mathcal S$ un opérateur unitaire d’échange qui satisfait
$\mathcal S\mathcal R\mathcal S^{-1}=-\mathcal R$. Les coordonnées angulaires
déjà construites s’écrivent $x=A_{\rm rad}$ et $y=C^*_{\rm rad}$, avec $x>|y|$.
Alors
\begin{equation}
 B=xI+y\mathcal R,\qquad T=e^{-B},\qquad
 \mathcal ST(y)\mathcal S^{-1}=T(-y).
 \label{iii:eq:channel-conjugation}
\end{equation}
En effet, la conjugaison transforme $B(y)$ en $B(-y)$ et commute avec sa série
exponentielle convergente. Les valeurs propres $q_\pm=e^{-x\mp y}$ sont échangées.
Pour fixer l’ordre des réponses, on utilise la chambre $x>y>0$, où
$0<q_+<q_-<1$. L’application $(x,y)\mapsto(x,-y)$ la transporte dans la
chambre opposée $x>-y>0$, où $0<q_-<q_+<1$. Toutes deux appartiennent au domaine
contractif $x>|y|$. À la frontière $y=0$, les deux canaux coïncident.
La conjugaison échange les chambres, elle ne conserve pas l’ordre strict de
leurs étiquettes ; elle agit sur les projecteurs au moyen de
$\mathcal SP_\pm\mathcal S^{-1}=P_\mp$, avec
$P_\pm=(I\pm\mathcal R)/2$.
Cette application concrète transmet l’orientation à l’opérateur constitutif
développé ci-après.

\section{Incidence hexade–octade et fonctionnelle orientée}
\subsection{Incidence dodécaphasique et fonctionnelle de retour}

La contribution de Barbero--Immirzi utilise le transport angulaire et
l’incidence hexade--octade. Soient $\mathsf H\subset\mathsf O$ une hexade
et une octade avec la paire marquée de l’hexade conservée. Les ensembles de
drapeaux pertinents sont
\begin{equation}
 \mathcal F_6=\mathsf H\times\binom{\mathsf H}{2}
 \lhook\joinrel\longrightarrow
 \mathcal F_8=\mathsf O\times\binom{\mathsf H}{2},
 \qquad |\mathcal F_6|=90,\quad |\mathcal F_8|=120.
 \label{iii:eq:barbero-flags}
\end{equation}
Les cardinaux sont $6\binom62$ et $8\binom62$ ; la même incidence
conserve le défaut normalisé
$(90-120)/(24\binom62)=-1/12$. Ce dénombrement identifie les deux secteurs
employés par le retour, tandis que le calendrier détermine leurs
multiplicités relatives.

Pour $m\in\{90,120\}$, le retour tritique sur un canal contractant est
\begin{equation}
 S_m(q)=\sum_{j\geq0}q^{m+3mj}
       =\frac{q^m}{1-q^{3m}},\qquad 0<q<1.
 \label{iii:eq:barbero-return-series}
\end{equation}
Chaque terme correspond à la hauteur initiale $m$ et à $j$ retours de
longueur $3m$. La série géométrique conserve donc tant la hauteur
initiale que la période du retour. Le calendrier dodécaphasique
$\mathbb Z/12\mathbb Z$ apporte douze secteurs et une unique incidence
orientée de frontière, notée $\partial_{\rm ret}$. Dans le module
libre des événements, sa chaîne fondamentale est
\begin{equation}
 c_{12}^{+}
 =\sum_{j\in\mathbb Z/12\mathbb Z}[j,\mathcal F_6]
  -[\partial_{\rm ret},\mathcal F_8].
 \label{iii:eq:barbero-cycle}
\end{equation}
Le lecteur linéaire $\mathscr L_q$ assigne $S_{90}(q)$ à chaque générateur
$[j,\mathcal F_6]$ et $S_{120}(q)$ au générateur de frontière. On obtient
\begin{equation}
 \mathscr L_q(c_{12}^{+})=12S_{90}(q)-S_{120}(q).
 \label{iii:eq:barbero-weights}
\end{equation}
Les poids $12$ et $-1$ sont les multiplicités orientées de la chaîne.
L’orientation conjuguée satisfait $c_{12}^{-}=-c_{12}^{+}$, avec les
mêmes types d’incidence.

\begin{proposition}[Coefficient singulier du retour fondamental]
\label{iii:prop:barbero-pole}
Le coefficient de $(1-q)^{-1}$ de
$\mathscr L_q(c_{12}^{+})$ est $1/24$.
\end{proposition}
\begin{proof}
Pour chaque entier $m>0$, la factorisation
$1-q^{3m}=(1-q)\sum_{j=0}^{3m-1}q^j$ donne
\[
 \lim_{q\uparrow1}(1-q)S_m(q)=\frac1{3m}.
\]
Appliquer cette identité à la fonctionnelle déjà déterminée dans
\eqref{iii:eq:barbero-weights} produit
\begin{equation}
 \lim_{q\uparrow1}(1-q)\mathscr L_q(c_{12}^{+})
 =\frac{12}{270}-\frac1{360}=\frac1{24}.
 \label{iii:eq:barbero-pole}
\end{equation}
Dans la coordonnée complexe $q-1$, le résidu est $-1/24$.
\end{proof}

L’évaluation singulière vérifie une conséquence de l’incidence et du
retour tritique. L’ordre démonstratif est fixé par
\eqref{iii:eq:barbero-flags}--\eqref{iii:eq:barbero-weights} ; le
coefficient singulier se calcule après la construction de la chaîne.


\section{Opérateur constitutif et reconstruction des canaux du vide}
\label{iii:sec:constitutiva}

Les grandeurs constitutives s'obtiennent au moyen d'une opération sur le transport angulaire. L'information de départ comprend les canaux, leurs projecteurs et l'orientation qui les distingue. Cette organisation permet de suivre séparément la construction de l'opérateur, l'évaluation de ses fonctions spectrales et sa réalisation dimensionnelle ultérieure.

\subsection{Réponse positive sur deux feuillets}

Dans la chambre orientée $x>y>0$ de la section~\ref{iii:sec:hojas}, soient
$P_\pm=(I\pm\mathcal R)/2$ et
$T=q_+P_++q_-P_-$, avec $0<q_+<q_-<1$. Les incidences construites par le
calendrier fixent les secteurs $90$ et $120$. La réponse constitutive est
\begin{equation}
 \mathcal V=(I-T^{90})(I-T^{120})^{-1}.
 \label{iii:eq:vacuum-op}
\end{equation}
Cette définition reçoit le transport provenant du TPK et conserve les
attributions de feuillet. La règle de réponse et son domaine font explicitement partie
de la construction.

Le quotient provient de la comparaison des deux incidences du calendrier sur un même transport, et non de l'ajustement séparé de quatre grandeurs. Si $S=T^{30}$ et $\Sigma_j(S)=I+S+\cdots+S^{j-1}$, la factorisation de $I-S^j$ donne
\begin{equation}
 \mathcal V\,\Sigma_4(S)=\Sigma_3(S),\qquad
 \mathcal V=\Sigma_3(S)\Sigma_4(S)^{-1}.
 \label{iii:eq:completion-response}
\end{equation}
Les facteurs $\Sigma_3$ et $\Sigma_4$ expriment respectivement les
secteurs $90=3\cdot30$ et $120=4\cdot30$. Comme le spectre de $S$ est
contenu dans $(0,1)$, $\Sigma_4(S)$ est inversible : l’équation de
complétion détermine une unique réponse pour le transport fixé.
Cette unicité concerne la règle constitutive déclarée et ses
incidences. L’identification électromagnétique ajoute l’attribution
orientée des feuillets et les droites dimensionnelles précisées plus
bas.

\begin{proposition}[Décomposition et positivité]
La réponse est bien définie et
\begin{equation}
 \mathcal V=r_+P_++r_-P_-,\qquad
 r_\pm=\frac{1-q_\pm^{90}}{1-q_\pm^{120}}.
 \label{iii:eq:rchannels}
\end{equation}
Ses valeurs propres satisfont $3/4<r_-<r_+<1$ dans cette chambre.
\end{proposition}
\begin{proof}
Chaque coefficient $1-q_\pm^{120}$ est positif. L'inverse de $I-T^{120}$ est la somme des projecteurs divisés par ces coefficients, ce qui prouve \eqref{iii:eq:rchannels}. Avec $s=q^{30}$, on obtient
\begin{equation}
 f(s)=\frac{1+s+s^2}{1+s+s^2+s^3},\qquad
 f'(s)=-\frac{s^2(3+2s+s^2)}{(1+s+s^2+s^3)^2}<0.
 \label{iii:eq:monotone}
\end{equation}
Les limites en $0$ et en $1$ sont $1$ et $3/4$, respectivement. La monotonie
et l’ordre des canaux concluent la preuve.
\end{proof}

\subsection{Quatre coordonnées d’une réponse commune}

On définit les évaluations normalisées
\begin{equation}
 \widehat\varepsilon=r_+^2,\qquad
 \widehat\mu=r_-^2,\qquad
 \widehat Z=\frac{r_-}{r_+},\qquad
 \widehat c=\frac1{r_+r_-}.
 \label{iii:eq:four}
\end{equation}
Le premier couple conserve les réponses quadratiques assignées à chaque feuillet.
Le produit recueille le mode commun ; le quotient distingue leur orientation relative.
On vérifie exactement
\begin{equation}
 \widehat\mu\widehat\varepsilon=\widehat c^{-2},\qquad
 \frac{\widehat\mu}{\widehat\varepsilon}=\widehat Z^2,
 \qquad
 r_+=\frac1{\sqrt{\widehat Z\widehat c}},\quad
 r_-=\sqrt{\frac{\widehat Z}{\widehat c}}.
 \label{iii:eq:four-identities}
\end{equation}
Les racines positives fixent l'inversion. Sous la conjugaison des feuillets, lue par rapport aux marques $P_\pm$ fixes, $\widehat\varepsilon$ et $\widehat\mu$ sont échangés, $\widehat Z$ est inversé et $\widehat c$ est conservé. La réponse est alors transportée vers la chambre $x>-y>0$ : les formules restent valables, avec les ordres $q_-<q_+$ et $r_+<r_-$ inversés. Par conséquent, le mode commun et l'orientation sont des données récupérables au moyen d'observables complémentaires.

\begin{proposition}[Détermination des évaluations quadratiques]
Soit $W=r_+P_++r_-P_-$ une réponse positive sur les deux feuillets.
Étant fixées les lectures de produit et de quotient
\[
 \widehat c^{-1}=\det W=r_+r_-,
 \qquad \widehat Z=r_-/r_+,
\]
il existe une unique paire positive $(\widehat\varepsilon,\widehat\mu)$
qui vérifie
$\widehat\mu\widehat\varepsilon=\widehat c^{-2}$ et
$\widehat\mu/\widehat\varepsilon=\widehat Z^2$.
C’est précisément $(r_+^2,r_-^2)$.
\end{proposition}
\begin{proof}
La positivité permet de prendre les racines sans ambiguïté :
\[
 \widehat\mu=\frac{\widehat Z}{\widehat c}
 =\frac{r_-}{r_+}r_+r_-=r_-^2,\qquad
 \widehat\varepsilon=\frac1{\widehat Z\widehat c}
 =\frac{r_+}{r_-}r_+r_-=r_+^2.
\]
Les expressions vérifient les deux relations et en sont les seules racines
positives.
\end{proof}

Ainsi, les carrés ne constituent pas deux choix indépendants ajoutés
aux canaux : ils sont déterminés par les lectures multiplicative et
orientée, une fois fixées les relations constitutives. Cette proposition
établit l’unicité au sein de ce système de règles. L’interprétation
des règles comme réponse électromagnétique conserve son contenu
supplémentaire de représentation physique.

La lecture déterminantale du transport et la réponse constitutive
doivent également conserver leur ordre. Pour $T=e^{-xI-y\mathcal R}$, on a
$\det T=e^{-2x}$, tandis que
\[
 \det\mathcal V=f(q_+^{30})f(q_-^{30})=\widehat c^{-1}.
\]
Les deux proviennent du même opérateur à deux feuillets, mais évaluent des fonctions différentes de son spectre. En général, appliquer d'abord le déterminant puis une fonction rationnelle n'équivaut pas à prendre le déterminant de cette fonction de l'opérateur. Maintenir les deux valeurs propres jusqu'à l'achèvement de la réponse conserve l'information d'orientation qu'un déterminant isolé ne distingue plus.

\subsection{Inversion cubique et récupération angulaire}

\begin{theorem}[Inversion constitutive dans le domaine contractif]
Pour chaque $r\in(3/4,1)$, il existe une unique $s\in(0,1)$ telle que $f(s)=r$.
Cette $s$ est l’unique racine dans cet intervalle de
\begin{equation}
 r s^3+(r-1)(1+s+s^2)=0.
 \label{iii:eq:cubic}
\end{equation}
La réponse étiquetée reconstruit $s_\pm$, $q_\pm=s_\pm^{1/30}$ et
\begin{equation}
 x=-\frac12\log(q_+q_-),\qquad
 y=\frac12\log\frac{q_-}{q_+}.
 \label{iii:eq:recover-angular}
\end{equation}
\end{theorem}
\begin{proof}
La dérivée et les limites de \eqref{iii:eq:monotone} prouvent que $f$ est une
bijection décroissante entre les intervalles indiqués. Multiplier $f(s)=r$
par son dénominateur produit \eqref{iii:eq:cubic}. La racine positive d’ordre
$30$ retrouve le canal. Enfin, la somme et la différence de
$\log q_\pm=-x\mp y$ donnent \eqref{iii:eq:recover-angular}.
\end{proof}

La reconstruction conserve les étiquettes de feuillet. Avec l’opérateur complet
$\mathcal V$, on conserve aussi les projecteurs et on applique la fonction
inverse par calcul spectral. Récupérer deux coordonnées angulaires est une
opération définie sur cette représentation ; la trace historique complète
du TPK conserve des données supplémentaires.

L’intervalle $r\in(3/4,1)$ et la fonction inverse de ce théorème appartiennent
à la normalisation interne de \eqref{iii:eq:four}. Si l’on dispose de
coordonnées exprimées dans une autre coordinatisation d’unités, on retrouve d’abord
cette normalisation au moyen de la transformation de la
section~\ref{iii:sec:evaluacion} ; on applique ensuite l’inversion cubique.

% Recuperación expositiva de INTERRELACIONES_ESTRUCTURALES_VACIO.md, §§6--8.
% Insertar después de la inversión cúbica y antes de las coordenadas logarítmicas.
\subsection{Composition transportée des réponses constitutives}
\label{amp-comp-sec}

Le transport angulaire et les incidences $90$ et $120$ proviennent de la
même structure APP--TRIT--TPK : les feuillets conservent les évaluations et
leurs retenues ; le régime tritique conserve l’orientation ; le calendrier
TPK transporte la phase et la mémoire jusqu’aux canaux angulaires. La réponse
constitutive applique à ces canaux la fonction $f$ de
\eqref{iii:eq:monotone}. Son inversion permet d’exprimer, dans les coordonnées
de réponse, la composition des transports du même repère. On conserve
ainsi l’ordre constructif : transport, composition et évaluation spectrale.
L’opération obtenue agit sur cette représentation de l’état enrichi ;
son inverse retrouve les canaux, tandis que l’histoire complète de mémoire
demeure dans la structure de provenance.

Soient $s=q^{30}$ et $\psi=f^{-1}$. L’exposant $30$ est le plus grand commun
diviseur des incidences $90$ et $120$. L’inversion cubique détermine
$\psi:(3/4,1)\longrightarrow(0,1)$. En adjoignant la réponse de frontière
$f(1)=3/4$, nous étendons la bijection à
\[
 \psi:D\longrightarrow(0,1],\qquad D=[3/4,1).
\]
Sur cette frontière, on utilise la fonction rationnelle $f$ ; le quotient
$(1-q^{90})/(1-q^{120})$ y a une singularité éliminable.

\begin{proposition}[Monoïde des réponses et coordonnée additive]
\label{amp-comp-monoid}
L’opération
\begin{equation}
 r\odot t=f\bigl(\psi(r)\psi(t)\bigr),\qquad r,t\in D,
 \label{amp-comp-law}
\end{equation}
fait de $D$ un monoïde commutatif et simplifiable d’élément neutre $3/4$.
L’application
\begin{equation}
 \ell(r)=-\frac1{30}\log\psi(r)
 \label{amp-comp-character}
\end{equation}
est un isomorphisme de ce monoïde avec $(\mathbb R_{\geq0},+)$.
L’intervalle $(3/4,1)$ est stable par $\odot$ ; le neutre appartient à
l’extension de frontière. Le seul élément inversible dans $D$
est le neutre.
\end{proposition}
\begin{proof}
Le produit de deux éléments de $(0,1]$ appartient au même intervalle et
\[
 \psi(r\odot t)=\psi(r)\psi(t).
\]
Pour trois réponses, les deux ordres d’association ont pour image
$\psi(r)\psi(t)\psi(u)$ ; l’injectivité de $\psi$ prouve
l’associativité. La commutativité provient du produit scalaire. Si
$r\odot t=r\odot u$, la positivité de $\psi(r)$ permet de simplifier et
d’obtenir $\psi(t)=\psi(u)$, donc $t=u$. Le paramètre $1$ correspond à
$3/4$ et fournit l’identité. Le logarithme transforme le produit en
somme, et son image dans $(0,1]$ prouve que $\ell$ est une bijection sur
$\mathbb R_{\geq0}$. En particulier, pour la réponse d’un canal $q$,
$\ell(r)=-\log q$. Si les deux paramètres sont strictement inférieurs à
$1$, leur produit l’est aussi. Enfin, deux nombres de $(0,1]$
ont pour produit $1$ uniquement s’ils sont tous deux $1$.
\end{proof}

\begin{proposition}[Simplification admissible et prolongement inverse]
\label{amp-comp-inverse}
Pour $r,u\in D$, l’équation $r\odot t=u$ a une solution $t\in D$
exactement lorsque $u\geq r$. Cette solution est unique et donnée par
\begin{equation}
 t=f\!\left(\frac{\psi(u)}{\psi(r)}\right).
 \label{amp-comp-cancellation}
\end{equation}
La même fonction $f$ est une bijection de $(0,\infty)$ sur $(0,1)$.
Avec cette extension de $\psi$, la loi~\eqref{amp-comp-law} transporte le
groupe multiplicatif des paramètres positifs sur l’intervalle $(0,1)$.
L’inversion du groupe satisfait
\begin{equation}
 r^{\ominus}=\psi(r)r,\qquad
 r\odot r^{\ominus}=\frac34.
 \label{amp-comp-group-inverse}
\end{equation}
\end{proposition}
\begin{proof}
L’équation équivaut à
$\psi(t)=\psi(u)/\psi(r)$. Ce quotient appartient à $(0,1]$
exactement lorsque $\psi(u)\leq\psi(r)$, ce qui, par monotonie décroissante,
équivaut à $u\geq r$. La bijection prouve l’existence et l’unicité.
La dérivée de \eqref{iii:eq:monotone} est négative pour tout $s>0$ ;
les limites de $f$ en $0$ et à l’infini sont $1$ et $0$, respectivement.
Cela établit la bijection étendue. En multipliant le numérateur et le
dénominateur par $s^3$, on obtient
\[
 f(s^{-1})=\frac{s^3+s^2+s}{s^3+s^2+s+1}=s f(s).
\]
L’inverse multiplicatif de $s=\psi(r)$ a donc pour réponse
$\psi(r)r$. Sa composition avec $r$ correspond au paramètre $1$.
\end{proof}

Le prolongement à $s>1$, ou de manière équivalente à $q>1$, est l’extension
algébrique du lecteur. Le secteur physique contractant du chapitre conserve
son domaine déclaré. La coordonnée $\ell$ se prolonge en un isomorphisme
avec $(\mathbb R,+)$ et change de signe sous l’inversion du groupe.

\subsubsection{Réalisation opératorielle dans un repère de feuillets commun}

\begin{proposition}[Composition avec projecteurs marqués]
\label{amp-comp-operator}
Soient $\mathcal R=\mathcal R^*$, $\mathcal R^2=I$ et
$P_\pm=(I\pm\mathcal R)/2$. Pour $j=a,b$, soient
\[
 T_j=q_{j,+}P_++q_{j,-}P_-,\qquad
 \mathcal V_j=f(T_j^{30}),\qquad 0<q_{j,\pm}<1.
\]
La composition $T_{a\circ b}=T_aT_b$ satisfait
\begin{equation}
 \mathcal V_{a\circ b}
 =f\bigl(\psi(\mathcal V_a)\psi(\mathcal V_b)\bigr).
 \label{amp-comp-operator-law}
\end{equation}
Ses deux réponses sont $r_{a,\pm}\odot r_{b,\pm}$. Si
$T_j=\exp[-(x_jI+y_j\mathcal R)]$ avec $x_j>|y_j|$, alors
\begin{equation}
 T_aT_b=
 \exp[-((x_a+x_b)I+(y_a+y_b)\mathcal R)],
 \label{amp-comp-angular-addition}
\end{equation}
et $x_a+x_b>|y_a+y_b|$.
\end{proposition}
\begin{proof}
L’orthogonalité $P_+P_-=0$ donne
\[
 T_aT_b=q_{a,+}q_{b,+}P_++q_{a,-}q_{b,-}P_-.
\]
En particulier, $T_a$ et $T_b$ commutent. Le calcul fonctionnel dans ces
projecteurs fournit
$\psi(\mathcal V_j)=T_j^{30}$ et
$(T_aT_b)^{30}=T_a^{30}T_b^{30}$ ; appliquer $f$ prouve
\eqref{amp-comp-operator-law}. Pour toute fonction $g$ définie sur les
deux valeurs propres, on a, plus explicitement,
\[
 P_\pm g(T_j)=g(q_{j,\pm})P_\pm.
\]
Comme $q_{j,\pm}=\exp[-(x_j\pm y_j)]$, leurs produits prouvent
\eqref{amp-comp-angular-addition}. L’inégalité triangulaire conclut
$x_a+x_b>|y_a|+|y_b|\geq|y_a+y_b|$.
\end{proof}

Les coordonnées retrouvées peuvent s’écrire directement comme
\[
 x=\frac{\ell(r_+)+\ell(r_-)}2,\qquad
 y=\frac{\ell(r_+)-\ell(r_-)}2.
\]
La composition additionne ces paires. L’échange simultané de feuillets
permute les deux canaux de chaque transport, agit comme
$(x,y)\mapsto(x,-y)$ et est un automorphisme de la loi composante par
composante. L’inversion des deux transports dans l’extension algébrique
agit comme $(x,y)\mapsto(-x,-y)$. Les deux involutions commutent et
conservent des fonctions distinctes : l’une échange les marques d’orientation ;
l’autre inverse le transport.

Pour représenter l’échange par conjugaison, on utilise un
opérateur $L$ satisfaisant $L\mathcal R L^{-1}=-\mathcal R$. Dans la
représentation
\[
 \mathcal R=\begin{pmatrix}0&1\\1&0\end{pmatrix},\qquad
 L=\begin{pmatrix}1&0\\0&-1\end{pmatrix},
\]
la multiplication directe donne $LP_\pm L^{-1}=P_\mp$ et, par conséquent,
$LT_jL^{-1}=q_{j,-}P_++q_{j,+}P_-$. La conjugaison par
$\mathcal R$ conserve ses propres projecteurs. Une trace qui réunit les
canaux ne remplace pas non plus les projecteurs dans cette identité fonctionnelle :
l’information de leurs étiquettes intervient dans la composition.

L’hypothèse de repère commun est substantielle. Par exemple, les opérateurs
positifs définis
\[
 A_0=\begin{pmatrix}1/2&0\\0&1/3\end{pmatrix},\qquad
 B_0=\begin{pmatrix}1/2&1/6\\1/6&1/2\end{pmatrix}
\]
ont des valeurs propres positives, mais
\[
 A_0B_0=\begin{pmatrix}1/4&1/12\\1/18&1/6\end{pmatrix}
\]
n’est pas autoadjoint. La proposition spécifie la composition des
transports dans les projecteurs communs ; son domaine ne comprend pas un mélange
arbitraire de milieux matériels. La réalisation d’une concaténation de
chemins HMT conserve en outre ses données de mémoire : la fermeture algébrique du
cône de paramètres décrit la représentation et ne sélectionne pas à elle seule
de nouveaux chemins primaires.

\subsubsection{Produit de canaux et composition de transports}

La vitesse normalisée d’un état reste déterminée par
\[
 \widehat c=(r_+r_-)^{-1}.
\]
Ici, le produit ordinaire réunit les deux réponses de ce même état.
La loi $\odot$ évalue une autre opération : la composition de deux transports
dans chaque feuillet. Le calcul exact suivant distingue leurs arguments :
\begin{equation}
 f(1/2)=\frac{14}{15},\qquad
 \frac{14}{15}\odot\frac{14}{15}=f(1/4)=\frac{84}{85}
 \ne\frac{196}{225}=\left(\frac{14}{15}\right)^2.
 \label{amp-comp-rational-example}
\end{equation}
Les fractions de cet exemple sont des paramètres algébriques de vérification.
L’expression de $\widehat c$ et les identités constitutives précédentes
restent inchangées. Le contenu ajouté est une loi explicite pour
transporter la composition et retrouver sa coordonnée angulaire additive.


\subsection{Coordonnées logarithmiques et raffinement}

Soient $\mathcal E=\log(r_+r_-)$ et $\mathcal B=\log(r_-/r_+)$. Alors
\begin{equation}
 \log\mathcal V=\tfrac12\mathcal E I-\tfrac12\mathcal B\mathcal R.
 \label{iii:eq:logvac}
\end{equation}
Cette décomposition exhibe la composante invariante et la composante orientée.
Dans l’amplification $\mathcal V_m=\mathcal V\otimes I_{9^m}$, les
espérances partielles normalisées vérifient
$E_m(\mathcal V_{m+1})=\mathcal V_m$. Pour tout polynôme $p$,
$p(\mathcal V_m)=p(\mathcal V)\otimes I_{9^m}$ ; d’où
$E_m(p(\mathcal V_{m+1}))=p(\mathcal V_m)$. La continuité uniforme
étend l’identité aux fonctions continues du spectre. Les projecteurs,
les puissances et le logarithme sont ainsi conservés dans la tour indiquée.

\subsection{Types dimensionnels des quatre évaluations}

Avec les droites d’action, de charge et de vitesse, munies de sections
$u_S,u_Q,u_C$, la réalisation dimensionnelle s’écrit
\begin{align}
 Z_{\rm vac}&=\widehat Z u_Su_Q^{-2},&
 c_{\rm vac}&=\widehat c u_C,\\
 \mu_{\rm vac}&=\widehat\mu u_Su_C^{-1}u_Q^{-2},&
 \varepsilon_{\rm vac}&=\widehat\varepsilon u_S^{-1}u_C^{-1}u_Q^2.
 \label{iii:eq:dimensions}
\end{align}
En multipliant et en divisant ces expressions, on préserve
$\mu_{\rm vac}\varepsilon_{\rm vac}=c_{\rm vac}^{-2}$ et
$\mu_{\rm vac}/\varepsilon_{\rm vac}=Z_{\rm vac}^2$.
La coordinatisation SI évalue ensuite les sections. La comparaison physique nécessite
d’exposer leur construction et leur normalisation avec l’évaluation ; les quatre
coordonnées normalisées de \eqref{iii:eq:four} conservent leur type adimensionnel.

% INSERCIÓN FOCAL AUTORIZADA: no introduce un capítulo ni el concepto de radión.
% Incluir dentro de la sección constitutiva, después de la inversión cúbica
% y de la declaración de las cartas dimensionales. Etiquetas fuera de grupos
% de prefijado automático. Conserva íntegramente el bloque geométrico anterior.
\subsection{Reconstruction de la forme d’action à partir de la réponse constitutive}
\label{iii:subsec:puente-forma-LC}

La réponse constitutive et la réalisation inductive--capacitive reçoivent
le même couple angulaire construit à partir d’APP--TRIT--TPK. La première évalue
des fonctions rationnelles des deux canaux ; la seconde utilise le rapport de
leurs exposants pour déterminer la forme de l’ellipse d’action.
L’inversion cubique permet de composer les deux lectures sans identifier leurs
quotients numériques.

Soient $x=A_{\rm rad}$ et $y=C^*_{\rm rad}$, avec $x>|y|$, les coordonnées
relevées déjà construites, et conservons leurs étiquettes de feuillet. Nous écrivons
\begin{equation}
 q_\pm=e^{-(x\pm y)},\qquad s_\pm=q_\pm^{30},\qquad
 r_\pm=f(s_\pm),\qquad
 f(s)=\frac{1+s+s^2}{1+s+s^2+s^3}.
 \label{iii:eq:forma-canales}
\end{equation}
Le domaine contractant donne $s_\pm\in(0,1)$ et $r_\pm\in(3/4,1)$. Dans la chambre $x>y>0$, on a $s_+<s_-$ et $r_+>r_-$ ; l'échange des feuillets transporte cette chambre vers l'orientation opposée.

L’image des deux coordonnées constitutives normalisées est le domaine
\begin{equation}
 \mathcal U=\left\{(z,v)\in(0,\infty)^2:
   (zv)^{-1/2}\in(3/4,1),\quad
   (z/v)^{1/2}\in(3/4,1)\right\}.
 \label{iii:eq:forma-dominio}
\end{equation}
Ici, $(z,v)=(\widehat Z,\widehat c)$. Les restrictions expriment l'appartenance des canaux reconstruits au domaine du théorème d'inversion ; elles ne constituent pas un choix ultérieur de leurs valeurs.

\Needspace{23\baselineskip}
\begin{proposition}[Composition constitutive de la forme et de l’impédance LC]
\label{iii:prop:forma-LC}
Pour $(\widehat Z,\widehat c)\in\mathcal U$, soient $s_\pm\in(0,1)$ les
racines uniques de
\begin{equation}
 r_\pm s_\pm^3+(r_\pm-1)(1+s_\pm+s_\pm^2)=0,
 \qquad
 r_+=\frac1{\sqrt{\widehat Z\widehat c}},\quad
 r_-=\sqrt{\frac{\widehat Z}{\widehat c}}.
 \label{iii:eq:forma-inversion}
\end{equation}
Le rapport angulaire, l’anisotropie de l’ellipse d’action et l’impédance
LC relative sont reconstruits exactement par
\begin{align}
 \frac{y}{x}
 &=\frac{\log s_+-\log s_-}{\log s_++\log s_-},
 \label{iii:eq:forma-razon}\\
 \eta_{\rm el}
 &=\frac12\log\frac{\log s_+}{\log s_-},
 \label{iii:eq:forma-eta}\\
 \frac{Z_{\rm LC}}{Z_*}
 &=e^{-\eta_{\rm el}}
  =\sqrt{\frac{\log s_-}{\log s_+}}.
 \label{iii:eq:forma-LC}
\end{align}
Dans la dernière identité, $Z_{\rm LC}=Z_{\eta_{\rm el}}$ est l’impédance
de la coordinatisation LC de \eqref{eq:ii-elipse-lc}.
\end{proposition}
\begin{proof}
L'inversion cubique de \eqref{iii:eq:cubic} détermine $s_\pm$ de manière univoque, car $f$ est une bijection décroissante de $(0,1)$ sur $(3/4,1)$. Ses racines positives d'ordre $30$ récupèrent $q_\pm$. Les identités
\[
 \log s_+=-30(x+y),\qquad \log s_-=-30(x-y)
\]
prouvent \eqref{iii:eq:forma-razon} par différence et somme. Les deux logarithmes
sont négatifs ; leur quotient est positif et leur somme n’est pas nulle. Par conséquent,
\[
 \frac12\log\frac{\log s_+}{\log s_-}
   =\frac12\log\frac{x+y}{x-y}
   =\operatorname{artanh}(y/x)=\eta_{\rm el}.
\]
La relation $Z_{\eta_{\rm el}}/Z_*=e^{-\eta_{\rm el}}$, démontrée dans
\eqref{eq:ii-elipse-impedancia}, donne \eqref{iii:eq:forma-LC}.
\end{proof}

La composition conserve l’information nécessaire à chaque reconstruction.
Avec $\hbar$, on retrouve
$a_S=\sqrt{2\hbar}\,e^{\eta_{\rm el}/2}$ et
$b_S=\sqrt{2\hbar}\,e^{-\eta_{\rm el}/2}$ ; avec la coordinatisation
$(\omega,Z_*)$, on retrouve l’inductance et la capacité déjà définies.
L’action, l’horloge et la référence dimensionnelle accompagnent la forme comme
données typées. L’inversion de deux coordonnées adimensionnelles ne
reconstruit ni leurs unités ni toute la mémoire de l’état.

La normalisation intervient avant de résoudre les équations cubiques.
Si l’on reçoit $Z_{\rm vac}$ et $c_{\rm vac}$ dans une coordinatisation dimensionnelle,
les relations \eqref{iii:eq:dimensions} déterminent d’abord
\[
 \widehat Z=\frac{Z_{\rm vac}}{u_Su_Q^{-2}},\qquad
 \widehat c=\frac{c_{\rm vac}}{u_C}.
\]
Ce sont les coordonnées auxquelles s'applique \eqref{iii:eq:forma-inversion}. La référence $Z_*$ appartient à la réalisation LC et est expressément conservée dans \eqref{iii:eq:forma-LC}.

Les deux quotients sont reliés au moyen des canaux, mais évaluent
des fonctions différentes :
\begin{equation}
 \widehat Z=\frac{f(s_-)}{f(s_+)},\qquad
 \frac{Z_{\rm LC}}{Z_*}
       =\sqrt{\frac{\log s_-}{\log s_+}}.
 \label{iii:eq:forma-dos-cocientes}
\end{equation}
La reconstruction précédente fournit l'application entre les deux lectures. En particulier, échanger les feuillets inverse les deux quotients, remplace $\eta_{\rm el}$ par $-\eta_{\rm el}$ et conserve $\widehat c$. Dans la limite symétrique $y=0$, les canaux coïncident et les deux quotients valent $1$. En dehors de ce cas, l'égalité d'origine n'impose pas l'égalité des deux fonctions. La forme d'action est récupérable à partir de la réponse constitutive en conservant l'orientation et la normalisation.

% Recuperación autocontenida: sector angular de árbol y reconstrucción Higgs--vacío.
% Propietarios y alcance: higgs_procedencia.md. Incorporar después del operador constitutivo.
\section{Reconstruction conjointe de la jauge, de Higgs et de la réponse constitutive}
\label{amp-higgs-seccion}

La réponse du vide et le secteur électrofaible admettent des lectures complémentaires
du transport angulaire déjà construit. La première conserve les deux canaux
marqués ; le déterminant conserve leur contraction commune ; la composante angulaire
de Higgs conserve en outre l’orientation qui disparaît dans ce déterminant.
La composition de ces lectures permet de reconstruire leur antécédent spectral
et d’exprimer l’impédance au moyen de l’autocouplage scalaire et des couplages
de jauge, dans une même réalisation à l’arbre.

\subsection{Provenance angulaire et signatures des chemins}

On conserve la chaîne APP--TRIT--TPK : le prolongement survivant détermine
le chemin, celui-ci produit son registre orienté et le lecteur du registre exprime
la signature angulaire. Les coordonnées $\pi=\pi_{\rm HMT}$ et
$\alpha=\alpha_{\rm HMT}$ sont des sorties préalables de la même génération ;
$A$ et $C^*$ sont leurs représentations angulaires ultérieures. Dans ce bloc, on
utilise des radians,
\begin{equation}
 x=\frac{\pi A_{\deg}}{180},\qquad
 y=\frac{\pi C^*_{\deg}}{180},\qquad
 A_{\deg}=1000\alpha,\qquad 0<y<x.
 \label{amp-higgs-angulos}
\end{equation}
Par conséquent, $x=50\pi\alpha/9$. Les canaux et les marques de feuillet sont ceux
de la section~\ref{iii:sec:hojas} :
\begin{equation}
 \begin{aligned}
 q_+&=e^{-x-y},& q_-&=e^{-x+y},& D&=q_+q_-=e^{-2x},\\
 P_\pm&=\frac{I\pm\mathcal R}{2},&
 T&=q_+P_++q_-P_-.
 \end{aligned}
 \label{amp-higgs-transporte}
\end{equation}
Ici, $\mathcal R$ est l’involution autoadjointe marquée et chaque projecteur
est de rang un dans la réalisation bicouche de base.

Le lecteur angulaire de durée--perdurance est
\begin{equation}
 \chi_{\rm dur}(n,s)=\exp\!\left(nx+\frac{s}{6}y\right).
 \label{amp-higgs-caracter}
\end{equation}
Pour conserver la provenance finie des coefficients qui seront utilisés,
nous réunissons les empreintes nominales des trois chemins du calendrier de longueur
$108$. Les colonnes de déplacements cardinaux se lisent comme $(N,E,S,O)$ ;
la paire suivante conserve le feuillet et le résidu exprimés par ce chemin.
Le tableau conserve, pour chaque chemin identifié, l’empreinte et la signature
angulaire exprimées par la chaîne précédente.
\begin{center}
\small
\begin{tabular}{c c c c c}
\hline
Secteur & Chemin et axe & $(N,E,S,O)$ & Feuillet/résidu & $(n,s)$\\
\hline
$W$ & $15/37$, E--O & $(23,44,11,30)$ & $(41,8)$ & $(94,-1)$\\
$Z$ & $20/23$, E--O & $(33,45,10,20)$ & $(57,7)$ & $(95,-1)$\\
$H$ & $24/29$, N--S & $(40,34,22,12)$ & $(47,10)$ & $(97,9)$\\
\hline
\end{tabular}
\end{center}
Les trois empreintes typées conservent l’identifiant de chemin, l’orientation
et le registre de provenance des signatures employées. L’évaluation angulaire
commence après l’extraction à partir du registre enrichi ; le tableau
réunit ses sorties pour composer les caractères dans cette réalisation.
Les marques $W_\pm$ qui suivent désignent des coordonnées d’enroulement,
distinctes du nom de l’espèce $W$ :
\begin{equation}
 \begin{aligned}
 W_\pm&=6n\pm s, &
 n&=\frac{W_++W_-}{12},& s&=\frac{W_+-W_-}{2},\\
 (n,s)_W&=(94,-1), & (W_+,W_-)_W&=(563,565),\\
 (n,s)_Z&=(95,-1), & (W_+,W_-)_Z&=(569,571),\\
 (n,s)_H&=(97,9), & (W_+,W_-)_H&=(591,573).
 \end{aligned}
 \label{amp-higgs-firmas}
\end{equation}
La matrice de la première application est
$\bigl(\begin{smallmatrix}6&1\\6&-1\end{smallmatrix}\bigr)$, de déterminant
$-12$ et d’image $\{(u,v)\in\mathbb Z^2:u+v\in12\mathbb Z\}$.
En effet, les formules inverses précédentes sont entières précisément sur
cette image. Le caractère conserve les deux orientations au moyen de
\begin{equation}
 \chi_{\rm dur}(n,s)
 =\left(q_+^{-W_+}q_-^{-W_-}\right)^{1/12}.
 \label{amp-higgs-caracter-canales}
\end{equation}
La preuve s’obtient en prenant les logarithmes : le numérateur est
$W_+(x+y)+W_-(x-y)=12nx+2sy$.
Ainsi sont identifiés la racine douzième et le diviseur $6$ du lecteur.

\begin{proposition}[Quotients angulaires déterminés par les signatures]
Dans une section d’échelle commune, avec les facteurs de raffinement partagés
qui permettent leur annulation dans les quotients, on a
\begin{equation}
 \frac{m_W^{\angle}}{m_Z^{\angle}}=e^{-x},\qquad
 \frac{m_H^{\angle}}{m_W^{\angle}}
 =e^{3x+5y/3},\qquad
 \left(\frac{m_H^{\angle}}{m_W^{\angle}}\right)^2
 =e^{6x+10y/3}.
 \label{amp-higgs-cocientes}
\end{equation}
\end{proposition}
\begin{proof}
Diviser les caractères de $W$ et $Z$ soustrait leurs signatures : $(94,-1)-(95,-1)=(-1,0)$.
Pour $H/W$, la différence est $(97,9)-(94,-1)=(3,10)$.
L’équation~\eqref{amp-higgs-caracter} donne $3x+(10/6)y$ ; le carré
multiplie cet exposant par deux. L’échelle commune et les raffinements
partagés s’annulent avant ces opérations.
\end{proof}
Ces quotients conservent l’identifiant de la coordinatisation angulaire. L’application
à des masses ayant des raffinements sectoriels distincts conserve leurs facteurs
supplémentaires ; la condition d’annulation fait partie de l’énoncé.

\subsection{Réalisation rationalisée à l’arbre}

La lecture électrofaible de cette coordinatisation définit
$c_W^2=D$, $s_W^2=1-D$ et $e_g^2=4\pi\alpha$. La racine positive fixe
$g=e_g/s_W$ et $g'=e_g/c_W$. On obtient
\begin{equation}
 g^2=\frac{4\pi\alpha}{1-D},\qquad
 g'^2=\frac{4\pi\alpha}{D}.
 \label{amp-higgs-calibre}
\end{equation}
La quantité $e_g$ est un couplage adimensionnel dans la normalisation
rationalisée, distinct d’une section dimensionnelle de charge.
La convention du potentiel scalaire est
\begin{equation}
 \begin{gathered}
 V(H)=-\mu_H^2H^\dagger H+\lambda_H(H^\dagger H)^2,\\
 m_H^2=2\lambda_Hv^2,\qquad m_W^2=\frac{g^2v^2}{4}.
 \end{gathered}
 \label{amp-higgs-potencial}
\end{equation}
La même $v$ intervient dans les deux relations. En les divisant, son échelle
s’annule et $m_H^2/m_W^2=8\lambda_H/g^2$. Avec
\eqref{amp-higgs-cocientes}, on obtient
\begin{equation}
 \boxed{\lambda_H=\frac{g^2}{8}\exp\!\left(6x+\frac{10}{3}y\right).}
 \label{amp-higgs-lambda}
\end{equation}
Le facteur $8$ provient des deux normalisations quadratiques de
\eqref{amp-higgs-potencial} ; les coefficients $6$ et $10/3$ proviennent des
signatures. La distinction conserve l’opération qui produit chaque facteur.
La composition reçoit les sorties HMT préalables et utilise cette réalisation physique
à l’arbre à échelle commune ; un changement de normalisation du champ, d’échelle ou
de schéma se transporte également dans ses relations de masses et de couplages.

\subsection{Inverse conjoint et domaine exact}

\begin{theorem}[Reconstruction angulaire au moyen de la jauge et de Higgs]
\label{amp-higgs-inversa-teorema}
Pour $\pi>0$ fixé, l’application de
\eqref{amp-higgs-transporte}, \eqref{amp-higgs-calibre} et
\eqref{amp-higgs-lambda}, sur $\alpha>0$, $x>y>0$, est injective et a
pour inverse
\begin{equation}
 \begin{aligned}
 D&=\frac{g^2}{g^2+g'^2},&
 x&=-\frac12\log D,&
 \alpha&=\frac{g^2g'^2}{4\pi(g^2+g'^2)},\\
 y&=\frac3{10}\left[\log\!\left(\frac{8\lambda_H}{g^2}\right)-6x\right].
 \end{aligned}
 \label{amp-higgs-inversa}
\end{equation}
Son image exacte est formée des triplets $g,g',\lambda_H>0$ tels que,
avec $D$ et $x$ reconstruits comme ci-dessus,
\begin{equation}
 \frac{g^2}{8}e^{6x}<\lambda_H<\frac{g^2}{8}e^{28x/3}.
 \label{amp-higgs-dominio}
\end{equation}
\end{theorem}
\begin{proof}
Les deux équations de jauge donnent
$g^2/(g^2+g'^2)=D$ et
$g^2g'^2/(g^2+g'^2)=4\pi\alpha$. Le logarithme réel retrouve $x$ et
l’équation scalaire retrouve $y$. Les conditions $g,g'>0$ impliquent
$0<D<1$ et $x>0$. L’inégalité~\eqref{amp-higgs-dominio} équivaut,
par le logarithme strictement croissant, à
$0<\log(8\lambda_H/g^2)-6x<10x/3$, c’est-à-dire $0<y<x$.
Réciproquement, les formules reconstruites produisent
$4\pi\alpha/(1-D)=g^2$, $4\pi\alpha/D=g'^2$ et
$g^2e^{6x+10y/3}/8=\lambda_H$. Les racines positives retrouvent les
couplages originaux. Cela prouve la nécessité, la suffisance et l’unicité.
\end{proof}

L’inversion précédente décrit exactement cette application entre lecteurs.
La section engendrée par HMT conserve en outre la relation de
\eqref{amp-higgs-angulos}, qui s’exprime comme
\begin{equation}
 -\frac12\log\!\left(\frac{g^2}{g^2+g'^2}\right)
 =\frac{25}{18}\frac{g^2g'^2}{g^2+g'^2}.
 \label{amp-higgs-compatibilidad-alpha}
\end{equation}
Elle s’obtient en substituant l’inverse de $\alpha$ dans $x=50\pi\alpha/9$.
Le contre-angle conserve également son équation de retour antérieure.
Par conséquent, l’image du lecteur et la section produite par le générateur
sont des objets distincts : les compatibilités préalables sélectionnent la section
au sein du domaine reconstructible. L’opération inverse agit après
la génération et permet de vérifier ses sorties conservées.

\subsection{Reconstruction Higgs--impédance}

Pour exprimer le contenu orienté sans altérer les canaux,
nous définissons
\begin{equation}
 \rho_{\rm or}=\frac{q_-}{q_+}=e^{2y},\qquad
 \Xi_H=\frac{8D^3\lambda_H}{g^2}.
 \label{amp-higgs-Xi}
\end{equation}
L’équation~\eqref{amp-higgs-lambda} et $D^3=e^{-6x}$ donnent
\begin{equation}
 \begin{aligned}
 \Xi_H&=e^{10y/3}=\rho_{\rm or}^{5/3},&
 \rho_{\rm or}&=\Xi_H^{3/5},\\
 q_-&=\sqrt D\,\Xi_H^{3/10},&
 q_+&=\sqrt D\,\Xi_H^{-3/10}.
 \end{aligned}
 \label{amp-higgs-canales}
\end{equation}
Les puissances sont les racines réelles positives. Le domaine exact
de la chambre orientée est
\begin{equation}
 0<D<1,\qquad 1<\Xi_H<D^{-5/3}.
 \label{amp-higgs-dominio-Xi}
\end{equation}
En effet, $0<y<x$ équivaut à $1<e^{10y/3}<e^{10x/3}$.
L’extrémité inférieure donne $q_+=q_-$ ; la supérieure donne $q_-=1$.
Pour les deux orientations ensemble, $|y|<x$ équivaut à
$D^{5/3}<\Xi_H<D^{-5/3}$.

Soit $F(q)=f(q^{30})=(1-q^{90})/(1-q^{120})$, la réponse déjà prouvée dans
\eqref{iii:eq:rchannels}--\eqref{iii:eq:monotone}.
\begin{corollary}[Confluence des lectures scalaire et constitutive]
\label{amp-higgs-impedancia-corolario}
Dans le domaine~\eqref{amp-higgs-dominio-Xi},
\begin{equation}
 \boxed{\widehat Z=
 \frac{F(\sqrt D\,\Xi_H^{3/10})}
 {F(\sqrt D\,\Xi_H^{-3/10})}.}
 \label{amp-higgs-impedancia}
\end{equation}
Les réponses complètes sont
\begin{equation}
 \begin{aligned}
 r_-&=F(\sqrt D\,\Xi_H^{3/10}),&
 r_+&=F(\sqrt D\,\Xi_H^{-3/10}),\\
 \widehat\mu&=r_-^2,& \widehat\varepsilon&=r_+^2,&
 \widehat c&=(r_-r_+)^{-1}.
 \end{aligned}
 \label{amp-higgs-vacio}
\end{equation}
De plus, les deux reconstructions satisfont
\begin{equation}
 \frac{F^{-1}(\sqrt{\widehat\mu})}
 {F^{-1}(\sqrt{\widehat\varepsilon})}
 =\rho_{\rm or}=\Xi_H^{3/5}.
 \label{amp-higgs-confluencia}
\end{equation}
\end{corollary}
\begin{proof}
Substituer les canaux de~\eqref{amp-higgs-canales} dans les quatre lectures
de~\eqref{iii:eq:four} prouve les premières identités. L’inverse
constitutif de~\eqref{iii:eq:cubic} retrouve $q_-$ et $q_+$ ; leur quotient
est $\rho_{\rm or}$. L’égalité avec $\Xi_H^{3/5}$ est déjà prouvée dans
\eqref{amp-higgs-canales}.
\end{proof}

Dans la section compatible avec $A_{\deg}=1000\alpha$, le couplage
$Q(D):=-9\log D/25$ satisfait $Q(D)=4\pi\alpha$ et
$g^2=Q(D)/(1-D)$. Par conséquent, la même coordonnée de Higgs s’écrit
\begin{equation}
 \Xi_H=\frac{8(1-D)D^3\lambda_H}{Q(D)},\qquad
 \frac1{g^2}+\frac1{g'^2}=\frac1{Q(D)}.
 \label{amp-higgs-curva}
\end{equation}
Ces expressions relient le couplage total, sa répartition rationalisée
et l’orientation scalaire dans la même section.

\subsection{Information conservée et transport d’orientation}

La paire de jauge retrouve $\alpha$ et $x$ ; le quotient
$\lambda_H/g^2$ retrouve ensuite $y$. Avec les marques $P_\pm$, on reconstruit
les opérateurs complets $T$ et $\mathcal V=r_+P_++r_-P_-$.
Les trois scalaires déterminent les valeurs propres ; les projecteurs marqués
déterminent en outre leur réalisation opératorielle. L’histoire enrichie du
chemin reste une information antérieure de résolution supérieure.

Sous $y\mapsto-y$ par rapport aux marques fixes, $D,g,g'$ restent
invariants et $\Xi_H\mapsto\Xi_H^{-1}$. Les réponses constitutives s’échangent,
$\widehat Z\mapsto\widehat Z^{-1}$ et $\widehat c$ est conservé.
Si $\lambda_H^+$ et $\lambda_H^-$ désignent les évaluations dans les deux
orientations conjuguées, on a
\begin{equation}
 \lambda_H^+\lambda_H^-
 =\frac{g^4}{64}D^{-6}.
 \label{amp-higgs-orientacion}
\end{equation}
La preuve consiste à multiplier~\eqref{amp-higgs-lambda} pour $y$ et $-y$ ;
les contributions orientées s’annulent. Cette comparaison conserve le
lecteur de signatures et change son argument angulaire. Une conjugaison simultanée
des chemins et des références doit également transporter les signatures.

Dans l’amplification $T_m=T\otimes I_{9^m}$, avec
$\tau_m=\operatorname{Tr}/(2\cdot9^m)$ et
$\mathcal R_m=\mathcal R\otimes I_{9^m}$, les coordonnées appropriées sont
\begin{equation}
 D=\exp\bigl(2\tau_m(\log T_m)\bigr),\qquad
 \rho_{\rm or}=\exp\bigl(-2\tau_m(\mathcal R_m\log T_m)\bigr).
 \label{amp-higgs-refinamiento}
\end{equation}
En effet, $\log T_m=(\log T)\otimes I_{9^m}$, de sorte que les deux traces
sont $-x$ et $-y$, respectivement. La normalisation élimine la multiplicité
de l’amplification. Les lectures de jauge, de Higgs et du vide restent ainsi
compatibles avec ces raffinements ; le déterminant ordinaire amplifié
serait $D^{9^m}$ et a une fonction distincte.

La conclusion de la composition est précise : l’invariant pair de jauge
et la coordonnée orientée scalaire reconstruisent conjointement la paire
spectrale qui gouverne le vide. Lorsque $\lambda_H$ est évalué au moyen de la
même formule, la confluence est une identité de cohérence. Une lecture
indépendante de $\lambda_H$, transportée vers des schéma et échelle identiques,
transforme cette égalité en comparaison supplémentaire. Dans les deux usages restent
explicites les signatures, les marques de feuillet et la normalisation à l’arbre.

\section{Décomposition cyclique, quart de tour et moment orienté de Catalan}
\label{iii:sec:catalan}

Le registre dodécaphasique permet de relier la réponse constitutive à
d’autres invariants du transport. Le lien s’établit sur des sous-espaces
et des applications déterminés : le défaut de rang mesure une différence
d’incidence, le déterminant conserve une réponse scalaire et un coefficient
de la résolvante retient l’orientation. La famille de moments qui culmine
en Catalan intervient à ce dernier niveau.

\subsection{Deux sous-espaces du cycle dodécaphasique}

Dans $\mathbb R^{12}$, soit $C e_j=e_{j+1\bmod12}$ et soient
\begin{equation}
 P_3^{\mathrm{cyc}}=\frac{I+C^3+C^6+C^9}{4},\qquad
 P_4^{\mathrm{cyc}}=\frac{I+C^4+C^8}{3},\qquad
 H_d=\operatorname{Ran}P_d^{\mathrm{cyc}}.
 \label{iii:eq:projectors34}
\end{equation}
Les moyennes sur les sous-groupes cycliques sont des projecteurs orthogonaux.
$H_3$ est formé des suites périodiques de période divisant $3$,
et $H_4$ de celles de période divisant $4$. Leurs dimensions sont $3$ et $4$.
La restriction de $C$ à chaque espace réalise le cycle correspondant.
L’exposant $\mathrm{cyc}$ identifie la moyenne du cycle dodécaphasique :
$P_3^{\mathrm{cyc}}$ est le projecteur sur les suites de période divisant
$3$. Son domaine et son action sont distincts de ceux du projecteur isotypique
tridimensionnel de la réalisation icosaédrique ; la coïncidence des rangs
n’identifie pas les deux opérateurs.

\begin{proposition}[Défaut de rang et réponse déterminantale]
On vérifie
\begin{equation}
 \frac{\operatorname{Tr}(P_3^{\mathrm{cyc}}-P_4^{\mathrm{cyc}})}{12}=-\frac1{12},\qquad
 \frac{\det_{H_3}(I-sC)}{\det_{H_4}(I-sC)}
 =\frac{1-s^3}{1-s^4}=f(s),\quad 0<s<1.
 \label{iii:eq:trace-det}
\end{equation}
\end{proposition}
\begin{proof}
La trace d’un projecteur orthogonal est son rang, de sorte que le numérateur
de la première expression est $3-4$. Le polynôme caractéristique d’un cycle
de longueur $d$ est $t^d-1$ ; son déterminant $\det(I-sC)$ est donc
$1-s^d$. Le quotient produit la seconde identité.
\end{proof}

La substitution $s=q^{30}$ transporte cette représentation vers le quotient
$90/120$. En $s=1$, le mode constant commun s’annule avant de prendre la
limite, qui vaut $3/4$. Cette réalisation de $-1/12$ comme défaut de trace
a un domaine fini explicite. Les réalisations régularisées et
modulaires du même scalaire conservent leurs opérateurs et domaines propres.

\subsection{Représentation fidèle du quart de tour}

Définissons le projecteur antipodal $P_{\rm ant}=(I-C^6)/2$ et
$Q=P_4^{\mathrm{cyc}}P_{\rm ant}$. Ce projecteur agit sur le cycle à douze phases et est
distinct du projecteur constitutif $P_-$ de la section~\ref{iii:sec:constitutiva}.
Les projecteurs $P_4^{\mathrm{cyc}}$ et $P_{\rm ant}$ commutent et
$Q$ est de rang $2$. Dans son image, on fixe la base orthonormée
\begin{equation}
 u=\frac1{\sqrt6}(1,0,-1,0,1,0,-1,0,1,0,-1,0)^{\mathsf T},
 \qquad v=Cu.
 \label{iii:eq:uv}
\end{equation}
On vérifie composante par composante que $Cu=v$ et $Cv=-u$. Par conséquent,
l’isométrie $W(a,b)=au+bv$ satisfait
\begin{equation}
 CW=WJ,\qquad
 J=\begin{pmatrix}0&-1\\1&0\end{pmatrix},\qquad J^2=-I.
 \label{iii:eq:quarter-intertwiner}
\end{equation}
La structure complexe du plan est réalisée par l’action effective du
cycle. Le signe dépend de l’orientation fixée : dans la même base,
$C^3W=-WJ$. Cette précision permet de composer le plan avec le lecteur orienté.

\subsection{Résolvante et profondeur illimitée}

Comme $C^2=-I$ sur $\operatorname{Ran}Q$, pour $0\leq s<1$, on a
\begin{equation}
 (I-sC)^{-1}\big|_{\operatorname{Ran}Q}=\frac{I+sC}{1+s^2},
 \qquad k_4(s):=\langle v,(I-sC)^{-1}u\rangle=\frac{s}{1+s^2}.
 \label{iii:eq:resolvent}
\end{equation}
La suite $\langle v,C^nu\rangle$ est $0,1,0,-1,\ldots$. Le lecteur
de profondeur conserve ce caractère au moyen du moment
\begin{equation}
 \mathcal C_2(s)=\sum_{k\geq0}\frac{(-1)^ks^{2k+1}}{(2k+1)^2},
 \qquad 0\leq s\leq1.
 \label{iii:eq:c2}
\end{equation}
La série converge absolument et uniformément sur cet intervalle par
comparaison avec $\sum(2k+1)^{-2}$. Pour $0<s<1$, les séries dérivées
convergent uniformément sur chaque intervalle compact intérieur. Avec
$\mathcal D=s\,d/ds$, on obtient
\begin{equation}
 \mathcal D^2\mathcal C_2(s)=k_4(s),\qquad
 \mathcal C_2(s)=\int_0^s\frac{dt}{t}\int_0^t\frac{d\xi}{\xi}\,k_4(\xi).
 \label{iii:eq:c2-integral}
\end{equation}
Les conditions $\mathcal C_2(0)=0$ et
$\lim_{s\downarrow0}\mathcal D\mathcal C_2(s)=0$ fixent les deux constantes
d’intégration. En effet, la différence de deux solutions de
$\mathcal D^2F=k_4$ a la forme $a\log s+b$ ; la seconde condition
annule $a$ et la première annule $b$.
La convergence uniforme permet d’évaluer l’extrémité $s=1$ : le résultat est
le moment de Catalan. La représentation finie du caractère conserve son
orientation ; l’indice de profondeur reste illimité.

\subsection{Composition avec la réponse constitutive}

\begin{theorem}[Relation différentielle constitutive]
Pour $0<s<1$, avec l’orientation de \eqref{iii:eq:uv},
\begin{equation}
 f(s)=\frac{1+s+s^2}{s(1+s)}\,\mathcal D^2\mathcal C_2(s).
 \label{iii:eq:catalan-vacuum}
\end{equation}
\end{theorem}
\begin{proof}
Substituer \eqref{iii:eq:resolvent} dans \eqref{iii:eq:c2-integral} et utiliser
$(1+s)(1+s^2)=1+s+s^2+s^3$ donne l’identité. Tous les dénominateurs sont
positifs sur le domaine. L’entrelaceur \eqref{iii:eq:quarter-intertwiner}
identifie le coefficient opératoriel utilisé avant son évaluation.
\end{proof}

En évaluant $s_\pm=q_\pm^{30}$, la même famille produit les deux réponses
de canal de \eqref{iii:eq:rchannels}. La limite de Catalan, le défaut de
trace et les valeurs propres constitutives maintiennent leurs fonctions différentes
au sein de cette composition. L’inversion du cycle avec les marques $u,v$
fixes change le signe de $k_4$ et conserve le déterminant $1+s^2$.
On identifie ainsi une information d’orientation que le déterminant
isolé ne distingue pas.

% Delta de §10.27. Insertar al final de 05_ciclo_catalan.tex.
% La actualización afín y su balance permanecen en sus propietarios originales.
\subsection{Traces de retour et moments du quotient constitutif}
\label{amp-afin-sec}

La mise à jour du registre et la réponse du cycle partagent le
transport dodécaphasique et conservent des informations de types différents.
L’événement signé s’obtient à partir de la direction et de l’orientation de chaque
fenêtre de l’histoire APP--TRIT--TPK. Son accroissement vectoriel et le
bilan~\eqref{mem:balance} retiennent l’interaction avec le registre préalable.
La construction suivante prolonge la lecture spectrale du même cycle,
une fois ses sous-espaces d’incidence fixés. Cet ordre permet de composer
le retour, le déterminant et les moments en conservant l’histoire qui les précède.

Nous conservons l’opérateur $C$ et les projecteurs de
\eqref{iii:eq:projectors34}, et définissons
\begin{equation}
 E_{\rm cyc}=P_3^{\mathrm{cyc}}-P_4^{\mathrm{cyc}},\qquad
 b_n=\operatorname{Tr}(C^n E_{\rm cyc}),\qquad n\geq0.
 \label{amp-afin-trazas}
\end{equation}
Le symbole $b_n$ distingue cette trace spectrale de l’événement signé
$a_m$ d’une fenêtre. La trace utilise l’opérateur de transport ;
l’événement utilise en outre la trajectoire effectivement parcourue.

\begin{proposition}[Tableau des traces produit par l’incidence]
\label{amp-afin-prop-tabla}
On a
\begin{equation}
 \begin{gathered}
 b_n=3\mathbf1_{3\mid n}-4\mathbf1_{4\mid n},\\
 (b_1,\ldots,b_{12})=(0,0,3,-4,0,3,0,-4,3,0,0,-1).
 \end{gathered}
 \label{amp-afin-tabla}
\end{equation}
La période minimale est douze et la somme sur une période est nulle. Pour
$B_N=\sum_{n=1}^{N}b_n$,
\begin{equation}
 B_N=3\lfloor N/3\rfloor-4\lfloor N/4\rfloor,\qquad |B_N|\leq3.
 \label{amp-afin-parciales}
\end{equation}
\end{proposition}
\begin{proof}
Dans $H_d$, l’opérateur $C$ permute cycliquement les $d$ classes de
coordonnées. Sa puissance $n$ a $d$ points fixes si $d$ divise $n$
et aucun dans le cas contraire. Comme $P_d^{\mathrm{cyc}}$ commute avec $C$,
$\operatorname{Tr}(C^nP_d^{\mathrm{cyc}})$ est la trace de cette restriction.
Soustraire les cas $d=3,4$ prouve la formule. La période positive minimale
est douze : $b_n=-1$ se produit exactement aux multiples de douze.
Sur une période, les contributions sont $3\cdot4-4\cdot3=0$.
La somme finie compte les multiples et donne~\eqref{amp-afin-parciales}.
Ses valeurs pour $N=0,\ldots,11$ sont
$(0,0,0,3,-1,-1,2,2,-2,1,1,1)$ et se répètent lorsqu’on ajoute douze.
\end{proof}

\begin{theorem}[Reconstruction logarithmique du lecteur constitutif]
\label{amp-afin-thm-logdet}
Pour $|s|<1$, le quotient $f$ de~\eqref{iii:eq:trace-det} vérifie
\begin{align}
 -\log f(s)&=\sum_{n\geq1}\frac{b_n}{n}s^n,\label{amp-afin-logdet}\\
 -\frac{s f'(s)}{f(s)}
 &=\frac{3s^3}{1-s^3}-\frac{4s^4}{1-s^4}.
 \label{amp-afin-logder}
\end{align}
La branche holomorphe est fixée par $\log f(0)=0$.
Sur $0<s<1$, elle coïncide avec le logarithme réel.
\end{theorem}
\begin{proof}
Les polynômes $1-s^3$ et $1-s^4$ ne s’annulent pas dans le disque ouvert.
Le développement $-\log(1-w)=\sum_{k\geq1}w^k/k$, obtenu en intégrant
la série géométrique depuis zéro, converge uniformément sur les disques
compacts. Avec $w=s^3,s^4$, leur différence donne
$-\log f(s)=\sum_{k\geq1}s^{3k}/k-\sum_{k\geq1}s^{4k}/k$.
Le coefficient de $s^n$ est $b_n/n$.
La convergence locale uniforme permet de dériver et de sommer les deux séries
géométriques, ce qui prouve~\eqref{amp-afin-logder}.
La condition en zéro fixe la constante d’intégration et la branche.
\end{proof}

Ainsi, $b_0/12=-1/12$ enregistre le défaut normalisé de dimension,
tandis que toute la suite $b_n$ reconstruit une fonction.
Ce sont deux lectures de portées différentes de la même paire de sous-espaces.
Lorsqu’on substitue $s=q_\pm^{30}$, les arguments restent les sorties
de la construction angulaire préalable ; le développement ne les sélectionne pas.

\begin{theorem}[Moments de profondeur et représentation de Mellin]
\label{amp-afin-thm-momentos}
La série
\begin{equation}
 M_b(z)=\sum_{n\geq1}\frac{b_n}{n^z}
 \label{amp-afin-momento}
\end{equation}
converge localement uniformément pour $\operatorname{Re}z>0$.
Si $\sigma=\operatorname{Re}z>0$, son reste satisfait
\begin{equation}
 \left|M_b(z)-\sum_{n=1}^{N}\frac{b_n}{n^z}\right|
 \leq3\left(1+\frac{|z|}{\sigma}\right)(N+1)^{-\sigma}.
 \label{amp-afin-cota}
\end{equation}
Dans le domaine de convergence absolue $\operatorname{Re}z>1$,
\begin{equation}
 M_b(z)=(3^{1-z}-4^{1-z})\zeta(z),\qquad
 \zeta(z)=\sum_{n\geq1}n^{-z}.
 \label{amp-afin-reconocimiento}
\end{equation}
Pour $\operatorname{Re}z>0$, on a en outre
\begin{equation}
 \Gamma(z)M_b(z)=
 \int_0^\infty t^{z-1}
 \left(\frac3{e^{3t}-1}-\frac4{e^{4t}-1}\right)dt,
 \label{amp-afin-mellin}
\end{equation}
où $\Gamma(z)=\int_0^\infty u^{z-1}e^{-u}\,du$.
\end{theorem}
\begin{proof}
La sommation par parties et $|B_n|\leq3$ donnent, en passant à la limite supérieure,
\[
 \sum_{n>N}b_n n^{-z}
 =-B_N(N+1)^{-z}
   +\sum_{n>N}B_n\bigl(n^{-z}-(n+1)^{-z}\bigr).
\]
La différence de puissances est bornée en module par
$|z|\int_n^{n+1}x^{-\sigma-1}dx$. Cela prouve
\eqref{amp-afin-cota}, la convergence et son uniformité sur les compacts.
Pour $\sigma>1$, la convergence absolue permet de regrouper
les multiples de trois et de quatre séparément, ce qui donne
\eqref{amp-afin-reconocimiento}.

La série géométrique du transport, évaluée en $s=e^{-t}$, produit
\[
 \sum_{n\geq1}b_n e^{-nt}
 =\frac3{e^{3t}-1}-\frac4{e^{4t}-1}=:K_b(t).
\]
Pour $\sigma>1$, on peut intégrer terme à terme parce que
$\sum|b_n|n^{-\sigma}<\infty$. La substitution $u=nt$ donne
$\Gamma(z)n^{-z}$ et prouve~\eqref{amp-afin-mellin} dans ce demi-plan.
À l’origine, les termes $1/t$ s’annulent :
$K_b(t)=1/2+O(t)$. À l’infini, $K_b(t)=O(e^{-3t})$.
L’intégrale est donc holomorphe pour $\operatorname{Re}z>0$
par convergence locale uniforme. La série l’est aussi par la borne
précédente ; le principe d’identité prolonge l’égalité à ce domaine.
\end{proof}

En $0<\operatorname{Re}z\leq1$, l’ordre croissant de profondeur
fixe la lecture convergente de la série. La réalisation comme trace
absolue de l’opérateur diagonal correspond à $\operatorname{Re}z>1$.

\begin{corollary}[Évaluations compatibles avec l’extrémité constitutive]
\label{amp-afin-cor-extremos}
On a
\begin{equation}
 M_b(2)=\frac{\zeta(2)}{12},\qquad
 M_b(1)=\log(4/3)=-\log f(1),\qquad f(1)=\frac34.
 \label{amp-afin-extremos}
\end{equation}
\end{corollary}
\begin{proof}
La première égalité résulte de $3^{-1}-4^{-1}=1/12$.
Pour la seconde, les sommes partielles exactes sont
$H_{\lfloor N/3\rfloor}-H_{\lfloor N/4\rfloor}$, avec
$H_j=\sum_{k=1}^j1/k$ et $H_0=0$.
La comparaison intégrale de $1/x$ entre ces deux indices donne la limite
$\log(4/3)$. Annuler $1-s$ avant d’évaluer le quotient déterminantal
donne $f(1)=3/4$.
\end{proof}

Le registre affine retient les événements et leur interaction avec la mémoire
préalable ; le tableau des traces retient l’incidence des deux sous-espaces
à toutes les puissances du transport ; le moment applique une pondération
de profondeur à ce tableau déjà produit. Sa périodicité spectrale
coexiste avec une histoire d’événements qui se poursuit après le retour
de phase. Inverser $C$ conserve $b_n$, parce que $d\mid n$ est invariant
sous $n\mapsto-n$ ; le coefficient orienté de la résolvante de
\eqref{iii:eq:resolvent}, avec ses marques fixes, change de signe.
La lecture conjointe conserve une distinction que le déterminant
scalaire n’enregistre pas à lui seul. La forme quadratique et la somme des
coordonnées du registre maintiennent leurs types mathématiques ; cette
composition ne leur assigne pas d’unités d’énergie ni de charge électrique.

\subsection{Composition résiduelle et produit de moments}
\label{amp-afin-sec-productos}

La composition ultérieure des lecteurs conserve également l’opération
effectuée sur la profondeur. Dans la réalisation
$\mathcal H=\ell^2(\mathbb N_{>0})$, nous fixons
$N\delta_n=n\delta_n$. Soient $a,c$ deux tableaux résiduels périodiques
déjà produits et $Q_a\delta_n=a(n)\delta_n$,
$Q_c\delta_n=c(n)\delta_n$. Lorsqu’on les prolonge tous deux à la période commune,
les deux lecteurs agissent sur le même indice. Pour
$\operatorname{Re}z>1$, nous écrivons
$M_a(z)=\operatorname{Tr}(N^{-z}Q_a)$ et de même $M_c$.

\begin{proposition}[Composition diagonale et composition tensorielle]
\label{amp-afin-prop-productos}
On a les identités
\begin{align}
 \operatorname{Tr}(N^{-z}Q_aQ_c)
 &=\sum_{n\geq1}\frac{a(n)c(n)}{n^z},
 \label{amp-afin-diagonal}\\
 M_a(z)M_c(z)
 &=\operatorname{Tr}_{\mathcal H\otimes\mathcal H}
 \bigl((N\otimes N)^{-z}(Q_a\otimes Q_c)\bigr)\notag\\
 &=\sum_{k\geq1}\frac{(a*c)(k)}{k^z},
 \label{amp-afin-tensor}\\
 (a*c)(k)&=\sum_{d\mid k}a(d)c(k/d),
 \qquad \operatorname{Re}z>1.
 \label{amp-afin-convolucion}
\end{align}
Tous les opérateurs figurant à l’intérieur de ces traces sont à trace.
\end{proposition}
\begin{proof}
Le produit diagonal agit sur $\delta_n$ au moyen de
$a(n)c(n)$, ce qui prouve~\eqref{amp-afin-diagonal}.
Dans $\delta_m\otimes\delta_n$, l’opérateur $N\otimes N$
a pour valeur propre $mn$ et $Q_a\otimes Q_c$ agit au moyen de
$a(m)c(n)$. Si $\sigma=\operatorname{Re}z>1$, la somme des
modules de ces coefficients spectraux est bornée par
$\|a\|_\infty\|c\|_\infty(\sum n^{-\sigma})^2<\infty$.
Cette borne prouve la condition de classe de trace et permet de factoriser
la somme double ou de la regrouper par $k=mn$. Les paires de produit $k$
sont exactement $(d,k/d)$, avec $d\mid k$.
La borne analogue avec une seule somme prouve le cas diagonal.
\end{proof}

Dans les lecteurs de cette section, prendre $a(n)=b_n$ et
$c(n)=\langle v,C^nu\rangle$ donne, pour $n=3$,
$a(3)c(3)=-3$, tandis que $(a*c)(3)=3$.
Cette différence de coefficients met en évidence l’information conservée par
chaque opération. La composition diagonale maintient une profondeur commune ;
la tensorielle conserve deux profondeurs avant de les regrouper par leur produit.
La réalisation tensorielle n’attribue pas d’indépendance statistique aux
histoires HMT qui ont produit les lecteurs. Les dérivées régularisées et
les changements de normalisation d’autres moments constituent des opérations
ultérieures distinctes des deux compositions démontrées ici.

\subsection{Factorisation à travers l’opérateur constitutif}

Conservons la coordinatisation angulaire relevée de la
section~\ref{iii:sec:hojas}. Pour rendre explicite la conversion angulaire,
écrivons
\begin{equation}
 A_{\deg}=\frac{180x}{\pi},\qquad
 C^*_{\deg}=\frac{180y}{\pi},\qquad
 q_\pm=e^{-x\mp y},\qquad s_\pm=q_\pm^{30}.
 \label{iii:eq:barbero-angular-lift}
\end{equation}
L’orientation considérée dans \eqref{iii:eq:rchannels} correspond à
$0<y<x$ ; l’échange de feuillets s’exprime par $y\mapsto-y$ sur le
domaine plus large $x>|y|$. Les exponentielles réelles agissent sur ces
relèvements, qui conservent les valeurs de $x\pm y$ pendant la
composition.

La fonctionnelle orientée de Barbero--Immirzi est
\begin{equation}
 \gamma=\frac{\pi A_{\deg}C^*_{\deg}}{180}
   +12\bigl[S_{90}(q_-)-S_{90}(q_+)\bigr]
   -\bigl[S_{120}(q_-)-S_{120}(q_+)\bigr].
 \label{iii:eq:barbero-functional}
\end{equation}
La première composante est l’évaluation angulaire ; la seconde est la
différence du lecteur de la chaîne fondamentale dans les canaux conjugués.
Les arguments sont reçus du transport angulaire précédemment construit.
La formule détermine un scalaire adimensionnel sur cette structure de
départ.

Soit $\psi:(3/4,1)\to(0,1)$ l’inverse de $f$ construit au moyen de
\eqref{iii:eq:cubic}. Définissons
\begin{equation}
 \mathcal H(s)=\frac{12s^3}{1-s^9}
             -\frac{s^4}{1-s^{12}}
             -\frac{(\log s)^2}{20\pi},\qquad 0<s<1.
 \label{iii:eq:barbero-h}
\end{equation}
Cette fonction est analytique réelle sur son domaine. Les exposants $3,4,9,12$
résultent de l’expression des hauteurs $90,120$ et des retours $270,360$ en
unités de $30$ ; le coefficient de la partie logarithmique est fixé lorsqu’on
transporte la contribution angulaire.

\begin{theorem}[Expression constitutive orientée de Barbero--Immirzi]
\label{iii:thm:barbero-factorization}
Dans la coordinatisation \eqref{iii:eq:barbero-angular-lift},
\begin{equation}
 \gamma=\mathcal H(s_-)-\mathcal H(s_+)
       =\mathcal H\bigl(\psi(r_-)\bigr)
        -\mathcal H\bigl(\psi(r_+)\bigr).
 \label{iii:eq:barbero-factorization}
\end{equation}
Sur la représentation bidimensionnelle des canaux,
\begin{equation}
 \gamma=-\operatorname{Tr}
    \left[\mathcal R\,\mathcal H\bigl(\psi(\mathcal V)\bigr)\right],
 \label{iii:eq:barbero-trace}
\end{equation}
où $\operatorname{Tr}$ est la trace ordinaire, avec
$\operatorname{Tr}I=2$.
\end{theorem}
\begin{proof}
La substitution $s=q^{30}$ dans \eqref{iii:eq:barbero-return-series} donne
\[
 S_{90}(q)=\frac{s^3}{1-s^9},\qquad
 S_{120}(q)=\frac{s^4}{1-s^{12}}.
\]
De plus, \eqref{iii:eq:barbero-angular-lift} implique
\[
 \log s_\pm=-30(x\pm y)
 =-\frac\pi6(A_{\deg}\pm C^*_{\deg}).
\]
Par différence de carrés,
\begin{align}
 \frac{(\log s_+)^2-(\log s_-)^2}{20\pi}
 &=\frac{900\bigl[(x+y)^2-(x-y)^2\bigr]}{20\pi}\\
 &=\frac{180xy}{\pi}
 =\frac{\pi A_{\deg}C^*_{\deg}}{180}.
 \label{iii:eq:barbero-log-term}
\end{align}
La somme de ces identités prouve la première égalité de
\eqref{iii:eq:barbero-factorization}. La seconde utilise l’inversion
$\psi(f(s_\pm))=s_\pm$ démontrée dans la
section~\ref{iii:sec:constitutiva}.

Pour l’expression opératorielle, les projecteurs de feuillet sont de rang un et
$\mathcal RP_\pm=\pm P_\pm$. Le calcul spectral fournit
\[
 \mathcal H\bigl(\psi(\mathcal V)\bigr)
 =\mathcal H(s_+)P_++\mathcal H(s_-)P_-.
\]
Multiplier par $\mathcal R$ et prendre la trace ordinaire donne
$\mathcal H(s_+)-\mathcal H(s_-)$. Son opposé est $\gamma$.
\end{proof}

La marque d’orientation $\mathcal R$ fait partie de la fonctionnelle. Si
l’on maintient cette marque fixe et que l’on échange les canaux de $\mathcal V$,
le signe de $\gamma$ change. L’évaluation du spectre sans ses étiquettes conserve
la paire non ordonnée, tandis que \eqref{iii:eq:barbero-trace} conserve la
contribution orientée. Si l’on emploie la trace normalisée
$\tau_2=\operatorname{Tr}/2$, la même formule s’écrit
$\gamma=-2\tau_2[\mathcal R\mathcal H(\psi(\mathcal V))]$.

En revanche, un changement simultané de représentation
$(\mathcal R,\mathcal V)\mapsto
(U\mathcal RU^{-1},U\mathcal VU^{-1})$ conserve la fonctionnelle :
le calcul spectral se transporte par conjugaison et la trace est
invariante. Échanger la réponse par rapport à une orientation fixe
et transporter conjointement la réponse et l’orientation sont donc
des opérations différentes.

Cette normalisation est compatible avec les amplifications de
\eqref{iii:eq:logvac}. Pour
$\mathcal V_m=\mathcal V\otimes I_{9^m}$ et
$\mathcal R_m=\mathcal R\otimes I_{9^m}$,
\begin{equation}
 \gamma=-\frac1{9^m}\operatorname{Tr}
 \left[\mathcal R_m\mathcal H\bigl(\psi(\mathcal V_m)\bigr)\right].
 \label{iii:eq:barbero-amplification}
\end{equation}
En effet, le calcul fonctionnel conserve le produit tensoriel et la trace
de $I_{9^m}$ est $9^m$. La multiplicité de la représentation se sépare
de l’évaluation orientée.

\subsection{Reconstruction à partir des coordonnées du vide}

Les identités \eqref{iii:eq:four-identities} permettent d’exprimer la même
fonctionnelle au moyen de l’impédance et de la vitesse normalisées :
\begin{align}
 \gamma={}&
 \mathcal H\!\left(\psi\!\left(
     \sqrt{\frac{\widehat Z}{\widehat c}}\right)\right)
 -\mathcal H\!\left(\psi\!\left(
     \frac1{\sqrt{\widehat Z\widehat c}}\right)\right).
 \label{iii:eq:barbero-zc}
\end{align}
Le domaine positif de cette reconstruction est décrit par
\begin{equation}
 1<\widehat Z\widehat c<\frac{16}{9},\qquad
 \frac9{16}<\frac{\widehat Z}{\widehat c}<1.
 \label{iii:eq:barbero-zc-domain}
\end{equation}
En effet, ces inégalités équivalent à
$3/4<r_\pm<1$ lorsqu’on substitue les racines positives de
\eqref{iii:eq:four-identities}. Pour l’orientation $y>0$, on conserve
en outre $r_-<r_+$, ou de manière équivalente $\widehat Z<1$. Les points du
domaine sont utilisés avec la compatibilité angulaire déjà établie par
le TPK ; l’inversion est une reconstruction ultérieure de ses canaux.
Les coordonnées sont les coordonnées internes de \eqref{iii:eq:four} : un changement
d’unités doit être défait au moyen de \eqref{iii:eq:undo-units} avant
d’appliquer $\psi$.

L’échange $y\mapsto-y$ agit comme
$(\widehat Z,\widehat c)\mapsto(\widehat Z^{-1},\widehat c)$ et
produit $\gamma\mapsto-\gamma$. Dans le cas équilibré $y=0$,
les deux canaux coïncident et \eqref{iii:eq:barbero-factorization} donne
$\gamma=0$. La vitesse normalisée conserve le produit des
canaux ; sa lecture conjointe avec l’impédance retrouve les deux
réponses et l’orientation nécessaire à l’évaluation de la fonctionnelle.
Dans ce cas équilibré, $\mathcal V=rI$ a un spectre dégénéré :
la marque $\mathcal R$ est conservée comme donnée de représentation, bien que
le scalaire $r$ ne distingue plus ses projecteurs. Pour un spectre simple,
en revanche,
$P_+=(\mathcal V-r_-I)/(r_+-r_-)$ et $P_-=I-P_+$.

\section{Évaluation constitutive et réalisation dimensionnelle}
\label{iii:sec:evaluacion}

L’évaluation réunit maintenant les constructions précédentes dans leur ordre effectif.
Les développements régionaux de $\pi$, $\varphi$ et $e$ proviennent du raffinement
nonadique ; le registre $K$ conserve ses douze positions. Sa lecture périodique
détermine la section positionnelle de $\alpha$. Sur celle-ci agissent le caractère
d’action et le retour régional, puis la paire angulaire et, enfin, la
réponse constitutive. Cette évaluation est postérieure à la sélection des
règles : aucun de ses chiffres n’est utilisé pour choisir un canal.

\subsection{Évaluation de la chaîne interne}

Pour reproduire le tableau, on emploie les développements régionaux à mille
décimales engendrés par le certificat nonadique et l’entier $N_K$ de
\eqref{eq:k-numerador}. La lecture utilisée est
\[
 \alpha=\pi+e-\varphi-4-\frac{N_K}{1000^{12}-1}.
\]
La section finie de douze blocs et la section prolongée ont leurs
définitions et leur différence exacte dans le chapitre du registre. Ici, on
maintient la seconde fixe pendant toute l’évaluation. On n’échange pas
les lectures lorsque la précision augmente, et on ne remplace pas l’opérateur analytique
de la seconde voie par une coïncidence décimale avec la première.

\begin{longtable}{@{}ll@{}}
\toprule
Objeto & Évaluation décimale (24 chiffres significatifs)\\
\midrule\endhead
$\alpha$ & $0.00729735256928380099728511\ldots$ \\
$H_5$ & $1.05457181764615446962582\ldots$ \\
$\eta_{\rm ret}$ & $1.05457181691503906113369\ldots$ \\
$A_{\deg}$ & $7.29735256928380099728511\ldots$ \\
$C^*_{\deg}$ & $2.12373833896864572426195\ldots$ \\
$q_+$ & $0.848377942576404374945737\ldots$ \\
$q_-$ & $0.913660151150557285295519\ldots$ \\
$r_+$ & $0.999999628548249363178695\ldots$ \\
$r_-$ & $0.999724137189518364131517\ldots$ \\
$\widehat\varepsilon$ & $0.999999257096636702760442\ldots$ \\
$\widehat\mu$ & $0.999448350479326935089982\ldots$ \\
$\widehat Z$ & $0.999724508538937215455554\ldots$ \\
$\widehat c$ & $1.00027631048615752610639\ldots$ \\
$\gamma$ & $0.274007672694459274747255\ldots$ \\
$G_{\rm Cat}$ & $0.915965594177219015054604\ldots$ \\
$\rho$ & $0.267949192431122706472554\ldots$ \\
\bottomrule
\end{longtable}


Le calcul s’effectue avec 120 chiffres de travail ; on montre 24 chiffres
significatifs pour faciliter la comparaison entre grandeurs de différents
ordres. Les 120 chiffres sont une précision d’évaluation, non une incertitude
expérimentale ni, à eux seuls, un certificat d’erreur par intervalles.
Le paquet conserve les entrées engendrées, leurs empreintes, les opérations et
les résidus des identités vérifiées. Les preuves algébriques du
corps du texte établissent ces identités indépendamment de ce tableau.

Le sens des quatre dernières coordonnées constitutives est conservé
lors de l’évaluation : $\widehat\varepsilon$ et $\widehat\mu$ sont les réponses
quadratiques par feuillet, $\widehat Z$ est leur orientation relative et
$\widehat c$ est l’inverse de leur produit. La proximité de certains chiffres avec
l’unité exprime le petit défaut du quotient $3$--$4$ sur ces
canaux. Son origine se voit exactement dans
\[
 1-f(s)=\frac{s^3}{1+s+s^2+s^3}.
\]
La formule permet d’étudier la réponse avant de lui attribuer une grandeur
dimensionnelle : le signe du défaut, son ordre cubique et sa monotonie sont
des conséquences du même opérateur.

\subsection{Coordinatisations d’unités et covariance}

Soient $\mathcal L_S$, $\mathcal L_Q$ et $\mathcal L_C$ les droites
d’action, de charge et de vitesse. Une coordinatisation physique positive évalue les sections
$u_S$, $u_Q$ et $u_C$ en unités d’action, de charge et de vitesse. La réalisation
\eqref{iii:eq:dimensions} produit alors une impédance, une vitesse,
une perméabilité et une permittivité, avec leurs dimensions respectives.
Les coordonnées du tableau précédent restent adimensionnelles.

\begin{proposition}[Covariance de la réalisation constitutive]
Si les sections de référence changent selon
$u_S'=a u_S$, $u_Q'=b u_Q$, $u_C'=d u_C$, avec $a,b,d>0$, les
coordonnées qui représentent les mêmes grandeurs satisfont
\[
 \widehat Z'=a^{-1}b^2\widehat Z,\qquad
 \widehat c'=d^{-1}\widehat c,\qquad
 \widehat\mu'=a^{-1}db^2\widehat\mu,\qquad
 \widehat\varepsilon'=adb^{-2}\widehat\varepsilon.
\]
Les relations $\widehat\mu'\widehat\varepsilon'=(\widehat c')^{-2}$
et $\widehat\mu'/\widehat\varepsilon'=(\widehat Z')^2$ sont conservées.
\end{proposition}
\begin{proof}
On égale chaque grandeur de \eqref{iii:eq:dimensions} à son expression
dans la nouvelle base et on isole sa coordonnée. Les facteurs du produit
sont $d^2$, et ceux du quotient sont $a^{-2}b^4$, exactement ceux requis
par les coordonnées transformées de $c$ et $Z$.
\end{proof}

Cette covariance distingue deux actes : construire une section et exprimer
cette section dans des unités choisies. Le théorème spectral fixe les coordonnées
dans la normalisation interne déclarée. Une comparaison SI doit également présenter
l’évaluation des trois sections de référence. Les ajuster
au moyen des grandeurs que l’on souhaite reproduire serait un étalonnage de
la coordinatisation, non une vérification indépendante de ces grandeurs.

Pour reconstruire les canaux à partir de coordonnées exprimées dans la coordinatisation
primée, on défait ce changement avant d’utiliser la racine cubique :
\begin{equation}
 \widehat Z=ab^{-2}\widehat Z',\quad
 \widehat c=d\widehat c',\quad
 \widehat\mu=ad^{-1}b^{-2}\widehat\mu',\quad
 \widehat\varepsilon=a^{-1}d^{-1}b^2\widehat\varepsilon'.
 \label{iii:eq:undo-units}
\end{equation}
Par exemple, les entrées de l’inverse spectral sont
$r_+=(\widehat Z\widehat c)^{-1/2}$ et
$r_-=(\widehat Z/\widehat c)^{1/2}$, calculées après
\eqref{iii:eq:undo-units}. Introduire directement les coordonnées primées
dans ces expressions changerait la normalisation de l’opérateur et pourrait
les faire sortir de l’intervalle $(3/4,1)$. La covariance dimensionnelle conserve
les grandeurs physiques ; l’inversion cubique reconstruit le transport
dans la coordinatisation interne dans laquelle il a été défini.

% Incorporación focal de dependencias anteriores; no sustituye al núcleo común.
% Procedencia detallada: technical/CIERRE_DEPENDENCIAS_ANTERIORES_VI.md.
% Los valores de prueba se publican después de sus operadores; no son entradas.
\section{Coordonnées engendrées, action et normalisation électronique}
\label{vi:dep:seccion}

La construction des collections utilise des résultats des étapes précédentes :
les caractères régionaux, le registre dodécaphasique, la droite d’action et la
section électronique. Nous réunissons ici les opérations qui les relient et les
preuves nécessaires à leur composition. Leur ordre est causal : les régions et leurs
registres sont produits sur APP par TRIT et TPK ; les évaluations réelles
sont effectuées ensuite ; le choix des unités intervient lors de la réalisation des
grandeurs dimensionnelles.

Les conditions initiales et la sélection régionale ont été développées dans
la section précédant l’arithmétique des histoires. Le domaine du registre
incidentiel, ses dénombrements et la reconstruction du registre sont définis
expressément ci-dessous. La sélection d’une collection concrète de
visites et de traces est une opération supplémentaire à l’évaluation d’un registre :
elle nécessite de fixer ses multiplicités, le transport de frontière et la
couverture des traces admises. Les formules de reconstruction qui sont
démontrées ici permettent d’évaluer et de vérifier ce registre ; son
inversibilité ne sélectionne pas à elle seule la collection incidentielle.

\subsection{Lecteurs régionaux et compatibilité de l’évaluation}
\label{vi:dep:regiones}

La clôture, la propagation et l’auto-échelle sont trois régions d’une même
génération. La région de clôture comporte cinq émissions et une multiplicité
\(432+4\cdot144=1008\). Dans la jauge orientée du catalogue, ses émissions
sont
\[
\begin{array}{c|r}
(501,614,498,169,272,272)&432\\
(810,923,870,169,272,272)&144\\
(810,983,810,169,272,272)&144\\
(870,923,810,169,272,272)&144\\
(870,923,810,418,674,521)&144.
\end{array}
\]
Le mot ternaire commun est \(010211\) ; les représentants de propagation
et d’auto-échelle sont \(201101\) et \(121200\). Le recensement et l’orientation
précèdent les évaluations qui suivent.

Aux cinq positions du lecteur de clôture, la projection uniforme
\(P_5=\mathbf1\mathbf1^{\mathsf T}/5\), avec
\(P_5\lambda=\lambda\) et \(\mathbf1^{\mathsf T}\lambda=1\), détermine
\(\lambda=\mathbf1/5\). Cette normalisation des positions du lecteur
est distincte de la distribution des germes indiquée dans le tableau.
Dans le régime elliptique \(J^2=-I\), la multiplication donne
\[
(I+aJ)(I+bJ)=(1-ab)I+(a+b)J.
\]
Pourvu que \(1-ab\ne0\), la pente composée est
\(a\oplus b=(a+b)/(1-ab)\). Deux duplications de \(1/5\) donnent
\(5/12\) et \(120/119\). Le compensateur positif \(1/q\), \(q>1\),
qui ramène la pente à un satisfait
\[
\frac{120/119-1/q}{1+(120/119)/q}=1,
\qquad q=239.
\]
Par conséquent, le caractère de clôture de ce lecteur est
\begin{equation}
 \Pi_{\rm cl}=4\left(4\arctan\frac15-\arctan\frac1{239}\right).
 \label{vi:dep:lector-clausura}
\end{equation}
La règle de composition, sa compensation et le choix de la branche angulaire
fixent cette valeur ; l’identité angulaire conventionnelle la reconnaît comme \(\pi\).

Le caractère de propagation se réalise dans une collection formelle commutative
\(P_t\), normalisée par \(P_0=1\), \(P'_0=1\) et
\(P_{t+s}=P_tP_s\). Si \(P_t=\sum_{n\ge0}a_nt^n\), la comparaison
du coefficient de \(s^nt\) donne
\[
 (n+1)a_{n+1}=a_n,\qquad a_0=1.
\]
Ainsi, \(a_n=1/n!\) ; la représentation à une unité de paramètre est
\(\Pi_{\rm pr}=\sum_{n\ge0}1/n!\), reconnue comme \(e\).
L’unité et le signe du générateur appartiennent au lecteur : la seule loi
de semi-groupe admettrait d’autres générateurs.

Pour l’auto-échelle, TRIT sélectionne les modes
\((\Sigma,\Pi,\Pi)\). Leurs trois transitions cycliques sont
\(\Sigma\to\Pi\), \(\Pi\to\Pi\) et \(\Pi\to\Sigma\). Leur matrice
d’incidence, dans l’ordre aréal--volumétrique, est
\begin{equation}
 F=\begin{pmatrix}0&1\\1&1\end{pmatrix}.
 \label{vi:dep:incidencia-hojas}
\end{equation}
Le vecteur propre positif normalisé comme \((1,x)^{\mathsf T}\)
satisfait \(x=\lambda\), \(1+x=\lambda x\) ; par conséquent,
\(x^2-x-1=0\). Sa racine positive est la coordonnée d’auto-échelle
\(\varphi\). Cette équation reconnaît le caractère de l’incidence déjà
produite par le calendrier.

\begin{proposition}[Évaluation régionale par intervalles compatibles]
\label{vi:dep:prefijos}
Les trois lecteurs précédents disposent d’intervalles rationnels de largeur
arbitrairement petite. Leurs préfixes décimaux finis sont déterminés par
raffinement et sont compatibles par troncature.
\end{proposition}
\begin{proof}
Pour \(0<x<1\), deux sommes alternées consécutives de
\(\sum_{j\ge0}(-1)^jx^{2j+1}/(2j+1)\) encadrent \(\arctan x\).
Le premier terme omis borne l’erreur. Cela fournit des intervalles
rationnels pour \eqref{vi:dep:lector-clausura}.
Si \(E_n=\sum_{j=0}^n1/j!\), alors, pour \(n\ge1\),
\[
 E_n<\Pi_{\rm pr}<E_n+
 \frac{n+2}{(n+1)(n+1)!}.
\]
La borne résulte de la majoration des quotients successifs de la queue par
\(1/(n+2)\). Enfin, \(F^n\) a pour entrées
\(F_{n-1},F_n,F_{n+1}\), où \(F_0=0,F_1=1\) et
\(F_{n+1}=F_n+F_{n-1}\). L’identité
\(F_{n+1}^2-F_nF_{n+2}=(-1)^n\) montre que deux quotients consécutifs
encadrent \(\varphi\) avec une largeur \(1/(F_nF_{n+1})\).
On peut prendre l’intersection des intervalles successifs pour conserver
l’emboîtement. Après la reconnaissance conventionnelle, l’irrationalité de
\(\pi,e,\varphi\) assure qu’aucune coordonnée ne se trouve sur une frontière
décimale rationnelle. Pour chaque profondeur finie, un intervalle
suffisamment étroit tient dans une unique cellule décimale. Son indice
détermine le préfixe et l’emboîtement démontre leur compatibilité.
\end{proof}

\subsection{Coordinatisation dodécaphasique signée de l’état TPK plein}
\label{vi:dep:estado-pleno}

L’évolution qui produit le registre est formulée sur l’état complet du
Topological Phase Kernel. Notons
\(\mathsf{TPKFullState}\) l’espace des états qui conserve conjointement
position APP, régime TRIT, route, feuillet, orientation, retenue, frontière,
mémoire, registre incidentiel et coordonnées terminales. Dans une coordinatisation canonique,
\[
 X=(X_{\rm APP},X_{\rm TRIT},X_{\rm ruta},X_{\rm mem},
       C_{\rm term},U_{12}^{\rm sgn},X_{\partial},\ldots).
\]
La dynamique de cent huit pas agit par la composition
\begin{equation}
 X_{\rm term}
 =\bigl(\mathcal U_{108}\circ\cdots\circ\mathcal U_1\bigr)(X_{\rm seed}),
 \qquad
 \mathcal U_t=\mathsf{Upd}_t\circ\mathsf{Tra}_t\circ\mathsf{Sel}_t,
 \label{vi:dep:estado-terminal-108}
\end{equation}
où \(\mathsf{Sel}_t\) applique le sélecteur tritique de phase,
\(\mathsf{Tra}_t\) transporte feuillet, route et orientation, et
\(\mathsf{Upd}_t\) incorpore la retenue et la mémoire produites à
l’instant \(t\). Chaque facteur agit donc sur l’état plein et la
composition conserve la généalogie des douze fenêtres nonadiques.

Le registre incidentiel de la section suivante est la coordonnée
\(\mathcal R_{12}(X_{\rm term})\). Son évaluation par les canaux
\(90/120\) définit
\[
 \mathcal E_{108}^{90,120}\circ\mathcal R_{12}:
 \mathsf{TPKFullState}\longrightarrow\mathbb Z^{25},
 \qquad
 X\longmapsto C(X)=(b^{90}(X),b^{120}(X),Q(X)).
\]
La projection dodécaphasique signée est
\[
 \mathcal S_{12}:=\operatorname{pr}_{U_{12}^{\rm sgn}}:
 \mathsf{TPKFullState}\longrightarrow\mathbb Z^{12}.
\]
Sur le sous-réseau compatible produit par la dynamique, les deux coordinatisations
satisfont l’identité opératoire
\begin{equation}
 \mathcal S_{12}
 =\Pi_H\circ\mathcal E_{108}^{90,120}\circ\mathcal R_{12}
 =\Pi_H\circ\operatorname{pr}_{C}.
 \label{vi:dep:identidad-productor-e108}
\end{equation}
Ainsi, la signature dodécaphasique est obtenue par projection et évaluation du même
état qui conserve le registre ; son domaine inclut l’orientation et la mémoire
qui interviennent dans chaque dénombrement. La projection brute
\(\rho_{108}:\mathsf{TPKFullState}\to\mathsf{CT108RawProjection}\)
désigne, en revanche, la coordinatisation qui oublie ces coordonnées après la
production. Cette séparation des types fixe l’ordre causal entre génération,
représentation signée et lecture réduite.

\subsection{Des incidences temporelles au registre dodécaphasique}
\label{vi:dep:registro}

\begin{definition}[Domaine du registre incidentiel orienté]
\label{vi:dep:libro-dominio}
Un registre \(\mathscr L\) contient une collection finie non vide de routes,
leurs visites identifiées et une collection finie de traces. Chaque route
conserve une visite par instant pendant les \(108\) premiers instants,
et les visites ultérieures qu’utilisent ses traces. Une visite conserve
\[
 v=(\gamma,t,i,j,s,\tau,w,\varepsilon_\partial),
\]
où \(\gamma\) identifie la route, \(t\) son instant, \(1\le i,j\le9\)
sa position APP, \(s\in\{\Sigma,\Pi,0\}\) son feuillet,
\(\tau\in\{-1,0,1\}\) son orientation tritique et \(w\) sa multiplicité
entière positive. La coordinatisation incidentielle évalue
\[
 n_v=
 \begin{cases}
 i+j,&s=\Sigma,\\
 ij,&s=\Pi,\\
 9,&s=0.
 \end{cases}
 \qquad
 r_v=1+((n_v-1)\bmod9),\quad q_v=(n_v-r_v)/9.
\]
La lecture \(n_v=9\) fixe la convention de seuil de cette coordinatisation.
Pour \(r_v=9\), l’étiquette \(\varepsilon_\partial\) distingue
couronne et intérieur ; elle ne s’obtient pas uniquement à partir du résidu.

Une trace \(\theta=(v_1,\ldots,v_N)\) parcourt des instants consécutifs d’une
même route. Elle conserve une multiplicité \(u(\theta)\in\mathbb N_{>0}\)
et une unité de longueur réduite
\(\lambda(\theta)\in\mathbb N_{>0}\), avec
\[
 N\lambda(\theta)=120(1+3j(\theta)),\qquad j(\theta)\in\mathbb N_0.
\]
Son canal est
\(\sigma(\theta)=\operatorname{sgn}(\sum_{a=1}^N\tau(v_a))\).
La collection dénombrée dans les canaux \(120\) comprend les traces de
signe non nul. La longueur réduite et le nombre d’instants sont
reliés par \(\lambda\) ; ils ne sont pas identifiés l’un à l’autre.
\end{definition}

La définition spécifie le domaine de l’évaluation, y compris les
données qui doivent être transmises au lecteur. Ses multiplicités et la collection
de traces appartiennent à la sélection dynamique du registre. Une énumération
du domaine ou une vérification de format ne remplace pas cette sélection.

La génération conserve un registre \(\mathscr L\) de visites et de traces.
La coordinatisation temporelle le divise en fenêtres
\(E_m=\{9m-8,\ldots,9m\}\), \(1\le m\le12\),
équivalentes aux positions \(m=\beta+1\) du registre
\(\mathcal R_{12}\). Chaque fenêtre retient sa sous-route,
son orientation, son incrément de frontière et ses incidences locales.
Chaque visite enregistre position APP, feuillet arithmétique, instant, multiplicité
et orientation de frontière. À chaque évaluation \(n_v\) est conservée la
division \(n_v=r_v+9q_v\), avec \(1\le r_v\le9\). Dans la fenêtre
\(m\), le lecteur incidentiel exprime
\begin{align}
 A_m&=\sum_{v\in\mathscr L_m}w(v)
      \mathbf1_{\{1,3,5,7\}}(r_v),\notag\\
 C_m&=\sum_{j\ge0}
       \bigl(\nu^-_{120,m}(j)-\nu^+_{120,m}(j)\bigr),\notag\\
 V_m&=\sum_{v\in\mathscr L_m}w(v)
      \mathbf1_{\{9\}}(r_v)\,\varepsilon_{\partial}(v).
 \label{vi:dep:incidencias}
\end{align}
Ici les traces de chaque fenêtre ont un support fini,
\(\varepsilon_{\partial}=+1\) dans la couronne et \(-1\) à l’intérieur,
et les étiquettes \(\pm,120\) appartiennent au transport conservé.
La multiplicité \(w\) et les traces \(\nu\) sont celles obtenues par ce
transport ; elles ne sont pas reconstruites en choisissant une valeur de \(\alpha\).

En termes de traces identifiées, la seconde somme a pour
expression finie
\[
 C_m=-\sum_{\theta:\,t(v_1)\in E_m}u(\theta)\sigma(\theta).
\]
En effet, regrouper par \(j(\theta)\) et par signe produit
\(\nu^-_{120,m}(j)-\nu^+_{120,m}(j)\). Cette égalité spécifie
l’opération effectivement appliquée à chaque registre ; elle conserve
séparés le canal, la multiplicité et la longueur.

\begin{proposition}[Mise à jour et concaténation du registre]
\label{vi:dep:libro-actualizacion}
La collection de visites et de traces admises par le transport étant fixée, leurs
dénombrements sont mis à jour par des incréments entiers. Lorsqu’on ajoute une
visite \(v\) de la fenêtre \(m\), l’incrément de \((A_m,C_m,V_m)\) est
\[
 \left(w(v)\mathbf1_{\{1,3,5,7\}}(r_v),\,0,\,
 w(v)\mathbf1_{\{9\}}(r_v)\varepsilon_\partial(v)\right).
\]
Lorsqu’une trace \(\theta\) dont l’origine appartient à \(E_m\) est achevée,
l’incrément est \((0,-u(\theta)\sigma(\theta),0)\).
Pour deux segments consécutifs, le dénombrement de leur concaténation est
la somme des dénombrements des visites et des traces contenues dans chaque
segment, plus celui des traces qui franchissent leur jonction.
\end{proposition}
\begin{proof}
Chaque somme de \eqref{vi:dep:incidencias} est indexée par des identifiants
de visites ou de traces ; ajouter un identifiant ajoute donc
exactement le terme écrit. Pour la concaténation, les visites
se répartissent entre les deux segments. Les traces se répartissent en
celles contenues dans le premier, celles contenues dans le second et celles
qui traversent la jonction. L’additivité des sommes entières prouve
la formule. Conserver les identifiants et les prolongements
en attente évite de compter deux fois une trace ou d’en perdre une qui
n’est pas encore terminée. La réduction décimale est appliquée après le
dénombrement ; les quotients de cette réduction restent dans la mémoire.
\end{proof}

La réduction décimale de ces trois dénombrements produit
\begin{equation}
 z_m=100[A_m]_{10}+10[C_m]_{10}+[V_m]_{10},
 \qquad 0\le z_m<1000.
 \label{vi:dep:bloque-incidencial}
\end{equation}
Soient \((D_3z)_i=z_i-z_{i+3}\),
\((D_4z)_i=z_i-z_{i+4}\), avec des indices dans \(\ZZ/12\ZZ\), et
\(Q(z)=\sum_i z_i\). Le lecteur conserve
\(\mathcal D(z)=(D_3z,D_4z,Q(z))\). La réduction décimale oublie des
multiples de dix de chaque dénombrement, tandis que le registre original conserve
les visites et les transports qui les ont engendrés.
La lecture des événements de code pair sur cette même histoire,
sa décomposition \(12=1+5+6\) et son inversion entière sont
développées dans la proposition \ref{vi:prop:eventos}. Ce vecteur
d’événements ne remplace pas le présent registre d’incidences :
les domaines et les coordonnées conservées restent distincts.

\begin{proposition}[Conjugaison incidentielle et réduction décimale]
\label{vi:dep:conjugacion-incidencial}
Soit \(\varrho(m)=13-m\). Considérons une conjugaison du registre
qui envoie les visites et les traces de la fenêtre \(m\) sur celles de
\(\varrho(m)\) par des bijections qui conservent leurs multiplicités.
Supposons qu’elle conserve la classe arithmétique et l’étiquette de frontière
des visites, et échange les canaux positif et négatif des
traces. Alors les dénombrements de \eqref{vi:dep:incidencias} satisfont
\begin{equation}
 (A_m^\dagger,C_m^\dagger,V_m^\dagger)
 =(A_{\varrho(m)},-C_{\varrho(m)},V_{\varrho(m)}).
 \label{vi:dep:conjugacion-recuentos}
\end{equation}
En écrivant \(u_m=[C_m]_{10}\), leurs blocs décimaux vérifient
\begin{equation}
 z_m^\dagger=z_{\varrho(m)}
       +100\mathbf1_{\{u_{\varrho(m)}\ne0\}}-20u_{\varrho(m)}.
 \label{vi:dep:conjugacion-bloque}
\end{equation}
\end{proposition}
\begin{proof}
Les bijections conservent chaque terme de \(A_m\) et \(V_m\).
L’échange des canaux permute \(\nu^-\) et \(\nu^+\),
ce qui change le signe de \(C_m\). On obtient
\eqref{vi:dep:conjugacion-recuentos}. Si \(C=10q+u\),
\(0\le u<10\), la division euclidienne de \(-C\) donne
\[
 [-C]_{10}=10\mathbf1_{\{u\ne0\}}-u,
 \qquad
 \left\lfloor\frac{-C}{10}\right\rfloor
 =-q-\mathbf1_{\{u\ne0\}}.
\]
Dans \eqref{vi:dep:bloque-incidencial}, seul le chiffre
intermédiaire change. Sa différence est
\(10([-C]_{10}-[C]_{10})=100\mathbf1_{\{u\ne0\}}-20u\),
ce qui démontre \eqref{vi:dep:conjugacion-bloque}. La seconde
identité conserve simultanément le quotient entier de la réduction.
\end{proof}

La proposition précise l’action sur le registre avant de transporter
sa représentation décimale. Son application à l’inversion d’une route
conserve également la convention de lecture des extrémités. Soit
\(\gamma=(x_0,\ldots,x_n)\) une route enrichie et soit \(f\) une
observable entière évaluée avant chaque avancée. La route inverse
\(\gamma^{-1}=(x_n,\ldots,x_0)\) vérifie
\begin{equation}
 \sum_{k=0}^{n-1}f(x_{n-k})
 =\sum_{k=0}^{n-1}f(x_k)+f(x_n)-f(x_0).
 \label{vi:dep:conjugacion-frontera}
\end{equation}
En effet, le membre de gauche somme les indices \(1,\ldots,n\) ;
le premier terme de celui de droite somme \(0,\ldots,n-1\). Leur différence
est exactement le terme de frontière indiqué. L’identité s’applique
à chaque observable incidentielle avec ses données conservées. Les extrémités
sont comptabilisées dans les dénombrements entiers avant réduction modulo dix ;
sur une route fermée le terme s’annule, tandis que dans une fenêtre
ouverte il peut modifier son bloc. Les traces qui traversent une jonction
sont conservées au moyen de la partition de la
proposition~\ref{vi:dep:libro-actualizacion}.

Ainsi les dénombrements sont reliés à la conjugaison des routes de la
section~\ref{vi:sec:conjugacion}, en maintenant déclarée l’action
concrète sur chaque incidence. La négation d’un canal et l’inversion
d’une route sont des opérations sur le registre ; leur réduction décimale est la
représentation qui vient d’être calculée.

\begin{proposition}[Image entière et reconstruction du registre]
\label{vi:dep:imagen-integral}
Soient \(b,c\in\ZZ^{12}\), \(q\in\ZZ\), avec des indices de zéro à onze.
Définissons
\[
 w_i=c_i-b_{i+1},\qquad P_0=0,\qquad
 P_i=\sum_{j=0}^{i-1}w_j,\qquad T_w=\sum_{i=0}^{11}P_i.
\]
Il existe un unique \(k\in\ZZ^{12}\) tel que
\(D_3k=b\), \(D_4k=c\), \(Q(k)=q\), si et seulement si
\begin{equation}
 b_i=w_i+w_{i+1}+w_{i+2},\qquad
 \sum_iw_i=0,\qquad q+T_w\equiv0\pmod{12}.
 \label{vi:dep:condiciones-imagen}
\end{equation}
Dans ce cas,
\begin{equation}
 k_i=\frac{q+T_w}{12}-P_i.
 \label{vi:dep:inversa-incidencial}
\end{equation}
\end{proposition}
\begin{proof}
Si \(k\) existe, alors
\(c_i-b_{i+1}=k_i-k_{i+1}\). La somme cyclique de ces différences
est nulle, leur somme de longueur trois est \(b_i\), et
\(k_i=k_0-P_i\). En sommant les douze coordonnées,
\(q=12k_0-T_w\), ce qui prouve les conditions nécessaires et la formule.
Réciproquement, les conditions rendent la formule entière et
cyclique. Ses différences de longueur trois sont \(b_i\), et celles de
longueur quatre sont \(w_i+b_{i+1}=c_i\). Sa somme est \(q\).
La différence entre deux solutions a toutes ses différences
consécutives nulles et sa somme nulle ; elle est donc nulle.
\end{proof}

La reconstruction est une opération sur l’image typée, non une
sélection rétrospective de ses arguments. En particulier,
\(\mathcal D(z)=\mathcal D(z')\) implique \(z=z'\).
Comme l’expression d’un entier de \([0,999]\) en trois chiffres est unique,
deux triplets de dénombrements expriment le même registre si et seulement s’ils coïncident
composante par composante modulo dix.

Le changement de représentation signé utilise le bloc
\begin{equation}
 H_4=\begin{pmatrix}
 1&1&1&1\\1&1&-1&-1\\
 1&-1&1&-1\\1&-1&-1&1
 \end{pmatrix},\qquad H_4^2=4I_4.
 \label{vi:dep:hadamard}
\end{equation}
L’opérateur \(H_{12}\) applique \(H_4\) aux trois orbites ordonnées
\((1,4,7,10)\), \((2,5,8,11)\), \((3,6,9,12)\).
Sur le sous-réseau
\(\mathcal L_H=\{U\in\ZZ^{12}:H_{12}U\equiv0\pmod4\}\),
les applications
\begin{equation}
 U=H_{12}K,\qquad K=\tfrac14H_{12}U
 \label{vi:dep:K-reversible}
\end{equation}
sont inverses. La congruence exprime exactement le caractère entier de la
seconde application, et \(H_{12}^2=4I\) prouve les deux compositions.

Le registre signé est également calculé directement à partir des coordonnées
transversales, avant de reconstruire \(K\). Avec des indices de zéro à onze,
soient, pour \(r=0,1,2\),
\[
 S_r=\sum_{j=0}^3c_{r+3j},\quad
 a_0=\frac{q+2S_0+S_1}{3},\quad
 a_1=a_0-S_0,\quad a_2=a_1-S_1,
 \qquad d_{r,j}=b_{r+3j}.
\]
L’application harmonique exprime par orbites
\begin{equation}
 \Pi_H^{(r)}(b,c,q)=
 (a_r,d_{r,1}-d_{r,3},d_{r,0}+d_{r,2},d_{r,0}-d_{r,2}).
 \label{vi:dep:PiH}
\end{equation}
Pour prouver la formule sur l’image entière, soit \(k\) son
antécédent. Si \(s_r=\sum_jk_{r+3j}\), alors
\(S_r=s_r-s_{r+1}\) et \(q=s_0+s_1+s_2\), avec des indices d’orbite
modulo trois. La résolution donne \(s_r=a_r\). Pour les quatre coordonnées
\((x_0,x_1,x_2,x_3)\) d’une orbite,
\(d_j=x_j-x_{j+1}\). Leurs trois combinaisons dans
\eqref{vi:dep:PiH} sont respectivement
\(x_0+x_1-x_2-x_3\), \(x_0-x_1+x_2-x_3\) et
\(x_0-x_1-x_2+x_3\). Avec \(a_r=\sum_jx_j\), elles sont
exactement \(H_4(x_0,x_1,x_2,x_3)^{\mathsf T}\).
La composition effective est ainsi
\(\mathscr L\to(b,c,q)\to\Pi_H(b,c,q)=U\to H_{12}U/4=K\).
L’antécédent utilisé dans cette démonstration vérifie l’identité ;
ce n’est pas un argument de l’application \(\Pi_H\).

Pour l’état terminal de \eqref{vi:dep:estado-terminal-108}, la coordonnée
produite avant la représentation signée est
\begin{align}
 b^{90}={}&(-495,-116,-684,108,601,-90,
              -173,-88,313,560,-397,461),\notag\\
 b^{120}={}&(-425,-281,-481,671,-255,30,
               475,-543,680,251,6,-128),\notag\\
 Q_{\rm TPK}={}&6263.
 \label{vi:dep:canales-terminales}
\end{align}
En substituant ces vingt-cinq champs dans \eqref{vi:dep:PiH}, les trois
orbites signées sont
\begin{align*}
 U_{\rm term}^{(1)}&=(2378,-452,-668,-322),\\
 U_{\rm term}^{(2)}&=(1406,998,-204,-28),\\
 U_{\rm term}^{(3)}&=(2479,-551,-371,-997).
\end{align*}
L’entrelacement par colonnes produit
\begin{equation}
 U_{\rm term}=\mathcal S_{12}(X_{\rm term})
 =(2378,1406,2479,-452,998,-551,-668,-204,-371,-322,-28,-997).
 \label{vi:dep:U-terminal}
\end{equation}
La transformée orbitale \(H_{12}/4\) est définie sur ce sous-réseau par
la divisibilité exacte de chaque bloc de quatre. Par conséquent, la chaîne
\begin{equation}
 X_{\rm term}\xrightarrow{\mathcal R_{12}}\mathscr L(X_{\rm term})
 \xrightarrow{\mathcal E_{108}^{90,120}}C_{\rm term}
 \xrightarrow{\Pi_H}U_{\rm term}
 \xrightarrow{H_{12}/4}K
 \label{vi:dep:cadena-terminal-completa}
\end{equation}
est une composition typée allant de l’état plein à son sceau périodique.
Les propositions~\ref{vi:dep:libro-actualizacion} et
\ref{vi:dep:imagen-integral} fournissent, respectivement, la mise à jour
finie du registre et l’inverse entière de la dernière représentation.

Le registre engendré fournit le témoin ordonné
\begin{equation}
 K=(234,543,140,729,659,824,621,058,914,794,146,601).
 \label{vi:dep:K-testigo}
\end{equation}
Le zéro initial de \(058\) fait partie de la lecture par blocs.
Pour le vérifier au moyen de la proposition~\ref{vi:dep:imagen-integral}, ses différences
et sa somme sont
\[
\begin{split}
b={}&(-495,-116,-684,108,601,-90,-173,-88,313,560,-397,461),\\
c={}&(-425,-281,-481,671,-255,30,475,-543,680,251,6,-128),\\
q={}&6263.
\end{split}
\]
Dans les trois orbites précédentes, le registre signé est
\[
 \begin{aligned}
 U^{(1)}&=(2378,-452,-668,-322),\\
 U^{(2)}&=(1406,998,-204,-28),\\
 U^{(3)}&=(2479,-551,-371,-997).
 \end{aligned}
\]
Appliquer \(H_4/4\) à ces trois vecteurs reproduit
\eqref{vi:dep:K-testigo}. Ces identités vérifient la représentation des
incidences ; leur genèse reste celle du registre
\eqref{vi:dep:incidencias}.

\subsection{\texorpdfstring{Représentation périodique de \(K\) et normalisation de \(\alpha\)}{Représentation périodique de K et normalisation d’alpha}}
\label{vi:dep:alpha}

Fixons \(B=1000\), l’origine de phase et les douze positions. Écrivons
\[
 N_K=\sum_{j=1}^{12}K_jB^{12-j},\qquad
 \kappa=\frac{N_K}{B^{12}-1}.
\]
C’est l’évaluation périodique exacte du registre. La multiplication
par \(B^{12}-1\) et la division euclidienne successive récupèrent ses douze
blocs, y compris tout zéro initial. Pour
\(\rho K=(K_2,\ldots,K_{12},K_1)\), la même concaténation donne
\(\kappa(\rho K)=B\kappa(K)-K_1\). Par conséquent, la réversibilité
du registre conserve une origine de phase spécifiée.

\begin{proposition}[Période minimale de la représentation rationnelle]
\label{vi:dep:periodo-minimo-k}
Le prolongement périodique du registre \eqref{vi:dep:K-testigo}
a une période minimale de \(12\) blocs en base \(1000\) et de
\(36\) chiffres en base \(10\).
\end{proposition}
\begin{proof}
Les douze composantes de \(K\) sont distinctes. Toute translation
cyclique non nulle inférieure à douze identifierait deux composantes si elle était
une période. Par conséquent, la période minimale des blocs est douze.

La concaténation de leurs blocs de trois chiffres, y compris \(058\),
produit une période décimale de longueur \(36\). Ce développement est
canonique : il n’est pas ultimement constant au chiffre \(9\).
Soit \(d\) sa période décimale minimale. Alors \(d\mid36\),
car les décalages qui préservent un mot périodique forment
un sous-groupe cyclique. En regroupant les chiffres trois par trois, la période
minimale des blocs est le plus petit \(k>0\) tel que \(d\mid3k\), c’est-
à-dire \(d/\gcd(d,3)\). Aucun diviseur propre de \(36\)
ne satisfait \(d/\gcd(d,3)=12\). Il s’ensuit que \(d=36\).
\end{proof}

Cette période appartient à la coordonnée rationnelle du registre, avec
l’origine de phase fixée. La composition suivante conserve en outre les
coordonnées régionales et les quotients de normalisation de sa sortie.

Soient \(P_j,E_j,\Phi_j\) les blocs fractionnaires des trois
caractères régionaux déjà construits, et prolongeons \(K_j\) avec une période de
douze. La composition des registres est
\begin{equation}
 Z_j=P_j+E_j-\Phi_j-K_j,
 \qquad
 A_j=Z_j+c_{j+1}-Bc_j.
 \label{vi:dep:normalizacion-alpha}
\end{equation}
Le quotient entier \(c_j\) transporte la continuation vers la position
observée. Dans une fenêtre finie, on fixe d’abord \(c_{N+1}\), puis
\[
 c_j=\left\lfloor\frac{Z_j+c_{j+1}}B\right\rfloor,
 \qquad A_j=Z_j+c_{j+1}-Bc_j.
\]
La division euclidienne détermine de manière unique \(0\le A_j<B\).
Comme \(-1998\le Z_j\le1998\), l’ensemble
\(\{-2,-1,0,1,2\}\) est stable pour ces quotients.

\begin{theorem}[Composition exacte à profondeur illimitée]
\label{vi:dep:alpha-completa}
L’évaluation complète de la normalisation, dans la convention décimale
canonique, est
\begin{equation}
 \alpha_A=\pi+e-\varphi-4-\kappa.
 \label{vi:dep:alpha-A}
\end{equation}
Il existe un unique développement canonique \((A_j)\), avec des blocs entre zéro
et \(B-1\), et une suite bornée de quotients entiers qui satisfont
\eqref{vi:dep:normalizacion-alpha} et \(c_1=0\).
\end{theorem}
\begin{proof}
Les parties entières régionales sont \(3,2,1\). Par convergence absolue,
\[
 y:=\sum_{j\ge1}Z_jB^{-j}
   =(\pi-3)+(e-2)-(\varphi-1)-\kappa.
\]
Les premiers blocs sont \(141,718,618,234\), de sorte que
\(Z_1=7\). La somme de la queue a une valeur absolue au plus
égale à \(1998\sum_{j\ge2}B^{-j}=0.002\) ; en particulier \(0<y<1\).
Prenons le développement canonique \(y=\sum A_jB^{-j}\), en choisissant le
développement terminal et non celui ultimement constant \(B-1\) si une
frontière se présente. Posons
\[
 c_j=\sum_{r\ge0}(Z_{j+r}-A_{j+r})B^{-r-1}.
\]
L’égalité des valeurs complètes montre que chaque \(c_j\) est
l’opposé d’une somme finie d’entiers pondérés par des puissances de
\(B\) ; il est donc entier, et \(c_1=0\). Séparer le premier terme
donne \(Bc_j=Z_j-A_j+c_{j+1}\). Les bornes des queues rendent la
suite bornée. Plus précisément, la queue de \(Z\) est dans \([-2,2]\)
et la queue canonique de \(A\) dans \([0,1)\), de sorte que
\(-3<c_j\le2\). Étant entier, \(c_j\) appartient à l’ensemble
\(\{-2,-1,0,1,2\}\) indiqué.
Réciproquement, la somme de la récurrence jusqu’à \(N\) donne
\[
 \sum_{j=1}^N A_jB^{-j}
 =\sum_{j=1}^N Z_jB^{-j}-c_1+c_{N+1}B^{-N}.
\]
Le caractère borné annule le dernier terme à la limite. L’unicité du
développement canonique de \(y\) détermine \(A\), et la formule des queues
détermine alors tous les \(c_j\).
\end{proof}

L’évaluation finie \(\kappa_{36}=N_K/B^{12}\) diffère de
\(\kappa\) de \(B^{-12}\kappa\). En outre, imposer
\(c_{13}=0\) dans une fenêtre de douze blocs ne remplace pas le quotient
de la continuation complète. Dans le registre exprimé, la continuation
donne \(c_{13}=-1\) : le douzième bloc est \(662\), tandis que la
fenêtre tronquée avec un quotient terminal nul donne \(663\). L’évaluation
complète commence par
\[
 \alpha_A=0.007297352569283800997285105472380662999\ldots.
\]
La période de \(K\) n’impose pas de période à \((A_j)\), car dans
\eqref{vi:dep:normalizacion-alpha} interviennent également les trois
registres régionaux et les quotients entiers.

Dans les formules suivantes, \(\alpha\) désigne cette sortie positionnelle
complète \(\alpha_A\). La voie analytique des vacances possède son opérateur
propre ; un polynôme tronqué d’ordre neuf et la racine de ce polynôme
ne sont pas identifiés ici à l’opérateur complet. Les compositions de
masses qui suivent ne nécessitent pas de promouvoir cette approximation en
une égalité entre les deux voies infinies.

\subsection{Construction tridimensionnelle du préfacteur électronique}
\label{vi:dep:prefactor}

Soit \(\mathcal A_3=\{1,\ldots,9\}^3\), muni de la mesure uniforme, et
\(\rho_9(n)=1+((n-1)\bmod9)\). Sur ce domaine agissent
\[
 \Sigma(x)=\rho_9(x_1+x_2+x_3),\qquad
 \Pi(x)=\rho_9(x_1x_2x_3).
\]
Dans \(\ell^2(\mathcal A_3)\), soient \(P\) et \(Q\) les moyennes
conditionnelles sur les fibres de \(\Sigma\) et \(\Pi\),
respectivement. Chacune est idempotente et autoadjointe : dans chaque
fibre, elle remplace un vecteur par sa composante constante et conserve
orthogonalement cette composante. Ce sont donc des projecteurs orthogonaux.

Les fibres additives ont pour cardinal \(81\). Les multiplicatives ont
pour cardinaux
\[
m=(36,36,108,36,36,108,36,36,297).
\]
L’incidence conjointe
\(N_{ab}=\#\{x:\Sigma(x)=a,\Pi(x)=b\}\) est \(9N_0\), où
\begin{equation}
 N_0=\begin{pmatrix}
0&1&1&0&1&1&0&1&4\\1&0&1&1&0&1&1&0&4\\
1&0&2&0&1&2&0&0&3\\0&1&1&0&1&1&0&1&4\\
1&0&1&1&0&1&1&0&4\\0&0&2&1&0&2&0&1&3\\
0&1&1&0&1&1&0&1&4\\1&0&1&1&0&1&1&0&4\\
0&1&2&0&0&2&1&0&3
\end{pmatrix}.
 \label{vi:dep:incidencia-cubica}
\end{equation}
Cette matrice est obtenue en dénombrant les \(729\) triplets. Elle est également
reproduite, sans arrondi, avec
\(N_{ab}=\sum_{i,j,k=1}^9
\mathbf1_{\{a\}}(\rho_9(i+j+k))
\mathbf1_{\{b\}}(\rho_9(ijk))\).
Dans les bases normalisées de fonctions constantes sur chaque fibre,
l’entrelacement a pour matrice \(C_{ab}=N_{ab}/\sqrt{81m_b}\).

\begin{proposition}[Norme de la non-commutativité des deux feuillets]
\label{vi:dep:norma-PQ}
Les projecteurs précédents satisfont
\begin{equation}
 \|[P,Q]\|=\frac{\sqrt3}{4}.
 \label{vi:dep:prefactor-valor}
\end{equation}
\end{proposition}
\begin{proof}
Dans l’espace des neuf fibres additives, \(PQP\) a pour matrice
\(M=CC^{\mathsf T}\). La multiplication des matrices de
\eqref{vi:dep:incidencia-cubica} donne la valeur propre un sur le vecteur
constant, \(1/4\) sur \((1,-1,0,1,-1,0,1,-1,0)\), et
\(7/132\) sur \((1,1,-2,1,1,-2,1,1,-2)\).
Sur les coordonnées \(\{3,6,9\}\) de somme nulle, il agit par
\(1/18\), avec une multiplicité de deux. Sur chacun des ensembles
\(\{1,4,7\}\) et \(\{2,5,8\}\), les vecteurs de somme nulle
appartiennent au noyau ; ils donnent quatre dimensions. Ces sous-espaces
orthogonaux totalisent les neuf dimensions.

Pour un vecteur propre unitaire \(u\in\operatorname{ran}P\) de
\(PQP\) de valeur propre \(0<\lambda<1\), le vecteur
\(v=(Qu-\lambda u)/\sqrt{\lambda(1-\lambda)}\) est unitaire et
orthogonal à \(\operatorname{ran}P\). Dans la base \((u,v)\),
\[
 P=\begin{pmatrix}1&0\\0&0\end{pmatrix},\quad
 Q=\begin{pmatrix}\lambda&\sqrt{\lambda(1-\lambda)}\\
 \sqrt{\lambda(1-\lambda)}&1-\lambda\end{pmatrix}.
\]
Le commutateur a pour norme \(\sqrt{\lambda(1-\lambda)}\).
Les blocs correspondant à \(\lambda=0,1\) commutent. Parmi les
valeurs propres calculées, le maximum est atteint en \(\lambda=1/4\),
car \(\lambda(1-\lambda)\) est croissante sur \([0,1/2]\).
Sa racine carrée est \(\sqrt3/4\).
\end{proof}

Dans le bloc extrémal, \(S=2P-I\) et
\(J=4[P,Q]/\sqrt3\) satisfont
\[
 S^2=I,\qquad J^2=-I,\qquad SJ=-JS.
\]
Les quatre matrices \(I,S,J,SJ\) sont linéairement indépendantes,
de sorte qu’elles réalisent \(\mathrm{Cl}_{1,1}\simeq M_2(\RR)\).
Les éléments \(n_\pm=(S\pm J)/2\) vérifient
\(n_\pm^2=0\) et \(n_+n_-+n_-n_+=I\), par développement direct.
Le préfacteur de l’équation électronique provient ainsi de deux évaluations
d’APP en dimension trois et de leur commutateur, avant toute masse
de référence.

\subsection{Lecture de frontière et composition électronique basale}
\label{vi:dep:basal}

L’incidence \eqref{vi:dep:incidencia-hojas} possède un vecteur propre
positif \(r=(1,\varphi)^{\mathsf T}\). La normalisation
\(p_{ij}=F_{ij}r_j/(\varphi r_i)\) produit
\[
 P_{\rm av}=\begin{pmatrix}0&1\\\varphi^{-2}&\varphi^{-1}\end{pmatrix},
 \qquad \mu=\frac{(1,\varphi^2)}{1+\varphi^2}.
\]
L’identité \(\varphi^{-2}+\varphi^{-1}=1\) prouve que les lignes
ont pour somme un. Comme \(F\) est symétrique, les poids proportionnels à
\(r_i^2\) satisfont l’équilibre détaillé ; ainsi \(\mu P_{\rm av}=\mu\).
L’incidence est irréductible et comporte une autorétention, de sorte que
cette distribution stationnaire est unique. La fréquence chronologique
\((1/3,2/3)\) compte les trois instants du calendrier ; \(\mu\)
pondère ses trajectoires admissibles de mémoire un.

L’opérateur qui conserve les feuillets, fixe la normalisation volumétrique
à un et marque l’aréale par le caractère de clôture est
\(D_\pi=\operatorname{diag}(\pi,1)\). Par récurrence,
\[
 F^n=\begin{pmatrix}F_{n-1}&F_n\\F_n&F_{n+1}\end{pmatrix}.
\]
Le quotient des retours diagonaux marqués est
\[
 \Theta_n=\pi\frac{F_{n-1}}{F_{n+1}}
 \longrightarrow\frac{\pi}{\varphi^2}.
\]
La convergence se déduit de la valeur propre dominante \(\varphi\), ou des
quotients consécutifs déjà bornés. Le lecteur
\(\mathcal R(x,y)=x/y^2\), sous l’échange des arguments,
exprime également
\begin{equation}
 \rho_A=\frac{\pi}{\varphi^2},\qquad
 \rho_E=\frac{\varphi}{\pi^2},\qquad
 \varphi^3\rho_A^2\rho_E=1.
 \label{vi:dep:dos-orientaciones}
\end{equation}
La dernière égalité se vérifie en multipliant les fractions.
Elle relie deux orientations du même lecteur et ne sélectionne pas les
autres coefficients de l’équation électronique.

Le retour central typé a pour longueur \(27\), avec un support
\(\ZZ/9\ZZ\times\FF_3\). Les cinq positions de la microfibre
s’insèrent, dans la jauge régionale fixée, comme
\((4,0),(5,0),(6,0),(6,1),(6,2)\).
Le complément possède \(27-5=22\) positions. Son incidence avec les
trois états de phase tritique comporte \(5\cdot3=15\) paires.
On obtient ainsi les entiers \(-22\) et \(15\) utilisés par le
lecteur électronique ; l’orientation fixe le signe du premier terme.
La composition cubique utilise trois copies de la coordonnée de fermeture,
de sorte que son évaluation est \(\alpha^3\). Le support, l’insertion
régionale et cette règle de composition sont des données opératoires distinctes ;
un dénombrement isolé ne remplace pas leur déclaration.

Le défaut global et son retour nonadique s’expriment par
\begin{equation}
 d_4=\frac{\pi-e\log\pi}{270},\qquad
 \Delta_4=d_4\left(1+\frac{\pi}{729}\right).
 \label{vi:dep:delta4}
\end{equation}
Le dénominateur \(270\) appartient à la couronne privée de son unité de
clôture ; \(729=9^3\) au support cubique. La fonction
\(f(x)=x-e\log x\) atteint son minimum nul en \(x=e\), car
\(f'(x)=1-e/x\) y change de signe. Comme \(\pi\ne e\),
on obtient \(d_4>0\) et \(\Delta_4>0\).

\begin{definition}[Coordonnée électronique basale]
\label{vi:dep:electron-basal}
La composition des lecteurs précédents est
\begin{equation}
 \varepsilon_e^{(0)}=
 \frac{\sqrt3}{4}
 \left[\exp\left(\frac{\varphi}{\pi^2}\right)-22\alpha^3\right]
 (1+15\Delta_4).
 \label{vi:dep:formula-basal}
\end{equation}
\end{definition}
L’exponentielle agit sur l’orientation \(\rho_E\) de
\eqref{vi:dep:dos-orientaciones} ; la lecture de vacance corrige cette
composante et le dernier facteur conserve le retour du défaut global.
Pour \(0<\alpha<0.01\), tous les facteurs sont positifs :
\(\exp(\varphi/\pi^2)>1\), \(22\alpha^3<22\cdot10^{-6}<1\)
et \(\Delta_4>0\). C’est une coordonnée adimensionnelle. Sa présence
dans une coordinatisation de masse est définie après la construction de l’échelle.

\subsection{Échelle déterminantale et constante de Planck réduite}
\label{vi:dep:planck}

Le bloc local \eqref{vi:dep:hadamard} a pour déterminant \(-16\).
Sa forme normale de Smith est \(\operatorname{diag}(1,2,2,4)\).
En effet, le plus grand diviseur des entrées est un ; les mineurs
d’ordre deux sont pairs et l’un vaut deux en valeur absolue. Les mineurs
d’ordre trois sont, par l’identité de la comatrice, de valeur absolue quatre,
et le déterminant a pour valeur absolue seize. Les diviseurs
déterminantaux sont \(1,2,4,16\), dont les rapports successifs donnent les
quatre facteurs de Smith. Par conséquent,
\begin{equation}
 G_H=\operatorname{coker}H_4
 \simeq\ZZ/2\ZZ\oplus\ZZ/2\ZZ\oplus\ZZ/4\ZZ,
 \qquad |G_H|=16.
 \label{vi:dep:conucleo}
\end{equation}

Le lecteur normé local multiplie la même coordonnée \(\alpha\)
une fois pour chaque classe de ce conoyau et applique une fois la coordonnée
d’auto-échelle. Sa définition produit
\begin{equation}
 \Sigma_H=\varphi\prod_{g\in G_H}\alpha
         =\varphi\alpha^{16}.
 \label{vi:dep:sigma-H}
\end{equation}
L’ordre du conoyau détermine l’exposant de ce lecteur. La forme
normale de Smith et le choix du lecteur normé remplissent des fonctions
distinctes : la première ne transforme pas n’importe quelle fonctionnelle du bloc en
l’expression \eqref{vi:dep:sigma-H}.

\begin{proposition}[Décade d’action]
\label{vi:dep:decada}
La sortie \eqref{vi:dep:sigma-H} satisfait
\(10^{-34}<\Sigma_H<10^{-33}\) ; sa valuation décimale est \(-34\).
\end{proposition}
\begin{proof}
Les évaluations régionales et positionnelles donnent les intervalles
\(1.618<\varphi<1.619\) et \(0.007297<\alpha<0.007298\).
Par positivité,
\[
 \frac{1618\,7297^{16}}{10^{99}}
 <\Sigma_H<
 \frac{1619\,7298^{16}}{10^{99}}.
\]
Les inégalités entières
\(1618\,7297^{16}>10^{65}\) et
\(1619\,7298^{16}<10^{66}\) prouvent le résultat sans utiliser de
constante d’action externe.
\end{proof}

La décade et la mantisse sont des représentations d’étapes différentes.
Avant de construire la seconde, le lecteur régional sélectionne le retour
persistant de la microfibre de clôture. La projection
\(p_{\rm ret}(u_1,\ldots,u_6)=(u_4,u_5,u_6)\) appliquée aux
cinq émissions de \ref{vi:dep:regiones} produit le bloc
\(b_\pi=(169,272,272)\) quatre fois et
\(b_\pi^-=(418,674,521)\) une fois. Ces multiplicités dénombrent des
lignes régionales, non les multiplicités des germes \(432,144\).

\begin{proposition}[Sélecteur équivariant du retour régional]
\label{vi:dep:selector169}
Parmi les règles qui conservent les permutations des cinq régions,
sélectionnent le bloc terminal de multiplicité maximale, sont équivariantes
sous les permutations de ses trois positions et retiennent une position fixe
sous le stabilisateur du bloc, la coordonnée sélectionnée est unique
et vaut \(\rho_\pi=169\).
\end{proposition}
\begin{proof}
La projection des cinq lignes montre que le bloc de multiplicité
maximale est unique et est \(b_\pi\). Son stabilisateur dans
\(\mathfrak S_3\) échange les deux positions de valeur \(272\)
et fixe uniquement la position de valeur \(169\). Retenir cette position
est compatible avec toute permutation simultanée des coordonnées.
De manière équivalente, la valeur s’exprime sans choisir de position initiale
comme
\(\operatorname{sing}(b_\pi)=\sum_j(b_\pi)_j
-2\operatorname{moda}(b_\pi)=169\).
\end{proof}

La graduation par longueur exprime un mot de longueur \(n\)
au moyen du caractère \(\chi_z([w])=z^{|w|}\). La concaténation
additionne les longueurs, de sorte que
\(\chi_z([uv])=\chi_z([u])\chi_z([v])\). Comme les émissions
régionales appartiennent à \(\mathcal U_6\), l’impulsion du retour
apparaît comme \(\rho_\pi z^6\). Le degré six correspond à la
longueur de l’émission, non à la longueur d’une marche fermée qui
parcourt plusieurs ancrages du graphe.

Le poids de prolongement s’obtient dans la droite déterminant du
transport bicouche. Pour une involution réelle \(R\) en dimension deux,
avec \(\operatorname{tr}R=0\), soient
\[
 x_A=\frac{50\pi}{9}\alpha,\qquad
 \mathcal T(x_A,y)=\exp(-x_AI_2-yR).
\]
La coordonnée orientée \(y\) est encore formelle. Dans une base propre
de \(R\), les deux entrées sont \(e^{-x_A-y}\) et
\(e^{-x_A+y}\). Par conséquent,
\[
 \bigwedge^2\mathcal T=\det\mathcal T
 =e^{-2x_A}=D_A.
\]
La dépendance en \(y\) s’annule. La multiplicativité par
concaténation détermine le caractère de prolongement \(k\mapsto D_A^k\) :
tout caractère de valeur élémentaire \(D_A\) doit avoir cette valeur
en \(k=1+\cdots+1\). Cette construction précède l’évaluation
de la coordonnée angulaire conjuguée.

Le lecteur régional d’ordre cinq emploie les valeurs propres \(9,5\)
du bloc central \(\bigl(\begin{smallmatrix}7&2\\2&7\end{smallmatrix}\bigr)\),
l’orbite de quatre positions et le calendrier de douze fenêtres.
Sa règle graduée fixe
\begin{align}
 H_5={}&\frac12\sqrt{\frac{1000\alpha}{\varphi}}-\alpha
       +\frac9{16}\alpha^2-\frac59\alpha^3
       +\frac7{48}\alpha^4-\frac1{54}\alpha^5,\notag\\
 D_A={}&\exp\left(-\frac{100\pi\alpha}{9}\right),\notag\\
 \mathcal C_\pi={}&169\alpha^6+
                  \frac{D_A\alpha^7}{1-D_A\alpha},\notag\\
 \eta_{\rm ret}={}&H_5-\frac{90}{\pi}\mathcal C_\pi.
 \label{vi:dep:mantisas}
\end{align}
Dans cette règle, les quotients sont
\(9/16=9/4^2\), \(5/9\), \(7/48=(\operatorname{tr}B_c/2)/(4\cdot12)\)
et \(1/54=1944/104976\). Leur affectation aux degrés et les signes
font partie du lecteur gradué. L’entier \(169\) est la représentation
du sélecteur démontré dans la proposition \ref{vi:dep:selector169}.
Sa donnée régionale est antérieure à l’évaluation de \(\alpha\) ; le
lecteur analytique la compose ensuite avec la coordonnée de fermeture.

La queue de \(\mathcal C_\pi\) possède une construction par renouvellement.
L’impulsion \(\rho_\pi z^6\) et la continuation stricte appartiennent
à des termes distincts. Après avoir enregistré une seule fois l’impulsion,
la continuation est normalisée par l’unité de la droite déterminant ;
chaque nouveau pas la multiplie par \(D_A\). L’orientation de la
jauge régionale place ce retour dans le secteur compensateur qui
est soustrait du caractère de cinquième degré.
Avec un paramètre formel \(z\), la règle
\begin{equation}
 R(z)=D_Az^7+D_AzR(z)
 \label{vi:dep:renovacion}
\end{equation}
détermine de manière unique
\(R(z)=D_Az^7/(1-D_Az)\). Cela peut être prouvé par comparaison des
coefficients ou par inversion de la série formelle \(1-D_Az\).
La normalisation du premier retour, outre le rapport récurrent,
est essentielle à cette unicité. Comme \(0<D_A<1\) et \(\alpha<1\),
l’évaluation en \(z=\alpha\) converge absolument.

\begin{proposition}[Positivité et rapport de retour d’action]
\label{vi:dep:positividad-accion}
Sur la branche engendrée, \(0<\eta_{\rm ret}<H_5\). Le rapport
\begin{equation}
 R_{\rm act}=\frac{H_5}{\eta_{\rm ret}}
 \label{vi:dep:Ract}
\end{equation}
est positif, supérieur à un et adimensionnel.
\end{proposition}
\begin{proof}
Les bornes \(0.007<\alpha<0.008\) et \(\varphi<1.619\) donnent
\(\tfrac12\sqrt{1000\alpha/\varphi}>1\). La somme des valeurs absolues
des termes négatifs restants de \(H_5\) est inférieure à
\(0.008001\), de sorte que \(H_5>0.99\).
Comme \(D_A<1\), \(\pi>3\) et \(\alpha<0.008\),
\[
 0<\frac{90}{\pi}\mathcal C_\pi
 <30\left(169(0.008)^6+
                  \frac{(0.008)^7}{1-0.008}\right)<10^{-6}.
\]
On en déduit \(0<\eta_{\rm ret}<H_5\). Le rapport de deux coordonnées
d’action dans la même base est adimensionnel et ne change pas lors d’un changement d’échelle
de cette base.
\end{proof}

Soit \(\mathcal U_S\) une base de la droite d’action. Les deux
sections, antérieure et postérieure au retour, sont
\begin{equation}
 \hbar_{\rm pre}=H_5\,10^{-34}\mathcal U_S,
 \qquad
 \hbar_{\rm ret}=\eta_{\rm ret}\,10^{-34}\mathcal U_S.
 \label{vi:dep:hbar}
\end{equation}
Sa lecture physique est la construction proposée de la constante de Planck
réduite, avec la section de retour explicitée. L’unité
\(\mathcal U_S\mapsto\mathrm{J\,s}\) est une réalisation métrologique
postérieure. Le coefficient de \(\Sigma_H\) dans sa décade n’est pas
\(H_5\) : le déterminant exprime l’échelle et le lecteur régional
exprime la coordonnée qui occupe cette échelle. Pour chaque section,
\(h=2\pi\hbar\) exprime l’action d’un tour complet.

\subsection{Coordonnées angulaires et horloge de la réalisation absolue}
\label{vi:dep:reloj}

La paire angulaire utilise les coordonnées adimensionnelles déjà engendrées :
\begin{equation}
 A_{\deg}=1000\alpha,\qquad
 C^*_{\deg}=2(\eta_{\rm ret}+\alpha),\qquad
 x=\frac{\pi}{180}A_{\deg},\quad
 y=\frac{\pi}{180}C^*_{\deg}.
 \label{vi:dep:angulos}
\end{equation}
La première opération décale d’un bloc le développement en base mille.
La seconde est la représentation bidirectionnelle de la coordonnée régionale :
de manière équivalente,
\(y=\pi(H_5+\alpha)/90-\mathcal C_\pi\).
Cette égalité se déduit en substituant \eqref{vi:dep:mantisas} ; elle
n’introduit pas de règle supplémentaire de sélection du contre-angle.
La coordinatisation d’un tour comporte \(360\) unités angulaires ou \(2\pi\)
radians. Ainsi, pour une coordonnée relevée,
\[
 \hbar\theta_{\rm rad}
 =h\frac{\theta_{\deg}}{360}.
\]
Le facteur \(10^{-34}\) reste dans la droite d’action et ne
s’insère pas dans les coordonnées angulaires.

Les canaux positifs sont
\[
 q_+=e^{-x-y},\qquad q_-=e^{-x+y}.
\]
On récupère la paire par
\(x=-\tfrac12\log(q_+q_-)\),
\(y=\tfrac12\log(q_-/q_+)\). L’échange des canaux conserve
\(x\) et change le signe de \(y\). L’opérateur muni des projecteurs
orientés \(P_\pm\),
\(T=q_+P_++q_-P_-\) conserve ces étiquettes en exprimant ses deux
valeurs propres.

Pour une unité temporelle interne \(t_0>0\), la fermeture de \(108\)
pas détermine
\begin{equation}
 \omega_0=\frac{2\pi}{108t_0},\qquad
 E_{\rm clk}=\hbar_{\rm ret}\omega_0,
 \qquad m_{\rm clk}=\frac{E_{\rm clk}}{c_{\rm int}^{2}}.
 \label{vi:dep:reloj-absoluto}
\end{equation}
Ces formules produisent des dimensions de fréquence, d’énergie et de masse.
La vitesse et la longueur s’obtiennent par la réponse des
mêmes canaux, comme le précisent les identités suivantes. Les exprimer
en secondes, en joules ou en unités de masse nécessite les coordinatisations
correspondantes ; ces coordinatisations ne modifient ni les canaux ni la section
électronique qui est normalisée par rapport à l’horloge.

\subsubsection{Identité entre vitesse constitutive et vitesse de l’horloge}

Dans le domaine \(x>|y|\), les deux canaux satisfont \(0<q_\pm<1\).
Sur la réalisation à deux feuillets, munie de projecteurs de rang un
\(P_+,P_-\), on définit la réponse temporelle
\[
 \mathsf V=(I-T^{90})(I-T^{120})^{-1}
          =r_+P_++r_-P_-,
 \qquad r_\pm=\frac{1-q_\pm^{90}}{1-q_\pm^{120}}>0.
\]
L’inverse existe car les deux valeurs propres de \(T^{120}\) sont inférieures
à un. La lecture symétrique de la réponse est
\[
 \mathcal E=
 \log[(1-q_+^{90})(1-q_-^{90})]
 -\log[(1-q_+^{120})(1-q_-^{120})]
 =\log(r_+r_-).
\]
Soient \(u_S,u_C,u_T\) les bases positives d’action, de vitesse et de
temps, et \(u_L=u_Cu_T\) la base de longueur. Les lectures
cinématique et constitutive sont
\[
 c_{\rm int}=e^{-\mathcal E}u_C,\qquad
 c_{\rm vac}=(r_+r_-)^{-1}u_C.
\]

\begin{proposition}[Réalisation commune de propagation et de masse absolue]
\label{vi:dep:velocidad-comun}
Dans une même coordinatisation dimensionnelle,
\[
 c_{\rm int}=c_{\rm vac}=\widehat c\,u_C,
 \qquad \widehat c=(\det\mathsf V)^{-1}=(r_+r_-)^{-1}.
\]
Pour \(t_0=\tau_0u_T\), avec \(\tau_0>0\), la longueur élémentaire est
\[
 \ell_0=c_{\rm int}t_0=\tau_0\widehat c\,u_L.
\]
La masse chronologique de \eqref{vi:dep:reloj-absoluto} conserve, par
conséquent, la forme
\[
 m_{\rm clk}=
 \frac{2\pi\,\eta_{\rm ret}\,10^{-34}}
      {108\,\tau_0\,\widehat c^{\,2}}\,
 \frac{\mathcal U_S}{u_Tu_C^2}.
\]
\end{proposition}
\begin{proof}
L’identité \(\mathcal E=\log(r_+r_-)\) donne
\(e^{-\mathcal E}=(r_+r_-)^{-1}\) ; le rang un de chaque feuillet donne
\(\det\mathsf V=r_+r_-\). L’échange des feuillets permute les
facteurs et conserve leur produit. La longueur s’obtient en multipliant
la même vitesse par \(t_0\). Substituer ces expressions et
\eqref{vi:dep:hbar} dans \eqref{vi:dep:reloj-absoluto} produit la
formule massique. Le facteur \(\mathcal U_S/(u_Tu_C^2)\) a la
dimension d’une masse ; il n’intervient pas dans la sélection des canaux.
\end{proof}

Pour une route homogène de \(n>0\) transitions, les lectures additives
\(T(\gamma)=nt_0\) et \(L(\gamma)=n\ell_0\) donnent
\(L(\gamma)/T(\gamma)=c_{\rm vac}\). Si le canal varie selon
l’arête, on additionne leurs longueurs ; le quotient est alors la vitesse
moyenne de ces arêtes, pas nécessairement une vitesse locale commune.

L’égalité des sections demeure sous les changements de coordinatisation. Si
\(u_C'=du_C\), \(u_T'=\eta u_T\), avec \(d,\eta>0\), alors
\[
 \widehat c'=\widehat c/d,\qquad
 \tau_0'=\tau_0/\eta,\qquad u_L'=d\eta u_L,
 \qquad
 \tau_0'\widehat c'u_L'=\ell_0.
\]
En particulier, choisir la base adaptée \(u_C'=c_{\rm vac}\)
rend \(\widehat c'=1\), mais la coordonnée du lecteur dans la
coordinatisation interne reste \((r_+r_-)^{-1}\).

\subsection{Deux lecteurs de la trajectoire électronique et évaluation bêta}
\label{vi:dep:beta}

La trajectoire centrale commence en \((r,c)=(5,5)\), avec
l’orientation \((h,v)=(1,1)\). TRIT répète le cycle \((+1,-1,0)\).
Dans le premier régime, on lit \(r+c\) puis on avance de
\(c\mapsto1+((c-1+h)\bmod9)\) ; dans le second, on lit \(rc\) puis
on avance de \(r\mapsto1+((r-1+v)\bmod9)\) ; dans le régime
zéro, on lit \(9\), sans déplacement. On exprime la classe modulo
trois de la racine numérique de la lecture. Après les pas \(27,54,81\),
on inverse simultanément \(h,v\). La représentation de \(108\)
pas conserve la trajectoire complète et donne
\begin{equation}
 a=100000220,\quad b=120200000,\quad
 \gamma_e=(a^3b^3)^2,
 \quad\tau_e=(+,+,+,-,-,-,+,+,+,-,-,-).
 \label{vi:dep:trayectoria-electron}
\end{equation}
L’identité se vérifie en parcourant les neuf instants de chaque
fenêtre ; au troisième bloc, l’inversion remplace \(a\) par \(b\),
et le second parcours de \(54\) pas répète la représentation.
L’état conserve également ses coordonnées de mémoire.

Le premier lecteur mesure les désaccords cycliques de la représentation ternaire :
\[
 D_k=\#\{t\in\ZZ/108\ZZ:\gamma_e(t)\ne\gamma_e(t+k)\}.
\]
Comme le mot exprimé est de période \(54\), chaque désaccord
apparaît avec son translaté de \(54\) ; tous les \(D_k\) sont pairs.
Dans la moitié \(a^3b^3\), le dénombrement des désaccords pour
\(k=1,3,4,27,36\) donne respectivement \(27,28,34,24,16\).
En le doublant,
\[
 (D_1,D_3,D_4,D_{27},D_{36})=(54,56,68,48,32).
\]
Par conséquent,
\begin{equation}
 K_D=\left(D_{27}+D_{36},D_1,\frac{D_4-D_3}{2}\right)
     =(80,54,6).
 \label{vi:dep:KD}
\end{equation}

Le second lecteur agit sur l’orientation des douze fenêtres.
Définissons, avec des indices cycliques,
\[
 C_k=\sum_{r=0}^{11}\tau_r\tau_{r+k},\quad
 A_\ell=\sum_{r=0}^{11}\left(\sum_{j=0}^{\ell-1}\tau_{r+j}\right)^2,
 \quad P_6=\sum_{r=0}^{5}\tau_r\tau_{r+6}.
\]
La combinaison entière primitive orientée qui annule la partie diagonale
de \(A_3,A_4\) est \(\Omega=4A_3-3A_4\) : ses coefficients
satisfont \(3u+4v=0\), \(\gcd(u,v)=1\), et l’on fixe \(u>0\).
Le développement du carré prouve
\[
 A_\ell=\ell C_0+
     2\sum_{k=1}^{\ell-1}(\ell-k)C_k,
 \qquad\Omega=-2C_1-4C_2-6C_3.
\]
Pour \(\tau_e\), on obtient
\(C_0=12,C_1=4,C_2=-4,C_3=-12\),
\(A_3=44,A_4=32,P_6=6\). D’où
\begin{equation}
 K_\Omega=(\Omega,54,P_6)=(80,54,6)=K_D.
 \label{vi:dep:KOmega}
\end{equation}
La seconde composante conserve le demi-retour de \(108\) pas ;
elle n’est pas introduite par l’annulation diagonale. L’égalité des deux
lecteurs a été prouvée sur la trajectoire centrale spécifiée ;
elle ne s’étend pas par définition à toutes les collections de l’atlas.

\begin{definition}[Coordonnée électronique postérieure au retour]
La composition électronique complète est
\begin{equation}
 \varepsilon_e^\beta
  =\varepsilon_e^{(0)}
      R_{\rm act}^{80-54\alpha+6\Delta_4}.
 \label{vi:dep:electron-beta}
\end{equation}
\end{definition}
La positivité de \eqref{vi:dep:formula-basal} et
\eqref{vi:dep:Ract} permet de définir cette puissance au moyen du
logarithme réel. Le triplet entier provient indifféremment des deux
lecteurs centraux qui viennent d’être calculés. L’évaluation relative
préserve la constante de Planck réduite par le rapport de ses
sections, et la réalisation absolue utilise en outre
\eqref{vi:dep:reloj-absoluto} et la coordinatisation dimensionnelle déclarée.

Dans les tableaux hérités, \(\mathfrak e_0\) désigne la
coordonnée basale \(\varepsilon_e^{(0)}\) dans la coordinatisation qui y est
spécifiée. La section \(\hbar_{\rm int}\) utilisée pour une
évaluation de survie doit être déclarée dans la même coordinatisation :
pour la réalisation postérieure au retour du présent développement,
c’est \(\hbar_{\rm ret}\). Les lignes historiques conservent la
convention déclarée dans leur tableau d’origine et ne sont pas recalculées par un
changement silencieux de cette section.

Si \(\mathcal U_E\) est la base d’énergie de l’évaluation
électronique, \(E_e=\varepsilon_e^\beta\mathcal U_E\). Lorsqu’on emploie
le quantum chronologique comme base, le coefficient est
\[
 \widehat\varepsilon_e
   =\varepsilon_e^\beta\frac{\mathcal U_E}{E_{\rm clk}},
 \qquad E_e=\widehat\varepsilon_e E_{\rm clk}.
\]
L’égalité conserve l’objet et transforme sa coordonnée. Ainsi,
le nombre présenté dans une coordinatisation en MeV ne peut pas être multiplié
sans conversion par n’importe quel autre quantum d’énergie. L’électron
conserve sa construction structurelle lorsqu’il est exprimé par rapport à une
échelle chronologique : il n’est pas remplacé par la masse de ce quantum.

La comparaison métrologique s’applique au résultat de cette composition,
avec la version de \(\alpha\), la section d’action, le lecteur
\(\Delta_4\) et la coordinatisation des unités déclarées. En particulier,
\(d_4\) et \(\Delta_4\) ne sont pas interchangeables : les remplacer
simultanément dans le facteur basal et dans l’exposant modifie le
résultat selon le rapport positif
\[
 \frac{1+15\Delta_4}{1+15d_4}
 R_{\rm act}^{6(\Delta_4-d_4)}>1.
\]
Les formules conservées permettent de localiser cette différence sans
l’attribuer à une masse externe ni sélectionner à nouveau la trajectoire.

\section{La particule comme classe de sections persistantes
}
\label{vi:sec:particula}

La construction matérielle commence dans l’état produit par APP, TRIT et
TPK. APP fournit les positions et les deux opérations arithmétiques ;
TRIT distingue le régime et l’orientation de leurs transitions ; TPK
transporte cette information avec les quotients, les incidences de
frontière et la mémoire. Lors du prolongement de l’état, une évaluation finie
reste liée aux évaluations qui la précèdent. La notion de
particule s’appuie sur cette compatibilité et sur sa représentation matérielle.
Sa masse sera une observable de la section ainsi construite.


Cette séquence détermine le sens des trois objets qui seront
utilisés : l’état généalogique conserve l’extension complète ; la route
enregistre son histoire observable ; la représentation matérielle fait agir
cette histoire sur un espace d’états et permet d’évaluer des observables.
La distinction est essentielle même lorsque les trois objets partagent
une même étiquette électronique ou une même coordonnée de l’atlas.


\subsection{Extension, histoire observable et section compatible
}
\label{vi:subsec:extension-persistente}

Soient $\mathcal S_m$ les états admissibles survivants jusqu’à la profondeur
$m$, obtenus par le prolongement du TPK, et soient
$\rho_{n,m}:\mathcal S_n\to\mathcal S_m$, pour $n\geq m$, leurs
applications de restriction. La compatibilité du raffinement signifie

\begin{equation}
 \rho_{m,m}=I,\qquad
 \rho_{n,m}\rho_{p,n}=\rho_{p,m}\quad(p\geq n\geq m).
 \label{vi:eq:sistema-estados}
\end{equation}
Une extension survivante est un élément

\begin{equation}
 x=(x_m)_{m\geq0}\in\Omega^{\mathrm{surv}}_{\mathrm{HMT}}
 :=\varprojlim_m\mathcal S_m,
 \qquad \rho_{n,m}(x_n)=x_m.
 \label{vi:eq:extension-superviviente}
\end{equation}
L’existence et l’admissibilité de ces états proviennent de la
construction du noyau formel. L’expression de limite projective décrit
les compatibilités que conserve un état déjà engendré.


Soit $\mathcal R_m$ l’espace des histoires observables à la profondeur $m$.
Le lecteur de route $\operatorname{sh}_m:\mathcal S_m\to\mathcal R_m$
et la restriction $\operatorname{tr}_{n,m}$ satisfont

\begin{equation}
 \operatorname{tr}_{n,m}\operatorname{sh}_n
 =\operatorname{sh}_m\rho_{n,m}.
 \label{vi:eq:naturalidad-ruta}
\end{equation}
Ainsi, le lecteur transmet les restrictions de l’état à ses histoires.
La route de $x$ est la collection
$\gamma_x=(\gamma_x^{[m]})_m$, où
$\gamma_x^{[m]}=\operatorname{sh}_m(x_m)$.


\begin{proposition}[Compatibilité de l’histoire exprimée
]
\label{vi:prop:historia-compatible}
Pour toute extension \eqref{vi:eq:extension-superviviente}, ses histoires
finies satisfont
$\operatorname{tr}_{n,m}(\gamma_x^{[n]})=\gamma_x^{[m]}$.

\end{proposition}
\begin{proof}
En appliquant \eqref{vi:eq:naturalidad-ruta} à l’état $x_n$, on obtient

\[
 \operatorname{tr}_{n,m}(\gamma_x^{[n]})
 =\operatorname{sh}_m(\rho_{n,m}x_n)
 =\operatorname{sh}_m(x_m)=\gamma_x^{[m]}.
\]
La seconde égalité utilise précisément la condition de limite projective.

\end{proof}

La proposition établit la descente de l’extension à l’histoire.
Son contenu n’inclut pas l’injectivité du lecteur : deux extensions
peuvent exprimer une même histoire finie et conserver des mémoires différentes.
C’est pourquoi, lorsqu’on réutilise une route, son état de provenance est maintenu,
au lieu de le reconstruire à partir d’un chiffre ou d’un mot visible.


Sur chaque histoire $\gamma_x^{[m]}$, on considère le fibré discret
$\mathcal B_{\gamma_x^{[m]}}$ d’états matériels de la représentation
déclarée. Soient
$\mathscr E_m(x)=\Gamma(\mathcal B_{\gamma_x^{[m]}})$ ses espaces
de sections, avec les restrictions
$p_{n,m}:\mathscr E_n(x)\to\mathscr E_m(x)$ compatibles avec celles
de l’état. Une section persistante est une collection

\begin{equation}
 s=(s_m)_{m\geq m_0},\qquad s_m\in\mathscr E_m(x),\qquad
 p_{n,m}s_n=s_m\quad(n\geq m\geq m_0).
 \label{vi:eq:seccion-persistente}
\end{equation}
La profondeur initiale $m_0$ permet de décrire un état depuis toute
coupure à partir de laquelle sa représentation est définie. L’information
antérieure qui intervient dans ses observables demeure dans le registre
transporté par l’extension.


\subsection{Identité sous raffinement et équivalence des représentations
}
\label{vi:subsec:identidad-material}

Fixons les groupes de symétries internes $G_m$ qui agissent sur les fibres
matérielles et préservent les observables déclarées. Deux sections
persistantes $s$ et $s'$ représentent la même identité matérielle lorsque,
à partir d’un certain niveau $m_1$, il existe une collection $g_m\in G_m$ telle que

\begin{equation}
 s'_m=g_ms_m,\qquad
 p_{n,m}g_n=g_mp_{n,m},
 \label{vi:eq:equivalencia-persistente}
\end{equation}
et que les caractères observables du registre généalogique coïncident. Les
symétries admissibles sont fixées avant de former le quotient ; leur action
transporte le feuillet, l’orientation, la mémoire et la représentation avec
les lois correspondantes.


\begin{proposition}[Équivalence compatible des sections
]
\label{vi:prop:equivalencia-persistente}
La relation \eqref{vi:eq:equivalencia-persistente}, sur les collections qui
portent les mêmes types de représentation et d’observables, est une relation
d’équivalence. Ses classes sont indépendantes du retrait d’un nombre
fini de niveaux initiaux, pourvu que leur information
généalogique transportée soit conservée.

\end{proposition}
\begin{proof}
L’identité de chaque $G_m$ donne la réflexivité. L’inverse d’une collection
compatible est encore compatible : de
$p_{n,m}g_n=g_mp_{n,m}$ se déduit
$p_{n,m}g_n^{-1}=g_m^{-1}p_{n,m}$. Cela donne la symétrie.
Si $s'=gs$ et $s''=hs'$, à partir du plus grand des deux niveaux initiaux,
on a $s''_m=h_mg_ms_m$ et
$p_{n,m}h_ng_n=h_mg_mp_{n,m}$. La composition préserve également les
observables, et prouve la transitivité. Toutes ces relations se formulent
sur une queue commune ; changer le point initial dans cette queue
conserve la classe. La dernière affirmation concerne le changement de
présentation, non l’effacement de la mémoire de l’état.

\end{proof}

\begin{definition}[Particule matérielle persistante
]
\label{vi:def:particula}
Une particule HMT--MD est une classe $[s]_{\mathrm{pers}}$ de sections
\eqref{vi:eq:seccion-persistente} sous
\eqref{vi:eq:equivalencia-persistente}, accompagnée de son extension et de son
registre de provenance, qui admet une représentation matérielle
irréductible ou un bloc minimal invariant et une sortie spectrale stable
sous le couplage spécifié.

\end{definition}

La définition réunit des conditions de différents niveaux. Le prolongement
et le registre appartiennent à l’état HMT. La représentation et le
couplage appartiennent à sa réalisation matérielle. La stabilité
spectrale se rapporte à l’opérateur et au domaine utilisés par ce
couplage. Cette organisation permet d’étudier conjointement identité,
charge, spin et masse sans transformer une observable particulière en définition
de toutes les autres.


Le retour nonadique illustre la distinction entre identité et état
instantané. Pour le lecteur de phase $g$ et le registre de mémoire
$\mathcal B$, peuvent coexister

\begin{equation}
 g(k+9)=g(k),\qquad
 p_{k+9,k}(s_{k+9})=s_k,\qquad
 \mathcal B_{k+9}\ne\mathcal B_k.
 \label{vi:eq:persistencia-fase-memoria}
\end{equation}
La première égalité est périodique ; la seconde exprime la persistance ; la
troisième enregistre l’avancée. Le retour de phase appartient à l’histoire
d’une même identité matérielle, tandis que l’état conserve le fait qu’un
tour supplémentaire a été effectué.


\subsection{L’espèce comme fibre : portée de l’atlas fini
}
\label{vi:subsec:especie-fibra}

L’atlas construit par les opérations APP--TRIT--TPK a pour support
$\mathcal C_{81}=\{1,\ldots,9\}^2$. Ses quatre orientations par
cellule forment le domaine de $324$ routes orientées. L’application de
registre $\mathcal R$ distingue $56$ configurations de route, tandis que le
classificateur interne

\begin{equation}
 q:\mathcal C_{81}\longrightarrow\Sigma_{13}
 \label{vi:eq:clasificador-especies}
\end{equation}
réunit $13$ types de collection et de niveau. Ces recensements et leurs règles sont
développés dans le corps consacré à l’atlas ; ici, on précise quel objet
ils transmettent à la définition de la particule.


Pour $y\in\Sigma_{13}$, la fibre $q^{-1}(y)$ conserve toutes ses
cellules et tous ses registres. L’ordre interne de chaque fibre de
$\mathcal R$ étant fixé, soit $k(c)$ le rang de la cellule au sein de celle-ci. Sa
présentation finie est

\begin{equation}
 \mathfrak S_y^{\mathrm{atlas}}
 =\{(\mathcal R(c),k(c)):q(c)=y\}.
 \label{vi:eq:multiseccion-atlas}
\end{equation}
Le rang est ajouté pour séparer les cellules qui possèdent le même registre de
collection ; il ne sélectionne pas une masse. La paire $(\mathcal R(c),k(c))$
récupère la cellule dans la coordinatisation ordonnée parce que la seconde donnée
individualise sa position dans la fibre de la première.


Si $\pi_{81}:\Omega^{\mathrm{surv}}_{\mathrm{HMT}}\to
\mathcal C_{81}$ est le lecteur cellulaire, le relèvement de l’espèce est

\begin{equation}
 \widetilde{\mathfrak S}_y
 =\{x\in\Omega^{\mathrm{surv}}_{\mathrm{HMT}}:
                  q(\pi_{81}x)=y\}.
 \label{vi:eq:multiseccion-enriquecida}
\end{equation}
Cette préimage contient des extensions, tandis que
\eqref{vi:eq:multiseccion-atlas} contient des présentations finies. La
construction et le domaine du lecteur $\pi_{81}$ font partie du passage
entre les deux. Une particule concrète se réalise comme section persistante
au sein de la multisection enrichie correspondante.


\begin{theorem}[Point fixe central et sélection équivariante
]
\label{vi:thm:selector-central}
Soit $\rho_c(r,s)=(10-r,10-s)$ l’inversion cellulaire. Dans les fibres
$\rho_c$-stables du classificateur \eqref{vi:eq:clasificador-especies},
avec une action triviale sur l’étiquette $y$, une sélection ponctuelle
équivariante ne peut prendre de valeurs que dans la cellule $(5,5)$. La fibre
électronique de niveau zéro contient cet unique point fixe ; aucune autre
fibre n’admet une telle sélection.

\end{theorem}
\begin{proof}
Résoudre $\rho_c(r,s)=(r,s)$ donne $r=s=5$. Si $a(y)\in q^{-1}(y)$
est une sélection équivariante et si $y$ reste fixe, alors
$a(y)=a(\rho_c y)=\rho_c a(y)$. Sa valeur doit être le point fixe
$(5,5)$. L’atlas situe cette cellule dans la fibre électronique de niveau
zéro, dont le cardinal vaut un. Pour toute autre fibre, la condition
de sélection produirait un point fixe qui ne lui appartient pas.

\end{proof}

La conséquence est constructive : les autres espèces sont présentées par
orbites ou multisections, et une réalisation qui nécessite une droite
particulière doit déclarer sa représentation et son couplage. L’électron
fournit le cas central de rang un sans transformer toutes
les collections en copies de ce cas. En outre, la cellule $(5,5)$, la signature
centrale brute, la signature régionale $555555$ et l’origine relative d’une
droite d’action conservent leurs domaines propres. L’unicité du point
n’identifie pas ces registres entre eux.


\subsection{Du registre matériel à l’observable de masse
}
\label{vi:subsec:particula-observable}

La chaîne utilisée par l’observable de masse conserve l’ordre

\begin{equation}
 \begin{gathered}
 \text{extension}\longrightarrow\text{route}\longrightarrow
 \text{registre}\longrightarrow\text{signature entière}\\
 \longrightarrow\text{caractère}\longrightarrow\text{tours}\longrightarrow
 \text{opérateur matériel}.
 \end{gathered}
 \label{vi:eq:cadena-material}
\end{equation}
Chaque flèche retient les données nécessaires à la suivante. La signature est
une lecture du registre ; le caractère évalue cette signature avec les
coefficients internes déjà construits ; les tours incorporent son
raffinement ; l’opérateur agit finalement dans la représentation de la
fibre sélectionnée par le couplage. Sa sortie peut être un scalaire,
un multiplet ou une distribution spectrale, selon cette représentation.
La comparaison métrologique a lieu après cette construction.


Une espèce est décrite par la charge, la représentation de spin, la couleur,
la saveur, la génération, la configuration de route et l’opérateur matériel. Ces données
expliquent pourquoi deux préparations de masses coïncidentes peuvent
conserver des informations différentes, et pourquoi une même identité peut
exprimer des observables différentes sous des couplages distincts. L’identité
réside dans la classe persistante et dans ses transports admissibles ;
la mesure est une lecture de cette identité dans le codomaine déclaré.


La même hiérarchie sépare une extension persistante d’une section
d’horizon fini. Si $\operatorname{Ext}_L(s)\ne\varnothing$ et
$\operatorname{Ext}_{L+1}(s)=\varnothing$, l’horizon $L$ enregistre
une obstruction au prolongement. Une participation intermédiaire de ce
type dans une composition ne crée pas, à elle seule, une nouvelle espèce
asymptotique. De même, une résonance exige l’opérateur de
survie et son pôle, en plus de l’identité de la section qu’elle
transporte. Ainsi est préservée la différence entre persistance,
prolongement fini et évaluation spectrale.

\section{Conjugaison, orientation et retour des sections matérielles
}
\label{vi:sec:conjugacion}

L’identité persistante de la section~\ref{vi:sec:particula} admet
des opérations d’inversion à différents niveaux de sa structure.
L’inversion cellulaire agit dans l’atlas ; la conjugaison dodécaphasique
agit dans le registre d’orientation ; le relèvement d’une boucle agit
dans une fibre ; une rotation spinorielle agit dans un espace de
représentation. Spécifier ces domaines permet de composer les
opérations et d’expliquer quelles informations sont conservées à chaque retour.


\subsection{Involution dodécaphasique et parité des observables
}
\label{vi:subsec:conjugacion-registro}

Soit $\mathcal T_{12}=\{-1,0,1\}^{12}$ l’espace ambiant des registres
tritiques de douze positions. Le domaine admissible du TPK est le
sous-ensemble muni des compatibilités de route déjà construites. L’opération
sur l’espace ambiant est

\begin{equation}
 C_M:\mathcal T_{12}\longrightarrow\mathcal T_{12},\qquad
 (C_M\tau)_j=-\tau_{11-j},\quad 0\leq j<12.
 \label{vi:eq:cm}
\end{equation}
La restriction aux registres admissibles utilise la stabilité de ce
domaine sous la conjugaison. L’ordre inversé et le changement de signe
sont des composantes simultanées de l’opération.


\begin{proposition}[Parité linéaire et quadratique du registre
]
\label{vi:prop:paridad-registro}
On a $C_M^2=I$. Si $a_j=a_{11-j}$ et
$L_a(\tau)=\sum_ja_j\tau_j$, alors
$L_a(C_M\tau)=-L_a(\tau)$. Si une forme quadratique
$Q_b(\tau)=\sum_{j,k}b_{jk}\tau_j\tau_k$ satisfait
$b_{jk}=b_{11-j,11-k}$, alors $Q_b(C_M\tau)=Q_b(\tau)$.

\end{proposition}
\begin{proof}
À chaque position,
$[C_M(C_M\tau)]_j=-(-\tau_{11-(11-j)})=\tau_j$.
Pour la fonctionnelle linéaire, substituer $k=11-j$ dans la somme donne
$L_a(C_M\tau)=-\sum_ka_{11-k}\tau_k=-L_a(\tau)$.
Dans l’expression quadratique, les deux changements de signe s’annulent et
la substitution $(r,s)=(11-j,11-k)$ récupère $Q_b$ par la
symétrie de ses coefficients.

\end{proof}

L’orientation linéaire et les corrélations de mémoire ont donc
des lois différentes. Même une composition de ces observables doit
conserver leurs parités : si $L_a$ est impair, son caractère exponentiel
satisfait

\begin{equation}
 \exp(L_a(C_M\tau))=\exp(L_a(\tau))^{-1}.
 \label{vi:eq:caracter-inverso}
\end{equation}
L’égalité exprime une inversion multiplicative. L’invariance d’un
opérateur matériel complet s’étudie sur cet opérateur, avec ses
canaux et sa représentation, et se démontre par la covariance
correspondante.


L’inversion cellulaire du théorème~\ref{vi:thm:selector-central} est
$\rho_c:\mathcal C_{81}\to\mathcal C_{81}$, tandis que $C_M$ a pour
domaine \eqref{vi:eq:cm}. Toutes deux sont des involutions, mais pour
les comparer, on requiert l’application de registre de la route et sa loi de
transport. Cette donnée évite de remplacer la conjugaison d’une histoire
par l’inversion d’une coordonnée isolée de l’atlas.


\subsection{Relèvement aux antisections et conservation de la masse
}
\label{vi:subsec:antisecciones}

La conjugaison s’applique à la particule au moyen de deux applications
compatibles. L’application de base $C_m^{\mathrm b}$ transporte l’état
$x_m$ vers son état conjugué $\bar x_m$ ; le relèvement
$\widetilde C_m:\mathscr E_m(x)\to\mathscr E_m(\bar x)$
transporte ses sections matérielles. Soient $t_m$ le lecteur du
registre dodécaphasique et $\bar p_{n,m}$ les restrictions dans la
collection conjuguée. Leurs conditions sont la conservation de
l’admissibilité, l’involutivité et la compatibilité :

\begin{equation}
 \begin{aligned}
 \rho_{n,m}C_n^{\mathrm b}&=C_m^{\mathrm b}\rho_{n,m},
 &t_m(C_m^{\mathrm b}x_m)&=C_Mt_m(x_m),\\
 \bar p_{n,m}\widetilde C_n&=\widetilde C_mp_{n,m},
 &\widetilde C_{m,\bar x}\widetilde C_{m,x}&=I.
 \end{aligned}
 \label{vi:eq:levantamiento-cm}
\end{equation}
Dans la première ligne, le lecteur exprime le registre d’orientation
de l’histoire. Le transport restant de l’état accompagne cette
représentation. La seconde ligne donne les lois du relèvement ;
les indices $x,\bar x$ explicitent ses deux fibres de départ.
Si $s$ est persistante, son antisection est
$\bar s=(\widetilde C_ms_m)_m$.


\begin{proposition}[Persistance de l’antisection
]
\label{vi:prop:antiseccion-persistente}
Les conditions \eqref{vi:eq:levantamiento-cm} transportent les sections
persistantes vers des sections persistantes et restituent la section lorsqu’on
conjugue deux fois. Si elles transportent en outre les symétries internes par
$g_m\mapsto\widetilde C_mg_m\widetilde C_m^{-1}$, elles induisent une
involution sur les classes matérielles de la
définition~\ref{vi:def:particula}.

\end{proposition}
\begin{proof}
La compatibilité donne
$\bar p_{n,m}\bar s_n=\widetilde C_mp_{n,m}s_n
=\widetilde C_ms_m=\bar s_m$.
L’involutivité restitue $s$. Pour $s'_m=g_ms_m$, leurs
antisections sont reliées au moyen de
$\widetilde C_mg_m\widetilde C_m^{-1}$. La loi de restriction
dans \eqref{vi:eq:levantamiento-cm} conserve la compatibilité de cette
collection, qui appartient aux groupes internes par la condition de
l’énoncé. Ainsi, la construction descend aux classes.

\end{proof}

La représentation de charge utilise la partie impaire du registre
matériel. Si $\mathfrak q$ est le caractère de charge déclaré, sa
loi est $\mathfrak q(\bar s)=-\mathfrak q(s)$. La charge élémentaire
qui intervient lorsqu’on exprime cette coordonnée dans une droite dimensionnelle est
la sortie interne déjà construite ; la représentation s’applique ensuite
à l’état, sans utiliser de charge cible pour choisir son registre.


Pour la masse, on a besoin de la covariance de son opérateur. Soient
$\widehat M_s$ et $\widehat M_{\bar s}$ les opérateurs autoadjoints
des deux réalisations, et soit $\mathcal C$ leur opérateur d’entrelacement
unitaire ou antiunitaire. Leur domaine se transporte au moyen de
$\mathcal C\operatorname{Dom}(\widehat M_s)
=\operatorname{Dom}(\widehat M_{\bar s})$. La condition opératorielle est

\begin{equation}
 \mathcal C\widehat M_s\mathcal C^{-1}
 =\widehat M_{\bar s}.
 \label{vi:eq:covariancia-masa}
\end{equation}

\begin{proposition}[Conservation spectrale sous conjugaison
]
\label{vi:prop:masa-conjugada}
Sous \eqref{vi:eq:covariancia-masa}, le transport d’une section
propre de masse $m\in\mathbb R$ conserve cette valeur propre. La mesure
spectrale se transporte au moyen de
$E_{\bar s}(B)=\mathcal C E_s(B)\mathcal C^{-1}$ pour tout
borélien réel $B$. En particulier, une réalisation scalaire conserve
la masse lors du passage à l’antisection.

\end{proposition}
\begin{proof}
Si $\widehat M_s\psi=m\psi$, la covariance donne
$\widehat M_{\bar s}\mathcal C\psi
=\mathcal C(m\psi)=m\mathcal C\psi$ ; la dernière égalité
vaut aussi dans le cas antiunitaire parce que $m$ est réel. Le calcul
fonctionnel réel des opérateurs autoadjoints transporte leurs
projecteurs spectraux et prouve la seconde affirmation.

\end{proof}

Dans un registre de bande, la même condition peut se voir sous une forme
élémentaire. Étant donnée une lecture réelle $E_0(\tau)$, soient
$E_\pm=(E_0\pm E_0\circ C_M)/2$ et soit $\xi\in\{+1,-1\}$
l’orientation de feuillet. Alors

\begin{equation}
 E_{\mathrm{band}}(\tau,\xi)=E_+(\tau)+\xi E_-(\tau),\qquad
 E_{\mathrm{band}}(C_M\tau,-\xi)=E_{\mathrm{band}}(\tau,\xi).
 \label{vi:eq:banda-paridad}
\end{equation}
L’égalité s’obtient en utilisant le fait que $E_+$ est pair et $E_-$ impair.
Elle explique comment l’énergie de bande peut être conservée en transportant
simultanément registre et feuillet ; son identification à la masse d’une
réalisation concrète utilise \eqref{vi:eq:covariancia-masa}.


Le registre électronique
$\tau_e=(+,+,+,-,-,-,+,+,+,-,-,-)$ satisfait $C_M\tau_e=\tau_e$.
Cette égalité concerne le mot de douze positions. La paire
$(\tau_e,\xi)$ conserve une orientation supplémentaire que la
conjugaison change. La distinction permet de maintenir le centre
électronique unique et, à la fois, les réalisations conjuguées de charge ;
un mot fixe ne devient pas une preuve de neutralité.


\subsection{Holonomie de signe et bande de persistance
}
\label{vi:subsec:mobius}

Soit $\ell$ une boucle de retour du lecteur matériel d’une extension
survivante. Son groupe de tours est $\langle\ell\rangle\simeq
\mathbb Z$. Sur ce groupe, on déclare le caractère de signe

\begin{equation}
 \varepsilon:\langle\ell\rangle\longrightarrow\{+1,-1\},
 \qquad \varepsilon(\ell)=-1.
 \label{vi:eq:holonomia-signo}
\end{equation}
Le groupe de tours enregistre combien de fois la boucle est parcourue ; son
type est différent de la phase finie $C_9$. Le fibré d’intervalles
associé au caractère est

\begin{equation}
 \mathcal M_\ell=([0,1]\times[-1,1])/
                  \bigl((0,u)\sim(1,-u)\bigr).
 \label{vi:eq:banda-mobius}
\end{equation}

\begin{theorem}[Retour orienté sur la bande
]
\label{vi:thm:retorno-mobius}
Le fibré \eqref{vi:eq:banda-mobius} est un ruban de Möbius.
Son transport sur $\ell$ change le signe de la fibre et sur
$\ell^2$ le restitue. Par conséquent, un vecteur non nul de cette droite
requiert deux tours pour son retour orienté. Cela conserve la
distinction entre retour d’orientation et retour de l’état complet.

\end{theorem}
\begin{proof}
La relation d’extrémités de \eqref{vi:eq:banda-mobius} est l’application
linéaire $u\mapsto-u$, de déterminant négatif. Elle donne la droite réelle
non orientable sur le cercle et sa bande
d’intervalles correspondante. Composer deux transports produit $u\mapsto u$.
La mémoire de l’extension ne fait pas partie de ce quotient de signe ;
sa mise à jour continue d’obéir au transport du TPK et à
\eqref{vi:eq:persistencia-fase-memoria}.

\end{proof}

Un paramétrage d’un tour de la base au moyen de $2\pi$
représente la restitution du revêtement au moyen de $4\pi$. Si le
lecteur temporel attribue $108$ pas au tour pertinent, le
relèvement orienté se ferme à $216$ pas. Ces attributions
requièrent la boucle et le lecteur indiqués ; elles ne divisent pas l’histoire en
deux particules successives de $108$ pas.


\begin{proposition}[Décomposition des observables sur le revêtement double
]
\label{vi:prop:observables-paridad}
Soit $p:\widetilde X\to X$ un revêtement principal double, d’involution
$\iota$. Pour $f\in C^0(\widetilde X,\mathbb R)$,

\[
 f^+=\frac{f+f\circ\iota}{2},\qquad
 f^-=\frac{f-f\circ\iota}{2}
\]
donnent $f=f^++f^-$, avec $f^+\circ\iota=f^+$ et
$f^-\circ\iota=-f^-$. La première partie descend à $X$ ; la
seconde représente une section de la droite de signe associée. Sur
la droite de Möbius, toute section globale continue possède un zéro.

\end{proposition}
\begin{proof}
Les identités s’obtiennent en substituant $\iota^2=I$. La partie
invariante est constante dans chaque fibre et définit une fonction sur le
quotient. La partie anti-invariante satisfait la loi de transition
de la représentation de signe. Sur le cercle, une section de cette
droite est représentée par une fonction réelle continue $u$ sur $[0,1]$
telle que $u(1)=-u(0)$. Si l’extrémité vaut zéro, il y a déjà un zéro ; dans le cas
contraire, les extrémités ont des signes opposés et le théorème des
valeurs intermédiaires fournit un zéro intérieur.

\end{proof}

La décomposition distingue deux classes d’observables sur un
même revêtement. Elle permet de représenter une information impaire d’orientation
et des données paires conservées. Pour l’appliquer à un registre matériel, on
utilise le relèvement \eqref{vi:eq:levantamiento-cm} ; le caractère
de signe décrit la droite transportée, non toutes les données de la
particule.


\subsection{Réseau bidirectionnel et conjugaison des canaux
}
\label{vi:subsec:reticulo-bidireccional}

Le registre angulaire entier exprime deux coordonnées

\begin{equation}
 W_+=6n_A+s_C,\qquad W_-=6n_A-s_C,
 \qquad (n_A,s_C)\in\mathbb Z^2.
 \label{vi:eq:reticulo-ida-retorno}
\end{equation}
Les coefficients $n_A,s_C$ sont des lectures internes de route. Le facteur
$6$ appartient à la normalisation dodécaphasique de ce registre.


\begin{proposition}[Indice du réseau et échange des lectures
]
\label{vi:prop:indice-bidireccional}
L’image de \eqref{vi:eq:reticulo-ida-retorno} a pour indice $12$
dans $\mathbb Z^2$ et pour forme normale de Smith $\operatorname{diag}(1,12)$.
Elle est caractérisée par $W_++W_-\equiv0\pmod{12}$. L’opération
$s_C\mapsto-s_C$ échange $W_+$ et $W_-$.

\end{proposition}
\begin{proof}
La matrice de l'application est
$\bigl(\begin{smallmatrix}6&1\\6&-1\end{smallmatrix}\bigr)$.
Le plus grand commun diviseur de ses entrées est $1$ et son déterminant
est $-12$, ce qui donne les facteurs invariants indiqués. Les
coordonnées inverses sont
$n_A=(W_++W_-)/12$ et $s_C=(W_+-W_-)/2$.
Si la somme est divisible par $12$, les deux entiers ont la même
parité et les deux expressions sont entières. La condition est également
nécessaire par la somme de \eqref{vi:eq:reticulo-ida-retorno}.
L’affirmation sur l’échange se vérifie en substituant $-s_C$.

\end{proof}

Dans la représentation angulaire réelle, les coordonnées déjà engendrées
$a=A_{\mathrm{rad}}$ et $b=C^*_{\mathrm{rad}}$ produisent les
canaux $q_\pm=\exp[-(a\pm b)]$. Dans le domaine $a>|b|$, tous deux
appartiennent à $(0,1)$. Soit $R^*=R$, $R^2=I$ la graduation des
feuillets et soit $C$ un échange unitaire tel que $CRC^{-1}=-R$.
Pour $T(b)=\exp[-(aI+bR)]$, on obtient

\begin{equation}
 CT(b)C^{-1}=T(-b),\qquad
 CP_\pm C^{-1}=P_\mp,\quad P_\pm=(I\pm R)/2.
 \label{vi:eq:conjugacion-canales}
\end{equation}
En effet, la conjugaison transforme l’exposant et chaque terme
de sa série convergente. Cette preuve transmet l’échange aux
canaux. Le réseau entier et l’opérateur réel ont des domaines
différents et expriment la même opération d’échange à travers
leurs applications de lecture déclarées.


\subsection{Représentation spinorielle du bloc électronique
}
\label{vi:subsec:spin-electron}

La construction APP tridimensionnelle des projecteurs de somme et de
produit sélectionne un plan principal électronique par le mode
tritique $(1,-1,0)^3$. Sa dérivation et l’évaluation des
projecteurs sont incluses parmi les antécédents du présent
article. Dans une base orthonormale de ce plan, leurs restrictions
sont

\begin{equation}
 P=\begin{pmatrix}1&0\\0&0\end{pmatrix},\qquad
 Q=\begin{pmatrix}1/4&\sqrt3/4\\\sqrt3/4&3/4\end{pmatrix},
 \qquad S=2P-I,\quad J=\frac4{\sqrt3}[P,Q].
 \label{vi:eq:bloque-electronico}
\end{equation}
Cette représentation reçoit les deux projecteurs construits et
conserve leur orientation. La multiplication donne
$S^*=S$, $J^*=-J$, $S^2=I$, $J^2=-I$ et $SJ=-JS$.
Concrètement, $S=\operatorname{diag}(1,-1)$ et
$J=\bigl(
\begin{smallmatrix}0&1\\-1&0\end{smallmatrix}\bigr)$.

Soit $\mathscr S_e=E_e\otimes_{\mathbb R}\mathbb C$ la
complexification du plan réel. L’unité scalaire $i$ appartient
au corps de cette représentation ; $J$ reste
l’endomorphisme réel provenant des projecteurs.


\begin{proposition}[Triplet hermitien et représentation des rotations
]
\label{vi:prop:spin}
Dans $\mathscr S_e$, définissons

\begin{equation}
 \sigma_1=SJ,\qquad \sigma_2=-iJ,\qquad \sigma_3=S.
 \label{vi:eq:terna}
\end{equation}
Ces opérateurs sont hermitiens, de trace nulle, et satisfont

\begin{equation}
 \sigma_a\sigma_b=\delta_{ab}I+
 i\sum_c\epsilon_{abc}\sigma_c,
 \qquad L_a=\sigma_a/2,\qquad
 [L_a,L_b]=i\sum_c\epsilon_{abc}L_c,
 \quad\sum_aL_a^2=\tfrac34I.
 \label{vi:eq:spin-algebra}
\end{equation}
Ils constituent la représentation irréductible de spin $1/2$ sur
le bloc électronique.

\end{proposition}
\begin{proof}
Les relations de $S,J$ donnent $(SJ)^*=SJ$, $(-iJ)^*=-iJ$ et
$\sigma_a^2=I$. Les produits cycliques sont
$(SJ)(-iJ)=iS$, $(-iJ)S=iSJ$ et $S(SJ)=J=i(-iJ)$ ;
inverser l’ordre change le signe. Cela prouve l’identité de
produit et les commutateurs. Dans la base de
\eqref{vi:eq:bloque-electronico}, les trois matrices sont

\[
 \begin{pmatrix}0&1\\1&0\end{pmatrix},\quad
 \begin{pmatrix}0&-i\\i&0\end{pmatrix},\quad
 \begin{pmatrix}1&0\\0&-1\end{pmatrix}.
\]
Leur trace est nulle et, avec $I$, elles engendrent $M_2(\mathbb C)$.
Un sous-espace invariant par les trois est invariant par toute cette
algèbre et ne peut être que $0$ ou $\mathscr S_e$. Enfin,
sommer $L_a^2=I/4$ donne la valeur $3I/4$ de l’opérateur quadratique.

\end{proof}

La reconnaissance de ce triplet comme matrices de Pauli est postérieure
à \eqref{vi:eq:bloque-electronico}. La norme et le produit
$\langle X,Y\rangle=\tfrac12\operatorname{tr}(XY)$ dans
l’espace réel des matrices hermitiennes de trace nulle donnent son espace
de directions. Pour $n\in\mathbb R^3$, $|n|=1$, soient

\begin{equation}
 H_e(n)=\sum_an_a\sigma_a,\qquad
 h_e(n)=\tfrac12H_e(n),\qquad
 \Pi_\pm(n)=\tfrac12(I\pm H_e(n)).
 \label{vi:eq:helicidad}
\end{equation}

\begin{proposition}[Décomposition directionnelle et double retour
]
\label{vi:prop:helicidad}
Les opérateurs $\Pi_\pm(n)$ sont des projecteurs orthogonaux de
rang un et $h_e(n)\Pi_\pm(n)=\pm\tfrac12\Pi_\pm(n)$.
Le groupe de rotations autour de $n$ est

\begin{equation}
 U_n(\theta)=e^{-i\theta H_e(n)/2}
 =\cos(\theta/2)I-i\sin(\theta/2)H_e(n),\qquad
 U_n(2\pi)=-I,\quad U_n(4\pi)=I.
 \label{vi:eq:retorno-spin}
\end{equation}
\end{proposition}
\begin{proof}
La partie antisymétrique de \eqref{vi:eq:spin-algebra} s’annule
lorsqu’elle est contractée avec $n_an_b$, de sorte que $H_e(n)^2=I$.
Il en découle l’idempotence, l’orthogonalité et la somme
$\Pi_++\Pi_-=I$. Ils sont hermitiens et de trace un, donc leur
rang vaut un. Multiplier par $H_e(n)/2$ donne les valeurs propres.
Séparer les puissances paires et impaires de l’exponentielle produit
\eqref{vi:eq:retorno-spin} et sa loi de groupe. La normalisation
angulaire se vérifie par

\[
 U_n(\theta)H_e(v)U_n(\theta)^*
 =H_e\bigl(v\cos\theta+(n\times v)\sin\theta
       +n(n\cdot v)(1-\cos\theta)\bigr),
\]
qui résulte de la multiplication avec l’identité de produit du
triplet. L’action induite sur les directions est une rotation
d’angle $\theta$ ; le relèvement spinoriel porte sa moitié.

\end{proof}

Si $n$ se réalise comme direction de propagation, $h_e(n)$ est
l’hélicité adimensionnelle. L’inversion de cette direction satisfait
$H_e(-n)=-H_e(n)$ et $\Pi_\pm(-n)=\Pi_\mp(n)$.
Cette opération change les étiquettes directionnelles ; la conjugaison
de charge transporte le registre matériel au moyen de
\eqref{vi:eq:levantamiento-cm}. La chiralité se définit au moyen de
la graduation relativiste correspondant à la représentation et à
ses projecteurs, distincts de \eqref{vi:eq:helicidad}.


L’évaluation scalaire électronique agit comme
$\varepsilon_e^\beta I$ sur $\mathscr S_e$ et commute avec
le triplet et avec $\Pi_\pm(n)$. L’information rotationnelle conserve
la structure de cette évaluation sans introduire d’autre facteur dans
la masse. De son côté, le ruban de Möbius transporte une droite
réelle autour d’une boucle et l’holonomie nonadique transporte
phase et mémoire. Les trois retours peuvent intervenir dans une
même réalisation ; leurs domaines et leurs opérateurs d’entrelacement déterminent la
composition précise.


\section{Du registre de trajectoire à l’opérateur de masse
}
\label{vi:sec:operador-masa}

La construction distingue deux opérations. D’abord, la dynamique produit
un registre de la trajectoire. Ensuite, une fonctionnelle extrait des coordonnées
entières et les évalue au moyen des constantes internes déjà obtenues.
Conserver cet ordre explique comment la masse peut être une grandeur de la
persistance sans épuiser sa structure.


\subsection{Lecture chronologique et événements orientés
}

Soit \(d_1,\ldots,d_{108}\) la lecture numérique d’une route enrichie.
Dans la coordinatisation déclarée, le support d’action est \(\{1,3,5,7\}\), et

\[
 n_A=\#\{t:d_t\in\{1,3,5,7\}\}.
\]
On conserve en outre les \(12\) événements des blocs de \(9\) pas :

\[
 e_m=\sigma_m\#\{t:9m+1\le t\le9m+9,\ d_t\in\{2,4,6,8\}\},
 \quad 0\le m<12,
\]
où \(\sigma=(+,+,+,-,-,-,+,+,+,-,-,-)\). Le support de cet
extracteur est une partie explicite de la coordinatisation de lecture ; sa fonction ne
doit pas être confondue avec le sélecteur ternaire de phase du TPK.


La distribution des événements est pertinente même lorsque leurs sommes
coïncident. Pour la voir, nous introduisons

\[
 p_i=e_i+e_{i+6},\quad a_i=e_i-e_{i+6},\quad
 s_C=\sum_{i=0}^{5}p_i,\quad M_i=p_i-p_5\ (0\le i<5).
\]

\begin{proposition}[Reconstruction entière des événements
]
\label{vi:prop:eventos}
La transformation \((e_0,\ldots,e_{11})\mapsto(s_C,M_0,\ldots,M_4,a_0,\ldots,a_5)\)
est injective sur \(\mathbb Z^{12}\). Son image est caractérisée par

\[
 s_C-\sum_{i=0}^4M_i\equiv0\pmod6,
 \qquad p_i\equiv a_i\pmod2\quad(0\le i<6),
\]
où \(p_5=(s_C-\sum M_i)/6\) et \(p_i=p_5+M_i\) pour \(i<5\).

\end{proposition}
\begin{proof}
Sommer les six expressions de \(p_i\) donne
\(s_C=6p_5+\sum_{i<5}M_i\). Cela récupère les \(p_i\), puis

\[
 e_i=\frac{p_i+a_i}{2},\qquad
 e_{i+6}=\frac{p_i-a_i}{2}.
\]
Les congruences sont nécessaires pour que ces quantités soient entières et
suffisantes pour que les formules définissent une préimage entière. Substituer
dans la transformation directe récupère toutes les coordonnées.

\end{proof}

La décomposition \(12=1+5+6\) acquiert ainsi un sens opératoire :
le total, les différences et les orientations récupèrent les événements.
La division par \(6\) appartient à cette reconstruction. Lorsque le caractère
utilise \(s_C/6\), \(s_C\) reste la somme entière ; une coordonnée déjà
divisée n’est pas divisée une seconde fois.


\subsection{Mémoire torsionnelle et signature
}

Pour chaque occurrence de \(9\), la lecture conserve son prédécesseur.
Nous définissons, avec des indices cycliques,

\[
 j_2(d)=\begin{cases}0,&d=9,\\1,&d\in\{2,4,6,8\},\\-1,&d\in\{1,3,5,7\},\end{cases}
 \quad z_t=\mathbf1_{\{d_t=9\}}j_2(d_{t-1}).
\]
Alors
\[
 T_z=\sum_{t=1}^{108}z_t,\quad
 C_z=\sum_{t\in\{27,54,81,108\}}z_t,\quad
 k_{\Delta_4}=T_z-2C_z.
\]
La première somme conserve la mémoire accumulée et la seconde distingue les
fermetures de couronne. Éliminer les prédécesseurs avant d’évaluer ces
expressions altère l’objet lu.


Avec \(\tau_m=\operatorname{sgn}(e_m)\), les autres coordonnées sont

\[
 b_{120}=\sum_{a=0}^{3}\operatorname{sgn}(e_a+e_{a+4}+e_{a+8}),
 \qquad b_{\zeta}=\sum_{m=0}^{11}\tau_m\tau_{m+3},
\]
où le second indice est modulo \(12\). On obtient la signature

\[
 L_{\mathrm{tick}}(\gamma)
 =\sigma(\gamma)=(n_A,s_C,k_{\Delta_4},b_{120},b_{\zeta})\in\mathbb Z^5.
\]
Cette application est déterministe dans la coordinatisation déclarée. Les signes et les
corrélations expliquent pourquoi la composition doit se réaliser sur le
registre avant d’extraire les \(5\) entiers : la somme des signatures ne
représente pas, en général, la composition de deux histoires.


\subsection{Caractère et spécialisation interne
}

Soit \(X=(X_A,X_C,X_4,X_{120},X_{270})\). Le caractère formel est

\[
 \chi(\sigma;X)=\exp\!\left(n_AX_A+s_CX_C+k_{\Delta_4}X_4+
 b_{120}X_{120}+b_{\zeta}X_{270}\right).
\]
Pour chaque \(X\), il satisfait
\(\chi(\sigma+\sigma';X)=\chi(\sigma;X)\chi(\sigma';X)\)
parce que l’exposant est linéaire en la signature. C’est une loi sur les signatures ; elle
n’affirme pas que l’extracteur est additif sur les routes sans conserver l’interface.


La spécialisation emploie les coordonnées internes construites précédemment :

\[
 X_A=A_{\mathrm{rad}},\quad X_C=C^*_{\mathrm{rad}}/6,
 \quad X_4=\Delta_4,\quad X_{120}=s_{120},\quad
 X_{270}=\zeta_{270}:=135\Delta_4^2.
\]
Le caractère est adimensionnel et positif. Son effet est une homothétie sur
la droite de masse. La forme exponentielle détermine comment agissent les
coefficients ; leur provenance est fixée dans les registres et les constructions
antérieurs.


\subsection{Hiérarchie bicouche et conservation des deux généalogies
}

Nous écrivons \(x=A_{\mathrm{rad}}\), \(y=C^*_{\mathrm{rad}}\). Dans la
chambre \(x>y>0\), les canaux sont

\[
 q_+=e^{-x-y},\qquad q_-=e^{-x+y},\qquad0<q_+<q_-<1.
\]
Pour une involution \(R\), avec \(P_\pm=(I\pm R)/2\), l’opérateur
des deux canaux est

\[
 T=q_+P_++q_-P_-.
\]
Ses moments scalaire et orienté s’expriment sous la forme

\[
 c_n=q_-^n+q_+^n,\qquad s_n=q_-^n-q_+^n.
\]
Multiplier ces expressions prouve

\[
 c_n^2-s_n^2=4e^{-2nx},\quad
 c_{m+n}=\frac{c_mc_n+s_ms_n}{2},\quad
 s_{m+n}=\frac{s_mc_n+c_ms_n}{2}.
\]
Les deux suites satisfont

\[
 X_{n+2}=c_1X_{n+1}-e^{-2x}X_n,\qquad X\in\{c,s\}.
\]
Ces relations montrent que les hauteurs partagent un même opérateur.
On ne définit pas de collection différente de moments pour chaque espèce.


\begin{proposition}[Semi-groupe des hauteurs
]
\[
 \langle90,120\rangle=\{90,120\}\cup\{30r:r\ge6\}.
\]
\end{proposition}
\begin{proof}
Après division par \(30\), il suffit d’étudier \(\langle3,4\rangle\).
Les entiers \(6,7,8\) sont \(3+3,3+4,4+4\) ; ajouter \(3\) produit
tous les entiers suivants. En dessous de \(6\), les seuls
éléments positifs sont \(3,4\).

\end{proof}

Les hauteurs \(360=4\cdot90=3\cdot120\) illustrent la différence entre
égalité d’évaluation et mémoire de composition. Le quotient algébrique
identifie les deux exposants ; le registre conserve si le transport
a employé quatre étapes d’un type ou trois de l’autre.


La spécialisation raffinée a pour exposant

\[
 \log\chi_{\mathrm{ref}}=n_Ax+\frac{s_C}{6}y+k_{\Delta_4}\Delta_4
 +b_{\zeta}\zeta_{270}+\sum_{h\in\langle90,120\rangle}N_hs_h,
 \qquad N_h=\eta_h+b_{120}\mathbf1_{\{h=120\}}.
\]
La paire \((\eta_{120},b_{120})\) est conservée même si l’évaluation utilise
sa somme. On conserve également séparés \(s_{270}\) et
\(\zeta_{270}\) : le premier est un moment bicouche et le second une
expression quadratique du raffinement.


\begin{proposition}[Convergence de l’évaluation raffinée
]
Si \(\sum_h|N_hs_h|<\infty\), le caractère raffiné définit un scalaire
positif fini et ses évaluations tronquées convergent vers lui.

\end{proposition}
\begin{proof}
L’hypothèse donne la convergence absolue de l’exposant. La continuité de
l’exponentielle permet de passer à la limite ; l’exponentielle d’un réel fini est
strictement positive.

\end{proof}

La convergence évalue une suite de coefficients déjà spécifiée.
La sélection de ces coefficients par la route est une opération différente,
et doit accompagner toute évaluation attribuée à une espèce.


La notation \(\chi_{\mathrm{full}}\) conservée dans certaines formules
d’évaluation des secteurs désigne ce même caractère
\(\chi_{\mathrm{ref}}\), avec les mêmes arguments et coefficients.
Le changement de symbole n’ajoute ni autre spécialisation ni autre tour.


\subsection{Lecteurs de déplacement et d’autocorrélation
}

La même route admet deux lectures du retour :

\[
 K_D=\left(D_{27}+D_{36},D_1,\frac{D_4-D_3}{2}\right),\qquad
 K_\Omega=(\Omega,54,P_6).
\]
La première application conserve les défauts de déplacement ; la seconde, la mémoire de
retour et l’autocorrélation. Dans le recensement complet de \(324\) orientations,
il y a \(14\) profils de la première, \(9\) de la seconde et \(40\) paires
conjointes. Les \(18\) coïncidences ont pour valeur \((80,54,6)\).
On n’en infère pas l’égalité des deux opérateurs sur tout l’atlas.


Pour la route électronique, les déplacements sont

\[
 (D_1,D_3,D_4,D_{27},D_{36})=(54,56,68,48,32).
\]
La substitution produit \(K_D=(80,54,6)\), coïncidant avec
\(K_\Omega\). Son exposant est

\[
 \kappa_e=80-54\alpha+6\Delta_4.
\]
La trajectoire électronique est construite dans
\eqref{vi:dep:trayectoria-electron} ; ses lecteurs de déplacement et
d’autocorrélation sont évalués, respectivement, dans
\eqref{vi:dep:KD} et \eqref{vi:dep:KOmega}. La coïncidence
\(K_D=K_\Omega=(80,54,6)\) appartient au théorème de fermeture bêta de la
sous-section~\ref{vi:dep:beta}. Ici, cette exception centrale est conservée et
l’opération s’étend aux fibres de l’atlas.


\subsection{Opérateur sur une fibre et section dimensionnelle
}

Soit \(\mathscr H_M\) l’espace de Hilbert fini de base orthonormale
\((e_\gamma)\) étiquetée par les routes d’une multisection. Pour
\(r\in\{D,\Omega\}\), l’opérateur de retour est défini par
\(B_M^re_\gamma=\kappa_\gamma^re_\gamma\), avec des exposants réels ;
il est donc autoadjoint. Soit \(M_M^{(0)}\) l’opérateur basal positif,
dont la connexion avec le caractère est fixée ci-dessous.
Le rapport adimensionnel des
sections d’action est \(R_{\mathrm{act}}>0\). On définit

\[
 M_M^{\beta,r}=R_{\mathrm{act}}^{B_M^r/2}
 M_M^{(0)}R_{\mathrm{act}}^{B_M^r/2}.
\]

\begin{proposition}[Positivité et restriction scalaire
]
L’opérateur \(M_M^{\beta,r}\) est positif. Si les deux opérateurs sont
diagonaux dans la base de routes, chaque valeur basale \(m_\gamma^{(0)}\)
se transforme en

\[
 m_\gamma^\beta=m_\gamma^{(0)}R_{\mathrm{act}}^{\kappa_\gamma}.
\]
\end{proposition}
\begin{proof}
Soit \(S=R_{\mathrm{act}}^{B_M^r/2}\), autoadjoint et inversible.
Pour tout \(v\),
\(\langle v,SM_M^{(0)}Sv\rangle=\langle Sv,M_M^{(0)}Sv\rangle\ge0\).
Dans la base diagonale, les deux facteurs \(S\) multiplient chaque valeur
par \(R_{\mathrm{act}}^{\kappa_\gamma/2}\) ; leur produit donne la formule.

\end{proof}

La section absolue de masse s’obtient à partir de l’action et de l’horloge :

\[
 \omega_0=\frac{2\pi}{108t_0},\qquad
m_{\mathrm{clk}}=\hbar_{\mathrm{ret}}\omega_0c_{\mathrm{int}}^{-2}.
\]
Dans la coordinatisation normalisée par rapport à l’électron, l’opérateur basal précédent est

\[
 M_M^{(0)}e_\gamma
 =m_e^{(0)}\,
   \chi_{\mathrm{ref}}(\sigma_\gamma-\sigma_e,
                       \eta_\gamma-\eta_e)e_\gamma.
\]
Ici, \(m_e^{(0)}\) est la réalisation de la coordonnée électronique
basale, avec son préfacteur, son raffinement et sa base dimensionnelle.
Si \(E_e^{(0)}=\varepsilon_e^{(0)}\mathcal U_E\), alors

\[
 m_e^{(0)}=\varepsilon_e^{(0)}\mathcal U_Ec_{\mathrm{int}}^{-2}
 =b_e\,m_{\mathrm{clk}},\qquad
 b_e=\varepsilon_e^{(0)}\frac{\mathcal U_E}{E_{\mathrm{clk}}}.
\]
Dans la coordinatisation \(\mathcal U_E=E_{\mathrm{clk}}\), on obtient
\(b_e=\varepsilon_e^{(0)}\) ; dans une autre coordinatisation, le quotient
de bases indiqué est conservé. Si les coordonnées de tour sont déjà
relatives à l’électron, on omet une seconde soustraction de \(\eta_e\).
Cette égalité relie l’opérateur à la signature de chaque route et conserve
explicitement la normalisation électronique. Un choix différent de
section de référence doit également transporter cette normalisation.
L’échelle de la constante réduite de Planck appartient à cette construction
dimensionnelle. Le caractère et le facteur de retour agissent ensuite comme
scalaires. L’électron conserve sa structure : il ne s’identifie pas au
quantum chronologique \(m_{\mathrm{clk}}\).


Pour deux routes dans la même coordinatisation absolue,

\[
 \frac{m_Y^\beta}{m_X^\beta}
 =\chi_{\mathrm{ref}}(\sigma_Y-\sigma_X,\eta_Y-\eta_X)
 R_{\mathrm{act}}^{\kappa_Y-\kappa_X}.
\]
La section dimensionnelle commune s’annule. La différence des exposants de
retour ne s’annule pas sauf si elle est nulle. Cette identité conserve la
distinction entre échelle absolue et structure relative d’une collection.


\subsection{Le secteur de propagation nulle
}

L’exponentielle prend des valeurs strictement positives. La coordinatisation nulle se
construit avant ce caractère, dans la réalisation déployée

\[
 \mathcal A_{\mathrm{sp}}=\mathbb R[R]/(R^2-1),\quad
 z=r_+P_++r_-P_-,\quad N(z)=r_+r_-.
\]
Les directions \(\mathbb RP_+\) et \(\mathbb RP_-\) ont une norme
nulle. Une propagation transportée dans l’une d’elles appartient à ce secteur ;
la lecture matérielle bicouche emploie le couplage des deux. De cette manière,
le photon et le secteur gluonique conservent leur place dans la théorie même
s’ils ne constituent pas une ligne du caractère positif. La réalisation de leurs
représentations est exposée avec les secteurs correspondants.


\section{Action, gravité et information thermique}
\label{sec:gravity}
\subsection{Période dodécaphasique et réalisation circulaire de l'action}
Le calendrier déjà construit détermine $T_{108}=108t_0$.
Dans sa réalisation circulaire,
\[
\omega_{108}=\frac{2\pi}{108t_0},\qquad
R_{108}=\frac{c_{\rm int}}{\omega_{108}}
=\frac{54}{\pi}\ell_0,\qquad \ell_0=c_{\rm int}t_0.
\]
L’unité $\ell_0$ et les coordinatisations de temps et de vitesse restent
explicites ; la coordinatisation SI ne sélectionne pas le cycle.
La composition longitudinale de la section~\ref{md:rigidez:longitud}
fixe dans cette édition $R_{108}=L_\alpha$ et
$t_0=\pi_{\rm HMT}L_\alpha/(54c_{\rm int})$.
L’argument circulaire suivant est conservé comme spécialisation
du caractère unitaire du théorème~\ref{md:rigidez:teorema}.
Pour chaque section orientée d’action $a\in\{\mathrm{pre},\mathrm{ret}\}$,
\[
E_{108,a}=\hbar_a\omega_{108},\qquad
m_{108,a}=\frac{E_{108,a}}{c_{\rm int}^2}.
\]
Les longueurs réduite et complète sont respectivement
\[
\frac{\hbar_a}{m_{108,a}c_{\rm int}}=R_{108},
\qquad
\frac{h_a}{m_{108,a}c_{\rm int}}=2\pi R_{108}.
\]
La même action et la même durée interviennent dans les deux lectures.

\begin{proposition}[Sélecteur de coïncidence des rayons]
Dans la famille dimensionnelle
\[
G_a(\vartheta)=\vartheta\frac{c_{\rm int}^5t_0^2}{\hbar_a},
\qquad \vartheta>0,
\]
la condition de la réalisation circulaire
$G_a(\vartheta)m_{108,a}/c_{\rm int}^2=R_{108}$
admet l'unique solution
\[
\vartheta^\circ_{108}=(54/\pi)^2.
\]
\end{proposition}
\begin{proof}
Le premier rayon est
$r_g(\vartheta)=\vartheta(\pi/54)\ell_0$.
Le comparer à $(54/\pi)\ell_0$, avec $\ell_0>0$, réduit
l'égalité à une équation linéaire strictement croissante.
\end{proof}
La condition de coïncidence est une partie explicite de cette réalisation ;
le théorème résout son coefficient sans utiliser de valeur métrologique de $G$.
Le rayon employé ici est $Gm/c^2$, de sorte que le facteur deux
de la convention $2Gm/c^2$ n'est pas transféré par inadvertance.

\begin{proposition}[Action et gravité réciproques, aire conservée]
Dans les deux coordinatisations précédentes,
\[
\frac{G_{\rm pre}}{G_{\rm ret}}=R_{\rm act}^{-1},\qquad
\hbar_aG_a=\vartheta^\circ_{108}c_{\rm int}^5t_0^2,
\qquad
\ell_P^2=\frac{\hbar_aG_a}{c_{\rm int}^3}
=\vartheta^\circ_{108}\ell_0^2.
\]
\end{proposition}
\begin{proof}
La définition de $G_a$ contient $\hbar_a^{-1}$ et le reste des facteurs
est commun aux deux coordinatisations. Prendre le quotient et le produit donne
les trois égalités.
\end{proof}

\subsection{Définition structurelle et fonction de la constante gravitationnelle}
\label{sec:ii-definicion-gravedad}

\begin{definition}[Section gravitationnelle contragrédiente de l'action]
Pour une section positive \(\hbar_a\), une vitesse
\(c_{\rm int}>0\) et le rayon chronologique \(R_{108}\) de la
section~\ref{sec:gravity}, la section gravitationnelle circulaire est
\begin{equation}
 G_a=\frac{c_{\rm int}^{3}R_{108}^{2}}{\hbar_a},
 \qquad R_{108}=\frac{54}{\pi}\ell_0,
 \qquad \ell_0=c_{\rm int}t_0.
 \label{eq:ii-definicion-gravedad-seccion}
\end{equation}
C'est la section réciproque de l'action qui réalise la coïncidence
entre le rayon gravitationnel réduit du cycle et son rayon
chronologique, dans la convention \(r_g=Gm/c_{\rm int}^2\).
\end{definition}

La proposition de coïncidence des rayons de la
section~\ref{sec:gravity} démontre son unicité dans la collection
\(G_a(\vartheta)=\vartheta c_{\rm int}^{5}t_0^2/\hbar_a\).
La condition circulaire fixe \(\vartheta=(54/\pi)^2\) et conduit
à \eqref{eq:ii-definicion-gravedad-seccion}. Cette condition et les
sections dimensionnelles positives font partie de la réalisation.
La section~\ref{md:rigidez} construit sa longueur et réunit la
preuve d’unicité pour le caractère positif général, avec
$R_{108}=L_\alpha$. Ainsi le sélecteur circulaire et le lecteur
longitudinal conservent une même normalisation.
En les maintenant communes aux deux orientations, on obtient
\begin{equation}
 \frac{G_{\mathrm{pre}}}{G_{\mathrm{ret}}}
 =R_{\mathrm{act}}^{-1},
 \qquad
 \frac{G_a\hbar_a}{c_{\rm int}^{3}}
 =R_{108}^{2}=\ell_P^2.
 \label{eq:ii-definicion-gravedad-area}
\end{equation}
La preuve est la substitution directe de
\eqref{eq:ii-definicion-gravedad-seccion} ; le changement d'action
est compensé par le changement réciproque de \(G_a\).

\section{Information réduite, énergie et coordonnées thermiques}

\subsection{Lecteur, mémoire et récupération de l'état}

Fixons un niveau fini du développement et notons \(\Omega\) l'
ensemble des états enrichis admissibles. Le registre conserve feuillet,
orientation, résidu, quotient, retenue, route, phase et mémoire dans les
domaines où intervient chaque coordonnée. Soient
\[
 q:\Omega\longrightarrow Y,\qquad
 M:\Omega\longrightarrow\mathcal M
\]
la représentation visible et la mémoire retenue, respectivement. Pour une
distribution \(p\) sur \(\Omega\), nous écrivons
\(\mathsf H(p)=-\sum_xp_x\log p_x\), avec \(0\log0=0\).
Le logarithme s’applique à des probabilités adimensionnelles.

\begin{definition}[Récupération généalogique d’une représentation]
Le couple \(\widehat q=(q,M)\) est un lecteur récupérable sur le support
\(\Omega_p\) lorsqu'il existe
\[
 d:\widehat q(\Omega_p)\longrightarrow\Omega_p,
 \qquad d\circ\widehat q=\operatorname{id}_{\Omega_p}.
\]
L'information de fibre de \(q\) est
\(\mathsf H_{\mathrm{fib}}(q;p)=\mathsf H(X\mid q(X))\), où
\(X\) a pour loi \(p\).
\end{definition}

\begin{proposition}[Identité de récupération et coût de la réduction]
Pour un lecteur récupérable,
\begin{align}
 \mathsf H(X\mid q(X),M(X))&=0,\\
 \mathsf H(X)&=\mathsf H(q(X))+\mathsf H(X\mid q(X)),\\
 \mathsf H(X\mid q(X))&=\mathsf H(M(X)\mid q(X)).
 \label{eq:ii-kb-aut-fibra}
\end{align}
\end{proposition}
\begin{proof}
L'application \(d\) restitue \(X\) à partir du couple \((q(X),M(X))\),
de sorte que la première entropie conditionnelle est nulle. Comme \(q\) est
une fonction de \(X\), la règle de la chaîne donne la seconde égalité.
Pour la troisième, on développe de deux manières
\(\mathsf H(X,M(X)\mid q(X))\). L'une fournit
\(\mathsf H(X\mid q(X))\), puisque \(M\) est une fonction de \(X\) ;
l'autre fournit
\(\mathsf H(M(X)\mid q(X))+\mathsf H(X\mid q(X),M(X))\).
\end{proof}

La division euclidienne des feuillets d'APP montre le mécanisme
arithmétique élémentaire. Avec \(i,j\in\{1,\ldots,9\}\), on conserve
\[
 i+j=9q^+_{ij}+R^+_{ij},\qquad
 ij=9q^\times_{ij}+R^\times_{ij},\qquad
 1\leq R^\sigma_{ij}\leq9\quad(\sigma\in\{+,\times\}).
\]
La paire résidu--quotient récupère la valeur entière de l'opération. En additionnant les différences sur les quatre-vingt-une positions, on obtient
\[
 2025-810=54+9\cdot129.
\]
La mémoire du quotient restaure la contribution que n’exprime pas la
réduction résiduelle. La récupérabilité de l’état complet requiert,
en outre, de conserver les coordonnées de position et de route pertinentes ;
le couple résidu--quotient identifie ici le résultat de l’opération.

\subsection{Mise à jour réversible, décision tritique et effacement}

Soit \(U:\Omega\to\Omega\) la mise à jour TPK dans un domaine où le transport enrichi possède une inverse. La loi transportée est \(p'_y=p_{U^{-1}y}\). La bijection conserve le multiensemble des probabilités et, par conséquent,
\begin{equation}
 \mathsf H(U(X))=\mathsf H(X),\qquad
 \mathsf H(X\mid U(X))=0.
 \label{eq:ii-kb-aut-reversible}
\end{equation}
Si la sortie observée est \(q(U(X))\), l’information non exprimée
est quantifiée par \(\mathsf H(X\mid q(U(X)))\). L’état complet,
sa représentation et la réinitialisation d’un registre sont ainsi trois applications distinctes.

Pour les trois états orientés du TRIT, soit
\(p=(p_{-1},p_0,p_{+1})\) une distribution normalisée. Son information est
\begin{equation}
 I_{\mathrm{TRIT}}(p)=-\sum_{a=-1}^{1}p_a\log p_a,
 \qquad 0\leq I_{\mathrm{TRIT}}(p)\leq\log3.
 \label{eq:ii-kb-aut-trit}
\end{equation}
L'extrémité supérieure correspond à \(p_a=1/3\). En effet,
\[
 D(p\Vert(1/3,1/3,1/3))
 =\sum_ap_a\log(3p_a)=\log3-I_{\mathrm{TRIT}}(p)\geq0;
\]
l'inégalité s'obtient à partir de la convexité de \(t\log t\), et l'égalité caractérise la distribution uniforme. L'extrémité inférieure correspond à une distribution concentrée sur un état.

La réalisation idéale de Landauer utilise une mémoire aux états logiques
énergétiquement dégénérés, un environnement isotherme à température \(\Theta\)
et un bilan qui conserve la sortie de l'application logique et comptabilise les
registres auxiliaires. Dans ce modèle, la diminution de l'entropie logique est
\(k_B[\mathsf H(X)-\mathsf H(F(X))]
=k_B\mathsf H(X\mid F(X))\) pour une application déterministe \(F\).
Le bilan entropique exige de transférer au milieu au moins cette quantité ;
son expression énergétique est \(Q\geq k_B\Theta\mathsf H(X\mid F(X))\).
La distinction entre évolution logique réversible et effacement physique
correspond à la réalisation de Landauer--Bennett \cite{bennett2003}.
La mise à jour inversible
\eqref{eq:ii-kb-aut-reversible} a une contribution logique nulle ; la
réinitialisation d'un triplet uniforme exige le transfert de \(\log3\) nats
au milieu et satisfait \(Q\geq k_B\Theta\log3\).
L'énergie de contrôle et la dissipation d'une réalisation physique
concrète appartiennent à son bilan thermodynamique supplémentaire.

\subsection{Action, fréquence et énergie de décision}

Soient \(\mathcal L_{\rm act}\) et \(\mathcal L_T\) les droites
dimensionnelles d'action et de durée ; la droite énergétique est
\(\mathcal L_E=\mathcal L_{\rm act}\otimes\mathcal L_T^*\).
Le développement précédent fournit les sections positives
\(\hbar_a\in\mathcal L_{\rm act}\), avec
\(a\in\{\mathrm{pre},\mathrm{ret}\}\), et la durée élémentaire
\(t_0\in\mathcal L_T\). Le calendrier complet fournit
\[
 T_{108}=108t_0,\qquad
 \omega_{108}=\frac{2\pi}{T_{108}},\qquad
 h_a=2\pi\hbar_a.
\]
La contraction dimensionnelle entre action et fréquence définit
\begin{equation}
 \epsilon_a:=\hbar_a\omega_{108}
 =\frac{h_a}{108t_0}
 =\frac{\hbar_a}{t_P},\qquad
 t_P=\frac{54}{\pi}t_0.
 \label{eq:ii-kb-aut-energia}
\end{equation}
Le dernier membre emploie la même période circulaire réduite que la section gravitationnelle. Les trois expressions correspondent à une unique section d'énergie, avec une orientation \(a\) maintenue dans tous les facteurs. L'information normalisée d'une décision uniforme produit la section d'énergie par nat
\begin{equation}
 \varepsilon_a:=\frac{\epsilon_a}{\log3}\in\mathcal L_E^+.
 \label{eq:ii-kb-aut-energia-nat}
\end{equation}
L'action précédemment construite détermine l'échelle d'énergie lorsque
la durée est fixée ; le cardinal des alternatives du TRIT détermine sa
normalisation informationnelle. La valeur de \(\log3\) intervient par le
dénombrement et la mesure déclarés dans \eqref{eq:ii-kb-aut-trit}.

\subsection{Définition dimensionnelle de Boltzmann et sélection des coordonnées}

Soit \(\mathcal L_\Theta\) la droite orientée de température. L'espace dimensionnel d'entropie est \(\mathcal L_{\rm ent}=\mathcal L_E\otimes\mathcal L_\Theta^*\). Une application linéaire positive \(k_B:\mathcal L_\Theta\to\mathcal L_E\) est, de manière équivalente, un élément positif de \(\mathcal L_{\rm ent}\). C'est pourquoi le même objet peut agir sur une température ou multiplier une information sans dimension.

\begin{definition}[Normalisation thermique de la décision tritique]
Un couple thermique compatible avec la section orientée \(a\) se compose d'
un transducteur positif \(k_B\) et d'une section positive
\(\Theta_{{\rm clk},a}\in\mathcal L_\Theta\) tels que
\begin{equation}
 k_B(\Theta_{{\rm clk},a})=\varepsilon_a,
 \qquad
 k_B\Theta_{{\rm clk},a}\log3=\epsilon_a.
 \label{eq:ii-kb-aut-par-termico}
\end{equation}
La constante de Boltzmann remplit la fonction de transducteur
énergie--température de l'information tritique dans cette réalisation.
\end{definition}

\begin{theorem}[Unicité relative à la section thermique et covariance dimensionnelle]
Une section positive \(\Theta_{{\rm clk},a}\) étant fixée, il existe une unique application linéaire positive qui satisfait \eqref{eq:ii-kb-aut-par-termico}. Pour des bases positives \(u_E,u_\Theta\), si
\[
 \epsilon_a=\epsilon_a^{\rm num}u_E,\quad
 \Theta_{{\rm clk},a}=\theta_a u_\Theta,\quad
 k_B(u_\Theta)=b\,u_E,
\]
sa coordonnée est
\begin{equation}
 b=\frac{\epsilon_a^{\rm num}}{\theta_a\log3}.
 \label{eq:ii-kb-aut-coordenada}
\end{equation}
Sous \(u'_E=s_Eu_E\), \(u'_\Theta=s_\Theta u_\Theta\), avec
\(s_E,s_\Theta>0\), on a
\[
 b'=\frac{s_\Theta}{s_E}b,\qquad
 \theta'_a=\frac{\theta_a}{s_\Theta},\qquad
 (\epsilon_a^{\rm num})'=\frac{\epsilon_a^{\rm num}}{s_E}.
\]
\end{theorem}
\begin{proof}
Une section non nulle engendre une droite. Prescrire son image détermine
l'unique application linéaire ; la positivité des deux sections prouve
sa positivité. Substituer les coordonnées dans
\eqref{eq:ii-kb-aut-par-termico} donne
\(b\theta_a\log3=\epsilon_a^{\rm num}\). Les trois lois de
transformation s'obtiennent en exprimant les mêmes sections et la même
application dans les nouvelles bases.
\end{proof}

Une même constante de Boltzmann doit agir dans les deux orientations.
Puisque \(t_P\) est commun, la compatibilité simultanée équivaut à
\begin{equation}
 \frac{\Theta_{{\rm clk},{\rm pre}}}{\Theta_{{\rm clk},{\rm ret}}}
 =\frac{\epsilon_{\rm pre}}{\epsilon_{\rm ret}}
 =\frac{\hbar_{\rm pre}}{\hbar_{\rm ret}}.
 \label{eq:ii-kb-aut-orientaciones}
\end{equation}
En effet, un isomorphisme de droites conserve les quotients de sections. Réciproquement, si l'égalité précédente est satisfaite, l'isomorphisme déterminé dans une orientation satisfait également la normalisation de l'autre. Ainsi, le retour modifie la section thermique du cycle dans la même proportion que l'énergie, tandis que le transducteur reste commun aux deux lectures.

L’équation précédente explicite le choix requis pour exprimer une
valeur individuelle. Si la section thermique n’a pas encore été sélectionnée
par un lecteur indépendant, pour chaque \(\lambda>0\) le couple
\[
 (k_B,\Theta_{{\rm clk},a})
 \longmapsto(\lambda k_B,\lambda^{-1}\Theta_{{\rm clk},a})
\]
conserve \(\epsilon_a\). Par conséquent, la présentation d’une
coordonnée individuelle sélectionnée en interne exige de donner le lecteur de
\(\Theta_{{\rm clk},a}\), les bases et leur transport dimensionnel.
Les équations d’action et d’information restent déterminées avant
cette spécification. L’identification ultérieure des bases au
joule et au kelvin appartient à la représentation métrologique du résultat.

\subsection{Entropie dimensionnelle et information de l'état réduit}

Pour une distribution de routes admises \(p\), ou pour un état réduit de matrice de densité \(\rho\), soient
\[
 s(p)=-\sum_xp_x\log p_x,\qquad
 s(\rho)=-\operatorname{Tr}(\rho\log\rho).
\]
Ses représentations dimensionnelles sont
\begin{equation}
 S(p)=k_Bs(p),\qquad S(\rho)=k_Bs(\rho)
 \quad\text{en }\mathcal L_{\rm ent}.
 \label{eq:ii-kb-aut-entropia}
\end{equation}
L'évaluation en \(\Theta\) produit une énergie
\(\Theta S=k_B\Theta\,s\). Pour le triplet uniforme à la température
du cycle, \eqref{eq:ii-kb-aut-par-termico} donne exactement
\[
 S_{\rm TRIT}=k_B\log3,\qquad
 \Theta_{{\rm clk},a}S_{\rm TRIT}=\epsilon_a.
\]

% Recuperación documental para el artículo V; RESULTADO_RECUPERADO.
% Propietario: /Users/ruben/Documents/New project/output/REVISION_EDITORIAL_HMT_MD_20260905/01_FUENTES/INTEGRAL/tree/New project/output/ENSAMBLADO_CAUSAL_RESTAURADO_HMT_MD_20260905/fuente/manuscrito/sections/md/delta_127/04b_completaciones_cuanticas.tex
% Líneas de origen: 1--334.
% REV02: construcción algebraica; la aplicación estadística se reúne en 40.
% Correspondencia editorial: gestion/CONSOLIDACION_PAULI_FOCK_REV02.json.
\section{Construction algébrique des complétions quantiques sur APP}
\label{v:rec:chap:completaciones-cuanticas}
\label{v:sec:completaciones-algebraicas}

APP et TPK fournissent des états orientés, un transport et une parité de feuillet.
Leur réalisation linéaire permet de construire les algèbres symétrique et extérieure,
avec leurs opérateurs de création, d’annihilation et d’occupation. Cette section
établit ces constructions algébriques, le secteur cellulaire des champs et
le couplage réversible de mesure. L’application des spectres d’occupation à l’équilibre est développée dans la section~\ref{v:rec:ocupacion:section:02}.
Chaque partie conserve les conditions de la représentation qu’elle utilise.

\subsection{Espace linéaire des états associé à une réalisation de particule}
\label{v:rec:completaciones:section:02}

Soit \(S\) un ensemble fini de classes de chemins TPK à la profondeur choisie
et soit
\[
 \mathcal H_1=\ell^2(S)
\]
l’espace de Hilbert de base orthonormale \((e_s)_{s\in S}\). Un choix de
poids positifs \(E_s\) définit l’opérateur diagonal
\[
 He_s=E_se_s.
\]
Le choix de \(S\), de son quotient par les symétries et des poids fait partie
du modèle. Une fois ce choix fixé, les complétions multiparticulaires sont canoniques.

\subsection{Complétions symétrique et extérieure}
\label{v:rec:completaciones:section:03}

On définit
\[
 \mathcal F_+(\mathcal H_1)=
 \bigoplus_{n\geq0}\operatorname{Sym}^n\mathcal H_1,
 \qquad
 \mathcal F_-(\mathcal H_1)=
 \bigoplus_{n\geq0}\bigwedge^n\mathcal H_1.
\]
La première est la complétion bosonique ; la seconde, la fermionique. Dans les deux,
le terme de somme directe de degré zéro est l’espace du vide. Les puissances extérieures
sont munies du produit scalaire pour lequel les produits extérieurs ordonnés de vecteurs
d’une base orthonormale sont orthonormaux ; les puissances symétriques sont utilisées
avec la normalisation usuelle des opérateurs de création et d’annihilation.

\subsection{Graduation de feuillet et algèbre d’échange}
\label{v:rec:completaciones:section:04}

Supposons que le feuillet orienté définit un caractère additif
\[
 p:\operatorname{Mor}(\mathsf P_{\rm TPK})\longrightarrow\mathbb Z/2\mathbb Z,
 \qquad p(\gamma\eta)=p(\gamma)+p(\eta).
\]
Sur l’espace linéaire gradué correspondant, on introduit la symétrie
\[
 \tau(v\otimes w)=(-1)^{p(v)p(w)}w\otimes v.
\]

\begin{theorem}[Complétion statistique déterminée par la parité]
\label{v:rec:completaciones:statement:02}
Le quotient de l’algèbre tensorielle par les relations
\[
 v\otimes w-(-1)^{p(v)p(w)}w\otimes v
\]
est symétrique sur le secteur pair et extérieur sur le secteur impair. Deux états
impairs identiques ont un produit nul ; les états pairs admettent une occupation
arbitraire.
\end{theorem}

\begin{proof}
Si \(p(v)=p(w)=0\), la relation est \(vw=wv\). Si
\(p(v)=p(w)=1\), elle est \(vw=-wv\) ; en prenant \(v=w\) et en travaillant en
caractéristique zéro, on obtient \(v^2=0\). Les secteurs mixtes commutent sans
signe supplémentaire. C’est la propriété universelle de l’algèbre symétrique
graduée.
\end{proof}

\subsection{Opérateurs canoniques et idempotence de l’occupation}
\label{v:pauli:car-construccion}

Le caractère de parité détermine la complétion ; le produit de cette
complétion détermine ensuite l’algèbre d’opérateurs. Dans le
secteur extérieur, la création \(a_s^\dagger\) est le produit extérieur par
\(e_s\), et \(a_s\) est son adjoint.

\begin{theorem}[Principe d’exclusion de Pauli dans la complétion extérieure]
\label{v:rec:completaciones:statement:01}
\label{v:rec:thm:principio-exclusion-pauli}
Pour tout mode \(s\in S\), les opérateurs de création et d’annihilation de la
complétion extérieure satisfont
\begin{equation}
 \{a_s,a_t^\dagger\}=\delta_{st}I,
 \qquad
 \{a_s,a_t\}=\{a_s^\dagger,a_t^\dagger\}=0.
 \label{v:pauli:eq:car}
\end{equation}
En particulier, l’opérateur d’occupation
\(N_s=a_s^\dagger a_s\) est un projecteur :
\begin{equation}
 N_s^2=N_s,
 \qquad \Spec(N_s)\subseteq\{0,1\}.
 \label{v:pauli:eq:proyeccion-ocupacion}
\end{equation}
Les valeurs propres \(0\) et \(1\) correspondent, respectivement, à l’absence
et à la présence du mode \(s\).
\end{theorem}

\begin{proof}
Sur un produit extérieur ordonné, \(a_s\) supprime le facteur \(e_s\), lorsqu’il est
présent, avec le signe de sa position ; s’il est absent, il donne zéro. Composer
cette opération avec le produit extérieur démontre les relations
\eqref{v:pauli:eq:car} : pour \(s=t\), les deux compositions couvrent les cas
complémentaires d’absence et de présence ; pour \(s\ne t\), les signes
d’échange s’annulent. L’antisymétrie implique
\(e_s\wedge e_s=0\) et
\(a_s^2=(a_s^\dagger)^2=0\). Alors
\begin{equation}
 N_s^2=a_s^\dagger a_s a_s^\dagger a_s
 =a_s^\dagger(I-a_s^\dagger a_s)a_s
 =a_s^\dagger a_s=N_s.
 \label{v:pauli:eq:prueba-idempotencia}
\end{equation}
L’opérateur est autoadjoint, et toute valeur propre d’un idempotent satisfait
\(\lambda^2=\lambda\). Son spectre est donc contenu dans
\(\{0,1\}\). Le vide et l’état \(e_s\) réalisent les deux valeurs
\cite{Pauli1925,Dirac1926}.
\end{proof}

Dans le secteur symétrique, les opérateurs bosoniques \(b_s^\dagger,b_s\)
agissent sur la base normalisée des occupations par les facteurs
\(\sqrt{m_s+1}\) et \(\sqrt{m_s}\), respectivement. De leurs compositions,
on obtient
\begin{equation}
 [b_s,b_t^\dagger]=\delta_{st}I,
 \qquad [b_s,b_t]=[b_s^\dagger,b_t^\dagger]=0.
 \label{v:pauli:eq:ccr}
\end{equation}
Ces identités sont comprises sur le domaine algébrique des particules
finies, stable sous les opérateurs considérés.

Le théorème~\ref{v:rec:completaciones:statement:02} transforme une parité compositionnelle TPK en une règle exacte
d’occupation. Pour obtenir le théorème physique de spin--statistique, il faut
construire en outre une théorie locale lorentzienne positive et démontrer que la
parité de feuillet coïncide avec le spin sous les rotations. L’égalité ne se
déduit pas uniquement du fait que les deux objets prennent leurs valeurs modulo deux.

\section{Distributions de Bose--Einstein et de Fermi--Dirac}
\label{v:rec:ocupacion:section:02}

La formulation d’équilibre part du Hamiltonien et du nombre total
d’occupations,
\begin{equation}
 H=\sum_kE_k^{\rm occ} n_k,
 \qquad
 N=\sum_k n_k,
 \qquad
 K_\mu:=H-\mu N.
 \label{v:rec:eq:hamiltoniano-gran-canonico-rev3}
\end{equation}
Pour \(\beta>0\), l’état grand canonique est
\begin{equation}
 \rho_{\beta,\mu}
 =\frac{e^{-\beta K_\mu}}{\mathcal Z_{\beta,\mu}},
 \qquad
 \mathcal Z_{\beta,\mu}
 =\operatorname{Tr}(e^{-\beta K_\mu}).
 \label{v:rec:eq:estado-gran-canonico-rev3}
\end{equation}

\begin{theorem}[Factorisation grand canonique et lois d’occupation]
\label{v:rec:thm:factorizacion-gran-canonica-rev3}
Pour des niveaux d’énergie \(E_k^{\rm occ}\) de dégénérescences \(g_k\), la
trace fermionique se factorise comme
\begin{equation}
 \mathcal Z_{\rm FD}
 =\prod_k\bigl(1+e^{-\beta(E_k^{\rm occ}-\mu)}\bigr)^{g_k},
 \qquad
 \langle n_k\rangle_{\rm FD}
 =\frac{g_k}{e^{\beta(E_k^{\rm occ}-\mu)}+1}.
 \label{v:rec:eq:factorizacion-fd-rev3}
\end{equation}
Dans le secteur bosonique, sous la condition
\(\mu<\inf_kE_k^{\rm occ}\),
\begin{equation}
 \mathcal Z_{\rm BE}
 =\prod_k\bigl(1-e^{-\beta(E_k^{\rm occ}-\mu)}\bigr)^{-g_k},
 \qquad
 \langle n_k\rangle_{\rm BE}
 =\frac{g_k}{e^{\beta(E_k^{\rm occ}-\mu)}-1}.
 \label{v:rec:eq:factorizacion-be-rev3}
\end{equation}
\end{theorem}

\begin{proof}
L’opérateur \(K_\mu\) est une somme d’opérateurs d’occupation commutants, de
sorte que sa trace se factorise par modes. Dans chaque mode fermionique, on somme les
occupations \(0\) et \(1\), et le facteur correspondant est
\(1+u_k\), où
\(u_k=e^{-\beta(E_k^{\rm occ}-\mu)}\). Dans chaque mode bosonique, on somme toutes
les occupations non négatives et on obtient la série géométrique
\((1-u_k)^{-1}\) ; la condition sur \(\mu\) garantit \(u_k<1\). Les
puissances \(g_k\) incorporent la dégénérescence. Enfin,
\[
 \langle n_k\rangle
 =-\frac1\beta\frac{\partial}{\partial E_k^{\rm occ}}
   \log\mathcal Z
\]
produit les deux lois d’occupation
\cite{Bose1924,Einstein1924,Fermi1926,Dirac1926,Pathria2011}.
\end{proof}

% RESULTADO_RECUPERADO: funtor dimensional integral, eq:cZ-internos,
% eq:longitud-ruta y definición de ell_0; composición con artículo III.
% La prueba reunida y su control algebraico son específicos de esta contribución.
\subsection{Vitesse constitutive et réalisation cinématique de l’horloge}
\label{v:sec:enlace-velocidad-constitutiva}

La propagation radiative utilise la section de vitesse produite par le
transport angulaire du même état HMT. L’identification est établie
avant le choix des unités de représentation. Soient $x=A_{\mathrm{rad}}$ et
$y=C^*_{\mathrm{rad}}$ les coordonnées angulaires engendrées dans la
section~\ref{sec:ii-angulos}, avec les canaux
\eqref{eq:ii-canales}, et soit
$\mathsf R$ l’involution autoadjointe des deux feuillets, avec projecteurs
de rang un $P_+$ et $P_-$. Dans le domaine $x>|y|$, le transport et ses
canaux sont
\begin{equation}
 T=\exp(-xI-y\mathsf R)=q_+P_++q_-P_-,
 \qquad q_+=e^{-x-y},\quad q_-=e^{-x+y},\quad 0<q_\pm<1.
 \label{v:eq:velocidad-transporte}
\end{equation}
La réponse des blocs temporels de longueurs $90$ et $120$ est
\begin{equation}
 \mathsf V=(I-T^{90})(I-T^{120})^{-1}
       =r_+P_++r_-P_-,\qquad
 r_\pm=\frac{1-q_\pm^{90}}{1-q_\pm^{120}}.
 \label{v:eq:velocidad-respuesta}
\end{equation}
La construction part des canaux hérités ; les symboles $x,y$ et
$q_\pm$ désignent leurs coordonnées et non des paramètres ajustés à une vitesse
cible.

\begin{proposition}[Identité des réalisations constitutive et cinématique]
\label{v:prop:identidad-velocidad}
Dans un même système de coordonnées dimensionnelles, la section $c_{\mathrm{int}}$ de l’horloge
et la section $c_{\mathrm{vac}}$ de la réponse constitutive coïncident :
\begin{equation}
 c_{\mathrm{int}}=c_{\mathrm{vac}}
   =\widehat c\,u_C,
 \qquad
 \widehat c=(r_+r_-)^{-1}
            =\exp(-\mathcal E),
 \label{v:eq:identidad-velocidad}
\end{equation}
où
\begin{equation}
 \mathcal E=
 \log\!\bigl[(1-q_+^{90})(1-q_-^{90})\bigr]
 -\log\!\bigl[(1-q_+^{120})(1-q_-^{120})\bigr].
 \label{v:eq:velocidad-lector-simetrico}
\end{equation}
L’égalité conserve l’échange des feuillets.
\end{proposition}
\begin{proof}
Les facteurs qui apparaissent dans les logarithmes sont positifs. Par conséquent,
\[
 \mathcal E=\log(r_+r_-)
            =\log\det\mathsf V,
 \qquad \exp(-\mathcal E)=(\det\mathsf V)^{-1}.
\]
Le déterminant est calculé ici dans la réalisation à deux feuillets de rang
un. La définition cinématique $c_{\mathrm{int}}=\exp(-\mathcal E)u_C$
et la définition constitutive $c_{\mathrm{vac}}=(r_+r_-)^{-1}u_C$
sont, par conséquent, la même section. L’échange $P_+\leftrightarrow
P_-$ permute $r_+$ et $r_-$ et conserve leur produit. Enfin,
$\det T=e^{-2x}$ correspond au transport élémentaire, tandis que
$\det\mathsf V=r_+r_-$ correspond à sa réponse temporelle ; les deux
déterminants ont des fonctions distinctes dans cette composition.
\end{proof}

\subsubsection{Longueur de trajectoire et durée élémentaire}

Le repère dimensionnel positif $(u_S,u_C,u_T,u_Q)$ induit la base de longueur $u_L=u_Cu_T$. Pour une trajectoire enrichie $\gamma$ de $n=|\gamma|$ transitions et à canal constant, les lecteurs additifs sont
\begin{equation}
 T_{\mathrm{int}}(\gamma)=n u_T,
 \qquad L_{\mathrm{int}}(\gamma)=n\widehat c\,u_L.
 \label{v:eq:velocidad-longitud-reloj}
\end{equation}
Son additivité se rapporte au dénombrement et à la réalisation de chaque transition :
les chemins enrichis conservent en outre l’orientation et la mémoire.

\begin{corollary}[Normalisation de la longueur élémentaire]
\label{v:cor:longitud-elemental-constitutiva}
Pour $n>0$, $L_{\mathrm{int}}(\gamma)/T_{\mathrm{int}}(\gamma)
=c_{\mathrm{vac}}$. Dans la normalisation $t_0=u_T$,
\begin{equation}
 \ell_0=c_{\mathrm{int}}t_0=\widehat c\,u_L.
 \label{v:eq:longitud-elemental-constitutiva}
\end{equation}
Dans un système de coordonnées temporelles général $t_0=\tau_0u_T$, avec $\tau_0>0$,
l’expression correspondante est $\ell_0=\tau_0\widehat c\,u_L$.
\end{corollary}
\begin{proof}
La division de \eqref{v:eq:velocidad-longitud-reloj} simplifie l’entier
$n$ et utilise $u_L/u_T=u_C$. La formule élémentaire s’obtient
en multipliant la section de vitesse par $t_0$.
\end{proof}

La généralisation à un canal dépendant de l’arête conserve la formule
additive $L_{\mathrm{int}}(\gamma)=\sum_{e\in\gamma}\widehat c(e)u_L$.
Son quotient par $nu_T$ est la moyenne des vitesses de ces arêtes.
La réalisation radiative homogène emploie le canal constant déclaré
dans \eqref{v:eq:velocidad-longitud-reloj}.

\subsubsection{Changement de coordinatisation et normalisation adaptée}

\begin{proposition}[Covariance de l’identification]
\label{v:prop:covariancia-identidad-velocidad}
Soient $u_C'=d u_C$ et $u_T'=\eta u_T$, avec $d,\eta>0$.
La même vitesse et la même durée élémentaire ont pour coordonnées
\begin{equation}
 \widehat c'=d^{-1}\widehat c,
 \qquad \tau_0'=\eta^{-1}\tau_0,
 \qquad u_L'=d\eta u_L.
 \label{v:eq:velocidad-cambio-carta}
\end{equation}
En particulier, si $c_V=\widehat c_Vu_C^V$ et
$u_C^V=d u_C^{III}$, l’égalité des sections équivaut exactement à
$d\widehat c_V=\widehat c_{III}$.
\end{proposition}
\begin{proof}
On exprime chaque section dans les deux bases. On obtient
$\widehat c'u_C'=\widehat c u_C$ et
$\tau_0'u_T'=\tau_0u_T$. Leur produit satisfait
\[
 \tau_0'\widehat c' u_L'
 =\frac{\tau_0}{\eta}\frac{\widehat c}{d}
       d\eta u_L
 =\tau_0\widehat c u_L=\ell_0.
\]
La dernière équivalence résulte de la substitution de la relation entre les bases.
\end{proof}

Le choix adapté $d=\widehat c$ donne $u_C'=c_{\mathrm{vac}}$ et
$\widehat c'=1$. Ce choix exprime la section déjà construite dans une
base qui lui est adaptée. Le coefficient spectral dans le système de coordonnées interne reste
$(r_+r_-)^{-1}$ ; l’inversion spectrale utilise ce système de coordonnées interne.
Par exemple, en maintenant les bases d’action et de charge, les coordonnées
constitutives se transforment comme
\[
 \widehat\mu'=d\widehat\mu,
 \qquad \widehat\varepsilon'=d\widehat\varepsilon,
 \qquad
 \widehat\mu'\widehat\varepsilon'=(\widehat c')^{-2}.
\]
Dans le système de coordonnées adapté, $\widehat\mu'=r_-/r_+$ et
$\widehat\varepsilon'=r_+/r_-$, car les coordonnées internes sont
$\widehat\mu=r_-^2$ et $\widehat\varepsilon=r_+^2$.

\section{Restitution radiale orientée et transport symplectique}
\subsection{Récupération de Planck avec orientation conservée}

Avec les demi-axes étiquetés de l’ellipse, nous définissons
\(R_\star=\sqrt{b_S/a_S}\). Alors
\(q=R_\star^4=e^{-2\eta_{\rm el}}=(1-\xi)/(1+\xi)\),
où \(\xi=C^*/A\) et \(\eta_{\rm el}=\operatorname{artanh}\xi\).
Ce rapport radial conserve les marques des deux axes.

\begin{theorem}[Récupérateur radial de la section de retour]
\label{rec:thm:accion}
Dans la coordinatisation positive, si \(q=R_\star^4\), alors
\begin{align}
 \eta_{\rm ret}&=\alpha\frac{499-501q}{1+q},\label{rec:eq:eta-radio}\\
 \hbar_{\rm ret}&=\alpha\frac{499-501q}{1+q}\,10^{-34}\mathcal U_S,
 \qquad\hbar_{\rm pre}=H_5\,10^{-34}\mathcal U_S,\label{rec:eq:planck-radio}\\
 R_{\rm act}&=\frac{\hbar_{\rm pre}}{\hbar_{\rm ret}}
 =\frac{H_5(1+q)}{\alpha(499-501q)}.\label{rec:eq:ratio-radio}
\end{align}
La branche \(\eta_{\rm ret}>0\) équivaut à \(q<499/501\), à l’intérieur de \(q>0\).

\end{theorem}
\begin{proof}
Les identités angulaires donnent \(\eta_{\rm ret}=\alpha(500\xi-1)\). On substitue \(\xi=(1-q)/(1+q)\) ; le numérateur est \(499-501q\). Le dénominateur et \(\alpha\) sont positifs. La réalisation dimensionnelle et le quotient sont ceux des sections d'action préalablement construites. La décade \(-34\) et l'unité \(\mathcal U_S\) conservent leur dérivation déterminantale antérieure.
\end{proof}

Cette formule explicite la section de Planck dans le contre-angle. En termes de grandeurs d’action, la relation angulaire est
\[
 C^*=2\left(\frac{\hbar_{\rm ret}}{10^{-34}\mathcal U_S}+\alpha\right).
\]
La division par la base dimensionnelle extrait le coefficient adimensionnel ; elle n’ajoute pas une grandeur d’action à une constante sans unités.


Pour l’orientation angulaire \(\varepsilon\in\{-1,+1\}\), nous définissons
\[
 C^*_{\varepsilon}=2\varepsilon(\eta_{\rm ret}+\alpha),\quad
 \xi_\varepsilon=C^*_{\varepsilon}/A,\quad
 q_\varepsilon=\frac{1-\xi_\varepsilon}{1+\xi_\varepsilon}>0.
\]
L’opérateur de récupération sur les deux coordinatisations est
\begin{equation}
 \eta_{\rm ret}=\alpha\left(500\varepsilon
        \frac{1-q_\varepsilon}{1+q_\varepsilon}-1\right).
 \label{rec:eq:eta-orientada}
\end{equation}

\begin{corollary}[Invariance du récupérateur conjugué
]
La transformation \((\varepsilon,q)\mapsto(-\varepsilon,q^{-1})\) conserve \(\eta_{\rm ret}\), \(\hbar_{\rm ret}\) et \(R_{\rm act}\).

\end{corollary}
\begin{proof}
Le rapport \((1-q)/(1+q)\) change de signe lorsqu’on inverse \(q\). Le produit par \(\varepsilon\) reste invariant. On applique \eqref{rec:eq:eta-orientada} puis les représentations d’action.
\end{proof}

Si l’orientation est écartée et que \(\varepsilon=+1\) est indûment maintenu, la même expression produit \(\eta_{\rm ret}(q^{-1})=-\eta_{\rm ret}(q)-2\alpha\). Le registre distingue donc deux opérations algébriquement différentes. L’orientation angulaire \(\varepsilon\) ne s’identifie pas non plus sans autre précision au signe symplectique \(\sigma\) : l’échange des axes est antisymplectique, tandis que la rotation qui les échange préserve la forme d’aire.


\subsection{Entrelacement de phase et d'anisotropie}

Soient
\[
 S_\eta=\operatorname{diag}(e^{\eta/2},e^{-\eta/2}),\quad
 J_0=\begin{pmatrix}0&-1\\1&0\end{pmatrix},\quad
 J_\eta=S_\eta J_0S_\eta^{-1}
       =\begin{pmatrix}0&-e^\eta\\e^{-\eta}&0\end{pmatrix}.
\]
On a \(\det S_\eta=1\), \(J_\eta^2=-I\) et \(C_\eta(\theta)=e^{\theta J_\eta}=S_\eta e^{\theta J_0}S_\eta^{-1}\). La loi \(S_{\eta+\delta}=S_\delta S_\eta\) prouve l'entrelacement exact
\begin{equation}
 C_{\eta+\delta}(\theta)S_\delta=S_\delta C_\eta(\theta).
 \label{rec:eq:entrelazamiento}
\end{equation}
La transformation préserve \(dQ\wedge dP\) et transporte l’évolution de phase vers la nouvelle ellipse. Avec la normalisation de \eqref{eq:ii-elipse-area-fase}, l’action balayée est \(\hbar\theta\), ou \(\sigma\hbar\theta\) dans le feuillet orienté. La phase \(\theta\), l’anisotropie \(\eta_{\rm el}\) et le paramètre \(t\) du flux de \(K\) conservent ainsi des fonctions différentes, reliées par des opérateurs explicites.

\section{Domaine positif de la coordinatisation elliptique}
\begin{proposition}[Appartenance des représentations au système de coordonnées elliptique]
\label{iv:ellipse:pertenencia}
La branche positive de \eqref{eq:ii-par-angular}, avec les
intervalles \eqref{md:accion:intervalos}, vérifie
\(A>0\), \(0<C^*<A\) et \(\hbar_s>0\).
Le renversement de la coordonnée impaire conserve \(|C^*/A|<1\).
\end{proposition}
\begin{proof}
La proposition \ref{md:accion:positividad} donne
\(0<\eta_{\mathrm{ret}}<H_5\) et la positivité des deux
sections d’action. Pour une borne supérieure de \(H_5\),
\eqref{md:accion:intervalos} implique
\(\alpha<1/125\) et \(\varphi>8/5\). En conservant seulement
les termes positifs de \eqref{eq:el-h5}, on obtient
\[
 H_5<
 \frac{\sqrt5}{2}+
 \frac{9}{16\cdot125^2}+
 \frac{7}{48\cdot125^4}.
\]
Comme \(5<81/16\), on a \(\sqrt5/2<9/8\). De plus,

\[
 \frac{9}{16\cdot125^2}+
 \frac{7}{48\cdot125^4}
 =\frac{421882}{11718750000}<\frac1{1000}.
\]
Par conséquent, \(H_5<9/8+1/1000<2\). En utilisant maintenant
\(\alpha>7/1000\),
\[
 0<C^*=2(\eta_{\mathrm{ret}}+\alpha)
 <2\left(2+\frac1{125}\right)=\frac{502}{125}
 <7<1000\alpha=A.
\]
L’échange de canaux de
\ref{prop:ii-inversion-canales-angulares} remplace \(C^*\) par \(-C^*\)
et conserve \(A\) ; il conserve donc l’inégalité absolue.

\end{proof}

\section{Retours primitifs et fréquence relevée}
\subsection{Automate local d’extension}
Soit \(\mathscr X_\Gamma\) l’ensemble fini des états locaux compatibles avec la route. Un état conserve

\[
 x=(t,g,B,p,\varepsilon,c,w),
\]
où \(t\) est le temps discret, \(g\) la phase nonadique, \(B\) le bloc, \(p\) le préfixe, \(\varepsilon\) le feuillet, \(c\) la retenue et \(w\) l’enroulement. Il existe une arête \(x\to y\) lorsque \(y\) est une extension admissible de \(x\), respecte le registre généalogique et passe l’élagage \(G_9\). Le tableau de 36.864 compositions de retenue fournit la loi locale d’addition ; les tableaux de retour fixent la mémoire.

Soit \(U_\Gamma\) l’opérateur de transfert normalisé de ce graphe et \(R_\Gamma\) l’involution d’orientation. On exige

\[
 U_{C_M\Gamma}=J U_\Gamma J^{-1},\qquad
 R_{C_M\Gamma}=-J R_\Gamma J^{-1}.
\]
\Needspace{7\baselineskip}
\subsection{Dénombrements primitifs orientés}
Pour \(n\ge1\), la trace orientée

\[
 a_n(\Gamma)=\operatorname{Tr}(R_\Gamma U_\Gamma^n)
\]
dénombre les retours fermés avec signe. L’inversion de Möbius sépare les orbites primitives :

\[
 p_n(\Gamma)
 =\frac1n\sum_{d\mid n}\mu(d)a_{n/d}(\Gamma).
\]
On prend l’électron comme origine affine des corrections relatives et on définit

\[
 \boxed{
 \eta_h(\Gamma)=p_h(\Gamma)-p_h(\Gamma_e),
 \qquad h\in\mathfrak S.}
\]
On obtient ainsi l’application

\[
 \mathcal F_{\rm tower}:
 \Gamma^{\rm enr}\longmapsto
 \eta(\Gamma)=(\eta_h(\Gamma))_{h\in\mathfrak S}.
\]
Il n’y a pas de paramètres par espèce. Le même automate, la même involution et la même inversion de Möbius agissent sur les 324 routes.

\Needspace{7\baselineskip}
\subsection{Convergence du raffinement primitif}
Soit \(N_\Gamma=\dim\mathscr X_\Gamma\). Si \(U_\Gamma\) est une contraction,

\[
 |a_n(\Gamma)|\le N_\Gamma.
\]
Le nombre de diviseurs de \(n\) croît de manière sous-linéaire, de sorte que

\[
 |p_n(\Gamma)|
 \le\frac{N_\Gamma\tau(n)}n.
\]
Comme \(0<q_+<q_-<1\),

\[
 |s_h|\le q_-^h+q_+^h\le2q_-^h.
\]
Par comparaison,

\[
 \sum_{h\in\mathfrak S}|\eta_h(\Gamma)s_h|<\infty
\]
uniformément sur toute collection finie de fibres. Le caractère raffiné est donc bien défini.

\Needspace{7\baselineskip}
\subsection{Identités de contrôle}
La construction doit satisfaire simultanément :

\begin{enumerate}
\item invariance par renumérotation des états ;
\item covariance sous conjugaison de feuillet ;
\item compatibilité avec la somme directe de composantes ;
\item égalité lors du recalcul à partir du registre généalogique et du graphe certifié ;
\item conservation de l’origine \(\eta(\Gamma_e)=0\) ;
\item insensibilité absolue à une permutation des colonnes externes de masse ;
\item détermination des queues au moyen d’une borne antérieure à la confrontation.
\end{enumerate}
La représentation exacte d’un résiduel donné démontre la complétude du codomaine, mais ne remplace pas ces sept preuves.

\Needspace{7\baselineskip}
\section{Horloge de route et normalisation absolue}
\Needspace{7\baselineskip}
\subsection{Les retours entiers ne suffisent pas à distinguer les espèces}
Dans l’atlas de 324 routes, le premier retour de cellule est 27, le retour cinématique au sein de CT--108 est 54 et le retour cinématique étendu est 216. Étant communs, ces entiers fixent l’architecture temporelle globale, mais ne peuvent pas expliquer à eux seuls les différences de masse entre espèces.

\Needspace{7\baselineskip}
\subsection{Holonomie temporelle}
L’information spécifique réside dans l’holonomie de l’opérateur d’extension. Soit \(\lambda_j(\Gamma)\) une valeur propre non triviale de \(U_\Gamma\) sur le sous-espace persistant. Nous écrivons

\[
 \lambda_j(\Gamma)=r_j(\Gamma)e^{i\theta_j(\Gamma)}.
\]
La fréquence interne associée est

\[
 \omega_j(\Gamma)
 =\frac{\theta_j(\Gamma)+2\pi w_j(\Gamma)}{108t_0},
\]
où l’entier \(w_j\) est fixé par l’enroulement enregistré, non par un choix de branche ultérieur. Pour un mode stable \(r_j=1\) ; pour une résonance \(r_j<1\) et son atténuation apporte une largeur.

Le quotient

\[
 \rho_{\omega,j}(\Gamma)
 =\frac{\omega_j(\Gamma)}{\omega_0}
\]
est adimensionnel et peut modifier les rapports de masse sans altérer la décade commune \(-34\). L’énergie et la masse du mode sont

\[
 E_j(\Gamma)=\hbar_{\rm ret}\omega_j(\Gamma)
 \chi_{\rm full}(\Gamma)R_{\rm act}^{\kappa_j(\Gamma)},
\]
\[
 m_j(\Gamma)=E_j(\Gamma)c_{\rm int}^{-2}.
\]
Cette équation réunit action, horloge, caractère et lecteur sans transformer aucun d’eux en un autre. Son application numérique exige de déterminer \(U_\Gamma\), de fixer l’enroulement et de déclarer la coordinatisation dimensionnelle.


% Clausura focal de antecedentes: fuentes del artículo XI, edición 21-09-2026.
% Rangos, copias íntegras y huellas: gestion/mapa_restitucion_memoria.json.
\section{Antécédents de restitution opératoire et de mémoire}
\label{md:antmem:capitulo}

Le point de départ de ces antécédents se situe après la production conjointe
APP--TRIT--TPK de l’état enrichi et de la structure discrète du
continu. On travaille sur le registre dodécaphasique déjà engendré et sur des
réalisations déclarées de ses lecteurs. Le support linéaire ultérieur
et les candidats d’une reconstruction ne remplacent pas le domaine des
histoires engendrées. Une réalisation réversible ultérieure n’est pas non plus
identifiée, par sa seule notation, à l’holonomie nonadique \(\Gamma_9\).

Dans la coordinatisation et l’orientation du registre,
\begin{equation}
 \begin{gathered}
 K=(234,543,140,729,659,824,621,58,914,794,146,601),\\
 \mathbf1^{*}K=6263,\qquad
 (Sx)_i=x_{i+1},\quad i\pmod {12}.
 \end{gathered}
 \label{md:antmem:registro}
\end{equation}
Le registre est utilisé après sa génération, non comme entrée
rétrospective pour sélectionner une sortie. L’opérateur \(S\) est ici
le décalage positionnel marqué. Ses domaines seront maintenus
séparés des plans canoniques utilisés plus loin.

\subsection{L’orbite complète comme référence de son porteur}
\label{md:antmem:k:orbita}

Dans \(V=\mathbb R^{12}\), muni de son produit euclidien canonique, considérons
\[
 O_K=(K,SK,\ldots,S^{11}K):\mathbb R^{12}\longrightarrow V.
\]
Le décalage \(S\) agit sur les positions du registre. Il n’est pas
identifié par la seule périodicité à l’évolution complète TPK ou à
une symétrie exceptionnelle ultérieure.

\begin{theorem}[Référence cyclique complète et matrice de Gram exacte]
\label{md:antmem:k:marco-ciclico}
La matrice \(O_K\) est inversible. Sa matrice de Gram a pour valeurs propres
\[
\begin{array}{c|c}
 \text{valeur}&\text{multiplicité}\\\hline
 39225169=6263^2&1\\
 1248025-345420\sqrt3&2\\
 589609&2\\
 1407529&2\\
 1053157&2\\
 1248025+345420\sqrt3&2\\
 697225&1
\end{array}
\]
et satisfait
\begin{equation}
 589609\|c\|^2\leq\|O_Kc\|^2\leq39225169\|c\|^2,
 \quad
 \|O_K^{-1}\|=\frac1{\sqrt{589609}},\quad
 \operatorname{cond}(O_K)=\frac{6263}{\sqrt{589609}}.
 \label{md:antmem:k:cotas-orbita}
\end{equation}
L’orbite du registre centré a rang onze.
\end{theorem}
\begin{proof}
Les entrées de la matrice de Gram sont les autocorrélations \(\langle K,S^jK\rangle\). Dans l'ordre \(j=0,\ldots,11\), elles valent
\begin{align*}
 (&4251257,2999323,3163385,3287920,3216553,3344743,\\
  &2950064,3344743,3216553,3287920,3163385,2999323).
\end{align*}
Comme \(S\) est orthogonal, le gramien est circulant. Substituer les douze
caractères \(\exp(2\pi i m/12)\) dans cette ligne donne le tableau réel
précédent. Toutes les valeurs sont positives. La plus petite des deux valeurs
avec \(\sqrt3\) est supérieure à \(589609\), car
\(658416>345420\sqrt3\), inégalité qui se vérifie en élevant au
carré ses deux membres positifs. Les autres valeurs du tableau sont
supérieures ou égales à \(589609\). Le mode uniforme a pour valeur
\((\sum K_i)^2=6263^2\) ; l’inégalité triangulaire de la transformée
de Fourier, en utilisant \(K_i\geq0\), borne toutes les autres par cette même
valeur. Le théorème spectral donne \eqref{md:antmem:k:cotas-orbita}. Centrer
le registre élimine seulement le mode uniforme ; les onze autres conservent
les valeurs non nulles du tableau.
\end{proof}

Le déterminant, avec l’ordre des colonnes et l’orientation de \(S\)
fixés ici, est
\[
 \det O_K=5483119259214020183850397930008283625.
\]
L'absence de période propre du vecteur ne prouverait pas à elle seule l'inversibilité : le point décisif est qu'aucun de ses modes de Fourier ne s'annule. Le conditionnement numérique est approximativement \(8.15643\). Si l'on normalise \(K\) à une norme égale à un, les extrêmes de la matrice de Gram sont divisés par \(\|K\|^2=4251257\) ; l'amplitude de ses composantes ne doit pas être confondue avec une amélioration intrinsèque par rapport à toute référence de même norme.

\begin{corollary}[Séparation polaire de la forme et de l’orientation]
\label{md:antmem:k:polar}
Soient
\[
 G_K=O_K^*O_K,\qquad P_K=G_K^{1/2},\qquad
 F_K=O_KG_K^{-1/2}.
\]
Alors \(P_K\) est défini positif, \(F_K\) est orthogonal et
\(O_K=F_KP_K\). Une réalisation isométrique \(J:V\to\mathcal Y\)
conserve \(G_K\) et \(P_K\), et transporte la partie isométrique vers
\(JF_K\).
\end{corollary}
\begin{proof}
La positivité du Gram permet de définir sa racine et son inverse au moyen du
théorème spectral. L’identité
\(F_K^*F_K=G_K^{-1/2}G_KG_K^{-1/2}=I\) prouve l’assertion.
De plus, \((JO_K)^*(JO_K)=O_K^*O_K\).
\end{proof}

Cette polaire est définie à partir du gramien complet : elle ne s’obtient pas seulement à partir de la
liste de ses valeurs propres ni de la norme de \(K\). Le symbole
\(P_K\) de ce corollaire n’est pas le projecteur dimensionnel de rang dix
d’une autre réalisation incidentielle. Cette séparation est une conséquence de la
référence cyclique, non une identification de tous ses lecteurs.

\subsection{Identification d'opérateurs et holonomie relative}
\label{md:antmem:k:identificacion}

\begin{theorem}[Douze réponses générales ou une réponse cyclique]
\label{md:antmem:k:operadores}
Pour tout opérateur linéaire \(H:V\to V\), ses douze réponses \(Y=(HK,HSK,\ldots,HS^{11}K)\) déterminent
\[
 H=YO_K^{-1}.
\]
Si \(HS=SH\), une unique réponse vectorielle complète \(y=HK\) détermine
\begin{equation}
 H=O_yO_K^{-1},\qquad
 H=\sum_{j=0}^{11}h_jS^j,\qquad h=O_K^{-1}y.
 \label{md:antmem:k:filtro}
\end{equation}
Avec des erreurs \(E\) dans la matrice des réponses ou \(\delta y\) dans le
vecteur, respectivement,
\begin{align}
 \|\widehat H-H\|&\leq\frac{\|E\|}{\sqrt{589609}},
 \label{md:antmem:k:error-matriz}\\
 \|\widehat h-h\|&\leq\frac{\|\delta y\|}{\sqrt{589609}},\qquad
 \|\widehat H-H\|\leq\sqrt{\frac{12}{589609}}\,\|\delta y\|.
 \label{md:antmem:k:error-filtro}
\end{align}
\end{theorem}
\begin{proof}
La première égalité est \(Y=HO_K\). Si \(H\) commute avec \(S\), alors \(HS^jK=S^jy\), de sorte que \(Y=O_y\). La commutation avec le cycle détermine toutes les colonnes de \(H\) à partir de l'une d'entre elles ; les douze puissances \(I,S,\ldots,S^{11}\) sont indépendantes et engendrent cet espace d'opérateurs. Par conséquent, \(y=O_Kh\), ce qui donne la seconde formule. Les deux premières bornes utilisent \(\|O_K^{-1}\|=1/\sqrt{589609}\). La dernière utilise \(\|O_{\delta y}\|\leq\|O_{\delta y}\|_{\rm F}
=\sqrt{12}\|\delta y\|\).
\end{proof}

Une réponse signifie ici un vecteur à douze composantes, non une
mesure scalaire. Par exemple, si \(H=2I+3S-S^5\), la réponse
\(y=HK\) restitue exactement les coefficients \(2,3,-1\) à leurs
positions au moyen de \eqref{md:antmem:k:filtro} ; l’inverseur a besoin de \(y,K,S\),
non des coefficients du filtre. En revanche, l’opérateur non nul dont la première
ligne est \((543,-234,0,\ldots,0)\) et dont les autres lignes sont nulles
annule \(K\). Une seule réponse n’identifie pas un opérateur arbitraire.

La fonction de référence admet une généralisation exacte. Dans un cycle de
\(n\) positions, \(H\mapsto Hg\) identifie tous les opérateurs
qui commutent avec le cycle si et seulement si
\(O_g=(g,Sg,\ldots,S^{n-1}g)\) est inversible. La suffisance est la
preuve précédente. Si \(O_g\) est singulier, un vecteur non nul de son noyau
donne un polynôme de degré inférieur à \(n\) qui annule \(g\) ; l’opérateur
est non nul car les puissances du cycle sont indépendantes. L’impulsion
unitaire a également une orbite orthonormale et sert de référence. La
conséquence spécifique pour \(K\) est que le même objet déjà utilisé dans
la clôture et l’incidence vérifie cette condition, sans avoir été choisi de
nouveau pour optimiser l’essai. On n’en déduit pas une identification de
dynamiques non linéaires ou d’états cachés arbitraires ; une application à
une linéarisation exige que celle-ci ait été définie et justifiée.

\begin{corollary}[Récupération de l'holonomie relative à partir du complément]
\label{md:antmem:k:holonomia}
Soient \(H_B\) connu et inversible et \(H_C\) un autre transport du même type. Le complément
\[
 D_\gamma=\frac{\sqrt8}{9}(H_C-H_B)
\]
détermine
\[
 H_B^{-1}H_C=I+\frac9{\sqrt8}H_B^{-1}D_\gamma.
\]
Si \(D_\gamma\) commute avec \(S\), en particulier si les deux
transports le font, \(y=D_\gamma K\) suffit à récupérer
\begin{equation}
 H_B^{-1}H_C
 =I+\frac9{\sqrt8}H_B^{-1}O_yO_K^{-1}.
 \label{md:antmem:k:holonomia-una-respuesta}
\end{equation}
L'erreur dans cette reconstruction est bornée par
\[
 \frac9{\sqrt8}\|H_B^{-1}\|
 \sqrt{\frac{12}{589609}}\,\|\delta y\|.
\]
\end{corollary}
\begin{proof}
On isole \(H_C\) dans la définition du complément et l'on applique \cref{md:antmem:k:operadores}. La borne résulte de la sous-multiplicativité.
\end{proof}

Pour \(H_B=S^3\) et \(H_C=S^4\), la donnée rationnelle
\(y/\sqrt8=(S^4-S^3)K/9\) permet de reconstruire l’holonomie relative
\(S\) sans fournir \(H_C\) à l’inverseur. Si \(H_B\) est
unitaire, le facteur d’erreur est d’environ \(0.01435510\). Sans
commutation, on maintient les douze réponses générales, non une hypothèse
fictive de symétrie. La sélection d’une boucle physique et sa représentation
appartiennent toujours à leur dynamique génératrice ; ce protocole ne les
remplace pas.

\subsubsection{Conséquence de la référence marquée pour deux orientations}
\label{md:antmem:orientaciones}
La conséquence utilisée dans les notes du cahier conserve le tableau
précédent comme antécédent. Pour un opérateur réel \(H\) qui commute avec
\(S\), soient \(y_+=HK\) et \(y_-=H^{\mathsf T}K\).
Leurs douze intensités de Fourier coïncident, mais \(y_+=y_-\) si et seulement
si \(H=H^{\mathsf T}\). Chaque réponse vectorielle complète, avec
\(K\) et l’orientation de \(S\), récupère son opérateur.
\begin{proof}
L’opérateur \(H\) est circulant ; ses multiplicateurs de Fourier sont
\(\lambda_m\), et ceux de \(H^{\mathsf T}\) sont leurs conjugués.
Ainsi les deux réponses ont pour intensités
\(|\lambda_m|^2|\widehat K_m|^2\).
Si leurs vecteurs coïncident, \((H-H^{\mathsf T})K=0\).
La commutation implique que \(H-H^{\mathsf T}\) annule également \(S^jK\)
pour tout \(j\). Ces vecteurs forment une base, donc
\(H-H^{\mathsf T}=0\). La récupération est
\(H=O_{y_+}O_K^{-1}\), et de même pour \(H^{\mathsf T}\).
\end{proof}
En particulier, les décalages \(S\) et \(S^{-1}\) donnent des réponses
de même spectre de puissance, mais distinctes. Les autocorrélations
exactes du tableau précédent donnent l’identité déjà utilisée dans le cahier :
\begin{equation}
 \|SK-S^{-1}K\|^2
 =2\bigl(\|K\|^2-\langle K,S^2K\rangle\bigr)
 =2\,175\,744>0.
 \label{md:antmem:distancia-orientada}
\end{equation}
Ce nombre est une conséquence des autocorrélations, non une donnée
de départ supplémentaire. Le contrôle rationnel du cahier récupère à partir des deux
réponses les coefficients \(e_1\) et \(e_{11}\), respectivement, par
inversion de \(O_K\). Une intensité spectrale ou un bilan scalaire ne
remplace pas la réponse vectorielle.
Ce corollaire concerne le support cyclique à douze composantes :
il n’identifie pas \(S\) à la dynamique enrichie complète et ne transforme
pas à lui seul \(S\leftrightarrow S^{-1}\) en conjugaison physique de charge.
Pour \(H\) général, la transposition ne signifie pas non plus une inversion temporelle ;
dans le cas \(S\), elle coïncide bien avec \(S^{-1}\).

\subsection{Couplage réversible et transport non stationnaire}
\label{md:antmem:acoplamiento}
\subsubsection{Extension réversible du bilan entre deux formes}

Soient \(B,C:V\to V'\) des isomorphismes qui conservent une forme non dégénérée \(b\) entre deux réalisations :

\[
B^*b'B=C^*b'C=b.
\]

La forme peut être une métrique positive ou une forme alternée. Dans les espaces réels, \(*\) désigne la transposition par rapport aux coordonnées choisies ; l’égalité exprime le transport de la forme entre les réalisations spécifiées. Dans les espaces de Hilbert, les opérateurs et leurs inverses sont bornés.


Avec \(s=\sqrt8\) et \(Q_V=\frac13\left(\begin{smallmatrix}sI&-I\\I&sI\end{smallmatrix}\right)\), nous définissons

\begin{equation}
\boxed{
\mathcal U(B,C)=Q_{V'}^*\operatorname{diag}(B,C)Q_V
=\frac19\begin{pmatrix}8B+C&\sqrt8(C-B)\\\sqrt8(C-B)&B+8C\end{pmatrix}.}
\label{md:antmem:acop:eq:1}
\end{equation}

Nous écrivons \(T=(8B+C)/9\), \(D=\sqrt8(C-B)/9\) et \(R=(B+8C)/9\). La mise à jour complète est

\begin{equation}
\binom{x'}{m'}=\begin{pmatrix}T&D\\D&R\end{pmatrix}\binom{x}{m}.
\label{md:antmem:acop:eq:2}
\end{equation}

\begin{theorem}
L’application \eqref{md:antmem:acop:eq:1} est réversible et conserve la forme \(b\oplus b\) :

\begin{equation}
\mathcal U(B,C)^*(b'\oplus b')\mathcal U(B,C)=b\oplus b,
\quad\mathcal U(B,C)^{-1}=\mathcal U(B^{-1},C^{-1}).
\label{md:antmem:acop:eq:3}
\end{equation}
\end{theorem}

\begin{proof}
On a \(Q^*Q=I\) par \(8+1=9\). Ses blocs scalaires font que \(Q\) préserve \(b\oplus b\) pour toute forme \(b\). L’application diagonale conserve les formes par hypothèse. La composition prouve \eqref{md:antmem:acop:eq:3}, et l’inversion des trois facteurs donne l’inverse déclaré.
\end{proof}

Pour une mémoire entrante \(m=0\), on retrouve le bilan \(b=T^*b'T+D^*b'D\). La seconde colonne détermine la réponse à une mémoire entrante arbitraire. La conservation des normes positives et la conservation de l’aire orientée sont les spécialisations respectives de \eqref{md:antmem:acop:eq:3}.

\paragraph{Réalisation collective dans les neuf copies
}

Dans \(V^9\), soit \(W_E=\operatorname{diag}(I,\ldots,I,E)\), avec \(E\) en neuvième position. Soient \(A(a,b)=(a/\sqrt8,\ldots,a/\sqrt8,b)\) et \(L=AQ\). L’inclusion uniforme est \(Jx=(x,\ldots,x)/3\). Alors \(L(x,0)=Jx\) et

\begin{equation}
W_E L=L\mathcal U(I,E).
\label{md:antmem:acop:eq:4}
\end{equation}

\begin{proof}
\(W_EA=A\operatorname{diag}(I,E)\) ; multiplier par \(Q\) et utiliser \(AQ=L\) donne \eqref{md:antmem:acop:eq:4}. L’orthogonalité de la moyenne des huit premières positions avec leurs sept différences montre en outre que le complément reste fixe sous \(W_E\).
\end{proof}

Soit \(P_{\rm cyc}(v_0,\ldots,v_8)=(v_8,v_0,\ldots,v_7)\). Le pas cyclique original est \(S_E=P_{\rm cyc}W_E\). Ainsi \(S_EL=(P_{\rm cyc}L)U(I,E)\), avec des coordinatisations d’entrée et de sortie distinctes. La réduction collective conserve le motif correspondant : \(S_E^9=I_9\otimes E\), tandis que \(W_E^9=\operatorname{diag}(I_8,E^9)\). Lorsqu’un lecteur utilise les neuf étiquettes individuelles, celles-ci sont conservées dans leur représentation originale.

\subsubsection{Composition exacte et réentrée de mémoire
}

Pour \(B_1,C_1:V_0\to V_1\) et \(B_2,C_2:V_1\to V_2\),

\begin{equation}
\boxed{\mathcal U(B_2,C_2)\mathcal U(B_1,C_1)
=\mathcal U(B_2B_1,C_2C_1).}
\label{md:antmem:acop:eq:5}
\end{equation}

\begin{proof}
Dans \eqref{md:antmem:acop:eq:1}, \(Q_{V_1}Q_{V_1}^*\) s’annulent, et les deux blocs diagonaux se multiplient dans l’ordre indiqué. La récurrence étend la formule à tout nombre fini de transports, en conservant leur ordre.

\end{proof}

Le bloc de lecture contient l’identité

\begin{equation}
\boxed{T_2T_1+D_2D_1=\frac{8B_2B_1+C_2C_1}{9}.}
\label{md:antmem:acop:eq:6}
\end{equation}

À partir de \((x,0)\), la première étape produit \((T_1x,D_1x)\). La seconde lecture est \(T_2T_1x+D_2D_1x\). Le second terme est exactement la contribution de la mémoire réintroduite. En stockant \(D_1x\) hors du couplage suivant et en poursuivant seulement avec \(T_1x\), on obtient le procédé à registres séparés, dont la lecture est \(T_2T_1x\). L’égalité \eqref{md:antmem:acop:eq:6} précise la différence entre les deux procédés.

La réalisation \(U\) conserve l’état du couplage et les produits ordonnés \(B_N\cdots B_1\), \(C_N\cdots C_1\). La conservation de tous les préfixes inclut en outre le registre des émissions de HMT, qui maintient chaque facteur avec son ordre.

\paragraph{Équation covariante d’incrément
}

Pour le registre séparé \(x_{n+1}=T_nx_n\), \(m_n=D_nx_n\), on a


\begin{equation}
x_{n+1}-B_nx_n=m_n/\sqrt8.
\label{md:antmem:acop:eq:7}
\end{equation}

Avec \(G_0=I\) et \(G_{n+1}=B_nG_n\), le télescopage donne


\begin{equation}
x_0=G_N^{-1}x_N-\frac1{\sqrt8}\sum_{n<N}G_{n+1}^{-1}m_n.
\label{md:antmem:acop:eq:8}
\end{equation}

La mémoire est l’incrément par rapport au transport entre formes \(B_n\), avec la normalisation déclarée. L’égalité \eqref{md:antmem:acop:eq:8} est une reconstruction finie algébrique ; son passage à une série infinie requiert la convergence. Le théorème suivant construit l’inverse stable au moyen de la limite d’opérateurs positifs.


\subsubsection{Mémoire d’une suite non stationnaire}

Dans une réalisation de Hilbert \(H\), soient \(C_n\) des unitaires, avec leur ordre de composition déclaré. Nous écrivons

\[
T_n=(8I+C_n)/9,\quad D_n=\sqrt8(C_n-I)/9,
\quad P_0=I,\quad P_{n+1}=T_nP_n,\quad A_N=P_N^*P_N.
\]

L’indice \(N\) compte les compressions avec stockage séparé, à la différence du transport complet avec réentrée. Pour des métriques variables avec \(B_n\) des isomorphismes isométriques, la conjugaison par \(G_n\) transporte les transports vers un espace de Hilbert fixe et produit ce même cas, avec \(\widetilde C_n=G_{n+1}^{-1}C_nG_n\). L’assertion utilise les métriques des fibres et les isométries déclarées entre elles.

\begin{theorem}[Terminal positif et récupération non stationnaire]
\label{md:antmem:terminal-no-estacionario}
Il existe un opérateur positif \(0\leq A_\infty\leq I\), limite forte de \(A_N\), et

\begin{equation}
\boxed{I=A_\infty+\sum_{n\ge0}P_n^*D_n^*D_nP_n.}
\label{md:antmem:acop:eq:9}
\end{equation}

L'application

\begin{equation}
\boxed{\mathcal Vv=(A_\infty^{1/2}v,(D_nP_nv)_{n\ge0})}
\label{md:antmem:acop:eq:10}
\end{equation}

est isométrique et possède un inverse à gauche stable


\begin{equation}
\mathcal V^*(a,(m_n))=A_\infty^{1/2}a+
\sum_{n\ge0}P_n^*D_n^*m_n,\qquad\|\mathcal V^*\|=1.
\label{md:antmem:acop:eq:11}
\end{equation}

\end{theorem}
\begin{proof}
L’identité locale \(T_n^*T_n+D_n^*D_n=I\) donne
\(A_N-A_{N+1}=P_N^*D_N^*D_NP_N\geq0\). Les formes quadratiques de
la suite positive décroissante convergent vers celle d’un opérateur positif
\(A_\infty\). Pour \(B_N=A_N-A_\infty\), l’inégalité
\(0\leq B_N^2\leq B_N\) transforme cette convergence en convergence forte.
Le télescopage fini et la limite prouvent \eqref{md:antmem:acop:eq:9}. L’évaluation sur \(v\)
prouve \eqref{md:antmem:acop:eq:10}. Son adjoint est \eqref{md:antmem:acop:eq:11} ; la série converge en norme parce que
les troncatures d’un registre de carré sommable convergent et que l’adjoint est borné.
\end{proof}

La reconstruction avec le terminal limite et \(N\) registres a pour erreur
exacte \((A_N-A_\infty)v\). Le vecteur \(A_\infty^{1/2}v\) est la
réalisation positive canonique du terminal de norme. Son existence provient
de la limite de Gram, indépendamment de la convergence vectorielle de \(P_Nv\).


\subsubsection{Trois propriétés différentes : épuisement, distinction et stabilité}

Soit \(Mv=(D_nP_nv)_n\). De \eqref{md:antmem:acop:eq:9}, \(M^*M=I-A_\infty\). Par conséquent :

\begin{enumerate}
\item Toute la norme de \(v\) passe au registre exactement lorsque \(A_\infty v=0\).

\item Le registre distingue chaque état exactement lorsque \(\ker(I-A_\infty)=\{0\}\).
\item La récupération à partir de la mémoire seule possède un inverse borné exactement lorsque
\(I-A_\infty\geq\varepsilon I\) pour un certain \(\varepsilon>0\).
Dans ce cas, un inverse à gauche est \((I-A_\infty)^{-1}M^*\).

\end{enumerate}

De plus,

\begin{equation}
\boxed{\ker M=\bigcap_{n\ge0}\operatorname{Fix}(C_n).}
\label{md:antmem:acop:eq:12}
\end{equation}

\begin{proof}
Si \(Mv=0\), alors \(D_0v=0\), d’où \(C_0v=v\)
et \(P_1v=v\). Par récurrence, \(D_nP_nv=0\) donne \(C_nv=v\) et
\(P_{n+1}v=v\). La réciproque est immédiate. Les trois équivalences se
déduisent du Gram \(M^*M\) et de la caractérisation des opérateurs
bornés inférieurement.
\end{proof}

\paragraph{Exemple exact qui sépare les trois propriétés}

Dans l’exemple algébrique \(C_0=-I\) et \(C_n=I\) pour \(n\geq1\),

\begin{equation}
A_\infty=\frac{49}{81}I,\qquad M^*M=\frac{32}{81}I,
\qquad m_0=-\frac{2\sqrt8}{9}v,
\qquad v=-\frac9{2\sqrt8}m_0.
\label{md:antmem:acop:eq:13}
\end{equation}

La mémoire contient \(32/81\) du carré de la norme et permet de récupérer le vecteur complet. Le terminal retient \(49/81\) et est un opérateur positif distinct d’un projecteur. La récupérabilité dépend du noyau et du conditionnement du lecteur. Lecture et mémoire constituent des composantes corrélées d’une isométrie.


Cet exemple sépare les trois conséquences logiques du théorème. L’application
à une trajectoire physique conserve en outre la sélection et la
provenance effective de ses \(C_n\) au moyen du TPK.

La lettre \(\eta\) de la réalisation gaussienne suivante désigne la
déformation de l’ellipse, c’est-à-dire le paramètre \(\zeta\) des
développements de transport ; ce n’est pas le coefficient d’action
\(\eta_{\rm ret}\). On maintient une section positive \(\hbar_a\) et
une même unité pour les deux coordonnées du quotient angulaire.
Les demi-axes conservés sont
\[
 a_S=\sqrt{2\hbar_a}\,e^{\eta/2},\qquad
 b_S=\sqrt{2\hbar_a}\,e^{-\eta/2},\qquad
 \eta=\operatorname{artanh}(C^*/A).
\]

\subsection{Covariance canonique et réciprocité de l'ellipse d'action}
\label{md:antmem:covarianza-accion}

La construction précédente détermine deux invariants complémentaires :
le produit des demi-axes fixe l’aire d’action et leur quotient conserve
l’anisotropie de la paire angulaire. Précisons maintenant la réalisation canonique
dans laquelle ces demi-axes sont deux écarts-types. La structure de
départ est l’ellipse déjà obtenue à partir des représentations angulaires et
d’action d’APP--TRIT--TPK ; la réalisation opératorielle s’établit après
cette construction.

Dans ce paragraphe, on écrit \(\eta=\eta_{\rm el}
=\operatorname{artanh}(C^*/A)\), avec \(|C^*/A|<1\) et
\(\hbar_a>0\). Le quotient angulaire utilise la même unité pour ses deux
composantes. Dans les coordonnées équilibrées \((Q,P)\), toutes deux de
dimension racine carrée d'action, définissons
\begin{equation}
 V_\eta:=\frac14
 \begin{pmatrix}a_S^2&0\\0&b_S^2\end{pmatrix}
 =\frac{\hbar_a}{2}
 \begin{pmatrix}e^\eta&0\\0&e^{-\eta}\end{pmatrix}.
 \label{md:antmem:covarianza-geometrica}
\end{equation}
Cette matrice positive possède des entrées de dimension action. Son interprétation
comme covariance quantique est spécifiée par la réalisation suivante.

\begin{proposition}[Réalisation gaussienne de la covariance d'action]
\label{md:antmem:covarianza-gaussiana}
Considérons la représentation canonique dans \(L^2(\mathbb R,dQ)\), avec
\(\widehat Qf(Q)=Qf(Q)\) et
\(\widehat Pf(Q)=-i\hbar_a f'(Q)\), initialement sur l'espace de
Schwartz \(\mathcal S(\mathbb R)\). L'état normalisé
\begin{equation}
 \psi_\eta(Q)=\bigl(\pi\hbar_a e^\eta\bigr)^{-1/4}
 \exp\!\left(-\frac{Q^2}{2\hbar_a e^\eta}\right)
 \label{md:antmem:gaussiana-accion}
\end{equation}
a des moyennes nulles et pour matrice de covariance \(V_\eta\). En particulier,
\begin{equation}
 (\Delta Q)^2=\frac{\hbar_a e^\eta}{2},\qquad
 (\Delta P)^2=\frac{\hbar_a e^{-\eta}}{2},\qquad
 \operatorname{Cov}(Q,P)=0,\qquad
 \det V_\eta=\frac{\hbar_a^2}{4}.
 \label{md:antmem:varianzas-accion}
\end{equation}
Cette réalisation sature l'inégalité de Robertson--Schrödinger.
\end{proposition}

\begin{proof}
Les opérateurs précédents sont essentiellement autoadjoints sur
\(\mathcal S(\mathbb R)\), qui constitue un cœur commun invariant ;
sur cet espace, on a
\([\widehat Q,\widehat P]=i\hbar_a I\). Pour
\(\widehat Y=(\widehat Q,\widehat P)\), la covariance symétrique est
\[
 V_{jk}=\frac12\left\langle
 (\widehat Y_j-\langle\widehat Y_j\rangle)
 (\widehat Y_k-\langle\widehat Y_k\rangle)
 +(\widehat Y_k-\langle\widehat Y_k\rangle)
 (\widehat Y_j-\langle\widehat Y_j\rangle)
 \right\rangle.
\]
Avec \(s=\hbar_a e^\eta>0\), soit
\(I_s=\int_{\mathbb R}e^{-Q^2/s}\,dQ\). Le produit de deux
intégrales identiques, évalué en coordonnées polaires, donne
\(I_s^2=2\pi\int_0^\infty r e^{-r^2/s}\,dr=\pi s\).
En outre, l'intégrale de
\(\frac{d}{dQ}(Qe^{-Q^2/s})\) est nulle, par décroissance aux deux
extrémités ; d'où
\(\int_{\mathbb R}Q^2e^{-Q^2/s}\,dQ=sI_s/2\).
Ces identités et la parité démontrent la normalisation, la moyenne nulle et
\(\langle\widehat Q^2\rangle=s/2\). De plus,
\(\widehat P\psi_\eta=i e^{-\eta}Q\psi_\eta\).
Ainsi, \(\langle\widehat P\rangle=0\),
\(\|\widehat P\psi_\eta\|^2=\hbar_a e^{-\eta}/2\) et
\(\operatorname{Re}\langle\widehat Q\psi_\eta,
\widehat P\psi_\eta\rangle=0\).

Pour tout état normalisé de ce domaine, Cauchy--Schwarz appliquée
à \((\widehat Q-\langle\widehat Q\rangle)\psi\) et
\((\widehat P-\langle\widehat P\rangle)\psi\) donne
\[
 (\Delta Q)^2(\Delta P)^2
 \geq \operatorname{Cov}(Q,P)^2
       +\frac14\left|\langle[\widehat Q,\widehat P]\rangle\right|^2
 =\operatorname{Cov}(Q,P)^2+\frac{\hbar_a^2}{4}.
\]
Les variances calculées atteignent l'égalité. De manière équivalente,
l'opérateur
\[
 a_\eta=\frac{e^{-\eta/2}\widehat Q
                  +i e^{\eta/2}\widehat P}{\sqrt{2\hbar_a}}
\]
satisfait \([a_\eta,a_\eta^\dagger]=I\) et
\(a_\eta\psi_\eta=0\), par substitution directe.
\end{proof}

\subsection{Covariance, forme normale et spectre gaussien}
\label{md:antmem:gauss:gaussianas-espectro}

La section positive d’action \(\hbar_a\) et la réalisation canonique
précédente déterminent la normalisation employée ici. On démontre la
correspondance entre covariance, spectre, pureté et entropie pour un
nombre fini de modes ; l’espace de Fock conserve toutes ses occupations.
La réduction symplectique et les formules gaussiennes sont des résultats
établis qui sont reproduits pour rendre autonome leur application au
transport de mémoire.

\subsubsection{Coordonnées canoniques et fonction caractéristique}
Dans \(L^2(\mathbb R^d)\), soient \(\widehat q_j=\widehat Q_j/\sqrt{\hbar_a}\),
\(\widehat p_j=\widehat P_j/\sqrt{\hbar_a}\) et
\(\widehat z=(\widehat q_1,\widehat p_1,\ldots,\widehat q_d,\widehat p_d)\).
Sur Schwartz,

\[
 [\widehat z_j,\widehat z_k]=iJ_{jk}I,\qquad
 J=\bigoplus_{j=1}^d\begin{pmatrix}0&1\\-1&0\end{pmatrix}.
\]
Cette convention fixe le signe de la forme canonique. Les opérateurs de
Weyl \(\mathcal W(\xi)=\exp(i\xi^{\mathsf T}\widehat z)\) sont définis
par les fermetures auto-adjointes des combinaisons linéaires indiquées.
Pour un état centré \(\rho\) de moments d’ordre deux finis, la
covariance normalisée et sa fonction caractéristique sont
\[
 W_{jk}=\operatorname{Tr}\rho(\widehat z_j\widehat z_k+\widehat z_k\widehat z_j),
 \qquad \chi_\rho(\xi)=\operatorname{Tr}(\rho\mathcal W(\xi)),\qquad
 V=\frac{\hbar_a}{2}W.
\]
L’état est gaussien centré lorsque

\begin{equation}
 \chi_\rho(\xi)=\exp(-\xi^{\mathsf T}W\xi/4).
 \label{md:antmem:gauss:eq:caracteristica-gaussiana}
\end{equation}
L’espérance de \((c^{\mathsf T}\widehat z)^\dagger
(c^{\mathsf T}\widehat z)\) est \(c^\dagger(W+iJ)c/2\) ; ainsi,
\(W+iJ\geq0\). De plus, \(W>0\) : si \(v\) est réel et
\(v^{\mathsf T}Wv=0\), la positivité implique \((W+iJ)v=0\),
donc \(Jv=0\) et \(v=0\).


\begin{lemma}[Unicité de la fonction caractéristique
]
\label{md:antmem:gauss:lem:unicidad-weyl}
Deux opérateurs de classe trace ayant la même fonction caractéristique sont égaux.
\end{lemma}
\begin{proof}
En regroupant les positions et les moments sous la forme \(\xi=(u,v)\),

\[
 (\mathcal W(u,v)f)(x)=e^{iu\cdot(x+v/2)}f(x+v).
\]
Pour un opérateur de rang fini à vecteurs de Schwartz et de noyau \(K_A\),

\[
 \operatorname{Tr}(A\mathcal W(u,v))=\int K_A(x+v/2,x-v/2)e^{iu\cdot x}\,dx.
\]
Plancherel et le changement de variables de jacobien absolu un donnent
\[
 \|A\|_{\rm HS}^2=(2\pi)^{-d}\int_{\mathbb R^{2d}}
 |\operatorname{Tr}(A\mathcal W(\xi))|^2\,d\xi.
\]
Les opérateurs précédents sont denses dans la classe trace. La convergence
dans cette norme implique la convergence en Hilbert--Schmidt et la convergence
uniforme des fonctions caractéristiques, car leur différence est
bornée par \(\|A-B\|_1\). L’égalité s’étend dans \(L^2\).
L’appliquer à la différence des opérateurs prouve l’assertion.
\end{proof}

\subsubsection{Réduction symplectique de la forme positive
}
\begin{theorem}[Forme normale de Williamson]
\label{md:antmem:gauss:thm:williamson}
Pour une matrice réelle symétrique \(W>0\) d’ordre \(2d\), il existe
\(S\in\operatorname{Sp}(2d,\mathbb R)\) et \(\nu_j>0\) tels que
\[
 W=SDS^{\mathsf T},\qquad D=\bigoplus_j\nu_jI_2.
\]
Les valeurs propres de \(JW\) sont \(\pm i\nu_j\). La condition
\(W+iJ\geq0\) équivaut à \(\nu_j\geq1\) pour tout \(j\).

\end{theorem}
\begin{proof}
La matrice \(B=W^{1/2}JW^{1/2}\) est antisymétrique et inversible.
Si \(-B^2u=\nu^2u\) et \(\|u\|=1\), le vecteur
\(v=-Bu/\nu\) est unitaire et orthogonal à \(u\) ; il satisfait
\(Bu=-\nu v\), \(Bv=\nu u\). Le complément de cette paire est
invariant. L’itération produit une matrice orthogonale \(O\) telle que
\(O^{\mathsf T}BO=JD\). En définissant \(S=W^{1/2}OD^{-1/2}\),

\[
 S^{\mathsf T}JS=J,\qquad SDS^{\mathsf T}=W.
\]
La similitude \(W^{1/2}(JW)W^{-1/2}=B\) détermine les valeurs
\(\nu_j\), avec multiplicité. La congruence par \(S^{-1}\)
transforme \(W+iJ\) en \(D+iJ\) ; chaque bloc a pour valeurs propres
\(\nu_j+1\) et \(\nu_j-1\).
\end{proof}

\begin{lemma}[Implémentation unitaire dans un nombre fini de modes]
\label{md:antmem:gauss:lem:implementacion-simplectica}
Toute matrice symplectique \(S\) admet un unitaire \(U_S\) qui satisfait
\(U_S^\dagger\mathcal W(\xi)U_S=\mathcal W(S^{\mathsf T}\xi)\).
La conjugaison \(\rho\mapsto U_S\rho U_S^\dagger\) transporte
la covariance par \(W\mapsto SWS^{\mathsf T}\).
\end{lemma}
\begin{proof}
Dans la décomposition polaire \(S=OP\), l’identité
\(J^{-1}(S^{\mathsf T}S)J=(S^{\mathsf T}S)^{-1}\), appliquée par
calcul spectral à la racine positive, prouve \(PJP=J\).
Alors \(O\) est orthogonale symplectique. Les espaces propres de \(P\)
s’apparient au moyen de \(J\) avec les valeurs propres \(\lambda,\lambda^{-1}\) ;
celui de valeur propre un est \(J\)-invariant. Une matrice
orthogonale symplectique \(K\) réduit donc \(P\) à
\(\bigoplus_j\operatorname{diag}(e^{r_j},e^{-r_j})\).

La dilatation est mise en œuvre par

\[
 (\mathcal D_rf)(x)=e^{-\frac12\sum_jr_j}
 f(e^{-r_1}x_1,\ldots,e^{-r_d}x_d).
\]
Un changement de variables prouve son unitarité et la conjugaison de
\((\widehat q_j,\widehat p_j)\) en
\((e^{r_j}\widehat q_j,e^{-r_j}\widehat p_j)\).
Une transformation orthogonale symplectique commute avec \(J\) et
correspond à une matrice \(u\in U(d)\) sur
\(a_j=(\widehat q_j+i\widehat p_j)/\sqrt2\). Sa seconde
quantification \(\Gamma(u)=\bigoplus_{n\geq0}u^{\otimes_s n}\)
est unitaire dans \(\mathcal F_s(\mathbb C^d)\), et l’identité sur les
produits symétriques finis donne la transformation des \(a_j\).

La base de Hermite identifie cet espace à \(L^2(\mathbb R^d)\).
Pour justifier sa complétude, si \(f\) est orthogonal à tous les
polynômes multipliés par \(e^{-|x|^2/2}\), la fonction entière
\(F(z)=\int f(x)e^{-|x|^2/2}e^{z\cdot x}\,dx\) a toutes
ses dérivées nulles en zéro. Ainsi, \(F=0\) ; en restreignant à
\(z=i\xi\), l’unicité de Fourier donne \(f=0\). Les
normalisations et l’orthogonalité résultent de la création et de l’annihilation.
La composition des mises en œuvre précédentes produit \(U_S\).
L’égalité de Weyl implique
\(\chi_{U_S\rho U_S^\dagger}(\xi)=\chi_\rho(S^{\mathsf T}\xi)\),
qui démontre la loi de covariance.

\end{proof}

\subsubsection{Calcul du spectre de la densité isotrope}
Dans un mode, les états cohérents normalisés sont
\[
 |z\rangle=e^{-|z|^2/2}\sum_{n\geq0}\frac{z^n}{\sqrt{n!}}|n\rangle.
\]
Leur fonction de position,
\(\pi^{-1/4}\exp[-(x-q_0)^2/2+ip_0(x-q_0/2)]\), où
\(q_0=\sqrt2\Re z\), \(p_0=\sqrt2\Im z\), donne par intégration
\[
 \langle z|\mathcal W(u,v)|z\rangle
 =e^{-(u^2+v^2)/4}e^{i(uq_0+vp_0)}.
\]
Pour \(s>0\), l’intégrale
\[
 \rho^{(s)}=\frac1{\pi s}\int_{\mathbb C} e^{-|z|^2/s}
 |z\rangle\langle z|\,d^2z
\]
converge en classe trace, est positive et a pour trace un. L’intégration
angulaire annule ses éléments hors de la diagonale d’occupation ; l’intégration
radiale donne
\[
 \langle n|\rho^{(s)}|n\rangle=\frac{s^n}{(1+s)^{n+1}},\qquad
 \chi_{\rho^{(s)}}(u,v)=e^{-(2s+1)(u^2+v^2)/4}.
\]
Avec \(\nu=2s+1\), l’unique état gaussien de covariance
\(\nu I_2\) est, par le lemme~\ref{md:antmem:gauss:lem:unicidad-weyl},
\begin{equation}
 \rho_\nu=(1-t)\sum_{n\geq0}t^n|n\rangle\langle n|,
 \qquad t=\frac{\nu-1}{\nu+1}.
 \label{md:antmem:gauss:eq:estado-geometrico}
\end{equation}
Le cas \(\nu=1\) est le projecteur sur le vide.

\begin{theorem}[Spectre, pureté et entropie
]
\label{md:antmem:gauss:thm:espectro-gaussiano}
Un état gaussien centré de covariance \(W\) et de valeurs
symplectiques \(\nu_1,\ldots,\nu_d\) est unitairement équivalent
à \(\bigotimes_j\rho_{\nu_j}\). Ses valeurs propres sont
\[
 \lambda_{\mathbf n}=\prod_j(1-t_j)t_j^{n_j},\qquad
 t_j=\frac{\nu_j-1}{\nu_j+1},\quad \mathbf n\in\mathbb N^d.
\]
Avec \(0\log0=0\),
\begin{align}
 \operatorname{Tr}\rho^2&=\prod_j\nu_j^{-1}=(\det W)^{-1/2},
 \label{md:antmem:gauss:eq:pureza-general}\\
 -\operatorname{Tr}(\rho\log\rho)&=\sum_j
 \left[\frac{\nu_j+1}{2}\log\frac{\nu_j+1}{2}
 -\frac{\nu_j-1}{2}\log\frac{\nu_j-1}{2}\right].
 \label{md:antmem:gauss:eq:entropia-general}
\end{align}
\end{theorem}
\begin{proof}
Williamson donne \(W=SDS^{\mathsf T}\). Le produit des densités
isotropes a pour fonction caractéristique \(e^{-\xi^{\mathsf T}D\xi/4}\).
L’implémentation unitaire et l’unicité de Weyl prouvent
l’équivalence et la formule spectrale. Dans un mode,
\[
 \sum_{n\geq0}(1-t)^2t^{2n}=\frac{1-t}{1+t}=\nu^{-1},
\]
et
\[
 -\sum_{n\geq0}(1-t)t^n\log((1-t)t^n)
 =-\log(1-t)-\frac{t}{1-t}\log t.
\]
La substitution \(t=(\nu-1)/(\nu+1)\) prouve les formules ; la
limite en \(t=0\) est nulle. Dans le produit fini, Tonelli justifie
la somme des termes d’entropie positifs ou nuls, et
\(\det W=\prod_j\nu_j^2\) donne l’expression déterminantale.

\end{proof}

Si l’état est la réduction d’un vecteur pur bipartite, ses
valeurs propres sont les carrés des coefficients de Schmidt.
La décomposition en valeurs singulières de l’opérateur de Hilbert--Schmidt défini
par ce vecteur démontre l’égalité des spectres non nuls des deux
réductions. L’entropie calculée est alors celle d’intrication.

\subsubsection{Préparation et transport des deux covariances}
L’application ultérieure utilise deux entrées gaussiennes pures identiques.
En général, pour \(M\in\operatorname{Sp}(2d,\mathbb R)\), leur covariance est
\(V_0=(\hbar_a/2)MM^{\mathsf T}\), et l’entrée conjointe possède la
covariance \(V_0\oplus V_0\). Le transport physique s’obtient
en conjuguant \(U_H\) par \(M\oplus M\). Dans la coordinatisation équilibrée,
la covariance normalisée d’entrée est \(I_{4d}\) ; avec
\[
 T=\frac{8I+H}{9},\qquad D=\frac{\sqrt8(H-I)}9,
\]
la première réduction a
\[
 W=TT^{\mathsf T}+DD^{\mathsf T}=\frac{8I+HH^{\mathsf T}}9.
\]
Sa fonction caractéristique s’obtient en annulant les arguments du
second groupe de modes. Elle reste donc gaussienne et les théorèmes
précédents calculent tout son spectre, sans tronquer le registre d’occupation.

\subsection{Prolongement du transport de mémoire : réalisation quantique}

La réalisation quantique utilise le couplage réversible de la section~\ref{md:antmem:acoplamiento}. Nous considérons d’abord un plan canonique réel de matrice alternée \(\Omega=\left(\begin{smallmatrix}0&1\\-1&0\end{smallmatrix}\right)\). La partition nonadique détermine

\begin{equation}
Q=\frac13\begin{pmatrix}\sqrt8 I&-I\\I&\sqrt8 I\end{pmatrix},\qquad
U_H=Q^{\mathsf T}\operatorname{diag}(I,H)Q
=\frac19\begin{pmatrix}8I+H&\sqrt8(H-I)\\
\sqrt8(H-I)&I+8H\end{pmatrix}.
\label{md:antmem:cuantica:eq:5.1}
\end{equation}

L’opérateur \(H\in\operatorname{Sp}(2,\mathbb R)\) transporte ce plan. La matrice \(Q\) conserve le produit euclidien et la forme \(\Omega\oplus\Omega\) ; c’est pourquoi \(U_H\) est symplectique. La première composante définit le système observé et la seconde le registre de mémoire ; la transformation complète agit sur les deux.

Le lacet ordonné dont les identités sont reproduites dans
\ref{md:antmem:lazo-ep} est

\begin{equation}
C=\begin{pmatrix}0&-1\\1&-1\end{pmatrix},\quad
P=\begin{pmatrix}1&1\\0&1\end{pmatrix},\quad
H_n=C^{-1}P^{-n}CP^n
=\begin{pmatrix}1+n&n^2\\n&n^2-n+1\end{pmatrix}.
\label{md:antmem:cuantica:eq:5.2}
\end{equation}

Pour chaque \(n\in\mathbb Z\), on a \(\det H_n=1\). En particulier, pour \(n=1\),

\begin{equation}
H_1=\begin{pmatrix}2&1\\1&1\end{pmatrix}
=F_1^2,\quad F_1=\begin{pmatrix}1&1\\1&0\end{pmatrix}
=P F_{\rm av}P^{-1}.
\label{md:antmem:cuantica:eq:5.3}
\end{equation}

Son spectre est \(\{\varphi^2,\varphi^{-2}\}\), par la reconnaissance ultérieure du mode d’or déjà engendré. L’entier \(n\) spécifie le nombre orienté d’itérations dans le lacet choisi. La collection des transports est déterminée algébriquement ; sa réalisation comme protocole physique conserve en outre les données de préparation et d’évolution correspondantes.

\subsubsection{Identités algébriques du lacet et reconnaissance ultérieure}
\label{md:antmem:lazo-ep}
Pour conserver la preuve des identités utilisées, on reproduit
leur composition matricielle. Dans la même coordinatisation,
\[
 N_0=\begin{pmatrix}0&1\\0&0\end{pmatrix}.
\]
Soit
\[
 P=I+N_0=\begin{pmatrix}1&1\\0&1\end{pmatrix}.
\]
Comme \(N_0^2=0\), on a \(P^n=I+nN_0\) pour tout
\(n\in\mathbb Z\).
Le lacet ordonné de quatre transports donne

\begin{equation}
H_n=C^{-1}P^{-n}CP^n
=\begin{pmatrix}1+n&n^2\\n&n^2-n+1\end{pmatrix}.
\label{md:antmem:matriz:eq:28}
\end{equation}

Pour \(n\ne0\), sa différence normalisée est

\[
F_n=\frac{H_n-I}{n}
=\begin{pmatrix}1&n\\1&n-1\end{pmatrix}.
\]

En multipliant ces matrices, on obtient, sans introduire de valeurs spectrales,

\begin{equation}
\boxed{F_n^2=H_n,\qquad F_n^2-nF_n-I=0.}
\label{md:antmem:matriz:eq:29}
\end{equation}

Avec \(H_B=I\), la mémoire est
\(D_{\gamma,n}=(\sqrt8/9)nF_n\).
Par conséquent,

\begin{equation}
\boxed{H_n=\left(\frac9{n\sqrt8}D_{\gamma,n}\right)^2.}
\label{md:antmem:matriz:eq:30}
\end{equation}

La réponse matricielle de mémoire reconstruit l’auto-échelle
hyperbolique. Pour \(n>0\), la reconnaissance spectrale ultérieure
identifie

\begin{equation}
\rho_{{\rm ep},n}=\frac{n+\sqrt{n^2+4}}2,\quad
\operatorname{Spec}F_n=\{\rho_{{\rm ep},n},-\rho_{{\rm ep},n}^{-1}\},\quad
\operatorname{Spec}H_n=\{\rho_{{\rm ep},n}^2,\rho_{{\rm ep},n}^{-2}\}.
\label{md:antmem:matriz:eq:31}
\end{equation}

Le cas \(n=1\) produit le polynôme d’or.
La relation avec l’opérateur
\(F_{\rm av}=
\begin{pmatrix}0&1\\1&1\end{pmatrix}\)
engendré par le calendrier conserve le changement
de coordinatisation :

\begin{equation}
F_1=PF_{\rm av}P^{-1},\qquad
H_1=PF_{\rm av}^{\,2}P^{-1}.
\label{md:antmem:matriz:eq:32}
\end{equation}

Le transport \(P\), de déterminant un, entrelace les deux
matrices et conserve \(\Omega\).

\subsubsection{État initial, transport et lecture}

La proposition~\ref{md:antmem:covarianza-gaussiana} construit la réalisation gaussienne de l’ellipse d’action. Dans la section positive \(\hbar_a>0\), avec \(A>0\) et \(|C^*/A|<1\), sa covariance est


\begin{equation}
V_\eta=\frac{\hbar_a}{2}M_\eta M_\eta^{\mathsf T},\qquad
M_\eta=\operatorname{diag}(e^{\eta/2},e^{-\eta/2}),\qquad
\eta=\operatorname{artanh}(C^*/A).
\label{md:antmem:cuantica:eq:5.4}
\end{equation}

Dans la représentation de Schrödinger sur \(L^2(\mathbb R,dQ)\), avec \(\widehat Q\) la multiplication par \(Q\) et \(\widehat P=-i\hbar_a\,d/dQ\), la relation \([\widehat Q,\widehat P]=i\hbar_a I\) vaut dans l’espace de Schwartz. L’état pur normalisé est

\[
\psi_\eta(Q)=(\pi\hbar_a e^\eta)^{-1/4}
\exp[-Q^2/(2\hbar_a e^\eta)].
\]

La préparation est le produit \(\psi_\eta\otimes\psi_\eta\). Le transport de \eqref{md:antmem:cuantica:eq:5.1} s’exprime dans la même coordinatisation de covariance par

\begin{equation}
\widetilde U=(M_\eta\oplus M_\eta)U_H
(M_\eta^{-1}\oplus M_\eta^{-1}).
\label{md:antmem:cuantica:eq:5.5}
\end{equation}

Le lemme~\ref{md:antmem:gauss:lem:implementacion-simplectica} fournit une implémentation unitaire de cette matrice symplectique à deux modes. Les implémentations qui diffèrent d’une phase globale produisent la même conjugaison de l’état. La covariance se transforme par \(V\mapsto\widetilde U V\widetilde U^{\mathsf T}\) ; l’état conjoint reste pur et sa trace partielle sur le second mode détermine la lecture visible, conformément à la construction de la section~\ref{md:antmem:gauss:gaussianas-espectro}.

\begin{theorem}[Mélange réduit par déformation relative]

La covariance de la lecture visible est

\begin{equation}
\boxed{V_{\rm vis}=
\frac{\hbar_a}{2}M_\eta
\frac{8I+HH^{\mathsf T}}9
M_\eta^{\mathsf T}.}
\label{md:antmem:cuantica:eq:5.6}
\end{equation}

Sa valeur symplectique normalisée est

\begin{equation}
\boxed{\nu(H):=\frac{2\sqrt{\det V_{\rm vis}}}{\hbar_a}
=\sqrt{1+\frac8{81}\left[\operatorname{tr}(HH^{\mathsf T})-2\right]}.}
\label{md:antmem:cuantica:eq:5.7}
\end{equation}
\end{theorem}

\begin{proof}
Soient \(T=(8I+H)/9\) et \(D=\sqrt8(H-I)/9\). Les deux modes initiaux sont indépendants et ont la même covariance. Dans la coordinatisation équilibrée de l’état initial, la covariance visible normalisée est \(TT^{\mathsf T}+DD^{\mathsf T}\). Le développement donne

\[
\begin{aligned}
TT^{\mathsf T}+DD^{\mathsf T}
&=\frac{64I+8(H+H^{\mathsf T})+HH^{\mathsf T}
+8[HH^{\mathsf T}-H-H^{\mathsf T}+I]}{81}\\
&=\frac{8I+HH^{\mathsf T}}9.
\end{aligned}
\]

Cela démontre \eqref{md:antmem:cuantica:eq:5.6}. Puisque \(\det M_\eta=1\) et \(\det(HH^{\mathsf T})=1\),
\(\det(8I+HH^{\mathsf T})=65+8\operatorname{tr}(HH^{\mathsf T})\).
La division par \(81\) démontre \eqref{md:antmem:cuantica:eq:5.7}. Les deux valeurs propres positives de \(HH^{\mathsf T}\) sont réciproques ; leur somme est au moins deux. En conséquence, \(\nu\geq1\), avec égalité exactement lorsque \(HH^{\mathsf T}=I\).
\end{proof}

L’échelle \(\hbar_a\) et la déformation \(\eta\) sont conservées jusqu’au calcul du déterminant. Leur simplification est covariante car l’état et l’opérateur sont transportés conjointement. L’expression \(HH^{\mathsf T}\) utilise le produit euclidien de la coordinatisation équilibrée de l’état initial ; un changement de coordinatisation transporte également ce produit métrique.

Un transport \(H\) orthogonal et symplectique peut avoir \(D\neq0\) tout en conservant le produit gaussien isotrope, avec \(\nu=1\). Le complément opératoire et l’intrication de l’état ont donc des critères distincts. Dans cette préparation, la déformation relative non isométrique caractérise exactement le cas \(\nu>1\).

\subsubsection{Évaluation exacte du lacet d’or}

Pour \eqref{md:antmem:cuantica:eq:5.2},

\begin{equation}
\operatorname{tr}(H_nH_n^{\mathsf T})-2
=n^2(2n^2-2n+5),\qquad
\nu_n^2=1+\frac8{81}n^2(2n^2-2n+5).
\label{md:antmem:cuantica:eq:5.8}
\end{equation}

L’identité résulte de la somme des carrés des quatre entrées. Pour \(n=1\), la somme est sept et

\begin{equation}
\boxed{\nu_1=\frac{11}{9},\quad
\operatorname{Tr}\rho_{\rm vis}^2=\frac9{11},\quad
\bar n_{\rm W}=\frac19,\quad
\frac{S_{\rm vis}}{k_B}=\frac{10}{9}\log10-\log9.}
\label{md:antmem:cuantica:eq:5.9}
\end{equation}

Le théorème~\ref{md:antmem:gauss:thm:espectro-gaussiano} démontre que l’état gaussien centré d’un mode de valeur symplectique \(\nu\) est unitairement équivalent à la densité géométrique de \eqref{md:antmem:gauss:eq:estado-geometrico}, dont l’occupation moyenne est \(\bar n_{\rm W}=(\nu-1)/2\). Pour \(\nu=11/9\), ses valeurs propres sont

\begin{equation}
\lambda_j=\frac9{10}\,10^{-j},\qquad j=0,1,2,\ldots.
\label{md:antmem:cuantica:eq:5.10}
\end{equation}

La série a pour somme un. Son carré a pour somme
\((9/10)^2/(1-10^{-2})=9/11\). Enfin,


\[
-\sum_j\lambda_j\log\lambda_j
=-\log(9/10)+\log10\sum_jj\lambda_j
=\frac{10}{9}\log10-\log9.
\]

La réduction gaussienne et son spectre sont démontrés dans la section~\ref{md:antmem:gauss:gaussianas-espectro} ; les deux séries précédentes évaluent exactement la pureté et l’entropie de cette réalisation. Comme l’état conjoint est pur, \(S_{\rm vis}/k_B\) est également son entropie adimensionnelle d’intrication. La transformation globale est unitaire, tandis que la restriction au premier mode a une entropie positive.

Le résultat correspond à la préparation et au transport spécifiés. Son identification à une mesure du vide requiert la réalisation expérimentale de ce protocole. L’entropie d’une réduction et la chaleur transférée restent des grandeurs distinctes.


% Correspondencia editorial: gestion/mapa_transportes_auxiliares.md.
\section{Mémoire opérationnelle et invariants d'une structure évolutive}
\label{md:lectin:seccion}

\subsection{Valeur observable, capacité d'action et compatibilité}
\label{md:lectin:valor}

Connaître un objet exige de déterminer ce qui conserve sa valeur observable et quels aspects de sa capacité d'action requièrent une description supplémentaire. Dans la construction HMT, deux lectures présentes peuvent coïncider tandis que l'orientation et la mémoire conservées produisent des réponses distinctes à une même continuation. La généalogie acquiert un contenu physico-mathématique lorsqu'elle détermine ces différences de réponse.

Cette recherche a pour origine la composition APP--TRIT--TPK, son état enrichi et la structure discrète conjointe. Les coordonnées $\pi,\varphi,e,\alpha$ conservent leur fonction de sorties corrélées. Les lecteurs de capacité, de réciprocité et d'action interviennent ensuite, dans leurs domaines propres. Les compositions suivantes examinent ces lecteurs et maintiennent le générateur comme antécédent.

Une caractérisation conceptuelle des constantes consiste à les étudier comme invariants de compatibilité d'une structure qui évolue. Cette caractérisation exige de préciser, pour chaque constante, les transformations sous lesquelles elle demeure invariante, l'information qui peut varier et le lecteur qui réalise son identification. La réciprocité radiale fournit un cas démontré : la section d'action reste invariante lorsque le rayon est inversé et que le changement correspondant d'orientation est conservé. Une extension de cette invariance à l'holonomie complète requiert la composition explicite des opérations qui interviennent.

\subsection{Mémoire conservée et mémoire opératoirement visible}
\label{md:lectin:memoria-visible}

Dans la représentation conjointe de la lecture et de la mémoire, la mise à jour de la lecture est $x'=Tx+Dm$. Pour deux états de même lecture présente, on obtient
\begin{equation}
 \Delta x'=D\,\Delta m.
 \label{md:lectin:diferencia-memoria}
\end{equation}
La différence généalogique est visible à cette étape exactement lorsque $D\Delta m\ne0$. L'identité identifie la continuation qui transforme une différence de mémoire en différence de lecture.

De façon complémentaire, pour un transport stationnaire $C$ de cette représentation, on a
\begin{equation}
 D=\frac{\sqrt8}{9}(C-I),\qquad R=\frac{I+8C}{9}.
 \label{md:lectin:bloques-memoria}
\end{equation}
Si $C\Delta m=\Delta m$, alors $D\Delta m=0$ et $R\Delta m=\Delta m$. Cette différence de mémoire reste invisible lors des répétitions du même protocole. Une opération différente peut la transférer à la lecture.

Le caractère suffisant d'une description concerne donc les différences qu'elle doit conserver pour prédire les réponses aux opérations admises. La généalogie complète peut contenir davantage d'information que n'en requiert une observation particulière. La portée de l'observateur détermine quelle réduction conserve une capacité prédictive. Cette réduction est étudiée à partir des transports produits par HMT et maintient la provenance de ses états et opérateurs.

Dans le cas stationnaire, les différences de mémoire invisibles sont exactement $\operatorname{Fix}(C)$. Pour une collection d'opérations disponibles, les différences invisibles sous toutes les continuations forment l'intersection de leurs espaces fixes. La nécessité s'obtient en prenant chaque opération comme première étape. La suffisance se démontre par récurrence : chaque opération laisse fixe la différence considérée et son bloc de transfert annule cette différence, de sorte qu'elle n'apparaît jamais dans la lecture. Le résultat concerne cette représentation linéaire ; le compteur chronologique de l'holonomie conserve sa propre application vers elle.

\subsection{Régularité des opérations et prolongement de l'état}
\label{md:lectin:regularidad}

Les cadences de vacances distinguent plusieurs niveaux de clôture. Un régime TRIT peut revenir tandis que la phase de capacité demeure décalée. Une phase nonadique peut revenir tandis que le compteur de tours augmente. La répétition d'un type d'opération est compatible avec la production d'un nouvel historique.

Cette distinction permet d'étudier la régularité d'un prolongement illimité : une règle finie peut ordonner l'évolution sans que chaque nouvel état reproduise un état antérieur. La finitude de la règle et celle de l'ensemble des états accessibles sont des propriétés différentes. Une lecture décimale apériodique peut procéder d'une dynamique soumise à des relations strictes.

Dans les régions corrélées, l'objet étudié est la relation qui détermine conjointement le bloc suivant de chaque région, les propriétés conservées par le miroir et la mémoire qui rend compatibles les prolongements. Les cadences $21/22$ et $6/7$ sont des observations de capacité ; leurs applications concrètes les relient aux lecteurs régionaux. La preuve à profondeur arbitraire et l'explication de la structure partagée constituent des résultats distincts et complémentaires.

La composition des cadences sépare exactement les niveaux : achever les trois changements du sélecteur laisse un déplacement de capacité de trois ou six classes modulo neuf. Les durées $437$ et $459$ quantifient cette distinction dans la construction correspondante ; leur fonction est ici de décrire ces cadences, sans les promouvoir au rang de constantes fondamentales supplémentaires.

\subsection{Caractère d'action et restitution radiale orientée}
\label{md:lectin:h5}

Le caractère $H_5$ et la continuation régionale déterminent une section d'action. La forme radiale permet de restituer cette section lorsque l'orientation est conservée. À $\alpha$ fixé, la coordonnée $\chi_{\mathrm{rad}}=\varepsilon\log q_\varepsilon$ classe la réciprocité et entretient une relation monotone avec la section de retour :
\begin{equation}
 \frac{\eta_{\mathrm{ret}}}{\alpha}
 =-500\tanh\!\left(\frac{\chi_{\mathrm{rad}}}{2}\right)-1.
 \label{md:lectin:accion-radial}
\end{equation}
La forme radiale orientée et la section d'action sont soumises à une restriction conjointe. L'action incorpore une condition de compatibilité géométrique ; ses coefficients, sa correction régionale et son orientation doivent rester coordonnés.

Cette relation distingue deux niveaux de connaissance. La reparamétrisation démontre l'équivalence des descriptions et permet de retrouver l'une à partir de l'autre. Une confrontation entre lecteurs apporte une information supplémentaire lorsque les deux sont évalués à partir de la structure antérieure sans introduire en entrée la valeur vérifiée. Calculer le rayon à l'aide de l'inverse de la formule d'action puis le substituer à nouveau vérifie une identité algébrique. Produire la lecture géométrique par ses opérations antécédentes et la comparer au caractère d'action met à l'épreuve la composition causale. Les deux vérifications possèdent leurs fonctions propres.

Cette distinction identifie les développements à réunir et à exécuter ; elle maintient la généalogie HMT comme racine commune. L'indépendance de la confrontation est exigée par rapport à la valeur vérifiée, non par rapport à l'origine commune de ses lecteurs.

\subsection{Action, information et température}
\label{md:lectin:termica}

L'état des continuations admissibles et leurs poids déterminent une information sans dimension. Le transducteur $k_B$ en réalise la lecture en unités d'entropie. L'action et la période interviennent dans
\begin{equation}
 k_B\Theta_{\mathrm{clk}}\log3=\frac{h}{T_{108}}.
 \label{md:lectin:accion-informacion}
\end{equation}
Les trois possibilités équiprobables produisent $\log3$ ; une distribution non uniforme conserve ses poids dans l'évaluation de l'information.

La composition permet de suivre comment une même évolution modifie la mémoire opérationnelle, les continuations compatibles et la lecture énergétique. La longueur d'un historique et sa température correspondent à des lectures différentes. L'évaluation thermique requiert d'identifier l'état statistique et l'opération d'échange qui la réalisent. Le problème précis consiste à déterminer quelles transformations engendrées conservent l'action, lesquelles redistribuent l'information accessible et lesquelles produisent une différence thermique calculable.

Lorsqu'un parcours se ferme dans une projection tandis que sa mémoire change, l'évaluation de son travail exige de déclarer l'espace dans lequel le cycle est considéré comme complet. Tel est le problème du caractère conservatif. La non-commutativité caractérise, en revanche, la sensibilité à l'ordre des opérations. Conserver cette distinction permet d'étudier la relation entre les deux propriétés.

\subsection{Programme de confrontation et hiérarchie des vérifications}
\label{md:lectin:contrastes}

\paragraph{Historiques de lecture commune et de réponses futures distinctes.}
La première confrontation réunit deux parcours admissibles produits par TPK qui coïncident pour un lecteur, conserve leurs mémoires et applique la même continuation. La grandeur examinée est la séparation exacte de leurs lectures ultérieures et la composante de mémoire qui la détermine. Ainsi, la différence entre valeur et structure s'exprime par une différence de capacité prédictive. Les exemples déjà construits maintiennent leur provenance et peuvent être prolongés dans ce schéma.

\paragraph{Compatibilité entre action régionale et forme géométrique.}
La deuxième confrontation suit le lecteur de $H_5$ et sa continuation, ainsi que le lecteur géométrique avec ses antécédents propres. Leur compatibilité est examinée avant et après un retour, en distinguant la section limite de ses approximations finies. La confrontation identifie ce qui demeure constant tandis que la mémoire progresse. L'invariance radiale démontrée fournit le cas de référence et conserve sa portée concrète.

\paragraph{Ordre des opérations et réentrée de mémoire.}
La troisième confrontation emploie deux opérations effectives du TPK et maintient les coordinatisations intermédiaires. La répartition $1:8$ de la réalisation auxiliaire correspond au régime d'une fibre ayant l'identité pour référence. L'application doit vérifier ce régime ou transporter les coordinatisations correspondantes. L'expérience mathématique compare les deux ordres avec et sans réentrée de mémoire ; la réalisation physique doit spécifier la préparation, l'interaction et l'observable qui représentent ces opérations.

Les deux premières confrontations fixent la priorité conceptuelle : elles déterminent quelle information possède une efficacité dynamique et quels invariants soutiennent la compatibilité de l'architecture. La troisième fournit une réponse quantitative directe dans son régime démontré. Cette hiérarchie conserve le rôle auxiliaire de la réalisation énergétique.

\subsection{Comparaison scientifique et portée de l'interprétation}
\label{md:lectin:comparacion}

La mémoire d'une dynamique réduite possède des antécédents en physique mathématique. Le formalisme de Mori--Zwanzig permet d'étudier l'évolution et les termes de mémoire des observables réduites ; l'analyse de Hudson et Li en fournit un exemple.\footnote{Hudson et Li,
\emph{Coarse-graining of overdamped Langevin dynamics via the
Mori--Zwanzig formalism},
\href{https://arxiv.org/abs/1810.08175}{arXiv:1810.08175}.}
La comparaison pertinente porte sur la composition concrète examinée ici : origine APP--TRIT--TPK, régions corrélées, action, orientation et restitution. La seule mention de la mémoire ne détermine pas la priorité de cette composition.

L'apport de l'exposition consiste à préciser les conséquences et les confrontations au sein de la construction. Une évaluation historique de la nouveauté et une corroboration physique conservent leurs preuves spécifiques. Le programme mathématique établit des objets et des opérations pouvant être soumis à comparaison scientifique sans anticiper le résultat de ces évaluations.

La mémoire stockée et la mémoire opératoirement visible sont ici caractérisées au moyen du noyau de $D$ et de l'espace fixe de $C$. Cette caractérisation est une conséquence de la mise à jour déclarée et conserve la portée de sa représentation. Les trois confrontations de la section~\ref{md:lectin:contrastes} constituent un programme de travail ; leur formulation distingue la définition de l'expérience mathématique, l'exécution des parcours TPK et la réalisation d'une expérience physique.

% Fuente conservada íntegra: originales/cuaderno/NOTA.md, §§1--7.
\section{Mémoire chronologique, bilan et réponse des vacances}
\label{md:cron:seccion}

La construction suivante reçoit le facteur symbolique de capacité et étudie ses lecteurs ultérieurs. Le rapport interne $3^6/10^3=729/1000$ et son calendrier appartiennent à la construction précédente (sections~\ref{mono-ampl:relojes} et~\ref{vac-capacidad-trit}). L'architecture APP--TRIT--TPK, les régions, le retour nonadique et le registre $K$ conservent leur position causale antérieure. Les logarithmes employés dans le lecteur de capacité décrivent ce facteur et ne sélectionnent pas à nouveau les régions de $\pi,\varphi,e,\alpha$. Les propositions sont des conséquences focales concernant l'observation et la mémoire ; leur portée se distingue de l'état TPK complet et de l'attribution d'une priorité historique à des résultats généraux sur les suites mécaniques ou l'information.

\subsection{Objet, mesure et deux lecteurs}
\label{md:cron:lectores}

Soit $\delta=\log_{10}(10/9)$ la pente du calendrier de capacité, et soit $\theta$ une phase uniformément distribuée dans $[0,1)$. La mesure représente l'incertitude sur l'origine de lecture, tandis que le mécanisme demeure déterminé. Nous définissons
\begin{equation}
 \begin{aligned}
 z_j(\theta)&=\lfloor(j+1)\delta+\theta\rfloor
             -\lfloor j\delta+\theta\rfloor,\\
 W_n(\theta)&=(z_0(\theta),\ldots,z_{n-1}(\theta)),\\
 B_n(\theta)&=\sum_{j=0}^{n-1}z_j(\theta).
 \end{aligned}
 \label{md:cron:dos-lectores}
\end{equation}
Le bloc $W_n$ conserve la chronologie et $B_n$ conserve le bilan. Pour une phase exacte et une règle exacte connues, les deux objets sont déterministes ; les entropies se rapportent expressément à l'ensemble des phases possibles.

La sommation télescopique donne
\begin{equation}
 B_n(\theta)=\lfloor n\delta+\theta\rfloor.
 \label{md:cron:telescopia}
\end{equation}
Ainsi, $B_n$ n'admet que les valeurs $\lfloor n\delta\rfloor$ et $\lceil n\delta\rceil$. Les points $\{-j\delta\}$, avec $j=0,\ldots,n$, divisent le cercle en $n+1$ intervalles de phase. Chaque intervalle porte une suite distincte de longueur $n$. En effet, les préfixes cumulés $\lfloor k\delta+\theta\rfloor$, pour $k=1,\ldots,n$, sont monotones en $\theta$, et chaque frontière modifie l'un d'eux ; la suite reconstruit tous ces préfixes.

Les extrémités sont de mesure nulle. Une convention semi-ouverte résout leur assignation et conserve les entropies.

\subsection{Information croissante et taux entropique nul}
\label{md:cron:crecimiento}

Soient $\ell_0,\ldots,\ell_n$ les longueurs des intervalles de phase et
\begin{equation}
 H_n=-\sum_{j=0}^{n}\ell_j\log_2\ell_j.
 \label{md:cron:entropia-bloque}
\end{equation}

\begin{proposition}
\label{md:cron:proposicion-tasa}
On a simultanément
\begin{equation}
 H_n\longrightarrow+\infty,\qquad
 \frac{H_n}{n}\longrightarrow0.
 \label{md:cron:dos-limites}
\end{equation}
\end{proposition}
\begin{proof}
La borne supérieure est $H_n\leq\log_2(n+1)$. Si $\ell_{\max}(n)=\max_j\ell_j$, alors $H_n\geq-\log_2\ell_{\max}(n)$. La pente $\delta$ est irrationnelle : si $\delta=p/q$ avec $q>0$, on aurait $(10/9)^q=10^p$, ce qui contredit la valuation du nombre premier $3$. Par irrationalité, les points de l'orbite sont denses. Les partitions successives ont un pas maximal tendant vers zéro : pour chaque $\epsilon>0$, il existe un segment fini $\epsilon$-dense, et les segments ultérieurs ne font que subdiviser ses intervalles. La borne inférieure diverge donc. La borne supérieure divisée par $n$ tend vers zéro.
\end{proof}

Les longueurs des intervalles ne sont pas égales en général. La preuve établit les bornes précédentes sans affirmer $H_n=\log_2(n+1)$ ni une loi asymptotique logarithmique exacte qui requerrait des conditions diophantiennes supplémentaires. L'entropie topologique du facteur est également nulle parce que $\log(n+1)/n\to0$. L'information du bloc et son taux sont des quantités différentes. La croissance de l'archive d'une trajectoire connue n'équivaut pas, à elle seule, à une production d'information statistique indépendante.

\subsection{Information chronologique non déterminée par le bilan}
\label{md:cron:condicional}

\begin{corollary}
\label{md:cron:corolario-condicional}
Pour le même ensemble de phases,
\begin{equation}
 H(W_n\mid B_n)=H_n-H(B_n)
 \geq H_n-1\longrightarrow+\infty.
 \label{md:cron:entropia-condicional}
\end{equation}
\end{corollary}
\begin{proof}
$B_n$ est une fonction de $W_n$, de sorte que la règle de chaîne donne l'égalité. Son support possède deux éléments, d'où $H(B_n)\leq1$ bit. On applique la proposition~\ref{md:cron:proposicion-tasa}.
\end{proof}

L'entropie conditionnelle divergente est la moyenne sur les bilans ; on n'attribue pas la même incertitude à chaque bilan concret. Cette perte n'est pas non plus attribuée à l'état enrichi complet : lorsque la mémoire retenue restitue $W_n$, on a
\begin{equation}
 H(W_n\mid B_n,M_n)=0.
 \label{md:cron:restitucion-memoria}
\end{equation}
Par cardinalité, au moins une des deux fibres de $B_n$ contient $\lceil(n+1)/2\rceil$ suites ou davantage. Cette borne combinatoire se distingue de la borne entropique. Le résultat précise la relation entre valeur et généalogie dans ce lecteur : les deux bilans possibles n'épuisent pas les historiques temporels compatibles.

\subsection{Courant centré : bilan borné et activité persistante}
\label{md:cron:corriente}

Le lecteur centré s'exprime par
\begin{equation}
 J_j=z_j-\delta,\qquad
 \sum_{j=0}^{n-1}J_j=B_n-n\delta.
 \label{md:cron:corriente-definicion}
\end{equation}
Par conséquent, $\left|\sum_{j=0}^{n-1}J_j\right|<1$ pour toute fenêtre. Simultanément,
\begin{equation}
 \lim_{n\to\infty}\frac1n\sum_{j=0}^{n-1}J_j=0,\qquad
 \lim_{n\to\infty}\frac1n\sum_{j=0}^{n-1}J_j^2
 =\delta(1-\delta)>0.
 \label{md:cron:actividad}
\end{equation}
\begin{proof}
On a $z_j^2=z_j$ et $B_n/n\to\delta$ par la borne de discrépance. Le développement $(z_j-\delta)^2=z_j(1-2\delta)+\delta^2$, suivi de la substitution de la moyenne de $z_j$, donne la seconde identité. L'argument vaut pour toute phase et ne requiert pas de moyenne aléatoire.
\end{proof}

Le bilan accumulé borné, l'activité quadratique positive et le taux d'entropie nul sont donc compatibles. Les identités conservent la distinction entre le lecteur centré et une représentation de charge élémentaire ; elles n'établissent pas non plus de production énergétique sans bilan.

Si une réalisation dimensionnelle déjà établie prend $I_j=(u_Q/t_0)J_j$, on obtient
\begin{equation}
 I_{\mathrm{rms}}^2
 =\left(\frac{u_Q}{t_0}\right)^2\delta(1-\delta).
 \label{md:cron:eficaz}
\end{equation}
Si chaque $I_j$ est maintenu pendant un intervalle de durée $t_0$, cette moyenne discrète coïncide avec la moyenne temporelle. À titre de comparaison physique ultérieure, une charge résistive idéale $R_{\mathrm{res}}>0$ aurait pour puissance moyenne
\begin{equation}
 \overline P_{\mathrm{res}}
 =R_{\mathrm{res}}\left(\frac{u_Q}{t_0}\right)^2
 \delta(1-\delta)>0.
 \label{md:cron:potencia-resistiva}
\end{equation}
C'est une prédiction du modèle composé avec cette charge. La résistance et le dispositif conservent leur propre spécification ; ils ne se dérivent pas de la seule suite de capacité.

\subsection{Spectre du courant et orientation}
\label{md:cron:espectro}

La fonction d'observation est $f(\theta)=\mathbf1_{[1-\delta,1)}(\theta)-\delta$, avec $J_j=f(\theta+j\delta)$. Ses coefficients de Fourier, pour $m\neq0$, sont
\begin{equation}
 a_m=\frac{\exp(2\pi i m\delta)-1}{2\pi i m},\qquad
 |a_m|^2=\frac{\sin^2(\pi m\delta)}{\pi^2m^2},\qquad
 a_0=0.
 \label{md:cron:fourier}
\end{equation}
Ils s'obtiennent en intégrant l'indicatrice sur son intervalle. L'identité de Parseval donne
\begin{equation}
 \sum_{m\neq0}|a_m|^2=\delta(1-\delta),
 \label{md:cron:parseval}
\end{equation}
en accord avec l'activité quadratique. Les fréquences temporelles sont $m\delta$ modulo $1$. La corrélation stationnaire, avec phase uniforme, est déterminée par ces intensités ; celles-ci ne contiennent pas la phase initiale. Une translation multiplie $a_m$ par une phase unitaire, et l'inversion temporelle échange les fréquences opposées.

La puissance spectrale isolée ne restitue ni l'orientation ni l'origine temporelle. Une lecture sensible à la phase ou une référence marquée apporte l'information complémentaire. La déduction conserve la distinction entre cette référence générale et le registre $K$, ainsi qu'entre les rythmes symboliques et une réalisation expérimentale déterminée.

\subsubsection{Référence cyclique et transports de même intensité spectrale}
\label{md:cron:k-ciclico}

La propriété cyclique du registre $K$ permet une conséquence finie
indépendante du spectre précédent. Dans $V_K=\mathbb R^{12}$, soit $S$ le
décalage $(Sx)_i=x_{i+1}$ et
\begin{equation}
 O_K=(K,SK,\ldots,S^{11}K).
 \label{md:cron:matriz-ciclica}
\end{equation}
On utilise l’inversibilité de $O_K$ établie par la référence cyclique
du registre. Pour tout opérateur réel $H$ qui commute avec $S$,
posons $y_+=HK$ et $y_-=H^TK$.

\begin{proposition}
\label{md:cron:proposicion-intensidades}
Les douze intensités de Fourier de $y_+$ et $y_-$ coïncident.
Cependant, $y_+=y_-$ si et seulement si $H=H^T$.
Chaque réponse vectorielle complète, avec $K$ et l’orientation de $S$,
retrouve son opérateur.
\end{proposition}
\begin{proof}
$H$ est circulant. Ses multiplicateurs de Fourier sont $\lambda_m$ et ceux
de $H^T$ sont leurs conjugués, de sorte que les deux réponses ont
les intensités $|\lambda_m|^2|\widehat K_m|^2$.
Si leurs vecteurs coïncident, $(H-H^T)K=0$. La commutation implique que
$H-H^T$ annule également $S^jK$ pour tout $j$. Ces vecteurs forment une base ;
par conséquent, $H-H^T=0$. L’implication inverse est immédiate.
La récupération est exprimée par
\begin{equation}
 H=O_{y_+}O_K^{-1},\qquad
 H^T=O_{y_-}O_K^{-1}.
 \label{md:cron:restitucion-operador}
\end{equation}
\end{proof}

En particulier, les décalages $S$ et $S^{-1}$ produisent des réponses de
même spectre de puissance, mais distinctes. Les autocorrélations exactes
du registre donnent
\begin{equation}
 \|SK-S^{-1}K\|^2
 =2\bigl(\|K\|^2-\langle K,S^2K\rangle\bigr)
 =2175744>0.
 \label{md:cron:separacion-desplazamientos}
\end{equation}
L’inversion rationnelle de $O_K$ retrouve à partir des deux réponses les
coefficients $e_1$ et $e_{11}$, respectivement. $K$ est utilisé après
sa génération comme référence marquée. L’intensité spectrale ou le
bilan scalaire ne remplacent pas la réponse vectorielle.
Le corollaire décrit le support cyclique à douze composantes ;
il conserve la différence entre $S$ et la dynamique enrichie complète,
ainsi que les opérateurs supplémentaires d’une conjugaison physique de charge.
Pour un $H$ général, la transposition ne signifie pas non plus une inversion temporelle ;
dans le cas $S$, elle coïncide bien avec $S^{-1}$.

\subsection{Lien opératoriel électronique}
\label{md:cron:electron}

L’incidence entre les fibres de somme et de produit d’APP sur $D_9^3$
définit $C$. Son vecteur tritique
$m_{\mathrm{ph}}=(1,-1,0)^3$ satisfait
\begin{equation}
 Cm_{\mathrm{ph}}=C^Tm_{\mathrm{ph}}
 =-\frac12m_{\mathrm{ph}}.
 \label{md:cron:modo-tritico}
\end{equation}
Le calendrier induit parcourt les régimes du même sélecteur tritique ;
l’incidence sélectionne le plan principal de singularité $1/2$.
Dans celui-ci,
\begin{equation}
 P=\begin{pmatrix}1&0\\0&0\end{pmatrix},\qquad
 Q=\begin{pmatrix}1/4&\sqrt3/4\\\sqrt3/4&3/4\end{pmatrix},
 \qquad J=\frac4{\sqrt3}[P,Q],\qquad J^2=-I.
 \label{md:cron:plano-electronico}
\end{equation}
L’incidence et la norme du commutateur sont développées dans la
section~\ref{vi:dep:prefactor}. On retrouve ainsi une relation effective entre
l’organisation temporelle et le bloc électronique, en conservant le sélecteur,
ses étiquettes et l’incidence. La représentation de charge et la réalisation
dimensionnelle maintiennent leurs opérateurs propres : le lien n’assigne pas
directement les charges élémentaires opposées à $z_j-\delta$.

\subsection{Information retenue et composition énergétique}
\label{md:cron:informacion-energia}

La récupération de fibre établit
\begin{equation}
 H(X\mid q(X),M(X))=0
 \label{md:cron:inversa-enriquecida}
\end{equation}
lorsque le lecteur enrichi a un inverse ; une mise à jour bijective
conserve $H(X)$. Ces propriétés sont réunies dans
\eqref{eq:ii-kb-aut-fibra} et~\eqref{eq:ii-kb-aut-reversible}.
La coexistence de $H_n$ croissant et d’un taux nul est compatible avec ces
propriétés et avec les identités de Shannon.

Le travail d’effacement dépend de l’information accessible retenue.
Dans le régime idéal isotherme des mémoires logiques dégénérées et dans la
limite asymptotique appropriée, le coût par copie avec information latérale est
\begin{equation}
 W_{\mathrm{borrado,ideal}}
 =k_BT\ln(2)\,H(X\mid M).
 \label{md:cron:borrado}
\end{equation}
Si $X$ est retrouvé à partir de $M$, la contribution logique idéale peut être nulle.
Le transport physique, la lecture, le contrôle et le cycle complet conservent
des bilans propres ; l’égalité $h_{\mathrm{top}}=0$ ne les élimine pas.

La composition pertinente conserve la séquence
\[
 \begin{gathered}
 \text{chronologie}\longrightarrow
 \text{registre de frontière}\longrightarrow
 \text{lecteur de courant et d’état}\\
 \longrightarrow
 \text{mémoire conservée}\longrightarrow
 \text{bilan thermique}.
 \end{gathered}
\]
Le développement précédent construit les deux premiers lecteurs et leurs
conséquences informationnelles. Sa portée est distinguée de la certification
d’un dispositif physique complet et maintient intacte la généalogie
constructive des constantes.

\section{Répétition visible et différenciation généalogique dans une région}
\label{md:ejemplo:region}
Les biographies de \eqref{act21:eq:biografias-tpk-tres-caracteres}
permettent de préciser la différence entre une répétition de chiffres et un
retour de l’état. Pour le canal reconnu comme \(e\), le
calendrier \eqref{act21:eq:calendario-0011-interno} produit
\[
 201101\mid121221\mid102011\mid012222\mid102011.
\]
L’élévation précède la lecture décimale. Les mêmes matrices agissent
dans les deux autres canaux, en conservant leurs émissions initiales et leurs
états régionaux distincts. Cette structure commune permet de comparer
les biographies sans identifier leurs coordonnées.

Soit \(w\) une concaténation de \(L\) trits, et \(N\) sa valeur entière
en base trois. Le cylindre correspondant est
\[
 I_w=[N/3^L,(N+1)/3^L).
\]
Pour \(d\) positions décimales, le préfixe est stable sur tout le
cylindre exactement lorsque
\[
 \left\lfloor\frac{10^dN}{3^L}\right\rfloor
 =\left\lfloor\frac{10^d(N+1)-1}{3^L}\right\rfloor.
\]
C’est l’application du lemme~\ref{act21:censo:estabilidad} au lecteur
décimal. L’extrémité droite semi-ouverte explique le terme \(-1\)
de la division entière. En ajoutant la partie entière de la coordinatisation régionale,
les cinq lectures successives sont
\begin{center}
\begin{tabular}{@{}rl@{}}
\toprule
Longueur tritique & Préfixe stable de la lecture régionale\\
\midrule
6 & \(2.71\ldots\)\\
12 & \(2.71828\ldots\)\\
18 & \(2.7182818\ldots\)\\
24 & \(2.7182818284\ldots\)\\
30 & \(2.71828182845904\ldots\)\\
\bottomrule
\end{tabular}
\end{center}
À la longueur vingt-quatre, on a
\(N=202864003874\) et \(3^{24}=282429536481\).
À la longueur trente, on obtient
\(N=147887858824447\) et \(3^{30}=205891132094649\).
Les divisions entières précédentes certifient les préfixes du tableau.
En particulier, les douze premiers chiffres fractionnaires se regroupent comme
\[
 718281828459
 =7\,\underbrace{1828\,1828}_{\text{réitération positionnelle}}\,459.
\]
Le cylindre de vingt-quatre trits fixe déjà les deux copies et le chiffre
suivant. Cette localisation distingue l’épisode décimal de la
réapparition du bloc tritique \(102011\), situé aux positions
troisième et cinquième : leurs coupures et leurs bases de lecture sont différentes.

La répétition du bloc visible n’identifie pas non plus ses préétats.
Les préétats modulo vingt-sept conservés dans le développement de
cette biographie sont
\[
 (25,11,7,17,8,16),\qquad (7,8,4,23,26,7),
\]
avec pour quotients ultérieurs respectifs
\[
 (6,13,9,2,13,1),\qquad (15,0,3,19,23,25).
\]
L’égalité de l’émission coexiste ainsi avec une généalogie différenciée.
Sa fonction explicative est concrète : l’observation d’un bloc
répété conserve moins d’information que la continuation enrichie
qui le produit. Pour attribuer une différence future exclusivement
à la mémoire, il faut également conserver les phases et les coordinatisations de
comparaison.

Cet épisode relie trois niveaux : les transitions produisent la
chaîne tritique, les cylindres déterminent les chiffres stables et
l’état conserve la provenance d’une émission répétée. La
profondeur arbitraire appartient au prolongement démontré
précédemment ; cette lecture finie précise un épisode de cette
biographie. La régularité qu’il importe d’étudier réside dans la
compatibilité entre les niveaux, non dans le prolongement indéfini
d’un fragment décimal : une troisième copie de \(1828\) commencerait par
\(1\), tandis que la continuation déjà produite commence par \(4\).

% Fuente íntegra: originales/cuaderno/REORIENTACION_GENEALOGICA_Y_RESULTADOS.md,
% §§2--7. Gestión y procedencia conservadas aparte.
\Needspace{16\baselineskip}
\section{Régularité généalogique, restitution opératorielle et géométrie de l'action}
\label{md:restop:seccion}

\subsection{Production, compatibilité et restitution}
\label{md:restop:jerarquia}

Le phénomène primaire est une dynamique qui conserve la provenance en produisant de nouvelles observations. Trois opérations aux fonctions mathématiques différenciées sont réunies :
\begin{enumerate}
 \item \emph{Production.} La mise à jour APP--TRIT--TPK prolonge l'état et ses régions corrélées. La phase peut revenir tandis que la mémoire progresse.
 \item \emph{Compatibilité.} Les lectures régionales, angulaires, d'action et d'incidence obéissent à des relations communes. Leur variation doit respecter ces relations pour correspondre à un même état.
 \item \emph{Restitution.} La conservation de l'information complémentaire permet de retrouver les valeurs antérieures, les produits d'opérateurs et les différences entre ordres de composition.
\end{enumerate}
La comparaison entre procédures pour obtenir $\pi$ doit considérer ces trois niveaux. La thèse spécifique de HMT réside dans la construction conjointe et ses relations, outre la production de développements numériques. Une coïncidence décimale ou la périodicité d'une phase ne décrivent pas à elles seules la portée de la construction. Cette distinction identifie l'objet mathématique pertinent pour une comparaison historique ; le développement focal n'établit pas à lui seul une priorité mondiale.

\subsection{Régimes de vacance et retour de la phase de capacité}
\label{md:restop:vacancias}

\subsubsection{Calendrier de capacité et sélecteur tritique}
\label{md:restop:antecedentes-capacidad}

Le lecteur de capacité reçoit la compactification interne $729/1000$ et
distingue l’indice de raffinement de l’indice de capacité. Avec la
convention de la section~\ref{vac-capacidad-trit},
\begin{equation}
 \delta=\log_{10}(10/9),\qquad
 p_j=\left\lfloor\frac{j}{\delta}\right\rfloor,\qquad
 v_j=p_j-j,\qquad j\geq1.
 \label{md:restop:capacidad}
\end{equation}
Les intervalles entre événements sont
\begin{equation}
 h_j=p_{j+1}-p_j\in\{21,22\},\qquad
 s_j=22-h_j.
 \label{md:restop:intervalos}
\end{equation}
La classe tritique obéit à
\begin{equation}
 [v_{j+1}]_3=[v_j-s_j]_3.
 \label{md:restop:clase-tritica}
\end{equation}
Un intervalle court change la classe ; un intervalle long la conserve.
Les séparations $6/7$ décrivent les plateaux de ce régime induit.
La construction maintient son type de facteur d’observation de l’état TPK ;
elle ne remplace pas cet état et n’identifie pas ces séparations à des signes de
charge ou à des chiffres de $\pi$ ou de $e$.
Le sélecteur canonique est
$M_{\mathrm{ph}}(1+[v_j]_9)$, dont le motif dépend de la classe modulo trois.

\subsubsection{Restitution du régime local et déplacement modulo neuf}
\label{md:restop:consecuencia-capacidad}

Soit
\begin{equation}
 \beta_{\mathrm{cap}}=22-\delta^{-1},\qquad
 q_k=\left\lfloor\frac{k}{\beta_{\mathrm{cap}}}\right\rfloor,\qquad k\geq1.
 \label{md:restop:pendiente-inducida}
\end{equation}
Les positions $q_k$ sont les indices des intervalles courts. En effet,
\begin{equation}
 p_j=22j-\lceil j\beta_{\mathrm{cap}}\rceil,\qquad
 s_j=\lceil(j+1)\beta_{\mathrm{cap}}\rceil
     -\lceil j\beta_{\mathrm{cap}}\rceil.
 \label{md:restop:eventos-cortos}
\end{equation}
L’irrationalité de $\delta$ s’obtient à partir des valuations de $2,3,5$
dans l’égalité qu’imposerait $(10/9)^q=10^p$ ; elle implique celle de
$\beta_{\mathrm{cap}}$. Ainsi, les indices positifs ne présentent pas
d’ambiguïté aux extrémités. La discontinuité correspondant à l’entier $k$
se produit en $j=\lfloor k/\beta_{\mathrm{cap}}\rfloor$.

\begin{proposition}
\label{md:restop:proposicion-regimen}
Entre les événements d’indices $q_k$ et $q_{k+3}$, en prenant les intervalles dont
les indices appartiennent à $[q_k,q_{k+3})$, la classe TRIT revient, tandis que la phase
de capacité modulo neuf change.
\end{proposition}
\begin{proof}
Il y a exactement trois intervalles courts dans ce segment.
Les bornes $1/7<\beta_{\mathrm{cap}}<3/20$ donnent
$20<3/\beta_{\mathrm{cap}}<21$, de sorte que
\begin{equation}
 m_k=q_{k+3}-q_k\in\{20,21\}.
 \label{md:restop:duracion-tritica}
\end{equation}
La somme des longueurs et le gain de capacité sont, respectivement,
\begin{equation}
 \Delta p=22m_k-3\in\{437,459\},\qquad
 \Delta v=21m_k-3\in\{417,438\}.
 \label{md:restop:ganancias}
\end{equation}
Dans les deux cas, $\Delta v\equiv0\pmod3$. En revanche,
\begin{equation}
 \boxed{\Delta v\equiv3\ \text{ou}\ 6\pmod9.}
 \label{md:restop:desplazamiento-nonadico}
\end{equation}
L’énoncé est démontré.
Les bornes de $\beta_{\mathrm{cap}}$ disposent de certificats entiers :
\begin{equation}
 \begin{aligned}
 10^{146}>9^{153}&\ \Longrightarrow\
 \delta>\frac7{153},\\
 10^{417}<9^{437}&\ \Longrightarrow\
 \delta<\frac{20}{437}.
 \end{aligned}
 \label{md:restop:certificados-cotas}
\end{equation}
La monotonie de $22-1/x$ produit
$1/7<\beta_{\mathrm{cap}}<3/20$. Ainsi, les bornes reposent sur des
inégalités entières et non sur une substitution décimale arrondie.
\end{proof}

La valeur du sélecteur
$\tau_j=M_{\mathrm{ph}}(1+[v_j]_9)$ revient parce qu’elle dépend de la classe
modulo trois. Ce retour conserve sa distinction avec l’état TRIT
complet et avec un tour chronologique de $\Gamma_9$.
Les deux durées apparaissent une infinité de fois par la densité de la rotation
irrationnelle de pente $1/\beta_{\mathrm{cap}}$.

La régularité présente des niveaux distincts : le régime local est restauré
tandis que la phase de capacité reste décalée. Le retour nonadique de
l’état possède, en outre, son propre accroissement de mémoire. Chaque objet
conserve donc sa notion de retour et sa durée.
Les nombres $437$ et $459$ sont des durées de cette observation induite,
sans qu’on leur attribue le statut de constantes universelles supplémentaires.
Le résultat réunit explicitement les formules précédentes ; sa présentation
comme conséquence focale ne constitue pas une affirmation de priorité
historique par rapport à des développements qui n’ont pas été comparés exhaustivement.

\subsection{Participation de la mémoire à la composition opératorielle}
\label{md:restop:memoria}

\subsubsection{Représentation élargie de lecture et de mémoire}
\label{md:restop:estructura-partida}

La représentation de lecture et de mémoire s’obtient en séparant huit
composantes et une neuvième. Dans une même fibre, avec l’identité comme transport de
référence, le transport interne $C$ induit
\begin{equation}
 T_C=\frac{8I+C}{9},\qquad
 D_C=\frac{\sqrt8(C-I)}{9},\qquad
 R_C=\frac{I+8C}{9},
 \label{md:restop:bloques}
\end{equation}
\begin{equation}
 \mathcal U_C=
 \begin{pmatrix}
 T_C&D_C\\D_C&R_C
 \end{pmatrix}.
 \label{md:restop:transporte-ampliado}
\end{equation}
Le couplage provient du changement de base de la décomposition nonadique.
Pour les transports qui conservent la forme déclarée, il conserve la forme
élargie. L’identité de composition est
\begin{equation}
 \mathcal U_{C_2}\mathcal U_{C_1}
 =\mathcal U_{C_2C_1},\qquad
 T_2T_1+D_2D_1=\frac{8I+C_2C_1}{9},
 \label{md:restop:composicion-ampliada}
\end{equation}
avec $T_i=T_{C_i}$ et $D_i=D_{C_i}$.
Le second terme correspond à la mémoire produite lors de la première
opération et réintroduite dans la seconde. L’écarter change la
procédure : cela produit $T_2T_1$ au lieu de la lecture de la composition
élargie.

\subsubsection{Répartition de la sensibilité à l’ordre de composition}
\label{md:restop:conmutadores}

\begin{proposition}
\label{md:restop:proposicion-orden}
Avec la convention $[A,B]=AB-BA$,
\begin{equation}
 [T_2,T_1]=\frac1{81}[C_2,C_1],\qquad
 [D_2,D_1]=\frac8{81}[C_2,C_1].
 \label{md:restop:reparto-conmutadores}
\end{equation}
Par conséquent, la différence de lecture entre les deux ordres, avec la
mémoire réintroduite, est
\begin{equation}
 \boxed{(T_2T_1+D_2D_1)-(T_1T_2+D_1D_2)
 =\frac19[C_2,C_1].}
 \label{md:restop:orden-completo}
\end{equation}
\end{proposition}
\begin{proof}
En développant les produits de $T$, les termes constants et linéaires
s’annulent dans la différence ; il reste le commutateur divisé par $81$.
En développant les produits de $D$, la même annulation se produit et le
facteur est $8/81$. Leur somme est $1/9$.
\end{proof}

Le mécanisme de réentrée fait partie de la représentation élargie.
La comparaison des deux ordres explicite une conséquence :
dans cette représentation, la contribution de la mémoire à la réponse
antisymétrique est huit fois celle retenue par la composition de lectures
isolées. La réponse complète est neuf fois cette dernière. Les
proportions correspondent à des opérateurs de réponse, non à des pourcentages
d’énergie.

La restitution conserve la non-commutativité et son intensité de lecture.
Dans ce cas, la compression isolée atténue un commutateur non nul sans
l’annuler. Une application à des forces physiques conserve en outre les opérateurs
d’interaction et la fonctionnelle de travail correspondante ; l’identité
algébrique ne les identifie pas à elle seule.

\subsubsection{Confrontation opératorielle et normalisation nonadique}
\label{md:restop:contraste}

Pour le motif algébrique général $(n-1)+1$, les mêmes calculs donnent les
facteurs $1/n^2$, $(n-1)/n^2$ et $1/n$. L’architecture nonadique fixe
$n=9$. On distingue ainsi la loi de composition de la valeur concrète de
sa normalisation.

Lorsque les opérations effectives d’une application sont représentées par
$C_1,C_2$, les procédures qui écartent ou réintroduisent la mémoire
prédisent des réponses différentes à l’ordre. La confrontation porte sur la
composition opérationnelle et sur l’accessibilité de la mémoire, en conservant
les coefficients de la représentation.

\subsubsection{Normes de lecture et de mémoire dans l’inversion centrale}
\label{md:restop:treinta-y-dos}

Pour $C=-I$, on obtient
\begin{equation}
 T^*T=\frac{49}{81}I,\qquad
 D^*D=\frac{32}{81}I.
 \label{md:restop:normas}
\end{equation}
Le numérateur $32$ provient de $8|{-1}-1|^2$.
Dans les neuf composantes de la mémoire, la norme comprend
$8^2+8=72$, avec le dénominateur $27^2$. Le carré $64$ participe à
ce calcul avec les huit autres contributions.
Une identification à une région de $64$ cases exige l’application
d’incidence de la région et ne s’obtient pas en doublant le numérateur $32$.
La lecture géométrique d’un support $8\times8$ par rapport à un autre
$9\times9$ conserve donc son problème de connexion avec
l’application concrète ; l’égalité numérique isolée ne le remplace pas.

\clearpage
\subsection{Caractère d’action et coordonnée complète de réciprocité}
\label{md:restop:h5}

\subsubsection{Coefficients généalogiques et continuation régionale}
\label{md:restop:coeficientes}

La matrice centrale et ses invariants sont
\begin{equation}
 B_c=\begin{pmatrix}7&2\\2&7\end{pmatrix},\qquad
 \lambda_+=9,\qquad\lambda_-=5.
 \label{md:restop:matriz-central}
\end{equation}
Interviennent en outre l’orbite de longueur quatre du pas trois, le
calendrier de longueur douze et le quotient de recensement $104976/1944=54$.
Les relations entre coefficients sont
\begin{equation}
 \begin{aligned}
 \frac9{16}&=\frac{\lambda_+}{4^2},&
 \frac59&=\frac{\lambda_-}{\lambda_+},\\
 \frac7{48}&=\frac{\operatorname{tr}B_c/2}{4\cdot12},&
 \frac1{54}&=\frac{1944}{104976}.
 \end{aligned}
 \label{md:restop:coeficientes-h5}
\end{equation}
L’attribution des degrés et des signes appartient à la définition du caractère ;
leurs cardinaux isolés ne la déterminent pas. Le lecteur complet est
\begin{equation}
 H_5=\frac12\sqrt{\frac{1000\alpha}{\varphi}}-\alpha
 +\frac9{16}\alpha^2-\frac59\alpha^3
 +\frac7{48}\alpha^4-\frac1{54}\alpha^5,
 \label{md:restop:h5-completo}
\end{equation}
\begin{equation}
 \begin{aligned}
 D_A&=\exp(-100\pi\alpha/9),\\
 \mathcal C_\pi&=169\alpha^6+
 \frac{D_A\alpha^7}{1-D_A\alpha},\\
 \eta_{\mathrm{ret}}&=H_5-\frac{90}{\pi}\mathcal C_\pi.
 \end{aligned}
 \label{md:restop:retorno-regional}
\end{equation}
La fraction conserve la continuation infinie définie par
\begin{equation}
 \mathscr R(z)=D_Az^7+D_Az\mathscr R(z),
 \label{md:restop:renovacion}
\end{equation}
avec un poids initial unitaire. La continuation demeure intégrale dans le
lecteur ; un polynôme tronqué ne remplace pas cet objet. La construction
des coefficients et le renouvellement sont réunis dans la
section~\ref{sec:revision-planck}.

\subsubsection{Classification de la réciprocité radiale orientée}
\label{md:restop:geometria}

On conserve $A_{\deg}=1000\alpha$,
$C^*_{\deg}=2(\eta_{\mathrm{ret}}+\alpha)$ et le récupérateur radial du
théorème~\ref{rec:thm:accion}. Pour maintenir les deux orientations, soient
\begin{equation}
 \epsilon\in\{+1,-1\},\qquad
 q_{\mathrm{rad}}=R_\star^4>0,\qquad
 \Xi_{\mathrm{rad}}=
 \epsilon\frac{1-q_{\mathrm{rad}}}{1+q_{\mathrm{rad}}}.
 \label{md:restop:par-radial}
\end{equation}
Le récupérateur \eqref{rec:eq:eta-orientada} s’écrit
\begin{equation}
 \frac{\eta_{\mathrm{ret}}}{\alpha}=500\Xi_{\mathrm{rad}}-1.
 \label{md:restop:recuperador}
\end{equation}
L’involution
$(\epsilon,q_{\mathrm{rad}})
 \mapsto(-\epsilon,q_{\mathrm{rad}}^{-1})$
conserve cette action. La coordonnée qui réunit cette structure est
\begin{equation}
 \boxed{\chi_{\mathrm{rad}}=\epsilon\log q_{\mathrm{rad}}.}
 \label{md:restop:chi}
\end{equation}

\begin{proposition}
\label{md:restop:proposicion-orbitas}
$\chi_{\mathrm{rad}}$ classe exactement les orbites de cette involution
sur $\{\pm1\}\times\mathbb R_{>0}$.
\end{proposition}
\begin{proof}
Changer simultanément le signe et inverser $q_{\mathrm{rad}}$ conserve
$\chi_{\mathrm{rad}}$. Si deux paires ont la même valeur de
$\chi_{\mathrm{rad}}$, avec le même signe, elles ont le même
$q_{\mathrm{rad}}$ ; avec des signes opposés, leurs rayons quartiques sont
réciproques. Ce sont exactement les deux cas d’appartenance à une même
orbite. L’argument comprend $q_{\mathrm{rad}}=1$, dont l’orbite conserve
les deux orientations.
\end{proof}

L’identité
$(1-q)/(1+q)=-\tanh(\log q/2)$ donne
\begin{equation}
 \Xi_{\mathrm{rad}}
 =-\tanh(\chi_{\mathrm{rad}}/2),\qquad
 \boxed{\eta_{\mathrm{ret}}
 =\alpha[-500\tanh(\chi_{\mathrm{rad}}/2)-1].}
 \label{md:restop:accion-chi}
\end{equation}
La compatibilité avec le caractère complet prend la forme
\begin{equation}
 \boxed{\chi_{\mathrm{rad}}
 =-2\operatorname{artanh}
 \left(\frac{H_5+\alpha-(90/\pi)\mathcal C_\pi}{500\alpha}\right).}
 \label{md:restop:chi-caracter}
\end{equation}
Le terme $-\alpha$ de $H_5$ s’annule avec le $+\alpha$ du
contraste angulaire avant cette lecture. Le dénominateur de la
continuation régionale demeure complet. L’égalité reçoit les
sorties et les lecteurs HMT antérieurs et réunit leurs deux représentations ;
elle ne définit pas rétroactivement $\alpha$.

Dans la branche positive d’action et de rayon fini,
\begin{equation}
 0<\eta_{\mathrm{ret}}<499\alpha,\qquad
 \chi_{\mathrm{rad}}
 <-2\operatorname{artanh}(1/500).
 \label{md:restop:dominio-positivo}
\end{equation}
Avec $\alpha$ fija,
\begin{equation}
 \frac{\partial\eta_{\mathrm{ret}}}{\partial\chi_{\mathrm{rad}}}
 =-250\alpha\,\operatorname{sech}^2(\chi_{\mathrm{rad}}/2)<0.
 \label{md:restop:derivada-accion}
\end{equation}
Par conséquent, la section adimensionnelle d’action détermine de manière unique la
classe radiale orientée dans cette représentation, et réciproquement.
Pour des variations conjointes,
\begin{equation}
 d\eta_{\mathrm{ret}}
 =\frac{\eta_{\mathrm{ret}}}{\alpha}\,d\alpha
 -250\alpha\,\operatorname{sech}^2(\chi_{\mathrm{rad}}/2)
 \,d\chi_{\mathrm{rad}}.
 \label{md:restop:diferencial-accion}
\end{equation}

Le caractère $H_5$, la correction régionale et l’anisotropie radiale orientée
maintiennent ainsi une relation de compatibilité et ne constituent pas trois
paramètres indépendants. Le caractère analytique et la géométrie de
réciprocité représentent la même section de retour.
Ce niveau structurel de l’équation de Planck demeure explicite
avant de composer l’action avec une fréquence.

La décennie d’action conserve sa généalogie déterminantale, indépendante
de cette reparamétrisation :
\begin{equation}
 \hbar_{\mathrm{ret}}
 =10^{-34}\mathcal U_S\eta_{\mathrm{ret}},\qquad
 \hbar_{\mathrm{pre}}=10^{-34}\mathcal U_SH_5.
 \label{md:restop:secciones-accion}
\end{equation}
On conserve en outre l’action par tour $h=2\pi\hbar$ et la lecture
thermique
\begin{equation}
 k_B\Theta_{\mathrm{clk}}\log3=\frac{h}{T_{108}}
 \label{md:restop:lectura-termica}
\end{equation}
dans son domaine. L’horloge physique d’une réalisation conserve son constructeur :
la coordonnée $\chi_{\mathrm{rad}}$ classe cette réciprocité et non la
totalité des histoires TPK. La classification par
$\chi_{\mathrm{rad}}$ et la relation différentielle réunissent les représentations
existantes de l’action ; elles maintiennent leur caractère de formalisation focale et ne
constituent pas une seconde dérivation indépendante de Planck.

\subsection{Circuits de phase, mémoire et conservativité}
\label{md:restop:conservatividad}

L’holonomie nonadique restaure la phase et accroît la mémoire, selon
la section~\ref{mono-ampl:vueltas}. Si un parcours fermé dans les phases
augmente le compteur entier, cet accroissement ne peut pas s’écrire comme
différence d’un potentiel dépendant uniquement de la phase.

La démonstration s’obtient par télescopage : la somme
\begin{equation}
 \sum_j\bigl[F(g_{j+1})-F(g_j)\bigr]=0
 \label{md:restop:potencial-fase}
\end{equation}
sur un circuit de phases contraste avec l’accroissement de mémoire, égal
à un par tour. Sur l’état élargi, le compteur $m$ enregistre
cet accroissement. Le parcours fermé dans la projection conservait
donc une ouverture dans l’état complet.

L’étude de la conservativité commence par déclarer l’espace dans
lequel le parcours se ferme et les coordonnées qui participent à la
dynamique. Non-commutativité et non-conservativité sont des propriétés
différentes. La section~\ref{md:restop:memoria} contrôle l’ordre des
opérations ; cette section distingue un circuit projeté d’un circuit
complet. Pour étudier une force concrète, la fonctionnelle d’action ou de
travail doit agir sur le même espace d’états et le même transport.

La croissance de la mémoire n’attribue pas à elle seule une production
thermodynamique. La section~\ref{md:cron:crecimiento} démontre que
l’information des blocs peut croître avec un taux entropique nul.
La restitution opératorielle ajoute une précision : conserver l’information
peut signifier conserver les transformations qui peuvent être composées,
en plus de maintenir les symboles stockés.

\subsection{Compatibilité dynamique entre observations}
\label{md:restop:sintesis}

L’unité d’étude est la compatibilité dynamique entre observations,
la mémoire faisant partie de l’opération. La valeur d’une constante
constitue une lecture ; sa structure contient des contraintes sur d’autres
lectures et sur les transformations admissibles. Le prolongement de
l’analyse consiste à obtenir des relations qui empêchent de modifier de manière
indépendante la valeur, l’orientation, la mémoire, la forme et l’action.

Les trois développements focaux parcourent ce même axe :
\begin{itemize}
 \item Un cycle complet de régimes TRIT conserve encore un
 déplacement nonadique de capacité : la régularité locale n’épuise pas
 le retour.
 \item La réentrée de mémoire modifie exactement la réponse à l’ordre
 de deux opérations : l’historique intervient causalement dans la composition.
 \item Le caractère $H_5$ et la géométrie radiale orientée déterminent
 mutuellement leur section de retour : l’action admet une représentation
 structurelle récupérable.
\end{itemize}
Ces conséquences prolongent la généalogie de $\pi,\varphi,e,\alpha$ et
explicitent ce que signifie la conserver lors du passage aux opérations et à
l’action. La composition ultérieure réunit ces niveaux au moyen des applications
communes déjà construites, en maintenant la distinction entre leurs indices et
leurs mémoires même s’ils partagent une dénomination.

% Focal insertion; hierarchy inherited from the host.
\begingroup
% Fragmento maestro: incluir desde el anfitrion, con \HMTFGRoot fijada a la raiz del paquete.
% Requiere amsmath, amssymb/mathrsfs o unicode-math, amsthm, booktabs, longtable, hyperref, needspace.
% El anfitrion conserva sus entornos theorem, proposition, lemma y proof.
\providecommand{\HMTFGRoot}{.}
\begingroup
\providecommand{\ZZ}{}\renewcommand{\ZZ}{\mathbb Z}
\providecommand{\RR}{}\renewcommand{\RR}{\mathbb R}
\providecommand{\FF}{}\renewcommand{\FF}{\mathbb F}
\providecommand{\APP}{}\renewcommand{\APP}{\mathrm{APP}}
\providecommand{\TPK}{}\renewcommand{\TPK}{\mathrm{TPK}}
\providecommand{\TRIT}{}\renewcommand{\TRIT}{\{-1,0,+1\}}
% unicode-math usa lBrack/rBrack; newtxmath ya proporciona llbracket/rrbracket.
\providecommand{\llbracket}{\lBrack}
\providecommand{\rrbracket}{\rBrack}
% Fragmento de inserción; se incluye desde la raíz de este paquete editorial.
% Las fuentes completas y los cortes materiales constan en ANTECEDENTES_Y_DEPENDENCIAS.md.
% Procedencia: XI REV11 ES, sections/nucleo.tex, tpk_desarrollo_integrado.tex,
% ch_tres_vueltas_ext.tex, ch_elevacion_catalogal.tex, ch_cierre_cardinal.tex,
% ch_descenso_por_fibras.tex; perfil integral public_final/feigenbaum/c37_feigenbaum.tex;
% Grassmannianos REV07 ES, sections/jacobi_spectral.tex;
% TRANSVERSALIDAD_Y_ESCALADO_LOCAL_FEIGENBAUM.md, sección 14.
% Ampliación expositiva de demostraciones existentes; no introduce resultados nuevos.
\providecommand{\ZZ}{\mathbb Z}
\providecommand{\RR}{\mathbb R}
\providecommand{\FF}{\mathbb F}
\providecommand{\APP}{\mathrm{APP}}
\providecommand{\TPK}{\mathrm{TPK}}
\providecommand{\TRIT}{\{-1,0,+1\}}

\section{Antécédents operatoriels de la famille de criticité}
\label{hmtfg:antecedentes-operatorios}

La famille de criticité est construite à partir d’une lecture du calendrier TPK. Les deux feuilles arithmétiques, l’orientation tritique et le transport avec mémoire déterminent d’abord les états et leurs prolongements. La lecture ordonnée du continu permet ensuite d’évaluer des fonctions analytiques, de comparer des paramètres et de sélectionner des racines. Ce développement conserve l’état généalogique et ses deux lectures : la carte archimédienne fournit l’espace d’évaluation, tandis que la lecture d’incidence conserve les relations de frontière. La structure discrète conjointe du continu maintient ses cinq constructions consubstantielles. La présente exposition utilise son lecteur ordonné et conserve l’antécédent commun dont il provient.

\subsection{Cellules arithmétiques, orientation et transport}
\label{hmtfg:app-trit-transporte}

On utilise la cellule relevée \eqref{eq:app-levantada}, les coordonnées tritiques et leur loi arithmétique (\ref{trit-ampl-def-coordenadas}, \ref{trit-ampl-prop-aritmetica}), ainsi que le sélecteur~\ref{tpk-int:seleccion}. Fixons les notations : $D_9=\{1,\ldots,9\}$, $X_9=(\mathbb Z/9\mathbb Z)^2$ et, pour tout $n\in\mathbb Z$, $\rho_9(n)=1+((n-1)\bmod9)$ et $q_9(n)=(n-\rho_9(n))/9$. Le représentant équilibré est $r_3(n)\in\{-1,0,+1\}$ ; sa retenue est $\chi(a,b)=(a+b-r_3(a+b))/3$. Ces formules conservent l’extension aux entiers des évaluations positives du noyau.

Le transport relevé et sa réversibilité sont démontrés en \ref{tpk-int:prop-desplazamiento} ; la configuration observable et son retour de $108$ pas, en \ref{tpk-int:fobs} et \ref{tpk-int:periodo}. L’actualisation \eqref{eq:actualizacion} conserve l’état complet. L’holonomie de \ref{tpk-int:gamma} et les retours de \ref{hist:tres-retornos} distinguent restauration de phase et progression de la mémoire. La réalisation concrète des trois branches et son registre intégral sont construits ci-dessous à partir de ces opérations.

\input{inserciones_20260930/feigenbaum/00a_transiciones_nonadicas.tex}

\subsection{Raffinement des cylindres et lecture ordonnée}
\label{hmtfg:refinamiento-ordenado}

% XI REV11, ch_elevacion_catalogal.tex: bloques de 54 aristas, beta_blk,
% partición efectiva de 729 subcilindros; ch_cierre_cardinal.tex:375--497.
Pour \(\mathbf b=(b_1,\ldots,b_6)\in\FF_3^6\), soit
\(\upsilon:\FF_3\to\{-1,0,+1\}\) le représentant équilibré. Les trois prolongements construits arête par arête définissent
\[
 \Gamma^{[54]}_{\mathbf b,k}
 =\gamma_{\upsilon(b_6),k+45}\circ\cdots\circ
   \gamma_{\upsilon(b_1),k}.
\]
En partant de \(Y_0=\{x_\star\}\), on pose
\[
 Y_{N+1}
 =\{\Gamma^{[54]}_{\mathbf b,1+54N}(x):
          x\in Y_N,\ \mathbf b\in\FF_3^6\}.
\]
La troncature \(\rho^{\rm blk}_{N+1,N}\) supprime les cinquante-quatre dernières arêtes, en récupérant l’état antérieur au moyen du registre de transition.

\begin{proposition}[Codage effectif des blocs de prolongement]
\label{hmtfg:bloques-729}
La lecture des six couples tritiques de chaque bloc définit des bijections
\[
 \beta_N:Y_N\xrightarrow{\ \cong\ }(\FF_3^6)^N
\]
qui commutent avec la troncature. Chaque état de \(Y_N\) possède exactement \(729\) prolongements de bloc dans \(Y_{N+1}\).
\end{proposition}
\begin{proof}
À chaque tour, le registre conserve le couple tritique et sa retenue. Les trois couples produisent les trois sorties distinctes \(\chi(\varepsilon,\varepsilon)=\varepsilon\), et le registre permet de récupérer la branche utilisée, y compris celle de retenue nulle. Les six concaténations produisent donc tous les mots de \(\FF_3^6\). Deux mots différents se distinguent au premier couple enregistré où ils diffèrent. La récupération du préfixe prouve la compatibilité avec la troncature. Par récurrence sur \(N\), chaque mot de longueur \(N\) possède un unique état antécédent ; le mot suivant peut être n’importe laquelle des \(3^6=729\) possibilités.
\end{proof}

Soit \(\Omega_{\Gamma}^{\rm cal}=\varprojlim_NY_N\). Les bijections précédentes induisent
\[
 \Omega_{\Gamma}^{\rm cal}\cong(\FF_3^6)^{\mathbb N}.
\]
La correspondance et son inverse sont continues : chaque coordonnée finie dépend d’un préfixe fini. Les cylindres sont ouverts et fermés. Chaque cylindre de profondeur \(N\) se partitionne en ses \(729\) sous-cylindres de profondeur \(N+1\) ; chacun contient un prolongement infini, obtenu, par exemple, en poursuivant par la branche nulle dans tous les blocs ultérieurs.

% Antecedentes efectivos del continuo conjunto y del universo interno.
% Se imprimen en ambos destinos; la etiqueta de universo no reemplaza su regla.
\input{inserciones_20260930/feigenbaum/00b_continuo_universo.tex}

\subsection{Lecture ordonnée des cylindres de bloc}
\label{hmtfg:publicacion-cilindros-bloque}

Pour le représentant \([b]_3\in\{0,1,2\}\), on fixe
\[
 \nu(\mathbf b)=\sum_{i=1}^6[b_i]_3\,3^{6-i}\in\{0,\ldots,728\}.
\]
Dans la complétion ordonnée HMT, un préfixe
\(s=(m;\mathbf b_0,\ldots,\mathbf b_{N-1})\), \(m\in\ZZ\), détermine
\[
 a_N(s)=m+\sum_{j=0}^{N-1}\nu(\mathbf b_j)\,729^{-j-1},\qquad
 I_N(s)=[a_N(s),a_N(s)+729^{-N})_{\rm HMT}.
\]
La retenue conserve le raccord entre les deux représentants d’une frontière. Si \(C_N(s)=\overline{I_N(s)}\), les relations précédentes donnent
\[
 I_N(s)=\bigsqcup_{\mathbf b\in\FF_3^6}I_{N+1}(s;\mathbf b),
 \qquad \operatorname{diam}C_N(s)=729^{-N}\longrightarrow0.
\]

\begin{proposition}[Couverture de la lecture archimédienne]
\label{hmtfg:cobertura-arquimediana}
L'application
\[
 P_{\rm ar}:\ZZ\times\Omega_\Gamma^{\rm cal}
       \longrightarrow\mathbb R_{\rm HMT},\qquad
 \{P_{\rm ar}(m,\xi)\}
    =\bigcap_N C_N(m;\beta_N(\rho_{\infty,N}\xi))
\]
est bien définie, est surjective et détermine
\[
 (\ZZ\times\Omega_\Gamma^{\rm cal})/\!\sim_{P_{\rm ar}}
       \ \cong\ \mathbb R_{\rm HMT}.
\]
\end{proposition}
\begin{proof}
Les extrémités gauches forment une suite de Cauchy dans la complétion HMT, puisque les intervalles sont emboîtés et que leurs diamètres tendent vers zéro. Leur limite appartient à tous les intervalles fermés et est unique. Pour tout élément de la complétion, on choisit d’abord l’intervalle entier qui le contient, puis l’unique sous-intervalle semi-ouvert de chaque partition qui le contient. La proposition précédente réalise chacun de ces choix par un prolongement TPK effectif. Les préfixes compatibles définissent \(\xi\), dont la lecture est l’élément choisi. Aux frontières, la retenue identifie les écritures équivalentes après leur construction. Le quotient par égalité de lecture est, par définition et surjectivité, en bijection avec la complétion.
\end{proof}

La construction est également effectuée dans l’univers généalogique interne \(\mathfrak G=\mathfrak G_{\rm HMT}(h,m_\infty)\). Ici, \(h\) code les règles operatorielles et \(m_\infty\) l’histoire compatible ; l’univers s’obtient par clôture transitive initiale, définissabilité à chaque successeur et réunion aux stades limites :
\[
 G_0=\operatorname{TC}\{\omega,u_{h,m},h,m_\infty\},\qquad
 G_{\beta+1}=\operatorname{Def}_{\Sigma_{\rm HMT}}(G_\beta),\qquad
 G_\lambda=\bigcup_{\beta<\lambda}G_\beta .
\]
Les relations de \(\Sigma_{\rm HMT}\) sont celles des opérations déjà construites et de leurs codes. Le domaine interne est donc maintenu dans toutes les affirmations de couverture. Notons \(\mathbb R_{\rm HMT}^{\mathfrak G}\) sa complétion et \(\mathbb R^{\mathfrak G}\) la réalisation ordonnée au sein du même univers.

\begin{proposition}[Carte ordonnée interne]
\label{hmtfg:carta-interna}
Il existe un unique isomorphisme de corps ordonnés archimédiens complets
\[
 \iota:\mathbb R_{\rm HMT}^{\mathfrak G}
             \xrightarrow{\ \cong\ }\mathbb R^{\mathfrak G}
\]
qui fixe l’unité. Il conserve les extrémités rationnelles, les racines carrées positives et les limites de suites convergentes.
\end{proposition}
\begin{proof}
L’unité fixe les entiers et les rationnels. À chaque élément de la complétion HMT est associée la coupure des rationnels qui lui sont inférieurs ; la complétude de la réalisation ordonnée fournit l’unique élément ayant cette coupure. La densité rationnelle prouve l’injectivité et la conservation stricte de l’ordre. Le même procédé en sens inverse prouve la surjectivité. Les approximations rationnelles par excès et par défaut conservent la somme et le produit, et donc la structure de corps. L’élément positif dont le carré est \(a>0\) est transporté vers l’élément positif dont le carré est \(\iota(a)\). Si une suite converge, pour chaque tolérance rationnelle positive ses derniers termes sont dans le voisinage rationnel de la limite ; la conservation de l’ordre transporte cette condition. Toutes ces opérations s’effectuent dans \(\mathfrak G\). L’unicité découle de la conservation des coupures rationnelles.
\end{proof}

\subsection{Descente par les fibres et conservation du sélecteur}
\label{hmtfg:descenso-selector}

% XI REV11, ch_descenso_por_fibras.tex:10--39; prueba completa pertinente.
\begin{proposition}[Critère de descente d’une opération]
\label{hmtfg:criterio-descenso}
Soient \(q:\Omega\twoheadrightarrow J\) et
\(\star:D\subseteq\Omega^2\to\Omega\), avec \(D\) saturé par les fibres de \(q\times q\). Il existe une unique opération \(\bar\star\) sur \((q\times q)(D)\) telle que
\[
 q(x\star y)=q(x)\,\bar\star\,q(y)
\]
si et seulement si
\[
 q(x)=q(x'),\quad q(y)=q(y')
 \ \Longrightarrow\ q(x\star y)=q(x'\star y')
\]
pour les couples du domaine.
\end{proposition}
\begin{proof}
Si \(\bar\star\) existe, les deux résultats sont l’évaluation du même couple d’arguments. Réciproquement, on choisit des antécédents des arguments visibles et on définit le résultat comme la lecture de leur opération. La saturation fixe le domaine ; la condition sur les fibres rend le résultat indépendant des antécédents choisis. La surjectivité prouve l’unicité.
\end{proof}

Dans une application au moyen de \(P_{\rm ar}\), ce critère identifie la condition précise permettant de transporter une opération depuis les états complets jusqu’à leur lecture. Dans la carte \(\iota\), la structure de corps et les limites satisfont déjà cette condition. Le profil analytique est ensuite construit à partir du calendrier, et cette compatibilité est utilisée pour transporter son itération et ses racines.

\subsection{Lecture dodécaphasique et normalisation quadratique}
\label{hmtfg:perfil-dodecafasico}

% Propietario integral c37_feigenbaum.tex: comienzo hasta la familia f_lambda.
La section orientée de criticité sur les douze fenêtres du calendrier est
\[
 \mathcal P_{\rm crit}(m)=\varphi_m
   =(30m-10)\frac{\pi_{\rm HMT}}{180}
   =\pi_{\rm HMT}\frac{3m-1}{18},\qquad m=1,\ldots,12.
\]
Cette section fixe l’origine et l’orientation de ses représentants angulaires. Le lecteur uniforme des douze fenêtres est
\[
 \eta_0(F)=\frac1{12}\sum_{m=1}^{12}F(\varphi_m).
\]
On distingue la fonction \(\eta_0\), qui attribue des poids à ces fenêtres, de la carte ordonnée \(\iota\). La normalisation uniforme et la section angulaire font partie des données operatorielles de ce profil.

En posant \(a_m=(3m-1)/18\), on obtient
\[
 D_0=\frac1{12}\sum_{m=1}^{12}a_m^2=\frac{899}{648},
 \qquad
 s_0=\frac{36}{\pi_{\rm HMT}\sqrt{899}},\qquad
 k_m=s_0\varphi_m=\frac{2(3m-1)}{\sqrt{899}}.
\]
En effet, \(\sum_{m=1}^{12}(3m-1)^2=5394=6\cdot899\), de sorte que \(s_0^2\pi_{\rm HMT}^2D_0=2\). Le profil est
\begin{equation}
 u_0(x)=\frac1{12}\sum_{m=1}^{12}\bigl[1-\cos(k_mx)\bigr],
 \qquad f_\lambda(x)=1-\lambda u_0(x).
 \label{hmtfg:perfil-familia}
\end{equation}
La constante \(\pi_{\rm HMT}\) intervient comme lecture angulaire du noyau HMT et s’annule algébriquement dans la normalisation du profil. L’échelle \(\lambda\) est le paramètre de la famille, dont la valeur critique sera sélectionnée par la dynamique.

\begin{proposition}[Normalisation et forme fermée du profil]
\label{hmtfg:forma-perfil}
Le profil \(u_0\) est pair et analytique, satisfait
\(u_0(0)=u_0'(0)=0\) et \(u_0''(0)=2\), et, pour
\(y=x/\sqrt{899}\), admet l’expression
\[
 u_0(x)=1-\frac{\cos(37y)\sin(36y)}{12\sin(3y)}
       =1-\frac{\sin(73y)-\sin(y)}{24\sin(3y)}.
\]
Les singularités apparentes s’interprètent par continuité. Son développement initial est
\[
 u_0(x)=x^2-\frac{715769}{2424603}x^4+O(x^6).
\]
\end{proposition}
\begin{proof}
La somme finie de cosinus conserve la parité et l’analyticité. Sa dérivée seconde en zéro est \(12^{-1}\sum_m k_m^2=2\). La somme de la progression trigonométrique \(\sum_{m=1}^{12}\cos((6m-2)y)\) est \(\cos(37y)\sin(36y)/\sin(3y)\) ; l’identité produit--somme fournit la seconde expression. La série du cosinus et \(\sum_m(3m-1)^4=4294614\) donnent le coefficient de degré quatre indiqué. L’expression comme somme finie fixe le prolongement aux zéros du dénominateur de la forme fermée.
\end{proof}

\subsection{Transport des itérées, des racines et des minima}
\label{hmtfg:transporte-raices}

% Reunión existente: TRANSVERSALIDAD_Y_ESCALADO_LOCAL_FEIGENBAUM.md, sección 14.1.
Dans \(\mathbb R_{\rm HMT}^{\mathfrak G}\), la série convergente
\[
 \cos_{\rm HMT}(t)=\sum_{j=0}^{\infty}
                   \frac{(-1)^j t^{2j}}{(2j)!}
\]
définit la réalisation analytique correspondante. Soient
\[
 u_{\rm HMT}(x)=\frac1{12}\sum_{m=1}^{12}
  \left[1-\cos_{\rm HMT}
   \left(\frac{2(3m-1)x}{\sqrt{899}_{\rm HMT}}\right)\right],
 \qquad F_a(x)=1-a\,u_{\rm HMT}(x),
\]
\[
 S_n^{\rm HMT}(a)=F_a^{\,2^n}(0),\qquad
 S_n(\lambda)=f_\lambda^{\,2^n}(0).
\]

\begin{proposition}[Conjugaison ordonnée de la famille et du sélecteur de racines]
\label{hmtfg:conjugacion-selector}
Pour les arguments internes des expressions précédentes,
\[
 \iota\circ F_a=f_{\iota(a)}\circ\iota,\qquad
 \iota(S_n^{\rm HMT}(a))=S_n(\iota(a)).
\]
La carte transporte exactement les conditions de retour, de période exacte et d’ordre des paramètres. Si un ensemble de racines admissibles possède un minimum, l’image de ce minimum est le minimum de l’ensemble image.
\end{proposition}
\begin{proof}
La carte fixe les rationnels, transporte la racine carrée positive et conserve les limites des sommes partielles du cosinus. Ainsi, \(\iota(u_{\rm HMT}(x))=u_0(\iota(x))\). La compatibilité avec la somme et le produit prouve la première égalité ; la récurrence sur le nombre d’itérations prouve la seconde. Une racine se caractérise par l’égalité à zéro, conservée dans les deux sens. La période exacte \(2^n\) ajoute les inégalités \(F_a^{2^j}(0)\ne0\), \(0\le j<n\), également conservées par l’injectivité de \(\iota\). Si l’on exige que le paramètre dépasse un centre précédent, cette inégalité est transportée par conservation de l’ordre. Enfin, \(a\le b\) pour tout \(b\) de l’ensemble équivaut à \(\iota(a)\le\iota(b)\) pour tout élément de son image. La surjectivité sur les ensembles ainsi définis complète l’affirmation relative au minimum.
\end{proof}

Le domaine de ce transport est le continu interne indiqué. L’existence de chaque minimum et son identification à une branche dynamique précise sont des résultats de la construction de criticité qui suit ; la carte conserve ces propriétés lorsqu’elles ont été établies.

\subsection{Représentation spectrale finie du lecteur uniforme}
\label{hmtfg:jacobi-finito}

% Grassmannianos REV07, sections/jacobi_spectral.tex:35--136,
% mecanismo de Hankel, Gram--Schmidt, recurrencia y representación espectral.
% Especialización atómica ya reunida en la sección 14.2 del desarrollo Feigenbaum.
Les fréquences au carré du profil définissent la mesure positive
\[
 s_m=k_m^2=\frac{4(3m-1)^2}{899},\qquad
 \nu_0=\frac1{12}\sum_{m=1}^{12}\delta_{s_m},\qquad
 M_r=\int s^r\,d\nu_0(s).
\]
Ses douze nœuds sont distincts et positifs. Pour \(1\le r\le12\), la matrice de Hankel \(H_r=(M_{i+j})_{0\le i,j<r}\) est définie positive. En effet, pour \(p\) de degré inférieur à \(r\),
\[
 \mathbf p^{\mathsf T}H_r\mathbf p
   =\frac1{12}\sum_{m=1}^{12}p(s_m)^2>0
\]
sauf si \(p=0\), car un polynôme de degré au plus onze qui s’annule en douze nœuds distincts est nul. Au degré douze apparaît le polynôme \(\prod_m(s-s_m)\), de norme nulle pour cette mesure. La représentation spectrale correspondante est de dimension douze.

Soient \(p_0,\ldots,p_{11}\) les polynômes orthogonaux moniques,
\(h_j=\int p_j^2\,d\nu_0>0\) et
\(\psi_j=p_j/\sqrt{h_j}\). Comme \(\nu_0\) est une probabilité,
\(\psi_0=1\). La multiplication par \(s\) dans \(L^2(\nu_0)\) est autoadjointe et possède une matrice tridiagonale \(J_{12}\) dans cette base.

\begin{proposition}[Récurrence et entrelacement spectral exacts]
\label{hmtfg:jacobi-entrelazamiento}
Avec
\[
 a_j=\frac{\int s p_j(s)^2\,d\nu_0(s)}{h_j},\qquad
 \beta_j=\frac{h_j}{h_{j-1}}>0\quad(1\le j\le11),
\]
la matrice \(J_{12}\) a pour diagonale \(a_j\) et pour sous-diagonale
\(\sqrt{\beta_j}\). Si
\[
 Q_{mj}=\frac{\psi_j(s_m)}{\sqrt{12}}
 \quad(1\le m\le12,\ 0\le j\le11),\qquad
 D=\operatorname{diag}(s_1,\ldots,s_{12}),
\]
alors
\[
 Q^{\mathsf T}Q=I,\qquad DQ=QJ_{12},\qquad
 J_{12}=Q^{\mathsf T}DQ,
\]
et le profil satisfait
\begin{equation}
 u_0(x)=e_0^{\mathsf T}
          \bigl[I-\cos(x\sqrt{J_{12}})\bigr]e_0.
 \label{hmtfg:perfil-espectral}
\end{equation}
\end{proposition}
\begin{proof}
La positivité des systèmes de Hankel fournit de manière unique les polynômes orthogonaux moniques. En développant \(s p_j\) dans la base orthogonale, les coefficients de \(p_i\), \(i\le j-2\), s’annulent car
\(\langle sp_j,p_i\rangle=\langle p_j,sp_i\rangle=0\).
Le coefficient de \(p_{j+1}\) est un par monicité, celui de \(p_j\) est \(a_j\) et celui de \(p_{j-1}\) est \(h_j/h_{j-1}\). Au dernier degré, la même égalité s’entend dans \(L^2(\nu_0)\), où le polynôme monique de degré douze qui s’annule aux nœuds représente le vecteur nul. La normalisation donne la sous-diagonale positive indiquée.

L’orthonormalité équivaut à \(Q^{\mathsf T}Q=I\). L’égalité de récurrence à chaque nœud donne \(DQ=QJ_{12}\). Comme \(Q\) est carrée, on a également \(QQ^{\mathsf T}=I\), et l’on obtient la diagonalisation de \(J_{12}\). Ses valeurs propres positives permettent de définir sa racine carrée positive. Pour toute fonction définie en ces douze nœuds,
\[
 e_0^{\mathsf T}\Phi(J_{12})e_0
   =\frac1{12}\sum_{m=1}^{12}\Phi(s_m),
\]
car \(Qe_0=(1,\ldots,1)^{\mathsf T}/\sqrt{12}\).
Le choix \(\Phi(s)=1-\cos(x\sqrt{s})\) prouve
\eqref{hmtfg:perfil-espectral}.
\end{proof}

La représentation conserve exactement la fonctionnelle uniforme, ses moments et son profil. Ses premiers coefficients sont
\[
 a_0=M_1=2,\qquad
 \beta_1=M_2-M_1^2=\frac{2493348}{808201}.
\]
La généalogie complète demeure dans l’état TPK dont provient le lecteur ; \(J_{12}\) représente sa lecture spectrale particulière. La construction de Gram--Schmidt et la récurrence de Jacobi s’appliquent à cette mesure atomique, avec leur terminaison en dimension douze.

La chaîne operatorielle utilisée dans les sections suivantes est ainsi fixée : les cellules APP et l’orientation TRIT alimentent la transition TPK ; ses prolongements produisent les cylindres et leur lecture ordonnée ; les douze fenêtres orientées et la fonctionnelle uniforme produisent \(u_0\) ; la famille \(f_\lambda\) détermine les retours et les racines ; la carte interne conserve l’ordre de leurs sélecteurs. La réalisation spectrale décrit le même profil au moyen d’une matrice autoadjointe finie et d’un vecteur cyclique.

% Fragmento integrable. Preambulo anfitrion: amsmath, amssymb, longtable, hyperref.
\section{Profil dodécaphasique et retour quadratique}
\label{hmtfg:fragmento:perfil-dodecafasico-y-retorno-cuadratico}

\subsection{Généalogie et objet de la démonstration}
\label{hmtfg:desarrollo:1}

L’objet initial est la lecture critique de la roue dodécaphasique orientée déjà construite dans le corpus. APP évalue la somme et le produit sur les deux feuilles du réseau de chiffres, en conservant quotient et reste ; TRIT détermine le régime et l’orientation ; le transport TPK sélectionne, transporte et actualise l’état avec ses registres de retenue, de frontière et de mémoire. La connexion nonadique prolonge cet état, avec retour de phase et avancement de mémoire. Ses lecteurs agissent sur la même structure discrète du continu.

Le lecteur focal de la source est
\[
\mathcal P_{\mathrm{crit}}:U_{12}^{\mathrm{sign}}\longrightarrow
\mathbb R/2\pi\mathbb Z,\qquad
m\longmapsto\phi_m=\frac{(3m-1)\pi}{18},\quad 1\le m\le12.
\]
Pour évaluer le moment angulaire, on conserve le représentant orienté de cette carte : le remplacer par un autre représentant modulo \(2\pi\) modifierait le moment. La trace uniforme fixe \(w_m=1/12\). La chaîne focale est
\[\begin{aligned}
\text{APP–TRIT–TPK}&\longrightarrow
\text{structure discrète du continu}\\ &\longrightarrow U_{12}^{\mathrm{sign}}
\xrightarrow{\mathcal P_{\mathrm{crit}}}(\phi_m,w_m)
\\ &\longrightarrow u_0\longrightarrow f_\lambda
\longrightarrow\text{retours et renormalisation}.
\end{aligned}\]

L’écriture abrégée conserve comme antécédent l’état complet et ses prolongements ; la famille scalaire est une lecture de cet antécédent. Le lecteur déjà fixé, avec son ordre, son orientation et son secteur uniforme, est pris comme structure de départ explicite. Les conclusions portent sur ce lecteur précis ; l’unicité de ce lecteur parmi toutes les lectures possibles du TPK reste hors de ces preuves.

La coordonnée \(\pi_{\mathrm{HMT}}\) utilisée par la carte appartient aux sorties internes du corpus. Son facteur s’annule exactement lors de la normalisation du deuxième moment. Les valeurs de Feigenbaum interviennent exclusivement dans la reconnaissance ultérieure, jamais dans la définition de la famille ni dans la certification des racines.


\subsection{Profil exact et criticité quadratique}
\label{hmtfg:desarrollo:2}

Soit \(s_m=3m-1\). La somme entière
\[
\sum_{m=1}^{12}s_m^2=5394=6\cdot899
\]
détermine
\[
\kappa_0=\frac{36}{\pi\sqrt{899}},\qquad
k_m=\kappa_0\phi_m=\frac{2s_m}{\sqrt{899}},
\]
\[
u_0(x)=\frac1{12}\sum_{m=1}^{12}(1-\cos(k_mx)),\qquad
f_{\lambda,0}(x)=1-\lambda u_0(x).
\]
Par l’identité du deuxième moment,
\[
u_0(0)=u_0'(0)=0,\quad u_0''(0)=2,\quad
f_{\lambda,0}''(0)=-2\lambda.
\]
La représentation analytique conserve des coefficients rationnels :
\[
u_0(z)=\sum_{j\ge1}(-1)^{j+1}
\frac{4^j\sum_m s_m^{2j}}{12\,899^j(2j)!}z^{2j}
=z^2-\frac{715769}{2424603}z^4
+\frac{1353832978}{32695771455}z^6+\cdots.
\]
Ainsi, le registre fini de douze secteurs produit un profil entier, dont la criticité quadratique et l’échelle sont fixées avant l’itération.


\subsection{Robustesse structurelle avec bornes uniformes}
\label{hmtfg:desarrollo:3}

La déformation déjà présente dans le code d’origine est
\[
w_m(\eta)=\frac{1+\eta p_m}{12},\qquad
p_m=12\cos(3\phi_m)-\cos(4\phi_m).
\]
On définit
\[
M_{2,\eta}=\sum_mw_m(\eta)\phi_m^2,\quad
\kappa_\eta=\sqrt{2/M_{2,\eta}},\quad k_{m,\eta}=\kappa_\eta\phi_m,
\]
\[
u_\eta(x)=\sum_mw_m(\eta)(1-\cos(k_{m,\eta}x)),\qquad
f_{\lambda,\eta}=1-\lambda u_\eta.
\]

\textbf{Théorème de robustesse structurelle.} Pour \(|\eta|\le1/40\), les poids sont positifs et leur somme vaut un. Le profil est pair, entier, satisfait \(u_\eta''(0)=2\) et est strictement croissant sur \((0,1]\). Pour \(0<\lambda\le2/u_\eta(1)\), l’application envoie \([-1,1]\) dans lui-même, possède un unique point critique dans cet intervalle et satisfait
\[
S f_{\lambda,\eta}(x)\le-\frac{640}{47647}<0,\qquad0<|x|\le1.
\]

\textbf{Démonstration.} Les sommes des cosinus \(3\phi_m\) et \(4\phi_m\) s’annulent par cycles de longueurs quatre et trois. Comme \(|p_m|\le13\),
\[
\frac9{160}\le w_m\le\frac{53}{480},\qquad
\frac{27}{40}M_{2,0}\le M_{2,\eta}\le\frac{53}{40}M_{2,0}.
\]
Il en découle que
\[
k_{\max,\eta}^2\le\frac{196000}{24273}<9<\pi^2,\qquad
k_{\min,\eta}^2\ge\frac{640}{47647}.
\]
Pour \(0<x\le1\), tous les sinus \(\sin(k_{m,\eta}x)\) sont positifs, d’où
\[
u_\eta'(x)=\sum_mw_mk_m\sin(k_mx)>0,\qquad
u_\eta'''(x)=-\sum_mw_mk_m^3\sin(k_mx)<0.
\]
Le quotient \(u_\eta'''/u_\eta'\) est l’opposé d’une moyenne positive des \(k_m^2\). L’invariance de la schwarzienne sous les transformations affines de la valeur donne
\[
S f_{\lambda,\eta}
=\frac{u_\eta'''}{u_\eta'}-\frac32
\left(\frac{u_\eta''}{u_\eta'}\right)^2
\le-k_{\min,\eta}^2.
\]
La parité couvre \(x<0\), et l’image de l’intervalle est
\([1-\lambda u_\eta(1),1]\). \(\square\)

La largeur \(1/40\) est un choix permettant d’obtenir des bornes uniformes, distinct d’une constante fondamentale ou d’un seuil optimal. Le secteur original reste \(\eta=0\).


\subsection{Coordonnée holomorphe quadratique globale sur une bande}
\label{hmtfg:desarrollo:4}

\textbf{Théorème de factorisation holomorphe.} Pour \(|\eta|\le1/40\), dans
\[
\mathcal S=\{z\in\mathbb C:|\operatorname{Re}z|<\pi/3\}
\]
il existe une fonction holomorphe, impaire et univalente \(b_\eta\), avec \(b_\eta'(0)=1\), telle que
\[
\boxed{u_\eta(z)=b_\eta(z)^2,\qquad
f_{\lambda,\eta}(z)=1-\lambda b_\eta(z)^2.}
\]

\textbf{Démonstration.} Pour \(z=x+iy\) dans la demi-bande droite,
\[
\operatorname{Re}u_\eta'(z)
=\sum_mw_mk_m\sin(k_mx)\cosh(k_my)>0.
\]
L’intégrale de \(u_\eta'\) sur le segment entre deux points, divisée par la différence des extrémités, a une partie réelle positive. La convexité de la demi-bande y prouve l’injectivité. La parité donne le résultat correspondant à gauche.

De plus,
\[
\operatorname{Im}u_\eta(x+iy)
=\sum_mw_m\sin(k_mx)\sinh(k_my)
\]
a le signe de \(y\) pour \(x>0\). Sur l’axe réel positif, \(u_\eta>0\) ; sur l’axe imaginaire,
\[
u_\eta(iy)=\sum_mw_m(1-\cosh(k_my))
\]
est strictement négatif hors de zéro et décroît strictement avec \(|y|\). Ces séparations, l’injectivité de chaque demi-bande et la parité prouvent que les fibres de \(u_\eta\) dans \(\mathcal S\) sont exactement \(\{z,-z\}\), sauf la fibre de l’origine.

La fonction \(u_\eta(z)/z^2\), prolongée par la valeur un en zéro, est holomorphe et ne possède aucun zéro dans la bande simplement connexe. Elle possède un logarithme holomorphe \(L_\eta\) avec \(L_\eta(0)=0\), qui est pair par unicité de la normalisation. Définissons
\[
b_\eta(z)=z\exp\!\left(\frac12L_\eta(z)\right).
\]
Alors \(b_\eta^2=u_\eta\), \(b_\eta\) est impaire et \(b_\eta'(0)=1\). L’égalité \(b_\eta(z)=b_\eta(w)\) impose \(w=z\) ou \(w=-z\). Le second cas ne donne l’égalité qu’à l’origine. Cela prouve son univalence. \(\square\)

Les douze phases déterminent ainsi une déformation conforme explicite du pli quadratique. La conjugaison dynamique complète conserve cette déformation :
\[
b_\eta\circ f_{\lambda,\eta}\circ b_\eta^{-1}(w)
=b_\eta(1-\lambda w^2),
\]
sur leurs domaines de composition. La factorisation démontrée et une conjugaison à toute la famille quadratique sont des affirmations distinctes.


\subsection{Retour, changement d’échelle et mémoire}
\label{hmtfg:desarrollo:5}

Soient \(a=-f(1)>0\), \(H_a(x)=-ax\), \(J=H_a([-1,1])\). Lorsque \(J\subset[-1,1]\) est admissible pour le retour, \(f^2(J)\subset J\) et que les compositions sont définies,
\[
\mathcal Rf=H_a^{-1}f^2H_a,\qquad
\mathcal Rf(x)=-a^{-1}f(f(-ax)).
\]
Si, de plus, \(f(J)\subset(0,1]\), un unique maximum quadratique central est conservé :
\[
(\mathcal Rf)(0)=1,\qquad
(\mathcal Rf)''(0)=-a f'(1)f''(0)<0.
\]
La composition des schwarziennes donne, aux points réguliers,
\[
S(\mathcal Rf)(x)=a^2
\left[(Sf)(f(-ax))f'(-ax)^2+(Sf)(-ax)\right]<0.
\]
Les inclusions de domaines font partie des hypothèses et sont vérifiées à chaque échelle.

Soit \(v_\eta=(u_\eta|_{[0,1]})^{-1}\). Les branches inverses sont
\[
\beta_\sigma(y)=\sigma v_\eta((1-y)/\lambda),\quad
\sigma\in\{-1,+1\},
\]
avec le registre critique \(\beta_0(1)=0\). Le pas
\[
(x,w)\longmapsto(f(x),w\mathbin{\|}\sigma(x))
\]
s’inverse sur son image en lisant la dernière feuille et en appliquant \(\beta_\sigma\). Le mot conserve les décisions de branche ; la valeur antérieure se reconstruit par l’équation de l’application.

Le retour regroupe deux transitions. À partir de l’extrémité \(x'\) et de ses feuilles \((\sigma_0,\sigma_1)\), on récupère
\[
x=H_a^{-1}\beta_{\sigma_0}
\!\left(\beta_{\sigma_1}(H_ax')\right).
\]
Les points intermédiaires sont également reconstruits. Les registres APP–TRIT–TPK de feuille, de retenue, de route, de frontière et de mémoire sont concaténés en conservant leur identité. Le signe de la branche analytique est un registre supplémentaire : il ne remplace pas ces données par une étiquette binaire.

Si \(C_N\) regroupe les histoires de longueur \(2N\) par couples, les restrictions satisfont
\[
\pi^{\mathrm{ret}}_{N,M}C_N
=C_M\pi^f_{2N,2M},\qquad M\le N.
\]
Les deux membres suppriment les mêmes couples terminaux et conservent les mêmes branches inverses. Le transport compatible des histoires finies est ainsi démontré.


\subsection{Échelles accumulées et dérivée de l’opérateur}
\label{hmtfg:desarrollo:6}

Soient \(c_n=f^{2^n}(0)\), \(c_0=1\) et \(H_n(x)=c_nx\). Tant que les compositions sont admissibles et que \(c_n\ne0\),
\[
\boxed{\mathcal R^nf=H_n^{-1}f^{2^n}H_n},\qquad
(\mathcal R^nf)(1)=\frac{c_{n+1}}{c_n}=-a_n.
\]
La récurrence découle de \(H_nH_{a_n}=H_{n+1}\). La convention \(a_n>0\) exige l’alternance du signe de \(c_n\). À une racine superstable de période \(2^N\), \(c_N=0\) : cette normalisation devient singulière et il faut utiliser les niveaux antérieurs ou une limite contrôlée.

Pour \(z=-ax\), \(w=f(z)\) et une perturbation \(h\), la dérivée complète est
\[
\boxed{
D\mathcal R_f[h](x)=
-\frac{h(w)+f'(w)h(z)}a+
\frac{h(1)}a
\left[\mathcal Rf(x)-xf'(w)f'(z)\right].
}
\]
On dérive en utilisant \(\dot a=-h(1)\), \(\dot z=xh(1)\),
\(\dot w=h(z)+xf'(z)h(1)\) et la variation du dénominateur. Si \(h(0)=0\), on obtient \(D\mathcal R_f[h](0)=0\). La direction paramétrique initiale HMT est \(h=-u_\eta\).

Cette formule inclut la variation d’échelle ; figer \(a\) donnerait un autre opérateur linéarisé.


\subsection{Racines et itinéraires certifiés jusqu’à la période 256}
\label{hmtfg:desarrollo:7}

Soit \(F_n(\lambda)=f_{\lambda,0}^{2^n}(0)\). On calcule
\[
x_0=p_0=0,\qquad x_{j+1}=1-\lambda u_0(x_j),\qquad
p_{j+1}=-u_0(x_j)-\lambda u_0'(x_j)p_j.
\]
Alors \(F_n=x_{2^n}\) et \(F_n'=p_{2^n}\).

Le programme joint utilise des intervalles rationnels avec arrondi entier vers l’extérieur, une racine carrée encadrée par une racine entière et des sinus/cosinus avec reste de Taylor explicite. La borne exacte \(|u_0''|\le2\) permet une évaluation centrée avec reste quadratique.

Les approximations historiques proposent seulement des boîtes de rayon \(10^{-42}\). On certifie des signes opposés à leurs extrémités, \(F_n'\) séparé de zéro dans toute la boîte et l’exclusion des retours de périodes qui divisent proprement \(2^n\). Le théorème des valeurs intermédiaires et la monotonie prouvent l’existence et l’unicité locales ; l’exclusion des diviseurs prouve la période minimale.

\begingroup\small
\begin{longtable}{p{.12\linewidth}p{.17\linewidth}p{.62\linewidth}}
\hline
Niveau & Période minimale & Centre de la boîte, arrondi \\
\hline
1 & 2 & 1.345913990471832792210949 \\
2 & 4 & 1.763379787846910127290794 \\
3 & 8 & 1.850413796950136464530567 \\
4 & 16 & 1.868979756606798125150574 \\
5 & 32 & 1.872955034435659740004205 \\
6 & 64 & 1.873806367467030119412262 \\
7 & 128 & 1.873988695221365362307169 \\
8 & 256 & 1.874027744154450672897008 \\
\hline\end{longtable}
\endgroup

La première racine possède l’expression exacte \(\lambda_1=1/u_0(1)\). Les boîtes complètes et leurs vérifications figurent dans le certificat JSON.

Les 502 états de préclôture de ces huit niveaux, tous séparés de zéro, ainsi que leurs mots de signes complets, ont été encadrés. En particulier,
\[
\operatorname{sign}f_{\lambda_n,0}^{2^j}(0)=(-1)^j,\qquad0\le j<n.
\]
Les orientations finies des changements d’échelle antérieurs à la clôture sont établies. Les inclusions d’intervalles restrictifs à toutes les échelles constituent une vérification supplémentaire.

Pour \(\delta_{\mathrm F,n}=(\lambda_{n-1}-\lambda_{n-2})/(\lambda_n-\lambda_{n-1})\), un intervalle décimal extérieur abrégé est
\[
\boxed{\begin{gathered}
4.66921218915020415455609621588966392804<\delta_{\mathrm F,8},\\
\delta_{\mathrm F,8}<4.66921218915020415455609621588966392863.
\end{gathered}}
\]
La largeur de l’intervalle complet est inférieure à \(5.808\cdot10^{-37}\). Elle mesure l’erreur numérique du rapport fini \(\delta_{\mathrm F,8}\), distincte de l’erreur d’approximation de la limite universelle.


\subsection{Persistance analytique des racines déformées}
\label{hmtfg:desarrollo:8}

Pour chaque niveau certifié, \(F_n'(\lambda_n)\ne0\). Le théorème des fonctions implicites produit une branche analytique locale \(\lambda_n(\eta)\) avec
\[
f_{\lambda_n(\eta),\eta}^{2^n}(0)=0.
\]
Par continuité, les retours propres restent séparés de zéro pour \(\eta\) suffisamment petit ; la période minimale et l’itinéraire sont préservés. L’existence de ces voisinages locaux est démontrée. Leur largeur commune n’est pas encore identifiée à toute la bande structurelle \(1/40\).

Avec \(M=M_{2,\eta}\), la réponse du profil est
\[
v_\eta(x)=\partial_\eta u_\eta(x)
=\sum_m\frac{p_m}{12}(1-\cos(k_{m,\eta}x))
-\frac{M'}{2M}\,x u_\eta'(x).
\]
Le second terme conserve la normalisation du deuxième moment. Si \(q_j=\partial_\eta x_j\) à \(\lambda\) fixé,
\[
q_0=0,\qquad q_{j+1}=-\lambda v_\eta(x_j)-\lambda u_\eta'(x_j)q_j,
\qquad
\boxed{\lambda_n'(\eta)=-q_{2^n}/p_{2^n}}.
\]
La chaîne poids–échelle–réponse–paramètre superstable est ainsi explicitée. Cette non-dégénérescence finie est distincte de la transversalité asymptotique à la feuille stable du point fixe universel.


\subsection{Queues analytiques et propagation de l’erreur}
\label{hmtfg:desarrollo:9}

Pour \(\eta=0\), soient
\[
\mu_{2j}=\frac1{12}\sum_mk_m^{2j},\quad
p_N(z)=\sum_{j=1}^N(-1)^{j+1}\frac{\mu_{2j}}{(2j)!}z^{2j},
\quad\Omega=\frac{70}{\sqrt{899}}.
\]
Si \(q_N=\Omega^2r^2/[(2N+3)(2N+4)]<1\),
\[
\|u_0-p_N\|_r\le
\frac{\mu_{2N+2}r^{2N+2}}{(2N+2)!(1-q_N)}.
\]
L’inégalité \(\mu_{2j+2}\le\Omega^2\mu_{2j}\) majore la queue par une série géométrique. Pour \(0\le s\le2N+2\),
\[
\|(u_0-p_N)^{(s)}\|_r\le
\frac{\mu_{2N+2}r^{2N+2-s}}{(2N+2-s)!}
\left(1-\frac{\Omega^2r^2}{(2N+3-s)(2N+4-s)}\right)^{-1},
\]
lorsque le dénominateur est positif.

Pour transporter les erreurs par un retour, on contrôle également ses domaines : si \(\|f-g\|\le\varepsilon\), \(a_f,a_g\ge a_*>0\), que les compositions restent dans un domaine convexe commun et que l’on y a \(\|f'\|\le L,\ \|g\|\le M\), alors
\[
\|\mathcal Rf-\mathcal Rg\|_r
\le\varepsilon\left[\frac{1+L+L^2r}{a_*}+\frac{M}{a_*^2}\right].
\]
La preuve additionne la perturbation extérieure, la perturbation intérieure, la modification de l’argument due à \(|a_f-a_g|\le\varepsilon\) et la variation du dénominateur. Ces bornes évitent de réutiliser la somme initiale de cosinus comme si tous les profils renormalisés conservaient cette même forme.

% Fragmento integrable. Preambulo anfitrion: amsmath, amssymb, longtable, hyperref.
\section{Profondeur nonadique et persistance de l’itinéraire}
\label{hmtfg:fragmento:profundidad-nonadica-y-persistencia-del-itinerario}

\subsection{Croissance extérieure et résolution intérieure}
\label{hmtfg:prolongacion:2}

La loi de raffinement de l’article XI, de base entière \(B\ge2\), est
\[
L_c(x,y)=(Bx-c,By+c),\qquad 0\le c<B.
\]
Pour un mot \(c_1\ldots c_m\), son registre est
\[
C_m=\sum_{j=1}^m c_jB^{m-j},\qquad
(x_m,y_m)=(B^mx_0-C_m,B^my_0+C_m).
\]
La somme extérieure vaut \(B^mM_0\), avec \(M_0=x_0+y_0\). La lecture normalisée du registre appartient au cylindre
\[
I_m(C_m)=\left[\frac{C_m}{B^m},\frac{C_m+1}{B^m}\right],
\qquad \operatorname{diam}I_m=B^{-m}.
\]
Le couple entier et l’échelle restituent les quotients, les restes et les zéros initiaux du mot. Les extrémités à double développement conservent leur marque de frontière. Ainsi se coordonnent la croissance du support et le raffinement de la lecture sans identifier des histoires distinctes par leur seule image limite.

La composition est exacte :
\[
L_D^{[b]}\circ L_C^{[a]}=L_{B^bC+D}^{[a+b]}.
\]
Par associativité, regrouper une histoire de dix-huit événements en neuf blocs de deux ou en deux blocs de neuf produit le même registre. Au niveau de retour \(2^N\), l’horizon commun est \(9\,2^N\). La preuve est une égalité de composition sur la même histoire ordonnée ; elle conserve ses retenues et ses sous-blocs récupérables.

La restriction des états et le prolongement de leurs lecteurs s’expriment par les limites suivantes :
\[
X_\infty=\varprojlim X_m,\qquad
r^*a=a\circ r,\qquad
r^*e_x=\sum_{r(y)=x}e_y,
\qquad \mathfrak H_{\rm cyl}=\varinjlim\mathfrak H_m.
\]
Les états se restreignent ; les lecteurs se prolongent. L’image réciproque préserve les produits, l’unité et les idempotents. Une famille exactement compatible induit un lecteur sur la limite inductive. Les approximants qui ne sont compatibles qu’à une erreur près exigent en outre une estimation de convergence dans leur complétion. La section suivante fournit cette estimation pour le retour critique.

\subsubsection{Capacité et vacances du registre}

Le lecteur de capacité de XI compare des blocs de \(729\) et de \(1000\) possibilités. Avec \(\rho=\log_{1000}729=\log_{10}9\), il conserve
\[
\mathcal C_k=\lfloor(k+1)\rho\rfloor,\qquad
\mathcal C_k-\mathcal C_{k-1}=1-z_k,
\]
\[
\mathcal C_b-\mathcal C_a+\sum_{k=a+1}^b z_k=b-a,
\qquad T_{\rm cap}(a)=\lceil a/\rho\rceil-1\quad(a\ge1).
\]
Chaque vacance enregistre une transition pendant laquelle la capacité se maintient ; la longueur chronologique est conservée dans le bilan. La phase de capacité et la phase chronologique sont des lecteurs distincts.

Pour stocker un préfixe de \(m\) chiffres en base neuf, une capacité décimale suffisante est
\[
a(m)=\left\lceil\frac{m\rho}{3}\right\rceil.
\]
Le quotient et le reste permettent de convertir l’entier \(C_m\) entre bases, en conservant séparément la profondeur \(m\). On a
\[
9^m\le1000^{a(m)}\le729^{T_{\rm cap}(a(m))+1}.
\]
Ces inégalités se décident également au moyen de puissances entières, sans arrondir les logarithmes. Le budget de lecture de la section suivante se traduit ainsi en capacité du registre en conservant les unités. Ce calcul dimensionne le stockage ; les règles de survie déterminent en outre quelles histoires natives sont admissibles.


\subsection{Budget explicite de profondeur pour chaque retour}
\label{hmtfg:prolongacion:3}

\textbf{Théorème.} Pour \(|\eta|\le1/40\) fixé, normalisons
\[
t=\lambda u_\eta(1)/2\in[0,1],\qquad z=(x+1)/2,
\]
\[
F_{t,\eta}(z)=1-t\frac{u_\eta(2z-1)}{u_\eta(1)},\qquad z\in[0,1].
\]
Pour tout entier \(q\ge1\) et deux paramètres \(t,s\),
\[
\left|F_{t,\eta}^{q}(1/2)-F_{s,\eta}^{q}(1/2)\right|
\le S_q|t-s|,
\qquad S_q=\sum_{j=0}^{q-1}(6\sqrt2)^j<9^q.
\]
Par conséquent, un cylindre de paramètres en base neuf, de profondeur \(m=2^N+r\), enferme la lecture du retour \(q=2^N\) dans un intervalle de diamètre inférieur à \(9^{-r}\).

\textbf{Démonstration.} La fonction \((1-\cos v)/v^2\) décroît pour \(0<v<\pi\) : le signe de sa dérivée est celui de \(v\sin v-2(1-\cos v)\), négatif car \(\tan(v/2)>v/2\). Par le deuxième moment,
\[
u_\eta(1)\ge\frac{1-\cos3}{9}\sum_jw_jk_j^2
=\frac{2(1-\cos3)}9>\frac13.
\]
La dernière inégalité se vérifie, par exemple, en majorant \(\cos3\) par sa somme de Taylor alternée jusqu’au degré huit, qui est inférieure à \(-1/2\). De plus, Cauchy–Schwarz donne \(|u_\eta'(x)|\le\sqrt2\) sur la droite réelle. On en déduit
\[
\operatorname{Lip}_zF_{t,\eta}\le6\sqrt2<9,\qquad
\operatorname{Lip}_tF_{t,\eta}\le1.
\]
Si \(d_j\) est la différence entre les deux orbites, alors \(d_0=0\) et \(d_{j+1}\le6\sqrt2\,d_j+|t-s|\). La récurrence prouve la somme géométrique. Pour \(|t-s|\le9^{-m}\), le diamètre est inférieur à \(9^{q-m}\). La substitution de \(q=2^N\) et de \(m=q+r\) achève la preuve. \(\square\)

La borne contrôle l’incertitude du paramètre pour \(\eta\) fixé. Une implémentation ajoute et propage séparément l’erreur d’évaluation des cosinus et l’erreur d’arrondi ; le certificat rationnel antérieur fournit ces opérations. La borne est suffisante et conservatrice, avec un coût qui croît avec l’horizon.

\textbf{Lecteur limite.} Réalisons \(t\) à partir de \((x_0,y_0)=(1,0)\), \(B=9\), de sorte que
\[
(x_m,y_m)=(9^m-C_m,C_m),\qquad t_m=C_m/9^m.
\]
Les lecteurs \(E_{q,m}=F_{t_m,\eta}^{q}(1/2)\), ramenés à un niveau commun, satisfont
\[
\|E_{q,m+a}-r^*E_{q,m}\|_\infty\le S_q9^{-m}.
\]
Ils convergent uniformément vers un lecteur continu \(E_q\) dans la complétion de l’algèbre cylindrique. La version ensembliste
\[
J_{q,m}(C)=\{F_{t,\eta}^{q}(1/2):t\in I_m(C)\}
\]
est exactement emboîtée et ses diamètres tendent vers zéro. Ce résultat donne le passage à la limite de la \textbf{lecture à horizon fixé}, avec contrôle de profondeur à tout horizon demandé.


\subsection{Mot de doublement et lecteur nonadique de la chronologie}
\label{hmtfg:prolongacion:4}

Les huit itinéraires certifiés dans le travail antérieur obéissent à
\[
W_1=+,\qquad W_{n+1}=W_n\,(-1)^n\,W_n.
\]
Son unique prolongement préfixiel est
\[
\tau_j=(-1)^{v_2(j)},\qquad j\ge1;
\qquad \tau_{2j}=-\tau_j,\quad\tau_{2j+1}=+.
\]
La preuve utilise \(v_2(2^n+r)=v_2(r)\) pour \(0<r<2^n\). C’est une récurrence à toute profondeur, et la concordance avec les huit retours calculés constitue sa vérification finie dans la famille.

Dans l’ordre lexicographique à parité d’une application à maximum critique, le facteur d’orientation précédant la position \(k\) est \(\prod_{j<k}(-\tau_j)\). Le mot \(\tau\) est strictement supérieur à chacun de ses décalages positifs. En effet, le premier désaccord avec un décalage impair se produit en \(k=1\) ou \(2\) ; un décalage pair réduit la comparaison de moitié et double \(k\). Le premier désaccord se situe donc en \(k=2^r\). À cet endroit
\[
\prod_{j<2^r}(-\tau_j)=(-1)^r=\tau_{2^r},
\]
et la différence orientée de signe opposé est positive. Cette démonstration prouve également l’apériodicité.

Le registre chronologique conserve
\[
K=\rho_9(K)+9q_9(K).
\]
Son lecteur dyadique est
\[
\ell_n(B_K)=[\rho_9(K)+9q_9(K)]_{2^n}=[K]_{2^n}.
\]
Sur la carte des histoires où \(\Gamma_9 B_K=B_{K+9}\),
\[
\ell_n(\Gamma_9B_K)=\ell_n(B_K)+9\pmod{2^n},\qquad
\pi_{n+1,n}\ell_{n+1}=\ell_n.
\]
Comme neuf est impair, l’avancement de neuf visite les \(2^n\) positions. Sa conjugaison avec l’avancement unitaire est \(j\mapsto9j\). Pour \(0<j<2^n\),
\[
v_2(9j\bmod2^n)=v_2(j),
\]
de sorte que le rééchantillonnage nonadique conserve également le mot orienté dans chaque cycle certifié.

Ce lecteur utilise le reste \textbf{et} le quotient. Les phases de \(K=1\) et de \(K=10\) coïncident modulo neuf, tandis que leurs signes dyadiques sont opposés. La mémoire distingue les deux histoires. La propriété arithmétique du rééchantillonnage vaut pour tout multiplicateur impair ; le calendrier HMT fixe ici le choix de neuf.


\subsection{Mémoire bilatérale compatible à profondeur infinie}
\label{hmtfg:prolongacion:5}

Soient \(\mathcal H_n=\ell^2(\mathbb Z/2^n\mathbb Z)\),
\[
C_ne_j=e_{j+1},\qquad
J_ne_j=(e_j+e_{j+2^n})/\sqrt2.
\]
Une vérification sur la base donne \(J_n^*J_n=I\) et \(C_{n+1}J_n=J_nC_n\), y compris les termes de bord.

L’article XI fournit les projecteurs
\[
P_0=\frac19\begin{pmatrix}8&-\sqrt8\\-\sqrt8&1\end{pmatrix},\qquad
P_1=\frac19\begin{pmatrix}1&\sqrt8\\\sqrt8&8\end{pmatrix}.
\]
Ils sont orthogonaux et complémentaires. Par conséquent
\[
\mathbb U_n=P_0\otimes I+P_1\otimes C_n
=\begin{pmatrix}T_n&D_n\\D_n&R_n\end{pmatrix}
\]
est unitaire, avec
\[
T_n=(8I+C_n)/9,\quad D_n=\sqrt8(C_n-I)/9,\quad R_n=(I+8C_n)/9.
\]
La naturalité déjà démontrée implique
\[
\mathbb U_{n+1}(I_2\otimes J_n)=(I_2\otimes J_n)\mathbb U_n,
\quad \mathbb U_n^a=P_0\otimes I+P_1\otimes C_n^a
\quad(a\in\mathbb Z).
\]
La limite hilbertienne s’identifie à \(L^2(\mathbb Z_2,\mu_{\rm Haar})\), en envoyant \(e_j^{(n)}\) sur \(2^{n/2}\mathbf1_{j+2^n\mathbb Z_2}\). Les identités se prolongent par densité et bornitude. Pour tout vecteur \(v\) et tout temps entier \(a\),
\[
\boxed{\|T_{\infty,a}v\|^2+\|D_{\infty,a}v\|^2=\|v\|^2,}
\]
\[
\boxed{v=T_{\infty,a}^*(T_{\infty,a}v)+D_{\infty,a}^*(D_{\infty,a}v).}
\]
Ici \(T_{\infty,a}=(8I+C_\infty^a)/9\). La composition bilatérale conserve l’archive entre les pas ; composer seulement \(T_n\) décrit une autre opération.

La permutation \(S_ne_j=e_{9j\bmod2^n}\) satisfait
\[
S_{n+1}J_n=J_nS_n,\qquad S_nC_nS_n^{-1}=C_n^9.
\]
On obtient ainsi à la limite
\[
S_\infty C_\infty S_\infty^{-1}=C_\infty^9,
\quad (I_2\otimes S_\infty)\mathbb U_\infty
(I_2\otimes S_\infty^{-1})=\mathbb U_\infty^9.
\]
La représentation temporelle est fidèle : \(C_\infty^a=C_\infty^b\) impose \(2^n\mid a-b\) pour tout \(n\), d’où \(a=b\). Des préparations particulières peuvent conserver des symétries temporelles. L’identité operatorielle conserve cette distinction.

Le point projectif \(([K]_{2^n})_n\) et un vecteur de \(L^2\) sont de types différents : une histoire ponctuelle induit une mesure de Dirac, tandis que les amplitudes normalisables appartiennent à l’espace de Hilbert. Le résultat construit un lecteur équivariant de la chronologie ; les autres coordonnées natives du TPK accompagnent le registre complet.


\subsection{Existence et séparation des obligations asymptotiques}
\label{hmtfg:prolongacion:6}

\subsubsection{Réalisation du mot infini sur toute la bande}
\label{hmtfg:prolongacion:6.1}

\textbf{Théorème d’existence.} Pour chaque \(|\eta|\le1/40\) il existe au moins un
\[
\lambda_\infty(\eta)\in
\left[\frac1{2u_\eta(1)},\frac2{u_\eta(1)}\right]
\]
tel que
\[
\boxed{\operatorname{sign}f_{\lambda_\infty(\eta),\eta}^{j}(0)
=(-1)^{v_2(j)}\quad\text{pour tout }j\ge1.}
\]
La conclusion réalise l’itinéraire infini de doublement au sein de la famille fixée. Le quantificateur d’existence reste séparé d’une sélection unique et de la loi des distances entre paramètres superstables.

\textbf{Démonstration.} Nous fixons \(\eta\). Nous écrivons \(K_\lambda\) pour le mot de signes de l’orbite critique ; lorsque celle-ci revient pour la première fois à zéro, le mot se termine par ce zéro. Nous utilisons l’ordre unimodal à maximum, dont le facteur avant la comparaison du symbole \(j\) est le produit des signes \(-K_i\), \(i<j\).

À l’extrémité gauche, l’image de tout l’intervalle est dans \([1/2,1]\) ; donc \(K_-=+++\cdots\). À l’extrémité droite, l’orbite est \(0\mapsto1\mapsto-1\mapsto-1\), d’où \(K_+=+---\cdots\). La comparaison des deuxième et troisième symboles, respectivement, donne
\[
K_-<\tau<K_+.
\]

Si \(K_{\lambda_0}\) est infini et distinct de \(\tau\), la première différence se produit à une position finie. La continuité des itérées correspondantes rend localement constant le côté de la comparaison. Il reste à examiner un paramètre superstable de première période \(p\), de mot \(W0\), où \(|W|=p-1\).

Pour \(\lambda\) proche de \(\lambda_0\), soient \(g_\lambda=f_{\lambda,\eta}^p\) et \(a=g_\lambda(0)\). On a \(g_\lambda'(0)=0\), et une borne locale uniforme de la dérivée seconde donne
\[
g_\lambda(x)=a+O(x^2).
\]
Pour \(|a|\) suffisamment petit et non nul, \(g_\lambda\) envoie \([-2|a|,2|a|]\) dans \([a-|a|/2,a+|a|/2]\). Les retours conservent le signe de \(a\), et les \(p-1\) positions intermédiaires conservent \(W\). Les mots voisins sont donc \((W+)^\infty\) ou \((W-)^\infty\) ; les paramètres avec \(a=0\) conservent \(W0\).

Si \(\tau\) diffère de \(W\) avant la position \(p\), cette différence fixe le même côté pour les trois mots. Si elle partage \(W\), posons
\[
s=\tau_p,\qquad e=\prod_{j<p}(-\tau_j),\qquad A=Ws.
\]
Le signe de la comparaison de \(\tau\) avec \(W0\) et avec \((W,-s)^\infty\) est \(es\). Pour la comparer à \(A^\infty\), soit \(q\) la première position où \(\tau_{p+q}\ne\tau_q\). Elle existe par apériodicité. Jusqu’à la position \(p+q-1\), les deux mots coïncident, et le signe d’orientation se factorise en \(O_pO_{q-1}\), avec \(O_p=-es\). La maximalité stricte déjà démontrée donne
\[
O_{q-1}(\tau_q-\tau_{p+q})>0.
\]
Par conséquent, la comparaison de \(\tau\) avec \(A^\infty\) a pour signe \(-O_p=es\). Les trois possibilités locales sont du même côté de \(\tau\).

Si aucun paramètre ne réalisait \(\tau\), les ensembles \(\{\lambda:K_\lambda<\tau\}\) et \(\{\lambda:K_\lambda>\tau\}\) seraient ouverts, disjoints, non vides et couvriraient l’intervalle des paramètres. La connexité de l’intervalle exclut cette partition. Le paramètre annoncé existe donc. \(\square\)

Cet argument de continuité des itinéraires s’applique après fixation du générateur HMT et de sa famille analytique ; ses hypothèses sont vérifiées ici directement. La construction conserve le mot infini sans introduire comme donnée un paramètre de la famille quadratique ni la valeur universelle.

\subsubsection{Marges empêchant de perdre l’itinéraire lors du passage à la limite}
\label{hmtfg:prolongacion:6.2}

Écrivons \(c_p=f_{\lambda,\eta}^p(0)\). Une borne rationnelle uniforme suffisante est
\[
|f'|,|f''|\le16,\qquad
B_p:=16^p\frac{16^p-1}{15}\ge\|(f^p)''\|_{[-1,1]}.
\]
En effet, \(1-\cos v\ge v^2/2-v^4/24\ge v^2/8\) pour \(|v|\le3\), de sorte que \(u_\eta(1)\ge1/4\), \(\lambda\le8\) et les deux dérivées sont majorés par \(16\), en utilisant les moments. La règle de dérivation en chaîne prouve par récurrence la borne \(B_p\).

Le mot \(\tau\) satisfait \(\tau_{2p}=-\tau_p\) pour tout \(p\). Pour chacune de ses réalisations, \(c_{2p}\) et \(c_p\) sont de signes opposés. Comme \((f^p)'(0)=0\), Taylor donne
\[
|c_{2p}-c_p|\le\tfrac12B_p|c_p|^2.
\]
Le membre de gauche est supérieur à \(|c_p|\), et donc
\[
\boxed{|c_p|>2/B_p>0.}
\]
Pour chaque horizon fixé, il existe ainsi une marge uniforme qui empêche une suite de réalisations de perdre son signe en convergeant. Dans les coordonnées compactes \((\eta,b)\), \(b=\lambda u_\eta(1)\in[1/2,2]\), l’ensemble des réalisations de \(\tau\) est fermé et compact, et sa projection sur toute la bande de \(\eta\) est surjective.

Le même argument vaut pour un arbre de préfixes : si tous les préfixes de \(\tau\) sont réalisés, une sous-suite dans l’ensemble compact conserve chaque signe, car les préfixes suffisamment longs comprennent simultanément les positions \(p\) et \(2p\). Le théorème d’existence fournit maintenant la non-vacuité nécessaire à cet argument.

\subsubsection{Intervalles invariants de retour à tous les niveaux}
\label{hmtfg:prolongacion:6.3}

\textbf{Théorème des retours emboîtés.} Pour chacune des réalisations du théorème d’existence de la section \ref{hmtfg:prolongacion:6.1}, soient \(c_j=f^j(0)\) et
\[
J_n=\operatorname{conv}\{c_{2^n},c_{2^{n+1}}\},\qquad n\ge1.
\]
Alors \(0\in\operatorname{int}J_n\), \(J_{n+1}\subset J_n\),
\[
f^{2^n}(J_n)=J_n,
\]
et les \(2^n\) images \(J_n,f(J_n),\ldots,f^{2^n-1}(J_n)\) sont disjointes. Les retours normalisés sont unimodaux à maximum, de valeur critique un et de même mot \(\tau\). Ces intervalles sont des noyaux invariants de retour ; une normalisation exigeant des extrémités périodiques peut élargir le noyau à l’intervalle restrictif maximal correspondant et doit préciser cette étape.

\textbf{Démonstration du premier retour.} Soit \(F\) une application unimodale à maximum, avec \(F(0)=1\) et de mot \(\tau\), et \(d_j=F^j(0)\). Ses premiers signes sont \(+,-,+,+,+,-,+\). L’intervalle \(I=[d_2,1]\) est invariant : \(F(d_2)=d_3>0\), \(F(1)=d_2<0\), et le maximum vaut un. Posons \(J=[d_2,d_4]\). Comme \(F\) décroît sur les arguments positifs et
\[
F(d_3)=d_4>0> d_6=F(d_5),
\]
on a \(d_3<d_5\), et donc \(F(J)=[d_3,1]\).

De plus, \(d_3>d_4\). Pour le vérifier, \(F^2\) est strictement croissante entre les deux points positifs \(d_3,d_4\), car leurs images \(d_4,d_5\) sont positives. Ses valeurs sont \(F^2(d_3)=d_5>0>d_6=F^2(d_4)\), ce qui impose cet ordre. Ainsi \(J\) et \(F(J)\) sont séparés par un intervalle ouvert, et
\[
F^2(J)=F([d_3,1])=[d_2,d_4]=J.
\]
L’application \(F^2|_J\) possède un unique minimum critique : l’image \(F(J)\) reste positive. La conjugaison \(h(x)=d_2x\) inverse l’orientation et normalise le retour en un maximum de valeur un. Son itinéraire est
\[
\operatorname{sign}\frac{d_{2j}}{d_2}
=-\tau_{2j}=\tau_j.
\]
Le même argument peut être répété à toute profondeur.

\textbf{Disjonction par récurrence.} Supposons disjointes les \(p=2^n\) images du parent \(J_n\), et soit \(g=f^p\). Le premier retour appliqué à l’application normalisée produit \(B=J_{n+1}\), \(A=g(B)\), avec \(A\cap B=\varnothing\), \(g(B)=A\), \(g(A)=B\). Pour \(0\le j<p\), une intersection de \(f^j(A)\) et de \(f^j(B)\) produirait, par application de \(f^{p-j}\), une intersection de \(B\) et de \(A\). Ces deux images sont donc disjointes. Les images correspondant à des \(j\) distincts appartiennent à des images disjointes du parent. Les \(2p\) images disjointes et le retour \(f^{2p}(B)=B\) sont ainsi démontrés. La formule des extrémités résulte de la composition des normalisations \(x\mapsto c_{2^n}x\). \(\square\)

Les ensembles compacts
\[
\Lambda_n=\bigcup_{0\le j<2^n}f^j(J_n)
\]
sont non vides et emboîtés. Leur intersection \(\Lambda_\infty\) possède un lecteur continu sur l’odomètre dyadique : à chaque point est attribuée l’étiquette de l’unique composante qu’il occupe à chaque niveau. Toute suite compatible d’étiquettes possède une préimage par compacité, et l’avancement par \(f\) ajoute un aux étiquettes. Cette application est une semi-conjugaison surjective. Sa promotion en conjugaison ponctuelle exige de démontrer que les diamètres des composantes tendent vers zéro ; cette promotion métrique reste séparée.

% Fragmento integrable. Preambulo anfitrion: amsmath, amssymb, longtable, hyperref.
\section{Compatibilité des lecteurs et restitution du paramètre}
\label{hmtfg:fragmento:compatibilidad-de-lectores-y-restitucion-del-parametro}

\subsection{Compatibilité des lecteurs, des minima et restitution de l’histoire}
\label{hmtfg:escalado:13}

Les opérations de prolongement et les lecteurs ordonnés conservent l’information requise pour la sélection des racines. Les identités suivantes fixent leurs domaines et leur composition.

\subsubsection{Règle forward restituée}
\label{hmtfg:escalado:13.1}

L’actualisation déterministe part de l’état de frontière R36 avec le caractère \(\chi\) et le reste interne. L’actualisation est
\begin{equation}
\mathcal U_K^\chi=
\operatorname{Upd}_K^\chi\circ
\operatorname{Tra}_K^\chi\circ
\operatorname{Sel}_K^\chi,
\qquad
N_{K+1}=729N_K+d_K^\chi.
\label{hmtfg:escalado:eq:26}
\end{equation}
Le chiffre et l’actualisation dépendent de l’état précédent. Deux sections forward ayant le même état complet R36 coïncident par induction. La sous-relation \(\operatorname{Ext}^{\chi,\rightarrow}\) est le graphe de cette actualisation ; la relation totale \(\operatorname{Ext}\) admet une multiplicité de continuations selon les données et les caractères.

Les prolongements compatibles transportent les lecteurs régionaux et la signature conjointe sur les mêmes préfixes.

Ces règles transportent l’orientation, le reste, le quotient, la retenue, la frontière et la mémoire. Leur application à la cascade doit également conserver la sélection ordonnée du paramètre du lecteur critique.

\subsubsection{Indépendance des constantes antécédentes à l’égard du paramètre libre}
\label{hmtfg:escalado:13.2}

Soit \(B\) l’état HMT qui fournit la roue et ses constantes antécédentes. Notons \(\mathcal C(B)\) l’ensemble de ces lectures et \(u_B\) le profil critique. Dans le domaine du lecteur critique, la construction ultérieure est
\begin{equation}
(B,\lambda)\longmapsto f_{B,\lambda}(x)=1-\lambda u_B(x),
\qquad 0<\lambda\le 2/u_B(1).
\label{hmtfg:escalado:eq:27}
\end{equation}
Pour un même \(B\), la lecture des constantes antérieures se factorise par la projection sur la première composante :
\begin{equation}
\widetilde{\mathcal C}(B,\lambda)
=\mathcal C\circ\operatorname{pr}_B(B,\lambda)
=\mathcal C(B).
\label{hmtfg:escalado:eq:28}
\end{equation}
Par conséquent, deux valeurs distinctes de \(\lambda\) conservent les mêmes lectures antécédentes. L’égalité est une propriété de cette construction causale : le profil est fixé avant de faire varier \(\lambda\), et le facteur \(\pi_{\mathrm{HMT}}\) se simplifie dans sa normalisation quadratique.

\textbf{Lemme.} Supposons qu’une condition \(A\) soit évaluée uniquement sur ces lectures antécédentes. Pour tout \(B\) fixé,
\begin{equation}
\{\lambda\in\Lambda:A(\widetilde{\mathcal C}(B,\lambda))\}
=
\begin{cases}
\Lambda,&A(\mathcal C(B))\text{ est satisfaite},\\
\varnothing,&A(\mathcal C(B))\text{ échoue}.
\end{cases}
\label{hmtfg:escalado:eq:29}
\end{equation}
\textbf{Démonstration.} La substitution de \eqref{hmtfg:escalado:eq:28} rend le prédicat indépendant de \(\lambda\). Par conséquent, tous les paramètres du domaine le satisfont, ou aucun. Cela inclut une comparaison des constantes antécédentes avec des intervalles métrologiques.

L’unicité de la section \(B\) et la restriction ultérieure \(f_{B,\lambda}^{2^n}(0)=0\) produisent une fibre de solutions dans \(\lambda\). La règle de première racine agit sur cette fibre par le retour et son itinéraire. La sélection et son prolongement sont démontrés dans les résultats suivants.

Ce résultat maintient l’ordre HMT : la métrologie confronte des sorties préalablement construites et la sélection de la racine se démontre dans le lecteur qui la produit.

\subsubsection{Lemme de transport ordonné du sélecteur}
\label{hmtfg:escalado:13.3}

Soit \(P_n\) un ensemble ordonné de candidats généalogiques de minimum \(p_n^*\). Soit \(Z_n\) l’ensemble initial des racines candidates du corpus, muni de l’ordre réel, et soit \(L_n:P_n\to Z_n\) un lecteur monotone. Son image est dite \textbf{coinitiale} lorsque
\begin{equation}
\forall z\in Z_n\quad\exists p\in P_n:
L_n(p)\le z.
\label{hmtfg:escalado:eq:30}
\end{equation}
Alors
\begin{equation}
\boxed{L_n(p_n^*)=\min Z_n
\quad\Longleftrightarrow\quad
L_n(P_n)\text{ est coinitial dans }Z_n.}
\label{hmtfg:escalado:eq:31}
\end{equation}
\textbf{Démonstration.} Si l’image du minimum est le minimum global, on prend \(p=p_n^*\) dans \eqref{hmtfg:escalado:eq:30}. Réciproquement, étant donné \(z\in Z_n\), la coinitialité fournit \(p\) avec \(L_n(p)\le z\). La monotonie donne \(L_n(p_n^*)\le L_n(p)\le z\). Puisque \(L_n(p_n^*)\in Z_n\), c’est le minimum global. La surjectivité de \(L_n\) constitue une condition suffisante, plus forte que \eqref{hmtfg:escalado:eq:30}.

L’isomorphisme ordonné de la carte du continu satisfait la surjectivité du lecteur. L’application concrète au retour et à son ensemble complet de zéros est démontrée dans la section de transport ordonné ; la classification et l’isolement déterminent ensuite sa continuation dynamique.

\subsubsection{Sélection minimale par survie à profondeur arbitraire}
\label{hmtfg:escalado:13.6}

La composition avec la survie produit également un sélecteur global précis. On conserve la famille uniforme \(\eta=0\), construite dans la section~\ref{hmtfg:escalado:1}, et l’on utilise les théorèmes de réalisation de l’itinéraire infini et de conservation de ses signes. Soient
\[
I=\left[\frac1{2u_0(1)},\frac2{u_0(1)}\right],\qquad
c_j(\lambda)=f_\lambda^j(0),\qquad
\tau_j=(-1)^{v_2(j)}.
\]
L’ensemble
\[
\Lambda_\tau=\{\lambda\in I:\operatorname{sign}c_j(\lambda)=\tau_j
\text{ pour tout }j\ge1\}
\]
est compact et non vide. Le contrôle des dérivées de cette source fournit
\begin{equation}
B_j=\frac{16^j(16^j-1)}{15},\qquad
|c_j(\lambda)|>\frac2{B_j}\quad(\lambda\in\Lambda_\tau).
\label{hmtfg:escalado:eq:32}
\end{equation}
Définissons
\begin{equation}
\varepsilon_j=\frac1{B_j},\qquad
A_N=\{\lambda\in I:\tau_jc_j(\lambda)\ge\varepsilon_j,\
1\le j\le N\}.
\label{hmtfg:escalado:eq:33}
\end{equation}
Chaque \(A_N\) est compact par continuité des itérées et non vide parce qu’il contient \(\Lambda_\tau\). De plus,
\begin{equation}
A_{N+1}\subseteq A_N,\qquad
\bigcap_{N\ge1}A_N=\Lambda_\tau.
\label{hmtfg:escalado:eq:34}
\end{equation}
L’inclusion de droite à gauche découle de \eqref{hmtfg:escalado:eq:32}. L’inclusion inverse découle des marges strictement positives de \eqref{hmtfg:escalado:eq:33}, qui fixent chaque signe.

\textbf{Théorème du minimum survivant.} Les minima \(a_N=\min A_N\) satisfont
\begin{equation}
\boxed{a_N\uparrow a_*=\min\Lambda_\tau.}
\label{hmtfg:escalado:eq:35}
\end{equation}
\textbf{Démonstration.} L’inclusion \eqref{hmtfg:escalado:eq:34} rend croissante la suite des minima, bornée dans \(I\) ; elle admet donc une limite \(a\). Pour chaque \(r\), tous les termes \(a_N\), \(N\ge r\), appartiennent au fermé \(A_r\) ; ainsi \(a\in A_r\). Par \eqref{hmtfg:escalado:eq:34}, \(a\in\Lambda_\tau\). Pour tout \(\lambda\in\Lambda_\tau\), on a \(a_N\le\lambda\) à chaque niveau et, en passant à la limite, \(a\le\lambda\). Cela prouve \eqref{hmtfg:escalado:eq:35}.

Dans le lecteur cylindrique HMT, l’objet \(a_*\) possède des préfixes compatibles à toute profondeur, avec conservation de l’orientation, du reste et de la mémoire. Son choix utilise l’ordre de la lecture paramétrique et la survie complète. La compacité démontre l’existence et la convergence de ces minima.

Les paramètres \(a_N\) résolvent des contraintes de signes avec marge. Leurs frontières peuvent satisfaire \(\tau_jc_j=\varepsilon_j\), tandis que les racines superstables satisfont \(c_{2^n}=0\). Ainsi, \eqref{hmtfg:escalado:eq:35} établit un sélecteur global de réalisation infinie et préserve, comme objet distinct, le sélecteur de premières racines du corpus.

\subsubsection{Restitution du paramètre à partir de l’histoire critique complète}
\label{hmtfg:escalado:13.7}

Le deuxième terme de l’histoire critique fournit un lecteur inverse exact :
\begin{equation}
c_1(\lambda)=1,\qquad
c_2(\lambda)=1-\lambda u_0(1),\qquad
\boxed{\lambda=\frac{1-c_2(\lambda)}{u_0(1)}.}
\label{hmtfg:escalado:eq:36}
\end{equation}
La positivité \(u_0(1)>0\), démontrée pour le profil, rend la division valide. Par conséquent, deux histoires critiques complètes identiques ont le même paramètre. Le mot de signes \(\tau\) constitue une lecture réduite de cette histoire ; il conserve sa combinatoire et omet la grandeur de \(c_2\).

Les formules \eqref{hmtfg:escalado:eq:35}–\eqref{hmtfg:escalado:eq:36} unissent la sélection par survie et la restitution du paramètre à partir de l’histoire critique. L’égalité de sa limite avec le croisement stable est démontrée par première sortie, et l’identification des centres par le transport analytique de leurs classes hybrides.

% Fragmento integrable. Preambulo anfitrion: amsmath, amssymb, longtable, hyperref.
\section{Transport ordonné et représentation de Jacobi}
\label{hmtfg:fragmento:transporte-ordenado-y-representacion-de-jacobi}

\subsection{Transport complet depuis le continu et réalisation spectrale du profil}
\label{hmtfg:escalado:14}

Le transport ordonné du continu et la représentation spectrale du lecteur se composent avant la sélection paramétrique.

\subsubsection{Couverture de toutes les racines et transport de leur minimum}
\label{hmtfg:escalado:14.1}

La construction ordonnée du continu développée dans les travaux antérieurs produit la lecture archimédienne surjective à partir de cylindres compatibles et de leur quotient ordonné. Dans son univers déclaré \(\mathfrak G\), elle démontre l’isomorphisme de corps ordonnés complets
\begin{equation}
\iota:\mathbb R_{\rm HMT}^{\mathfrak G}
   \xrightarrow{\cong}\mathbb R^{\mathfrak G}.
\label{hmtfg:escalado:eq:37}
\end{equation}
Le nom \(\eta\) du document de référence est remplacé ici par \(\iota\), afin de le distinguer du paramètre de pondération, fixé à zéro. La seconde projection conserve l’incidence, la frontière et la mémoire sur le même antécédent ; le quotient d’évaluation et l’histoire complète conservent leurs domaines respectifs.

Construisons dans le corps HMT, avant de sélectionner une racine,
\[
\cos_{\rm H}(t)=\sum_{k=0}^{\infty}\frac{(-1)^kt^{2k}}{(2k)!},
\qquad
u_{\rm H}(x)=\frac1{12}\sum_{m=1}^{12}
\left[1-\cos_{\rm H}\!\left(\frac{2(3m-1)x}{\sqrt[\rm H]{899}}\right)\right],
\]
\begin{equation}
F_a(x)=1-a\,u_{\rm H}(x),\qquad S_n^{\rm H}(a)=F_a^{\,2^n}(0).
\label{hmtfg:escalado:eq:38}
\end{equation}
Les entiers et le deuxième moment sont ceux déjà dérivés du lecteur dodécaphasique. La racine carrée est la racine positive du corps ordonné. La série convergente réalise analytiquement ce lecteur ; elle conserve le profil et le coefficient normalisateur de la section~\ref{hmtfg:escalado:1}.

\textbf{Théorème de transport du retour.} Pour tout paramètre du domaine interne de \eqref{hmtfg:escalado:eq:38},
\begin{equation}
\iota(F_a(x))=f_{\iota(a)}(\iota(x)),\qquad
\iota(S_n^{\rm H}(a))=S_n(\iota(a)).
\label{hmtfg:escalado:eq:39}
\end{equation}
\textbf{Démonstration.} L’isomorphisme ordonné fixe les rationnels et préserve la racine positive. Il préserve les limites convergentes : chaque voisinage de rayon rationnel positif est transporté vers un autre de même rayon. Appliqué aux sommes partielles de \eqref{hmtfg:escalado:eq:38}, il transporte le cosinus et \(u_{\rm H}\). La première égalité résulte de la conservation de la somme et du produit ; la seconde, d’une récurrence sur le nombre d’itérations.

Un centre antérieur \(a_{n-1}\) étant fixé, définissons \(Z_n^{\rm H}\) par un retour nul, une période critique exacte \(2^n\) et \(a>a_{n-1}\), dans le domaine paramétrique déclaré. Soit \(Z_n\) l’ensemble défini par les mêmes conditions pour \(f_\lambda\), à droite de \(\iota(a_{n-1})\). Alors
\begin{equation}
\boxed{\iota[Z_n^{\rm H}]=Z_n,\qquad
\iota(\min Z_n^{\rm H})=\min Z_n.}
\label{hmtfg:escalado:eq:40}
\end{equation}
La seconde égalité est affirmée lorsque le minimum existe ; la première préserve également l’absence de candidats.

\textbf{Démonstration.} \eqref{hmtfg:escalado:eq:39} préserve le retour nul et les retours antérieurs qui déterminent la période exacte. L’ordre conserve la position par rapport au centre précédent. La surjectivité de \eqref{hmtfg:escalado:eq:37}, appliquée à tout paramètre racine, prouve la couverture inverse. Si \(a_*\) est le minimum, tout \(z\in Z_n\) est de la forme \(\iota(a)\), avec \(a\in Z_n^{\rm H}\), d’où \(\iota(a_*)\le z\). L’affirmation réciproque utilise \(\iota^{-1}\).

La surjectivité couvre tous les paramètres de la carte, y compris ceux dont les racines n’ont pas été isolées par un calcul fini. L’isomorphisme conserve l’univers d’interprétation déclaré et l’ordre de sélection.

\subsubsection{Représentation de Jacobi du profil dodécaphasique complet}
\label{hmtfg:escalado:14.2}

Définissons la mesure atomique du lecteur,
\begin{equation}
s_m=\frac{4(3m-1)^2}{899},\qquad
\nu_0=\frac1{12}\sum_{m=1}^{12}\delta_{s_m},\qquad
M_r=\int s^r\,d\nu_0(s).
\label{hmtfg:escalado:eq:41}
\end{equation}
Les douze nœuds sont distincts et positifs. Pour un polynôme \(p\ne0\) de degré inférieur à \(r\le12\),
\begin{equation}
\sum_{i,j=0}^{r-1}p_iM_{i+j}p_j
=\frac1{12}\sum_{m=1}^{12}p(s_m)^2>0.
\label{hmtfg:escalado:eq:42}
\end{equation}
Un polynôme de degré inférieur à douze ne s’annule pas aux douze nœuds. Au degré douze apparaît \(\prod_m(s-s_m)\), de norme nulle.

L’orthogonalisation par rapport à cette mesure produit des polynômes orthonormaux \(\psi_0,\ldots,\psi_{11}\), avec \(\psi_0=1\). Soient
\[
Q_{mj}=\psi_j(s_m)/\sqrt{12},\quad
D=\operatorname{diag}(s_m),\quad J_{12}=Q^{\mathsf T}DQ.
\]
Les relations suivantes sont satisfaites
\begin{equation}
Q^{\mathsf T}Q=I,\quad DQ=QJ_{12},\qquad
\boxed{u_0(x)=e_0^{\mathsf T}
\bigl[I-\cos(x\sqrt{J_{12}})\bigr]e_0.}
\label{hmtfg:escalado:eq:43}
\end{equation}
\textbf{Démonstration.} L’orthogonalité découle de \eqref{hmtfg:escalado:eq:42}. La multiplication par \(s\) produit une récurrence à trois termes : les éléments séparés de plus d’une position s’annulent par orthogonalité et degré. Fixer des coefficients dominants positifs rend les sous-diagonales positives. Par calcul fonctionnel,
\(\cos(x\sqrt{J_{12}})=Q^{\mathsf T}\cos(x\sqrt D)Q\).
Le vecteur \(Qe_0\) est constant, de valeur \(1/\sqrt{12}\) ; sa substitution restitue la somme de la section~\ref{hmtfg:escalado:1}.

L’article sur les grassmanniennes développe ce mécanisme pour les mesures marquées et les raffinements. Il est appliqué ici à \eqref{hmtfg:escalado:eq:41} ; le rang douze appartient à ce lecteur. Les marques généalogiques sont conservées avec le lecteur, conformément à la fidélité marquée du document de référence. \(J_{12}\) représente le profil analytique ; l’état complet conserve en outre les routes, l’orientation et la mémoire.

Le vérificateur de cette extension contrôle, avec des fractions exactes, les douze degrés de norme positive, l’annulation au degré douze, l’entrelaceur \(VT=DV\) dans la base monique et l’autoadjonction pondérée \(HT=T^{\mathsf T}H\). Il vérifie les moments de degré zéro à trente ; l’identité à tout degré découle de l’entrelaceur. Ses premiers coefficients sont
\begin{equation}
a_0=2,\qquad \beta_1=\frac{2493348}{808201}.
\label{hmtfg:escalado:eq:44}
\end{equation}

\subsubsection{Sensibilité orientée du retour}
\label{hmtfg:escalado:14.3}

Pour un retour exact de période \(p=2^n\), soient
\[
c_j=f_\lambda^j(0),\quad q_j=\partial_\lambda c_j,\quad
D_j=\prod_{k=1}^j f_\lambda'(c_k),\quad D_0=1.
\]
Avant le retour critique, les facteurs de \(D_{p-1}\) sont non nuls. La différentiation et le télescopage donnent
\begin{equation}
q_{j+1}=-u_0(c_j)+f_\lambda'(c_j)q_j,\qquad
\boxed{\frac{q_p}{D_{p-1}}
=-\sum_{j=1}^{p-1}\frac{u_0(c_j)}{D_j}.}
\label{hmtfg:escalado:eq:45}
\end{equation}
Dans le mot canonique, \(\operatorname{sign}c_j=(-1)^{v_2(j)}\). Le signe de \(f_\lambda'(c_j)\) est opposé ; par conséquent,
\begin{equation}
\operatorname{sign}D_j=(-1)^{j+\sum_{k=1}^jv_2(k)}
=(-1)^{s_2(j)},
\label{hmtfg:escalado:eq:46}
\end{equation}
en utilisant \(\sum_{k=1}^jv_2(k)=j-s_2(j)\), où \(s_2(j)\) compte les uns binaires. La sensibilité normalisée est
\begin{equation}
\mathcal T_n=-\sum_{j=1}^{2^n-1}(-1)^{s_2(j)}t_j,\qquad
t_j=\frac{u_0(c_j)}{|D_j|}>0.
\label{hmtfg:escalado:eq:47}
\end{equation}
La positivité nodale de \eqref{hmtfg:escalado:eq:42} détermine \(u_0\) ; \eqref{hmtfg:escalado:eq:47} conserve en outre les produits orbitaux et leurs orientations.

La représentation supplémentaire par moments exigerait une mesure positive \(\rho\) sur \([0,1]\) avec \(t_j=\int z^{j-1}\,d\rho\). Sous cette condition,
\begin{equation}
\mathcal T_n=\int_0^1
\frac{1-\prod_{r=0}^{n-1}(1-z^{2^r})}{z}\,d\rho(z)>0
\label{hmtfg:escalado:eq:48}
\end{equation}
pour \(\rho\ne0\), en prolongeant l’intégrande par sa valeur \(1\) en zéro. L’identité s’obtient en développant le produit binaire.

Le contrôle rationnel à la racine de période quatre précédemment certifiée obtient
\[
t_1\approx0.40425224267600139730,\quad
t_2\approx0.04925249167432646403,\quad
t_3\approx0.15740870098294833737
\]
et l’intervalle extérieur rigoureux
\begin{equation}
\begin{gathered}
0.296096033367379523967764444238832345360107<\mathcal T_2,\\
\mathcal T_2<0.296096033367379523967764444238832345360112.
\end{gathered}
\label{hmtfg:escalado:eq:49}
\end{equation}
Il certifie également \(t_3>t_2\). Une mesure positive sur \([0,1]\) imposerait \(t_3\le t_2\) ; ce relèvement précis des poids orbitaux en moments positifs est exclu. Le signe positif de \eqref{hmtfg:escalado:eq:49} reste certifié. Ce contrôle distingue l’échec d’une preuve proposée de l’échec de la propriété que l’on cherchait à prouver.

% Fragmento integrable. Preambulo anfitrion: amsmath, amssymb, longtable, hyperref.
\section{Classification et prolongement du sélecteur des premières racines}
\label{hmtfg:fragmento:clasificacion-y-continuacion-del-selector-de-primeras-raices}

\subsection{Objet et provenance}
\label{hmtfg:clasificacion:1}

Cette annexe conserve la famille issue de la lecture dodécaphasique HMT :

\[
u_0(x)=\frac1{12}\sum_{m=1}^{12}\left(1-\cos\frac{2(3m-1)x}{\sqrt{899}}\right),
\qquad f_\lambda(x)=1-\lambda u_0(x).
\]

La famille provient des antécédents operatoriels et analytiques exposés dans ce développement. L’existence et la compacité de ses réalisations de l’itinéraire infini ont été démontrées dans les sous-sections consacrées à la réalisation du mot et à la conservation de ses signes.

L’argument suivant est propre à ce développement et utilise exclusivement ces propriétés et les opérations de retour de la famille fixée. Il n’utilise comme entrée ni un paramètre d’une autre famille ni la valeur d’une constante universelle. Les feuilles des itinéraires utilisés ici sont les deux branches monotones du lecteur critique ; les autres coordonnées et la mémoire HMT conservent leur résidence dans l’état dont provient le lecteur.

Écrivons
\[
c_j(\lambda)=f_\lambda^j(0),\qquad
\tau_j=(-1)^{v_2(j)},\qquad
W_n=\tau_1\cdots\tau_{2^n-1}.
\]
L’itinéraire d’une orbite critique qui revient pour la première fois à l’origine se termine par le symbole zéro. L’ordre des itinéraires est l’ordre unimodal à maximum : avant la comparaison du symbole d’indice \(j\), le facteur d’orientation est
\[
O_{j-1}(K)=\prod_{i=1}^{j-1}(-K_i).
\]
L’ordre ordinaire des symboles est \(-1<0<1\).


\subsection{Classification avant le mot infini}
\label{hmtfg:clasificacion:2}

\textbf{Théorème de classification.} Soit \(F:I\to I\), avec \(I=[\ell,1]\), \(\ell<0\), une application continue, strictement croissante à gauche de zéro et strictement décroissante à sa droite, dont le maximum satisfait \(F(0)=1\). Supposons que zéro soit de période exacte \(2^n\), avec \(n\ge1\), et que son itinéraire critique \(K\) satisfasse \(K<\tau\). Alors
\[
\boxed{K=W_n0.}
\]
La conclusion porte sur l’itinéraire et permet que des paramètres distincts d’une famille possèdent ce même itinéraire.

\textbf{Démonstration.} Notons \(d_j=F^j(0)\). Pour la période deux, l’itinéraire est \(+0=W_10\). Considérons une période exacte \(p=2^n\ge4\).

On a \(d_2<0\). En effet, \(d_2=0\) donnerait la période deux. Si \(d_2>0\), la monotonie de la branche positive donne
\[
F([d_2,1])=[d_2,d_3]\subset[d_2,1],
\]
et l’orbite critique demeure dans un intervalle strictement positif, incompatible avec le retour à zéro.

Les signes initiaux sont alors \(+,-\). Le facteur d’orientation pour comparer le troisième signe est \(-1\) ; par conséquent, \(K<\tau\) impose \(d_3>0\). Un zéro à cette position produirait la période trois.

On ne peut pas non plus avoir \(d_4<0\). Sous cette hypothèse,
\[
d_2<d_4<0,
\qquad F([d_2,d_4])=[d_3,d_5]\subset(0,1],
\]
car \(d_4=F(d_3)>F(1)=d_2\) et \(d_5=F(d_4)>F(d_2)=d_3\). La branche positive est décroissante, de sorte que
\[
F^2([d_2,d_4])=[d_6,d_4]\subset[d_2,d_4].
\]
Ici \(d_6\ge d_2\), car \(d_5\le1\), et \(d_6<d_4\), car \(d_5>d_3\). Cet intervalle strictement négatif piège les retours pairs de l’orbite critique et exclut le retour à zéro. Donc \(d_4\ge0\). Si \(p=4\), on obtient précisément \(+-+0=W_20\).

Supposons maintenant \(p\ge8\). Alors \(d_4>0\). Les facteurs d’orientation pour les cinquième et sixième signes sont, respectivement, \(-1\) et \(+1\). Comme \(\tau_5=+1\) et \(\tau_6=-1\), l’inégalité \(K<\tau\) impose
\begin{equation}
\operatorname{sign}(d_1,\ldots,d_6)=(+,-,+,+,+,-).
\label{hmtfg:clasificacion:eq:1}
\end{equation}
Les signes nuls à ces positions sont exclus par la période exacte.

Nous construisons l’intervalle de retour
\[
J=[d_2,d_4].
\]
L’inégalité \(F(d_3)=d_4>0>d_6=F(d_5)\), jointe à \(d_3,d_5>0\), implique \(d_3<d_5\) ; donc
\[
F(J)=[d_3,1].
\]
De plus, \(d_4<d_3\). Pour le vérifier, l’image par \(F\) de l’intervalle d’extrémités \(d_3,d_4\) est strictement positive, puisque les images de ses extrémités sont \(d_4,d_5>0\). Sur cet intervalle, \(F^2\) est strictement croissante. Ses valeurs aux extrémités satisfont
\[
F^2(d_3)=d_5>0>d_6=F^2(d_4),
\]
ce qui impose \(d_3>d_4\). Par conséquent,
\begin{equation}
J\cap F(J)=\varnothing,
\qquad F^2(J)=F([d_3,1])=[d_2,d_4]=J.
\label{hmtfg:clasificacion:eq:2}
\end{equation}

L’application \(F^2|_J\) possède un unique minimum en zéro. La conjugaison par \(h(x)=d_2x\), qui inverse l’orientation, définit
\begin{equation}
G=h^{-1}\circ F^2\circ h:
\left[\frac{d_4}{d_2},1\right]\longrightarrow
\left[\frac{d_4}{d_2},1\right].
\label{hmtfg:clasificacion:eq:3}
\end{equation}
Cette application satisfait les mêmes hypothèses de maximum que \(F\), avec \(G(0)=1\), et son point critique a pour période exacte \(p/2\). Son itinéraire est
\begin{equation}
K'_j=-K_{2j}.
\label{hmtfg:clasificacion:eq:4}
\end{equation}
Tous les symboles d’indice impair de \(K\) sont positifs, car les positions paires appartiennent à \(J\) et les positions impaires à \(F(J)\subset(0,1]\).

La transformation \eqref{hmtfg:clasificacion:eq:4} conserve l’ordre de comparaison avec \(\tau\). En effet, \(\tau_{2j-1}=+1\) et \(\tau_{2j}=-\tau_j\). Si \(q\) est la première position où \(K'\) et \(\tau\) diffèrent, la première position de différence entre \(K\) et \(\tau\) est \(2q\), et
\begin{equation}
O_{2q-1}(K)=-O_{q-1}(K'),
\qquad K_{2q}-\tau_{2q}=-(K'_q-\tau_q).
\label{hmtfg:clasificacion:eq:5}
\end{equation}
Le produit qui détermine l’ordre est identique. Les identités restent valides lorsque le premier symbole distinct est le zéro final. Ainsi \(K'<\tau\).

La récurrence appliquée à \(G\) donne \(K'=W_{n-1}0\). Les positions impaires positives et \eqref{hmtfg:clasificacion:eq:4} reconstruisent exactement \(K=W_n0\). Le théorème est démontré. \(\square\)

Le théorème démontre également que chaque application classifiée possède les retours restrictifs de doublement antérieurs à son retour superstable, avec les domaines de \eqref{hmtfg:clasificacion:eq:2}–\eqref{hmtfg:clasificacion:eq:3}. Il ne présuppose pas de monotonie du paramètre d’une famille.

\textbf{Normalisation symétrique de la famille fixée.} Lorsque \(F\) est paire et définie sur \([-1,1]\), les inégalités précédentes donnent en outre
\[
F(d_4)=d_5>d_3=F(d_2)=F(-d_2).
\]
Comme les deux arguments \(d_4,-d_2\) sont positifs et que \(F\) décroît sur cette branche,
\begin{equation}
0<d_4<-d_2<1.
\label{hmtfg:clasificacion:eq:5a}
\end{equation}
Donc \(F^2([d_2,-d_2])=[d_2,d_4]\), et la même conjugaison \(h(x)=d_2x\) définit sur tout \([-1,1]\) une application paire à maximum dont l’image est \([d_4/d_2,1]\subset[-1,1]\). Elle coïncide exactement avec l’opérateur \(RF=-F(F(-ax))/a\), \(a=-F(1)\), utilisé dans le corpus. Cette extension symétrique conserve le retour critique, sa période et son itinéraire.


\subsection{Premier paramètre d’itinéraire infini}
\label{hmtfg:clasificacion:3}

Nous travaillons sur l’intervalle compact déjà utilisé dans la démonstration d’existence,
\[
L=\left[\frac1{2u_0(1)},\frac2{u_0(1)}\right].
\]
Soit
\[
\Lambda_\tau=\{\lambda\in L:
\operatorname{sign}c_j(\lambda)=\tau_j\text{ pour tout }j\ge1\}.
\]
D’après les démonstrations d’existence et de séparation des signes du prolongement, cet ensemble est non vide et compact. Il existe donc
\begin{equation}
a=\min\Lambda_\tau.
\label{hmtfg:clasificacion:eq:6}
\end{equation}

\textbf{Lemme de comparaison.} Pour tout \(\lambda\in[1/(2u_0(1)),a)\), on a \(K_\lambda<\tau\).

\textbf{Démonstration.} À l’extrémité gauche, l’itinéraire est \(+^\infty<\tau\). L’argument de séparation de la réalisation du mot infini démontre que les ensembles de comparaison stricte avec \(\tau\) sont ouverts : dans un itinéraire infini distinct de \(\tau\), la première différence finie persiste par continuité ; à un paramètre superstable, les deux mots périodiques voisins et le mot à zéro final se situent du même côté de \(\tau\), par sa maximalité stricte. L’intervalle antérieur à \(a\) ne contient, par définition, aucun paramètre d’itinéraire \(\tau\). Sa connexité maintient la comparaison initiale sur tout cet intervalle. \(\square\)


\subsection{Existence de chaque première nouvelle racine}
\label{hmtfg:clasificacion:4}

Nous conservons le sélecteur original. La première racine est
\[
\lambda_1=\frac1{u_0(1)},
\]
et, pour \(n\ge2\), \(\lambda_n\) est la première racine de \(c_{2^n}\) située à droite de \(\lambda_{n-1}\) qui satisfait
\begin{equation}
c_{2^j}(\lambda_n)\ne0\quad(1\le j<n).
\label{hmtfg:clasificacion:eq:7}
\end{equation}
Comme toute première période divisant \(2^n\) est une puissance de deux, \eqref{hmtfg:clasificacion:eq:7} équivaut à exiger une période critique exacte \(2^n\).

\textbf{Lemme d’existence.} Toutes les racines de ce sélecteur existent et satisfont
\begin{equation}
\lambda_1<\lambda_2<\cdots<\lambda_n<a.
\label{hmtfg:clasificacion:eq:8}
\end{equation}

\textbf{Démonstration.} Le signe \(c_2(a)<0\) implique \(\lambda_1<a\). Supposons \(\lambda_{n-1}<a\) construite, et posons \(p=2^{n-1}\). La fonction analytique \(c_p(\lambda)\) n’est pas identiquement nulle, puisque \(c_p(a)\ne0\). Son ensemble de zéros dans \([\lambda_{n-1},a]\) est fini, non vide et ne contient pas \(a\). Soit \(z<a\) son dernier zéro.

Sur \((z,a]\), le signe de \(c_p\) est constant et égal à \(s=\tau_p\). Pour \(g_\lambda=f_\lambda^p\), on a
\[
g_\lambda(0)=c_p(\lambda),\qquad g'_\lambda(0)=0.
\]
Une borne locale uniforme de la dérivée seconde donne
\begin{equation}
c_{2p}(\lambda)
=g_\lambda(c_p(\lambda))
=c_p(\lambda)+O(c_p(\lambda)^2).
\label{hmtfg:clasificacion:eq:9}
\end{equation}
Lorsque \(\lambda\) tend vers \(z\) par la droite, le terme du premier ordre domine ; donc \(c_{2p}(\lambda)\) a pour signe \(s\). En \(a\), son signe est
\[
\tau_{2p}=-\tau_p=-s.
\]
Le théorème des valeurs intermédiaires fournit un zéro \(\beta\in(z,a)\) de \(c_{2p}\). Comme \(c_p\ne0\) sur cet intervalle, la période critique en \(\beta\) est exactement \(2p\).

L’analyticité et \(c_{2p}(a)\ne0\) impliquent que les zéros de \(c_{2p}\) dans \([\lambda_{n-1},a]\) sont en nombre fini. Parmi ses zéros de période exacte \(2p\) situés à droite de \(\lambda_{n-1}\), il en existe donc un premier. C’est le sélecteur original \(\lambda_n\), et il satisfait \(\lambda_n\le\beta<a\). \(\square\)

L’utilisation du dernier zéro auxiliaire \(z\) appartient uniquement à la preuve d’existence. Le paramètre sélectionné reste la \textbf{première nouvelle racine} \(\lambda_n\).


\subsection{Convergence du sélecteur global vers le premier itinéraire infini}
\label{hmtfg:clasificacion:5}

\textbf{Théorème de prolongement global.} Pour la famille HMT fixée et le sélecteur original, tous les itinéraires sélectionnés sont canoniques et
\begin{equation}
\boxed{K_{\lambda_n}=W_n0,\qquad
\lambda_n\nearrow\min\Lambda_\tau.}
\label{hmtfg:clasificacion:eq:10}
\end{equation}

\textbf{Démonstration.} Les lemmes de comparaison et d’existence placent chaque paramètre sélectionné au-dessous de \(a\), avec un itinéraire inférieur à \(\tau\). Le théorème de classification identifie cet itinéraire à \(W_n0\). La suite est croissante et majorée par \(a\) ; soit \(b\le a\) sa limite.

Fixons \(p\ge1\). Pour \(n\) suffisamment grand, \(2p<2^n\). Alors les signes de \(c_p(\lambda_n)\) et de \(c_{2p}(\lambda_n)\) coïncident avec \(\tau_p\) et \(\tau_{2p}=-\tau_p\), respectivement. La borne uniforme des dérivées secondes,
\[
\|(f_\lambda^p)''\|\le B_p,
\qquad B_p=16^p\frac{16^p-1}{15},
\]
et l’identité \((f_\lambda^p)'(0)=0\) impliquent
\[
|c_{2p}(\lambda_n)-c_p(\lambda_n)|
\le\tfrac12B_p|c_p(\lambda_n)|^2.
\]
Comme les deux termes du membre de gauche sont de signes opposés,
\begin{equation}
|c_p(\lambda_n)|>2/B_p.
\label{hmtfg:clasificacion:eq:11}
\end{equation}
Par continuité, \(c_p(b)\) conserve le signe \(\tau_p\) et satisfait \(|c_p(b)|\ge2/B_p>0\). Cela vaut pour chaque \(p\), de sorte que \(b\in\Lambda_\tau\). La minimalité de \(a\) donne \(b\ge a\), et donc \(b=a\). \(\square\)

La preuve complète conserve le sélecteur global et produit ses itinéraires et sa limite. L’information des huit premiers niveaux est compatible avec le résultat, mais la démonstration de \eqref{hmtfg:clasificacion:eq:10} n’extrapole pas ces huit niveaux : elle utilise la récurrence, des intervalles invariants et la séparation uniforme de chaque signe.


\subsection{Contrôles logiques et falsificateurs}
\label{hmtfg:clasificacion:7}

\par\noindent\textbf{1.}\enspace Le théorème de classification exige un maximum unimodal strict et un intervalle invariant. Si l’une des deux branches perd sa monotonie, les intervalles \eqref{hmtfg:clasificacion:eq:2} peuvent cesser d’être restrictifs et la classification doit être vérifiée à nouveau.
\par\noindent\textbf{2.}\enspace La condition \(K<\tau\) est essentielle. Au-delà du mot limite, il peut exister des centres de période \(2^n\) avec d’autres combinatoires ; le théorème ne les identifie pas à \(W_n0\).
\par\noindent\textbf{3.}\enspace L’analyticité en le paramètre assure l’isolement et la finitude des zéros sur le compact. Pour une famille seulement continue, l’existence d’une première racine strictement postérieure exige un argument supplémentaire.
\par\noindent\textbf{4.}\enspace Une limite de racines critiques peut, en général, perdre des signes en atteignant un retour de période inférieure. L’inégalité \eqref{hmtfg:clasificacion:eq:11} exclut précisément ce phénomène pour cette suite et chaque horizon fixé.
\par\noindent\textbf{5.}\enspace La convergence des premières racines vers le premier paramètre d’itinéraire \(\tau\) n’équivaut pas à l’unicité de tous les paramètres de cet itinéraire et ne détermine pas à elle seule un taux métrique. Le croisement transverse et l’identification de la branche apportent ces conclusions ultérieures.

% Fragmento integrable. Preambulo anfitrion: amsmath, amssymb, longtable, hyperref.
\section{Entrée transverse et changement d’échelle local}
\label{hmtfg:fragmento:entrada-transversal-y-escalado-local}

\subsection{Provenance et résultat}
\label{hmtfg:escalado:1}

APP évalue la somme et le produit sur les deux feuilles du réseau, en conservant quotient et reste. TRIT détermine le régime et l’orientation. Le transport TPK sélectionne, transporte et actualise l’état avec ses registres de retenue, de frontière et de mémoire. Ses lecteurs agissent sur la même structure discrète du continu. Le prolongement \texttt{w6→w12→w18→w24→w30→R36→G9} conserve cette provenance et distingue le retour de phase de l’incrément de mémoire.

Le lecteur critique de la roue dodécaphasique, construit dans les antécédents operatoriels et transporté par le raffinement compatible, produit
\[
\phi_j=\frac{(3j-1)\pi_{\mathrm{HMT}}}{18},\qquad
w_j=\frac1{12},\qquad
k_j=\sqrt{\frac2{\sum_{\ell=1}^{12}w_\ell\phi_\ell^2}}\,\phi_j
=\frac{2(3j-1)}{\sqrt{899}}.
\]
La coordonnée \(\pi_{\mathrm{HMT}}\) conserve son origine APP–TRIT–TPK ; son facteur commun s’annule dans ce deuxième moment. Le lecteur résultant est
\begin{equation}
u_0(x)=\frac1{12}\sum_{j=1}^{12}(1-\cos k_jx),\qquad
f_\lambda(x)=1-\lambda u_0(x).
\label{hmtfg:escalado:eq:1}
\end{equation}
Le secteur focal est le secteur uniforme, \(\eta=0\). L’extension antérieure conserve séparément la famille pondérée \(|\eta|\le1/40\). Les nouvelles bornes numériques correspondent à \eqref{hmtfg:escalado:eq:1}.

La famille est paire, entière, avec \(f_\lambda(0)=1\), un point critique quadratique et un deuxième moment \(\sum_jw_jk_j^2=2\). Le paramètre \(\lambda\) parcourt une famille postérieure au générateur HMT. La constante cible de Feigenbaum reste extérieure aux entrées de \eqref{hmtfg:escalado:eq:1}.

\textbf{Résultat principal.} Il existe un unique
\begin{equation}
\lambda_*\in[1.874038,1.874039]
\label{hmtfg:escalado:eq:2}
\end{equation}
pour lequel le quatrième retour normalisé de \eqref{hmtfg:escalado:eq:1} appartient à la feuille stable locale du point fixe réel de doublement. Le croisement de la famille avec cette feuille est transverse, avec une séparation quantifiée. Pour tout \(n\ge0\),
\begin{equation}
\boxed{\|R^{4+n}f_{\lambda_*}-g\|_{\mathcal A}
<0.001666\,(0.83996)^n.}
\label{hmtfg:escalado:eq:3}
\end{equation}
La cascade superstable locale associée à ce croisement possède des paramètres \(\mu_j\) et des constantes \(C\ne0\), \(0<\theta<1\), tels que
\begin{equation}
\mu_j-\lambda_*=C\delta_{\mathrm F}^{-j}\bigl[1+O(\theta^j)\bigr],
\qquad
\frac{\mu_{j-1}-\mu_{j-2}}{\mu_j-\mu_{j-1}}\longrightarrow\delta_{\mathrm F}.
\label{hmtfg:escalado:eq:4}
\end{equation}
Ici \(\delta_{\mathrm F}>1\) est la valeur propre instable de l’opérateur de doublement. La sélection de \(\mu_j\) s’effectue par les retours locaux de la section~\ref{hmtfg:escalado:8}. L’identification avec le sélecteur global est démontrée dans les sections~\ref{hmtfg:escalado:15}--\ref{hmtfg:escalado:17} : tous les termes globaux conservent la combinatoire et les centres coïncident à tout niveau suffisamment élevé lorsque leurs périodes sont égalisées.


\subsection{Carte analytique et retour effectif}
\label{hmtfg:escalado:2}

On utilise la carte de fonctions paires
\[
s=x^2,\qquad z=\frac{s-1}{5/2},\qquad
f(x)=1-x^2V(z),\qquad V(z)=\sum_{j\ge0}v_jz^j.
\]
Les coordonnées de l’espace de Banach sont
\begin{equation}
u=10v_0,\qquad \nu=(v_1,v_2,\ldots),\qquad
\|(\Delta u,\Delta\nu)\|_{\mathcal A}
=|\Delta u|+\sum_{j\ge1}|\Delta v_j|.
\label{hmtfg:escalado:eq:5}
\end{equation}
En particulier, le coefficient constant de \(V\) a un poids de dix. Cette normalisation coïncide avec la carte \(u/10\) de Hertling–Spandl, section 2 de la source.

Soit \(a=-f(1)=v_0-1\). Pour les retours réels admissibles,
\begin{equation}
(Rf)(x)=-\frac{f(f(-ax))}{a}.
\label{hmtfg:escalado:eq:6}
\end{equation}
Le programme évalue une factorisation exacte de \eqref{hmtfg:escalado:eq:6}. En définissant
\[
S=\frac{a^2s-1}{5/2},\qquad b=V(S),\qquad
h=1-a^2sb,\qquad t=\frac{h^2-1}{5/2},
\]
\[
\operatorname{shift}V(z)=\sum_{j\ge0}v_{j+1}z^j,
\]
on obtient
\begin{equation}
\boxed{V_R=a\,b\,(h+1)\left[V(t)+\frac25\operatorname{shift}V(t)\right].}
\label{hmtfg:escalado:eq:7}
\end{equation}
Pour la vérifier, on utilise \(v_0=1+a\),
\(h^2-1=-a^2sb(h+1)\) et
\(V(t)-v_0=t\operatorname{shift}V(t)\). Ainsi,
\(a+1-h^2V(t)=a^2sb(h+1)[V(t)+(2/5)\operatorname{shift}V(t)]\),
ce qui est précisément \(asV_R\).

Les moments initiaux sont rationnels :
\begin{equation}
\frac1{12}\sum_jk_j^{2m}
=\frac1{12}\left(\frac4{899}\right)^m
\sum_{j=1}^{12}(3j-1)^{2m}.
\label{hmtfg:escalado:eq:8}
\end{equation}
Ainsi, tant \eqref{hmtfg:escalado:eq:7} que la construction initiale utilisent des intervalles rationnels ; les fonctions trigonométriques initiales sont représentées par leurs séries avec une borne extérieure du reste.


\subsection{Bornes externes utilisées comme reconnaissance ultérieure}
\label{hmtfg:escalado:3}

Une fois \eqref{hmtfg:escalado:eq:1} construite, le théorème d’hyperbolicité de Lanford est utilisé comme certificat analytique de l’opérateur \eqref{hmtfg:escalado:eq:6}. Soient \(g_0\) le polynôme central publié et \(g\) le point fixe. Les estimations utilisées sont
\[
\|g-g_0\|_{\mathcal A}<\varepsilon_0=4\,10^{-5},
\]
et, dans la boule \(\|f-g_0\|_{\mathcal A}<0.01\), les blocs de \(DR\), notés \(R=(U,V)\), satisfont
\begin{equation}
U_u\ge a_0=4.521,\quad
\|U_\nu\|\le b_0=0.560,\quad
\|V_u\|\le c_0=0.756,\quad
\|V_\nu\|\le d_0=0.719.
\label{hmtfg:escalado:eq:9}
\end{equation}
Le point fixe possède une unique direction instable, de valeur propre réelle positive \(\delta_{\mathrm F}>1\) ; le reste du spectre se situe à l’intérieur du disque unité. Les chiffres de \eqref{hmtfg:escalado:eq:9} proviennent des estimations publiées, et non d’une nouvelle reproduction de la preuve de Lanford dans ce dossier. Les coefficients de \(g_0\) sont explicitement incorporés au polynôme central de référence du programme, conformément au tableau 2 de Hertling–Spandl.

Sources : \href{https://repo-archives.ihes.fr/A_GARDER/20180409/Test/I_Prepublications/LANFORD/1968-2013/P_81_17/P_81_17_web.pdf}{Lanford, estimations d’hyperbolicité}; \href{https://arxiv.org/pdf/1410.3277}{Hertling–Spandl, carte analytique et tableau 2}.


\subsection{Certificat rationnel d’entrée et de direction}
\label{hmtfg:escalado:4}

Le certificat d’entrée dans la renormalisation divise \eqref{hmtfg:escalado:eq:2} en seize intervalles rationnels égaux. Il propage quatre retours de \eqref{hmtfg:escalado:eq:7} et la dérivée exacte par rapport à \(\lambda\), en utilisant la différentiation de la composition et l’arrondi extérieur à cent décimales.

Chaque série se compose d’un polynôme de degré 32 à coefficients intervalles et d’une erreur analytique arbitraire \(E\), \(\|E\|_1\le\epsilon\). Cette erreur peut également modifier les coefficients de bas degré. Pour \(\|q\|_1<1\),
\begin{equation}
\|E\circ q\|_1\le\epsilon,\qquad
\|E'\circ q\|_1\le\frac{\epsilon}{(1-\|q\|_1)^2},\qquad
\|\operatorname{shift}E\|_1\le\epsilon.
\label{hmtfg:escalado:eq:10}
\end{equation}
Les produits comprennent les termes polynôme–erreur et erreur–erreur. La queue initiale est majorée par son premier terme en valeur absolue et une raison géométrique, en utilisant \(\|1+(5/2)z\|_1=7/2\). Chaque argument composé reste strictement à l’intérieur de la boule unité de cette algèbre.

Pour \(\Gamma(\lambda)=R^4f_\lambda=(u_\lambda,\nu_\lambda)\), le certificat donne les bornes conservatrices suivantes, uniformes en \eqref{hmtfg:escalado:eq:2} :

\begingroup\small
\begin{longtable}{p{0.465\linewidth}p{0.465\linewidth}}
\hline
Quantité & Borne extérieure certifiée \\
\hline
\(\|\Gamma(\lambda)-g_0\|_{\mathcal A}\) & \(<0.004993128\) \\
\(\|\nu_\lambda-\nu_{g_0}\|_1\) & \(<0.001396077\) \\
\(u'_\lambda\) & \(>3821.289115\) \\
\(\|\nu'_\lambda\|_1\) & \(<551.757000\) \\
\(\|\nu'_\lambda\|_1/u'_\lambda\) & \(<0.144386<1/4\) \\
Erreur analytique de la fonction & \(<8.482\,10^{-19}\) \\
Erreur analytique de la tangente & \(<4.569\,10^{-15}\) \\
\hline\end{longtable}
\endgroup

Les décimales complètes, les seize sous-intervalles et chaque retour intermédiaire sont conservés dans le certificat rationnel d’entrée.

Le cône \(\|\dot\nu\|_1\le|\dot u|/4\) est invariant par \eqref{hmtfg:escalado:eq:9} :
\[
|\dot u_{\rm nuevo}|\ge(4.521-0.560/4)|\dot u|
=4.381|\dot u|,
\]
\begin{equation}
\|\dot\nu_{\rm nueva}\|_1\le(0.756+0.719/4)|\dot u|
=0.93575|\dot u|<\tfrac14(4.381)|\dot u|.
\label{hmtfg:escalado:eq:11}
\end{equation}
La direction paramétrique certifiée entre donc dans un cône uniformément expansif de l’opérateur.


\subsection{Construction quantitative de la feuille stable}
\label{hmtfg:escalado:5}

Translatons \(g\) à l’origine et écrivons les coordonnées sous la forme \((x,y)=(u-u_g,\nu-\nu_g)\). Fixons
\begin{equation}
L=0.16,\qquad r=0.004,\qquad
A=a_0-Lc_0=4.40004,\qquad q=d_0+c_0L=0.83996.
\label{hmtfg:escalado:eq:12}
\end{equation}
Le cylindre \(|x|\le Lr\), \(\|y\|_1\le r\) se trouve à l’intérieur de la boule de \eqref{hmtfg:escalado:eq:9}, puisque
\((1+L)r+\varepsilon_0=0.00468<0.01\).

Considérons l’espace complet des fonctions \(\sigma:B_r\to\mathbb R\) avec \(\sigma(0)=0\) et de constante de Lipschitz au plus \(L\), muni de la distance uniforme. Pour chaque \(y\), on définit \((\mathcal G\sigma)(y)\) comme la racine dans \([-L\|y\|_1,L\|y\|_1]\) de
\begin{equation}
F_\sigma(x,y)=U(x,y)-\sigma(V(x,y)).
\label{hmtfg:escalado:eq:13}
\end{equation}
Sur cet intervalle, \(\|V(x,y)\|_1\le q\|y\|_1<r\). Lorsque \(x\) augmente, \eqref{hmtfg:escalado:eq:9} montre que \(F_\sigma\) croît avec une pente au moins égale à \(A\). À ses extrémités,
\[
F_\sigma(L\|y\|_1,y)
\ge[AL-b_0-Ld_0]\|y\|_1,
\]
et l’inégalité opposée vaut à l’extrémité négative. La marge est
\begin{equation}
AL-b_0-Ld_0=0.0289664>0.
\label{hmtfg:escalado:eq:14}
\end{equation}
On obtient une racine unique. En comparant \eqref{hmtfg:escalado:eq:13} pour deux valeurs de \(y\),
\[
\operatorname{Lip}(\mathcal G\sigma)
\le\frac{b_0+Ld_0}{A}
=\frac{0.67504}{4.40004}<0.16.
\]
En comparant deux graphes,
\begin{equation}
\|\mathcal G\sigma-\mathcal G\widetilde\sigma\|_\infty
\le A^{-1}\|\sigma-\widetilde\sigma\|_\infty,
\qquad A^{-1}<0.228.
\label{hmtfg:escalado:eq:15}
\end{equation}
Le théorème de contraction produit un unique graphe fixe \(x=\sigma_*(y)\). Sur celui-ci,
\begin{equation}
\|y_n\|_1\le q^n\|y_0\|_1,
\qquad
\|(x_n,y_n)\|_{\mathcal A}\le(1+L)q^n\|y_0\|_1.
\label{hmtfg:escalado:eq:16}
\end{equation}
De plus, la distance verticale \(d(x,y)=x-\sigma_*(y)\) satisfait
\[
|d(R(x,y))|\ge A|d(x,y)|
\]
tant que l’orbite reste dans le cylindre. Toute orbite qui y demeure indéfiniment appartient au graphe. Celui-ci identifie exactement la feuille stable locale et fournit l’estimation \eqref{hmtfg:escalado:eq:16}. Sa régularité locale coïncide avec celle de la variété stable analytique du point fixe hyperbolique.


\subsection{Existence, unicité et transversalité du paramètre}
\label{hmtfg:escalado:6}

Pour \(\lambda\) dans \eqref{hmtfg:escalado:eq:2}, définissons
\[
H(\lambda)=u_\lambda-u_g-
\sigma_*(\nu_\lambda-\nu_g).
\]
Le certificat assure \(\|\nu_\lambda-\nu_g\|_1<0.001436077<r\). Pour \(\lambda_2>\lambda_1\),
\begin{equation}
H(\lambda_2)-H(\lambda_1)
\ge\int_{\lambda_1}^{\lambda_2}
\bigl[u'_t-L\|\nu'_t\|_1\bigr]dt
>3733.007995\,(\lambda_2-\lambda_1).
\label{hmtfg:escalado:eq:17}
\end{equation}
Pour déterminer les signes aux extrémités, on utilise \(\|g-g_0\|_{\mathcal A}<\varepsilon_0\) et \(|\sigma_*(y)|\le L\|y\|_1\). L’évaluation ponctuelle rationnelle produit
\begin{equation}
H(1.874038)<-0.0011856110945,
\qquad H(1.874039)>0.0021866053399.
\label{hmtfg:escalado:eq:18}
\end{equation}
Par continuité, il existe une racine ; \eqref{hmtfg:escalado:eq:17} prouve son unicité et la transversalité du croisement. À cette racine, \eqref{hmtfg:escalado:eq:16} et
\((1+L)(0.001396077+0.00004)<0.001666\)
démontrent \eqref{hmtfg:escalado:eq:3}.

La borne est uniforme en le nombre de retours : elle convertit l’avancement en profondeur en une erreur operatorielle géométrique explicite.


\subsection{Admissibilité réelle de tous les retours}
\label{hmtfg:escalado:7}

Les quatre retours initiaux s’accompagnent, sur chaque sous-intervalle rationnel, des inégalités
\begin{equation}
0<a=-f(1)<1,\qquad b=f(a)>a,\qquad f(b)<a.
\label{hmtfg:escalado:eq:19}
\end{equation}
Le programme conserve \(a,b,f(b)\) avant chaque application de \eqref{hmtfg:escalado:eq:7}. L’unimodalité, la parité et le point critique quadratique sont conservés par ces retours normalisés.

Pour les itérations suivantes, une vérification uniforme dans la boule de \eqref{hmtfg:escalado:eq:9} suffit. Si \(E=V-V_{g_0}\), la norme \eqref{hmtfg:escalado:eq:5} donne, pour \(-0.4\le z\le0\),
\[
|E(z)|\le0.004,\qquad |E'(z)|\le0.01.
\]
La seconde inégalité utilise \(\sup_{j\ge1}j(0.4)^{j-1}=1\). L’évaluation rationnelle des coefficients centraux fournit
\begin{equation}
\frac65<V(z)<\frac85,\qquad |V'(z)|<\frac12,
\qquad 0.39<a<0.41.
\label{hmtfg:escalado:eq:20}
\end{equation}
En conséquence,
\[
f'(x)=-2x\left[V(z)+\frac25x^2V'(z)\right],
\]
et l’expression entre crochets dépasse un pour \(0\le x\le1\). L’application est unimodale et préserve \([-1,1]\). De plus,
\[
b=f(a)>1-(0.41)^2\frac85=\frac{4569}{6250}>a,
\]
\begin{equation}
f(b)<1-\left(\frac{4569}{6250}\right)^2\frac65
=\frac{35028967}{97656250}<0.39<a.
\label{hmtfg:escalado:eq:21}
\end{equation}
Les orbites de la feuille stable satisfont \eqref{hmtfg:escalado:eq:19} à toute profondeur. Ainsi, les itérations analytiques de l’opérateur correspondent à des retours réels de doublement, dont les domaines et les orientations sont conservés.


\subsection{Changement d’échelle de la cascade superstable locale}
\label{hmtfg:escalado:8}

La conclusion métrique exige deux croisements : celui de la famille avec la feuille stable, démontré dans la section~\ref{hmtfg:escalado:6}, et celui de la variété instable avec la cible superstable. Le second est démontré par Eckmann–Wittwer pour
\(\Sigma_2=\{\psi:\psi(1)=0\}\), correspondant au cycle superstable de période deux. La normalisation \eqref{hmtfg:escalado:eq:6} coïncide avec celle de l’opérateur réel pair utilisé dans ce travail.

Le théorème général de Collet–Eckmann–Lanford, section 6 de la source citée, proposition 6.1 et théorèmes 6.2–6.3, s’applique à une application d’un espace de Banach \(C^2\), un point fixe possédant une unique valeur propre instable \(\delta_{\mathrm F}>1\), une courbe transverse à sa feuille stable et une cible transverse à sa variété instable. Les préimages locales de la cible rencontrent la courbe en des points qui satisfont
\begin{equation}
\delta_{\mathrm F}^j(\mu_j-\lambda_*)\longrightarrow C\ne0.
\label{hmtfg:escalado:eq:22}
\end{equation}
L’indice \(j\) compte les retours locaux. Comme la courbe initiale est ici \(R^4f_\lambda\), la période physique est \(2^{j+5}\) lorsque la cible est prise dans \(\Sigma_2\) ; tout nombre fixe de pas utilisé pour ramener la cible dans le voisinage local est incorporé à un décalage fixe de l’indice et à \(C\).

Les conditions réelles de la section~\ref{hmtfg:escalado:7} et la branche réelle de doublement de la cible fixent la période exacte et l’itinéraire. La suite est sélectionnée par ce prolongement local orienté. En prenant les différences dans \eqref{hmtfg:escalado:eq:22}, on obtient la limite des rapports de \eqref{hmtfg:escalado:eq:4}.

Sources des deux résultats utilisés : \href{https://link.springer.com/article/10.1007/BF01013368}{Eckmann–Wittwer, \emph{A complete proof of the Feigenbaum conjectures}}; \href{https://people.math.harvard.edu/~knill/history/lanford/papers/ColletEckmannLanford.pdf}{Collet–Eckmann–Lanford, section 6}.

\subsubsection{Reste asymptotique en puissance}

L’analyticité de notre famille permet de préciser \eqref{hmtfg:escalado:eq:22}. Dans la construction de coordonnées normales de Collet–Eckmann–Lanford, seules la variété stable, la variété instable et l’hypersurface cible sont aplaties. Ces trois sous-variétés admettent des changements locaux \(C^2\) à dérivées bornées. La coupure lisse employée s’effectue dans la coordonnée scalaire instable ; le feuilletage complet s’obtient ensuite par le changement construit.

Dans ces coordonnées, la correction s’écrit comme une série d’incréments \(w_j\) dont la norme \(C^1\) satisfait
\[
\|w_j\|_{C^1}\le C_1\rho^j,\qquad 0<\rho<1.
\]
La règle de dérivation en chaîne pour l’application munie d’une coupure, avec des dérivées d’ordres un et deux uniformément bornées, fournit
\[
\|w_j\|_{C^2}\le C_2B^j
\]
pour un certain \(B\ge1\). La norme pondérée de la construction contrôle la norme \(C^1\) ordinaire sur le domaine considéré. Par interpolation,
\[
[Dw_j]_\beta\le C_3
\left(\rho^{1-\beta}B^\beta\right)^j.
\]
On choisit
\begin{equation}
0<\beta<\min\left\{1,
\frac{-\log\rho}{\log B-\log\rho}\right\}.
\label{hmtfg:escalado:eq:23}
\end{equation}
La série converge dans \(C^{1,\beta}\). Soit \(z\) la coordonnée normale obtenue, normalisée de sorte que les préimages de la cible satisfassent \(z=\delta_{\mathrm F}^{-j}\), en absorbant le décalage fixe de l’indice. Comme notre courbe est analytique et transverse,
\[
z(\Gamma(\lambda))=c(\lambda-\lambda_*)
+O(|\lambda-\lambda_*|^{1+\beta}),\qquad c\ne0.
\]
L’inversion locale donne
\begin{equation}
\mu_j-\lambda_*=C\delta_{\mathrm F}^{-j}
+O\left(\delta_{\mathrm F}^{-(1+\beta)j}\right).
\label{hmtfg:escalado:eq:24}
\end{equation}
Cela démontre le reste de \eqref{hmtfg:escalado:eq:4}, avec \(\theta=\delta_{\mathrm F}^{-\beta}\). L’énoncé affirme l’existence des constantes du reste paramétrique. Le taux operatoriel explicite \(0.83996\) de \eqref{hmtfg:escalado:eq:3} contrôle, quant à lui, la convergence des applications renormalisées.

\subsubsection{Conversion du reste en erreur sur les rapports}

Si \(\mu_j-\lambda_*=C\delta_{\mathrm F}^{-j}(1+e_j)\), avec \(|e_j|\le M\theta^j\), on écrit
\[
\mu_{j-1}-\mu_j=C(\delta_{\mathrm F}-1)\delta_{\mathrm F}^{-j}(1+r_j),
\quad
r_j=\frac{\delta_{\mathrm F} e_{j-1}-e_j}{\delta_{\mathrm F}-1}.
\]
Avec \(K=M(\delta_{\mathrm F}+\theta)/(\delta_{\mathrm F}-1)\), on a \(|r_j|\le K\theta^{j-1}\). Par conséquent,
\begin{equation}
\boxed{
\left|\frac{\mu_{j-2}-\mu_{j-1}}{\mu_{j-1}-\mu_j}-\delta_{\mathrm F}\right|
\le\frac{\delta_{\mathrm F} K(1+\theta)\theta^{j-2}}
{1-K\theta^{j-1}}}
\label{hmtfg:escalado:eq:25}
\end{equation}
lorsque \(K\theta^{j-1}<1\). Cette formule sépare la précision d’évaluation de chaque racine de l’erreur de troncature de la cascade.


\subsection{Prolongement certifié des huit racines initiales}
\label{hmtfg:escalado:9}

Le certificat de prolongement fini conserve la famille \eqref{hmtfg:escalado:eq:1} et le certificat antérieur de huit racines. Pour \(2\le n\le8\), il étudie
\[
S_n(\lambda)=f_\lambda^{2^n}(0)
\]
depuis la boîte certifiée de \(\lambda_{n-1}\) jusqu’à \(1.874039\). Le recouvrement rationnel complet contient exactement la racine héritée \(\lambda_{n-1}\) et la nouvelle racine \(\lambda_n\). En particulier, \(\lambda_n\) est l’unique nouvelle racine dans \((\lambda_{n-1},1.874039]\).

Les boîtes sont exclues par le signe du retour ou par une dérivée locale de signe certifié. Les boîtes d’ancrage conservent l’existence et l’unicité par les signes aux extrémités et la monotonie locale. Leurs extrémités couvrent exactement chaque intervalle ; les adjacences sont vérifiées. L’évaluation utilise les moments \eqref{hmtfg:escalado:eq:8}, l’évaluation de Horner d’ordre 32, des queues géométriques extérieures et \(|u_0''|\le2\). La queue uniforme du potentiel est inférieure à \(1.294\,10^{-58}\) sur \(|x|\le3/2\).

Le certificat final contient 650 boîtes traitées et 332 feuilles de recouvrement. Le premier niveau possède l’expression exacte \(\lambda_1=1/u_0(1)\). Les approximations suivantes sont présentées uniquement à titre indicatif ; leurs boîtes complètes sont conservées dans les certificats :

\begingroup\small
\begin{longtable}{p{0.465\linewidth}p{0.465\linewidth}}
\hline
Niveau & Paramètre superstable \\
\hline
1 & 1.345913990471832792210949095… \\
2 & 1.763379787846910127290794030… \\
3 & 1.850413796950136464530566778… \\
4 & 1.868979756606798125150574309… \\
5 & 1.872955034435659740004205128… \\
6 & 1.873806367467030119412262311… \\
7 & 1.873988695221365362307168780… \\
8 & 1.874027744154450672897007573… \\
\hline\end{longtable}
\endgroup

% Fragmento integrable. Preambulo anfitrion: amsmath, amssymb, longtable, hyperref.
\section{Isolement de la limite et changement d’échelle du sélecteur global}
\label{hmtfg:fragmento:aislamiento-del-limite-y-escalado-del-selector-global}

\subsection{Continuation indéfinie de la première racine et exclusion de l’intervalle intermédiaire}
\label{hmtfg:escalado:15}

La classification complète démontrée précédemment conserve la famille et le sélecteur de première racine nouvelle. Le transport ordonné conserve les zéros et leur minimum ; le mot de retours conserve la lecture chronologique nonadique.

\subsubsection{Classification et existence à toute profondeur}
\label{hmtfg:escalado:15.1}

Écrivons \(W_n=\tau_1\cdots\tau_{2^n-1}\), avec \(\tau_j=(-1)^{v_2(j)}\), et \(a_*=\min\Lambda_\tau\). L’existence et la compacité de \(\Lambda_\tau\) sont démontrées dans le développement de prolongation, \ref{hmtfg:prolongacion:6.1}–\ref{hmtfg:prolongacion:6.2} ; son minimum existe par compacité et ordre.

\textbf{Classification.} Toute application strictement unimodale à maximum, sur un intervalle invariant, dont le point critique a pour période exacte \(2^n\) et dont l’itinéraire \(K\) satisfait \(K<\tau\), a pour itinéraire \(W_n0\).

L’induction s’effectue au moyen du retour effectif. Pour une période d’au moins huit, la comparaison avec \(\tau\), jointe aux intervalles qui piégeraient des orbites de signes incompatibles, impose le préfixe \((+,-,+,+,+,-)\). On en déduit
\[
J_1=[d_2,d_4],\quad F(J_1)=[d_3,1],\quad
d_4<d_3,\quad F^2(J_1)=J_1.
\]
Le retour normalisé par \(h(x)=d_2x\) a pour maximum un, pour période la moitié de la période initiale et pour itinéraire \(K'_j=-K_{2j}<\tau\). L’ordre est conservé parce que
\[
O_{2q-1}(K)=-O_{q-1}(K'),\qquad
K_{2q}-\tau_{2q}=-(K'_q-\tau_q).
\]
Les cas initiaux sont \(+0\) et \(+-+0\) ; l’induction reconstruit \(W_n0\). Pour les applications paires, on a aussi \(0<d_4<-d_2<1\), de sorte que le retour s’étend à \([-1,1]\) et coïncide avec \eqref{hmtfg:escalado:eq:6}.

Avant \(a_*\), tous les itinéraires sont inférieurs à \(\tau\) : les ensembles de comparaison stricte sont ouverts par la séparation démontrée dans \ref{hmtfg:prolongacion:6.1} ; l’intervalle précédant \(a_*\) est connexe et commence par \(+^\infty<\tau\).

\textbf{Existence de chaque première racine.} Supposons \(\lambda_{n-1}<a_*\) construite, et soit \(p=2^{n-1}\). Les zéros de \(c_p\) dans \([\lambda_{n-1},a_*]\) forment un ensemble fini non vide : la fonction est analytique et \(c_p(a_*)\ne0\). Soit \(z<a_*\) son dernier zéro. Dans \((z,a_*]\), \(c_p\) est de signe \(\tau_p\). Pour \(g_\lambda=f_\lambda^p\),
\[
g_\lambda'(0)=0,\qquad
c_{2p}(\lambda)=c_p(\lambda)+O(c_p(\lambda)^2).
\]
Au voisinage de \(z\), à droite, \(c_{2p}\) est de signe \(\tau_p\) ; en \(a_*\), son signe est opposé. Il existe un zéro de \(c_{2p}\) dans \((z,a_*)\), et sa période est exactement \(2p\), car on y a \(c_p\ne0\). La finitude des zéros fournit la première racine nouvelle dans \((\lambda_{n-1},a_*)\). Le dernier zéro auxiliaire prouve l’existence ; le choix effectif reste la première racine. Le cas de base est \(\lambda_1=1/u_0(1)<a_*\).

\subsubsection{Limite de la suite initiale}
\label{hmtfg:escalado:15.2}

\textbf{Théorème.} Le sélecteur initial existe à tous les niveaux et satisfait
\begin{equation}
\boxed{K_{\lambda_n}=W_n0,\qquad
\lambda_n\nearrow a_*=\min\Lambda_\tau.}
\label{hmtfg:escalado:eq:50}
\end{equation}
\textbf{Démonstration.} Le résultat d’existence précédent donne une suite croissante majorée par \(a_*\). La classification fixe tous ses itinéraires. Soit \(b\le a_*\) sa limite. Pour chaque \(p\), les termes suffisamment avancés conservent les signes opposés \(\tau_p,\tau_{2p}\). La formule de Taylor et la borne uniforme \(B_p=16^p(16^p-1)/15\) donnent
\[
|c_{2p}(\lambda_n)-c_p(\lambda_n)|
\le\tfrac12 B_p|c_p(\lambda_n)|^2,\qquad
|c_p(\lambda_n)|>2/B_p.
\]
Le passage à la limite conserve \(\operatorname{sign}c_p(b)=\tau_p\), avec une valeur absolue d’au moins \(2/B_p>0\). Pour tout \(p\), cela implique \(b\in\Lambda_\tau\) ; la minimalité donne \(b=a_*\).

Le résultat établit l’existence indéfinie, la conservation de la cascade symbolique et la sélection de son premier paramètre limite pour la règle globale initiale. La complétion HMT transporte cette suite et sa limite en conservant l’ordre et la lecture.

\subsubsection{Certificat d’exclusion de l’intervalle intermédiaire}
\label{hmtfg:escalado:15.3}

Le programme de pont certifie
\begin{equation}
\Lambda_\tau\cap[\lambda_{8,\rm inf},1.874038]=\varnothing,
\label{hmtfg:escalado:eq:51}
\end{equation}
où
\[
\lambda_{8,\rm inf}
=1.874027744154450672897007573595092408435133414523352839687903472138975
\]
est l’extrémité inférieure rationnelle de l’intervalle certifié de \(\lambda_8\).

Le recouvrement contient 374 intervalles exactement adjacents et aucune région en suspens : 319 avec \(c_{512}>0\), 10 avec \(c_{1024}<0\), deux avec \(c_{2048}>0\), tous opposés au signe de \(\tau\), et 43 proches de centres de périodes 256, 512 ou 1024, exclus par leur deuxième retour.

Dans ce dernier cas, \(g=f_\lambda^p\), \(d=g(0)\) et une borne uniforme \(|g''|\le M\) entre zéro et \(d\) satisfont \(M|d|<2\). Alors
\[
|g(d)-d|\le\tfrac12M|d|^2<|d|\quad(d\ne0).
\]
Les deux retours ont le même signe ; pour \(d=0\), ils sont tous deux nuls. Les deux situations excluent \(\tau_{2p}=-\tau_p\) avec des signes stricts.

L’évaluation utilise des moments rationnels, des queues géométriques du potentiel et de sa dérivée seconde, et \(|u_0'''|<6\). Le centrage paramétrique intersecte l’enveloppe naturelle avec
\[
c_j(\lambda_{\rm medio})+
\partial_\lambda c_j(I)\,(I-\lambda_{\rm medio});
\]
la dérivée est propagée sur la boîte entière. La propriété d’autoapplication de \([-1,1]\) est démontrée avant d’y restreindre les enveloppes.

Le certificat complet conserve le recouvrement et vérifie son pavage exact, ainsi que chaque inégalité de signe ou de Taylor.

\subsubsection{Localisation conjointe}
\label{hmtfg:escalado:15.4}

Le croisement local \(\lambda_*\) réalise \(\tau\) par ses retours réels à toute profondeur, de sorte que \(a_*\le\lambda_*\). Par \eqref{hmtfg:escalado:eq:50}, \(\lambda_8<a_*\) ; avec \eqref{hmtfg:escalado:eq:51},
\begin{equation}
\boxed{1.874038<a_*\le\lambda_*<1.874039.}
\label{hmtfg:escalado:eq:52}
\end{equation}
La localisation conserve l’itinéraire canonique et situe son premier paramètre limite dans le même intervalle que le croisement stable.

\subsection{Isolement par première sortie et coïncidence des limites}
\label{hmtfg:escalado:16}

Cette section complète l’isolement identifié dans la section~\ref{hmtfg:escalado:15.4}. Elle conserve la famille HMT \eqref{hmtfg:escalado:eq:1}, la sélection minimale \eqref{hmtfg:escalado:eq:50}, la carte \eqref{hmtfg:escalado:eq:5} et la feuille stable des sections~\ref{hmtfg:escalado:5}--\ref{hmtfg:escalado:7}. La preuve utilise l’ordre \(a_*\le\lambda_*\) : il suffit d’exclure une séparation vers le côté négatif du cône expansif.

\subsubsection{Bornes supérieure et inférieure du transport transversal}
\label{hmtfg:escalado:16.1}

Outre \eqref{hmtfg:escalado:eq:9}, le lemme 2.1 de \href{https://arxiv.org/pdf/1410.3277}{Hertling–Spandl, section 2} fournit \(\|D\Phi\|<0.9\) dans la même boule de rayon \(0.01\), où
\[
\Phi_u(u,\nu)=u-\frac{U(u,\nu)-u}{3.669}.
\]
Par conséquent,
\begin{equation}
\left|1-\frac{U_u-1}{3.669}\right|<0.9,\qquad
U_u<1+1.9(3.669)=7.9711.
\label{hmtfg:escalado:eq:53}
\end{equation}
Le nombre \(3.669\) appartient au préconditionneur publié de ce certificat analytique ; la famille générée \eqref{hmtfg:escalado:eq:1} conserve ses coefficients antérieurs.

Prenons \(\kappa=0.21\). Pour deux points de la boule joints par une sécante avec
\(\Delta u<0\), \(\|\Delta\nu\|_1\le\kappa|\Delta u|\), l’intégration des blocs de \(DR\) sur le segment donne
\begin{equation}
4.4034|\Delta u|
\le|\Delta u'|\le8.0887|\Delta u|,\qquad \Delta u'<0,
\label{hmtfg:escalado:eq:54}
\end{equation}
\begin{equation}
\|\Delta\nu'\|_1
\le(0.756+0.719\kappa)|\Delta u|
=0.90699|\Delta u|
<\kappa(4.4034)|\Delta u|.
\label{hmtfg:escalado:eq:55}
\end{equation}
Ici \(4.4034=4.521-0.560\kappa\) et \(8.0887=7.9711+0.560\kappa\). Ainsi, le cône des sécantes conserve son orientation et se dilate tant que le segment reste dans la boule.

\subsubsection{Réduction d’une séparation infinie à une région finie}
\label{hmtfg:escalado:16.2}

Supposons \(a_*<\lambda_*\), et posons
\[
f_n=R^{4+n}f_{a_*},\qquad
s_n=R^{4+n}f_{\lambda_*},\qquad
(\Delta u_n,\Delta\nu_n)=f_n-s_n.
\]
La borne positive de \(u'\) et \(\|\nu'\|_1/u'<0.144386<\kappa\) sur \(\Gamma(J)\) placent la sécante initiale dans le cône négatif. La feuille stable est écrite par rapport à \(g\), avec
\begin{equation}
s_n=g+e_n,\quad
\|e_n\|_{\mathcal A}<0.001666(0.83996)^n,\quad
|(e_n)_u|\le0.16\|(e_n)_\nu\|_1.
\label{hmtfg:escalado:eq:56}
\end{equation}
Les contrôles d’entrée impliquent
\begin{equation}
|\Delta u_0|
<0.004993128+0.00004
 +0.16(0.001396077+0.00004)
=0.00526290032< h,\qquad h=0.006.
\label{hmtfg:escalado:eq:57}
\end{equation}
Tant que \(|\Delta u_n|<h\), \eqref{hmtfg:escalado:eq:55}–\eqref{hmtfg:escalado:eq:56} donnent
\begin{equation}
\|f_n-g_0\|_{\mathcal A}
<(1+\kappa)h+0.001666+0.00004
=0.008966<0.01.
\label{hmtfg:escalado:eq:58}
\end{equation}
Le point \(s_n\) et le segment entre les deux restent également dans cette boule convexe. L’expansion inférieure \eqref{hmtfg:escalado:eq:54} impose une première sortie finie \(N\), et la borne supérieure appliquée à l’étape précédente fournit
\begin{equation}
0.006\le-\Delta u_N<0.0485322,\qquad
\|\Delta\nu_N\|_1\le0.21|\Delta u_N|.
\label{hmtfg:escalado:eq:59}
\end{equation}
On utilise l’inégalité \(8.0887h=0.0485322\) ; le saut complet, et non seulement la surface de première sortie, est inclus.

Par conséquent, \(f_N\) appartient à la région suivante :
\begin{equation}
\begin{split}
\mathcal C_-=\{\,g+(d_u,d_\nu)+e:\;&
-0.0485322\le d_u\le-0.006,\\
&\|d_\nu\|_1\le0.21|d_u|,\quad
|e_u|+\|e_\nu\|_1\le0.001666,\\
&|e_u|\le0.16\|e_\nu\|_1\,\}.
\end{split}
\label{hmtfg:escalado:eq:60}
\end{equation}
Puisque \(a_*\) réalise \(\tau\), tous ses retours normalisés sont des autoapplications unimodales de \([-1,1]\) ayant le même itinéraire. Cette propriété découle de l’induction des retours, indépendamment de la distance de \(f_N\) à \(g_0\). C’est pourquoi le certificat suivant traite \(\mathcal C_-\) intersectée avec cette classe d’autoapplications.

\subsubsection{Exclusion certifiée de la région de première sortie}
\label{hmtfg:escalado:16.3}

Le certificat rationnel d’exclusion par première sortie, reproduit dans le supplément computationnel, couvre exactement l’intervalle de \eqref{hmtfg:escalado:eq:59} par 64 bandes rationnelles adjacentes. Toutes sont exclues ; aucune subdivision supplémentaire ni région en suspens n’intervient dans le résultat.

Parmi les 64 bandes, 54 présentent un signe contraire et dix un retour contractant de période 16 ou 32. La plus grande borne de Lipschitz de ces retours est inférieure à \(0.475146491\) ; l’horizon maximal d’exclusion est l’itérée 64.

L’évaluation conserve comme variables partagées les quarante premiers coefficients de la perturbation, ainsi qu’une borne explicite de leur queue infinie. Pour chaque bande, le centre est \(g_0\) déplacé dans la coordonnée \(u\). Si \(a,b\ge0\) bornent les sensibilités duales dans \(u,\nu\), le support de l’incertitude du point fixe et du résidu stable est
\begin{equation}
0.00004\max(a,b)
+0.001666\max\left(b,\frac{0.16a+b}{1.16}\right).
\label{hmtfg:escalado:eq:61}
\end{equation}
La deuxième expression s’obtient en maximisant \(a|e_u|+b\|e_\nu\|_1\) sous les deux contraintes de \eqref{hmtfg:escalado:eq:56}. On utilise ainsi l’orientation de la feuille stable démontrée précédemment, en plus de sa norme.

Soit \(x_j\) l’orbite du centre et \(L_j\) sa sensibilité linéaire à tous les coefficients partagés. Le modèle de Taylor prend la forme
\[
c_j=x_j+L_j(\Delta v)+r_j,\qquad |r_j|\le E_j.
\]
Si \(\rho_j\) borne \(|L_j(\Delta v)|+E_j\), la récurrence extérieure employée est
\begin{equation}
E_{j+1}\le |f_0'(x_j)|E_j
+\tfrac12\sup|f_0''|\rho_j^2
+\sup|\Delta f'|\rho_j.
\label{hmtfg:escalado:eq:62}
\end{equation}
Les dérivées sont bornées sur le segment des deux états. L’orbite centrale est vérifiée à l’intérieur de \([-1,1]\) ; les autres orbites appartiennent à cet intervalle par la propriété d’autoapplication incluse dans le domaine certifié. Les calculs utilisent un arrondi rationnel vers l’extérieur, y compris pour les queues géométriques.

Chaque bande est exclue par l’une de deux inégalités :

\par\noindent\textbf{1.}\enspace Une valeur critique a un signe strictement opposé à \(\tau_j\).
\par\noindent\textbf{2.}\enspace Pour \(p=2^k\), le retour \(H=f^p\) a pour constante de Lipschitz \(L_p<1\) sur le segment entre \(0\) et \(H(0)\). Alors
\[
|H(H(0))-H(0)|\le L_p|H(0)|.
\]
Si \(H(0)\ne0\), les deux valeurs ont le même signe ; si \(H(0)=0\), elles s’annulent toutes deux. Dans les deux cas, l’itinéraire strict \(\tau_{2p}=-\tau_p\) est exclu.

La constante \(L_p\) s’obtient en propageant le segment initial \([-r,r]\), \(r\ge|H(0)|\), avec les boîtes de l’orbite critique et en multipliant les bornes extérieures des dérivées. On borne les applications de la bande entière, en maintenant \eqref{hmtfg:escalado:eq:61}.

\subsubsection{Théorème de coïncidence du paramètre d’accumulation}
\label{hmtfg:escalado:16.4}

\textbf{Théorème.} Pour la famille uniforme HMT \eqref{hmtfg:escalado:eq:1}, la limite de la sélection initiale des premières racines coïncide avec le croisement stable certifié :
\begin{equation}
\boxed{\lambda_n\nearrow a_*=\lambda_*,
\qquad 1.874038<\lambda_*<1.874039.}
\label{hmtfg:escalado:eq:63}
\end{equation}

\textbf{Démonstration.} \eqref{hmtfg:escalado:eq:50}–\eqref{hmtfg:escalado:eq:52} donnent \(\lambda_n\uparrow a_*\le\lambda_*\), les deux paramètres étant dans \(J\). Si l’inégalité était stricte, \eqref{hmtfg:escalado:eq:54}–\eqref{hmtfg:escalado:eq:59} produiraient un retour \(f_N\in\mathcal C_-\) réalisant \(\tau\). Le recouvrement complet de la section~\ref{hmtfg:escalado:16.3} exclut précisément cette possibilité. Donc \(a_*=\lambda_*\).

La preuve contrôle toutes les profondeurs au moyen d’une première sortie finie et d’un recouvrement uniforme de sa région. L’identité \eqref{hmtfg:escalado:eq:63} relie la limite du sélecteur global au croisement transversal.


\subsection{Identification des centres initiaux et changement d’échelle global}
\label{hmtfg:escalado:17}

La dernière étape conserve la sélection de \eqref{hmtfg:escalado:eq:50}. La continuation locale et les paramètres initiaux sont comparés par leur période exacte et leur classe hybride ; la coïncidence s’obtient par une unicité dans un voisinage paramétrique fixe.

\subsubsection{Réalisation analytique de degré deux et transversalité}
\label{hmtfg:escalado:17.1}

Le point fixe \(g\) de la section~\ref{hmtfg:escalado:3} est pair, analytique, à point critique quadratique et de combinatoire stationnaire de doublement. Le théorème 2.1 de \href{https://annals.math.princeton.edu/wp-content/uploads/annals-v164-n3-p01.pdf}{de Faria–de Melo–Pinto}, appliqué à cette unique combinatoire, donne un ensemble limite à un seul élément, avec une extension holomorphe propre de degré deux entre disques topologiques emboîtés. Son attraction sur l’intervalle identifie cet élément à \(g\), puisque \(Rg=g\).

Le lemme 3.2 de la même source étend un nombre fixe de retours à un voisinage analytique complet de \(g\), avec des images de degré deux et un module positif. La norme \eqref{hmtfg:escalado:eq:5} contrôle la norme uniforme sur un voisinage complexe suffisamment petit de \([-1,1]\), car on y a \(|(x^2-1)/2.5|<1\). La convergence \eqref{hmtfg:escalado:eq:3} entre dans ce voisinage ; la dépendance analytique des retours fournit, pour un certain entier fixe \(d\), une famille
\begin{equation}
\mathcal F_\lambda=R^d f_\lambda
\label{hmtfg:escalado:eq:64}
\end{equation}
d’applications de degré deux, analytique dans un voisinage complexe de \(\lambda_*\). Les disques et le nombre de retours sont fixés avant de faire varier \(\lambda\) dans ce voisinage.

La tangente paramétrique appartient au cône expansif, tandis que la feuille stable satisfait la borne de pente de la section~\ref{hmtfg:escalado:5}. Les retours conservent cette séparation. La réalisation de degré deux identifie localement la feuille stable à la classe hybride de l’application infiniment renormalisable ; \eqref{hmtfg:escalado:eq:64} est donc transversale à cette classe.

Le changement d’espace analytique est justifié en détail dans \ref{hmtfg:clasificacion:8.2}, consacrée au transport de la transversalité. Une direction tangente à la classe hybride produirait des tangentes renormalisées décroissantes, par la convergence uniforme sur un disque hybride et l’estimation de Cauchy. La direction certifiée satisfait, en revanche, \(|\dot f(1)|=|\dot u|/10\) et croît sous les retours grâce au cône expansif. Cette évaluation commune des deux représentants prouve la transversalité sans supposer l’équivalence de leurs normes.

\subsubsection{Unicité dans un voisinage fixe}
\label{hmtfg:escalado:17.2}

Pour \(n>d\), la classification de la section~\ref{hmtfg:escalado:15} donne
\begin{equation}
K(\mathcal F_{\lambda_n})=W_{n-d}0.
\label{hmtfg:escalado:eq:65}
\end{equation}
Le redressement d’une application de degré deux conserve l’orbite critique ; \eqref{hmtfg:escalado:eq:65} identifie sa classe hybride à celle du centre quadratique canonique de période \(2^{n-d}\). Les bornes a priori de la suite réelle de doublement sont celles employées dans le théorème d’universalité de Lyubich.

Le \href{https://www.maths.tcd.ie/EMIS/journals/Annals/149_2/lyubich.pdf\#page=69}{théorème 7.4 de Lyubich, p. 387–388} donne une seule intersection avec chacune de ces classes, aux niveaux suffisamment élevés, sur la transversale analytique et dans un voisinage fixe du paramètre d’accumulation.

Par \eqref{hmtfg:escalado:eq:63}, \(\lambda_n\) entre finalement dans ce voisinage. La continuation locale de la section~\ref{hmtfg:escalado:8}, indexée par la même période, appartient à la même classe hybride. L’unicité implique
\begin{equation}
\boxed{\lambda_n=\widehat\mu_n\quad(n\ge n_0),}
\label{hmtfg:escalado:eq:66}
\end{equation}
où \(\widehat\mu_n\) désigne le centre local de période physique \(2^n\). Il s’agit d’une réindexation de \(\mu_j\) par un décalage fixe déterminé par les retours et la cible de la section~\ref{hmtfg:escalado:8}. La sélection initiale de \(\lambda_n\) reste intacte.

L’entier \(n_0\) est existentiel dans cette preuve. La continuation à tous les niveaux et le mot \(W_n0\) sont déjà démontrés dans \eqref{hmtfg:escalado:eq:50} ; \eqref{hmtfg:escalado:eq:66} apporte l’identité métrique nécessaire à la limite.

\subsubsection{Théorème de changement d’échelle pour les premières racines HMT}
\label{hmtfg:escalado:17.3}

\textbf{Théorème.} Pour la famille uniforme \eqref{hmtfg:escalado:eq:1}, soient \(\lambda_n\) les premières racines nouvelles successives de période critique exacte \(2^n\). Il existe \(C>0\), \(0<\theta<1\) et \(M<\infty\) tels que, pour tout \(n\) suffisamment grand,
\begin{equation}
\lambda_*-\lambda_n=C\delta_{\mathrm F}^{-n}(1+e_n),
\qquad |e_n|\le M\theta^n,
\label{hmtfg:escalado:eq:67}
\end{equation}
et
\begin{equation}
\boxed{\lim_{n\to\infty}
\frac{\lambda_{n-1}-\lambda_{n-2}}
{\lambda_n-\lambda_{n-1}}=\delta_{\mathrm F}.}
\label{hmtfg:escalado:eq:68}
\end{equation}

\textbf{Démonstration.} \eqref{hmtfg:escalado:eq:66} transporte le changement d’échelle local \eqref{hmtfg:escalado:eq:24} aux racines initiales ; le décalage fini des indices est absorbé dans \(C,M\). Son signe est positif parce que \(\lambda_n<\lambda_*\).

Le reste en puissance admet aussi une justification directe par l’holonomie des classes hybrides. Le lemme 7.3 de Lyubich donne une régularité \(C^{1+\beta}\), avec une dérivée non nulle, sur la transversale au point d’accumulation. La dynamique instable unidimensionnelle est analytique et a pour multiplicateur \(\delta_{\mathrm F}>1\) ; sa coordonnée linéarisante place les centres successifs à une distance proportionnelle à \(\delta_{\mathrm F}^{-n}\). La composition avec l’inverse de l’holonomie donne un terme linéaire non nul et un reste \(O(\delta_{\mathrm F}^{-(1+\beta)n})\). On peut ainsi prendre \(\theta=\delta_{\mathrm F}^{-\beta}\).

En écrivant
\[
r_n=\frac{\delta_{\mathrm F} e_{n-1}-e_n}{\delta_{\mathrm F}-1},
\qquad K=\frac{M(\delta_{\mathrm F}+\theta)}{\delta_{\mathrm F}-1},
\]
on obtient
\[
\lambda_n-\lambda_{n-1}
=C(\delta_{\mathrm F}-1)\delta_{\mathrm F}^{-n}(1+r_n),\qquad
|r_n|\le K\theta^{n-1}.
\]
Par conséquent,
\begin{equation}
\left|
\frac{\lambda_{n-1}-\lambda_{n-2}}
{\lambda_n-\lambda_{n-1}}-\delta_{\mathrm F}
\right|
\le
\frac{\delta_{\mathrm F} K(1+\theta)\theta^{n-2}}
{1-K\theta^{n-1}},
\quad K\theta^{n-1}<1.
\label{hmtfg:escalado:eq:69}
\end{equation}
Le membre de droite tend vers zéro, ce qui démontre \eqref{hmtfg:escalado:eq:68}.

\subsubsection{Enchaînement et portée de la clôture}
\label{hmtfg:escalado:17.4}

La composition effective est : lecteur dodécaphasique APP–TRIT–TPK ; profil \eqref{hmtfg:escalado:eq:1} ; transport ordonné des retours \eqref{hmtfg:escalado:eq:38}–\eqref{hmtfg:escalado:eq:40} ; classification et existence des premières racines \eqref{hmtfg:escalado:eq:50} ; localisation \eqref{hmtfg:escalado:eq:52} ; isolement par première sortie \eqref{hmtfg:escalado:eq:63} ; identité des centres à partir d’un certain rang \eqref{hmtfg:escalado:eq:66} ; changement d’échelle \eqref{hmtfg:escalado:eq:67}–\eqref{hmtfg:escalado:eq:69}. La représentation de Jacobi \eqref{hmtfg:escalado:eq:43} conserve le profil complet de ce lecteur et ses marques restent dans l’antécédent.

L’expansion spatiale du lecteur HMT et le raffinement nonadique apportent la généalogie et l’évaluation à profondeur arbitraire. Le changement d’échelle entre paramètres se démontre par la dynamique du retour, la transversalité et l’isolement que nous venons de composer. Les preuves analytiques citées interviennent sur la famille déjà construite ; leurs valeurs de comparaison ne sélectionnent ni ses phases ni ses coefficients.

La clôture concerne le sélecteur global initial du secteur uniforme \(\eta=0\), avec les dépendances théorématiques énumérées. La loi asymptotique conserve ses quantificateurs existentiels. La largeur d’une boîte de racines finies mesure la précision d’évaluation ; l’estimation \eqref{hmtfg:escalado:eq:69} exprime la convergence vers la limite.


Les programmes et certificats du supplément de reproduction permettent de reproduire les parties finies de la preuve. La continuation à toute profondeur repose sur les inductions, l’argument uniforme de première sortie et les théorèmes d’identification, avec leurs domaines explicites.

% Fragmento integrable. Preambulo anfitrion: amsmath, amssymb, longtable, hyperref.
\section{Identification analytique des centres initiaux}
\label{hmtfg:fragmento:identificacion-analitica-de-los-centros-originales}

\subsection{Réalisation de degré deux et identification métrique des centres}
\label{hmtfg:clasificacion:8}

L’égalité entre le premier paramètre de l’itinéraire infini et le croisement stable, démontrée par première sortie, fournit
\begin{equation}
\boxed{a_*:=\min\Lambda_\tau=\lambda_*.}
\label{hmtfg:clasificacion:eq:H}
\end{equation}

La famille, le profil dodécaphasique, les retours et le sélecteur restent exactement tels qu’ils sont dans les sections~\ref{hmtfg:clasificacion:1}--\ref{hmtfg:clasificacion:5}. Les résultats externes employés ci-dessous reconnaissent la dynamique de cet objet déjà construit. Leurs valeurs universelles n’interviennent pas dans la sélection des coefficients ni des paramètres de la famille HMT.

\textbf{Théorème de transfert métrique.} L’identification \eqref{hmtfg:clasificacion:eq:H}, jointe aux bornes d’entrée, d’admissibilité et de transversalité démontrées précédemment, implique que la suite initiale des premières racines nouvelles satisfait, pour certaines constantes \(C>0\) et \(0<\theta<1\),
\begin{equation}
\boxed{\lambda_*-\lambda_n=C\delta_{\mathrm F}^{-n}
\bigl(1+O(\theta^n)\bigr),}
\label{hmtfg:clasificacion:eq:14}
\end{equation}
où \(\delta_{\mathrm F}>1\) est la valeur propre instable de l’opérateur de doublement au point fixe reconnu. En particulier,
\begin{equation}
\frac{\lambda_{n-1}-\lambda_{n-2}}
     {\lambda_n-\lambda_{n-1}}
=\delta_{\mathrm F}+O(\theta^n).
\label{hmtfg:clasificacion:eq:15}
\end{equation}
La suite coïncide à partir d’un certain rang avec les centres de la continuation locale, une fois que les deux indices désignent la même période physique. La preuve conserve le sélecteur de première racine de la section~\ref{hmtfg:clasificacion:4}.

\subsubsection{Obtention d’un domaine commun de type quadratique}
\label{hmtfg:clasificacion:8.1}

Notons \(g\) le point fixe de la renormalisation et
\(\Gamma(\lambda)=R^4f_\lambda\). La carte analytique utilisée dans le certificat est
\[
f(x)=1-x^2V\!\left(\frac{x^2-1}{2.5}\right),\qquad
V(z)=\frac{u}{10}+\sum_{k\ge1}\nu_kz^k,
\]
avec la norme des différences \(|\Delta u|+\|\Delta\nu\|_1\). Les inégalités d’admissibilité prouvent que \(g\) est paire, analytique, unimodale, quadratique et de type doublement ; de plus, \(g(x)=\varphi(x^2)\) avec \(\varphi'<0\) dans \([0,1]\).

Le théorème 2.1 de de Faria–de Melo–Pinto, appliqué à l’ensemble de combinatoires constitué uniquement du doublement, fournit un ensemble limite à un seul point avec une extension de type quadratique. Puisque \(Rg=g\), la convergence réelle de ce théorème identifie ce point à \(g\). Le lemme 3.2 fournit, dans un voisinage analytique de \(g\), une itérée finie de renormalisation avec un domaine de type quadratique et un module uniformément positif. Les références précises sont les sections 2.1–2.3, p. 736–738, et la section 3.1, lemme 3.2, p. 745, de \href{https://annals.math.princeton.edu/wp-content/uploads/annals-v164-n3-p01.pdf}{de Faria–de Melo–Pinto, \emph{Global hyperbolicity of renormalization for \(C^r\) unimodal mappings}}.

Pour vérifier l’appartenance au voisinage de ce lemme, soit \(\Omega_\epsilon\) un voisinage complexe suffisamment petit de \([-1,1]\), avec
\[
\rho_\epsilon:=\sup_{x\in\Omega_\epsilon}
\left|\frac{x^2-1}{2.5}\right|<1,
\qquad M_\epsilon:=\sup_{x\in\Omega_\epsilon}|x|^2<\infty.
\]
La restriction de la carte certifiée satisfait
\begin{equation}
\|\Delta f\|_{\Omega_\epsilon}
\le M_\epsilon\left(\frac{|\Delta u|}{10}
          +\rho_\epsilon\|\Delta\nu\|_1\right)
\le C_\epsilon\|\Delta f\|_{\mathcal A}.
\label{hmtfg:clasificacion:eq:16}
\end{equation}
Le certificat donne
\[
\|R^m\Gamma(\lambda_*)-g\|_{\mathcal A}
<0.001666(0.83996)^m.
\]
On choisit donc un entier fixe \(m\) qui place cette image dans le voisinage du lemme 3.2. La continuité de la composition finie fournit un intervalle \(I_*\) autour de \(\lambda_*\) où les mêmes opérations sont définies. Si \(N\) est l’itérée fournie par ce lemme, nous posons
\begin{equation}
d=4+m+N,\qquad
\mathcal F_\lambda=R^df_\lambda,\quad \lambda\in I_*.
\label{hmtfg:clasificacion:eq:17}
\end{equation}
La famille \(\mathcal F\) possède des représentants de type quadratique de degré deux, sur un domaine commun et avec un module uniformément positif. Sa dépendance analytique réelle admet une complexification locale du paramètre. Le nombre \(d\) est fixe et compte les doublements effectués avant d’appliquer le résultat paramétrique.

Cette construction fixe un domaine commun et un nombre fini de retours \(d\), suffisants pour appliquer le théorème paramétrique.

\subsubsection{Transport de la transversalité certifiée}
\label{hmtfg:clasificacion:8.2}

La direction de \(\Gamma\) satisfait la condition de cône
\[
\|\dot\nu\|_1\le\tfrac14|\dot u|,\qquad \dot u\ne0.
\]
Les bornes différentielles certifiées impliquent l’invariance de ce cône et
\begin{equation}
|\dot u_{k+1}|\ge4.381|\dot u_k|
\label{hmtfg:clasificacion:eq:18}
\end{equation}
le long de l’orbite de \(\Gamma(\lambda_*)\). Tous ses points restent dans le domaine certifié. En particulier, la direction transportée par les étapes finies de \eqref{hmtfg:clasificacion:eq:17} conserve une composante expansive non nulle.

La classe hybride du point fixe coïncide avec sa variété stable dans l’espace des germes de type quadratique : \href{https://www.maths.tcd.ie/EMIS/journals/Annals/149_2/lyubich.pdf}{Lyubich, \emph{Feigenbaum–Coullet–Tresser universality and Milnor's hairiness conjecture}, théorème 6.1, p. 371–372}. La transversalité reste valable lors du passage à cet espace, comme on peut le vérifier sans identifier des normes de Banach différentes. En effet, si \(\partial_\lambda\mathcal F_{\lambda_*}\) était tangente à la classe hybride, elle serait la tangente d’un disque analytique contenu dans cette classe. La convergence exponentielle uniforme de la renormalisation sur un sous-disque compact, suivie de l’estimation de Cauchy en son paramètre, ferait converger vers zéro l’évaluation en \(x=1\) de ses tangentes renormalisées. Cependant, dans la carte précédente,
\[
\dot f(1)=-\dot u/10,
\]
et \eqref{hmtfg:clasificacion:eq:18} fait croître sa valeur absolue. Les deux dérivées représentent le même germe et ont la même évaluation réelle. La contradiction démontre
\begin{equation}
\partial_\lambda\mathcal F_{\lambda_*}
\notin T_{\mathcal F_{\lambda_*}}\mathcal H_{c_\infty}.
\label{hmtfg:clasificacion:eq:19}
\end{equation}
Ainsi, \(S=\{\mathcal F_\lambda\}\) est une transversale analytique complexe à la classe hybride de la limite de doublement.

\subsubsection{Identification à partir d’un certain rang des premières racines par leur période}
\label{hmtfg:clasificacion:8.3}

Notons \(c_r\) le centre quadratique canonique de période \(2^r\), obtenu par \(r\) doublements successifs du centre de période un. Ce symbole désigne le paramètre de la famille quadratique de reconnaissance ; il se distingue des fonctions orbitales \(c_j(\lambda)=f_\lambda^j(0)\) des sections précédentes. On écrit \(c_r^{\rm quad}\) pour conserver cette distinction.

Le théorème de continuation globale et \eqref{hmtfg:clasificacion:eq:H} donnent \(\lambda_n\to\lambda_*\), d’où \(\lambda_n\in I_*\) pour \(n\) suffisamment grand. La classification et les intervalles restrictifs de la section~\ref{hmtfg:clasificacion:2} montrent que, lorsque \(n>d\),
\begin{equation}
\mathcal F_{\lambda_n}=R^df_{\lambda_n}
\quad\text{a une période critique exacte }2^{n-d}
\quad\text{et itinéraire }W_{n-d}0.
\label{hmtfg:clasificacion:eq:20}
\end{equation}
Son orbite critique reste dans l’intervalle réel invariant, contenu dans le domaine de type quadratique ; son ensemble de Julia rempli est donc connexe. Le redressement conserve l’orbite postcritique et les branches réelles. En parcourant ses retours restrictifs, on obtient \(n-d\) doublements canoniques et, à la fin, un point critique fixe. Ce dernier se redresse au centre \(0\) ; les inverses successifs du redressement des copies de doublement donnent l’unique centre \(c_{n-d}^{\rm quad}\). Il s’agit de l’application concrète du paramétrage par combinatoire et de l’homéomorphisme de redressement des copies, décrits dans Lyubich, sections 5.1–5.3, p. 362–364. On conclut
\begin{equation}
\mathcal F_{\lambda_n}\in\mathcal H_{c_{n-d}^{\rm quad}}.
\label{hmtfg:clasificacion:eq:21}
\end{equation}

Le théorème 7.4 de Lyubich, p. 387–388, donne une unique intersection locale de \(S\) avec chaque classe \(\mathcal H_{c_r^{\rm quad}}\) pour \(r\) suffisamment grand. Ses hypothèses de bornes complexes a priori sont satisfaites pour la cascade quadratique réelle de doublement, conformément au théorème 5.6, p. 367, du même travail. Notons \(\zeta_r\) le paramètre de cette intersection. Les équations \eqref{hmtfg:clasificacion:eq:20}–\eqref{hmtfg:clasificacion:eq:21}, la convergence vers le croisement et l’unicité locale impliquent
\begin{equation}
\boxed{\lambda_n=\zeta_{n-d}\quad\text{pour tout }n\text{ suffisamment grand}.}
\label{hmtfg:clasificacion:eq:22}
\end{equation}
L’unicité utilisée ici concerne \textbf{chaque classe hybride canonique} dans un voisinage fixe du croisement. Elle ne se déduit ni de la seule existence d’une cascade locale ni de l’unicité du croisement avec la feuille stable. L’équation \eqref{hmtfg:clasificacion:eq:22} identifie le sélecteur initial, déjà défini globalement, aux intersections locales ; elle n’introduit pas de sélecteur de substitution.

En particulier, soit \(\mu_j\) la cascade locale de la section \ref{hmtfg:escalado:8}, et soit \(s\) l’entier fixe déterminé par sa période physique exacte :
\[
\operatorname{per}_{\rm crit}(f_{\mu_j})=2^{j+s}.
\]
Dans la présentation avec la courbe \(R^4f_\lambda\) et la cible de période deux, \(s=5\) ; les étapes fixes supplémentaires utilisées pour localiser la cible modifient \(s\). Par \eqref{hmtfg:clasificacion:eq:21} et la même unicité,
\begin{equation}
\mu_j=\zeta_{j+s-d},\qquad
\boxed{\lambda_n=\mu_{n-s}\quad(n\text{ suffisamment grand}).}
\label{hmtfg:clasificacion:eq:23}
\end{equation}
Le décalage \(d\) de la préparation de type quadratique et le décalage \(s\) des périodes de la cible sont des objets distincts. L’identification s’effectue par la période physique, et non en égalant ces deux entiers.

\subsubsection{Reste en puissance par holonomie transversale}
\label{hmtfg:clasificacion:8.4}

Le lemme 7.3 de Lyubich, p. 384, avec démonstration p. 384–387, apporte une régularité \(C^{1+\beta}\)-conforme, pour un certain \(\beta>0\), de l’holonomie entre transversales à la classe hybride du point de Feigenbaum. Sa définition p. 384 concerne les lieux de connexité sur les transversales ; ils contiennent notamment les centres considérés ici. Nous pouvons réduire l’exposant de sorte que \(0<\beta\le1\).

Soit \(h\) l’holonomie de \(S\) vers la variété instable et soit \(z\) une coordonnée analytique qui linéarise la restriction unidimensionnelle de la renormalisation :
\begin{equation}
z(g)=0,\qquad z(Rq)=\delta_{\mathrm F} z(q).
\label{hmtfg:clasificacion:eq:24}
\end{equation}
L’existence de cette coordonnée s’obtient en appliquant la linéarisation d’un germe holomorphe répulsif, de multiplicateur \(\delta_{\mathrm F}>1\), à la variété instable analytique. La direction expansive de \eqref{hmtfg:clasificacion:eq:18}, l’invariance du cône et le retour stationnaire identifient ce multiplicateur à la valeur propre instable positive employée dans le certificat.

Si \(q_r\) est le centre de la classe \(\mathcal H_{c_r^{\rm quad}}\) situé sur la variété instable, la compatibilité du redressement et de la renormalisation donne \(Rq_r=q_{r-1}\). Cette identité est celle utilisée dans la démonstration du théorème 7.4, p. 388. Par \eqref{hmtfg:clasificacion:eq:24}, il existe \(b\ne0\), qui absorbe tout indice initial fixe, tel que
\begin{equation}
z(q_r)=b\delta_{\mathrm F}^{-r}
\label{hmtfg:clasificacion:eq:25}
\end{equation}
pour tout \(r\) suffisamment grand.

La régularité de \(h\) et la transversalité \eqref{hmtfg:clasificacion:eq:19} donnent, sur le lieu de connexité proche de \(\lambda_*\),
\begin{equation}
H(\lambda):=z\bigl(h(\mathcal F_\lambda)\bigr)
=A(\lambda-\lambda_*)
+O\!\left(|\lambda-\lambda_*|^{1+\beta}\right),
\qquad A\ne0.
\label{hmtfg:clasificacion:eq:26}
\end{equation}
Puisque \(H(\zeta_r)=b\delta_{\mathrm F}^{-r}\), l’estimation inférieure
\(|H(\lambda)|\ge |A|\,|\lambda-\lambda_*|/2\) dans un voisinage suffisamment petit implique d’abord
\(|\zeta_r-\lambda_*|=O(\delta_{\mathrm F}^{-r})\). En substituant cette borne dans le reste de \eqref{hmtfg:clasificacion:eq:26}, on obtient
\begin{equation}
\zeta_r-\lambda_*=\frac bA\delta_{\mathrm F}^{-r}
+O\!\left(\delta_{\mathrm F}^{-(1+\beta)r}\right).
\label{hmtfg:clasificacion:eq:27}
\end{equation}
Nous appliquons \eqref{hmtfg:clasificacion:eq:22} avec \(r=n-d\). Le décalage fini modifie le coefficient principal et la constante du reste. Puisque \(\lambda_n<\lambda_*\), on obtient \eqref{hmtfg:clasificacion:eq:14}, avec
\[
C=-\frac bA\delta_{\mathrm F}^d>0,
\qquad \theta=\delta_{\mathrm F}^{-\beta}\in(0,1).
\]
Enfin,
\[
\lambda_n-\lambda_{n-1}
=C(\delta_{\mathrm F}-1)\delta_{\mathrm F}^{-n}\bigl(1+O(\theta^n)\bigr),
\]
ce qui démontre \eqref{hmtfg:clasificacion:eq:15}. \(\square\)


\subsubsection{Portée quantitative du théorème}
\label{hmtfg:clasificacion:8.5}

L’identification des limites démontrée par première sortie active le transfert métrique pour les premières racines globales du secteur uniforme. Le taux opératoriel \(0.83996\) contrôle les applications renormalisées sur la feuille stable. Le taux paramétrique \(\theta=\delta_{\mathrm F}^{-\beta}\) découle de la régularité de l’holonomie transversale. Les deux estimations conservent leur variable et leur domaine.

Le théorème établit l’existence de \(\beta,\theta,C\), du voisinage \(I_*\), du nombre fini de retours \(d\) et d’un indice \(n_0\) à partir duquel les centres coïncident. Leurs valeurs numériques individuelles ne sont pas attribuées dans cet énoncé asymptotique. Cette formulation complète la limite et son reste en puissance ; l’évaluation numérique de ces témoins constitue une spécialisation quantitative distincte.

\section{Retour, mémoire et universalité de la cascade dodécaphasique}
\label{sec:hmt-feig-interpretacion}

\subsection{La structure génératrice et la lecture critique}

La construction de Feigenbaum en Holographie Modulaire Triadique relie une
organisation arithmétique orientée à une loi asymptotique de changement d’échelle.
Son point de départ est la composition effective d’APP, TRIT et TPK. APP évalue
la somme et le produit sur les deux feuilles du réseau de chiffres et conserve les
quotients et les restes de ces opérations. TRIT détermine le régime local et
l’orientation de sa lecture. TPK sélectionne et transporte les états,
actualise leurs registres et prolonge leurs routes avec retenue, frontière et mémoire.
La structure discrète conjointe du continu fournit la résidence commune
de ces états et des lecteurs qui agissent sur eux.

Le prolongement
\[
w_6\longrightarrow w_{12}\longrightarrow w_{18}
\longrightarrow w_{24}\longrightarrow w_{30}
\longrightarrow R_{36}\longrightarrow G_9
\]
organise cette persistance. La graine initiale est relevée en conservant ses
préfixes ; les relèvements successifs changent le régime sans perdre l’orientation
ni la mémoire ; la frontière coinductive permet de poursuivre la construction, et les
neuf cartes réunissent ses lectures locales. L’holonomie nonadique exprime le
retour de phase accompagné d’un accroissement de mémoire. Ainsi, un tour
visible conserve sa position dans l’histoire complète. L’opération de retour
acquiert dès l’origine deux composantes inséparables :
réitération d’une règle et conservation de sa provenance.

Le lecteur critique fixé sur la roue dodécaphasique orientée convertit cette
structure en douze phases avec une trace uniforme. Leurs représentants sont
\[
\phi_m=\frac{(3m-1)\pi_{\mathrm{HMT}}}{18},
\qquad w_m=\frac1{12},\qquad 1\le m\le12.
\]
Le choix du représentant orienté importe parce que le deuxième moment
utilise sa grandeur. Le quotient angulaire enregistre la phase, tandis que le
représentant conserve l’information nécessaire à sa normalisation. La
coordonnée \(\pi_{\mathrm{HMT}}\) provient des lecteurs coinductifs du même
noyau APP--TRIT--TPK. Dans son image corrélée, les coordonnées
\(\pi,\varphi,e\) et la clôture dodécaphasique de \(\alpha_{\mathrm{HMT}}\)
conservent leur ordre constructif et leurs domaines propres. Dans le lecteur critique
considéré ici, on utilise la coordonnée angulaire déjà générée.

La normalisation est déterminée par une somme entière :
\[
\sum_{m=1}^{12}(3m-1)^2=5394=6\cdot899,
\qquad k_m=\frac{2(3m-1)}{\sqrt{899}}.
\]
Le facteur commun \(\pi_{\mathrm{HMT}}\) se simplifie lors de la fixation du deuxième moment.
Cette simplification conserve la provenance des phases et produit une
représentation particulièrement précise de leur lecture normalisée. On obtient
\[
u_0(x)=\frac1{12}\sum_{m=1}^{12}\bigl(1-\cos(k_mx)\bigr),
\qquad f_\lambda(x)=1-\lambda u_0(x),
\qquad \sum_mw_mk_m^2=2.
\]
Le profil est pair et entier ; il satisfait \(u_0(0)=u_0'(0)=0\) et
\(u_0''(0)=2\). Pour \(0<\lambda\le2/u_0(1)\), la famille envoie
\([-1,1]\) dans lui-même et possède un maximum critique quadratique avec
\(f_\lambda''(0)=-2\lambda\). Le paramètre \(\lambda\) parcourt cette famille
après fixation du lecteur, de l’orientation, des poids et de l’échelle. Le nombre
universel qui sera reconnu à la fin est ainsi séparé de la sélection
des coefficients.

\subsection{Douze phases, moments et représentation spectrale}

Le profil admet une deuxième expression qui réunit ses phases dans un opérateur
fini. Les nœuds
\[
s_m=k_m^2=\frac{4(3m-1)^2}{899}
\]
déterminent la mesure positive \(\mu_0=12^{-1}\sum_m\delta_{s_m}\). Ses
moments \(M_r=\int s^r\,d\mu_0(s)\) définissent des formes de Hankel strictement
positives jusqu’au rang douze. Cette positivité provient de l’évaluation
d’un polynôme sur douze nœuds distincts : un polynôme de degré inférieur à
douze conserve au moins une évaluation non nulle. Au degré douze, le produit
des douze facteurs nodaux annule la norme et détermine exactement le rang
du lecteur.

L’orthogonalisation de ces moments produit une matrice de Jacobi
\(J_{12}\), réelle symétrique, à coefficients sous-diagonaux positifs. Son calcul fonctionnel
restitue le profil complet :
\[
u_0(x)=e_0^{\mathsf T}
\bigl[I-\cos(x\sqrt{J_{12}})\bigr]e_0.
\]
La somme de cosinus et l’opérateur de Jacobi sont deux réalisations du même
lecteur. Les moments maintiennent toutes les fréquences et leurs poids ; les
marques d’orientation, de route et de mémoire restent dans l’état qui les
engendre. Le rang douze identifie la dimension spectrale de cette lecture
particulière. La profondeur de ses prolongements relève d’une autre dimension
de la construction : la longueur des histoires compatibles et la résolution
avec laquelle elles sont évaluées.

La géométrie analytique est elle aussi explicitée. Pour la déformation
pondérée \(u_\eta\), dans la bande \(|\eta|\le1/40\), les poids restent
positifs et le deuxième moment demeure normalisé. Dans
\(|\operatorname{Re}z|<\pi/3\), il existe une fonction holomorphe impaire et univalente
\(b_\eta\), normalisée par \(b_\eta'(0)=1\), telle que
\(u_\eta=b_\eta^2\). La roue produit donc une déformation conforme
du pli quadratique. Sa conjugaison dynamique conserve cette déformation
à travers \(b_\eta(1-\lambda w^2)\). Le théorème de changement d’échelle paramétrique
développé ici concerne le membre uniforme \(\eta=0\) ; les
propriétés structurelles de la bande pondérée et la persistance locale de
ses racines conservent leurs quantificateurs spécifiques.

\subsection{Profondeur de lecture et restitution des histoires}

La profondeur nonadique affine la lecture d’un paramètre, tandis que le
doublement dynamique change l’horizon de retour. Le raffinement avec
retenue exprime la première opération par
\[
L_c(x,y)=(Bx-c,By+c),\qquad
C_m=\sum_{j=1}^m c_jB^{m-j},
\]
\[
(x_m,y_m)=(B^mx_0-C_m,B^my_0+C_m).
\]
Le support extérieur croît comme \(B^m\), et le cylindre normalisé du registre
a pour diamètre \(B^{-m}\). La paire entière, la profondeur et la marque de
frontière conservent la reconstruction des quotients, des restes et des préfixes.
Croissance extérieure et résolution intérieure décrivent ainsi des opérations
compatibles sur une même histoire.

La composition exacte
\(L_D^{[b]}L_C^{[a]}=L_{B^bC+D}^{[a+b]}\) permet de regrouper une histoire en
blocs sans altérer ses données. Un retour de doublement réunit deux
transitions ; le calendrier nonadique en réunit neuf. Les deux regroupements sont
comparés sur un horizon commun de \(9\,2^N\) événements, en conservant
l’ordre des sous-blocs et leurs retenues.

La contribution de la mémoire peut s’observer directement dans les branches
inverses du profil. Si \(v=(u_0|_{[0,1]})^{-1}\), les deux branches sont
\[
\beta_\sigma(y)=\sigma v\!\left(\frac{1-y}{\lambda}\right),
\qquad \sigma\in\{-1,+1\}.
\]
La valeur finale, jointe à la dernière feuille, reconstruit la valeur précédente. Une
suite de feuilles permet de restituer toutes les étapes intermédiaires. Le mot
des branches accompagne les registres APP--TRIT--TPK de route, de frontière, de retenue et de
mémoire ; chaque registre conserve sa fonction. Cette restitution concrète
explique quelle information exige l’inversion de la lecture critique.

La résolution de chaque retour dispose également d’une borne explicite. Avec
\(t=\lambda u_\eta(1)/2\), \(z=(x+1)/2\) et
\[
F_{t,\eta}(z)=1-t\frac{u_\eta(2z-1)}{u_\eta(1)},
\]
on obtient
\[
\left|F_{t,\eta}^{q}(1/2)-F_{s,\eta}^{q}(1/2)\right|
\le S_q|t-s|,\qquad S_q=\sum_{j=0}^{q-1}(6\sqrt2)^j<9^q.
\]
Un cylindre de base neuf et de profondeur \(m=2^N+r\) permet donc de lire
le retour d’horizon \(2^N\) avec un diamètre inférieur à \(9^{-r}\). Cette
estimation donne un contenu quantitatif à la profondeur arbitraire : pour chaque
horizon et chaque précision demandés, elle détermine une profondeur suffisante.
L’évaluation des fonctions analytiques et l’arrondi possèdent leurs propres
restes, qui se propagent avec l’incertitude du paramètre.

Les vacances interviennent dans le bilan de capacité du registre. La
comparaison entre blocs de 729 et 1000 possibilités distingue la capacité
accumulée de la longueur chronologique, en conservant les étapes où la première
reste constante. Cette comptabilité détermine l’espace requis pour
conserver une lecture et sa profondeur. L’admissibilité des histoires
découle, quant à elle, des règles de survie. Les deux opérations
collaborent à la construction et maintiennent des objets mathématiques distincts.

\subsection{Chronologie nonadique et combinatoire de doublement}

Le mot critique de doublement possède une expression à toute profondeur :
\[
\tau_j=(-1)^{v_2(j)},\qquad
\tau_{2j}=-\tau_j,\qquad \tau_{2j+1}=+1.
\]
Ses préfixes satisfont \(W_{n+1}=W_n(-1)^nW_n\). La règle détermine comment
les deux copies du retour sont orientées et où apparaît leur séparation. La
lecture nonadique de la chronologie conserve l’entier complet
\(K=\rho_9(K)+9q_9(K)\). Lorsqu’on le réduit modulo \(2^n\), l’holonomie qui
avance de neuf événements induit la translation de neuf sur ce cycle. Puisque
neuf est impair, cette translation visite toutes ses positions. La
permutation \(j\mapsto9j\pmod{2^n}\), qui la conjugue à l’avancée unitaire,
conserve la valuation dyadique de chaque position non nulle. Le reste et le quotient
chronologiques soutiennent conjointement cette compatibilité.

La mémoire bilatérale possède également une réalisation opératorielle compatible
avec le raffinement. Sur les cycles de longueur \(2^n\), l’opérateur
d’avancée \(C_n\) est incorporé dans
\[
\mathbb U_n=P_0\otimes I+P_1\otimes C_n,
\]
avec les projecteurs complémentaires du lecteur nonadique. À la limite, ses
composantes visible et complémentaire satisfont
\[
\|T_{\infty,a}v\|^2+\|D_{\infty,a}v\|^2=\|v\|^2,
\qquad
v=T_{\infty,a}^*T_{\infty,a}v+D_{\infty,a}^*D_{\infty,a}v.
\]
Ces identités décrivent la conservation et la restitution dans l’espace
et la représentation déclarés. La contraction des applications renormalisées,
qui apparaîtra ensuite, mesure une autre opération : le rapprochement de leurs profils
vers une forme stable. L’histoire conservée par le relèvement et la
convergence d’une lecture fonctionnelle s’articulent à des niveaux différents.

\subsection{La carte ordonnée et la sélection globale des racines}

La lecture archimédienne du continu HMT transporte les opérations
par un isomorphisme de corps ordonnés complets, dans l’univers
déclaré \(\mathfrak G\) :
\[
\iota:\mathbb R_{\mathrm{HMT}}^{\mathfrak G}
\longrightarrow\mathbb R^{\mathfrak G}.
\]
La construction interne du cosinus par sa série, de la racine positive de 899,
du profil et de ses itérées précède la sélection des zéros. La conservation
de la somme, du produit, de l’ordre et des limites convergentes donne
\[
\iota\bigl(S_n^{\mathrm H}(a)\bigr)=S_n\bigl(\iota(a)\bigr),
\qquad S_n(\lambda)=f_\lambda^{2^n}(0).
\]
La surjectivité couvre toutes les racines du domaine de cette carte ; l’ordre
conserve les retours antérieurs, la période exacte et la condition de
se trouver à droite de la racine précédente. On transporte ainsi
également le minimum de chaque ensemble de candidats. La généalogie atteint
la sélection effective du paramètre, la deuxième projection et ses
marques étant conservées sur le même antécédent.

Le sélecteur est fixé par \(\lambda_1=1/u_0(1)\) et par la première racine
nouvelle de période critique exacte \(2^n\) située à droite de
\(\lambda_{n-1}\). La classification des itinéraires démontre que chaque racine
sélectionnée a pour mot \(W_n0\). Les retours restrictifs réduisent sa
période de moitié en conservant l’ordre unimodal, et l’induction restitue
le mot complet. L’analyticité des retours non identiquement nuls
isole leurs zéros dans chaque compact ; la
comparaison des signes avant le premier paramètre d’itinéraire infini
garantit l’existence d’une racine nouvelle à tous les niveaux.

Soit \(\Lambda_\tau\) l’ensemble compact non vide des paramètres qui réalisent
le mot infini, et soit \(a_*\) son minimum. On démontre
\[
\lambda_n\nearrow a_*.
\]
Le passage à la limite conserve chaque signe au moyen d’une borne uniforme à horizon
fixe. Si \(c_p=f_\lambda^p(0)\), les signes opposés de \(c_p\) et
\(c_{2p}\), avec \((f_\lambda^p)'(0)=0\), impliquent
\[
|c_p|>\frac2{B_p},\qquad
B_p=16^p\frac{16^p-1}{15}.
\]
Ainsi, les préfixes de périodes croissantes préservent le mot infini en
s’accumulant. La preuve utilise la suite initiale des premières racines dès
sa définition globale.

\subsection{La feuille stable et l’identification de la limite sélectionnée}

Le retour normalisé
\[
(Rf)(x)=-\frac{f(f(-ax))}{a},\qquad a=-f(1)>0,
\]
réunit deux transitions, change l’échelle et conserve l’orientation nécessaire
pour restituer un maximum de valeur un. Les domaines restrictifs établissent
son interprétation comme retour réel. La dérivée de cet opérateur inclut
la variation de l’échelle \(a\) ; cette contribution détermine la sensibilité
paramétrique de la transformation complète.

Dans la carte \(f(x)=1-x^2V((x^2-1)/(5/2))\), de coordonnées \(u=10v_0\)
et \(\nu=(v_1,v_2,\ldots)\), le quatrième retour de la famille entre dans un
voisinage quantifié du point fixe \(g\). La courbe traverse transversalement
sa feuille stable en un unique paramètre local
\[
1.874038<\lambda_*<1.874039,
\]
et satisfait
\[
\|R^{4+n}f_{\lambda_*}-g\|_{\mathcal A}
<0.001666(0.83996)^n.
\]
La feuille stable est un graphe de pente au plus \(0.16\). Les
sécantes paramétriques appartiennent à un cône transversal expansif. Les deux
propriétés distinguent la diminution des déformations stables de la
croissance de la séparation entre paramètres.

L’identification de \(a_*\) et de \(\lambda_*\) réunit cette géométrie locale
et le minimum global. On exclut d’abord l’intervalle entre la huitième racine
et l’extrémité inférieure de la boîte locale. Ensuite, l’ordre
\(a_*\le\lambda_*\) permet d’étudier une séparation hypothétique au moyen
du cône négatif. Tant que sa coordonnée transversale a une valeur absolue inférieure
à \(0.006\), la sécante reste dans le domaine analytique et croît d’un
facteur compris entre \(4.4034\) et \(8.0887\). Une séparation persistante
produirait donc une première sortie d’amplitude comprise entre
\(0.006\) et \(0.0485322\).

Le résidu de l’orbite stable conserve une information plus précise que sa
norme : ses composantes satisfont \(|e_u|\le0.16\|e_\nu\|_1\). Cette
relation permet de borner conjointement leurs effets sur l’orbite critique.
Le recouvrement de toute la région de première sortie montre dans chaque cas un
signe incompatible avec \(\tau\), ou un retour contractant qui impose à deux
valeurs critiques dyadiques successives d’avoir le même signe. Le mot
infini exige des signes opposés. La contradiction identifie les paramètres :
\[
\lambda_n\nearrow a_*=\lambda_*.
\]
La démonstration combine une induction valable à toute profondeur avec une
exclusion finie uniforme de la région qu’atteindrait toute
séparation. La portée infinie repose sur cette composition.

\subsection{Changement d’échelle universel et signification de la précision}

La réalisation holomorphe de degré deux des retours et la transversalité
établie permettent de comparer leurs classes hybrides. Pour chaque période
suffisamment élevée, la classe canonique possède une unique intersection locale
avec la famille. Les premières racines globales convergent déjà vers le croisement et
conservent la combinatoire de doublement ; elles coïncident donc à partir d’un certain rang
avec les centres locaux lorsque leurs périodes physiques sont égalées. La preuve
transporte la sélection initiale jusqu’à la loi métrique universelle.

Il existe \(C>0\), \(M<\infty\) et \(0<\theta<1\) tels que
\[
\lambda_*-\lambda_n=C\delta_{\mathrm F}^{-n}(1+e_n),
\qquad |e_n|\le M\theta^n,
\]
et, par conséquent,
\[
\lim_{n\to\infty}
\frac{\lambda_{n-1}-\lambda_{n-2}}
{\lambda_n-\lambda_{n-1}}=\delta_{\mathrm F}.
\]
Le facteur \(\delta_{\mathrm F}\) est la valeur propre expansive de l’opérateur de doublement
reconnu à la fin de la construction. La famille particulière conserve ses
douze phases, ses moments et son histoire de provenance ; la limite de ses
rapports paramétriques appartient à la même classe universelle de criticité
quadratique. Cette relation explique conjointement la spécificité du
générateur et l’universalité de la sortie.

La précision possède ici trois résidences. Les intervalles des racines
contrôlent leur évaluation finie. L’estimation
\(0.001666(0.83996)^n\) contrôle la convergence opératorielle sur la feuille
stable. Le terme \(M\theta^n\) contrôle le transfert de cette géométrie à
l’échelle paramétrique au moyen de l’holonomie transversale. Chaque estimation
correspond à sa variable et à son opération. Le huitième quotient est encadré
autour de \(4.669212189150204154556096215889663928\), avec une largeur inférieure
à \(5.808\cdot10^{-37}\) ; cette largeur mesure l’évaluation de \(\delta_{\mathrm F,8}\).
La distance de \(\delta_{\mathrm F,8}\) à la limite exige de quantifier les constantes
du reste paramétrique, dont l’existence est déjà démontrée.

La notation maintient séparés \(\delta_{\mathrm F}\), facteur de changement d’échelle paramétrique
de Feigenbaum ; \(\alpha_{\mathrm F}\), constante spatiale associée à son
problème de renormalisation ; et \(\alpha_{\mathrm{HMT}}\), constante de
structure fine du corpus. Le présent résultat identifie \(\delta_{\mathrm F}\)
au moyen des premières racines du secteur uniforme \(\eta=0\). La bande
pondérée conserve ses théorèmes structurels et d’existence, et son extension
métrique quantifiée constitue un énoncé spécifique distinct.

L’apport de l’enchaînement réside dans la continuité entre génération,
lecture, mémoire, sélection et limite. L’arithmétique orientée détermine le
profil ; la représentation spectrale conserve ses douze composantes ; les
cylindres permettent de le résoudre avec une précision arbitraire à chaque horizon ;
le transport des histoires restitue ses parcours ; la carte ordonnée
conserve toutes les racines et leurs minima ; la classification prolonge le
sélecteur ; et l’expansion transversale, jointe à la contraction stable,
détermine sa loi asymptotique. L’universalité apparaît comme une propriété
démontrée de cette construction complète et de sa réalisation analytique.

\endgroup

% Fragmento compartido por VII y CRISTAL; incluir despues del cuerpo Feigenbaum.
% Procedencia: resultados recuperados, con demostraciones elementales reunidas.
% Perfil, Jacobi y escalado: 00_antecedentes_operatorios.tex, 03, 06 y 07
% del presente paquete; no se cambia su selector ni su dominio uniforme.
% Constitutiva: CRISTAL_ES/main_autosuficiente.tex:58738--58840,58864--58894.
% Inversion: CRISTAL_ES/sections/cr28_restitucion_constitutiva.tex:18--59.
% Catalan: el mismo propietario, 277--373; identidad integral, unicidad y serie.
% Corte material: CUMPLIMIENTO_EDITORIAL_20260928/fuentes/CRISTAL_ES.
% No contiene una nueva calibracion ni una nueva campana de calculo.
\section{Articulation des lecteurs critique, spectral et constitutif}
\label{hmtfg:articulacion:seccion}

La composition APP--TRIT--TPK précède les lecteurs de cette section.
APP conserve dans chaque cellule les évaluations additive et multiplicative,
avec leurs restes et leurs quotients ; la réduction tritique détermine le régime,
l’orientation et la retenue de la transition. La composition effective
de la sélection, du transport et de l’actualisation de la mémoire produit l’état
enrichi et ses prolongements, développés dans
\ref{hmtfg:app-trit-transporte} et
\ref{hmtfg:subsubsec:tres-vueltas-ext-ch}.
L’holonomie nonadique compose ces transitions : le retour de phase
fait avancer la mémoire et conserve l’histoire. Les lecteurs agissent sur la
même structure discrète conjointe du continu, avec ses cinq
constructions consubstantielles et les deux lectures du même
antécédent : évaluation ordonnée et incidence de frontière.

La roue orientée et les canaux angulaires sont des sorties de cette
composition. Leurs coordonnées HMT déjà générées sont utilisées dans des opérations
ultérieures, avec leurs marques et leurs domaines ; ce ne sont pas des nombres métrologiques
introduits pour sélectionner une réponse. La lecture ordonnée
permet d’évaluer les formules suivantes, tandis que l’état de provenance
conserve la feuille, la route, le reste, le quotient, la retenue, la frontière et la mémoire.
Les restitutions démontrées retrouvent l’objet de leur domaine
déclaré ; la représentation évaluée reste accompagnée de ces
registres.

\subsection{La séparation critique comme lecture asymptotique intrinsèque}
\label{hmtfg:articulacion:separacion}

Dans le secteur uniforme de la roue dodécaphasique, les représentants
orientés et leurs poids sont
\[
 \phi_m=\frac{(3m-1)\pi_{\mathrm{HMT}}}{18},\qquad
 w_m=\frac1{12},\qquad 1\le m\le12.
\]
Le deuxième moment est normalisé par
\(\sum_{m=1}^{12}(3m-1)^2=5394=6\cdot899\).
Ainsi, le facteur angulaire commun se simplifie et il reste
\begin{equation}
 k_m=\frac{2(3m-1)}{\sqrt{899}},\qquad
 u_0(z)=\frac1{12}\sum_{m=1}^{12}\bigl(1-\cos(k_mz)\bigr),
 \qquad f_\lambda(z)=1-\lambda u_0(z).
 \label{hmtfg:articulacion:perfil}
\end{equation}
La variable \(\lambda\) parcourt la famille après fixation du
lecteur, de ses poids et de la normalisation. Dans
\(0<\lambda\le2/u_0(1)\), le domaine dynamique est \([-1,1]\).
La première racine est \(\lambda_1=1/u_0(1)\) ; pour \(n\ge2\),
\(\lambda_n\) est la plus petite racine supérieure à \(\lambda_{n-1}\)
dont le point critique a pour période exacte \(2^n\).
La classification, l’existence de chaque minimum et l’identification
métrique de ces centres ont été démontrées dans
\ref{hmtfg:escalado:15}--\ref{hmtfg:escalado:17}.

\begin{proposition}[Lecture structurelle de la constante de changement d’échelle]
\label{hmtfg:articulacion:delta}
Pour le lecteur uniforme \eqref{hmtfg:articulacion:perfil} et son sélecteur
de premières racines, la constante de séparation critique est
\begin{equation}
 \delta_{\mathrm{HMT,crit}}
 :=\lim_{n\to\infty}
 \frac{\lambda_n-\lambda_{n-1}}{\lambda_{n+1}-\lambda_n}
 =\delta_{\mathrm F}.
 \label{hmtfg:articulacion:limite}
\end{equation}
C’est l’invariant asymptotique de séparation des paramètres des retours
critiques de ce lecteur, dont l’ordre de sélection est conservé.
\end{proposition}
\begin{proof}
Le théorème de changement d’échelle \eqref{hmtfg:escalado:eq:67} fournit
\(\lambda_*-\lambda_n=C\delta_{\mathrm F}^{-n}(1+e_n)\),
avec \(C>0\) et \(e_n\to0\). En soustrayant deux termes consécutifs,
\[
 \lambda_n-\lambda_{n-1}
 =C\delta_{\mathrm F}^{-n}
   \bigl[\delta_{\mathrm F}-1+
          \delta_{\mathrm F}e_{n-1}-e_n\bigr].
\]
Le quotient de cette égalité et de celle d’indice \(n+1\) converge vers
\(\delta_{\mathrm F}\). L’identité à partir d’un certain rang des centres locaux
et des premières racines, démontrée dans
\eqref{hmtfg:escalado:eq:66}, permet d’appliquer le changement d’échelle précisément
à la suite sélectionnée, et non à une autre cascade choisie après coup.
\end{proof}

Cette formulation conserve plus d’information que le numéral terminal :
elle identifie la roue, sa lecture, la famille et le sélecteur dont la limite
est évaluée. Son interprétation est intrinsèque à cette composition HMT ;
la reconnaissance de la constante universelle utilise ensuite le théorème
analytique appliqué au lecteur déjà construit.

\subsection{La mesure finie et la restitution spectrale du profil}
\label{hmtfg:articulacion:jacobi}

Les fréquences au carré produisent la mesure de probabilité atomique
\begin{equation}
 s_m=k_m^2=\frac{4(3m-1)^2}{899},\qquad
 \nu_0=\frac1{12}\sum_{m=1}^{12}\delta_{s_m},\qquad
 M_j=\int s^j\,d\nu_0(s).
 \label{hmtfg:articulacion:medida-finita}
\end{equation}
Ici, \(\delta_{s_m}\) est la masse de Dirac en un nœud, distincte
de la constante \(\delta_{\mathrm F}\). La mesure conserve les
douze fréquences et leurs poids. Pour un polynôme non nul de degré inférieur
à \(r\le12\),
\[
 \int p(s)^2\,d\nu_0(s)
 =\frac1{12}\sum_{m=1}^{12}p(s_m)^2>0,
\]
car un polynôme de degré inférieur à douze ne peut s’annuler aux
douze nœuds distincts. Ainsi, les matrices de Hankel
\((M_{i+j})_{0\le i,j<r}\) sont définies positives.
Le polynôme \(\prod_m(s-s_m)\) est de norme nulle ; l’espace
\(L^2(\nu_0)\) est de dimension douze.

\begin{proposition}[Représentation exacte du lecteur critique]
\label{hmtfg:articulacion:representacion}
Soient \(\psi_0,\ldots,\psi_{11}\) les polynômes orthonormaux
à coefficient principal positif pour \(\nu_0\), avec
\(\psi_0=1\). La multiplication par \(s\) a pour matrice de
Jacobi \(J_{12}\), réelle symétrique à coefficients sous-diagonaux positifs.
Avec \(e_0=(1,0,\ldots,0)^{\mathsf T}\),
\begin{equation}
 u_0(z)=\int\bigl(1-\cos(z\sqrt{s})\bigr)\,d\nu_0(s)
 =e_0^{\mathsf T}\bigl[I-\cos(z\sqrt{J_{12}})\bigr]e_0.
 \label{hmtfg:articulacion:perfil-jacobi}
\end{equation}
\end{proposition}
\begin{proof}
La positivité précédente permet d’orthogonaliser jusqu’au degré onze.
Si \(i<j-1\), l’autoadjonction de la multiplication donne
\(\langle s\psi_j,\psi_i\rangle
=\langle\psi_j,s\psi_i\rangle=0\).
Par symétrie, il ne reste que la diagonale et ses deux voisines ; le rapport
positif des coefficients principaux fixe la positivité de la
sous-diagonale. Définissons
\[
 Q_{mj}=\frac{\psi_j(s_m)}{\sqrt{12}},\qquad
 D=\operatorname{diag}(s_1,\ldots,s_{12}).
\]
L’orthonormalité donne \(Q^{\mathsf T}Q=I\), et l’évaluation de la
récurrence aux nœuds donne \(DQ=QJ_{12}\). Puisque \(Q\) est
carrée, \(J_{12}=Q^{\mathsf T}DQ\). Ses valeurs propres sont
positives, de sorte que la racine positive et le cosinus se calculent dans
cette diagonalisation. Enfin,
\(Qe_0=(1,\ldots,1)^{\mathsf T}/\sqrt{12}\), ce qui prouve
\eqref{hmtfg:articulacion:perfil-jacobi}.
\end{proof}

La réalisation par \((J_{12},e_0)\) restitue le même profil ;
ainsi, en conservant la famille et le sélecteur, elle restitue également
ses retours et ses premières racines. La positivité spectrale
justifie cette représentation. La transversalité paramétrique
utilise en outre la dérivée du retour et ses produits orbitaux,
développés dans \ref{hmtfg:escalado:14.3} et
\ref{hmtfg:escalado:6}. Ce sont des propriétés d’opérations différentes :
la démonstration du changement d’échelle conserve les deux et ne remplace pas le
contrôle de la tangente par la positivité des moments.

\subsection{Les canaux constitutifs et leur inversion étiquetée}
\label{hmtfg:articulacion:constitutiva}

Le lecteur constitutif agit sur un autre objet typé du même
support : le transport angulaire à deux feuilles. Ses coordonnées déjà
générées sont \(x=A_{\rm rad}\) et \(y=C^*_{\rm rad}\), dans la
chambre contractante \(x>|y|\). L’involution des feuilles
\(R=R^*=R^{-1}\) détermine les projecteurs marqués
\(P_\pm=(I\pm R)/2\). La représentation du transport est
\begin{equation}
 q_+=\exp[-(x+y)],\qquad q_-=\exp[-(x-y)],\qquad
 T=q_+P_++q_-P_-,\qquad 0<q_\pm<1.
 \label{hmtfg:articulacion:canales}
\end{equation}
L’évaluation exponentielle réalise ces canaux après leur
construction angulaire. Les incidences \(90\) et \(120\) du
transport fixent la réponse
\[
 \mathcal V=(I-T^{90})(I-T^{120})^{-1}.
\]
L’inverse existe parce que \(1-q_\pm^{120}>0\).
En écrivant \(s=q^{30}\), avec \(30=\gcd(90,120)\), on obtient
\begin{equation}
 f(s)=\frac{1-s^3}{1-s^4}
     =\frac{1+s+s^2}{1+s+s^2+s^3},\qquad
 r_\pm=f(q_\pm^{30}),\qquad
 \mathcal V=r_+P_++r_-P_-.
 \label{hmtfg:articulacion:respuesta}
\end{equation}
Dans cette sous-section, \(f\) désigne la réponse constitutive ;
\(f_\lambda\) continue de désigner l’application dynamique de
\eqref{hmtfg:articulacion:perfil}. Leurs arguments et leurs domaines
sont distincts.

\begin{proposition}[Restitution des canaux et de l’orientation]
\label{hmtfg:articulacion:inversion}
La fonction constitutive \(f:(0,1)\to(3/4,1)\) est une bijection
décroissante. Son inverse \(\psi\) associe à \(r\) l’unique racine
\(s\in(0,1)\) de
\begin{equation}
 r s^3+(r-1)(1+s+s^2)=0.
 \label{hmtfg:articulacion:cubica}
\end{equation}
Les lectures constitutives normalisées, avec leurs feuilles étiquetées,
\[
 \widehat\varepsilon=r_+^2,\qquad
 \widehat\mu=r_-^2,\qquad
 \widehat Z=\frac{r_-}{r_+},\qquad
 \widehat c=\frac1{r_+r_-},
\]
avec \(\widehat\varepsilon,\widehat\mu\in(9/16,1)\), permettent de restituer
\begin{equation}
 \begin{aligned}
 r_+&=\sqrt{\widehat\varepsilon},&
 r_-&=\sqrt{\widehat\mu},\\
 s_\pm&=\psi(r_\pm),& q_\pm&=s_\pm^{1/30}.
 \end{aligned}
 \label{hmtfg:articulacion:recuperacion-canales}
\end{equation}
\begin{equation}
 x=-\frac12\log(q_+q_-),\qquad
 y=\frac12\log\frac{q_-}{q_+}.
 \label{hmtfg:articulacion:recuperacion-angular}
\end{equation}
\end{proposition}
\begin{proof}
La dérivation du quotient rationnel donne
\begin{equation}
 f'(s)=-\frac{s^2(3+2s+s^2)}{(1+s+s^2+s^3)^2}<0
 \qquad(0<s<1).
 \label{hmtfg:articulacion:monotonia}
\end{equation}
Les limites \(f(0^+)=1\) et \(f(1^-)=3/4\), jointes à la
continuité, prouvent la bijectivité. Multiplier \(f(s)=r\) par
le dénominateur positif produit \eqref{hmtfg:articulacion:cubica}.
La positivité de \(r_\pm\) détermine les racines carrées
et celle de \(s_\pm\) détermine les racines trentièmes.
Par \eqref{hmtfg:articulacion:canales},
\(\log q_+=-x-y\) et \(\log q_-=-x+y\) ;
l’addition et la soustraction restituent \(x,y\).
Réciproquement, les racines restituées appartiennent à \((0,1)\),
de sorte que \(x\pm y=-\log q_\pm>0\) ; la paire obtenue
reste dans la chambre \(x>|y|\).
\end{proof}

En particulier,
\(\widehat\mu\widehat\varepsilon=\widehat c^{-2}\) et
\(\widehat\mu/\widehat\varepsilon=\widehat Z^2\).
L’échange des étiquettes permute \(q_+\) et \(q_-\),
conserve \(x\) et change le signe de \(y\). La restitution
exige cette orientation : le produit isolé \(r_+r_-\)
conserve la propagation commune, mais ne détermine pas les deux réponses
étiquetées. Les racines positives et l’inversion cubique restituent
les canaux à l’intérieur du lecteur fixé ; le registre complet de leur
génération reste dans l’état enrichi.

\subsection{Le moment de Catalan et la restitution fonctionnelle}
\label{hmtfg:articulacion:catalan}

Les mêmes incidences déterminent une relation entre fonctions,
avant l’évaluation de l’un ou l’autre canal. Pour \(0<s<1\), soit
\(\mathcal D=s\,d/ds\), et définissons le moment orienté
\begin{equation}
 \mathcal C_2(s)
 =\int_0^s\log\frac{s}{t}\,
     \frac{1+t}{1+t+t^2}f(t)\,dt
 =\int_0^s\frac{\log(s/t)}{1+t^2}\,dt.
 \label{hmtfg:articulacion:integral}
\end{equation}
La deuxième égalité découle de la factorisation
\(1+t+t^2+t^3=(1+t)(1+t^2)\). Le poids continu et positif
\(dt/(1+t^2)\) est ainsi produit par la composition rationnelle
du lecteur constitutif, et non par une identification à
\(\nu_0\).

\begin{proposition}[Inversion différentielle, unicité et lecture à l’extrémité]
\label{hmtfg:articulacion:restitucion-catalan}
Le moment \(\mathcal C_2\) est l’unique fonction
\(F\in C^2((0,1))\) qui satisfait
\begin{equation}
 \mathcal D^2F(s)=\frac{s(1+s)}{1+s+s^2}f(s),\qquad
 \lim_{s\downarrow0}F(s)=0.
 \label{hmtfg:articulacion:problema-diferencial}
\end{equation}
La réponse constitutive est restituée par
\begin{equation}
 \mathcal D^2\mathcal C_2(s)=\frac{s}{1+s^2},\qquad
 f(s)=\frac{1+s+s^2}{s(1+s)}\,
          \mathcal D^2\mathcal C_2(s).
 \label{hmtfg:articulacion:inversa-diferencial}
\end{equation}
Son extension continue à l’intervalle fermé possède la série
\begin{equation}
 \mathcal C_2(s)=\sum_{j=0}^{\infty}
       \frac{(-1)^j s^{2j+1}}{(2j+1)^2},\qquad 0\le s\le1,
 \qquad
 \mathcal C_2(1)=\sum_{j=0}^{\infty}\frac{(-1)^j}{(2j+1)^2}
 =G_{\rm Cat}.
 \label{hmtfg:articulacion:serie-catalan}
\end{equation}
\end{proposition}
\begin{proof}
L’intégrande de \eqref{hmtfg:articulacion:integral} est positif ou nul
et
\[
 0\le\mathcal C_2(s)\le\int_0^s\log(s/t)\,dt=s.
\]
Cela prouve l’intégrabilité en zéro et la limite requise.
L’identité \(\log(s/t)=\int_t^s du/u\), suivie de
l’interversion d’intégrales positives, fournit
\[
 \mathcal C_2(s)
 =\int_0^s\frac1u\left(\int_0^u\frac{dt}{1+t^2}\right)du.
\]
Sur tout intervalle compact du domaine ouvert, on applique le
théorème fondamental du calcul intégral. Par conséquent,
\[
 \mathcal D\mathcal C_2(s)=\int_0^s\frac{dt}{1+t^2},
 \qquad
 \mathcal D^2\mathcal C_2(s)=\frac{s}{1+s^2}.
\]
La factorisation rationnelle utilisée dans
\eqref{hmtfg:articulacion:integral} donne
\eqref{hmtfg:articulacion:problema-diferencial} ; isoler \(f\)
donne \eqref{hmtfg:articulacion:inversa-diferencial}, dont les
dénominateurs sont positifs dans \((0,1)\).

Si deux solutions ont la limite nulle requise, leur différence
\(H\) satisfait \(\mathcal D^2H=0\). Le changement \(z=\log s\)
transforme cette équation en \(d^2H(e^z)/dz^2=0\), d’où
\(H(s)=a\log s+b\). L’existence de la limite nulle impose
\(a=b=0\) et prouve l’unicité.

Pour \(s<1\), le développement géométrique de \((1+t^2)^{-1}\)
peut être intégré terme à terme avec le poids logarithmique, car
\[
 \int_0^s t^{2j}\log(s/t)\,dt
 =\frac{s^{2j+1}}{(2j+1)^2},\qquad
 \sum_{j\ge0}\frac{s^{2j+1}}{(2j+1)^2}<\infty.
\]
Cela prouve la série à l’intérieur. La borne
\(1/(2j+1)^2\) donne la convergence absolue et uniforme de cette
série sur \([0,1]\). Dans l’intégrale, l’intégrande prolongé
par zéro pour \(t>s\) est dominé par \(\log(1/t)\)
sur \((0,1)\). La convergence dominée permet d’atteindre
\(s=1\) et coïncide avec le prolongement de la série. Sa valeur
à l’extrémité est précisément la série de Catalan indiquée.
\end{proof}

L’inversion différentielle utilise la famille \(\mathcal C_2(s)\),
et non seulement sa valeur à l’extrémité \(G_{\rm Cat}\). En particulier,
\(\mathcal D^2\mathcal C_2(s)\to1/2\) lorsque \(s\uparrow1\).
La série dérivée deux fois avec \(\mathcal D\) y possède
une limite radiale d’Abel ; la série du moment lui-même est
absolument convergente à la frontière. Cette distinction maintient
les conditions précises de chaque passage à la limite.

\subsection{Composition des restitutions et portée conjointe}
\label{hmtfg:articulacion:conclusion}

Les interrelations précédentes ont des opérations et des domaines
explicites. Les moments de la mesure atomique \(\nu_0\)
produisent \(J_{12}\), et son calcul fonctionnel restitue le profil
\(u_0\). Le profil et le sélecteur de premières racines déterminent
la suite dont le changement d’échelle fournit la lecture \(\delta_{\mathrm F}\).
D’autre part, la fonction constitutive \(f\) produit
\(\mathcal C_2\) au moyen du moment intégral, et
\eqref{hmtfg:articulacion:inversa-diferencial} restitue \(f\)
à partir de cette famille. Ses évaluations étiquetées en
\(s_\pm=q_\pm^{30}\) déterminent les réponses constitutives ;
l’inversion cubique restitue les deux canaux et leurs coordonnées
angulaires. La lecture du moment à l’extrémité donne \(G_{\rm Cat}\).

La mesure \(\nu_0\) est positive, atomique et de rang douze.
Le moment de Catalan intègre un poids positif continu et possède
un développement alterné. Les signes de ce développement sont des
coefficients de la série géométrique ; ce ne sont pas des poids négatifs de
\(\nu_0\), et ils ne convertissent pas l’intégrale positive en la même
représentation spectrale. Chaque lecteur conserve son support,
son orientation, sa normalisation et son opération d’évaluation.

L’unité réside dans la structure APP--TRIT--TPK qui engendre et
transporte leurs antécédents ; les relations fonctionnelles sont
établies par les compositions démontrées. L’égalité
d’origine n’identifie pas des opérateurs de domaines différents.
En particulier, cette articulation ne postule pas de formule scalaire
qui détermine \(\delta_{\mathrm F}\) à partir de la constante de Catalan et
d’autres constantes physiques. Elle situe le changement d’échelle critique aux côtés de
restitutions effectives entre lecteurs du même support,
en conservant l’information requise par chaque direction.

\endgroup

\begingroup\allowdisplaybreaks[0]% Formalización reunida el 26 de septiembre de 2026.
% Procedencia y condiciones: COMUN/procedencia_20260926/.
\clearpage
\section{Normalisation conjointe de l’action, de la longueur et du couplage gravitationnel}
\label{md:rigidez}

La détermination d’une grandeur est conservée lors de la composition des lectures du
même état. On réunit ici le retour inscrit, la section d’action, la
vitesse constitutive et le caractère positif de route. L’antécédent commun
est la génération APP--TRIT--TPK, avec ses feuillets, son orientation, sa retenue, sa mémoire
et ses prolongements ; les coordonnées régionales et $\alpha$ interviennent comme
sorties de cette génération. La construction suivante utilise ces sorties
avant d’effectuer la reconnaissance dimensionnelle de $G$.

Le problème de normalisation comprend trois opérations distinctes : produire
une longueur à partir du retour enregistré, transporter le caractère sans perdre
son complément et comparer les lectures d’énergie et de longueur qui en résultent.
Les identités de cette section fixent conjointement ces opérations. Leurs
preuves portent sur des fibres finies ou sur des opérateurs bornés uniformément
positifs ; le secteur nul conserve son traitement séparé.

\subsection{Retour inscrit et section longitudinale}
\label{md:rigidez:longitud}

Dans le support marqué du transport, on conserve le support
\[
 \Omega=\{(j,k,q,u,v):j\in\mathbb Z/12\mathbb Z,\ 
 k-j\in\{0,3,6,9\},\ 1\le q\le8,\
 u,v\in\{1,2,4,5,7,8\},\ u<v\}.
\]
L’indice $j$ conserve la phase dodécaphasique ; les quatre décalages
conservent le quart de tour et les indices restants l’incidence marquée.
Par conséquent, $N=|\Omega|=12\cdot4\cdot8\binom62=5760$ compte les positions
d’incidence. L’espace de ce registre est transporté avec la mémoire
généalogique ; sa dimension finie correspond à ce support particulier.

Soient $B:V_U\to V_K$ l’inverseur de Hadamard transporté,
$B^*B=BB^*=I/4$, $u_K=N^{-1/2}(1,\ldots,1)$,
$P_K=u_Ku_K^*$ et $Q_K=I-P_K$. Nous définissons
\[
 C_0=Q_KB,\qquad M_0=P_KB,\qquad U=2B.
\]
L’orthogonalité des deux projecteurs démontre
\[
 C_0^*C_0+M_0^*M_0=I/4,\qquad C_0^*M_0=0,\qquad U^*U=UU^*=I.
\]
Ainsi, $x\mapsto(2C_0x,2M_0x)$ conserve l’archive complète de ce
transport. $U$ est sa normalisation unitaire ; l’évolution enrichie
$U_t$ maintient son propre domaine et sa mise à jour.

Soit $E_i=e_ie_i^*$ la résolution du registre et
$\Delta_\Omega(A)=\sum_iE_iAE_i$ son espérance diagonale.

\begin{lemma}[Retour longitudinal inscrit]
\label{md:rigidez:retorno}
Le retour et sa lecture inscrite satisfont
\[
 C_0C_0^*=Q_K/4,\qquad
 \Delta_\Omega(C_0C_0^*)=r_NI,\qquad
 r_N=\frac{N-1}{4N}=\frac{5759}{23040}.
\]
Avec la norme scalaire locale $n_H=\alpha^{16}$ du conoyau de $H_4$,
la composition longitudinale sur une cochaîne de longueur $\beta_*$ est
\[
 \mathscr T_L(\beta_*)=
 n_H\bigl(\Delta_\Omega(C_0C_0^*)\otimes I\bigr)\beta_*
 =q_\alpha\beta_*,
 \qquad q_\alpha=\alpha^{16}r_N.
 \label{md:rigidez:qalpha}
\]
\end{lemma}
\begin{proof}
L’égalité $BB^*=I/4$ donne $C_0C_0^*=Q_K/4$. Chaque entrée diagonale
de $P_K$ vaut $1/N$, ce qui prouve l’espérance. Multiplier une
cochaîne par ce scalaire conserve son type longitudinal et donne la dernière
égalité. Si $\beta_*(\bar a)=-U_a\beta_*(a)$, l’image satisfait la
même inversion. De même, à partir de
$\beta_m(a)=\sum_jU_j^{-1}\beta_{m+1}(a_j)$, on obtient l’identité
raffinée en multipliant les deux membres par $q_\alpha$.
\end{proof}

L’espérance possède une réalisation inscrite explicite :
$We_i=e_i\otimes e_i$ et
$W^*(A\otimes I)W=\Delta_\Omega(A)$. Pour
$R_K=Q_K/4$ et $P_W=WW^*$, le complément
$Z_\perp=(I-P_W)(R_K\otimes I)W$ vérifie
\[
 (R_K\otimes I)W=Wr_N+Z_\perp,\qquad W^*Z_\perp=0,\qquad
 Z_\perp^*Z_\perp=\frac{N-1}{16N^2}I.
\]
En effet, $R_K^2=R_K/4$, de sorte que le carré de la norme complète
est $r_N/4$, tandis que la composante inscrite a pour carré $r_N^2$.
La différence est précisément le complément écrit. La conservation de
ce complément accompagne la lecture longitudinale ; l’ajouter à une autre lecture
changerait l’observable.

L’espérance est unique parmi les applications bimodulaires vers
l’algèbre diagonale qui fixent ses projecteurs : sur une unité matricielle
$E_{ij}$, la bimodularité exige
$\mathcal E(E_{ij})=E_i\mathcal E(E_{ij})E_j$, qui vaut zéro si $i\ne j$,
et $\mathcal E(E_{ii})=E_{ii}$. Cette unicité conserve la résolution marquée.

Avec la section longitudinale de référence $L_*$ déjà spécifiée, il reste
\begin{equation}
 L_\alpha=q_\alpha L_*,\qquad D=L_\alpha U.
 \label{md:rigidez:longitud-def}
\end{equation}
Le lecteur utilise le retour comme opérateur linéaire sur une longueur. Une
lecture quadratique de la source produirait un autre type de grandeur. La
section $L_*$, la résolution inscrite et la norme locale sont des antécédents
explicites de cette composition.

\subsection{Transport polaire et réciprocité du même caractère}
\label{md:rigidez:polar}

Soit $X$ le caractère opératoriel du secteur positif non nul, en conservant la route et
ses normalisations. En dimension finie, on exige $X>0$ ; en dimension
infinie, on exige $X$ borné, autoadjoint et $X\ge\epsilon I$ pour un certain
$\epsilon>0$. Nous écrivons $Y=UXU^*$. La représentation radiale polaire est
\[
 R=(DX^2D^*)^{1/2}.
\]
La racine positive est unique et, en substituant $D=L_\alpha U$, on obtient
$R=L_\alpha Y$. De manière équivalente,
$\|Rv\|^2=\|XD^*v\|^2$ pour tout $v$ : la polarisation retrouve
$R^2=DX^2D^*$ et la positivité détermine $R$. Cette condition spécifie
la métrique induite qui est conservée.

L’action $\hbar$ et la vitesse constitutive $c$ appartiennent à la même
section. La réalisation circulaire compatible fixe
\begin{equation}
 \omega_0=\frac{c}{L_\alpha},\qquad
 t_0=\frac{\pi_{\rm HMT}L_\alpha}{54c},\qquad
 R_{108}=\frac{54}{\pi_{\rm HMT}}ct_0=L_\alpha.
 \label{md:rigidez:reloj}
\end{equation}
La longueur a été construite avant d’exprimer l’horloge circulaire. Cette
égalité permet de conserver la notation $R_{108}$ dans les développements
ultérieurs et fixe sa normalisation. Les bases dimensionnelles sont transportées
conjointement ; $L_*$ et la base d’arête $u_L$ ont des fonctions différentes.

\begin{theorem}[Détermination conjointe du couplage radial]
\label{md:rigidez:teorema}
Les lecteurs sur le même caractère,
\[
 H=\frac{\hbar c}{L_\alpha}Y,\qquad
 \Lambda=L_\alpha Y^{-1},\qquad R=L_\alpha Y,\qquad M=H/c^2,
\]
satisfont
\begin{equation}
 H\Lambda=\hbar cI,\qquad R\Lambda=L_\alpha^2I,\qquad
 R=\frac{L_\alpha^2}{\hbar c}H.
 \label{md:rigidez:identidades}
\end{equation}
Dans l’identification dimensionnelle de $R$ comme rayon gravitationnel réduit,
$R=GM/c^2$, le coefficient commun est unique :
\begin{equation}
 \boxed{G=\frac{c^3L_\alpha^2}{\hbar}.}
 \label{md:rigidez:G}
\end{equation}
\end{theorem}
\begin{proof}
La multiplication des lecteurs donne les deux premières identités.
Les multiplier à droite par $\Lambda^{-1}$ prouve la troisième.
Substituer $M=H/c^2$ identifie le coefficient. Si $G'$ conserve les mêmes
lectures, $(G'-G)M/c^2=0$ ; dans la fibre non nulle, l’inversibilité de $M$
implique $G'=G$.
Tous les produits sont définis dans la fibre complète sous les
hypothèses de l’énoncé.
\end{proof}

La preuve s’applique à tous les modes positifs de cette réalisation, sans choisir
la masse électronique pour normaliser $G$. L’identification du rayon
réduit est explicite ; le rayon $2GM/c^2$ correspond à une autre convention
de lecture de la même constante. Le secteur nul est maintenu en dehors de
l’inverse. L’extension aux opérateurs non bornés requiert un développement de
domaines et de limites distinct du présent théorème.

Le carré de $L_\alpha$ procède du produit des longueurs conjuguées.
L’action apparaît inversement lorsqu’on élimine $\Lambda$ entre l’énergie et la
longueur. Le passage de l’énergie à la masse et de la masse au rayon apporte $c^4$,
qui, combiné avec le $c$ de $H\Lambda=\hbar cI$, produit $c^3$.
Les deux normalisations fixent le facteur adimensionnel à un. L’analyse
dimensionnelle vérifie les exposants ; la composition détermine le coefficient.

\subsection{Échelles de Planck et spécialisation du caractère unitaire}
\label{md:rigidez:escalas}

La longueur construite, la section d’action et la vitesse déterminent
\[
 t_P=\frac{L_\alpha}{c},\qquad
 m_{P,a}=\frac{\hbar_a}{cL_\alpha},\qquad
 E_{P,a}=\frac{\hbar_a c}{L_\alpha}.
\]
Substituer \eqref{md:rigidez:G} dans
$\sqrt{\hbar_aG_a/c^5}$, $\sqrt{\hbar_ac/G_a}$ et
$\sqrt{\hbar_ac^5/G_a}$ démontre respectivement ces identités ;
la positivité fixe chaque racine. La longueur précède la reconnaissance
conventionnelle de ces échelles.
Pour un mode de caractère $y>0$,
\[
 m(y)=m_{P,a}y,\qquad r_g(y)=L_\alpha y,\qquad
 \bar\lambda(y)=L_\alpha/y.
\]
L’égalité des deux longueurs équivaut à $y=y^{-1}$ et, par conséquent,
à $y=1$. Le produit surfacique est conservé pour tous les modes positifs ;
la coïncidence des rayons caractérise le mode unitaire.

Le pas chronologique satisfait $t_0=(\pi_{\rm HMT}/54)t_P$ ;
la période de phase est $108t_0=2\pi_{\rm HMT}t_P$ et
$\omega_{108}=1/t_P$. Un horizon $R_h=\zeta_hL_\alpha$ conserve
sa condition géométrique propre. Pour la spécialisation de Schwarzschild
de la même masse, $R_h=2r_g=2L_\alpha y$.
En comparant des sections avec $L_\alpha,c$ fixes et
$\rho_a=\hbar_b/\hbar_a$, on obtient
$G_b=G_a/\rho_a$, $m_{P,b}=\rho_a m_{P,a}$ et
$E_{P,b}=\rho_a E_{P,a}$. Cette comparaison de sections conserve
$L_\alpha$ et $t_P$.

\subsection{Aire et restitution des secteurs de mémoire}
\label{md:rigidez:area}

L’incidence de bord conserve le regroupement de six trits.
Si $x_j=r_{2j-1}+3r_{2j}$, $r_i\in\{0,1,2\}$, l’application
\[
 (x_1,x_2,x_3)\longmapsto
 (r_1+3r_2+9r_3,\ r_4+3r_5+9r_6)
\]
est une bijection entre $(\mathbb Z/9\mathbb Z)^3$ et
$(\mathbb Z/27\mathbb Z)^2$. Son inverse extrait les six chiffres dans l’ordre,
en conservant l’endroit où la retenue est enregistrée. Dans le tore de bord,
le bloc $A=\{0,\ldots,8\}^2$ a deux collections de liens :
\begin{align*}
 \partial A_{90}&=\{((-1,j),(0,j)),((8,j),(9,j)):0\le j\le8\},\\
 \partial A_{120}&=\{((i,-1),(i,0)),((i,8),(i,9)):0\le i\le8\}.
\end{align*}
Chaque collection contient dix-huit liens. Leurs étiquettes conservent la
provenance de canal ; chaque lien porte un opérateur nombre de spectre
$\{0,1,2\}$. Soient $N_{90},N_{120}$ les sommes correspondantes.

L’échelle angulaire surfacique conserve le contre-angle déjà construit :
\[
 \theta_{\rm clk}=\frac{\pi C^*_{\deg}}{1080},\qquad
 a_0=\frac{1+2\cos(\pi/18)}3\,\theta_{\rm clk}.
\]
Le préfacteur est la trace réelle normalisée de
$\operatorname{diag}(1,e^{i\pi/18},e^{-i\pi/18})$.
Ainsi, $a_0$ conserve la signification de normalisation adimensionnelle d’aire.
La lecture d’aire est prise sur la branche positive
$\theta_{\rm clk}>0$ et $\gamma>0$.
Avec la fonctionnelle $\gamma$ de Barbero--Immirzi et la longueur déjà déterminée,
\[
 \mathcal A=\gamma a_0L_\alpha^2(N_{90}+N_{120}).
\]
Les trente-six opérateurs locaux agissent dans des facteurs tensoriels
distincts. Par conséquent, le spectre de $N=N_{90}+N_{120}$ est
$\{0,\ldots,72\}$ et la multiplicité de $n$ est le coefficient de $z^n$
dans $(1+z+z^2)^{36}$. La preuve consiste à multiplier les fonctions
génératrices des facteurs. La nouvelle longueur transporte l’unité
surfacique et conserve cette combinatoire.

\begin{proposition}[Lecture conjointe de l’aire et de la provenance sectorielle]
Le couple
$N=N_{90}+N_{120}$ et $\mathcal D_{\partial}=90N_{90}+120N_{120}$
restitue les deux secteurs :
\[
 N_{90}=4N-\frac{\mathcal D_{\partial}}{30},\qquad
 N_{120}=\frac{\mathcal D_{\partial}}{30}-3N.
\]
\end{proposition}
\begin{proof}
La matrice de la transformation est
$\left(\begin{smallmatrix}1&1\\90&120\end{smallmatrix}\right)$,
de déterminant $30$. Son inverse donne les formules. La commutativité
des deux opérateurs nombre implique celle des deux lecteurs.
\end{proof}
Dans la réalisation modulaire avec
$\rho_A=Z_A^{-1}\exp(-\Delta_{\rm mod}\mathcal D_{\partial})$,
$\Delta_{\rm mod}>0$, le calcul fonctionnel donne
$-\log\rho_A=(\log Z_A)I+\Delta_{\rm mod}\mathcal D_{\partial}$.
Les normalisations étant conservées, le couple formé par l’aire et le
Hamiltonien modulaire retrouve les deux occupations. La preuve de restitution
vaut pour chaque valeur construite de $\Delta_{\rm mod}$ ; le nouveau lecteur
longitudinal ne redéfinit ni ce coefficient ni l’état modulaire.
La somme des occupations et la provenance sectorielle expriment des observables
distincts. De même, les trente-six liens de réponse et les
quatre-vingt-une capacités unitaires de la coupure cubique conservent leurs fonctions.

\subsection{Variation conjointe et réduction compatible de mémoire}
\label{md:rigidez:variaciones}

Pour des collections différentiables en norme et uniformément positives dans un
voisinage, soit $\kappa=L_\alpha^2/(\hbar c)$. Alors
\[
 \delta R=\delta\kappa\,H+\kappa\,\delta H,\qquad
 \frac{\delta G}{G}=3\frac{\delta c}{c}
 +2\frac{\delta L_\alpha}{L_\alpha}-\frac{\delta\hbar}{\hbar}.
\]
La première identité résulte de la différentiation de
$R=L_\alpha Y$ et $H=(\hbar c/L_\alpha)Y$ ; elle ne requiert pas
$[Y,\delta Y]=0$. Une variation admissible de connexion qui agit dans $Y$
et conserve les trois sections maintient $\delta G=0$.
L’affirmation comprend les variations de cette réalisation radiale ;
l’équivalence entre des fonctionnelles de champ complètes a ses propres
domaines et conditions aux limites.

Soit maintenant
$Y=\left(\begin{smallmatrix}A&C\\C^*&D_i\end{smallmatrix}\right)>0$
dans la séparation entre lecture et mémoire, avec
$S=A-CD_i^{-1}C^*>0$. L’extension minimisante d’un vecteur $b$ est
$(b,-D_i^{-1}C^*b)$, comme le montre la complétion du carré :
\[
 \left\langle\binom b m,Y\binom b m\right\rangle
 =\langle b,Sb\rangle+
 \langle m+D_i^{-1}C^*b,D_i(m+D_i^{-1}C^*b)\rangle.
\]
Par conséquent, la réduction cohérente donne
\[
 H_{\rm eff}=\frac{\hbar c}{L_\alpha}S,\qquad
 R_{\rm eff}=L_\alpha S,\qquad\Lambda_b=L_\alpha S^{-1}.
\]
La compression de $Y^{-1}$ sur le bloc de lecture est $S^{-1}$ ;
la compression directe de $Y$ est $A$. La mémoire distingue ces opérations.

Si l’élimination dynamique produit une métrique temporelle bornée
$Z\ge z_0I$, $z_0>0$, on prend $T=Z^{1/2}$, borné et inversible.
La normalisation canonique exige les transports duaux
\[
 H_c=T^{-*}H_{\rm eff}T^{-1},\qquad
 R_c=T^{-*}R_{\rm eff}T^{-1},\qquad
 \Lambda_c=T\Lambda_bT^*.
\]
Leur multiplication donne $H_c\Lambda_c=\hbar cI$ et
$R_c\Lambda_c=L_\alpha^2I$, sans hypothèse de commutativité.
Dans le développement résolvant d’une énergie de blocs $H_{bb},V,H_{ii}$,
avec $H_{ii}$ positif inversible et
$|\hbar\omega|<\|H_{ii}^{-1}\|^{-1}$,
\[
 K_b(\omega)=H_{bb}-\hbar\omega I
 -V(H_{ii}-\hbar\omega I)^{-1}V^*
 =H_{\rm stat}-\hbar\omega Z+O(\omega^2),
 \qquad Z=I+VH_{ii}^{-2}V^*.
\]
La dernière expression est locale autour de zéro et conserve l’ordre
d’approximation ; les produits duaux précédents sont des identités exactes
des lecteurs normalisés.

La proportionnalité est également conservée dans l’élimination dynamique
complète. Pour $z$ tel que $H_{ii}-zI$ soit inversible, nous définissons
\[
 \mathcal K_H(z)=H_{bb}-zI-V(H_{ii}-zI)^{-1}V^*,\qquad
 \kappa=\frac{L_\alpha^2}{\hbar c}.
\]
La même décomposition de lecture et de mémoire appliquée à $R=\kappa H$
produit
\begin{equation}
 \mathcal K_R(\kappa z)=\kappa\mathcal K_H(z).
 \label{md:rigidez:resolvente-exacto}
\end{equation}
En effet, ses blocs sont $R_{bb}=\kappa H_{bb}$,
$R_{bi}=\kappa V$ et $R_{ii}=\kappa H_{ii}$, et
\[
 (R_{ii}-\kappa zI)^{-1}
 =\kappa^{-1}(H_{ii}-zI)^{-1}.
\]
La substitution dans le complément de Schur démontre \eqref{md:rigidez:resolvente-exacto}, en conservant l'ordre des facteurs. Lorsque les deux compléments sont inversibles, leurs inverses satisfont $\mathcal K_R(\kappa z)^{-1}=\kappa^{-1}\mathcal K_H(z)^{-1}$. Le paramètre énergétique est $z=\hbar\omega$ et son image dans le lecteur radial est $\kappa z$ : tous deux conservent leurs dimensions respectives. L'identité vaut sur tout ce domaine résolvant et transporte donc conjointement tous les ordres de la mémoire dynamique. Le développement local précédent décrit ses premiers coefficients ; cette égalité explique la conservation exacte du rapport radial lors de la réduction.

\subsection{Rigidité temporelle et mémoire de frontière}
\label{md:rigidez:memoria}

La conservation de la seule aire admet $R\mapsto\zeta R$, $\Lambda\mapsto\Lambda/\zeta$. Conserver également $H,\hbar,c$ et $H\Lambda=\hbar cI$ impose $\zeta=1$. Si $H\mapsto\zeta H$ est également transformé, le quotient $R/H$ et le coefficient $G$ restent égaux. Les comparaisons doivent conserver les mêmes observations conjointes.

La phase relevée apporte une autre restriction effective. Dans une réalisation bornée $H=H_{\rm clk}+D_0^*WD_0$ avec $W>0$, l'évolution $P(t)=\exp(-itH/\hbar)$ détermine $H=i\hbar\dot P(0)$. Une fois $H_{\rm clk},D_0,\hbar$ et l'horloge fixés, l'égalité des évolutions implique $D_0^*(W'-W)D_0=0$. Pour un parcours d'extrémité initiale fixée, $D_0$ est triangulaire inversible et l'on en déduit $W'=W$. Pour un graphe général, on détermine la forme sur $\operatorname{Ran}D_0$ ; étendre ensuite le domaine des résidus modifierait le problème d'élimination. Un seul retour $P(\tau)$ admet des alias spectraux, tandis que la phase relevée ou les temps compatibles tendant vers zéro les distinguent.

Dans la coordinatisation de capacité nonadique, soient $\rho=\log_{10}9$, $\delta=1-\rho$, $p_j=\lfloor j/\delta\rfloor$ et $v_j=p_j-j$. La capacité d'une fenêtre est
\[
 C(j,n)=p_{j+n}-p_j-n.
\]
Pour $n=35$, ses valeurs sont $729$ ou $730$ selon l'origine ; le retour de capacité modulo neuf distingue la première au sein de cette classe. La fenêtre $j=20\to55$ satisfait $p_j:437\to1201$ et $v_j:417\to1146$. Sa phase de capacité revient à $3$ et le quotient augmente de $46$ à $127$. Les six capacités de $20$ et les vingt-neuf de $21$ ont pour somme $729$ ; la mémoire accumulée fait toujours partie de l'état.

Pour la première arrivée $T(k)=\lceil k/\rho\rceil-1$, le dénombrement est $N(k,C)=T(k+C)-T(k)-C$. Écrire $\epsilon(k)=\lceil k/\rho\rceil-k/\rho$ donne exactement
\[
 N(k,C)=(\rho^{-1}-1)C+\epsilon(k+C)-\epsilon(k).
\]
Les fenêtres de capacité $729$ dont les dénombrements sont $34$ et $35$ conservent des termes de frontière différents. Même $k=21$ et $k=2289$ partagent $(k\bmod108,T(k)\bmod108)=(21,22)$ et accumulent respectivement $1$ et $109$ vacances. La coïncidence des calendriers finis conserve moins d'information que l'état complet. Ces calculs interprètent le retour ; le lecteur longitudinal de \eqref{md:rigidez:longitud-def} possède son propre constructeur et ne sélectionne pas son échelle à l'aide de ces fenêtres.

\subsection{Normalisation géométrique et évaluation}
\label{md:rigidez:evaluacion}

Le bloc central d'APP
$B_c=\left(\begin{smallmatrix}7&2\\2&7\end{smallmatrix}\right)$
satisfait $B_c^{\mathsf T}g_cB_c=45g_c$, $g_c=\operatorname{diag}(1,-1)$. Sa normalisation positive unimodulaire est $B_c/\sqrt{45}$ : un facteur positif supplémentaire $a$ donnerait le déterminant $a^2$. L'incidence aréale--volumétrique
$T_{\rm av}=\left(\begin{smallmatrix}1&1\\1&2\end{smallmatrix}\right)$
conserve
$G_{\rm av}=\left(\begin{smallmatrix}2&1\\1&-2\end{smallmatrix}\right)$
par multiplication directe. Dans une réalisation de Minkowski, la forme causale et la soudure dimensionnelle jouent des rôles distincts. La cochaîne $\beta_m=\ell_m\sigma_mW_\square e_\lambda$, $\ell_m=\ell_0/9^m$, et la relation $\ell_0=ct_0$ conservent l'échelle avec le cône. Modifier la soudure en ne conservant que la signature change cette réalisation.

L'unité expansive de Pell de norme un et de trace quatre satisfait $x+x^{-1}=4$, $x>1$, d'où $x=2+\sqrt3$. L'orbite $r_n=r_0x^n$ conserve $r_nr_{-n}=r_0^2$. Lorsque $r_0=L_\alpha$, le produit conserve la même section aréale. La norme, la forme marquée, le retour et la section longitudinale restreignent des objets différents ; leurs fonctions sont maintenues lors de leur composition.

\Needspace{10\baselineskip}
Dans la section SI déclarée, $L_*=1\,\mathrm m$, $S_*=10^{-34}\,\mathrm{J\,s}$ et $c=299792458\,\mathrm{m\,s^{-1}}$, la substitution structurelle produit
\[
 L_\alpha=1.616254982625126330\ldots\,10^{-35}\,\mathrm m,\qquad
 \hbar_{\rm ret}=1.054571816915039061\ldots\,10^{-34}\,\mathrm{J\,s},
\]
\begin{equation}
 G_{\rm ret}=6.674299659414022706\ldots\,10^{-11}
 \,\mathrm{m^3\,kg^{-1}\,s^{-2}}.
 \label{md:rigidez:valor}
\end{equation}
L'évaluateur substitue les coordonnées régionales et le registre $K$, reconstruit $\alpha$, le retour de $H_5$ et $L_\alpha$ ; la comparaison métrologique intervient ensuite. La référence CODATA 2022 est $6.67430(15)\,10^{-11}$ dans ces unités \cite{codata2022}. La différence est $-3.40586\,10^{-18}$, soit $-0.00227$ incertitudes-types. Les chiffres calculés décrivent l'évaluation mathématique dans cette section, et non une incertitude expérimentale. Les sections d'action antérieure et retournée conservent leurs étiquettes et $\hbar_aG_a=c^3L_\alpha^2$.

La signification structurelle de $G$ est ainsi déterminée : c'est le coefficient qui convertit, dans la réalisation radiale commune, la lecture énergétique en longueur conjuguée en préservant simultanément l'action et l'aire. La mémoire, les marques de feuillet et les bases font partie de la construction qui fixe ce coefficient. Leur conservation conjointe explique l'unicité démontrée.
\endgroup
% Preparación autónoma de VII, 2026-09-30. Fragmento sin inputs adicionales.
% RESULTADOS_RECUPERADOS, pruebas reunidas sin modificación de los originales:
% /Users/ruben/Documents/excelencia academica/output/CUMPLIMIENTO_EDITORIAL_20260928/fuentes/VIII_ES/manuscrito/29_retorno_y_memoria_traslacional.tex
% /Users/ruben/Documents/excelencia academica/output/CUMPLIMIENTO_EDITORIAL_20260928/fuentes/VIII_ES/manuscrito/30_pantalla_y_respuesta.tex
% /Users/ruben/Documents/excelencia academica/output/CUMPLIMIENTO_EDITORIAL_20260928/fuentes/VIII_ES/manuscrito/30b_variacion_energia_memoria.tex
% /Users/ruben/Documents/excelencia academica/output/CUMPLIMIENTO_EDITORIAL_20260928/fuentes/VIII_ES/manuscrito/30c_composicion_corriente_conexion.tex
% /Users/ruben/Documents/excelencia academica/output/CUMPLIMIENTO_EDITORIAL_20260928/fuentes/VIII_ES/manuscrito/30d_realizacion_geometrica.tex
% /Users/ruben/Documents/excelencia academica/output/CUMPLIMIENTO_EDITORIAL_20260928/fuentes/VIII_ES/manuscrito/31_corriente_y_respuesta_cartan.tex
% /Users/ruben/Documents/excelencia academica/output/CUMPLIMIENTO_EDITORIAL_20260928/fuentes/VIII_ES/manuscrito/33_corriente_espinorial_y_densidad.tex
% Dependencias internas ya presentes: núcleo común, rutas, lectores y unidades.
% Las inversas de 31 están demostradas en 10, hmtcuatro:gea:torsion.
% Legendre, fuente completa y restricciones permanecen en 20.
% Se incluye de 33 sólo lo utilizado por 10+20, no sus estados cosmológicos.
\section{Fondements géométriques et matériels de l’action conjointe}
\label{hmtcuatro:ant:seccion}

Les graines et les opérations d’APP, l’orientation et le régime du TRIT et le transport du TPK produisent les routes enrichies sur lesquelles agissent les lecteurs du noyau commun. Feuillet, orientation, retenue, frontière et mémoire sont conservés conjointement ; le retour de phase ne réinitialise pas cet état. Les constructions du continu appartiennent au même objet et ne sont pas transformées ici en secteurs indépendants. La réalisation qui suit transporte vers la géométrie et la matière cette structure déjà construite. Ses unités et ses couplages sont des sorties HMT antérieures : aucune comparaison numérique n’intervient comme générateur.

Pour que l’action et sa réduction puissent être lues au sein de ce manuscrit, nous réunissons les opérations d’écran, la variation effective du transport et le courant spinoriel qu’elles utilisent. L’élimination de Cartan--Holst est démontrée dans \ref{hmtcuatro:gea:torsion} ; l’action totale et sa réduction canonique, dans \ref{hmtcuatro:accion-retroaccion-restricciones}. Les formules de cette section précisent les données et les calculs qui interviennent dans les deux.

\subsection{Incidence, écran et soudure}
\label{hmtcuatro:ant:pantalla}

La représentation cubique de la frontière d’une route possède les six faces \(F=\{(i,\epsilon):1\leq i\leq3,\ \epsilon=\pm1\}\). Sur \(H_F=\mathbb R^F\), l’opposition est \(Re_{i,\epsilon}=e_{i,-\epsilon}\). Avec
\[
 u=\sum_{i,\epsilon}e_{i,\epsilon},\qquad
 P_u=\frac{uu^{\mathsf T}}6,\qquad
 P_-=\frac{I-R}2,\qquad
 P_0=\frac{I+R}2-P_u ,
\]
les trois projecteurs sont orthogonaux, leur somme vaut \(I\) et leurs rangs sont \(1,3,2\). L’incidence entre faces et sommets \(\nu\in\{-1,1\}^3\) est \(M_{(i,\epsilon),\nu}=\mathbf1_{\{\nu_i=\epsilon\}}\).

\begin{proposition}[Écran de l’incidence orientée]
\label{hmtcuatro:ant:prop-pantalla}
On a
\[
 MM^{\mathsf T}=12P_u+4P_-,\qquad
 \ker M^{\mathsf T}=P_0H_F .
\]
Le quotient \(Q=H_F/P_0H_F\), identifié à \((P_u+P_-)H_F\), est de dimension quatre. La forme \(B=P_--P_u\) sur \(Q\) a pour signature \((-+++)\).
\end{proposition}
\begin{proof}
Chaque face contient quatre sommets ; deux faces opposées n’en partagent aucun et deux faces d’axes distincts en partagent deux. Par conséquent,
\[
 MM^{\mathsf T}=4I+2(uu^{\mathsf T}-I-R)
              =12P_u+4P_- .
\]
L’identité \(\|M^{\mathsf T}x\|^2=\langle x,MM^{\mathsf T}x\rangle\) détermine le noyau. Dans la base orthonormée
\[
 t=u/\sqrt6,\qquad d_i=(e_{i,+}-e_{i,-})/\sqrt2
\]
de \(Q\), la matrice de \(B\) est \(\operatorname{diag}(-1,1,1,1)\).
\end{proof}

Soit \(W_\square=\sqrt6P_u+\sqrt2P_-\). Alors
\[
 W_\square e_{i,\epsilon}=t+\epsilon d_i=:n_{i,\epsilon},
 \qquad B(n_{i,\epsilon},n_{i,\epsilon})=0,\qquad
 MM^{\mathsf T}|_Q=2W_\square^*W_\square|_Q .
\]
Ces égalités s’obtiennent en appliquant les projecteurs à la base des faces. Les rotations propres du cube préservent les projecteurs, l’orientation temporelle \(t\) et la forme \(B\).

La route enrichie fournit à chaque arête \(a\) sa face, son signe \(\sigma_m(a)\) et le transport \(U_m(a):Q_{s(a)}\to Q_{t(a)}\). Sa soudure cellulaire est
\begin{equation}
 \beta_m(a)=\ell_m\sigma_m(a)n_{\lambda(\underline a)}^{s(a)},
 \quad \ell_m=\ell_0\,9^{-m},\quad
 \beta_m(\bar a)=-U_m(a)\beta_m(a).
 \label{hmtcuatro:ant:soldadura}
\end{equation}
L’inversion \(U_m(\bar a)=U_m(a)^{-1}\) implique qu’inverser deux fois restitue la soudure. Le registre ordonné de la retenue est également conservé : il n’est pas remplacé par le vecteur d’arrivée.

\begin{proposition}[Naturalité de la soudure]
\label{hmtcuatro:ant:refinamiento-soldadura} Soit \(a=a_1\cdots a_9\) un raffinement compatible : les incidences filles transportées à l’origine du prédécesseur conservent face et signe, et \(\ell_{m+1}=\ell_m/9\). Si \(\mathcal U_j:Q_{m,s(a)}\to Q_{m+1,s(a_j)}\) inclut l’identification par raffinement et le transport jusqu’à \(s(a_j)\), alors
\[
 \beta_m(a)=\sum_{j=1}^9\mathcal U_j^{-1}\beta_{m+1}(a_j).
\]
\end{proposition}
\begin{proof}
Chaque terme dans la fibre d’origine est \((\ell_m/9)\sigma_m(a)n_{\lambda(\underline a)}\). Leur somme donne \eqref{hmtcuatro:ant:soldadura}. Inverser la route inverse l’ordre des transports et chaque transport ; la formule d’inversion de la soudure donne le signe global. La concaténation conserve dans le même ordre la mémoire des neuf incidences. L’égalité s’itère par composition de raffinements compatibles.
\end{proof}

\subsection{Représentation du transport et réalisation géométrique}
\label{hmtcuatro:ant:geometria}

Le registre d’une route \(\gamma\) compte les inversions et les changements effectifs de feuillet au moyen du cocycle \(c(\gamma)=(c_g(\gamma),c_s(\gamma))\in\mathbb Z^2\). Le décompte est cumulatif sur les routes vers l’avant ; son extension au groupoïde est orientée et attribue des incréments opposés à la route inverse. Pour des routes composables, \(c(\gamma_2\gamma_1)=c(\gamma_2)+c(\gamma_1)\), et \(c(\bar\gamma)=-c(\gamma)\). Soit \(R_4\) un élément d’ordre quatre dans \(\operatorname{Spin}^+(3,1)\) dont la représentation vectorielle \(\Lambda(R_4)\) fixe le vecteur temporel unitaire \(u_0\). La représentation affine reçue est
\[
 \rho_{\rm C}^{\rm disc}(\gamma)
 =\bigl(R_4^{c_g(\gamma)},\ell_0c_s(\gamma)u_0\bigr).
\]
Avec le produit \((R,a)(S,b)=(RS,a+\Lambda(R)b)\), la fixation de \(u_0\) et l’additivité de \(c\) prouvent directement que \(\rho_{\rm C}^{\rm disc}\) conserve l’identité, la concaténation et l’inversion. Lorsque la composante rotationnelle revient à l’identité, la translation de mémoire reste déterminée par \(c_s(\gamma)\), qui n’est pas réinitialisé.

La normalisation
\[
 N_{\ell_0}(R,a)=(\Lambda(R),a/\ell_0),\qquad
 \rho_{\rm C}=N_{\ell_0}\circ\rho_{\rm C}^{\rm disc}
\]
est également un homomorphisme, car \((a+\Lambda(R)b)/\ell_0=a/\ell_0+\Lambda(R)(b/\ell_0)\). Le relèvement \(R_4^{c_g(\gamma)}\) est conservé lorsque des spineurs agissent ; il ne se récupère pas à partir du seul \(\Lambda(R_4^{c_g(\gamma)})\), dont la représentation a pour noyau \(\{1,-1\}\).

Dans une réalisation d’écran \(J_{{\rm scr},m}\) qui transporte la forme et la soudure, \(e_m=J_{{\rm scr},m}\beta_m\). La réalisation lisse considérée consiste en une région orientée \(\mathcal U\) de dimension quatre, une cotétrade non dégénérée \(e\), une connexion de Lorentz \(\omega\) et une réalisation des routes \(\mathfrak R_{\rm C}\) avec la condition de compatibilité
\begin{equation}
 \operatorname{Hol}_{\mathcal A_{\ell_0}}
          (\mathfrak R_{\rm C}(\gamma))
 =\rho_{\rm C}(\operatorname{Hol}_{\rm TPK}(\gamma)),\qquad
 \mathcal A_{\ell_0}=
 \begin{pmatrix}\omega&\ell_0^{-1}e\\0&0\end{pmatrix}.
\label{hmtcuatro:ant:compatibilidad-holonomia}
\end{equation}
L’argument de \(\rho_{\rm C}\) conserve le cocycle de l’holonomie enrichie ; il n’est pas remplacé par sa seule phase réduite. L’échelle \(\ell_0>0\) est constante dans la coordinatisation. Ce sont les données et la condition de la réalisation lisse utilisée ; préserver la composition des chemins n’identifie pas entre eux des chemins distincts de mêmes extrémités et n’élimine pas leur mémoire. La métrique est \(g=\eta_{IJ}e^I\otimes e^J\), \(\eta=\operatorname{diag}(-1,1,1,1)\).

\begin{proposition}[Courbure, covariance et Bianchi]
\label{hmtcuatro:ant:curvatura} Avec \(F_\omega=\mathrm d\omega+\omega\wedge\omega\) et \(T=\mathrm de+\omega\wedge e\),
\[
 \mathcal F_{\ell_0}=
 \begin{pmatrix}F_\omega&\ell_0^{-1}T\\0&0\end{pmatrix},
 \qquad D_\omega T=F_\omega\wedge e,\quad D_\omega F_\omega=0 .
\]
Un changement de repère \(g_L:\mathcal U\to SO^+(3,1)\) agit par \(e'=g_Le\), \(\omega'=g_L\omega g_L^{-1}-\mathrm dg_L\,g_L^{-1}\) ; alors \(T'=g_LT\), \(F_{\omega'}=g_LF_\omega g_L^{-1}\).
\end{proposition}
\begin{proof}
La multiplication des matrices de formes donne le bloc diagonal \(\mathrm d\omega+\omega\wedge\omega\) et le bloc translationnel \(\ell_0^{-1}(\mathrm de+\omega\wedge e)\). En substituant les champs transformés, les termes contenant \(\mathrm dg_L\) s’annulent. Développer \(D_\omega T\) laisse \((\mathrm d\omega+\omega\wedge\omega)\wedge e\). Dans \(D_\omega F_\omega\), les dérivées de \(\omega\) et les produits cubiques s’annulent deux à deux.
\end{proof}

Le feuillet tritique \(\tau\in\{+1,0,-1\}\) conserve son régime dans la réalisation de Cartan avec les générateurs \(J_{ab},P_a^{(\tau)}\), l’action vectorielle de Lorentz et \([P_a^{(\tau)},P_b^{(\tau)}]=-\tau J_{ab}\). Pour \(\mathcal A_\tau=\frac12\omega^{ab}J_{ab}+\ell^{-1}e^aP_a^{(\tau)}\),
\[
 \mathcal F_\tau=
 \tfrac12(F_\omega^{ab}-\tau\ell^{-2}e^a\wedge e^b)J_{ab}
 +\ell^{-1}T^aP_a^{(\tau)} .
\]
En effet, le crochet de Lorentz donne \(F_\omega\), le crochet mixte donne \(D_\omega e\) et le crochet translationnel donne le terme \(-\tau e\wedge e/(2\ell^2)\). La condition \eqref{hmtcuatro:ant:compatibilidad-holonomia} correspond au feuillet central ; les feuillets extrêmes utilisent leur représentation \(\rho_{{\rm C},\tau}\) et leur condition d’holonomie propres.

\subsection{Différentielle du transport et courant matériel}
\label{hmtcuatro:ant:corriente-transporte}

Dans une réalisation cellulaire finie, soit \(U_a:\mathcal H_{s(a)}\to\mathcal H_{t(a)}\) un transport inversible entre fibres hermitiennes. On ne présuppose pas que tout transport de Lorentz soit unitaire pour le produit énergétique. On choisit une orientation par arête. Pour
\[
 (Df)_a=f_{t(a)}-U_af_{s(a)},\quad r=Df,\quad
 q=\mathsf W r,\quad v_a=U_af_{s(a)},\quad
 E=\langle r,\mathsf W r\rangle ,
\]
le poids \(\mathsf W=\mathsf W^*>0\) peut coupler des arêtes. Ce poids n’est pas \(W_\square\). Les variations correspondent à la représentation de la connexion et au problème de bord choisi.

\begin{proposition}[Variation complète]
\label{hmtcuatro:ant:variacion-completa} Pour une courbe différentiable de données, avec \(A_a=\dot U_aU_a^{-1}\), on a
\begin{equation}
 \dot E=2\operatorname{Re}\sum_a
 \langle q_a,\dot f_{t(a)}-U_a\dot f_{s(a)}-A_av_a\rangle
 +\langle r,\dot{\mathsf W}r\rangle .
 \label{hmtcuatro:ant:variacion}
\end{equation}
À champ et poids fixés, le gradient complet du transport est \(G_a=-2q_av_a^*\) :
\[
 \delta_U E=\sum_a\operatorname{Re}\operatorname{Tr}(G_a^*A_a).
\]
Dans la sous-classe unitaire \(A_a^*=-A_a\), sa partie antihermitienne \(J_a=v_aq_a^*-q_av_a^*\) suffit.
\end{proposition}
\begin{proof}
La règle du produit donne \(\dot E=2\operatorname{Re}\langle\mathsf W r,\dot r\rangle+
\langle r,\dot{\mathsf W}r\rangle\). Dériver chaque résidu produit \eqref{hmtcuatro:ant:variacion}. Comme \(\operatorname{Tr}(v_aq_a^*A_a)=q_a^*A_av_a\), on obtient \(G_a\). La partie hermitienne de \(G_a\) est orthogonale, pour \(\operatorname{Re}\operatorname{Tr}(X^*Y)\), aux variations antihermitiennes. Cette réduction ne s’applique que dans cette sous-classe.
\end{proof}

Pour une route avec \(U_\gamma=U_N\cdots U_1\), soit \(B_k=U_N\cdots U_{k+1}\). La règle du produit démontre
\begin{equation}
 \dot U_\gamma U_\gamma^{-1}
 =\sum_k B_kA_kB_k^{-1}.
 \label{hmtcuatro:ant:variacion-ruta}
\end{equation}
Par conséquent, un gradient \(G_\gamma\) est transporté vers l’arête \(k\) sous la forme \(G_k=B_k^*G_\gamma(B_k^{-1})^*\). La cyclicité de la trace prouve cette formule même si \(B_k\) n’est pas unitaire.

Si les coordonnées réelles de variation de la connexion sont \(\theta_b\), sa différentielle effective s’écrit \(A_a(\theta)=\sum_b B_{ab}\theta_b\). Pour \(S_{\rm mem}=\mathfrak a E\), à champ, poids et section d’action \(\mathfrak a\) fixés,
\begin{equation}
 \delta S_{\rm mem}=-\tfrac12\sum_b\theta_b s_b,\qquad
 s_b=4\mathfrak a\,\operatorname{Re}
          \sum_a\langle q_a,B_{ab}v_a\rangle .
 \label{hmtcuatro:ant:corriente-celular}
\end{equation}
Substituer \(A_a(\theta)\) dans la variation démontre l’égalité et fixe le facteur \(4\). Si l’appariement cellulaire orienté est une matrice inversible \(\mathsf M\), définie par \(\delta S_{\rm mem}=-\frac12\theta^{\mathsf T}\mathsf M\sigma\), le courant est \(\sigma=\mathsf M^{-1}s\). Cette matrice conserve les volumes et l’orientation ; elle ne peut être remplacée par une division scalaire sans fixer une base où celle-ci convient. Si \(f,\mathsf W,\mathfrak a\) varient également, tous les termes de \eqref{hmtcuatro:ant:variacion} et de \(E\delta\mathfrak a\) sont conservés.

Pour un poids par arête, l’inversion conserve la réponse complète en transformant conjointement
\[
 U_{\bar a}=U_a^{-1},\qquad
 r_{\bar a}=-U_a^{-1}r_a,\qquad
 \mathsf W_{\bar a}=U_a^*\mathsf W_aU_a .
\]
La substitution donne \(r_{\bar a}^*\mathsf W_{\bar a}r_{\bar a}=r_a^*\mathsf W_ar_a\). Sa dérivée conserve nécessairement \(\dot{\mathsf W}_{\bar a}
=U_a^*(A_a^*\mathsf W_a+\dot{\mathsf W}_a+\mathsf W_aA_a)U_a\) : éliminer ce terme altérerait le courant. Si le poids couple plusieurs arêtes, la même congruence est appliquée à la matrice complète avec le transport diagonal par blocs ; ses blocs croisés sont également conservés.

\begin{proposition}[Compatibilité variationnelle par raffinement]
\label{hmtcuatro:ant:refinamiento-corriente} Pour des applications fixées \(J_m,K_m\), supposons que, tout au long de la variation considérée,
\[
 D_{m+1}J_m=K_mD_m,\qquad
 K_m^*\mathsf W_{m+1}K_m=\mathsf W_m .
\]
Alors \(E_{m+1}(J_mf)=E_m(f)\) et leurs différentielles coïncident. Si les actions coïncident également et si \(P_m\) transporte les variations de connexion vers le raffinement, leurs courants vérifient
\[
 s_m=P_m^{\mathsf T}s_{m+1},\qquad
 \mathsf M_m\sigma_m=P_m^{\mathsf T}\mathsf M_{m+1}\sigma_{m+1}.
\]
\end{proposition}
\begin{proof}
Substituer \(D_{m+1}J_m=K_mD_m\) dans l’énergie, puis la congruence du poids, prouve la première égalité sur toute la courbe. La dériver et utiliser \(-\frac12\theta^{\mathsf T}s_m
=-\frac12(P_m\theta)^{\mathsf T}s_{m+1}\) prouve les autres. Si \(J_m\) varie, sa différentielle apporte \(2\operatorname{Re}\langle D_{m+1}J_mf,
\mathsf W_{m+1}D_{m+1}\dot J_m f\rangle\) ; elle n’est pas écartée.
\end{proof}

Le courant de mémoire ainsi obtenu provient de son action. Le courant spinoriel qui suit provient de l’action affine des champs. Ils se composent lorsqu’ils sont des termes d’une même action ; ils ne s’identifient pas du seul fait qu’ils partagent le transport. En particulier, éliminer une contorsion dont dépendent également le champ préparé, l’incidence ou son minimum exige de conserver ces différentielles. La formule linéaire de Cartan utilisée ci-après correspond au problème matériel affine expressément déclaré.

\subsection{Représentation spinorielle et courant affine}
\label{hmtcuatro:ant:espinores}

On fixe la structure de spin de la réalisation géométrique. Les matrices héritées, après complexification du module réel, sont
\begin{equation}
 \begin{gathered}
 J=\begin{pmatrix}0&-1\\1&0\end{pmatrix},\quad
 R=\begin{pmatrix}0&1\\1&0\end{pmatrix},\quad
 Z=\begin{pmatrix}1&0\\0&-1\end{pmatrix},\\
 g^0=J\otimes I_2,\quad g^1=R\otimes I_2,\quad
 g^2=Z\otimes R,\quad g^3=Z\otimes Z .
 \end{gathered}
 \label{hmtcuatro:ant:clifford}
\end{equation}
Leurs carrés sont \(-I_4,I_4,I_4,I_4\) et les matrices distinctes anticommutent. Il en résulte \(\{g^I,g^J\}=2\eta^{IJ}I_4\). Nous définissons
\[
 g_I=\eta_{IJ}g^J,\quad
 \mathbb B_{IJ}=\tfrac14[g_I,g_J],\quad
 A_D=-ig^0,\quad \bar\psi=\psi^\dagger A_D,\quad
 \Omega=\tfrac12\omega^{IJ}\mathbb B_{IJ}.
\]
Le développement du commutateur de trois matrices donne \([\mathbb B_{IJ},g^K]=\delta_J^Kg_I-\delta_I^Kg_J\). Par conséquent, \(\Omega\) représente la même connexion métrique, avec \(D\psi=\mathrm d\psi+\Omega\psi\) et \(D\bar\psi=\mathrm d\bar\psi-\bar\psi\Omega\). La matrice \(A_D\) est hermitienne et \(A_Dg^I\) est antihermitienne ; elles fixent l’adjoint et le facteur de l’action en signature \((-+++)\).

Soit \(\eta_I=\iota_{E_I}\operatorname{vol}_e\), de sorte que \(e^J\wedge\eta_I=\delta_I^J\operatorname{vol}_e\). L’action affine est
\begin{equation}
 S_D=\int_M\left\{
 \frac{\hbar}{2}
 [\bar\psi g^ID\psi-(D\bar\psi)g^I\psi]\wedge\eta_I
 -mc\,\bar\psi\psi\operatorname{vol}_e\right\}.
 \label{hmtcuatro:ant:accion-dirac}
\end{equation}
La masse et les sections positives \(c,\hbar\) sont les lectures reçues, et \([\psi]=L^{-3/2}\). La notation alternative \(\Gamma^I=-ig^I\), avec le terme principal \(i\hbar\Gamma^ID_I\), passe à la signature de Clifford opposée ; les deux conventions ne sont pas mélangées.

\begin{proposition}[Courant spinoriel normalisé]
\label{hmtcuatro:ant:corriente-spin} Avec \(B^{IJK}=\bar\psi g^{[I}g^Jg^{K]}\psi\), où l’antisymétrisation est normalisée, la convention \(\delta_\omega S_D=-\frac12\int\delta\omega^{JK}\wedge\sigma_{JK}\) donne
\begin{equation}
 \sigma_{JK}
 =-\frac{\hbar}{2}\bar\psi\{g^I,\mathbb B_{JK}\}\psi\,\eta_I
 =-\frac{\hbar}{2}B^I{}_{JK}\eta_I .
 \label{hmtcuatro:ant:sigma-spin}
\end{equation}
Le courant est indépendant de la contorsion lorsque \(e,\psi\) est fixé.
\end{proposition}
\begin{proof}
On a \(\delta\Omega=\frac12\delta\omega^{JK}\mathbb B_{JK}\). Les deux termes cinétiques de \eqref{hmtcuatro:ant:accion-dirac} donnent
\[
 \delta_\omega S_D=\frac{\hbar}{4}\int
 \delta\omega^{JK}\wedge\eta_I\,
          \bar\psi\{g^I,\mathbb B_{JK}\}\psi .
\]
La comparaison avec la convention variationnelle détermine le signe et le facteur. En développant l’anticommutateur à l’aide de Clifford, les contractions de degré un s’annulent et il reste \(\{g^I,\mathbb B_{JK}\}=g^{[I}g_Jg_{K]}\). L’action est linéaire en \(\Omega\), ce qui prouve l’indépendance.
\end{proof}

Si des connexions internes sont incluses, elles agissent sur le facteur matériel du module spinoriel, tandis que \(\mathbb B_{IJ}\) agit sur le facteur de spin. Ainsi, \([\,\mathbb B_{IJ}\otimes I,I\otimes\rho(A)\,]=0\). Faire varier \(\omega\) produit la même formule, contractée sur les composantes internes et sommée sur les multiplets présents. Cela préserve leurs représentations chirales, de couleur et de famille ; aucun multiplet n’est ajouté à ceux déjà construits.

\subsection{Contraction torsionnelle et termes croisés entre espèces}
\label{hmtcuatro:ant:contraccion}

L’inversion réelle de Cartan--Holst et de l’application \(T^I\mapsto T^I\wedge e^J-e^I\wedge T^J\), démontrée dans \ref{hmtcuatro:gea:torsion}, fixe sans ambiguïté la contorsion \(\mathfrak k_*\) pour le courant affine \eqref{hmtcuatro:ant:sigma-spin}. Ces formules conservent \(\gamma\in\mathbb R\setminus\{0\}\), une cotétrade non dégénérée et les sections \(G,c,\hbar,\gamma\) fixées pendant la variation. L’élimination utilise une seule fois l’interaction linéaire et donne \(-\frac14\mathfrak k_*^{IJ}\wedge\sigma_{IJ}\).

\begin{theorem}[Coefficient de la contraction spinorielle]
\label{hmtcuatro:ant:coeficiente} Soit \(\kappa=8\pi G/c^4\). L’interaction effective est
\[
 \mathcal L_{\rm spin}^{\rm ef}
 =C_\gamma q(B)\operatorname{vol}_e,\qquad
 C_\gamma=\frac{3\kappa c\hbar^2}{16}
                    \frac{\gamma^2}{1+\gamma^2},
\]
où
\[
 q(B)=\frac1{3!}B^{IJK}B_{IJK}
 =-(B^{012})^2-(B^{013})^2-(B^{023})^2+(B^{123})^2 .
\]
\end{theorem}
\begin{proof}
On calcule d’abord
\[
 \mathcal Y^{IJ}
 =\kappa c\frac{\gamma^2}{1+\gamma^2}
       (-\star+\gamma^{-1}I)\sigma^{IJ},\quad
 \mathcal B^I=\sum_J\iota_{E_J}\mathcal Y^{IJ},\quad
 T^I=\mathcal B^I-\tfrac14e^I\wedge\sum_L\iota_{E_L}\mathcal B^L,
\]
puis \(\mathfrak k_{IJK}=(T_{KIJ}-T_{IJK}-T_{JKI})/2\). Ce sont des compositions linéaires des inverses déjà démontrés. Pour expliciter tous les coefficients de la contraction, on ordonne le trivecteur sous la forme \((012,013,023,123)\) et l’on extrait le facteur \(\kappa c\hbar^2\gamma^2/(1+\gamma^2)\). La substitution de \(\sigma_{JK}/\hbar=-\frac12B^I{}_{JK}\eta_I\) dans ces compositions et dans \(-\frac14\mathfrak k^{IJ}\wedge\sigma_{IJ}\) donne
\[
 \begin{array}{c|rrrr|c}
  \text{opérateur sur }\sigma&012&013&023&123&
           \text{six coefficients de polarisation}\\ \hline
  -\star&-3/16&-3/16&-3/16&3/16&0\\
  I&0&0&0&0&0
 \end{array}
\]
Chaque entrée utilise \(e^I\wedge\eta_J=\delta_J^I\operatorname{vol}_e\) et les contractions écrites ci-dessus. Les quatre entrées diagonales et les six polarisations énumèrent la forme quadratique complète. La partie \(\gamma^{-1}I\) s’annule et la partie restante est \((3/16)q(B)\), ce qui prouve la formule.
\end{proof}

Pour plusieurs espèces, la variation s’effectue sur leur action commune : \(\sigma=\sum_s\sigma_s\). Comme la réponse \(\mathfrak k_*=\mathsf C_e(\sigma)\) est linéaire,
\begin{equation}
 -\tfrac14\mathsf C_e\!\left(\sum_s\sigma_s\right)
       \wedge\sum_t\sigma_t
 =-\tfrac14\sum_{s,t}\mathsf C_e(\sigma_s)\wedge\sigma_t .
 \label{hmtcuatro:ant:cruces-especies}
\end{equation}
Ce développement démontre que les termes \(s\ne t\) subsistent. Calculer des réponses séparées et ne conserver que \(s=t\) changerait l’action totale.

\subsection{Ordre normal du courant total}
\label{hmtcuatro:ant:car}

Dans la réalisation fermionique finie d’une cellule, on utilise \(\mathcal F=\bigoplus_n\Lambda^n V\), avec les relations \(\{a_r,a_s^\dagger\}=\delta_{rs}\), \(\{a_r,a_s\}=0\) et le vide d’occupation comme référence de l’ordre normal. Nous écrivons \(d\Gamma(M)=\sum_{r,s}a_r^\dagger M_{rs}a_s\) ; ce \(\Gamma\) désigne la seconde quantification, non l’holonomie nonadique. Pour le module spinoriel à quatre composantes, les matrices sont
\[
 \begin{array}{c|cccc}
 abc&012&013&023&123\\ \hline
 M_{abc}=A_Dg^ag^bg^c&
 iJ\otimes R&iJ\otimes Z&iI_2\otimes J&iZ\otimes J
 \end{array}.
\]
Multiplier \eqref{hmtcuatro:ant:clifford} donne ces matrices ; elles sont hermitiennes, de trace nulle et de carré \(I_4\).

\begin{proposition}[Contraction normale et conservation des termes croisés]
\label{hmtcuatro:ant:orden-normal} Pour toute matrice \(M\),
\[
 :d\Gamma(M)^2:
 =d\Gamma(M)^2-d\Gamma(M^2).
\]
Dans le module à quatre composantes,
\[
 \widehat{\mathscr Q}_{\rm sp}
 =\sum_{a=1}^4s_a d\Gamma(M_a)^2+2\widehat N,
 \qquad s=(-1,-1,-1,+1),\quad \widehat N=d\Gamma(I_4).
\]
L’opérateur est autoadjoint et conserve l’occupation. Dans une somme de multiplets, on emploie la même première identité avec les matrices totales \(\widetilde M_a\) ; les contractions entre espèces sont conservées.
\end{proposition}
\begin{proof}
Dans \(d\Gamma(M)^2\), la relation \(a_sa_t^\dagger=\delta_{st}-a_t^\dagger a_s\) sépare la contraction \(d\Gamma(M^2)\) du terme à deux créateurs et deux annihilateurs. Ce dernier est, par définition, le produit normal. Comme \(M_a^2=I_4\) et \(\sum_a s_a=-2\), la soustraction des contractions vaut \(+2\widehat N\). L’hermiticité des matrices rend les termes autoadjoints et chaque monôme conserve l’occupation.

Si \(\widetilde M=\bigoplus_s M_s\), alors \(d\Gamma(\widetilde M)=\sum_s d\Gamma(M_s)\) et \(d\Gamma(\widetilde M^2)=\sum_s d\Gamma(M_s^2)\). Les bilinéaires de multiplets distincts commutent ; en élevant la première somme au carré, il reste \(2d\Gamma(M_s)d\Gamma(M_t)\) pour \(s<t\). La soustraction d’une occupation n’élimine pas ces termes croisés. C’est la version opératorielle de \eqref{hmtcuatro:ant:cruces-especies}.
\end{proof}

L’invariance sous les changements unitaires internes s’obtient parce qu’ils agissent sur les indices matériels et entrelacent les matrices de courant : la seconde quantification transporte la formule complète par conjugaison. La densité et son évaluation sur un état utilisent en outre le volume de la cellule et l’état matériel déclarés. Le courant moyen et son second moment ne sont pas identifiés. Ces données complètent les fondements utilisés par l’action conjointe, sans introduire une seconde élimination de la torsion ni modifier ses conditions de bord.

% Preparación separada del 30-09-2026; no modifica manuscritos publicados.
% Propietarios: VI_ES/sections/04_atlas_familias.tex:140--163;
% md/04d_ocupacion_estadistica.tex:305--319;
% md/06e_color_qutrit_gauge_finito.tex:1--194.
% Carta qutrit y forma compacta recibidas con sus hipótesis explícitas.
% Inserción: después de 00_antecedentes_geometricos_materiales y antes de 10.
\section{Résidu chromatique et représentation interne de l’action conjointe}
\label{hmtcuatro:color:seccion}

La chaîne APP--TRIT--TPK conserve l’état enrichi, ses routes, orientations, feuillets et mémoire au sein de la structure discrète conjointe du continu. Le résidu lu ci-après est une lecture finie de cet état ; il ne le remplace pas et ne réinitialise pas son histoire. Les coordonnées de phase, y compris l’unité complexe et la coordonnée régionale \(\pi\), sont des sorties HMT reçues avant cette réalisation algébrique. La représentation est construite sur le support résiduel avec la carte, l’orientation et le caractère de phase déclarés. Aucun groupe de jauge n’est inféré du seul nombre d’étiquettes.

\subsection{Support résiduel et orientation conservée}
\label{hmtcuatro:color:residuo}

Dans la copie positionnelle d’APP, soient
\[
 I_9=\{1,\ldots,9\},\qquad
 \mathcal C=I_9\times I_9,\qquad
 \mathcal C_q=\{(r,c)\in\mathcal C:r\text{ est pair}\}.
\]
L’observable résiduel et l’inversion cellulaire sont
\[
 \kappa_{\rm col}(r,c)=r-c\pmod3,\qquad
 \rho_{\rm c}(r,c)=(10-r,10-c).
\]
Le symbole \(\kappa_{\rm col}\) ne désigne pas le coefficient gravitationnel de l’action. Le relèvement de cet observable aux sections persistantes s’effectue par la projection sur l’atlas et conserve les autres registres ; il ne sélectionne pas une cellule par sa masse évaluée.

\begin{proposition}[Fibres ternaires et bilan chromatique]
\label{hmtcuatro:color:balance} La fonctionnelle \(\kappa_3:\mathbb F_3^2\to\mathbb F_3\), \(\kappa_3(r,c)=r-c\), possède trois fibres affines de cardinal trois. Sur \(\mathcal C_q\), les trois résidus ont pour cardinaux \(12,12,12\). L’inversion cellulaire change le signe du résidu.
\end{proposition}
\begin{proof}
Pour chaque \(j\in\mathbb F_3\), la fibre est \(\{(c+j,c):c\in\mathbb F_3\}\), un translaté du noyau \(\{(c,c)\}\). Dans \(\mathcal C_q\), il existe quatre valeurs paires de \(r\) ; pour chacune, \(c=1,\ldots,9\) parcourt trois fois chaque classe modulo trois. Chaque résidu possède \(4\cdot3=12\) représentants. Enfin,
\[
 \kappa_{\rm col}(\rho_{\rm c}(r,c))
 =(10-r)-(10-c)=-(r-c)\pmod3.
\]
La classe zéro reste fixe et les classes un et deux sont échangées.
\end{proof}

Ce bilan n’identifie pas saveur, profondeur et couleur. Il n’affirme pas non plus que les trois droites qui seront construites sont trois états complets : les routes et la mémoire qui partagent une lecture résiduelle demeurent dans la fibre enrichie.

\subsection{Carte qutrit, paire de Weyl et algèbre de couleur}
\label{hmtcuatro:color:weyl}

On fixe la carte
\[
 V_{\rm col}=\mathbb C[\mathbb F_3]\simeq\mathbb C^3,
 \qquad (e_0,e_1,e_2),\qquad
 \omega_3=e^{2\pi i/3},
\]
avec une base delta orthonormée, l’orientation \(j\mapsto j+1\) et le caractère \(j\mapsto\omega_3^j\). Les opérateurs sont
\[
 Xe_j=e_{j+1\bmod3},\qquad Ze_j=\omega_3^j e_j,
 \qquad X^3=Z^3=I,\quad ZX=\omega_3XZ.
\]
La dernière égalité se vérifie en appliquant les deux membres à chaque \(e_j\). L’application de l’atlas vers cette carte est
\[
 (r,c)\longmapsto\mathbb C e_{\kappa_{\rm col}(r,c)}.
\]
Ainsi, \(X\) permute les droites et \(Z\) enregistre leurs phases. Changer la carte affine et transporter conjointement base, orientation et caractère conjugue les opérateurs par la matrice unitaire de changement de base correspondante. Cela ne revient pas à changer une seule étiquette en laissant les autres données fixes.

\begin{theorem}[Algèbre matricielle et forme réelle compacte]
\label{hmtcuatro:color:algebra} Les neuf opérateurs \(W_{ab}=X^aZ^b\), \((a,b)\in\mathbb F_3^2\), forment une base orthogonale de \(\operatorname{End}_{\mathbb C}(V_{\rm col})\). Les huit indices non nuls engendrent linéairement \(\mathfrak{sl}_3(\mathbb C)\), dont la forme réelle de matrices antihermitiennes est \(\mathfrak{su}(3)\).
\end{theorem}
\begin{proof}
Si \(a\ne0\), la matrice \(X^aZ^b\) n’a pas d’entrées diagonales. Si \(a=0\), sa trace est \(1+\omega_3^b+\omega_3^{2b}\), qui vaut zéro pour \(b\ne0\) et trois pour \(b=0\). La relation de Weyl donne alors
\[
 \operatorname{tr}(W_{ab}^{\dagger}W_{cd})
 =3\delta_{ac}\delta_{bd}.
\]
L’indépendance et la dimension neuf prouvent la première assertion ; les huit autres forment une base du sous-espace de trace nulle. De plus,
\[
 W_{ab}^{\dagger}=\omega_3^{ab}W_{-a,-b}.
\]
Les huit indices non nuls forment quatre paires adjointes. En prenant \(\mathcal R=\{(1,0),(0,1),(1,1),(1,2)\}\), les matrices
\[
 A_u=W_u-W_u^{\dagger},\qquad
 B_u=i(W_u+W_u^{\dagger}),\qquad u\in\mathcal R,
\]
sont antihermitiennes et de trace nulle. Leur indépendance réelle découle de l’indépendance complexe des huit \(W_u,W_{-u}\) : dans chaque paire, les coefficients d’une combinaison réelle nulle imposent l’annulation des deux coefficients réels. Elles forment une base réelle de dimension huit. Par conséquent,
\[
 \{A\in\mathfrak{sl}_3(\mathbb C):-A^{\dagger}=A\}
 =\mathfrak{su}(3).
\]
\end{proof}

La forme hermitienne et l’orientation complexe fixées déterminent le groupe de réalisation
\[
 SU(V_{\rm col})=\{U:U^{\dagger}U=I,\ \det U=1\}.
\]
Son algèbre tangente est exactement la précédente : différencier les deux équations à l’identité donne \(A^{\dagger}=-A\) et \(\operatorname{tr}A=0\) ; réciproquement, \(e^{tA}\) satisfait les deux. Tout élément du groupe s’écrit comme l’exponentielle d’une matrice de cette algèbre. En effet, diagonalisons unitairement \(U\) avec les valeurs propres \(e^{i\theta_j}\) ; comme le déterminant vaut un, \(\sum_j\theta_j=2\pi n\). Soustraire \(2\pi n\) à l’un des angles ne modifie pas \(U\) et produit un logarithme antihermitien de trace nulle. Cette réalisation continue ne s’identifie pas au groupe fini engendré uniquement par \(X,Z\). Elle conserve la forme compacte et n’introduit aucune valeur pour le couplage fort.

\subsection{Transport chromatique et facteurs internes indépendants}
\label{hmtcuatro:color:transporte}

Dans une cellularisation finie orientée compatible avec les routes, on attribue \(U_e\in SU(V_{\rm col})\) et \(U_{\bar e}=U_e^{-1}\). Pour \(g_v\in SU(V_{\rm col})\),
\[
 U_e^g=g_{t(e)}U_eg_{s(e)}^{-1},\qquad
 \psi_v^g=g_v\psi_v.
\]
Par substitution,
\[
 U_e^g\psi_{s(e)}^g-\psi_{t(e)}^g
 =g_{t(e)}(U_e\psi_{s(e)}-\psi_{t(e)}).
\]
Le signe opposé, utilisé dans \(D_T\), possède la même covariance. La transformation d’un produit orienté autour d’une face est une conjugaison par la jauge du sommet de base : les facteurs des sommets intermédiaires s’annulent successivement. En particulier, pour \(\beta\ge0\),
\[
 S_{\rm col}^{(\mathcal K,\beta)}[U]
 =\frac\beta3\sum_f(3-\operatorname{Re}\operatorname{tr}U_f)
\]
est invariante et positive ou nulle. L’invariance utilise la trace cyclique ; la positivité résulte de \(\operatorname{Re}\operatorname{tr}U_f\le3\). Le paramètre et la cellularisation conservent leur statut de données de cette réalisation. Cette preuve ne les détermine pas par comparaison avec des valeurs chromodynamiques.

La matière de l’action conjointe utilise ce facteur, non une autre partition de l’atlas. Pour les quarks gauches, sa fibre interne possède les facteurs \(V_{\rm col}\otimes\mathbb C^2\otimes F\), où \(F\) conserve les familles ; pour chaque type droit, on utilise \(V_{\rm col}\otimes F\). Le facteur de couleur est trivial pour les leptons. Le facteur spinoriel est conservé séparément. Pour des matrices de couleur \(A\) et des matrices faibles \(B\), on a
\[
 [A\otimes I,I\otimes B]=0,
\]
car les deux produits valent \(A\otimes B\). Les caractères d’hypercharge agissent comme des scalaires sur chaque multiplet et commutent également. Cette identité permet de composer les connexions internes et celle de spin dans la même dérivée sans confondre couleur, chiralité et saveur.

Les applications de Yukawa de l’action ultérieure sont l’identité sur le facteur chromatique et des applications sur les facteurs Higgs--famille. Par conséquent, \(Y_q(H)(g_{\rm col}\otimes I)
=(g_{\rm col}\otimes I)Y_q(H)\). La compatibilité faible et d’hypercharge y est démontrée avec leurs colonnes et leurs poids explicites ; elle ne résulte pas du recensement \(12+12+12\). On conserve ainsi la chaîne résidu--qutrit--algèbre--transport requise par l’action commune, sans affirmer dans ces fondements une limite quantique de Yang--Mills ni remplacer ses documents sources.

% Fragmento integrable en VII: Relaciones estructurales entre las constantes físicas.
% Insertar después de desarrollos/normalizacion_conjunta.tex.
% Fuentes y cobertura íntegra: editorial/INSERCION_VII.md.
% Sin preámbulo ni macros; usa amsmath y los entornos ya definidos por VII.
\section{Gravité, énergie et autoinertie du système conjoint}
\label{hmtcuatro:gea:seccion}

La normalisation conjointe permet de formuler une question opératorielle : comment la géométrie, la matière et la mémoire interviennent dans une même réalisation. Trois liens sont réunis. Le caractère positif conserve simultanément ses lectures énergétique, massique et radiale ; la variation du transport conjoint produit des courants géométriques et internes d’une même énergie ; et la réponse autoinertielle s’exprime au moyen du rayon gravitationnel de ce caractère. On démontre la conservation de ces liens sous réduction de mémoire et la propagation des contraintes de l’action couplée au cours de son évolution régulière.

Le point de départ est le noyau commun de l’article. APP produit les deux feuillets avec leurs résidus, quotients et incidences ; TRIT conserve régime et orientation : $+1$ correspond au retour et à la mémoire, $0$ au seuil présent et $-1$ à l’ouverture. TPK sélectionne, transporte, plie, retourne et conserve l’état enrichi et son histoire. Le prolongement $w_6\to w_{12}\to w_{18}\to w_{24}\to w_{30}\to R_{36}\to G_9$ maintient préfixe, retenue, survie, route et frontière. La dynamique, son holonomie nonadique et la lecture réduite suivent l’ordre $U_t\to\Gamma_9\to K_{\rm ph}$ : la phase revient, la mémoire avance et l’état n’est pas réinitialisé. Les cinq constructions consubstantielles de la structure discrète du continu ne sont pas séparées lors de la réalisation de leurs lecteurs.

Les constantes et les échelles employées sont des sorties de cette généalogie. La quadrirelation $(\pi,\varphi,e,\alpha)$ conserve son origine commune et son ordre constructif interne ; action, vitesse constitutive et Barbero--Immirzi sont reçus avec leurs lecteurs. Le langage opératoriel et variationnel compose et reconnaît ces sorties, sans introduire de valeurs métrologiques pour sélectionner des états ou des coefficients.

\subsection{Énergie, masse et rayon du même caractère}
\label{hmtcuatro:gea:caracter}

On conserve la réalisation marquée de la section précédente :
\[
\Omega=\{(j,k,q,u,v):j,k\in\mathbb Z/12\mathbb Z,\ 
k-j\in\{0,3,6,9\},\ 1\leq q\leq8,\ 
u<v,\ u,v\in\{1,2,4,5,7,8\}\}.
\]
Son cardinal $N_\Omega=12\cdot4\cdot8\binom62=5760$ compte des incidences, non des états de matière ni des laps. Pour l’inverseur transporté $B:V_{\rm in}\to V_{\rm out}$, avec $B^*B=BB^*=I/4$, le projecteur $P=uu^*$ du mode commun a $|u_i|^2=1/N_\Omega$. L’archive conserve
\[
C_0=(I-P)B,\qquad M_0=PB,\qquad C_0+M_0=B,\qquad\mathsf U=2B.
\]
Les images de $C_0,M_0$ sont orthogonales et $\mathsf U$ est unitaire. L’espérance marquée $\Delta_\Omega(A)=\sum_iP_iAP_i$, $P_i=e_ie_i^*$, est unique parmi les espérances bimodulaires vers l’algèbre diagonale qui fixent les $P_i$ : sur une unité matricielle, la bimodularité exige $\mathcal E(E_{ij})=P_i\mathcal E(E_{ij})P_j$, qui vaut zéro si $i\ne j$ et vaut $E_{ii}$ si $i=j$.

De $C_0C_0^*=(I-P)/4$ et de $P_iPP_i=P_i/N_\Omega$, il résulte
\begin{equation}
\Delta_\Omega(C_0C_0^*)=r_\Omega I,\qquad
\Delta_\Omega(M_0M_0^*)=\frac{I}{4N_\Omega},\qquad
r_\Omega=\frac{5759}{23040}.
\label{hmtcuatro:gea:retorno}
\end{equation}
La somme des deux lectures vaut $I/4$. La règle constitutive longitudinale applique linéairement le premier coefficient à une section de longueur :
\begin{equation}
L=\alpha^{16}r_\Omega L_*,\qquad D=L\mathsf U.
\label{hmtcuatro:gea:longitud}
\end{equation}
La norme locale $\alpha^{16}$ et la section $L_*$ sont des antécédents explicites ; prendre une racine d’intensité constituerait un autre lecteur. La règle métrique radiale est $\|R_Xv\|^2=\|XD^*v\|^2$. Pour $X=X^*>0$ dans la fibre finie, la polarisation et l’unique racine positive donnent $R_X^2=DX^2D^*=L^2\mathsf UX^2\mathsf U^*$ et $R_X=LY$, avec $Y=\mathsf UX\mathsf U^*>0$. Le spectre n’est pas sélectionné à partir d’une réponse gravitationnelle souhaitée.

La même section $\hbar>0$ et l’horloge conjuguée $\omega L=c$ déterminent $E_L=\hbar c/L$ et les lecteurs
\begin{equation}
H_X=E_LY,\qquad M_X=H_X/c^2,\qquad
R_X=LY,\qquad\Lambda_X=LY^{-1}.
\label{hmtcuatro:gea:lectores}
\end{equation}
La notation $H_X$ est réservée à ce lecteur du caractère. L’hamiltonien conjoint $\mathbb H$ introduit ensuite possède un autre domaine et ne s’identifie pas à $H_X$ du seul fait qu’il porte aussi le nom d’énergie.

\begin{proposition}[Coefficient radial commun]
\label{hmtcuatro:gea:coeficiente} Dans l’identification explicite de $R_X$ comme rayon gravitationnel réduit, $R_X=GM_X/c^2$, il existe un unique coefficient compatible :
\begin{equation}
G=\frac{c^3L^2}{\hbar},\qquad
R_X=\frac{G}{c^4}H_X=\frac{G}{c^2}M_X,\qquad
R_X\Lambda_X=L^2I,\quad H_X\Lambda_X=\hbar cI.
\label{hmtcuatro:gea:radial}
\end{equation}
\end{proposition}
\begin{proof}
Multiplier les lecteurs prouve les deux produits. L’identification radiale équivaut à $LY=G\hbar Y/(c^3L)$. Simplifier par le caractère inversible détermine $G$ ; la substitution inverse vérifie l’égalité pour chaque caractère positif. Si deux coefficients la satisfont, leur différence multipliée par $Y$ est nulle et ils coïncident.
\end{proof}

Les règles longitudinale et métrique fixent la réalisation ; la convention du rayon réduit fixe la signification physique de $G$. Le rayon de Schwarzschild de la même masse est $2R_X$. L’égalité conserve les superpositions et les éléments non diagonaux dans toute base transportée, et non les seules valeurs moyennes. Pour une famille différentiable et une connexion de repère commune $\mathcal A_t$, avec $G,c$ fixés,
\begin{equation}
\mathcal D_tR_X=\frac{G}{c^4}\mathcal D_tH_X,\qquad
\mathcal D_tM_X=c^{-2}\mathcal D_tH_X,\qquad
\mathcal D_tO=\dot O+[\mathcal A_t,O].
\label{hmtcuatro:gea:derivada}
\end{equation}
Cela s’obtient en différenciant l’identité complète, repère compris. Cela n’implique pas l’immobilité de l’état.

Pour des modes positifs, ou pour des opérateurs sur des facteurs tensoriels distincts,
\begin{equation}
m(y)=\frac{\hbar}{cL}y,\qquad
\frac{Gm_1m_2}{\hbar c}=y_1y_2,\qquad
\frac{R(y)}{\bar\lambda(y)}=y^2,\quad \bar\lambda(y)=L/y.
\label{hmtcuatro:gea:adimensional}
\end{equation}
La substitution annule les sections dimensionnelles ; les opérateurs sur des facteurs distincts commutent. La taille adimensionnelle dépend du caractère, sans autre paramètre libre par espèce. Il n’en découle ni loi spatiale de force ni évaluation numérique sans évaluer le caractère.

Le calendrier conserve $t_P=L/c$, $t_0=\pi L/(54c)$ et $108t_0=2\pi t_P$. Par conséquent, $L=(54/\pi)\ell_0$, avec $\ell_0=ct_0$ ; les deux longueurs ne sont pas identifiées. Entre les sections $\hbar_b=\rho\hbar_a$, $\rho>0$, avec $L,c,y_i$ fixés, on a $G_b=G_a/\rho$, $m_{i,b}=\rho m_{i,a}$ et la quantité $Gm_1m_2/(\hbar c)$ demeure invariante. On compare des sections conjointes, non une variation temporelle dont les autres coordonnées dimensionnelles seraient artificiellement maintenues fixes.

\subsection{Mémoire statique, dynamique et domaine spectral}
\label{hmtcuatro:gea:memoria}

En dimension finie, ou avec des blocs bornés et coercifs, écrivons
\[
Y=\begin{pmatrix}A&F\\F^*&C\end{pmatrix}>0,\quad
Y_{\rm ef}=A-FC^{-1}F^*,\quad \eta=y+C^{-1}F^*b.
\]
La complétion du carré démontre
\begin{equation}
\left\langle\binom b y,Y\binom b y\right\rangle
=\langle b,Y_{\rm ef}b\rangle+\langle\eta,C\eta\rangle.
\label{hmtcuatro:gea:schur-prueba}
\end{equation}
D’où $Y_{\rm ef}>0$, l’extension minimisante $y=-C^{-1}F^*b$ et la reconstruction $y=\eta-C^{-1}F^*b$. Le résidu $\eta$ est conservé. Comme $\operatorname{Schur}(sY)=sY_{\rm ef}$ pour $s>0$, on obtient $H_{X,\rm ef}=E_LY_{\rm ef}$ et $R_{X,\rm ef}=LY_{\rm ef}$. Résoudre $Y(x,y)=(f,0)$ montre que le bloc de frontière de $Y^{-1}$ est $Y_{\rm ef}^{-1}$, donc $\Lambda_b=LY_{\rm ef}^{-1}$. La compression directe de $Y$ donnerait $A$ : ce n’est pas la même opération.

Soient $\chi=L/E_L=G/c^4>0$ et $R_X=\chi H_X$. Dans une décomposition compatible avec les domaines, nous définissons
\[
\mathcal K_{H_X}(z)=H_{bb}-zI-V(H_{ii}-zI)^{-1}V^*.
\]
\begin{proposition}[Mémoire dynamique complète]
\label{hmtcuatro:gea:resolvente-prop} Dans le domaine résolvant du bloc intérieur,
\begin{equation}
\mathcal K_{R_X}(\chi z)=\chi\mathcal K_{H_X}(z).
\label{hmtcuatro:gea:resolvente}
\end{equation}
Si les deux compléments sont inversibles, leurs inverses portent le facteur $\chi^{-1}$.
\end{proposition}
\begin{proof}
Chaque bloc radial vaut $\chi$ fois le bloc énergétique et $(\chi H_{ii}-\chi zI)^{-1}=\chi^{-1}(H_{ii}-zI)^{-1}$. La substitution laisse un facteur $\chi$ dans chaque terme, en conservant l’ordre. L’inversion prouve la dernière assertion.
\end{proof}
Le paramètre énergétique $z$ est transporté vers la longueur $\chi z$ ; l’identité conserve chaque ordre de mémoire spectrale. Pour $|z|<\|H_{ii}^{-1}\|^{-1}$, la série de la résolvante donne
\[
\mathcal K_{H_X}(z)=H_{\rm est}-zZ+O(z^2),\quad
H_{\rm est}=H_{bb}-VH_{ii}^{-1}V^*,\quad
Z=I+VH_{ii}^{-2}V^*\geq I.
\]
La positivité découle de $Z-I=(H_{ii}^{-1}V^*)^*(H_{ii}^{-1}V^*)$. L’extension statique $(b,-H_{ii}^{-1}V^*b)$ a pour norme au carré $\langle b,Zb\rangle$ : énergie et norme temporelle proviennent de la même extension.

Si $T=Z^{1/2}$ est borné et inversible, le transport est dual :
\begin{equation}
H_{X,c}=T^{-*}H_{X,\rm ef}T^{-1},\quad
R_{X,c}=T^{-*}R_{X,\rm ef}T^{-1},\quad
\Lambda_c=T\Lambda_bT^*.
\label{hmtcuatro:gea:dual}
\end{equation}
La multiplication prouve $R_{X,c}=\chi H_{X,c}$, $R_{X,c}\Lambda_c=L^2I$ et $H_{X,c}\Lambda_c=\hbar cI$, sans commutativité de $T$ avec le caractère. Si $\psi=T(t)b$, $\dot\psi=\dot Tb+T\dot b$ est conservé. Le terme temporel de connexion du générateur n’est pas incorporé rétrospectivement au caractère radial sans déclarer son application.

\begin{proposition}[Extension spectrale et raffinement compatible]
\label{hmtcuatro:gea:limite} Soit $X$ autoadjoint, positif et injectif, et $\mathsf U$ unitaire. Les lecteurs de \eqref{hmtcuatro:gea:lectores}, définis par le calcul spectral de $Y=\mathsf UX\mathsf U^*$, satisfont $R_X=\chi H_X$ sur $\operatorname{Dom}(Y)$. Leurs produits avec $\Lambda_X$ sont valables sur $\operatorname{Dom}(Y^{-1})$ et possèdent les prolongements bornés de \eqref{hmtcuatro:gea:radial}. Si des plongements isométriques satisfont $X'J=JX$, $\mathsf U'J=K\mathsf U$ et conservent $L,E_L$, alors $R'_XK=KR_X$ et $H'_XK=KH_X$ sur les domaines transportés.
\end{proposition}
\begin{proof}
Sur $P_{\varepsilon,M}=\mathbf1_{[\varepsilon,M]}(Y)$, lecteurs et inverses sont bornés et les identités sont des multiplications de fonctions de $Y$. Pour $\psi\in\operatorname{Dom}(Y)$,
\[
\|Y(P_{\varepsilon,M}-I)\psi\|^2
=\int_{(0,\infty)\setminus[\varepsilon,M]}
\lambda^2\,d\langle\psi,E_Y(\lambda)\psi\rangle\longrightarrow0.
\]
L’injectivité élimine la masse spectrale en zéro ainsi que $P_{\varepsilon,M}\psi\to\psi$. Il s’agit d’une convergence en norme de graphe. Si $\psi\in\operatorname{Dom}(Y^{-1})$, alors $Y^{-1}\psi\in\operatorname{Dom}(Y)$ et $YY^{-1}\psi=\psi$, ce qui prouve les produits et leurs prolongements. Substituer les entrelacements dans $R'_X=L\mathsf U'X'\mathsf U'^*$ prouve $R'_XK=KR_X$ ; le calcul énergétique est identique.
\end{proof}
Le secteur nul demeure séparé. On précise ainsi l’extension non bornée de la normalisation précédente. Une projection orthogonale arbitraire n’est pas pour autant admise dans les formules par blocs : pour les opérateurs non bornés, des domaines compatibles ou des formes fermées à blocs contrôlés sont requis.

\subsection{Rétroaction dans le même transport}
\label{hmtcuatro:gea:retroaccion}

Le registre géométrique $x$, les liens internes $u$ et les champs $f$ proviennent du même état enrichi. Dans une cellularisation orientée,
\begin{equation}
E(f,x,u)=\langle r,W(x,u)r\rangle,\quad r=D_Tf,\quad
(D_Tf)_a=f_{t(a)}-T_af_{s(a)},\quad
T_a=U_a^{\rm sp}(x)\otimes R_{\rm int}(u_a).
\label{hmtcuatro:gea:energia-memoria}
\end{equation}
$W=W^*>0$ conserve les blocs entre incidences. Spin et représentation interne possèdent des facteurs distincts ; la réalisation chirale conserve ses blocs et l’application Higgs--Yukawa équivariante.

\begin{proposition}[Différentielle d’une même énergie]
\label{hmtcuatro:gea:variacion} À champs fixés et avec $q=Wr$,
\begin{align}
\delta E&=-2\operatorname{Re}\sum_a
\langle q_a,\delta T_af_{s(a)}\rangle+\langle r,\delta Wr\rangle,
\label{hmtcuatro:gea:diferencial}\\
\delta T_a&=\delta U_a^{\rm sp}\otimes R_{\rm int}(u_a)
+U_a^{\rm sp}\otimes\delta R_{\rm int}(u_a).\nonumber
\end{align}
\end{proposition}
\begin{proof}
Dériver la forme quadratique donne deux termes conjugués par $W=W^*$ et le terme $\langle r,\delta Wr\rangle$. Utiliser $\delta r_a=-\delta T_af_{s(a)}$ prouve la première formule ; la règle du produit prouve la seconde. Les deux variations agissent sur les mêmes $r,q$, sans courants choisis indépendamment.
\end{proof}

\subsubsection{Un espace quantique et un opérateur conjoint}
\label{hmtcuatro:gea:cuantica}

On fixe une cellularisation finie, un ensemble fini de registres géométriques $\mathcal X$ et un espace de Fock fermionique fini $\mathcal F$. Les liens parcourent le produit compact $\mathcal G^E$ reçu ; les doublets $\Phi_v\in\mathbb C^2$ n’ont pas de coupure d’amplitude. Les mesures invariantes $d\mu_x=Z_x^{-1}e^{-S_x}d\mu_0$ possèdent des densités positives bornées à inverse borné et ne dépendent pas de $\Phi$ dans cette carte. L’espace est
\[
\mathscr H=\bigoplus_{x\in\mathcal X}
L^2(\mathcal G^E\times\mathbb C^{2|V|},
\mu_x\otimes d\Phi;\mathcal F).
\]
Les identifications $A_xf=\sqrt{Z_x}e^{S_x/2}f$ sont unitaires à partir de la mesure commune et transportent l’horloge vers $H_{\rm clk}^{\mu}=A(H_{\rm clk}\otimes I)A^*$. L’horloge mélange des registres géométriques du même espace. L’opérateur est réalisé au moyen de la forme de
\begin{equation}
\mathbb H=H_{\rm clk}^{\mu}+\hbar cK_{\rm gauge}
+d\Gamma(D_T^*WD_T)+H_{\Phi,\rm cin}+V_\Phi
+\widehat H_Y+\widehat H_{\rm tor}.
\label{hmtcuatro:gea:total}
\end{equation}
L’application de Yukawa est linéaire en $\Phi,\bar\Phi$ et satisfait $Y(g\Phi)\rho_R(g)=\rho_L(g)Y(\Phi)$. La contorsion éliminée n’est pas conservée simultanément comme variable indépendante dans les termes de base de cette carte effective.

La cinétique des moments scalaires a la forme $\int\langle\nabla_\Phi\Psi,A(x,u)\nabla_\Phi\Psi\rangle$, avec $A$ bornée, uniformément positive et indépendante de $\Phi$. On exige $A(x,u^g)=O_gA(x,u)O_g^{\mathsf T}$ pour l’action réelle $O_g$ des doublets ; la positivité seule ne prouve pas cette covariance. L’énergie spatiale scalaire est un potentiel quadratique positif majoré par $C_{\rm grad}r^2$, $r^2=\sum_v|\Phi_v|^2$, non un opérateur borné sur tout $L^2$. Les volumes satisfont $0<w_{\min}\leq w_v(x)\leq w_{\max}$ et $\lambda_\Phi>0$. L’horloge, le bloc de jauge, les transports, $W$ et les coefficients torsionnels sont bornés dans la carte. Le potentiel conserve sa valeur absolue :
\begin{equation}
V_\Phi=\sum_vw_v\lambda_\Phi(|\Phi_v|^2-v_0^2/2)^2
-\sum_vw_v\lambda_\Phi v_0^4/4.
\label{hmtcuatro:gea:potencial}
\end{equation}
La dernière contribution dépend du volume et n’est pas soustraite.

\begin{theorem}[Forme fermée et évolution conjointe]
\label{hmtcuatro:gea:forma} La forme de \eqref{hmtcuatro:gea:total} est fermée et minorée sur $\mathcal Q=\{\Psi\in\mathscr H:\nabla_\Phi\Psi\in L^2,\
r^2\Psi\in L^2\}$. Elle détermine un unique opérateur autoadjoint associé $\mathbb H$ et l’évolution unitaire $e^{-it\mathbb H/\hbar}$.
\end{theorem}
\begin{proof}
Cauchy--Schwarz donne $\sum_vw_v\lambda_\Phi|\Phi_v|^4\geq ar^4$, $a=\lambda_\Phi w_{\min}/|V|>0$. L’espace de Fock fini et Yukawa linéaire donnent $\|\widehat H_Y(\Phi)\|\leq C_Yr$ ; la torsion est bornée. Pour $\varepsilon>0$,
\[
C_Yr\leq\varepsilon r^4+
\frac{3C_Y^{4/3}}{4^{4/3}\varepsilon^{1/3}},\qquad
br^2\leq\varepsilon r^4+\frac{b^2}{4\varepsilon}.
\]
La première borne maximise $C_Yr-\varepsilon r^4$ ; la seconde complète le carré en $r^2$. Les appliquer avec $\varepsilon=a/4$ à Yukawa et au terme quadratique négatif donne
\[
q(\Psi)\geq c_1\|\nabla_\Phi\Psi\|^2+
(a/2)\|r^2\Psi\|^2-C\|\Psi\|^2,\qquad c_1>0.
\]
La norme $\|\Psi\|^2+\|\nabla_\Phi\Psi\|^2+\|r^2\Psi\|^2$ est complète par fermeture de la dérivée faible et de la multiplication. Ces bornes et les majorations rendent la norme de forme, après ajout d’une constante, équivalente à cette norme complète. Le théorème de représentation des formes fermées produit l’opérateur associé et son calcul spectral donne l’évolution. L’unicité concerne cette forme et ce domaine, non chaque domaine formel alternatif.
\end{proof}

La covariance de $D_T,W$, l’invariance des mesures, l’identité de Yukawa et la contraction torsionnelle interne conservent $q$ et $\mathcal Q$. L’unicité de l’opérateur associé implique que la représentation compacte commute avec $\mathbb H$ ; sa moyenne de Haar réduit l’opérateur et son évolution. Ce n’est pas une moyenne corrective ultérieure.

On retire la coupure computationnelle au moyen de blocs complets de Peter--Weyl et de couches isotropes complètes d’Hermite. La régularisation et la troncature lisse approchent les dérivées et le poids $r^2$ ; les densités bornées préservent des normes équivalentes. Leur réunion est dense dans $\mathcal Q$. Pour les projections orthogonales $P_n$ dans les mesures $\mu_x$ et les opérateurs de Galerkin $\mathbb H_n$,
\[
(\mathbb H_n+\lambda)^{-1}P_nf\longrightarrow
(\mathbb H+\lambda)^{-1}f\qquad(\lambda>C).
\]
L’équation coercive de forme et sa version de Galerkin donnent l’orthogonalité de l’erreur ; la continuité et la coercivité la majorent par la meilleure erreur d’approximation, qui tend vers zéro. On retire les coupures scalaire et fonctionnelle des liens de cette carte, non les coupures spatiale, géométrique ou fermionique. L’autoadjonction n’implique pas non plus que $e^{-t\mathbb H}$ soit à trace.

\subsubsection{Échange et conservation}
\label{hmtcuatro:gea:intercambio}

Pour $V=\sum_xP_x\otimes B_x$ et un élément $h_{xy}$ de l’horloge, la multiplication par blocs donne
\begin{equation}
[H_{\rm clk}\otimes I,V]_{xy}=h_{xy}(B_y-B_x).
\label{hmtcuatro:gea:conmutador}
\end{equation}
Si $h_{xy}\ne0$ et $B_y\ne B_x$, il y a échange géométrique--matériel. Si $\mathbb H=A+B+V$, pour des opérateurs bornés ou un cœur commun justifiant les dérivées,
\[
\frac d{dt}\langle A\rangle=\frac i\hbar\langle[\mathbb H,A]\rangle,
\quad
\frac d{dt}\langle B\rangle=\frac i\hbar\langle[\mathbb H,B]\rangle,
\quad
\frac d{dt}\langle V\rangle=\frac i\hbar\langle[\mathbb H,V]\rangle.
\]
La somme est nulle par $[\mathbb H,\mathbb H]=0$. L’énergie est conservée, interaction comprise, non chaque secteur séparément. La géométrie est dans l’espace quantique, non ajoutée comme fond extérieur. Les histoires étendues conservent mémoire et retour ; leur raffinement n’équivaut pas par définition au raffinement spatial.

\subsection{La même action et ses courants}
\label{hmtcuatro:gea:accion}

La réalisation géométrique emploie une cotétrade non dégénérée $e$, une connexion de Lorentz $\omega$ et des connexions internes de la même matière. Soient $\Sigma^{IJ}=e^I\wedge e^J$, $P_\gamma=\star+\gamma^{-1}I$, $\gamma\in\mathbb R\setminus\{0\}$ et $\star^2=-I$ sur les bivecteurs. Les couplages restent fixes pendant la variation. Avec $\kappa=8\pi G/c^4$, la même action est
\begin{equation}
S_{\rm tot}=\frac{\hbar}{16\pi L^2}
\int\Sigma_{IJ}\wedge(P_\gamma F_\omega)^{IJ}
-\frac{\Lambda\hbar}{8\pi L^2}\int\operatorname{vol}_e+S_m,
\label{hmtcuatro:gea:accion-total}
\end{equation}
où $S_m=S_{\rm YM}+S_\Phi+S_{\rm Weyl}+S_Y$ contracte ses indices avec la même $e$, que l’on fait varier. Les coefficients découlent de $(2\kappa c)^{-1}=\hbar/(16\pi L^2)$.

Avec $\delta_\omega S_m=-\tfrac12\int
\delta\omega^{IJ}\wedge\sigma_{IJ}$, $\delta F=D_\omega\delta\omega$ et l’intégration par parties,
\begin{equation}
D_\omega(P_\gamma\Sigma)=\kappa c\sigma.
\label{hmtcuatro:gea:cartan}
\end{equation}
On utilise des variations intérieures ou la complétion de bord compatible. Le courant de spin et la source métrique sont des dérivées de la même matière, non deux entrées indépendantes.

\subsubsection{Élimination torsionnelle et termes croisés}
\label{hmtcuatro:gea:torsion}

L’inverse réel est $P_\gamma^{-1}=\gamma^2(-\star+\gamma^{-1}I)/(1+\gamma^2)$ : le produit des deux polynômes en $\star$ avant normalisation vaut $(1+\gamma^{-2})I$. Notons $\mathcal Y^{IJ}=\kappa c(P_\gamma^{-1}\sigma)^{IJ}$ la forme à valeurs bivectorielles de degré trois, distincte du caractère $Y$. Pour $T^I=D_\omega e^I$, $D_\omega\Sigma^{IJ}=T^I\wedge e^J-e^I\wedge T^J$. Avec le repère dual $E_I$, son inverse est
\[
\mathcal B^I=\sum_J\iota_{E_J}\mathcal Y^{IJ},\qquad
\vartheta=\tfrac14\sum_I\iota_{E_I}\mathcal B^I,\qquad
T^I=\mathcal B^I-e^I\wedge\vartheta.
\]
En effet, si $t=\sum_J\iota_{E_J}T^J$, les contractions donnent $\mathcal B^I=T^I+e^I\wedge t$ puis $4t$. L’application a un noyau nul entre espaces de dimension vingt-quatre ; c’est donc un isomorphisme. Pour $\omega=\mathring\omega(e)+\mathfrak k$,
\[
T_{IJK}=\mathfrak k_{IKJ}-\mathfrak k_{IJK},\qquad
\mathfrak k_{IJK}=\tfrac12(T_{KIJ}-T_{IJK}-T_{JKI}).
\]
La combinaison cyclique et l’antisymétrie dans les deux premiers indices prouvent ces formules. Pour une matière affine en $\mathfrak k$, avec un courant qui en est indépendant, il existe une unique $\mathfrak k_*[e,\sigma]$ linéaire en le courant total $\sigma=\sum_s\sigma_s$.

Le développement de la courbure conserve le bord
\[
S_\partial[e,\mathfrak k]=\frac1{2\kappa c}
\int_{\partial M}\Sigma_{IJ}\wedge(P_\gamma\mathfrak k)^{IJ}.
\]
La densité intérieure dépendant de la contorsion est $\mathcal Q_{e,\gamma}(\mathfrak k)
-\tfrac12\mathfrak k^{IJ}\wedge\sigma_{IJ}$, où $\mathcal Q$ est homogène de degré deux. Sa stationnarité évaluée en $\delta\mathfrak k=\mathfrak k_*$ donne $2\mathcal Q(\mathfrak k_*)=
\tfrac12\mathfrak k_*^{IJ}\wedge\sigma_{IJ}$. Par conséquent,
\begin{equation}
\mathcal L_{\rm spin}^{\rm ef}
=-\tfrac14\mathfrak k_*^{IJ}\wedge\sigma_{IJ}
=-\mathcal Q_{e,\gamma}(\mathfrak k_*).
\label{hmtcuatro:gea:spin}
\end{equation}
La linéarité conserve les termes croisés entre espèces. La contorsion est éliminée une fois ; on n’ajoute pas simultanément son interaction linéaire éliminée et le terme quartique effectif.

Dans la spécialisation spinorielle, le coefficient reçu $C_\gamma=3\kappa c\hbar^2\gamma^2/[16(1+\gamma^2)]$ donne
\begin{equation}
\frac{C_\gamma}{\hbar}
=\frac{3\pi}{2}L^2\frac{\gamma^2}{1+\gamma^2}.
\label{hmtcuatro:gea:coeficiente-torsion}
\end{equation}
Il suffit de substituer $\kappa=8\pi G/c^4$ et $\hbar G=c^3L^2$, sans changer la branche de $\gamma$. Dans l’espace de Fock de la carte, le courant total conserve la contraction
\[
\widehat Q_{{\rm tot},v}=\sum_{a=1}^4s_a
\{d\Gamma(\widetilde M_{a,v})^2-d\Gamma(\widetilde M_{a,v}^2)\},
\quad s=(-1,-1,-1,+1),
\]
\[
\widehat H_{\rm tor}
=-\sum_v\frac{N_vcC_\gamma}{\mathcal V_v}\widehat Q_{{\rm tot},v}.
\]
Les matrices $\widetilde M_{a,v}$ représentent les composantes du courant. $N_v$ est le lapse, distinct de $N_\Omega$ et d’une occupation. L’identité d’ordre normal retire la contraction à une particule, non les termes croisés entre espèces. La signature demeure dans $s_a$.

\subsubsection{Source métrique et transformation de Legendre de la même action}
\label{hmtcuatro:gea:legendre}

L’élimination stationnaire préserve la variation des autres variables : la règle de la chaîne contient un terme en $\delta\mathfrak k_*$ multiplié par l’équation de connexion, qui est nulle. L’action réduite est
\[
S_{\rm red}=\frac1{2\kappa c}\int(R[g]-2\Lambda)\operatorname{vol}_e
+S_m^{\rm LC}+S_{\rm spin}^{\rm ef}+S_\partial.
\]
Nous définissons la source d’action par $\delta_g(S_m^{\rm LC}+S_{\rm spin}^{\rm ef})
=-\tfrac12\int\mathcal T^S_{\mu\nu}\delta g^{\mu\nu}
\operatorname{vol}_e$. La variation gravitationnelle intérieure produit
\begin{equation}
G_{\mu\nu}+\Lambda g_{\mu\nu}
=\kappa c\mathcal T^S_{\mu\nu}
=\kappa(T^{\rm LC}_{\mu\nu}+U^{\rm spin}_{\mu\nu}),
\quad T^{\rm LC}+U^{\rm spin}=c\mathcal T^S.
\label{hmtcuatro:gea:metrica}
\end{equation}
La variation de volume du terme constant du potentiel est incluse. Le facteur $c$ correspond à la cotétrade dans la droite de longueur ; les sources d’action et les sources physiques ne sont pas mélangées.

Avec $a=(2\kappa c)^{-1}$ et $\mathcal B_{IJ}=(P_\gamma\Sigma)_{IJ}$, le potentiel géométrique est $\theta_{\rm geom}=a\mathcal B_{IJ}\wedge\delta\omega^{IJ}$. Décomposer $\omega=\omega_0dt+\underline\omega$ et $F_{0a}=\dot\omega_a-D_a\omega_0$ identifie le terme cinétique ; intégrer par parties conserve Gauss et le bord. Avec $e_0=Nn+N^ae_a$ et la transformation de Legendre de la même matière,
\begin{equation}
S_{\rm tot}=\int dt\,[\Theta_{\rm red}(\dot z)
-\mathcal C[N]-\mathcal V[N^a]-\mathcal G_L[\omega_0]
-\mathcal G_{\rm int}[A_0]]+S_\partial.
\label{hmtcuatro:gea:canonica}
\end{equation}
Lapse, shift et jauge conservent leurs directions dégénérées. Les contraintes de seconde classe sont réduites avant d’utiliser $\Theta_{\rm red}$ ou sont expressément conservées. Les spineurs du premier ordre ne reçoivent pas d’inverse fictif des vitesses.

Pour la contorsion et son moment, les contraintes sont $\pi_{\mathfrak k}=0$ et $\partial\mathcal Q/\partial\mathfrak k-\sigma/2=0$. Leur matrice possède les blocs $\left(\begin{smallmatrix}0&-M\\M^{\mathsf T}&B\end{smallmatrix}\right)$, avec $M$ inversible par la reconstruction précédente ; elle est donc inversible. Le crochet de Dirac effectue cette élimination et conserve séparément les contraintes fermioniques graduées. Dans une direction bosonique régulière,
\[
H_\Phi=N\left(\frac{p^2}{2w}+wV+H_{\rm grad}\right),\quad
p=\frac wN D_t\phi,\quad
L_\Phi=\frac{w}{2N}|D_t\phi|^2-NwV-NH_{\rm grad}.
\]
Le volume $w=\sqrt{\det q}$ n’est pas figé. Pour $K_{ij}=(\dot q_{ij}-\mathcal L_{N^a}q_{ij})/(2N)$,
\[
\Pi^{ij}=\frac{\sqrt q}{2\kappa c}(K^{ij}-q^{ij}K),\qquad
K_{ij}=\frac{2\kappa c}{\sqrt q}
(\Pi_{ij}-\tfrac12q_{ij}\Pi).
\]
La trace et la partie sans trace prouvent l’inverse. Si le moment total contient une contribution matérielle de repère, on utilise $\Pi=p-p_m$, en conservant $p_m$. On retrouve ainsi la même action et sa source. La transformation de Legendre ne démontre pas à elle seule qu’une quantification arbitraire de \eqref{hmtcuatro:gea:canonica} coïncide avec \eqref{hmtcuatro:gea:total}.

\subsection{Propagation des contraintes avec la source totale}
\label{hmtcuatro:gea:restricciones}

Soit $E_{\mu\nu}=G_{\mu\nu}+\Lambda g_{\mu\nu}
-\kappa c\mathcal T^S_{\mu\nu}$. Sur les équations matérielles, l’invariance par difféomorphismes de la fonctionnelle réduite implique $\nabla^\mu\mathcal T^S_{\mu\nu}=0$ : une variation intérieure engendrée par $\xi$ laisse l’intégrale de $\mathcal T^S_{\mu\nu}\nabla^\mu\xi^\nu$, tandis que les équations matérielles annulent les autres termes. Intégrer par parties et faire varier $\xi$ prouve la divergence nulle, avant d’imposer $E=0$. Bianchi et les couplages fixes donnent $\nabla_\mu E^{\mu\nu}=0$.

\begin{proposition}[Conservation des données de contrainte]
\label{hmtcuatro:gea:propagacion} Une solution régulière des équations matérielles et d’évolution métrique conserve les contraintes satisfaites initialement pendant son intervalle d’existence régulière.
\end{proposition}
\begin{proof}
En coordonnées gaussiennes $g=-d\tau^2+q_{ij}(\tau,x)dx^idx^j$, on impose $E_{ij}=0$ et l’on écrit $C=E_{00}$, $C_i=E_{0i}$, $K_{ij}=\dot q_{ij}/2$ et $\mathcal K=q^{ij}K_{ij}$. On a $E^{00}=C$, $E^{0i}=-C^i$, $E^{ij}=0$, $\Gamma^i{}_{0j}=K^i{}_j$ et $\Gamma^\mu{}_{\mu0}=\mathcal K$. Les composantes temporelle et spatiale de la divergence donnent
\[
\partial_\tau C-D_iC^i+\mathcal KC=0,\qquad
\partial_\tau C^i+\mathcal KC^i+2K^i{}_jC^j=0.
\]
En abaissant l’indice, $\dot q_{ij}=2K_{ij}$ annule le dernier terme spatial. Il reste exactement
\begin{equation}
\partial_\tau C-D_iC^i+\mathcal KC=0,\qquad
\partial_\tau C_i+\mathcal KC_i=0.
\label{hmtcuatro:gea:propagacion-sistema}
\end{equation}
La seconde équation homogène maintient $C_i=0$ et la première devient $\dot C+\mathcal KC=0$, en conservant $C=0$. $q$ et $\mathcal K$ sont ceux de la solution couplée, non un fond prescrit.
\end{proof}
La carte gaussienne prouve localement une assertion tensorielle, sans figer $q^{-1}$ ni remplacer la source. On n’affirme pas d’existence globale sans singularités pour chaque donnée. La propagation variationnelle n’est pas non plus à elle seule une identité de commutateurs quantiques pour des déformations normales arbitraires.

\subsection{Frontière et opérateur autoinertiel}
\label{hmtcuatro:gea:autoinercia}

Soient $b:\mathcal X_\infty^{\rm enr}\to B_{\rm ef}$, $B_{\rm ef}=\operatorname{im}(b)$. Une réponse $I_v$ se factorise de manière unique sous la forme $I_v=\overline I_v\circ b$ si et seulement si elle est constante sur les fibres de $b$. La nécessité est immédiate ; pour la suffisance, $\overline I_v(\beta)=I_v(x)$ avec $b(x)=\beta$ est indépendant du représentant. L’image effective donne l’unicité. La dépendance machienne forte ajoute des états ayant la même restriction locale et des réponses distinctes.

La charnière du même état produit $r_{\rm av}=\pi/\varphi^2$ et $r_{\rm av}^J=\varphi/\pi^2$. Dans le secteur d’amplitude commune, $\mathcal L_{\rm ar}=\pi\mu_{\rm ar}\mathcal M(x)$ et $\mathcal L_{\rm vol}=\mu_{\rm vol}\mathcal M(x)$, $\mathcal M(x)>0$, avec $\mu_{\rm ar}/\mu_{\rm vol}=\varphi^{-2}$. Diviser par $\mathcal L_{\rm vol}$ donne des lecteurs du même type $A_\partial=r_{\rm av}$ et $V_\partial=1$. Le quotient annule l’amplitude, non l’histoire conservée dans l’archive.

Pour une orientation $a$, $q_a(x)=h_a(x)/(108t_0(x))=\hbar_a(x)/t_P(x)$ descend à $Q:B_{\rm ef}\to\mathcal L_E^+$ si et seulement si $b(x)=b(x')$ implique $h_a(x)/t_0(x)=h_a(x')/t_0(x')$. C’est le critère de factorisation appliqué à l’horloge. Les cartes fixes et les cartes variables dont le quotient descend le satisfont. Dans la carte radiale, $Q=\hbar_ac/L$. Avec des énergies de cette même droite, $E_1E_2/Q^2$ est adimensionnel et indépendant de la base.

Pour des projections orthogonales complémentaires,
\begin{align}
w_A&=\frac{A_\partial r_{\rm av}}
{A_\partial r_{\rm av}+V_\partial r_{\rm av}^J}
=\frac{\pi^4}{\pi^4+\varphi^5},\qquad w_V=1-w_A,
\label{hmtcuatro:gea:pesos}\\
W_{\rm av}&=w_AP_A+w_VP_V,\qquad
\overline I_v=\frac{E_1E_2}{Q^2}W_{\rm av}.
\label{hmtcuatro:gea:respuesta}
\end{align}
Multiplier le numérateur et le dénominateur de $r_{\rm av}^2/(r_{\rm av}^2+r_{\rm av}^J)$ par $\varphi^4\pi^2$ prouve l’évaluation. La distinction entre lecteur aréal et pondération est conservée : la marque régionale est générée une fois. $0<w_A,w_V<1$ rend $W_{\rm av}$ positif et inversible, et
\[
\langle z,\overline I_vz\rangle
=\frac{E_1E_2}{Q^2}
(w_A\|P_Az\|^2+w_V\|P_Vz\|^2)>0\quad(z\ne0).
\]
Échanger complètement $(A_\partial,r_{\rm av},P_A)$ et $(V_\partial,r_{\rm av}^J,P_V)$ permute les termes. Un transport unitaire conjugue la réponse et conserve son spectre ; réaliser l’échange au sein d’un espace exige des dimensions sectorielles égales. Une remise à l’échelle commune de $A_\partial,V_\partial$ ne change pas les poids et ne prouve pas de dépendance globale. À énergies et poids fixés, $Q\mapsto qQ$ modifie la réponse de $q^{-2}$ ; le témoin fort conserve en outre la même restriction locale.

\begin{theorem}[Autoinertie et rayon du même caractère]
\label{hmtcuatro:gea:autoinercial-radial} Avec la même horloge et la même orientation,
\begin{equation}
\overline I_v=\frac{G_aE_1E_2}{\hbar_ac^5}W_{\rm av}
=\frac{G_am_1m_2}{\hbar_ac}W_{\rm av}.
\label{hmtcuatro:gea:mach-gravedad}
\end{equation}
Pour une sonde $y_p>0$, $E_p=Qy_p$ et $\bar\lambda_p=L/y_p$, l’extension spectrale satisfait
\begin{equation}
\widehat I_{p,Y}=y_pY\otimes W_{\rm av}
=\frac{E_p}{Q^2}H_X\otimes W_{\rm av}
=\frac{R_X}{\bar\lambda_p}\otimes W_{\rm av}
=\frac{G_am_p}{\hbar_ac}M_X\otimes W_{\rm av}.
\label{hmtcuatro:gea:identidad-autoinercial}
\end{equation}
\end{theorem}
\begin{proof}
$Q^2=\hbar_a^2c^2/L^2$ et $\hbar_aG_a=c^3L^2$ donnent $Q^{-2}=G_a/(\hbar_ac^5)$ ; on utilise ensuite $E_i=m_ic^2$. Pour $Y=\sum_jy_jP_j$, appliquer le lecteur à chaque énergie $Qy_j$ produit $\sum_jy_py_jP_j\otimes W_{\rm av}$, indépendamment de la base propre. Le calcul spectral étend le lecteur linéaire au caractère autoadjoint positif général. La décomposition par $P_A,P_V$ donne les opérateurs $y_pw_AY$ et $y_pw_VY$, avec leurs multiplicités : ils sont autoadjoints positifs sur le domaine de $Y\otimes I$. Substituer les lecteurs radiaux prouve les autres expressions.
\end{proof}

Si une réalisation supplémentaire identifie effectivement son opérateur complet à $H_X=QY$, elle conserve cette identification et ses domaines. Pour $H_X=\sum_rH_r$ sur un cœur commun, $\widehat I_{p,Y}=(E_p/Q^2)\sum_rH_r\otimes W_{\rm av}$. La positivité de chaque terme n’est pas requise ; celle du caractère de l’opérateur identifié l’est. Aucun terme n’est supprimé et l’origine de l’énergie n’est pas modifiée pour l’imposer. Le théorème n’attribue pas automatiquement ce lecteur à $\mathbb H$.

Dans les blocs de \eqref{hmtcuatro:gea:schur-prueba}, l’inverse intérieur de $y_pY\otimes W_{\rm av}$ est $y_p^{-1}C^{-1}\otimes W_{\rm av}^{-1}$. La multiplication donne
\begin{equation}
\operatorname{Schur}(\widehat I_{p,Y})
=y_pY_{\rm ef}\otimes W_{\rm av}
=\frac{R_{X,\rm ef}}{\bar\lambda_p}\otimes W_{\rm av}.
\label{hmtcuatro:gea:schur-autoinercia}
\end{equation}
Les blocs croisés ne sont pas effacés. Pour $a_s=y_pw_s>0$, $s=A,V$, le même calcul donne $\mathcal K_{a_sY}(a_sz)=a_s\mathcal K_Y(z)$ dans le domaine résolvant intérieur, en conservant la mémoire complète. Si $Y'K=KY$ et $W'_{\rm av}J=JW_{\rm av}$, avec $y_p$ et la section intensive fixés,
\begin{equation}
\widehat I'_{p,Y'}(K\otimes J)
=(K\otimes J)\widehat I_{p,Y}.
\label{hmtcuatro:gea:naturalidad}
\end{equation}
Cela se prouve par substitution. Les coupures spectrales de la proposition \ref{hmtcuatro:gea:limite} convergent en norme de graphe ; les deux poids fixes transmettent la convergence à la réponse. La naturalité est celle de ces entrelaceurs effectifs, non celle d’une famille distincte d’applications spatiales identifiée par son seul nom.

\subsection{Autoinertie de l’état et accélération de sa projection}
\label{hmtcuatro:gea:proyeccion}

Dans une réalisation différentiable, l’état conjoint $\Gamma=(x,\mathcal O,m)$ possède la connexion $\nabla^{\rm tot}$ sur $\mathcal N$. Une projection lisse $\pi_S:\mathcal N\to\mathcal N_S$, avec la connexion $\nabla^S$, conserve les transports et les lecteurs qu’elle réalise. Le repère observé est $\mathfrak F_{\mathcal O}(\Gamma)
=\operatorname{Res}_{\mathcal O}(\Gamma,\nabla^{\rm tot})$. Pour une courbe deux fois différentiable et $\gamma_S=\pi_S\circ\Gamma$, l’autoinertie signifie $\nabla^{\rm tot}_{\dot\Gamma}\dot\Gamma=0$. Le défaut est
\[
\mathcal K_{\pi_S}(\dot\Gamma,\dot\Gamma)
=\nabla^S_{\dot\gamma_S}\dot\gamma_S
-d\pi_S(\nabla^{\rm tot}_{\dot\Gamma}\dot\Gamma).
\]
\begin{proposition}[Accélération issue de la lecture]
\label{hmtcuatro:gea:aceleracion}
Pour une trajectoire auto-inertielle,
\begin{equation}
a_S=\mathcal K_{\pi_S}(\dot\Gamma,\dot\Gamma),\qquad
(\mathcal K_{\pi_S})^a{}_{BC}
=\partial_B\partial_C\pi_S^a
+(\Gamma^S)^a{}_{bc}\partial_B\pi_S^b\partial_C\pi_S^c
-\partial_D\pi_S^a(\Gamma^{\rm tot})^D{}_{BC}.
\label{hmtcuatro:gea:defecto}
\end{equation}
\end{proposition}
\begin{proof}
La règle de la chaîne donne $\ddot\gamma_S^a=\partial_B\pi_S^a\ddot\Gamma^B
+\partial_B\partial_C\pi_S^a\dot\Gamma^B\dot\Gamma^C$. Ajouter la connexion du sous-système et soustraire l’image de l’accélération totale annule les termes en $\ddot\Gamma$. Les termes restants sont le tenseur affiché contracté avec les vitesses. L’autoinertie annule l’accélération totale et laisse l’égalité.
\end{proof}
La réponse géométrique se factorise par $b$ lorsque la connexion, la projection et la donnée cinématique descendent à la frontière, ou lorsque cette dernière est fixée dans la comparaison. Si la vitesse change, elle est conservée comme argument. L’ensemble autoinertiel peut ainsi donner en lecture des sous-systèmes accélérés sans repère absolu extérieur.

L’autoinertie n’annule pas à elle seule la courbure ni l’holonomie. Une donnée transportée le long d’une boucle revient si et seulement si elle appartient à l’espace fixe de son holonomie. Le retour de l’état complet exige en outre une réalisation injective dans le secteur comparé : retour visible et retour de l’histoire demeurent distincts.

\subsection{Résultat et contrôles de portée}
\label{hmtcuatro:gea:conclusion}

Le même caractère fixe masse, énergie et rayon ; leur proportionnalité conserve mémoire statique, mémoire spectrale et normalisation duale. Les transports géométrique et interne interviennent dans la différentielle d’une seule énergie. L’action conjointe produit les sources géométrique et torsionnelle et conserve les contraintes sur sa solution régulière. L’opérateur autoinertiel est le rayon du caractère divisé par la longueur conjuguée de la sonde, tandis que l’accélération est le défaut de projection de l’état.

Les deux derniers objets ont des types distincts : $\widehat I_{p,Y}$ ne s’identifie ni à $\mathcal K_{\pi_S}$ ni au tenseur d’Einstein. L’incorporation démontrée n’est pas non plus remplacée par l’exigence différente d’annuler une courbure quantique totale pour des déformations normales arbitraires. Cet énoncé possède ses générateurs et ses domaines ; il n’est pas affirmé ici, ni déclaré absent du corpus au motif qu’il ne constitue pas cette conclusion.

Les falsificateurs sont concrets. Changer seulement $G$ ou $\hbar$ avec $L,Y$ fixés rompt \eqref{hmtcuatro:gea:radial} ; remplacer $L$ par $\ell_0$ sans $54/\pi$ altère l’horloge ; effacer $F$ modifie la réduction ; remettre conjointement à l’échelle les lecteurs aréal et volumétrique ne prouve pas la séparation machienne ; une application qui n’entrelace pas $Y$ ne satisfait pas \eqref{hmtcuatro:gea:naturalidad}. Les essais rationnels sont des contrôles finis reproductibles, non des valeurs cibles du générateur ni des substituts aux preuves de domaine et de limite.

% Fragmento integrable; no documento autónomo y no orden de compilación.
% Fecha: 2026-09-30. No modifica fuentes publicadas.
% Procedencia: reunión de RETROACCION_PCH_Y_RESTRICCIONES_VARIACIONALES.md,
% PRECUANTIZACION_PCH_Y_FIDELIDAD.md y CIERRE_CANONICO_TOTAL_Y_REPRESENTACION.md.
% Prerrequisitos de inserción: núcleo común APP--TRIT--TPK; VIII/30d,
% VIII/31 y VIII/33; lectores de acción, gamma y acoplamientos.
% Los nuevos labels usan exclusivamente el prefijo hmtcuatro:.
\section{Action conjointe, rétroaction et clôture canonique des contraintes}
\label{hmtcuatro:accion-retroaccion-restricciones}

La séquence causale de ce développement est
\[
 \mathrm{APP}\longrightarrow\mathrm{TRIT}\longrightarrow\mathrm{TPK}
 \longrightarrow\text{état enrichi}
 \longrightarrow\text{structure discrète conjointe du continu}.
\]
Les feuillets, l’orientation, les résidus, la retenue, les routes, la frontière et la mémoire de cette construction sont conservés. L’échelle d’action et les couplages sont des sorties HMT reçues avant l’emploi du langage variationnel et symplectique. Ils ne sont pas sélectionnés par des valeurs cibles ni par la conclusion de courbure à démontrer. Le retour de phase ne réinitialise pas non plus l’état enrichi.

La réalisation géométrique précédemment construite fournit une cotétrade non dégénérée \(e\), sa connexion métrique et le relèvement spinoriel du transport. Ici, cette géométrie varie avec les champs matériels. Nous réunirons en une seule démonstration l’action commune, sa réduction torsionnelle, la transformation de Legendre, le caractère de première classe, la propagation des contraintes et la platitude caractéristique de la connexion canonique issue de la même action. Cette dernière réalisation n’est pas identifiée par son nom à un autre hamiltonien en interaction.

\subsection{Structure de départ et action à source complète}
\label{hmtcuatro:datos-accion-total}

Soit \(M\) une région lisse orientée de dimension quatre, munie d’une structure de spin, de signature \((-+++)\) et d’un feuilletage spatial régulier. La cotétrade prend ses valeurs dans la droite de longueur de la réalisation reçue. Nous écrivons
\[
 \Sigma^{IJ}=e^I\wedge e^J,\qquad
 F_\omega=\mathrm d\omega+\omega\wedge\omega,\qquad
 \mathsf P_\gamma=\star+\gamma^{-1}I,\qquad \star^2=-I.
\]
L’étoile de cette formule agit sur les bivecteurs internes ; ce n’est pas la dualité extérieure des formes. Pendant la variation, les sections reçues \(c>0,\hbar>0,G>0,\gamma\in\mathbb R\setminus\{0\}\) et \(\Lambda\) restent fixes. Avec \(\kappa=8\pi G/c^4\), l’action totale est
\begin{equation}
 \begin{split}
 S_{\rm tot}={}&\frac1{2\kappa c}
   \int_M\Sigma_{IJ}\wedge(\mathsf P_\gamma F_\omega)^{IJ}
 -\frac{\Lambda}{\kappa c}\int_M\operatorname{vol}_e\\
 &+S_{\rm YM}[e,A]+S_H[e,A,H]
   +S_W[e,\omega,A,\Psi]+S_Y[e,H,\Psi]+S_{\partial}.
 \end{split}
 \label{hmtcuatro:accion-total}
\end{equation}

\begin{definition}[Unification variationnelle et canonique]
\label{hmtcuatro:def-unificacion}
L’unification variationnelle et canonique HMT est la composition des réalisations géométrique, chromatique et électrofaible du même état enrichi dans l’action commune \eqref{hmtcuatro:accion-total}, dont les variations produisent des sources conjointes et dont la réduction conserve une même algèbre de contraintes ainsi que sa propagation. La composante électromagnétique appartient à la réalisation électrofaible. Le théorème~\ref{hmtcuatro:teorema-canonico-total} établit cette composition dans le domaine régulier et sous les conditions au bord déclarées.
\end{definition}
Les couplages électrofaibles $g,g'$ et l’autocouplage $\lambda_H$ proviennent de la construction de \ref{amp-higgs-seccion} et du théorème~\ref{amp-higgs-inversa-teorema}; la variation présente maintient ces sections fixes. Leur comparaison entre sections constitue une autre opération, conservée dans le développement de l’horizon de compatibilité.


Le symbole \(H\) sans indice désigne ici le champ de Higgs, non un hamiltonien. Toutes les contractions spatiotemporelles de la matière utilisent la même \(e\). Le problème de bord est celui reçu de la variation de Cartan--Holst : variations intérieures, données de bord fixes compatibles ou complétion annulant le même terme de frontière.

Pour fixer le contenu de \eqref{hmtcuatro:accion-total}, les connexions internes agissent sur le fibré chiral
\[
 (S_L\otimes E_L)\oplus(S_R\otimes E_R).
\]
Ce fibré n’est pas remplacé par un couplage vectoriel. Dans la normalisation entière \(y=6Y\), les multiplets reçus ont pour poids \(1,4,-2,-3,-6\), respectivement pour \(q_L,u_R,d_R,\ell_L,e_R\) ; le doublet \(H\) a pour poids \(3\). La couleur agit sur le facteur \(\mathbb C^3\) des quarks et trivialement sur les leptons. Le facteur de familles et ses applications \(Y_u,Y_d,Y_e\) sont conservés. Ce bloc n’ajoute pas de partenaire droit à un multiplet qui n’en possède pas.

Avec \(\widetilde H=i\sigma_2\overline H\), l’application matérielle est
\begin{equation}
 \begin{split}
 \mathsf Y_q(H)(u,d)&=
       \widetilde H\otimes Y_u u+H\otimes Y_d d,\\
 \mathsf Y_\ell(H)e&=H\otimes Y_e e .
 \end{split}
 \label{hmtcuatro:yukawa}
\end{equation}
L’identité sur le facteur de couleur est sous-entendue à la première ligne. Le terme d’action est \(-(\bar\Psi_L\mathsf Y(H)\Psi_R+\mathrm{h.c.})\), avec les sections dimensionnelles reçues. L’équivariance n’est pas seulement nominale : pour \(U\in SU(2)\), \(i\sigma_2\overline U=U i\sigma_2\), de sorte que \(\widetilde H\mapsto U\widetilde H\). Les poids des trois colonnes sont \(-3+4=1\), \(3-2=1\) et \(3-6=-3\). La couleur est entrelacée par l’identité et les applications de famille commutent avec l’action interne. Par conséquent, pour chaque transformation interne \(u\),
\begin{equation}
 \mathsf Y(uH)\rho_R(u)=\rho_L(u)\mathsf Y(H).
 \label{hmtcuatro:equivariancia-yukawa}
\end{equation}
Si la réalisation matérielle précédente contient en outre d’autres multiplets, leurs applications sont incorporées avec leur identité d’équivariance propre ; cette preuve n’invente pas de nouveaux états.

Dans une carte d’unités \(\hbar=c=1\), la densité matérielle qui résume ces actions a la forme
\begin{equation}
 \begin{split}
 \mathcal L_m={}&-\frac14 F_{B,\mu\nu}F_B^{\mu\nu}
 -\frac14\sum_a F_{W,\mu\nu}^aF_W^{a,\mu\nu}
 -\frac14\sum_a F_{C,\mu\nu}^aF_C^{a,\mu\nu}\\
 &-(\nabla_\mu H)^\dagger\nabla^\mu H-V(H)
 +\mathcal L_W(\Psi_L,\nabla_L)+\mathcal L_W(\Psi_R,\nabla_R)\\
 &-\bigl(\bar\Psi_L\mathsf Y(H)\Psi_R+\mathrm{h.c.}\bigr).
 \end{split}
 \label{hmtcuatro:densidad-material}
\end{equation}
Les courbures incluent leurs couplages ; ceux-ci et les normalisations des composantes sont ceux de leurs documents sources. L’écriture dans cette carte n’élimine pas la section d’action et ne recalibre pas les coefficients. En particulier, aucune nouvelle valeur numérique du couplage fort n’est introduite ici. La connexion de spin de \(\nabla_L,\nabla_R\) est celle de \(e,\omega\) ; les connexions \(B,W,C\) agissent dans les représentations qui viennent d’être décrites. La constante de \(V\) est conservée, car sa variation de volume appartient à la source gravitationnelle.

La partie spinorielle conserve les conventions explicites du développement précédent : \(\{g^I,g^J\}=2\eta^{IJ}I\), \(\mathbb B_{IJ}=\frac14[g_I,g_J]\), \(A_D=-ig^0\), \(\bar\psi=\psi^\dagger A_D\) et \(D\psi=\mathrm d\psi+\frac12\omega^{IJ}\mathbb B_{IJ}\psi+A\psi\). Sa cinétique symétrique est
\[
 \frac{\hbar}{2}
 [\bar\psi g^I D\psi-(D\bar\psi)g^I\psi]\wedge\eta_I,
 \qquad \eta_I=\iota_{E_I}\operatorname{vol}_e .
\]
La somme porte sur les champs réellement présents, avec leurs projections chirales, non sur une duplication vectorielle du catalogue. Les termes de Yukawa sont algébriques et ne dépendent pas de \(\omega\). Le courant total est fixé par
\begin{equation}
 \delta_\omega S_m=-\frac12\int_M\delta\omega^{IJ}\wedge\sigma_{IJ},
 \qquad
 \sigma_{JK}=\sum_s-\frac{\hbar}{2}
   \bar\psi_s\{g^I,\mathbb B_{JK}\}\psi_s\,\eta_I.
 \label{hmtcuatro:corriente-total}
\end{equation}

\subsection{Une démonstration de compatibilité, de propagation et de clôture}
\label{hmtcuatro:demostracion-unica}

\begin{theorem}[Compatibilité canonique de l’action conjointe]
\label{hmtcuatro:teorema-canonico-total} Considérons \eqref{hmtcuatro:accion-total} sur le domaine non dégénéré précédent, avec une matière affine en la contorsion, et dans une carte régulière de la transformation de Legendre restreinte. Conservons les contraintes fermioniques du premier ordre et éliminons celles de seconde classe au moyen de leur crochet de Dirac. Pour des paramètres à support intérieur, ou avec le bord compatible sans charges impropres, les propriétés suivantes sont conjointement satisfaites : l’équation métrique possède la source matérielle totale de cette même action ; ses générateurs de difféomorphismes et de jauge sont de première classe ; les contraintes se propagent sur toute solution couplée régulière ; et le potentiel canonique total induit une connexion préquantique plate dans les directions caractéristiques de la surface régulière des contraintes.

L’assertion porte sur cette action et cette réalisation canonique. Elle n’identifie pas automatiquement sa connexion à un hamiltonien en interaction distinct et ne prouve pas sa limite spatiale--chirale.
\end{theorem}

\begin{proof}
\par\medskip\noindent\textbf{1. Élimination torsionnelle sans duplication.}\quad
Écrivons \(\omega=\mathring\omega(e)+\mathfrak k\). Le développement \(F_\omega=F_{\mathring\omega}
+D_{\mathring\omega}\mathfrak k+\mathfrak k\wedge\mathfrak k\) et \(D_{\mathring\omega}\Sigma=0\) séparent la contribution de bord
\[
 S_{\partial,\mathfrak k}
 =\frac1{2\kappa c}\int_{\partial M}
                  \Sigma_{IJ}\wedge(\mathsf P_\gamma\mathfrak k)^{IJ}
\]
de la densité algébrique
\begin{equation}
 \mathcal Q_{e,\gamma}(\mathfrak k)
       -\frac12\mathfrak k^{IJ}\wedge\sigma_{IJ},\qquad
 \mathcal Q_{e,\gamma}(\mathfrak k)=\frac1{2\kappa c}
       \Sigma_{IJ}\wedge
       \mathsf P_\gamma(\mathfrak k\wedge\mathfrak k)^{IJ}.
 \label{hmtcuatro:cuadratica-contorsion}
\end{equation}
La variation originale donne \(D_\omega(\mathsf P_\gamma\Sigma)=\kappa c\,\sigma\). L’inverse de Holst et l’inverse torsion--contorsion du développement précédent sont inversibles pour le corepère et \(\gamma\) déclarés. Ils déterminent ainsi une unique \(\mathfrak k_*[e,\sigma]\), linéaire en \(\sigma\). La variation radiale de \eqref{hmtcuatro:cuadratica-contorsion} en son point stationnaire donne \(2\mathcal Q(\mathfrak k_*)=\frac12\mathfrak k_*^{IJ}\wedge\sigma_{IJ}\). Par conséquent,
\begin{equation}
 S_{\rm red}=\frac1{2\kappa c}\int_M(R-2\Lambda)\operatorname{vol}_e
       +S_m^{\rm LC}
       -\frac14\int_M\mathfrak k_*^{IJ}\wedge\sigma_{IJ}
       +S_{\partial,\rm red}.
 \label{hmtcuatro:accion-reducida}
\end{equation}
La première identité de Bianchi annule le terme de Holst de Levi--Civita dans cette écriture. La règle de la chaîne conserve les équations de \(e,A,H,\Psi\), car le coefficient de \(\delta\mathfrak k_*\) s’annule par stationnarité. Comme \(\sigma=\sum_s\sigma_s\), la contribution effective inclut \(\mathfrak k_*[\sigma_s]\wedge\sigma_t\) pour \(s\ne t\). Elle n’est pas remplacée par des carrés d’espèces isolées, et la même \(\mathfrak k\) n’est pas conservée simultanément comme variable indépendante.

\par\medskip\noindent\textbf{2. Source complète et rétroaction.}\quad
Définissons le tenseur d’action par la variation de tous les termes matériels de \eqref{hmtcuatro:accion-reducida} :
\begin{equation}
 \delta_g S_{m,\rm red}=-\frac12\int_M
       \mathcal T^S_{\mu\nu}\,\delta g^{\mu\nu}\operatorname{vol}_e,
 \qquad
 E_{\mu\nu}:=G_{\mu\nu}+\Lambda g_{\mu\nu}
                   -\kappa c\,\mathcal T^S_{\mu\nu}.
 \label{hmtcuatro:fuente-completa}
\end{equation}
L’équation métrique est \(E_{\mu\nu}=0\). Le facteur \(\kappa c\) correspond à la cotétrade de longueur ; si \(T^{\rm phys}=c\mathcal T^S\), on écrit \(\kappa T^{\rm phys}\). On fait également varier les champs matériels et les connexions internes. Les champs déterminent donc la source, et la géométrie qui résout l’équation détermine à son tour leurs connexions, leur volume et leur évolution. Il s’agit d’une rétroaction, non d’une évolution sur un fond fixé. La variation inclut le volume du potentiel de Higgs et la dépendance métrique du courant torsionnel, ainsi que ses termes croisés.

\par\medskip\noindent\textbf{3. Transformation de Legendre restreinte et réduction canonique.}\quad
Le potentiel canonique total s’obtient à partir de \(\delta L=\mathcal E\delta z+\mathrm d\theta(\delta z)\), en intégrant \(\theta\) sur la feuille spatiale et en effectuant la réduction de seconde classe. Sa partie géométrique avant cette réduction est
\[
 \theta_{\rm geom}(\delta)=
   \frac1{2\kappa c}(\mathsf P_\gamma\Sigma)_{IJ}
                       \wedge\delta\omega^{IJ}.
\]
La décomposition \(F_{0a}=\dot\omega_a-D_a\omega_0\), avec la même décomposition des actions matérielles, produit
\begin{equation}
 S_{\rm can}=\int d\tau\,
 \bigl[\Theta_{\rm red}(\dot z)-\mathcal C[N]-\mathcal D[\xi]
       -\mathcal G_{\rm L}[\lambda]-\mathcal G_{\rm int}[\chi]\bigr]
       +S_{\partial,\rm red}.
 \label{hmtcuatro:legendre-total}
\end{equation}
Les vitesses absentes du lapse, du shift ou des connexions temporelles ne sont pas artificiellement inversées.

La contorsion algébrique possède les contraintes \(\pi_{\mathfrak k}=0\) et \(\partial\mathcal Q/\partial\mathfrak k-\sigma/2=0\). Leur matrice de crochets possède le bloc
\[
 \begin{pmatrix}0&-M\\ M^T&B\end{pmatrix},
 \qquad M=\frac{\partial^2\mathcal Q}{\partial\mathfrak k^2}.
\]
L’inverse torsionnel rend \(M\) inversible, et donc cette matrice. Ce sont des contraintes de seconde classe. On utilise
\begin{equation}
 \{F,G\}_D=\{F,G\}
       -\{F,\chi_a\}(C^{-1})^{ab}\{\chi_b,G\}.
 \label{hmtcuatro:corchete-dirac}
\end{equation}
Pour les variables qui ne dépendent pas de la paire contorsion--moment, le bloc inférieur droit de \(C^{-1}\) est nul. Cette élimination ne fait apparaître aucun crochet supplémentaire entre elles. Les contraintes spinorielles du premier ordre sont conservées séparément ; dans l’espace de phase gradué, on emploie le crochet gradué et les générateurs pairs. \(\Theta_{\rm red}\) désigne le résultat complet, non sa seule partie géométrique.

L’inversion de Legendre dans les directions régulières restitue la même action. Par exemple, \(N[p^2/(2w)+wV+H_{\rm grad}]\), avec \(w=\sqrt{\det q}\), donne \(p=(w/N)D_\tau\phi\), et par substitution
\[
 L_H=\frac{w}{2N}|D_\tau\phi|^2-NwV-NH_{\rm grad}.
\]
\(D_\tau\) conserve le shift et la connexion temporelle ; \(w\) n’est pas figé. Pour \(K_{ij}=(\dot q_{ij}-\mathcal L_\xi q_{ij})/(2N)\), la contribution gravitationnelle au moment et son inverse sont
\begin{equation}
 \Pi^{ij}=\frac{\sqrt q}{2\kappa c}(K^{ij}-q^{ij}K),\qquad
 K_{ij}=\frac{2\kappa c}{\sqrt q}
                  (\Pi_{ij}-\tfrac12q_{ij}\Pi).
 \label{hmtcuatro:legendre-geometrica}
\end{equation}
Il en résulte
\[
 \mathcal C_{\rm grav}
 =\frac{2\kappa c}{\sqrt q}
       (\Pi^{ij}\Pi_{ij}-\tfrac12\Pi^2)
  -\frac{\sqrt q}{2\kappa c}(R^{(3)}-2\Lambda).
\]
Si la carte des spineurs apporte un décalage matériel du moment total \(p\), on utilise \(\Pi=p-p_m\) ; \(p_m\) n’est pas écarté. Le caractère indéfini de la forme de DeWitt n’empêche pas cet inverse algébrique. On n’exige pas d’inverse des vitesses pour la cinétique fermionique du premier ordre.

\par\medskip\noindent\textbf{4. Identité de Noether--Legendre et caractère de première classe.}\quad
La coïncidence des générateurs avec les contraintes de la même action peut être constatée avant d’imposer ses équations. Posons \(a=(2\kappa c)^{-1}\), \(B=\mathsf P_\gamma\Sigma\), \(u^I=\iota_\zeta e^I\) et \(\lambda^{IJ}=\iota_\zeta\omega^{IJ}\). De \(\mathcal L_\zeta\omega=\iota_\zeta F+D_\omega\lambda\), on obtient
\begin{equation}
 \begin{split}
 j_\zeta^{\rm geom}
 ={}&\mathrm d(a B_{IJ}\lambda^{IJ})
   -a(D_\omega B_{IJ})\lambda^{IJ}\\
 &-2a u^Ie^J\wedge(\mathsf P_\gamma F)_{IJ}
       +2a\Lambda u^I\eta_I .
 \end{split}
 \label{hmtcuatro:noether-legendre}
\end{equation}
La partie matérielle ajoute exactement ses contributions aux équations de \(e,\omega,A\), ses générateurs verticaux et son bord. Décomposer la même expression sur une feuille spatiale donne \eqref{hmtcuatro:legendre-total}. La réduction stationnaire et la réduction de Dirac transportent ces transformations au lieu de les remplacer par les moments d’une action extérieure.

Avec la convention canonique usuelle pour les contraintes, les crochets totaux sont, modulo les générateurs de Gauss conservés,
\begin{equation}
 \begin{aligned}
 \{\mathcal D[\xi],\mathcal D[\eta]\}_D
    &\simeq\mathcal D[[\xi,\eta]],\\
 \{\mathcal D[\xi],\mathcal C[N]\}_D
    &\simeq\mathcal C[\mathcal L_\xi N],\\
 \{\mathcal C[N],\mathcal C[M]\}_D
    &\simeq\mathcal D[
       q^{ij}(N\partial_jM-M\partial_jN)\partial_i].
 \end{aligned}
 \label{hmtcuatro:algebra-restricciones}
\end{equation}
En effet, les transformations tangentielles ont le crochet de Lie et transportent le lapse comme un scalaire. Pour les transformations normales, différencier leur orthogonalité aux vecteurs tangents et \(g(n,n)=-1\) détermine la variation tangentielle de \(n\) ; avec \(K_{ij}\) de \eqref{hmtcuatro:legendre-geometrica}, on obtient \(\delta_N n=\operatorname{grad}_qN\). Le commutateur des déformations, converti en crochet canonique avec ces conventions, donne la dernière ligne de \eqref{hmtcuatro:algebra-restricciones}. L’usage du transport covariantisé conserve en outre les rotations de repère et de jauge. Les transformations sont tangentes à la surface des contraintes de l’action invariante, d’où le caractère de première classe. Les fonctions \(q^{ij}\) dépendent de l’état et ne sont pas figées. Avec un support intérieur, aucune charge centrale de bord n’est ajoutée ; un générateur impropre de charge non nulle n’est pas déclaré contrainte nulle.

\par\medskip\noindent\textbf{5. Propagation dans la géométrie qui répond à la matière.}\quad
L’invariance par difféomorphismes de \(S_{m,\rm red}\), évaluée sur ses équations matérielles, donne par intégration par parties \(\mathring\nabla^\mu\mathcal T^S_{\mu\nu}=0\). L’identité utilise le tenseur complet de \eqref{hmtcuatro:fuente-completa} ; elle ne postule pas la conservation séparée de chaque contribution. Bianchi implique alors \(\mathring\nabla_\mu E^{\mu\nu}=0\), avant même d’imposer toutes les équations métriques.

Dans une carte gaussienne locale, \(ds^2=-d\tau^2+q_{ij}dx^idx^j\), imposons les équations d’évolution \(E_{ij}=0\) et posons \(C=E_{00}\), \(C_i=E_{0i}\), \(K_{ij}=\dot q_{ij}/2\), \(K=q^{ij}K_{ij}\). Ici, \(K\) est seulement la trace de la seconde forme fondamentale. Les deux composantes de l’identité de divergence sont
\begin{equation}
 \partial_\tau C-D_iC^i+KC=0,\qquad
 \partial_\tau C_i+KC_i=0.
 \label{hmtcuatro:propagacion}
\end{equation}
Pour vérifier les signes, on utilise \(E^{00}=C\), \(E^{0i}=-C^i\), \(E^{ij}=0\), \(\Gamma^i{}_{0j}=K^i{}_j\) et \(\Gamma^\mu{}_{\mu0}=K\). En abaissant l’indice de l’équation spatiale, \(\dot q_{ij}=2K_{ij}\) annule les deux termes de connexion. Si initialement \(C_i=0\), son équation homogène maintient \(C_i=0\) ; la première devient \(\dot C+KC=0\) et conserve \(C=0\). \(q\) et \(K\) sont ceux de la solution couplée. Le choix gaussien prouve l’assertion locale ; sa forme tensorielle permet de changer le lapse et le shift. De même, l’identité de jauge de Noether relie la divergence covariante de l’équation de jauge aux équations matérielles et propage Gauss lorsque les équations spatiales sont satisfaites. On n’en infère ni existence globale ni absence de singularités pour toute solution.

\par\medskip\noindent\textbf{6. La même action induit la clôture caractéristique.}\quad
Soit \(\mathcal P_{\rm red}\) la carte canonique régulière obtenue et \(\mathcal Z=\{C_A=0\}\) sa surface régulière de contraintes totales, Gauss compris. Nous utilisons pour cette surface un symbole distinct de \(\Sigma^{IJ}\). Le potentiel \(\Theta=\Theta_{\rm red}\) de \eqref{hmtcuatro:legendre-total} est conservé, termes matériels et mixtes compris. Nous fixons
\begin{equation}
 \Omega=-\mathrm d_{\mathcal P}\Theta,\qquad
 \iota_{X_f}\Omega=\mathrm d_{\mathcal P}f,\qquad
 \{f,g\}_D=\Omega(X_f,X_g).
 \label{hmtcuatro:convenciones-simpecticas}
\end{equation}
Dans une paire canonique, \(\Theta=p\,\mathrm dq\) donne \(X_f=f_p\partial_q-f_q\partial_p\), \(\{q,p\}_D=1\), \(X_f(g)=-\{f,g\}_D\) et \([X_f,X_g]=-X_{\{f,g\}_D}\). Dans l’espace de phase gradué, cette écriture s’applique aux fonctions et contraintes paires pertinentes.

Sur la droite de phase trivialisée au-dessus de cette carte, nous construisons
\begin{equation}
 \nabla_X=X-\frac{i}{\hbar}\Theta(X)
       =\delta_X+\frac{i}{\hbar}H_X^{\rm can},
 \qquad H_X^{\rm can}=-\Theta(X).
 \label{hmtcuatro:conexion-canonica}
\end{equation}
On ne choisit pas de conjugaison arbitraire pour la rendre plate : son coefficient provient du potentiel de l’action. Pour tous champs \(X,Y\) de la carte, le calcul donne
\begin{equation}
 \begin{aligned}
 \mathcal F^{\rm can}_{X,Y}
 &=XH_Y^{\rm can}-YH_X^{\rm can}-H_{[X,Y]}^{\rm can}
       +\frac{i}{\hbar}[H_X^{\rm can},H_Y^{\rm can}]\\
 &=-X\Theta(Y)+Y\Theta(X)+\Theta([X,Y])\\
 &=-\mathrm d_{\mathcal P}\Theta(X,Y)=\Omega(X,Y).
 \end{aligned}
 \label{hmtcuatro:curvatura-calculada}
\end{equation}
Les \(H_X^{\rm can}\) sont des multiplicateurs de la droite, de sorte que leur commutateur est nul ; les dérivées des coefficients et des champs ne sont pas supprimées.

Le caractère de première classe déjà prouvé signifie \(\{C_A,C_B\}_D=f_{AB}{}^D C_D\). Pour \(V\) tangent à \(\mathcal Z\), \(\Omega(X_{C_A},V)=\mathrm d_{\mathcal P}C_A(V)=0\). Les champs des contraintes sont donc caractéristiques. Sous la régularité déclarée, ils engendrent la distribution caractéristique de la surface coïsotrope. La tensorialité de \(\Omega\) étend le calcul à toutes leurs combinaisons lisses et prouve
\begin{equation}
 \left.\mathcal F^{\rm can}_{X,Y}\right|_{\mathcal Z}=0
 \quad\hbox{pour }X,Y\hbox{ caractéristiques}.
 \label{hmtcuatro:cierre-caracteristico}
\end{equation}
Pour deux lapses, \([X,Y]\) inclut la déformation tangentielle et les actions de Gauss de \eqref{hmtcuatro:algebra-restricciones}. Si l’on passe au quotient de jauge, l’égalité s’exprime modulo ces actions verticales. On n’affirme pas que \(\Omega\) soit nul sur tout l’espace de phase ni que toute holonomie globale d’une feuille soit triviale.

\par\medskip\noindent\textbf{7. Expression préquantique de la même clôture.}\quad
La réalisation différentielle associée à \eqref{hmtcuatro:conexion-canonica}, sur des sections lisses où les produits sont définis, est
\begin{equation}
 Q_{\rm pre}(f)=-i\hbar\nabla_{X_f}+f
            =-i\hbar X_f+f-\Theta(X_f).
 \label{hmtcuatro:operador-precuantico}
\end{equation}
Avec \(R^\nabla=(i/\hbar)\Omega\), le développement du commutateur et \(X_f(g)=-\{f,g\}_D\) donnent
\begin{equation}
 [Q_{\rm pre}(f),Q_{\rm pre}(g)]
           =i\hbar Q_{\rm pre}(\{f,g\}_D).
 \label{hmtcuatro:representacion-precuantica}
\end{equation}
La règle du produit qui conserve les fonctions de structure est
\begin{equation}
 Q_{\rm pre}(aC)
   =aQ_{\rm pre}(C)+C\bigl(Q_{\rm pre}(a)-a\bigr).
 \label{hmtcuatro:funciones-estructura}
\end{equation}
Elle découle de \(X_{aC}=aX_C+CX_a\) et de \eqref{hmtcuatro:operador-precuantico}. Le second terme n’est pas éliminé hors de \(\mathcal Z\). Sur \(\mathcal Z\), \(Q_{\rm pre}(C_A)=-i\hbar\nabla_{X_{C_A}}\) ; ainsi, \eqref{hmtcuatro:representacion-precuantica} et \eqref{hmtcuatro:cierre-caracteristico} sont les expressions opératorielle et géométrique de la même clôture canonique, non deux incorporations indépendantes de la gravité.
\end{proof}

\subsection{Domaine du résultat et comparaison des réalisations}
\label{hmtcuatro:alcance-canonico}

La preuve réunit une action à gravité dynamique et source conjointe, sa transformation de Legendre et sa connexion caractéristique. Le qualificatif \emph{arbitraire} pour les déformations se rapporte aux lapses et déplacements admissibles qui engendrent des directions caractéristiques dans cette région régulière ; non à tout vecteur de phase, tout bord ou toute solution singulière. La platitude locale conserve l’éventuelle monodromie globale et n’identifie pas deux histoires TPK du seul fait qu’elles ont les mêmes extrémités.

Le coefficient \(H_X^{\rm can}\) d’une connexion de droite n’est pas par définition l’opérateur d’horloge et de réponse \(H_{\rm clk}+D^*WD\), ni son extension en interaction. \(Q_{\rm pre}(C)\) inclut en outre une dérivation sur l’espace de phase ; ce n’est pas non plus un autre nom de ces opérateurs. Pour identifier une réalisation concrète au moyen d’une application \(I\), l’égalité typée est requise
\[
 \nabla_N^{\rm op}I=I\nabla_N^{\rm can}
\]
sur leurs domaines effectifs. Si \(I\) dépend de la géométrie, sa dérivée fait partie de cette égalité.

La distinction dispose d’un contrôle exact : dans la carte \(\Theta=p\,\mathrm dq\),
\begin{equation}
 Q_{\rm pre}(p^2/2)=-i\hbar p\,\partial_q-p^2/2.
 \label{hmtcuatro:control-polarizacion}
\end{equation}
En agissant sur la fonction \(1\), il produit \(-p^2/2\). Il ne préserve donc pas par simple restriction la polarisation des fonctions indépendantes de \(p\) et n’y est pas \(-\hbar^2\partial_q^2/2\). Le contrôle ne réfute pas l’hamiltonien reçu : il empêche de remplacer une représentation par une autre sans son application.

Aucune mesure de Liouville de dimension infinie n’a été supposée et l’autoadjonction ne se déduit pas des identités différentielles précédentes. Si les inclusions de raffinement entrelacent les connexions et conservent un cœur commun, leurs courbures s’entrelacent par composition ; une limite exige en outre ses propres contrôles de domaine et de convergence. Cette condition n’est pas présentée ici comme vérifiée pour l’opérateur en interaction et la limite spatiale--chirale conjointe. Les constructions quantiques antérieures sont conservées ; le résultat présent précise exactement la clôture canonique qui vient d’être démontrée.

% Exposición académica de las dos narrativas conservadas en el cuaderno.
% Procedencia y correspondencia de contenidos: gestion/mapa_narrativa.md.
% H_lazo y rho_lazo no designan la holonomía nonádica Gamma_9.
\section{Compatibilité structurelle de l’action, de la masse et de l’information opérationnelle}
\label{sec:narrativa-estructural}

La composition des équations permet d’identifier une relation plus précise
que l’équivalence entre masse et énergie ou la proportionnalité entre énergie
et fréquence : une même structure de route détermine une fréquence propre,
une relation entre longueurs et une pondération thermique. Le couple angulaire
participe à ces trois lectures parce qu’il contient le coefficient d’action et
détermine les canaux de transport. L’exposé qui suit conserve
conjointement la construction des objets, leurs opérations de lecture et
les conditions de leurs réalisations physiques.

Cette composition permet, en outre, de caractériser l’empreinte laissée par l’ordre
d’une séquence d’opérations. Cette empreinte admet une expression au moyen
du couple angulaire et une application inverse qui restitue le coefficient
d’action. Dans le domaine constitutif, la réponse du vide et le moment
spectral de Catalan sont réunis au moyen d’une même collection fonctionnelle. Le
contenu explicatif de ces relations réside dans les restrictions qu’elles
imposent conjointement, et non dans la simple accumulation d’identités.

La généalogie APP--TRIT--TPK fournit les objets qui sont composés. Les
lectures ultérieures conservent des aspects distincts d’une même structure :
l’échelle, l’orientation, la phase, la mémoire, l’incidence et la réponse. Leur analyse
conjointe établit quelles relations demeurent déterminées lors du changement de
représentation, quelle information une projection élimine et quelles observations
complémentaires permettent de la restituer. Les démonstrations des compositions
et de leurs domaines sont développées dans les sections démonstratives.

\subsection{L’action dans la composition angulaire et dans la structure électronique}
\label{subsec:narrativa-accion-angular}

La section de Planck intervient avant l’évaluation énergétique finale.
Son coefficient participe à la géométrie angulaire utilisée ultérieurement
par la loi des masses. Le lien, avec les unités conservées, est
\begin{equation}
 \hbar_{\mathrm{ret}}
   =\eta_{\mathrm{ret}}\,10^{-34}\mathcal U_S,
 \qquad
 A_{\deg}=1000\alpha,
 \qquad
 C^*_{\deg}=2(\eta_{\mathrm{ret}}+\alpha).
 \label{eq:narrativa-par-angular}
\end{equation}
Le coefficient $\eta_{\mathrm{ret}}$ est adimensionnel et peut donc
être composé avec $\alpha$ dans l’expression angulaire. La constante dimensionnelle
$\hbar_{\mathrm{ret}}$ conserve son unité d’action. Cette distinction permet de
suivre une même dépendance sans additionner des grandeurs de dimensions différentes.

L’action remplit ainsi deux fonctions liées. D’une part, elle fixe une
échelle dimensionnelle ; d’autre part, elle participe à la structure qui détermine la
répartition du transport entre les deux orientations. Après la conversion
unique des degrés en radians, l’angle et le contre-angle produisent les canaux
$q_+$ et $q_-$. Ces canaux interviennent dans les incidences hexadiques--octadiques,
dans la fonctionnelle de Barbero--Immirzi, dans la réponse constitutive et dans les
termes de la loi des masses.

La substitution de son expression structurelle à une constante rend visibles
des dépendances qui demeurent implicites lorsque l’équation est écrite au moyen
de constantes isolées. L’expression développée transmet l’opération qui
détermine la valeur, les coordonnées dont elle dépend et les conditions
dans lesquelles elle peut être réutilisée. Sa fonction explicative procède de cette
information transmise à l’étape suivante.

Dans l’électron, l’expression complète conserve l’exposant
$\varphi/\pi^2$, les corrections construites avec $\alpha$ et $\Delta_4$,
et le facteur de retour
\begin{equation}
 \left(\frac{H_5}{\eta_{\mathrm{ret}}}\right)^{80-54\alpha+6\Delta_4}.
 \label{eq:narrativa-factor-electronico}
\end{equation}
Le même $\eta_{\mathrm{ret}}$ qui intervient dans le contre-angle participe
à la lecture électronique complète. La géométrie angulaire et la correction
d’action de l’électron partagent donc une dépendance effective.
Étudier conjointement les deux équations permet de déterminer quelles relations
sont imposées par l’intervention d’une même coordonnée structurelle
dans des opérations distinctes.

\subsection{Fréquence propre et normalisation structurelle de la masse}
\label{subsec:narrativa-masa-frecuencia}

La loi générale des masses conserve son préfacteur électronique. On désigne par
$\Xi_\Gamma$ le facteur structurel complet d’une route $\Gamma$ : il réunit la
normalisation électronique, le caractère de la route avec sa signature et ses tours,
et la puissance correspondante du rapport d’action
$H_5/\eta_{\mathrm{ret}}$. Cette notation abrège une composition
d’opérations déjà définies ; son contenu reste lié à ces opérations.

L’horloge commune a pour fréquence angulaire
$\omega_0=2\pi/(108t_0)$, avec $t_0>0$. La composition de la loi des masses avec la section
d’action donne
\begin{equation}
 \boxed{
 m_\Gamma=\frac{\hbar_{\mathrm{ret}}\omega_0}{c^2}\,\Xi_\Gamma,
 \qquad
 E_\Gamma=\hbar_{\mathrm{ret}}\omega_0\Xi_\Gamma,
 \qquad
 \Omega_\Gamma=\frac{E_\Gamma}{\hbar_{\mathrm{ret}}}
             =\omega_0\Xi_\Gamma.}
 \label{eq:narrativa-masa-energia-frecuencia}
\end{equation}
Ici, $E_\Gamma$ est l’énergie au repos et $\Omega_\Gamma$ sa fréquence
énergétique associée.

Cette lecture sépare le rythme commun de la contribution propre à chaque
structure. Deux sections peuvent partager la période élémentaire et présenter
des masses et des fréquences propres différentes : leurs signatures, retours et mémoires
produisent des facteurs $\Xi_\Gamma$ différents. La diversité des lectures
spécifiques se trouve rapportée à la structure de la section, tandis que l’horloge
commune établit leur normalisation temporelle.

Dans le quotient $E_\Gamma/\hbar_{\mathrm{ret}}$, le facteur
absolu d’action qui apparaît explicitement dans
\eqref{eq:narrativa-masa-energia-frecuencia} se simplifie. Toutefois, sa
participation structurelle aux canaux et au rapport de retour contenu
dans $\Xi_\Gamma$ demeure. La simplification d’un facteur dans une lecture ultérieure
doit être distinguée de l’élimination des relations au moyen desquelles
ce facteur a été construit et est intervenu de nouveau.

Les rapports entre masses coïncident avec les rapports entre ces fréquences
propres. Une modification commune des unités d’action, de temps ou de vitesse
laisse ces rapports inchangés. La structure relative a donc
une exigence de confrontation propre : ses relations entre sections doivent
rester cohérentes, indépendamment de la normalisation commune choisie.

La fréquence associée à une phase relevée conserve, en outre, le nombre de
tours de l’histoire. Deux phases égales après leur réduction modulo
$2\pi$ peuvent correspondre à des relèvements différents. Le retour de phase
conserve une régularité, tandis que l’histoire distingue le parcours qui
a produit ce retour. L’identification de l’état complet à la phase
réduite éliminerait précisément cette distinction.

\subsection{Dilatation réciproque des lectures spatiales de la masse}
\label{subsec:narrativa-reciprocidad-espacial}

Dans la réalisation gravitationnelle circulaire, l’action et la gravitation
satisfont
\begin{equation}
 G=\frac{c^3L_\alpha^2}{\hbar_{\mathrm{ret}}},
 \qquad R_{108}=L_\alpha=\frac{c}{\omega_0}.
 \label{eq:narrativa-gravedad-accion}
\end{equation}
La composition avec la masse de
\eqref{eq:narrativa-masa-energia-frecuencia} détermine
\begin{equation}
 \boxed{
 r_{g,\Gamma}=R_{108}\Xi_\Gamma,
 \qquad
 \bar\lambda_\Gamma=\frac{R_{108}}{\Xi_\Gamma},
 \qquad
 r_{g,\Gamma}\bar\lambda_\Gamma=R_{108}^{2}.}
 \label{eq:narrativa-longitudes-reciprocas}
\end{equation}
On utilise la convention du rayon gravitationnel réduit du développement ;
la construction de $L_\alpha$ et la normalisation conjointe sont
démontrées dans la section~\ref{md:rigidez}. L’horloge et le rayon sont
deux expressions compatibles de cette même longueur.
$\bar\lambda_\Gamma$ désigne la longueur de Compton réduite. Ces
identités réciproques sont formulées pour $\Xi_\Gamma>0$ ; la lecture
$R_{108}/\Xi_\Gamma$ requiert ce domaine massique positif.

Le même facteur qui augmente la fréquence propre dilate la lecture
gravitationnelle et contracte la longueur associée à l’action et à la masse. Les deux
transformations conservent une aire commune. Dans cette réalisation, le facteur
massique agit comme paramètre de dilatation réciproque : une longueur croît avec
lui et l’autre diminue inversement.

Cette relation fournit un contenu mathématique précis au lien
entre masse et configuration spatiale. La hiérarchie de la composition est
essentielle : on détermine d’abord la route et son facteur structurel ; on obtient
ensuite ses lectures temporelles, énergétiques et spatiales. La structure
est transmise aux lectures au moyen d’opérations déterminées, au lieu
d’être reconstruite rétrospectivement à partir d’une masse mesurée.

La condition circulaire appartient au résultat
\eqref{eq:narrativa-longitudes-reciprocas}. Les deux longueurs sont des lectures
définies par cette réalisation ; leur interprétation comme rayons matériels
d’une particule requerrait l’identification physique correspondante.
L’égalité établit ici une réciprocité exacte entre les grandeurs
définies, avec le système de coordonnées d’action, l’horloge et la vitesse conservés.

L’incorporation de la quantité de mouvement maintient cette masse structurelle dans l’équation
de propagation. Dans la réalisation libre correspondante, la fréquence
propre détermine le terme massique de la dispersion, tandis que la quantité de mouvement
décrit la variation spatiale. Ainsi sont différenciés la structure
persistante de la section et son état de mouvement.

\subsection{Orientation des canaux et réponse constitutive du vide}
\label{subsec:narrativa-vacio-orientacion}

La composition de la permittivité, de la perméabilité et de la fonctionnelle de Barbero--Immirzi
permet de distinguer la valeur d’une réponse de la structure qui la
détermine. Le lecteur constitutif utilise deux canaux étiquetés, de
réponses $r_+$ et $r_-$. La vitesse normalisée dépend de leur produit ;
l’impédance dépend de leur quotient. La vitesse conserve une relation
commune entre les canaux, tandis que l’impédance distingue la répartition de la
réponse entre leurs orientations.

Au sein de la collection normalisée, la vitesse et la fonctionnelle de
Barbero--Immirzi permettent de retrouver conjointement les deux réponses
étiquetées. La fonctionnelle apporte une information orientée que la vitesse
isolée ne détermine pas. La restitution s’appuie sur la composition complète
de leurs lecteurs et conserve la distinction entre chaque canal et leur agrégat.

Soient $x$ et $y$ les coordonnées en radians de l’angle et du contre-angle.
Pour $x>0$ et $|y|<x$, la réponse satisfait
\begin{equation}
 \widehat c(x,y)\geq\widehat c(x,0),
 \label{eq:narrativa-convexidad-velocidad}
\end{equation}
avec égalité uniquement lorsque $y=0$. Plus précisément, en maintenant $x$
fixe, $\log\widehat c$ est une fonction paire et strictement convexe de $y$.
La démonstration de cette propriété appartient au développement constitutif
ultérieur.

La vitesse perd le signe de l’orientation, mais conserve un effet
de la séparation entre les canaux. Près de la configuration symétrique,
cet effet apparaît au second ordre ; l’impédance distingue l’orientation
dès le premier ordre. Par conséquent, remplacer les deux canaux par leur moyenne avant
d’évaluer la réponse change le résultat. Une lecture scalaire peut masquer
une orientation et conserver, en même temps, une conséquence de celle-ci.

Ce comportement spécifie quelle information chaque opération transmet.
La perte du signe, la conservation d’une séparation et la récupération
de la répartition constituent des propriétés différentes. La structure maintient
des conséquences observables même lorsque l’une de ses lectures est exprimée
au moyen d’un seul nombre.

L’inégalité \eqref{eq:narrativa-convexidad-velocidad} décrit la collection
mathématique de réponse dans son domaine. Une évolution physique doit se situer
dans les états et les trajectoires effectivement produits par TPK. La
paramétrisation analytique permet d’étudier la sensibilité du lecteur ;
l’identification d’une variation paramétrique à une trajectoire physique
conserve son obligation de réalisation.

\subsection{Boltzmann, fréquence structurelle et pondération thermique ternaire}
\label{subsec:narrativa-boltzmann}

La section thermique relie l’énergie du cycle à la constante de
Boltzmann et à sa température correspondante :
\begin{equation}
 k_B\Theta_{108}\log3=E_{\mathrm{clk}}.
 \label{eq:narrativa-boltzmann-ciclo}
\end{equation}
Comme l’énergie d’une section est
$E_\Gamma=E_{\mathrm{clk}}\Xi_\Gamma$, le poids thermique évalué à cette
température est
\begin{equation}
 \boxed{\exp\!\left(-\frac{E_\Gamma}{k_B\Theta_{108}}\right)
       =3^{-\Xi_\Gamma}.}
 \label{eq:narrativa-peso-termico}
\end{equation}
La base ternaire apparaît par composition de la normalisation informationnelle
avec l’action et l’horloge. La structure de route qui détermine la masse
détermine également l’exposant du poids thermique.

Le poids et la probabilité remplissent des fonctions distinctes. Les probabilités
s’obtiennent en normalisant les poids et en conservant les multiplicités des
états. De même, plusieurs histoires équivalentes requièrent une identification
d’états avant de recevoir une multiplicité physique : le nombre d’histoires
et la dégénérescence de l’état correspondent à des objets différents.

La composition permet de réunir les lectures dans une identité commune :
\begin{equation}
 \boxed{
 \Xi_\Gamma
 =\frac{\Omega_\Gamma}{\omega_0}
 =\frac{r_{g,\Gamma}}{R_{108}}
 =\frac{R_{108}}{\bar\lambda_\Gamma}
 =\frac{E_\Gamma}{k_B\Theta_{108}\log3}.}
 \label{eq:narrativa-cuatro-lecturas}
\end{equation}
On obtient quatre manières coordonnées de lire une même composante
structurelle : temporelle, spatiale, gravitationnelle et thermique. Chaque égalité
conserve les conditions de sa réalisation, en particulier la condition
circulaire gravitationnelle et la température du cycle.

L’identité \eqref{eq:narrativa-cuatro-lecturas} spécifie la relation
que partagent les grandeurs et fournit un critère pour confronter leur
compatibilité. La masse relative peut être caractérisée au moyen de la fréquence
relative ; la même coordonnée détermine une dilatation réciproque des
longueurs et une pondération thermique. Les fonctions de lecture sont distinctes,
mais leur dépendance se trouve liée par la généalogie commune.

La relation thermique identifie l’état et la normalisation dans lesquels
intervient $\Theta_{108}$. Cette identification maintient séparées l’énergie
de la section et son éventuel transfert sous forme de chaleur ; elle maintient également
localisée la signification de la température du cycle, sans l’étendre à
l’univers comme température unique. La capacité explicative procède des
conditions explicites de la composition.

\subsection{Ellipse d’action et transport de la forme géométrique}
\label{subsec:narrativa-elipse-transporte}

L’ellipse d’action a des demi-axes dont le produit est $2\hbar$. L’aire
orientée accumulée pendant un tour est $h=2\pi\hbar$. Une déformation
elliptique redistribue les demi-axes en conservant leur produit. Dans la
réalisation électrique et magnétique, on obtient deux contributions
proportionnelles à $\cos^2\theta$ et $\sin^2\theta$, dont la somme est
$\hbar\omega$. La structure détermine donc la répartition de
l’énergie pendant le cycle, en plus de sa valeur totale.

Le transport d’une ellipse dont la forme change exige de conserver la relation
entre la déformation et l’action. Soit $\zeta$ la coordonnée de déformation,
déterminée par le rapport entre le contre-angle et l’angle. Pour une
déformation différentiable, le transport exact introduit
\begin{equation}
 H_{\mathrm{tot}}
 =\frac{\omega}{2}
   \left(e^{-\zeta}Q^2+e^\zeta P^2\right)
  +\frac{\dot\zeta}{2}QP.
 \label{eq:narrativa-hamiltoniano-transporte}
\end{equation}
Le premier terme décrit l’évolution sur une ellipse de forme fixe.
Le second procède du transport d’une forme variable. La démonstration
de l’identité de transport établit que ce terme compense
exactement la variation géométrique et conserve l’action quadratique
correspondante. Son omission ou l’inversion de son signe détruisent cette
compensation.

L’interprétation exige de distinguer deux situations. Lorsqu’on change
uniquement de coordonnées, le terme supplémentaire est une connexion de
représentation. Un échange physique requiert, en revanche, une variation
relative entre système et référence. La réexpression d’un Hamiltonien
conserve le contenu de la représentation ; l’attribution d’une interaction
nécessite de déterminer les opérations relatives qui agissent sur le système.

Le couple angulaire directeur conserve une autre dépendance décisive : fixer $\alpha$
et la section d’action fixe également $\zeta$. Une évolution de la forme
doit donc être extraite des états et des opérations effectifs. La
déformation paramétrique à elle seule décrit une collection de représentations ;
son identification comme évolution HMT transmet les relations qui produisent
l’angle et le contre-angle.

Le transport conjoint de la forme et de la section d’action conserve
l’action normalisée, tandis que l’aire non normalisée change selon
le quotient entre les deux sections. Cette construction distingue trois
opérations : redistribuer une action conservée, transporter une section
d’action et changer la représentation ou les unités. La distinction permet
d’examiner les interactions en conservant un bilan qui inclut l’origine de la
variation, au lieu d’attribuer une génération d’énergie à une simple déformation
de coordonnées.

\subsection{Compatibilité conjointe et confrontation relationnelle}
\label{subsec:narrativa-compatibilidad-conjunta}

L’axe de recherche est la détermination des conditions de
compatibilité qu’une même structure impose à ses différentes
manifestations physiques. Deux lignes de composition se complètent.

La première réunit des relations statiques : les observables complémentaires
restituent les canaux, et la structure retrouvée détermine l’action,
la réponse constitutive et les facteurs de masse. La confrontation pertinente
est conjointe. Une modification qui préserve un chiffre doit également conserver
les autres relations qu’utilise la même construction. La vitesse,
Barbero--Immirzi et l’impédance fournissent un système de vérification ;
les rapports de masses et de fréquences en fournissent un autre.

La seconde ligne étudie les histoires effectives et compare leurs résultats
après avoir transporté de manière cohérente le système et sa référence. Les
différences de phase relative, de réponse constitutive relative ou de transfert
entre modes ont un contenu distinct des différences dues au
nom des coordonnées. La mémoire, l’orientation et les retours
interviennent au moyen des opérations qu’une lecture scalaire peut éliminer.

Une relation structurelle acquiert une capacité explicative et de confrontation
lorsque plusieurs lectures indépendantes doivent concorder parce qu’elles procèdent d’une
même structure. Sa valeur réside dans la rigidité de cette concordance :
l’action, le rythme, la réponse spatiale et la pondération informationnelle cessent de
pouvoir être choisis séparément au sein de la réalisation déclarée. La
composition réunit ces restrictions et permet de démontrer leurs conséquences
sans remplacer leur généalogie par des ajustements individuels des résultats.

\subsection{Clôture de phase et holonomie relative de l’état}
\label{subsec:narrativa-holonomia-relativa}

Le cycle dodécaphasique contient le plan dans lequel est réalisée la racine
interne de $-1$. L’inclusion de ce plan dans les douze composantes est
explicite ; le quart de tour utilisé ci-dessous procède de cette
construction. Sur celui-ci agit la forme de l’ellipse d’action déterminée
par le couple angulaire, dont la coordonnée est
\begin{equation}
 \xi=\frac{C^*}{A},\qquad
 \zeta=\operatorname{artanh}\xi.
 \label{eq:narrativa-coordenada-forma}
\end{equation}
Le quotient utilise les deux angles exprimés dans la même unité et a
la même valeur en degrés et en radians. La coordonnée réelle de forme est
définie pour $|\xi|<1$. La forme conjuguée correspond au
changement de signe de $\zeta$.

On compose le quart de tour d’une forme avec celui de sa conjuguée et avec
leurs inverses. Si $U_+$ et $U_-$ désignent les deux opérations, la
composition exacte est
\begin{equation}
 \boxed{
 H_{\mathrm{lazo}}=U_+U_-U_+^{-1}U_-^{-1}
 =\begin{pmatrix}\rho_{\mathrm{lazo}}&0\\[2pt]
                  0&\rho_{\mathrm{lazo}}^{-1}\end{pmatrix},
 \qquad
 \rho_{\mathrm{lazo}}=
 \left(\frac{1+\xi}{1-\xi}\right)^2.}
 \label{eq:narrativa-holonomia-lazo}
\end{equation}
La composition dilate une coordonnée et contracte la conjuguée par le facteur
inverse : elle conserve l’aire et modifie l’état. Pour $\xi\ne0$, son
résultat diffère de l’identité.

Les quatre opérations comprennent leurs inverses et leurs accroissements de phase
déclarés ont une somme nulle. L’ordre de composition conserve toutefois un
effet. Une description qui ne retient que les accroissements qui s’ajoutent et
se compensent élimine cet effet relatif. La clôture de la comptabilité
de phase et le retour de l’état complet sont donc des propriétés
distinctes de la séquence.

L’égalité admet un relèvement à l’espace initial à douze
composantes. L’opérateur agit par la dilatation réciproque sur le
plan considéré et comme l’identité sur son complément. Le lien avec
le cycle dodécaphasique est réalisé au moyen d’une transformation explicite ;
le plan conserve ainsi sa provenance au sein de l’espace des phases.

La notation $H_{\mathrm{lazo}}$ désigne cette composition ultérieure et
maintient distincte l’holonomie nonadique $\Gamma_9$. La réalisation
illustre de manière concrète une distinction structurelle : la clôture d’une
lecture du parcours peut coexister avec une transformation persistante
de l’état. La provenance des opérateurs et la portée de leur
relèvement demeurent exposées dans le développement démonstratif.

\subsection{Restitution du coefficient d’action à partir du transport relatif}
\label{subsec:narrativa-restitucion-accion}

Le couple angulaire de \eqref{eq:narrativa-par-angular} peut être exprimé avec
$s_0=10^{-34}\mathcal U_S$ et
$\hbar_{\mathrm{ret}}=s_0\eta_{\mathrm{ret}}$. Le coefficient
$\eta_{\mathrm{ret}}$ conserve son caractère adimensionnel, tandis que $s_0$
conserve l’échelle dimensionnelle. En incorporant le contre-angle complet
dans \eqref{eq:narrativa-holonomia-lazo}, on obtient
\begin{equation}
 \boxed{
 \rho_{\mathrm{lazo}}=
 \left(\frac{501\alpha+\eta_{\mathrm{ret}}}
              {499\alpha-\eta_{\mathrm{ret}}}\right)^2.}
 \label{eq:narrativa-lazo-alfa-accion}
\end{equation}
Cette expression utilise la branche positive
$0<\eta_{\mathrm{ret}}<499\alpha$. Les nombres $501$ et $499$ procèdent
de $500\pm1$, et $500$ procède de $1000/2$ : ce sont des conséquences des
normalisations qui interviennent déjà dans le couple angulaire.

La relation entre $\alpha$ et le coefficient d’action détermine la
séparation entre les réponses conjuguées après le parcours
composé. L’action participe aussi bien à l’échelle dimensionnelle qu’à
la forme qui gouverne le transport. La formule rend explicite cette
double intervention sans identifier les deux fonctions.

La relation admet l’inversion
\begin{equation}
 \boxed{
 \eta_{\mathrm{ret}}=
 \alpha\left[
 500\,\frac{\sqrt{\rho_{\mathrm{lazo}}}-1}
               {\sqrt{\rho_{\mathrm{lazo}}}+1}-1\right].}
 \label{eq:narrativa-inversa-accion}
\end{equation}
Avec $\alpha$ et les axes orientés conservés, la réponse de la boucle
permet de restituer le coefficient d’action. La lecture est strictement
croissante par rapport à $\rho_{\mathrm{lazo}}$ sur la branche déclarée.

On obtient une caractérisation opérationnelle : le coefficient d’action
peut être restitué au moyen de la dilatation relative produite par les
formes conjuguées du couple angulaire. Sa définition génératrice demeure
dans la construction précédente ; la caractérisation opérationnelle établit
comment le retrouver dans une représentation distincte.

L’échelle $s_0$ reste nécessaire pour retrouver la constante
dimensionnelle $\hbar_{\mathrm{ret}}$. La forme du transport conserve
une relation adimensionnelle ; l’échelle d’action détermine la
normalisation physique complète. Cette séparation identifie avec précision
quelle donnée l’inversion restitue et quelle structure supplémentaire transmet la
lecture dimensionnelle.

\subsection{Itération de l’holonomie et information opérationnelle de l’histoire}
\label{subsec:narrativa-informacion-operacional}

La répétition de la boucle conserve le même programme d’opérations et
produit une loi explicite d’accumulation :
\begin{equation}
 H_{\mathrm{lazo}}^n=
 \begin{pmatrix}\rho_{\mathrm{lazo}}^n&0\\
                 0&\rho_{\mathrm{lazo}}^{-n}\end{pmatrix}.
 \label{eq:narrativa-iteracion-lazo}
\end{equation}
Une préparation $(Q_0,P_0)$ se transforme en
\begin{equation}
 Q_n=\rho_{\mathrm{lazo}}^nQ_0,
 \qquad
 P_n=\rho_{\mathrm{lazo}}^{-n}P_0.
 \label{eq:narrativa-iteracion-estado}
\end{equation}
Le produit $Q_nP_n$ demeure constant et le transport conserve
l’aire symplectique, tandis que la proportion entre les deux
coordonnées change.

Une règle finie répétée peut laisser une empreinte cumulative qui
distingue ses itérations. Une préparation connue et une référence
conservée permettent de retrouver $n$ au moyen de l’accroissement logarithmique
d’une coordonnée non nulle, pour $\rho_{\mathrm{lazo}}\ne1$.
Lorsque seul l’état final est
conservé et que la préparation est inconnue, la même
restitution reste indéterminée.

L’information accessible dépend de la lecture retenue : bilan,
amplitude, proportion orientée, préparation ou opérateur complet. Chaque
choix peut identifier des histoires qu’une autre lecture distingue.
La comparaison de
\begin{equation}
 U_+U_-U_+^{-1}U_-^{-1}
 \qquad\text{et}\qquad
 U_+U_+^{-1}U_-U_-^{-1}
 \label{eq:narrativa-orden-multiplicidades}
\end{equation}
fournit une confrontation directe. Les deux produits contiennent les mêmes
opérations avec les mêmes multiplicités. Le second donne l’identité ;
le premier donne $H_{\mathrm{lazo}}$. Une statistique des fréquences
d’apparition est insuffisante pour prédire cette réponse : il faut
également conserver l’ordre.

Ce résultat précise quelle description statistique est insuffisante
pour l’objet étudié et quelles observations restituent la différence.
La théorie de l’information de Shannon admet des distributions sur des
séquences ordonnées et peut représenter cet ordre. La distinction
opérationnelle se réfère au contenu retenu par l’observable choisi.

L’information opérationnelle d’une histoire est caractérisée par
sa capacité à produire des réponses distinguables pour des préparations
et des détecteurs spécifiés. Cette définition permet de préciser combien
d’information une lecture conserve. Sa portée correspond à la représentation
utilisée : des histoires distinctes peuvent induire un même opérateur
réduit et conserver des différences dans d’autres composantes de la généalogie
TPK.

La trace de $H_{\mathrm{lazo}}$ coïncide avec celle de
$H_{\mathrm{lazo}}^{-1}$. Elle détecte l’intensité de l’effet, mais élimine
son orientation. Une lecture qui conserve les axes et compare leurs
réponses distingue les deux sens. Le choix du détecteur fait donc
partie de la restitution du parcours.

La finitude du programme, la croissance d’une mémoire observable et
un taux entropique décrivent des propriétés différentes. Un taux de
croissance combinatoire nul peut coexister avec des états distinguables.
L’entropie de la dynamique complète requiert de spécifier son espace
d’états et, le cas échéant, sa mesure. La composition démontre
l’accumulation et sa restitution dans la représentation déclarée ;
l’évaluation d’une entropie thermodynamique conserve ses variables
et ses conditions physiques propres.

\subsection{Potentiel spectral et réciprocité différentielle des réponses}
\label{subsec:narrativa-potencial-espectral}

Les canaux constitutifs conservent les coordonnées
$x=A_{\mathrm{rad}}$, $y=C^*_{\mathrm{rad}}$ et les incidences $90/120$.
Leur composition détermine l’identité
\begin{equation}
 \boxed{
 \frac{\partial\log\widehat c}{\partial y}
 =\frac{\partial\log\widehat Z}{\partial x}.}
 \label{eq:narrativa-reciprocidad-diferencial}
\end{equation}
Les grandeurs $\widehat c$ et $\widehat Z$ désignent la vitesse et
l’impédance normalisées du lecteur.

La sensibilité de la propagation à la séparation orientée se trouve
liée à la sensibilité de l’impédance à la coordonnée commune.
Les deux réponses partagent une fonction spectrale et leurs variations
doivent maintenir cette relation pour continuer d’appartenir à la même
construction.

Il existe une fonction $\mathscr F(x,y)$ dont les dérivées sont
\begin{equation}
 \partial_x\mathscr F=\log\widehat c,
 \qquad
 \partial_y\mathscr F=\log\widehat Z.
 \label{eq:narrativa-gradiente-potencial}
\end{equation}
Il s’agit d’un potentiel spectral défini à partir des canaux $90/120$,
avec une série convergente et une normalisation explicite. Son hessien est
défini négatif dans $x>|y|$, propriété qui permet de démontrer
l’unicité et la sensibilité de la restitution.

L’égalité entre les dérivées croisées établit une restriction
sur la variation conjointe des réponses au sein de cette collection
de réalisation. Une modification qui détruit l’identité
\eqref{eq:narrativa-reciprocidad-diferencial} change le lecteur
constitutif. La confrontation porte ainsi sur les relations de variation,
en plus des valeurs évaluées.

Le potentiel est univoque : l’intégrale de sa différentielle sur
une boucle fermée dans les deux coordonnées s’annule. La composition
ordonnée des transports peut conserver, en revanche, une holonomie
non triviale. La projection constitutive et le transport opératoriel
retiennent des aspects différents de l’histoire ; leur comparaison exige de
conserver leurs domaines et leurs opérations de lecture.

\subsection{Moment de Catalan et réponse contractive des incidences}
\label{subsec:narrativa-catalan-potencial}

La connexion entre le moment orienté de Catalan et le potentiel
constitutif est exprimée au moyen du moment de profondeur deux
\begin{equation}
 D_2(z)=\sum_{n\geq1}\frac{z^n}{n^2}.
 \label{eq:narrativa-momento-dos}
\end{equation}
Le moment de Catalan s’obtient comme partie imaginaire de
$D_2(is)$ ; en $s=1$, il retrouve la constante de Catalan. Le potentiel
constitutif utilise la même fonction sur les contractions des
canaux, en $q_\pm^{90}$ et $q_\pm^{120}$, avec les normalisations
correspondantes.

Le quart de tour orienté et la réponse contractive $90/120$
sont réunis au moyen d’un même moment spectral, évalué sur des
supports distincts. Catalan acquiert ici sa signification par
l’opération spectrale qu’il représente et par la relation de cette
opération avec la réponse constitutive. La collection fonctionnelle
conserve une information qui dépasse celle de sa valeur extrême.

La composition maintient également les différences entre ses objets.
La valeur isolée de Catalan ne détermine pas tous les canaux ;
la relation utilise le moment spectral complet, ses arguments
et l’orientation. La convergence de la série, les identités
de ses dérivées et sa compatibilité avec l’amplification font
partie du développement démonstratif de cette connexion.

\subsection{Restitution exacte et stabilité face à une précision finie}
\label{subsec:narrativa-estabilidad-restitucion}

Le couple $(\widehat c,\widehat Z)$ retrouve de manière unique les deux
paramètres du lecteur dans le domaine déclaré. Les bornes de
sensibilité relient, en outre, l’erreur des réponses à
l’erreur des paramètres restitués. L’unicité exacte et la
stabilité de l’inversion constituent des résultats distincts et
complémentaires.

Cette distinction sépare une information éliminée par la projection
d’une information conservée mais difficile à résoudre avec une précision
finie. Dans les régions de contractions intenses, des variations appréciables
des paramètres peuvent produire des différences très petites dans les
réponses. L’inversion reste unique pour des données exactes,
tandis qu’une incertitude finie peut empêcher de distinguer ces
variations.

La restitution généalogique exige de spécifier aussi bien quelle structure
peut être retrouvée que la stabilité avec laquelle elle peut l’être.
La marge de restitution désigne la borne qui quantifie cette stabilité
dans un domaine. Le choix de l’observable doit tenir compte de la différence
que l’on cherche à résoudre : une lecture peut être correcte et s’avérer
peu sensible à celle-ci.

La trace de la boucle offre un exemple de cette sélection. Sa lecture
élimine le sens du parcours ; la réponse sur des axes orientés
le conserve. La détermination du détecteur adéquat appartient à
la démonstration de récupérabilité et à la conception de la confrontation.

\subsection{Synthèse : coordonnées de la section et mémoire des opérations}
\label{subsec:narrativa-sintesis}

L’identité \eqref{eq:narrativa-cuatro-lecturas} conserve comme axe
la relation entre le facteur structurel d’une section, sa fréquence
propre, les deux longueurs réciproques et sa lecture thermique. Les conditions
demeurent explicites : le même système de coordonnées d’action, la réalisation
gravitationnelle circulaire et la température du cycle.

La composition opératorielle ajoute une distinction centrale. Le facteur
de la section détermine ses lectures coordonnées ; le transport
détermine le résultat de la composition d’opérations sur celle-ci et
l’empreinte que conserve leur ordre. Le facteur massique et la dilatation
$\rho_{\mathrm{lazo}}$ sont des objets différents. Leur relation doit
passer par les lecteurs correspondants pour distinguer une
transformation de l’état d’une identification massique.

Planck intervient dans la normalisation de l’action et dans la forme
angulaire du transport. Boltzmann intervient dans la relation entre
l’énergie et la pondération thermique. La réciprocité gravitationnelle conserve
le produit des longueurs dans sa réalisation déclarée. La vitesse
et l’impédance fournissent des lectures complémentaires des canaux.
Cette architecture admet conjointement une description des
compatibilités et une description de la mémoire opérationnelle.

La capacité explicative de la structure réside dans la détermination de
plusieurs réponses, l’imposition de relations entre elles et la spécification de
l’histoire qui peut être restituée à partir de ses observations. Les
coïncidences de valeurs sont accompagnées d’opérations,
de conditions de compatibilité et de critères de restitution.

La composition distingue le changement passif de coordonnées, qui
se télescope lorsque le système de coordonnées et la phase reviennent, du transport relatif
dont l’opérateur final est $H_{\mathrm{lazo}}\ne I$. Un échange
physique requiert de réaliser les opérations, leurs inverses et la
référence dans un même système, en incluant le contrôleur dans
le bilan. Ce critère conserve le contenu mathématique
d’une variation énergétique sans lui attribuer de production d’énergie
par le seul changement de lecture.

Les preuves de composition, d’itération, de restitution du coefficient
d’action, de réciprocité constitutive et de sensibilité soutiennent
ces caractérisations dans leurs domaines respectifs. La narration,
les définitions et les démonstrations forment ainsi un exposé continu :
la première identifie la signification des relations, les
deuxièmes fixent les objets et les troisièmes établissent les
opérations qui permettent de les composer et de les retrouver.

% Procedencia íntegra: CONTRAANGULO_ELECTRON_BARBERO_CATALAN.md, §§2--10.
% La gestión, el historial y los localizadores se conservan en gestion/.
\section{Composition et restitution des lectures angulaire, électronique et constitutive}
\label{md:compos:seccion}

\subsection{Coefficient d’action et construction du contre-angle}
\label{md:compos:contraangulo}

L’abréviation fonctionnelle $B(\alpha)$ réunit une évaluation en $\alpha$ ; son
argument appartient à la fonction et ne représente pas le produit d’une constante
indépendante par $\alpha$. Pour maintenir explicites les opérations qui
interviennent dans cette abréviation, on écrit
\begin{equation}
 B(\alpha)=\frac{9}{16}\alpha^2-\frac{5}{9}\alpha^3
 +\frac{7}{48}\alpha^4-\frac{1}{54}\alpha^5
 -\frac{90}{\pi}\mathcal C_\pi(\alpha),
 \label{md:compos:b-alpha}
\end{equation}
\begin{equation}
 \mathcal C_\pi(\alpha)=169\alpha^6+
 \frac{D_A\alpha^7}{1-D_A\alpha},
 \qquad D_A=\exp\!\left(-\frac{100\pi\alpha}{9}\right).
 \label{md:compos:correccion-regional}
\end{equation}
L’identité complète est
\begin{equation}
 C^*_{\deg}=\sqrt{\frac{1000\alpha}{\varphi}}+2B(\alpha).
 \label{md:compos:contraangulo-abreviado}
\end{equation}
En particulier, $\alpha$ appartient au radicand, et le terme radical est
conservé avec $2B(\alpha)$. Cette identité réécrit l’application construite
auparavant et maintient son objet et sa généalogie.

La définition originale conserve séparés le coefficient adimensionnel et la
section d’action :
\begin{equation}
 \eta_{\mathrm{ret}}=H_5-\frac{90}{\pi}\mathcal C_\pi,
 \qquad
 \hbar_{\mathrm{ret}}=\eta_{\mathrm{ret}}\,10^{-34}\mathcal U_S.
 \label{md:compos:seccion-accion}
\end{equation}
Par conséquent,
\begin{equation}
 \boxed{C^*_{\deg}=2\left(
 \frac{\hbar_{\mathrm{ret}}}{10^{-34}\mathcal U_S}+\alpha\right)
 =2(\eta_{\mathrm{ret}}+\alpha).}
 \label{md:compos:contraangulo-accion}
\end{equation}
L’expression utilise le coefficient adimensionnel de la section de Planck
réduite. La constante $h$ est obtenue ensuite par $h=2\pi\hbar$.
La composition consiste à extraire le coefficient d’action, à le composer avec
$\alpha$ et à appliquer la clôture bidirectionnelle. Ainsi est maintenu le type
dimensionnel de chaque opération.

La construction des masses exprime particule et antiparticule comme deux
sections orientées d’une même bande généalogique. L’opération
dodécaphasique est
\begin{equation}
 C_M(\tau)=-\operatorname{rev}(\tau),
 \qquad \tau_e=+++---+++---,
 \qquad C_M(\tau_e)=\tau_e,
 \label{md:compos:conjugacion-dodecafasica}
\end{equation}
prolongée au registre et à la représentation de charge. La route centrale est
fixe sous $C_M$, tandis que la représentation de charge distingue électron et
positron. Cette composition donne un contenu précis à la lecture bidirectionnelle.
Le facteur deux angulaire conserve sa fonction de clôture ; la conjugaison physique
conserve en outre les opérateurs indiqués. Le retour à $216$ restitue le
relèvement orienté : sa structure ne consiste pas à attribuer $108$
mises à jour à la particule et $108$ autres à l’antiparticule.

\subsubsection{Restitution radiale du coefficient d’action}
\label{md:compos:quinientos}

Nous écrivons $q_{\mathrm{rad}}=q_\varepsilon$ et $\epsilon=\varepsilon$ pour
le rapport radial et l’orientation de \eqref{rec:eq:eta-orientada}.
L’invariant radial orienté $\chi_{\mathrm{rad}}=\epsilon\log q_{\mathrm{rad}}$
permet d’écrire la restitution préalablement construite comme
\begin{equation}
 \frac{\eta_{\mathrm{ret}}}{\alpha}
 =-500\tanh\!\left(\frac{\chi_{\mathrm{rad}}}{2}\right)-1.
 \label{md:compos:restitucion-radial}
\end{equation}
Avec $\xi=C^*_{\deg}/A_{\deg}$ et $A_{\deg}=1000\alpha$, on obtient également
\begin{equation}
 \eta_{\mathrm{ret}}=\alpha(500\xi-1).
 \label{md:compos:restitucion-xi}
\end{equation}
Ce sont des opérateurs de récupération ultérieurs de la même section ; sa construction demeure
dans $H_5$ et dans le retour régional.

\subsection{Coordinatisations angulaires et conservation de l’action}
\label{md:compos:cartas}

La coordinatisation en degrés comporte un tour de $360$ unités et conserve les
partitions $12\cdot30=4\cdot90=3\cdot120=360$. La paire angulaire construite est
\begin{equation}
 A_{\deg}=1000\alpha,
 \qquad C^*_{\deg}=2(\eta_{\mathrm{ret}}+\alpha).
 \label{md:compos:par-angular}
\end{equation}
La conversion s’effectue une seule fois :
\begin{equation}
 x=\frac{\pi}{180}A_{\deg},\qquad
 y=\frac{\pi}{180}C^*_{\deg},\qquad
 q_\pm=\exp(-x\mp y).
 \label{md:compos:canales}
\end{equation}
Les exposants utilisent des représentants réels relevés et conservent le
nombre de tours. L’action satisfait
\begin{equation}
 \hbar\theta_{\mathrm{rad}}
 =\frac{h\theta_{\deg}}{360}.
 \label{md:compos:accion-cartas}
\end{equation}
Dans la fonctionnelle bilinéaire de Barbero, le transport de coordinatisation donne
\begin{equation}
 \frac{\pi A_{\deg}C^*_{\deg}}{180}
 =\frac{180xy}{\pi}.
 \label{md:compos:bilineal-cartas}
\end{equation}
La transformation bilinéaire fait partie de la même conversion qui est appliquée
aux arguments exponentiels ; l’omettre modifierait la fonctionnelle.

\subsection{Lecture électronique complète et retour d’action}
\label{md:compos:electron}

Les deux étapes de la même coordonnée électronique sont conservées :
\begin{equation}
 \varepsilon_e^{(0)}=\frac{\sqrt{3}}{4}
 \left[\exp\!\left(\frac{\varphi}{\pi^2}\right)-22\alpha^3\right]
 (1+15\Delta_4),\qquad
 \Delta_4=\frac{\pi-e\log\pi}{270}
 \left(1+\frac{\pi}{729}\right),
 \label{md:compos:electron-basal}
\end{equation}
\begin{equation}
 \boxed{\varepsilon_e^\beta=\varepsilon_e^{(0)}
 \left(\frac{H_5}{\eta_{\mathrm{ret}}}\right)^{80-54\alpha+6\Delta_4}.}
 \label{md:compos:electron-completo}
\end{equation}
Le centre $(5,5)$, son inversion, les déformations $(8,2)/(2,8)$ et le
défaut unitaire de quotient sont des antécédents de la valeur. La norme
$\sqrt{3}/4$ provient du commutateur entre projections additives et
multiplicatives ; le caractère exponentiel utilise l’orientation
$\varphi/\pi^2$ ; $22=27-5$ conserve un complément régional et $15=3\cdot5$
les incidences correspondantes. Les lecteurs de la route produisent
$(80,54,6)$. Le quotient $H_5/\eta_{\mathrm{ret}}$ incorpore le retour
d’action à la même section centrale.

La fonction du contre-angle est explicitée en substituant son inverse démontrée :
\begin{equation}
 \boxed{\varepsilon_e^\beta=\varepsilon_e^{(0)}
 \left(\frac{H_5}{C^*_{\deg}/2-\alpha}\right)^{80-54\alpha+6\Delta_4}.}
 \label{md:compos:electron-contraangulo}
\end{equation}
Cette expression conserve la génération électronique précédente. La même
section d’action qui participe au contre-angle détermine le facteur qui
fait passer la lecture basale à la lecture complète. Le quotient simplifie la base dimensionnelle
commune ; la réalisation absolue de masse conserve son lecteur temporel et sa coordinatisation
dimensionnelle. La substitution est une composition de dépendances construites,
sans introduire d’angle externe comme générateur électronique.

\subsection{Fonctionnelle de Barbero et transition hexade--octade}
\label{md:compos:barbero}

L’incidence maintient la même paire marquée lors de l’agrandissement du support :
\begin{equation}
 \mathcal F_6=\mathsf H\times\binom{\mathsf H}{2}
 \hookrightarrow
 \mathcal F_8=\mathsf O\times\binom{\mathsf H}{2},
 \qquad |\mathcal F_6|=90,\qquad |\mathcal F_8|=120.
 \label{md:compos:incidencia}
\end{equation}
Les cardinaux sont $6\cdot15$ et $8\cdot15$ ; le référent partagé permet de
comparer les deux incidences. La différence normalisée est $-1/12$. La chaîne
dodécaphasique conserve douze contributions hexadiques et une contribution
octadique de frontière de signe contraire. Le lecteur
\begin{equation}
 S_m(q)=\frac{q^m}{1-q^{3m}}
 =\sum_{j\geq0}q^{m+3mj}
 \label{md:compos:serie-retorno}
\end{equation}
conserve hauteur et retours. Ainsi se forment $12S_{90}-S_{120}$ et sa différence
entre les deux orientations. La fonctionnelle complète est
\begin{equation}
 \gamma=\frac{\pi A_{\deg}C^*_{\deg}}{180}
 +12\bigl[S_{90}(q_-)-S_{90}(q_+)\bigr]
 -\bigl[S_{120}(q_-)-S_{120}(q_+)\bigr].
 \label{md:compos:funcional-barbero}
\end{equation}
Barbero est, dans cette réalisation, l’évaluation orientée conjointe de la paire
angulaire et des retours de la transition d’incidences. Sa participation
à l’échelle d’aire se compose ensuite avec $\ell_P^2=G\hbar/c^3$ et avec
l’opérateur d’occupation. L’entropie d’intrication conserve en outre
l’état réduit et ses probabilités ; sa reconstruction nécessite ces données
et n’est pas déterminée par l’aire scalaire isolée.

\subsubsection{Élimination de la coordonnée commune électron--Barbero}
\label{md:compos:eliminacion}

Définissons $\kappa_e=80-54\alpha+6\Delta_4$. Sur la branche positive, avec
$\kappa_e\neq0$, l’équation électronique implique exactement
\begin{equation}
 \eta_{\mathrm{ret}}=H_5
 \left(\frac{\varepsilon_e^{(0)}}{\varepsilon_e^\beta}\right)^{1/\kappa_e}.
 \label{md:compos:eliminacion-eta}
\end{equation}
Pour exprimer la composition en maintenant visibles ses opérations, posons
$x=50\pi\alpha/9$, $y=\pi(\eta_{\mathrm{ret}}+\alpha)/90$ et
$\mathcal R(q)=12S_{90}(q)-S_{120}(q)$. Alors
\begin{equation}
 \gamma=\frac{100\pi}{9}\alpha(\eta_{\mathrm{ret}}+\alpha)
 +\mathcal R(e^{-x+y})-\mathcal R(e^{-x-y}).
 \label{md:compos:barbero-eta}
\end{equation}
La substitution de \eqref{md:compos:eliminacion-eta} dans cette identité produit
une relation explicite entre la lecture électronique complète et Barbero,
en conservant $\alpha$, $H_5$, $\Delta_4$ et les incidences. C’est une conséquence
de compatibilité de la construction. La branche physique angulaire utilisée
maintient $0<y<x$ ; l’opération n’introduit pas de masse métrologique comme donnée
génératrice.

Dans l’extension munie de l’étiquette d’orientation $\epsilon=\pm1$,
$C^*_{\deg,\epsilon}=2\epsilon(\eta_{\mathrm{ret}}+\alpha)$, l’échange
$\epsilon\mapsto-\epsilon$ laisse fixe $\eta_{\mathrm{ret}}$, conserve
l’expression électronique et permute $q_+$ et $q_-$. La fonctionnelle de Barbero change
de signe. Cette covariance distingue une lecture scalaire d’action d’une
évaluation orientée. L’identification à la conjugaison de charge conserve
en outre $C_M$ et la représentation électromagnétique des feuillets ; le signe angulaire
isolé ne remplace pas ces opérateurs.

\subsection{Moment orienté de Catalan et réponse constitutive}
\label{md:compos:catalan}

Le cycle dodécaphasique réalise un plan orienté de quart de tour. Son
coefficient résolvant est $k_4(s)=s/(1+s^2)$. Le moment de profondeur deux
est fixé par
\begin{equation}
 \mathcal C_2(s)=\sum_{k\geq0}
 \frac{(-1)^ks^{2k+1}}{(2k+1)^2},\qquad
 \mathcal D=s\frac{d}{ds},\qquad
 \mathcal D^2\mathcal C_2(s)=k_4(s),\qquad
 \mathcal C_2(1)=G_{\mathrm{Cat}}.
 \label{md:compos:momento-catalan}
\end{equation}
L’identité différentielle est utilisée pour $0<s<1$ ; la valeur extrême est
entendue avec les conditions de convergence et de restitution du moment déjà
établies. La définition intrinsèque de Catalan est la valeur extrême de ce
moment orienté de profondeur. Sa collection conserve plus d’information que
la valeur extrême. La réponse constitutive utilise cette collection :
\begin{equation}
 f(s)=\frac{1+s+s^2}{s(1+s)}\mathcal D^2\mathcal C_2(s)
 =\frac{1-s^3}{1-s^4},\qquad
 r_\pm=f(s_\pm),\qquad s_\pm=q_\pm^{30}.
 \label{md:compos:catalan-respuesta}
\end{equation}
Les puissances $3$ et $4$ correspondent à $90=3\cdot30$ et $120=4\cdot30$.
\begin{proposition}[Restitution fonctionnelle avec condition de frontière]
\label{md:compos:inversa-catalan}
La collection constitutive $f$ détermine le moment orienté par
\[
 \mathcal C_2(s)=\int_0^s\log\frac{s}{t}\,
 \frac{1+t}{1+t+t^2}f(t)\,dt,\qquad 0<s<1.
\]
C’est l’unique solution de
$\mathcal D^2\mathcal C_2(s)=s(1+s)f(s)/(1+s+s^2)$
qui tend vers zéro à l’origine.
\end{proposition}
\begin{proof}
Pour la collection construite,
$(1+t)f(t)/(1+t+t^2)=1/(1+t^2)$. L’intégrale converge, et sa
dérivée logarithmique est
$\mathcal D\mathcal C_2(s)=\int_0^s(1+t^2)^{-1}\,dt$.
Appliquer à nouveau $\mathcal D$ donne $s/(1+s^2)$. L’intégrale est
$O(s)$ à l’origine. Deux solutions diffèrent de $a+b\log s$ ;
la limite nulle impose $a=b=0$.
Pour $s<1$, intégrer la série géométrique sur des compacts et utiliser
$\int_0^s t^{2k}\log(s/t)\,dt=s^{2k+1}/(2k+1)^2$
retrouve la série de \eqref{md:compos:momento-catalan}.
Sa convergence absolue dominée permet $s\uparrow1$ et restitue Catalan.
\end{proof}
La condition de frontière fixe les deux constantes d’intégration.
La restitution utilise la fonction constitutive avec son domaine ; la
valeur extrême de Catalan et les deux valeurs de canal sont des lectures de
cette fonction déjà construite. À l’extrémité, l’identité doublement
dérivée conserve son interprétation par limite d’Abel.

L’opérateur constitutif $V=r_+P_++r_-P_-$ détermine
\begin{equation}
 \widehat\varepsilon=r_+^2,\qquad \widehat\mu=r_-^2,
 \qquad \widehat Z=\frac{r_-}{r_+},\qquad
 \widehat c=(r_+r_-)^{-1}.
 \label{md:compos:lectores-constitutivos}
\end{equation}
Avec $\psi=f^{-1}$ et
\begin{equation}
 \mathcal H_B(s)=\frac{12s^3}{1-s^9}
 -\frac{s^4}{1-s^{12}}-\frac{(\log s)^2}{20\pi},
 \label{md:compos:h-barbero}
\end{equation}
la factorisation construite prend la forme
\begin{equation}
 \gamma=\mathcal H_B(\psi(r_-))-\mathcal H_B(\psi(r_+)).
 \label{md:compos:factorizacion-barbero}
\end{equation}
La partie logarithmique restitue exactement la contribution angulaire de
Barbero, conservée en degrés ou transportée intégralement en radians.

\subsection{Restitution par la vitesse normalisée et Barbero}
\label{md:compos:difeomorfismo}

\paragraph{Proposition.}
Dans la collection constitutive normalisée à deux canaux étiquetés, l’application
\begin{equation}
 (r_+,r_-)\longmapsto(\widehat c,\gamma),\qquad
 (3/4,1)^2\longrightarrow(1,16/9)\times\mathbb R
 \label{md:compos:mapa-global}
\end{equation}
est un difféomorphisme. Son inverse récupère les deux valeurs propres étiquetées.
Les projecteurs de feuillet appartiennent à la représentation de départ : deux
scalaires ne récupèrent pas une base spatiale arbitraire.

\paragraph{Démonstration.}
L’expression régulière
$f(s)=(1+s+s^2)/(1+s+s^2+s^3)$ donne
\begin{equation}
 f'(s)=-\frac{s^2(s^2+2s+3)}{(1+s+s^2+s^3)^2}<0.
 \label{md:compos:derivada-f}
\end{equation}
Ainsi, $f$ est analytique et strictement décroissante, avec $f(0^+)=1$ et
$f(1^-)=3/4$. Pour $0<s<1$,
\begin{equation}
 \mathcal H_B'(s)=
 \frac{36s^2(1+2s^9)}{(1-s^9)^2}
 -\frac{4s^3(1+2s^{12})}{(1-s^{12})^2}
 -\frac{\log s}{10\pi s}>0.
 \label{md:compos:derivada-h}
\end{equation}
En effet, le quotient du premier terme positif par le second terme
positif avant de le soustraire est
\begin{equation}
 \frac{9}{s}\frac{1+2s^9}{1+2s^{12}}
 \left(\frac{1-s^{12}}{1-s^9}\right)^2>9.
 \label{md:compos:cota-derivada-h}
\end{equation}
Le dernier terme de \eqref{md:compos:derivada-h} est également positif.
Par conséquent, $F_B=\mathcal H_B\circ\psi$ possède une dérivée strictement négative,
$F_B(3/4^+)=+\infty$ et $F_B(1^-)=-\infty$.

$c=\widehat c$ étant fixé, nous écrivons $z=\log\widehat Z$ et
$r_\pm=c^{-1/2}\exp(\mp z/2)$. L’intervalle admissible est
\begin{equation}
 |z|<a(c):=\min\left\{\log c,\log\frac{16}{9c}\right\}.
 \label{md:compos:dominio-z}
\end{equation}
La fonction $\Gamma_c(z)=F_B(r_-)-F_B(r_+)$ satisfait
\begin{equation}
 \Gamma_c'(z)=\frac12
 \bigl[r_-F_B'(r_-)+r_+F_B'(r_+)\bigr]<0.
 \label{md:compos:derivada-gamma}
\end{equation}
Lorsque $z$ augmente jusqu’à $a(c)$, soit $r_-$ atteint $1$, soit $r_+$ atteint $3/4$,
soit les deux se produisent ; $\Gamma_c$ tend vers $-\infty$. À l’extrémité opposée,
elle tend vers $+\infty$. Il existe donc une unique solution pour chaque $\gamma$
réel. La dérivée non nulle et le théorème des fonctions implicites fournissent
une inverse lisse globale. Cela prouve l’énoncé. Dans l’orientation
originale $y>0$, on a $z<0$ et $\gamma>0$.

La vitesse normalisée conserve le produit des réponses ; Barbero
distingue de manière monotone leur répartition orientée. Ensemble,
ils restituent l’information spectrale que chaque lecture isolée ne détermine pas.

L’énoncé concerne la collection des lecteurs constitutifs. L’image
effectivement produite par le générateur HMT complet est un sous-ensemble de
cette collection : la restriction à cette image conserve l’injectivité et l’ordre
causal. La normalisation et les marques de feuillet sont des données nécessaires de la
restitution. La preuve n’introduit pas de valeurs CODATA pour choisir les canaux ni
n’affirme que chaque point synthétique de la collection soit un état engendré.

\subsection{Restitution sur l’image engendrée}
\label{md:compos:restitucion-completa}

Une fois $r_\pm$ récupérés à partir de $(\widehat c,\gamma)$, on obtient
$s_\pm=\psi(r_\pm)$. Alors
\begin{equation}
 A_{\deg}=-\frac{3}{\pi}\log(s_+s_-),\qquad
 C^*_{\deg}=\frac{3}{\pi}\log\frac{s_-}{s_+}.
 \label{md:compos:restitucion-angular}
\end{equation}
Sur l’image engendrée et dans son orientation positive,
\begin{equation}
 \alpha=\frac{A_{\deg}}{1000},\qquad
 \eta_{\mathrm{ret}}=\frac{C^*_{\deg}}{2}-\alpha>0.
 \label{md:compos:restitucion-alpha-eta}
\end{equation}
Avec les coordonnées régionales $\pi,\varphi,e$ et les sections dimensionnelles
déjà construites, cette restitution redonne $H_5$, le facteur
$R_{\mathrm{act}}=H_5/\eta_{\mathrm{ret}}$ et la lecture électronique complète
de la section~\ref{md:compos:electron}. Elle redonne également la section de
Planck. Si l’on conserve en outre la base positionnelle à douze blocs, l’origine
et la retenue, l’identité
\begin{equation}
 \kappa_{\mathrm{ag}}=\pi+e-\varphi-4-\alpha
 \label{md:compos:restitucion-k}
\end{equation}
permet la restitution du registre $K$ décrite par son lecteur. Il s’agit
d’un parcours de reconstruction ; son sens inverse ne constitue pas une nouvelle
flèche génératrice à partir de valeurs physiques externes.

La gravité et la conversion thermique conservent les sections du même
état. Le théorème~\ref{md:rigidez:teorema} et
\eqref{md:rigidez:reloj} déterminent d’abord $L_\alpha$ et sa
représentation circulaire $R_{108}=L_\alpha$ :
\begin{equation}
 G_a=\frac{c_{\mathrm{int}}^3R_{108}^2}{\hbar_a},\qquad
 \frac{G_a\hbar_a}{c_{\mathrm{int}}^3}=R_{108}^2,
 \qquad
 k_B\Theta_{\mathrm{clk},a}\log3=\frac{h_a}{108t_0}.
 \label{md:compos:gravedad-termica}
\end{equation}
Le rayon chronologique, la vitesse physique et la base thermique conservent leurs
constructeurs. En particulier, l’horloge de la représentation circulaire est
liée à $L_\alpha$ et $c_{\mathrm{int}}$. La première
composition explicite la réciprocité action--gravité ; la seconde réunit
action, durée et information tritique, avec individualisation de la température
et de $k_B$ dans leurs bases déclarées.

\subsection{Définitions structurelles des lectures composées}
\label{md:compos:definiciones}

\begin{description}
 \item[$\alpha$.] Coordonnée scalaire de clôture dodécaphasique compatible avec
 les régions et avec le registre arithmétique--géométrique $K$ ; elle transmet cette
 compatibilité au secteur angulaire et d’action.
 \item[Contre-angle.] Coordonnée angulaire de la clôture bidirectionnelle qui
 combine le coefficient de la section d’action de retour avec $\alpha$,
 dans la coordinatisation de $360$ unités.
 \item[Planck réduite.] Section d’action dont le coefficient régional et l’échelle
 déterminantale sont construits avant le contre-angle ; dans l’ellipse, mesure
 d’aire orientée par unité de phase paramétrique.
 \item[Électron.] Section centrale persistante ; sa lecture scalaire compose
 incompatibilité additive--multiplicative, transfert régional, défaut
 de clôture et retour d’action. L’équation complète conserve la section,
 en plus de sa valeur de masse.
 \item[Barbero--Immirzi.] Fonctionnelle orientée du transport angulaire et des
 retours de l’incidence hexade--octade ; dans la collection constitutive,
 coordonnée qui, avec la vitesse normalisée, restitue le spectre
 étiqueté.
 \item[Catalan.] Valeur extrême du moment orienté de profondeur deux du
 quart de tour dodécaphasique. Sa collection est reliée différentiellement à la
 réponse du vide.
 \item[$G$.] Coefficient commun de la réalisation radiale qui relie
 énergie et longueur conjuguée en conservant simultanément
 $H\Lambda=\hbar cI$ et $R\Lambda=L_\alpha^2I$.
 La composition longitudinale et polaire détermine sa normalisation ;
 son expression dimensionnelle est $c^3L_\alpha^2/\hbar$.
 \item[Boltzmann.] Transducteur entre la droite thermique et l’énergie de
 décision, avec une normalisation informationnelle tritique et une base thermique déclarées.
\end{description}
Ces formulations intrinsèques se fondent sur les objets et équations
construits. Leur portée est celle de cette réalisation ; les identifications
physiques conservent leurs démonstrations correspondantes et ne se déduisent pas de
la seule dénomination structurelle.

\subsection{Transport du bloc électronique nilpotent}
\label{md:compos:nilpotentes}

\paragraph{Proposition.}
Soient
\begin{equation}
 S=\begin{pmatrix}1&0\\0&-1\end{pmatrix},\qquad
 J=\begin{pmatrix}0&1\\-1&0\end{pmatrix},\qquad
 n_\pm=\frac{S\pm J}{2},\qquad
 D=\operatorname{diag}(r,r^{-1}),\quad r>0.
 \label{md:compos:bloque-nilpotente}
\end{equation}
Définissons $n_{\pm,D}=Dn_\pm D^{-1}$ et $g=D^{-T}D^{-1}$. Alors
\begin{equation}
 n_{\pm,D}^2=0,\qquad
 \{n_{+,D},n_{-,D}\}=I,\qquad
 (n_{+,D})^{\dagger_g}=n_{-,D},
 \label{md:compos:identidades-nilpotentes}
\end{equation}
où $A^{\dagger_g}=g^{-1}A^Tg$. L’ellipse transportée depuis
$Y^TY=2\hbar$ au moyen de $X=DY$ satisfait $X^TgX=2\hbar$.

\paragraph{Démonstration.}
Le calcul matriciel direct donne
$S^2=I$, $J^2=-I$ et $SJ+JS=0$. Ainsi,
\[
 n_\pm^2=\tfrac14(S^2+J^2\pm SJ\pm JS)=0,
 \qquad
 n_+n_-+n_-n_+=\tfrac12(S^2-J^2)=I.
\]
En outre, $S^T=S$ et $J^T=-J$ impliquent $n_+^T=n_-$. La conjugaison par $D$
préserve produits, carrés et identité. Pour l’adjoint métrique,
\[
 g^{-1}(Dn_+D^{-1})^Tg
 =DD^T D^{-T}n_-D^T D^{-T}D^{-1}
 =Dn_-D^{-1}=n_{-,D}.
\]
Enfin, $D^TgD=I$, d’où
$X^TgX=Y^TD^TgDY=Y^TY=2\hbar$. Les affirmations sont ainsi prouvées.

La vérification algébrique exacte a lieu dans
$\mathbb Q[r,r^{-1}]$ ; la restriction $r>0$ fixe la réalisation elliptique
orientée utilisée. Cette composition conserve les fonctions particulières
des différents facteurs $1/2$ du développement : une coïncidence numérique
entre eux ne constitue pas une identification de leurs opérateurs.

% Procedencia íntegra: HIBRIDACION_ACCION_MASA_PROPAGACION.md, §§2--7.
% La gestión, el historial y los localizadores se conservan en gestion/.
\section{Action, masse et propagation : compatibilité structurelle et transport}
\label{md:hib:seccion}

\subsection{Section d’action partagée et loi des masses}
\label{md:hib:accion-masas}

Soient
\begin{equation}
 s_0=10^{-34}\mathcal U_S,\qquad
 \hbar_{\mathrm{ret}}=\eta_{\mathrm{ret}}s_0,\qquad
 R_{\mathrm{act}}=\frac{H_5}{\eta_{\mathrm{ret}}}.
 \label{md:hib:accion}
\end{equation}
Les coordonnées construites en degrés sont maintenues :
\begin{equation}
 A_{\deg}=1000\alpha,\qquad
 C^*_{\deg}=2(\eta_{\mathrm{ret}}+\alpha).
 \label{md:hib:angulos}
\end{equation}
Ce n’est qu’ensuite qu’elles sont transportées en radians au moyen de
$x=\pi A_{\deg}/180$ et $y=\pi C^*_{\deg}/180$. Les canaux sont
$q_\pm=\exp(-x\mp y)$, avec $x>|y|$. La forme elliptique s’exprime par
\begin{equation}
 \zeta=\operatorname{artanh}(y/x).
 \label{md:hib:forma}
\end{equation}
La composition utilise d’abord le coefficient adimensionnel
$\eta_{\mathrm{ret}}$ et conserve ainsi les types d’action et d’angle.
Dans cette notation, $\eta_{\mathrm{ret}}$ est le coefficient d’action,
$\eta_h(\Gamma)$ appartient à la tour d’une route et $\zeta$ exprime la
déformation de l’ellipse ; leurs fonctions restent séparées.

On travaille avec $t_0,c_{\mathrm{int}},\eta_{\mathrm{ret}}>0$, un caractère
convergent et un facteur $\Xi_\Gamma>0$. Le lecteur massique construit détermine
\begin{equation}
 \omega_0=\frac{2\pi}{108t_0},\qquad
 E_{\mathrm{clk}}=\hbar_{\mathrm{ret}}\omega_0,\qquad
 m_{\mathrm{clk}}=\frac{E_{\mathrm{clk}}}{c_{\mathrm{int}}^2}.
 \label{md:hib:reloj}
\end{equation}
Pour une route à réalisation scalaire,
\begin{equation}
 m_\Gamma^\beta=m_{\mathrm{clk}}\Xi_\Gamma,\qquad
 \Xi_\Gamma=b_e\,
 \chi_{\mathrm{ref}}(\sigma_\Gamma-\sigma_e,\eta_\Gamma-\eta_e)
 R_{\mathrm{act}}^{\kappa_\Gamma}.
 \label{md:hib:factor-ruta}
\end{equation}
$\Xi_\Gamma$ abrège le facteur structurel de route. On conserve expressément
\begin{equation}
 b_e=\frac{\varepsilon_e^{(0)}\mathcal U_E}{E_{\mathrm{clk}}}.
 \label{md:hib:normalizacion-electron}
\end{equation}
L’état central conserve sa normalisation propre par rapport au quantum
chronologique :
\begin{equation}
 \varepsilon_e^{(0)}=\frac{\sqrt3}{4}
 \left[\exp\!\left(\frac{\varphi}{\pi^2}\right)-22\alpha^3\right]
 (1+15\Delta_4),\qquad
 \varepsilon_e^\beta=\varepsilon_e^{(0)}
 R_{\mathrm{act}}^{80-54\alpha+6\Delta_4}.
 \label{md:hib:electron}
\end{equation}
Le caractère raffiné inclut les termes
\begin{equation}
 n_Ax+\frac{s_C}{6}y+k_{\Delta_4}\Delta_4+b_\zeta\zeta_{270}
 +\sum_h N_h(q_-^h-q_+^h),\qquad
 N_h=\eta_h+b_{120}\mathbf 1_{\{h=120\}},
 \label{md:hib:caracter}
\end{equation}
avec $h$ dans le semi-groupe engendré par $90$ et $120$. Le terme de hauteur
$120$ est inclus une fois. Les coefficients primitifs de retour sont
construits au moyen de
\begin{equation}
 \begin{aligned}
 a_n(\Gamma)&=\operatorname{Tr}(R_\Gamma U_\Gamma^n),\\
 p_n(\Gamma)&=\frac1n\sum_{d\mid n}\mu_{\mathrm{Mob}}(d)
                         a_{n/d}(\Gamma),\\
 \eta_h(\Gamma)&=p_h(\Gamma)-p_h(\Gamma_e),
 \end{aligned}
 \label{md:hib:retornos-primitivos}
\end{equation}
où $\mu_{\mathrm{Mob}}$ est la fonction de Möbius arithmétique.
L’action participe à l’échelle absolue, au rapport $R_{\mathrm{act}}$
et à $y$, qui détermine les canaux. Par conséquent, la simplification d’une
échelle dimensionnelle conserve les autres dépendances structurelles.
Lorsqu’on fait varier l’horloge, $b_e=\varepsilon_e^{(0)}\mathcal U_E/E_{\mathrm{clk}}$
varie avec elle : conserver la lecture électronique exige de transporter le
caractère complet, au lieu de fixer séparément ce quotient.
Dans la réalisation longitudinale intégrée, on utilise
$t_0=\pi_{\rm HMT}L_\alpha/(54c_{\mathrm{int}})$ de
\eqref{md:rigidez:reloj}.

\subsubsection{Fréquence et longueur de chaque section massive}
\label{md:hib:frecuencia-longitud}

En définissant $E_\Gamma=m_\Gamma c_{\mathrm{int}}^2$ comme énergie de repos,
$\Omega_\Gamma=E_\Gamma/\hbar_{\mathrm{ret}}$,
$\nu_\Gamma=\Omega_\Gamma/(2\pi)$ et
$R_{108}=c_{\mathrm{int}}/\omega_0$, on obtient
\begin{equation}
 \begin{aligned}
 E_\Gamma&=\hbar_{\mathrm{ret}}\omega_0\Xi_\Gamma,&
 \Omega_\Gamma&=\omega_0\Xi_\Gamma,\\
 \nu_\Gamma&=\frac{\Xi_\Gamma}{108t_0},&
 \bar\lambda_\Gamma&=
 \frac{\hbar_{\mathrm{ret}}}{m_\Gamma c_{\mathrm{int}}}
 =\frac{R_{108}}{\Xi_\Gamma}.
 \end{aligned}
 \label{md:hib:frecuencias}
\end{equation}
\begin{proof}
On substitue $m_\Gamma=m_{\mathrm{clk}}\Xi_\Gamma$ et l’on effectue les
simplifications dans une même coordinatisation dimensionnelle. Pour des routes positives $X,Y$,
la propriété de caractère et la normalisation commune donnent
\begin{equation}
 \frac{\Omega_Y}{\Omega_X}=\frac{m_Y}{m_X}
 =\chi_{\mathrm{ref}}(\sigma_Y-\sigma_X,\eta_Y-\eta_X)
 R_{\mathrm{act}}^{\kappa_Y-\kappa_X}.
 \label{md:hib:cocientes}
\end{equation}
\end{proof}
La simplification est dimensionnelle ; les canaux et la mémoire
relative demeurent. Le secteur nul reste en dehors du domaine de l’inverse de masse.
Dans une fibre non scalaire, on conserve l’opérateur ou sa mesure spectrale,
au lieu de choisir un représentant ponctuel arbitraire.

\subsubsection{Fréquence relevée et compatibilité des normalisations}
\label{md:hib:tres-frecuencias}

La construction chronologique produit également la fréquence relevée
d’holonomie
\begin{equation}
 \omega_j=\frac{\theta_j+2\pi w_j}{108t_0},\qquad
 \lambda_j=r_j\exp(i\theta_j),
 \label{md:hib:frecuencia-levantada}
\end{equation}
où $w_j$ conserve l’enroulement de l’histoire. Sa lecture énergétique est
\begin{equation}
 E_j=\hbar_{\mathrm{ret}}\omega_j
 \chi_{\mathrm{full}}R_{\mathrm{act}}^{\kappa_j}.
 \label{md:hib:energia-levantada}
\end{equation}
Cette fréquence, $\omega_0$ et $\Omega_\Gamma$, correspondent à des lecteurs
distincts. Si les deux expressions énergétiques représentent le même mode et
le même exposant, leur compatibilité exacte exige
\begin{equation}
 \frac{\omega_j}{\omega_0}\chi_{\mathrm{full}}
 =b_e\chi_{\mathrm{ref,rel}}.
 \label{md:hib:compatibilidad-normalizaciones}
\end{equation}
Les deux normalisations sont identifiées au moyen de cette égalité et ne sont pas
multipliées entre elles. La phase réduite $\theta_j$ ne récupère pas à elle seule
$w_j$. Ainsi, un retour de phase reste distinct de l’histoire
chronologique complète. L’enroulement conserve son constructeur ; cette distinction
n’autorise pas à le choisir librement ni à produire de l’énergie sans dynamique.

\subsection{Réciprocité spatiale et composition gravitationnelle}
\label{md:hib:gravedad}

La section gravitationnelle circulaire exprime le résultat du
théorème~\ref{md:rigidez:teorema} :
$G_a=c_{\mathrm{int}}^3L_\alpha^2/\hbar_a$ et $R_{108}=L_\alpha$.
La longueur provient du retour inscrit et précède cette
identification circulaire. En maintenant
$\hbar_a=\hbar_{\mathrm{ret}}$ et la convention réduite
$r_g=Gm/c_{\mathrm{int}}^2$, la masse précédente détermine
\begin{equation}
 \boxed{r_{g,\Gamma}=R_{108}\Xi_\Gamma,\qquad
 \bar\lambda_\Gamma=\frac{R_{108}}{\Xi_\Gamma},\qquad
 r_{g,\Gamma}\bar\lambda_\Gamma=R_{108}^2.}
 \label{md:hib:reciprocidad-radial}
\end{equation}
\begin{proof}
La substitution de $G_a$ et de $m_\Gamma$ donne
\[
 \frac{G_am_\Gamma}{c_{\mathrm{int}}^2}
 =\frac{c_{\mathrm{int}}^3R_{108}^2}{\hbar_{\mathrm{ret}}}
 \frac{\hbar_{\mathrm{ret}}\omega_0\Xi_\Gamma}
 {c_{\mathrm{int}}^4}
 =R_{108}\Xi_\Gamma.
\]
La seconde égalité provient de \eqref{md:hib:frecuencias} ; leur produit
prouve la troisième.
\end{proof}
La même coordonnée de route augmente une lecture radiale et contracte l’autre.
En logarithmes,
\begin{equation}
 \log\frac{r_{g,\Gamma}}{R_{108}}=\log\Xi_\Gamma,
 \qquad
 \log\frac{\bar\lambda_\Gamma}{R_{108}}=-\log\Xi_\Gamma.
 \label{md:hib:radios-logaritmicos}
\end{equation}
Le centre géométrique des deux lectures est fixé par l’aire $R_{108}^2$.

Dans la version opératoire bornée, soit $X$ autoadjoint et borné, avec
$X\geq\delta I$ pour un certain $\delta>0$, dans une fibre massive. Les opérateurs
\begin{equation}
 M=m_{\mathrm{clk}}X,\qquad
 \Lambda=R_{108}X^{-1},\qquad R_g=R_{108}X
 \label{md:hib:operadores-radiales}
\end{equation}
satisfont $R_g\Lambda=\Lambda R_g=R_{108}^2I$ sur tout l’espace.
L’affirmation utilise le calcul fonctionnel du même opérateur. Le noyau
nul reste en dehors de l’inverse. Le présent énoncé conserve le régime
borné uniformément positif ; l’extension non bornée exige de démontrer
séparément les domaines, les coupures compatibles et le passage à la limite.
La section~\ref{md:rigidez:variaciones} démontre en outre la conservation
des deux identités sous des variations non commutatives et une réduction
compatible de mémoire.

L’interprétation proposée caractérise le facteur massique relatif comme
paramètre de dilatation réciproque de ces deux lectures. La condition circulaire
appartient au résultat. Les longueurs ainsi reliées ne sont pas identifiées
par cette seule égalité à des rayons matériels mesurés ni à un horizon
électronique. L’identité radiale isolée est utilisée pour sa fonction dans la
composition, sans lui attribuer de priorité historique indépendante.

\subsection{Réponse constitutive, orientation et grandeur}
\label{md:hib:constitutivas}

Avec $T=\exp(-xI-yR)$, $R=P_+-P_-$ et les projecteurs de feuillet conservés,
le lecteur temporel est
\begin{equation}
 \begin{aligned}
 V&=(I-T^{90})(I-T^{120})^{-1}=r_+P_++r_-P_-,\\
 r_\pm&=f(s_\pm),\qquad
 s_\pm=\exp[-30(x\pm y)],\qquad
 f(s)=\frac{1-s^3}{1-s^4}.
 \end{aligned}
 \label{md:hib:operador-constitutivo}
\end{equation}
On récupère $\widehat\varepsilon=r_+^2$, $\widehat\mu=r_-^2$,
$\widehat Z=r_-/r_+$ et $\widehat c=(r_+r_-)^{-1}$.
Le lien constitutif de vitesse établit
$c_{\mathrm{int}}=c_{\mathrm{vac}}=\widehat c\,u_C$ dans la même coordinatisation.
La collection orientée de Catalan participe à $f$ par son moment
résolvant, en conservant la structure antérieure à la valeur extrême.

La décomposition exacte
\begin{equation}
 V=\widehat c^{-1/2}
 \exp\!\left[-\frac12(\log\widehat Z)R\right]
 \label{md:hib:descomposicion-v}
\end{equation}
distingue la dilatation commune et la répartition entre feuillets. L’échange des feuillets
conserve $\widehat c$ et inverse $\widehat Z$ ; la fonctionnelle de Barbero change
de signe. La section~\ref{md:compos:difeomorfismo} démontre que
$(\widehat c,\gamma)$ récupère $r_+,r_-$ dans la collection normalisée.
Sa restriction à l’image engendrée conserve cette récupérabilité ;
l’énoncé d’inversion ne déclare pas engendrée toute paire synthétique.

\subsubsection{Convexité et évaluation des canaux séparés}
\label{md:hib:convexidad}

\begin{proposition}
\label{md:hib:proposicion-convexidad}
Pour $x>0$ fixé et $|y|<x$, la fonction
$y\mapsto\log\widehat c(x,y)$ est paire et strictement convexe.
Par conséquent,
\begin{equation}
 \widehat c(x,y)\geq\widehat c(x,0)
 =\frac{1}{f(e^{-30x})^2},
 \label{md:hib:cota-velocidad}
\end{equation}
avec égalité exactement pour $y=0$.
\end{proposition}
\begin{proof}
Soit $g(u)=\log f(e^{-u})$, pour $u>0$. On a
\begin{equation}
 g'(u)=\frac{3}{e^{3u}-1}-\frac{4}{e^{4u}-1}>0,
 \label{md:hib:derivada-g}
\end{equation}
\begin{equation}
 g''(u)=-\frac{9}{4\sinh^2(3u/2)}
 +\frac{16}{4\sinh^2(2u)}<0.
 \label{md:hib:segunda-g}
\end{equation}
La première inégalité provient de la décroissance de
$a/(\exp(au)-1)$ sur $a>0$. Pour la seconde, la fonction
$a/\sinh(au/2)$ décroît strictement : sa dérivée a le signe de
$\sinh v-v\cosh v$, négatif pour $v>0$. Comme
\[
 \log\widehat c(x,y)=-g(30(x+y))-g(30(x-y)),
\]
sa dérivée seconde par rapport à $y$ est
\begin{equation}
 -900\bigl[g''(30(x+y))+g''(30(x-y))\bigr]>0.
 \label{md:hib:segunda-velocidad}
\end{equation}
La parité fixe l’unique minimum en zéro.
\end{proof}
On obtient également
\begin{equation}
 \begin{aligned}
 \log\widehat c(x,y)
 &=-2g(30x)-900g''(30x)y^2+O(y^4),\\
 \log\widehat Z(x,y)
 &=-60g'(30x)y+O(y^3).
 \end{aligned}
 \label{md:hib:desarrollos-canales}
\end{equation}
La vitesse perd le signe d’orientation, mais conserve une réponse
quadratique à la séparation des canaux. L’impédance conserve une information
orientée dès le premier ordre. Faire la moyenne de $x+y$ et $x-y$ avant d’appliquer le
lecteur élimine cette contribution. L’affirmation est une conséquence
démontrée de la réponse constitutive, distincte d’une nouvelle mesure.

Sur la branche positive d’action,
$\eta_{\mathrm{ret}}=(90/\pi)(y-x/500)$ ; par conséquent,
\begin{equation}
 \hbar_{\mathrm{ret}}=s_0\frac{90}{\pi}(y-x/500),\qquad
 E_k=s_0\frac{90u_C|k|}{\pi}(y-x/500)\widehat c(x,y).
 \label{md:hib:energia-angular}
\end{equation}
Pour $x$ fixé, $|k|>0$ et $x/500<y<x$, cette dernière fonction croît
strictement en $y$. C’est une propriété de la collection de réalisation
radiative homogène. Son application à une trajectoire engendrée conserve le
constructeur de $y$ ; la variation d’un paramètre de cette collection n’équivaut pas
à elle seule à une variation expérimentale d’une constante universelle.
Sous inversion d’orientation, on transporte également l’étiquette d’action,
en maintenant la branche positive correspondante.

\subsection{Action, horloge et information thermique}
\label{md:hib:termica}

La section thermique conserve
$k_B\Theta_{108}\log3=E_{\mathrm{clk}}$, avec
$\Theta_{108}=\Theta_{\mathrm{clk,ret}}$ dans la section de retour.
Nous définissons
$\beta_{108}=(k_B\Theta_{108})^{-1}$. Pour
$E_\Gamma=E_{\mathrm{clk}}\Xi_\Gamma$, on déduit exactement
\begin{equation}
 \beta_{108}E_\Gamma=(\log3)\Xi_\Gamma,
 \qquad
 \boxed{\exp(-\beta_{108}E_\Gamma)=3^{-\Xi_\Gamma}.}
 \label{md:hib:peso-termico}
\end{equation}
L’expression est un poids de Gibbs non normalisé. Pour une collection finie
d’états, ou une collection ayant une somme de partition convergente, et des
multiplicités $g_\Gamma$, on obtient
\begin{equation}
 Z_{\mathrm{part}}=\sum_\Gamma g_\Gamma3^{-\Xi_\Gamma},\qquad
 p_\Gamma=\frac{g_\Gamma3^{-\Xi_\Gamma}}{Z_{\mathrm{part}}},\qquad
 \frac{p_Y/g_Y}{p_X/g_X}=3^{-(\Xi_Y-\Xi_X)}.
 \label{md:hib:particion}
\end{equation}
On suppose une réalisation canonique à $\Theta_{108}$ et une origine énergétique
commune. Lorsqu’un potentiel chimique intervient, il est incorporé à l’énergie
correspondante. L’indice $\Gamma$ énumère des états ou des groupes physiques
distincts ; plusieurs histoires équivalentes ne produisent pas automatiquement une
multiplicité. Les entropies adimensionnelles sont
\begin{equation}
 \begin{aligned}
 H_{\mathrm{micro}}
 &=\log Z_{\mathrm{part}}+(\log3)\sum_\Gamma p_\Gamma\Xi_\Gamma,\\
 H_{\mathrm{grupos}}
 &=H_{\mathrm{micro}}-\sum_\Gamma p_\Gamma\log g_\Gamma.
 \end{aligned}
 \label{md:hib:entropias}
\end{equation}
Pour affirmer la finitude de l’entropie, il faut en outre
\begin{equation}
 \sum_\Gamma g_\Gamma\Xi_\Gamma3^{-\Xi_\Gamma}<\infty.
 \label{md:hib:condicion-entropia}
\end{equation}
La finitude de $Z_{\mathrm{part}}$ garantit à elle seule la normalisation.
Dans les collections finies, les deux conditions sont vérifiées.

La formule relie la structure massique à la fréquence propre et à la pondération
thermique dans le même cadre. Elle conserve la distinction entre énergie et chaleur,
entre coordonnée réelle $\Xi_\Gamma$ et dénombrement entier de trits, et entre une
réalisation à $\Theta_{108}$ et une affirmation sur la température de
l’univers. Pour des occupations bosoniques ou fermioniques libres par mode, avec
un potentiel chimique nul, la même substitution produit
\begin{equation}
 \overline n_\Gamma=\frac{1}{3^{\Xi_\Gamma}\mp1}.
 \label{md:hib:ocupaciones}
\end{equation}
Le signe moins correspond à Bose et le signe plus à Fermi. Le total du
groupe dégénéré est $g_\Gamma\overline n_\Gamma$. Dans Bose, on exige
$\Xi_\Gamma>0$ et le mode nul est traité séparément. Ce sont des réalisations
statistiques ultérieures appliquées à un spectre déjà construit.

\subsubsection{Compatibilité conjointe de quatre lectures}
\label{md:hib:cuatro-lecturas}

En composant les sections précédentes, sous les mêmes hypothèses,
\begin{equation}
 \boxed{\Xi_\Gamma=\frac{\Omega_\Gamma}{\omega_0}
 =\frac{r_{g,\Gamma}}{R_{108}}
 =\frac{R_{108}}{\bar\lambda_\Gamma}
 =\frac{E_\Gamma}{k_B\Theta_{108}\log3}.}
 \label{md:hib:compatibilidad-cuatro}
\end{equation}
Dans la collection circulaire, $\Xi=1$ identifie la section de fréquence
$\omega_0$, les rayons coïncidents $R_{108}$ et l’énergie $E_{\mathrm{clk}}$.
Son poids thermique non normalisé est $1/3$. C’est un point de compatibilité des
lectures. L’existence d’une espèce de l’atlas avec $\Xi=1$, ou d’une
seconde espèce engendrée pour chaque section réciproque $1/\Xi$, nécessite
son appartenance produite à l’atlas ; elle ne se déduit pas de cette égalité.

\subsection{Identité énergétique de l’ellipse et transport exact}
\label{md:hib:elipse}

Dans l’ellipse d’action,
\begin{equation}
 a_S=\sqrt{2\hbar}\,e^{\zeta/2},\qquad
 b_S=\sqrt{2\hbar}\,e^{-\zeta/2},\qquad
 Q=a_S\cos\theta,\qquad P=b_S\sin\theta,
 \label{md:hib:elipse-coordenadas}
\end{equation}
avec $\zeta=\operatorname{artanh}(C^*_{\deg}/A_{\deg})$, on a
$a_Sb_S=2\hbar$, une aire orientée
$\mathcal A(\theta)=\hbar\theta$ et $\mathcal A(2\pi)=h$.
La phase paramétrique $\theta$ et l’angle polaire d’un point de l’ellipse
conservent leurs significations géométriques respectives.

La réalisation LC déclarée utilise
\begin{equation}
 L_\zeta=\frac{Z_*}{\omega}e^{-\zeta},\qquad
 C_\zeta=\frac{1}{Z_*\omega}e^\zeta.
 \label{md:hib:lc}
\end{equation}
En coordonnées d’action $Q=\sqrt{Z_*}\,q$ et
$P=\Phi/\sqrt{Z_*}$,
\begin{equation}
 H_0=\frac{\omega}{2}
 \left(e^{-\zeta}Q^2+e^\zeta P^2\right),\qquad
 U_E=\hbar\omega\cos^2\theta,\qquad
 U_B=\hbar\omega\sin^2\theta.
 \label{md:hib:energia-elipse}
\end{equation}
On obtient $H_0=\hbar\omega$ par composition des deux contributions.
L’horloge et le transducteur appartiennent à la réalisation déclarée ; la section
$\hbar$ conserve sa construction antérieure au circuit. L’impédance
$Z_{\mathrm{LC}}/Z_*=e^{-\zeta}$ et l’impédance constitutive du vide
$f(s_-)/f(s_+)$ sont des lecteurs distincts. Leur antécédent angulaire commun
n’établit pas d’égalité entre les deux.

\subsubsection{Transport discret de la forme quadratique}
\label{md:hib:transporte-discreto}

\begin{proposition}
\label{md:hib:proposicion-discreta}
Soient $S_\zeta=\operatorname{diag}(e^{\zeta/2},e^{-\zeta/2})$,
$g_\zeta=S_\zeta^{-T}S_\zeta^{-1}$ et
$\mathsf R(-\theta)$ une rotation. Pour des valeurs de forme $\zeta_n$ et des
phases $\theta_n$, le transport
\begin{equation}
 X_{n+1}=S_{\zeta_{n+1}}\mathsf R(-\theta_n)S_{\zeta_n}^{-1}X_n
 =T_nX_n
 \label{md:hib:transporte}
\end{equation}
conserve exactement
$I_n=X_n^Tg_{\zeta_n}X_n/2$ et a pour déterminant un.
\end{proposition}
\begin{proof}
La substitution et l’identité $\mathsf R^T\mathsf R=I$ donnent
\begin{equation}
 T_n^Tg_{\zeta_{n+1}}T_n=g_{\zeta_n}.
 \label{md:hib:isometria-transportada}
\end{equation}
La forme quadratique est conservée lorsqu’on évalue les deux membres sur $X_n$.
Les facteurs de $T_n$ ont pour déterminant un, d’où
$\det T_n=1$.
\end{proof}
C’est une composition du transport radial préalablement construit. Si
$\zeta_n,\theta_n$ sont extraits d’une route TPK, l’identité s’applique à
cette séquence ; la proposition algébrique maintient séparée la sélection
de la route.

\subsubsection{Transport différentiable et terme de connexion}
\label{md:hib:conexion}

Si $\zeta(t)$ est de classe $C^1$ et $\omega(t)>0$ est continue, la relation
$X=S_\zeta Y$ et l’évolution $\dot Y=-\omega J_0Y$, avec
\begin{equation}
 J_0=\begin{pmatrix}0&-1\\1&0\end{pmatrix},
 \label{md:hib:j-canonico}
\end{equation}
donnent
\begin{equation}
 \dot Q=\frac{\dot\zeta}{2}Q+\omega e^\zeta P,\qquad
 \dot P=-\omega e^{-\zeta}Q-\frac{\dot\zeta}{2}P.
 \label{md:hib:ecuaciones-transporte}
\end{equation}
Avec les conventions $\dot Q=\partial_PH$ et
$\dot P=-\partial_QH$, l’hamiltonien est, à un terme dépendant
seulement du temps près,
\begin{equation}
 \boxed{H_{\mathrm{tot}}=
 \frac{\omega}{2}\left(e^{-\zeta}Q^2+e^\zeta P^2\right)
 +\frac{\dot\zeta}{2}QP.}
 \label{md:hib:hamiltoniano-conexion}
\end{equation}
L’invariant
\begin{equation}
 I=\frac12\left(e^{-\zeta}Q^2+e^\zeta P^2\right)
 \label{md:hib:invariante}
\end{equation}
est conservé exactement, sans nécessiter d’approximation adiabatique.
\begin{proof}
En dérivant, on obtient
\[
 \dot I=\frac{\dot\zeta}{2}
 \left(e^\zeta P^2-e^{-\zeta}Q^2\right)
 +e^{-\zeta}Q\dot Q+e^\zeta P\dot P.
\]
La substitution de \eqref{md:hib:ecuaciones-transporte} annule les termes
proportionnels à $\dot\zeta$ et les deux termes $\omega QP$.
Ainsi, $\dot I=0$. Si l’on omet le terme mixte de
\eqref{md:hib:hamiltoniano-conexion}, il reste
\begin{equation}
 \dot I=\frac{\dot\zeta}{2}
 \left(e^\zeta P^2-e^{-\zeta}Q^2\right).
 \label{md:hib:defecto-sin-conexion}
\end{equation}
Le signe mixte contraire double ce défaut.
\end{proof}
Sur l’ellipse, $QP=I\sin(2\theta)$ et $\dot\theta=-\omega$, de sorte que
\begin{equation}
 H_{\mathrm{tot}}=I\omega+\frac{I\dot\zeta}{2}\sin(2\theta).
 \label{md:hib:hamiltoniano-elipse}
\end{equation}
La matrice de la forme quadratique a pour déterminant
$\omega^2-\dot\zeta^2/4$. Pour $\omega>0$, elle est définie positive
exactement lorsque $|\dot\zeta|<2\omega$, et semi-définie positive dans le cas
d’égalité. Ce seuil concerne la positivité : l’existence du
transport et la stabilité physique conservent leurs conditions respectives.
La conservation de $I$ n’implique pas la conservation de $H_{\mathrm{tot}}$,
qui dépend explicitement du temps.

La version quantique exige la symétrisation du terme mixte,
$\dot\zeta(QP+PQ)/4$. Cette expression indique l’opération requise sans
présupposer de représentation quantique supplémentaire.

\subsubsection{Compatibilité avec la section d’action génératrice du contre-angle}
\label{md:hib:compatibilidad-accion}

Dans la construction directrice,
\begin{equation}
 \zeta=\operatorname{artanh}
 \left[\frac{2(\eta_{\mathrm{ret}}+\alpha)}{1000\alpha}\right].
 \label{md:hib:forma-genealogica}
\end{equation}
Fixer $\alpha$ et $\eta_{\mathrm{ret}}$ fixe également $\zeta$.
Le résultat avec $\zeta(t)$ variable est donc une extension de
transport entre formes. Pour représenter une trajectoire physique de la paire
directrice, la variation doit provenir de la généalogie correspondante, et non
d’une déformation choisie de manière indépendante.

Le transport peut également être composé exactement lorsque la
section positive d’action change entre deux états, en conservant la même coordinatisation
dimensionnelle. Soient
\begin{equation}
 \hbar_n>0,\qquad
 \rho_{\mathrm{act},n}=\frac{\hbar_{n+1}}{\hbar_n},\qquad
 \widetilde T_n=\sqrt{\rho_{\mathrm{act},n}}\,
 S_{\zeta_{n+1}}\mathsf R(-\theta_n)S_{\zeta_n}^{-1}.
 \label{md:hib:transporte-conforme}
\end{equation}
Alors
\begin{equation}
 \det\widetilde T_n=\rho_{\mathrm{act},n},\qquad
 \widetilde T_n^Tg_{\zeta_{n+1}}\widetilde T_n
 =\rho_{\mathrm{act},n}g_{\zeta_n},\qquad
 \frac{I_{n+1}}{\hbar_{n+1}}=\frac{I_n}{\hbar_n}.
 \label{md:hib:accion-normalizada}
\end{equation}
Les égalités découlent de \eqref{md:hib:isometria-transportada} en multipliant
par le facteur scalaire. C’est un transport conformément symplectique : il conserve
l’action normalisée et conserve l’aire d’action non normalisée
exactement lorsque $\rho_{\mathrm{act},n}=1$.
Sa limite différentiable a pour trace $\dot\hbar/\hbar$. Un générateur
hamiltonien linéaire sur un plan à forme symplectique fixe a une trace nulle ;
ainsi, une variation matérielle de $\hbar$ exige de spécifier la forme
transportée ou les degrés de liberté avec lesquels l’action est échangée.
Le terme mixte de déformation conserve à lui seul une trace nulle.

Le résultat est un critère de compatibilité des réalisations. Pour
interpréter des variations numériques dues à un changement d’unités, on
transporte d’abord les bases : le changement de coordinatisation maintient la section
physique. La composition n’affirme pas une variabilité expérimentale de la
constante de Planck.

\subsubsection{Transformation de représentation et déformation matérielle}
\label{md:hib:interpretacion-transporte}

Si $S_\zeta$ est un changement passif de coordonnées, le terme mixte est une
connexion de représentation. Les observables sont également transportées,
\begin{equation}
 L_X(S_\zeta Y)=L_Y(Y).
 \label{md:hib:observables-transportados}
\end{equation}
Cette égalité conserve la lecture physique lorsqu’on change l’expression de
l’hamiltonien, sans introduire d’interaction supplémentaire ni de travail physique.

Une déformation matérielle exige d’obtenir $\zeta$ à partir de la mise à jour TPK et
de montrer un changement relatif entre système et référence, tous deux transportés
de manière cohérente. On peut confronter une phase relative, une réponse
$Z_{\mathrm{sistema}}/Z_{\mathrm{referencia}}$ ou un transfert entre
modes définis par le détecteur. Il faut distinguer deux histoires TPK non
équivalentes par changement de représentation et conserver leurs conditions de
lecture. Lorsqu’un échange énergétique est attribué, le bilan comprend
champ, mémoire et dispositif. Le théorème de transport délimite cette
distinction et maintient séparée la loi dynamique qui la rend physique.

\subsection{Quantité de mouvement, confrontations et portée de la composition}
\label{md:hib:momento}

La réalisation de Clifford--Dirac reçoit la masse de route. Dans une même
section d’action, son coefficient s’écrit
\begin{equation}
 \frac{m_\Gamma c}{\hbar}
 =\frac{\omega_0\Xi_\Gamma}{c}.
 \label{md:hib:coeficiente-dirac}
\end{equation}
Dans une réalisation libre et homogène, avec la relation relativiste de
dispersion déclarée, celle-ci se réexprime comme
\begin{equation}
 \Omega^2=c^2|k|^2+\omega_0^2\Xi_\Gamma^2.
 \label{md:hib:dispersion}
\end{equation}
La réécriture conserve les hypothèses de cette réalisation et ses antécédents
de connexion spatio-temporelle. Les moments de retour $q_-^h-q_+^h$
sont des évaluations des tours et se distinguent des moments mécaniques.

Le facteur $\Xi_\Gamma$ réapparaît comme poids massique, fréquence relative,
dilatation radiale et exposant thermique. La réponse constitutive détermine
la conversion spatiale commune. Barbero conserve une information orientée
complémentaire à la vitesse. L’action fixe la mesure aréale ; l’horloge
introduit le taux temporel ; la forme et son transport spécifient la répartition
et son évolution.

Les confrontations correspondantes comprennent les quotients de masses et de fréquences ;
la compatibilité de la vitesse, de l’impédance et de Barbero ; la conservation de la forme
quadratique sous transport ; et les poids thermiques une fois la température et les
multiplicités fixées. Ces relations préservent plus de données qu’une coïncidence
décimale isolée. Leurs preuves maintiennent le bilan énergétique et les
conditions des entropies considérées ; elles n’impliquent pas une entropie
thermodynamique nulle.

% Ampliación para insertar después del bloque térmico de hibridacion.tex.
% Procedencia, dependencias y alcance de los controles: INFORME.md.
\section{Lecture marquée et conjugaison thermométrique de l'action}
\label{md:r27:termica}

L'action orientée, le calendrier, les lecteurs de capacité et l'archive de mémoire procèdent de la même dynamique enrichie APP--TRIT--TPK. Les deux feuillets conservent résidu et quotient ; le régime tritique conserve l'orientation ; le transport $U_t$ et son holonomie $\Gamma_9$ précèdent la lecture réduite $K_{\rm ph}$. Le retour de phase conserve un historique qui continue de croître. La présente construction compose ces résultats dans un registre de référence et d'essai, au sein de la structure discrète conjointe du continu déjà construite. Sa règle thermométrique est spécifiée par deux opérations sur ce registre.

\subsection{Capacité, référence spectrale et origine conservée}
\label{md:r27:capacidad}
La capacité d'un préfixe d'indice $t\ge0$ et son reste sont
\begin{equation}
 C_t=\lfloor\log_{10}(9^{t+1})\rfloor,\qquad
 9^{t+1}=10^{C_t}R_t,\qquad 1\le R_t<10.
 \label{md:r27:capacidad-resto}
\end{equation}
Le dénombrement s'évalue par puissances entières. Les inégalités $10^{20}\le9^{21}<10^{21}$ et $10^{20}\le9^{22}<10^{21}$, avec $10^{j-1}\le9^j<10^j$ pour $1\le j\le21$, prouvent que la première vacance positive a pour indice $21$ et pour capacité $20$. Dans la lecture tritique de capacité, les résidus $0,1,2$ correspondent respectivement à $+1,-1,0$. Les capacités en $t=21,22,23,24$ sont $20,21,22,23$ ; le premier retour ultérieur du régime est donc $t=24$. Le préfixe contient $25$ positions et exprime la capacité $23$. Les fonctions différenciées de $(20,24,25,23)$ sont ainsi conservées. Le retour du régime conserve une phase nonadique différente ; la phase revient en $t=30$, avec la capacité $29$.

Le porteur du lecteur est $\mathcal D=\mathcal D_0\oplus\mathcal D_1
\oplus\mathcal D_2$, de rangs $(1,380,649)$ et de graduation $N=0,1,2$. La construction conserve deux graduations : $N$ distingue le préfixe et $L$ l'occupation de l'espace extérieur ou symétrique. Pour $m=20$, $n=24$, soient $\mathcal F_2^J$ le sous-espace fixe extérieur de degré deux de deux feuillets échangés et $\mathcal B_{\le2}^J$ le sous-espace fixe symétrique de degré au plus deux. Le porteur est $\mathbb C\Omega_{\rm ext}\oplus\mathcal F_2^J\oplus\mathcal B_{\le2}^J$. L'involution $J$ échange les deux feuillets par l'échange gradué ; dans le secteur extérieur, chacun de ses espaces propres a pour dimension $m$. Les caractères des sous-espaces fixes sont
\begin{equation}
 \begin{aligned}
 \chi_F(z)&=\frac{(1+z)^{2m}+(1-z^2)^m}{2},\\
 \chi_B(z)&=\frac{P_n(z)^2+P_n(z^2)}{2},\\
 P_n(z)&=\prod_{j\ge1}(1-z^j)^{-n}.
 \end{aligned}
 \label{md:r27:caracteres-fock}
\end{equation}
En effet, le projecteur $(I+J)/2$ prend la moyenne du caractère total et de la trace de l'échange. Dans l'espace extérieur, les valeurs propres $+1,-1$ de $J$ produisent $(1+z)^m(1-z)^m$ ; dans le produit symétrique, seules les paires identiques contribuent à la trace de l'échange, produisant $P_n(z^2)$. Pour les coefficients jusqu'au degré deux, le polynôme $P_n(z)=1+nz+n(n+3)z^2/2+O(z^3)$ suffit ; les produits formels sont finis en chaque degré. L'extraction des coefficients démontre
\begin{equation}
 \begin{split}
 \dim\mathcal F_2^J&=m(m-1)=380,\\
 (\dim\mathcal B_0^J,\dim\mathcal B_1^J,\dim\mathcal B_2^J)
 &=(1,n,n(n+2))=(1,24,624).
 \end{split}
 \label{md:r27:grados-fock}
\end{equation}
La décomposition simultanée $(N,L;\text{multiplicité})$ est donc
\[
 \begin{gathered}
 (0,0;1),(1,2;380),(2,0;1),\\
 (2,1;24),(2,2;624).
 \end{gathered}
\]
Le rang $649=1+24+624=25^2+24$ conserve trois occupations ; le vide externe et le vide du terme symétrique restent orthogonaux. Le caractère conjoint est $1+380xy^2+x^2(1+24y+624y^2)$. L'assignation de $N=0,1,2$ aux coupes $24,27,30$, de capacités $23,26,29$, spécifie le lecteur utilisé dans cette réalisation ; elle maintient $L$ comme observable concomitante. En particulier, $\operatorname{Tr}q^L=2+24q+1004q^2$ et $\operatorname{Tr}q^N=1+380q+649q^2$ conservent des opérations distinctes.

L'origine de la référence spectrale réside dans \eqref{md:antmem:cuantica:eq:5.1}--\eqref{md:antmem:cuantica:eq:5.10}. Le couplage nonadique $8:1$ et la boucle ordonnée
$H_1=\left(\begin{smallmatrix}2&1\\1&1\end{smallmatrix}\right)$, avec la préparation de deux modes purs identiques et la coordinatisation de covariance transportée conjointement, donnent $\nu^2=1+(8/81)[\operatorname{tr}(H_1H_1^{\mathsf T})-2]=121/81$. La somme des carrés des entrées de $H_1$ vaut sept, de sorte que $\nu=11/9$. Le théorème~\ref{md:antmem:gauss:thm:espectro-gaussiano} donne alors $r=(\nu-1)/(\nu+1)=1/10$ et les valeurs propres $(1-r)r^j$. On utilise cette préparation démontrée, y compris sa réduction de mémoire. La division $j=3n+b$, $0\le b<3$, conserve quotient et résidu et satisfait
\begin{equation}
 (1-r)r^{3n+b}=(1-r^3)(r^3)^n\frac{r^b}{1+r+r^2}.
 \label{md:r27:factorizacion}
\end{equation}
Par conséquent, le quotient a pour raison $q=r^3=1/1000$ et le résidu conserve la distribution $\operatorname{diag}(100,10,1)/111$. L'augmentation du quotient après achèvement des trois résidus conserve la retenue. Si l'énergie élémentaire est $\epsilon$, l'énergie de l'indice est $\epsilon(3n+b+1/2)$ : le saut entre blocs est $3\epsilon$ à la même température que la référence.

Nous définissons l'opérateur positif de référence et son intensité :
\begin{equation}
 \begin{gathered}
 \eta=\bigoplus_{n=0}^{2}r^{23+3n}I_{d_n},\qquad
 (d_0,d_1,d_2)=(1,380,649),\\
 b_*:=\operatorname{Tr}\eta=\frac{1380649}{10^{29}}.
 \end{gathered}
 \label{md:r27:eta}
\end{equation}
En particulier, $b_*=1{,}380649\times10^{-23}$ et $0<b_*<1$. L'origine $23$ est conservée dans $\eta$ ; normaliser uniquement son état élimine ce facteur des probabilités relatives.

La section d'action de \eqref{md:hib:reloj} fournit $\epsilon_a=h_a/(108t_0)$ et $\beta_{\rm clk}=\ln3/\epsilon_a$. Pour le relèvement énergétique additif de ce lecteur de capacité,
\begin{equation}
 \begin{split}
 H_{\rm tot}(t)&=2(t+1)\epsilon_a,\\
 H_{\rm mem}(t)&=\frac{\epsilon_a}{\ln3}\ln R_t,\\
 H_{\rm cap}(t)&=H_{\rm tot}(t)-H_{\rm mem}(t)
 =\frac{\epsilon_a\ln10}{\ln3}C_t.
 \end{split}
 \label{md:r27:energia-capacidad}
\end{equation}
La preuve consiste à prendre les logarithmes dans \eqref{md:r27:capacidad-resto}. Par substitution, $e^{-\beta_{\rm clk}H_{\rm cap}(t)}=10^{-C_t}$. Sur $\mathcal D$, cela donne $H_{\rm cap}=(\epsilon_a\ln10/\ln3)(23I+3N)$ et $\operatorname{Tr}e^{-\beta_{\rm clk}H_{\rm cap}}=b_*$. Le reste demeure enregistré comme contribution énergétique ; la compensation est une lecture différente de l'élimination d'un facteur par trace partielle.

\subsection{Identité tensorielle et spécification des lecteurs}
\label{md:r27:lectores}
\begin{lemma}[Information incrémentale du registre produit]
Pour toute matrice positive $\eta$ et toute matrice de densité $\rho$, avec $\mathscr H(X)=-\operatorname{Tr}(X\ln X)$ et $0\ln0=0$, on a
\begin{equation}
 \begin{gathered}
 \mathscr H(\eta\otimes\rho)-\mathscr H(\eta)
 =\operatorname{Tr}(\eta)s(\rho),\\
 s(\rho)=-\operatorname{Tr}(\rho\ln\rho).
 \end{gathered}
 \label{md:r27:tensor}
\end{equation}
\end{lemma}
\begin{proof}
Sur les supports, le logarithme du produit tensoriel est $\ln\eta\otimes I+I\otimes\ln\rho$. La factorisation de la trace et $\operatorname{Tr}\rho=1$ donnent $\mathscr H(\eta\otimes\rho)=\mathscr H(\eta)+\operatorname{Tr}(\eta)s(\rho)$. La continuité de $x\ln x$ en zéro étend l'égalité aux noyaux non triviaux.
\end{proof}

Pour \eqref{md:r27:eta}, nous formons un registre normalisé et sa marque :
\begin{equation}
 \widehat\eta=\eta\oplus(1-b_*)|f\rangle\langle f|,
 \qquad\Omega_\rho=\widehat\eta\otimes\rho,\qquad
 P=I_{\mathcal D}\oplus0,
 \label{md:r27:registro}
\end{equation}
où $P$ agit par l'identité sur le facteur d'essai. Le drapeau $f$ complète la probabilité de cette réalisation. Toute archive microscopique supplémentaire est transportée comme telle ; le drapeau conserve exclusivement cette distinction de sélection. Pour un Hamiltonien fixe $H$, on spécifie
\begin{equation}
 \begin{split}
 \mathcal S_P(\rho)&=-\operatorname{Tr}
 [P\Omega_\rho(I\otimes\ln\rho)]=b_*s(\rho),\\
 \mathcal E_P(\rho)&=
 \frac{\operatorname{Tr}[P\Omega_\rho(I\otimes H)]}
 {\operatorname{Tr}(P\Omega_\rho)}=\operatorname{Tr}(\rho H).
 \end{split}
 \label{md:r27:par-lector}
\end{equation}
La première opération lit l'information avant le conditionnement ; la seconde lit l'énergie conditionnée par la marque. Les deux opèrent sur le même registre, et la probabilité de la marque reste $b_*$. La référence demeure fixe lorsque l'essai varie.

\begin{theorem}[Conjugaison thermométrique du lecteur marqué]
Soit $H$ un Hamiltonien autoadjoint fixe sur un espace fini et soit $\rho_\beta=e^{-\beta H}/Z(\beta)$, avec $\beta>0$ et $\operatorname{Var}_{\rho_\beta}(H)>0$. Pour le couple \eqref{md:r27:par-lector}, la température conjuguée satisfait
\begin{equation}
 \Theta_P:=\left(\frac{d\mathcal S_P}{d\mathcal E_P}\right)^{-1},
 \qquad b_*\Theta_P=\beta^{-1}.
 \label{md:r27:transductor}
\end{equation}
\end{theorem}
\begin{proof}
Écrivons $E(\beta)=\operatorname{Tr}(\rho_\beta H)$. La dérivation de la somme spectrale finie donne $Z'/Z=-E$ et $E'=-\operatorname{Var}_{\rho_\beta}(H)$. Comme $s(\rho_\beta)=\beta E+\ln Z$, on a $s'=\beta E'$. La référence fixe permet de dériver $\mathcal S_P=b_*s$ sans terme supplémentaire, et la variance positive permet de diviser par $E'$. Ainsi, $d\mathcal S_P/d\mathcal E_P=b_*\beta$, ce qui prouve l'énoncé.
\end{proof}

La normalisation fait partie de la définition opératoire. Si l'information et l'énergie sont toutes deux lues conditionnellement, leurs valeurs sont $(s,E)$ et leur rapport différentiel est $\beta$. Si toutes deux sont lues avant le conditionnement, leurs valeurs sont $(b_*s,b_*E)$ et ce même rapport vaut $\beta$. Le couple \eqref{md:r27:par-lector} produit $b_*\beta$ parce qu'il conserve l'intensité dans la première opération et utilise l'énergie par essai sélectionné dans la seconde. Maintenir ce couple fixe son coefficient : le remplacer par $cb_*$, avec $c\ne1$, change l'équation. Si la référence variait, $d(b_*s)=b_*\,ds+s\,db_*$ ; ce protocole aurait une autre réponse et exige de conserver le second terme.

\subsection{Sélection variationnelle et bilan d'information}
\begin{theorem}[Principe variationnel dans la réalisation marquée]
Pour $\Theta>0$, la fonctionnelle
\begin{equation}
 \Phi_\Theta(\rho)=\mathcal E_P(\rho)-\Theta\mathcal S_P(\rho)
 =\operatorname{Tr}(\rho H)-b_*\Theta s(\rho)
 \label{md:r27:funcional}
\end{equation}
possède un unique minimum sur les matrices de densité, donné par
\begin{equation}
 \rho_\Theta^*=
 \frac{e^{-H/(b_*\Theta)}}{\operatorname{Tr}e^{-H/(b_*\Theta)}}.
 \label{md:r27:gibbs}
\end{equation}
De plus,
\begin{equation}
 \Phi_\Theta(\rho)-\Phi_\Theta(\rho_\Theta^*)
 =b_*\Theta D(\rho\Vert\rho_\Theta^*)\ge0.
 \label{md:r27:variacional}
\end{equation}
\end{theorem}
\begin{proof}
L'état \eqref{md:r27:gibbs} est fidèle en dimension finie. Substituer $\ln\rho_\Theta^*=-H/(b_*\Theta)-\ln Z_\Theta$ dans $D(\rho\Vert\rho_\Theta^*)=\operatorname{Tr}\rho(\ln\rho-\ln\rho_\Theta^*)$ donne l'égalité. Pour vérifier la positivité, on diagonalise $\rho$ et l'état fidèle $\sigma=\rho_\Theta^*$, avec les valeurs propres $r_i,s_j$ et la matrice bistochastique de recouvrements $p_{ij}$. La concavité de $\ln$ donne $D(\rho\Vert\sigma)\ge\sum_i r_i\ln(r_i/t_i)$, où $t_i=\sum_jp_{ij}s_j>0$ et $\sum_i t_i=1$. L'inégalité $-\ln u\ge1-u$ démontre que cette dernière somme est non négative. L'égalité exige $r_i=t_i$ pour tous les indices et que chaque vecteur propre de $\rho$ appartienne à un espace propre de $\sigma$ de valeur $t_i$ ; elle exige donc $\rho=\sigma$. Cela prouve le minimum et son unicité. La convention sur le noyau inclut les états singuliers ; la fidélité de $\sigma$ empêche l'un d'eux d'atteindre l'égalité.
\end{proof}
Dans cette réalisation, le même coefficient intervient dans la conjugaison différentielle et dans l'équilibre variationnel. La généralisation aux systèmes infinis conserve séparément les conditions de trace, de différentiabilité et de finitude de l'énergie et de l'entropie ; les théorèmes précédents ont un domaine fini.

\subsection{Composition avec l'action, l'aire et la réciprocité gravitationnelle}
\label{md:r27:composicion}
On utilise les sections orientées déjà construites, dans une même coordinatisation :
\begin{equation}
 \begin{gathered}
 \epsilon_a=\frac{h_a}{108t_0}=\frac{\hbar_a c_{\rm int}}{L_\alpha},
 \qquad q_A=\gamma a_0L_\alpha^2,\\
 \mathcal T_{I/A}=\frac{\epsilon_a}{q_A},\qquad
 G_a=\frac{c_{\rm int}^3L_\alpha^2}{\hbar_a}.
 \end{gathered}
 \label{md:r27:secciones}
\end{equation}
L'égalité longitudinale conserve l'horloge de \eqref{md:rigidez:reloj} ; $q_A>0$ fixe l'orientation aréale utilisée. La fonctionnelle $\gamma$ est celle de \eqref{md:compos:funcional-barbero} et maintient les incidences hexade--octade et les canaux $90/120$. Sa présence dans l'aire conserve cette dépendance complète. La section~\ref{md:rigidez:area} construit $a_0$ comme caractère tritique angulaire normalisé et démontre $\mathcal A=q_A(N_{90}+N_{120})$ sur trente-six liens étiquetés. Ainsi, $q_A$ est l'incrément d'aire par unité d'occupation dans cette section, et non un facteur déduit d'une entropie préfixée.

\begin{corollary}[Compatibilité des lecteurs thermique, aréal et gravitationnel]
Avec la référence tritique $\beta_{\rm clk}=\ln3/\epsilon_a$ et le même couple \eqref{md:r27:par-lector}, on a
\begin{equation}
 \begin{gathered}
 b_*\Theta_{{\rm clk},P}
 =\frac{\epsilon_a}{\ln3}
 =\frac{\mathcal T_{I/A}\gamma a_0L_\alpha^2}{\ln3},\\
 G_a b_*\Theta_{{\rm clk},P}\ln3=c_{\rm int}^4L_\alpha.
 \end{gathered}
 \label{md:r27:area-gravedad}
\end{equation}
Pour la section d'horizon $\tau_R=\hbar_a c_{\rm int}/(2\pi R)$,
\begin{equation}
 b_*T_{R,P}=\tau_R,\qquad
 \frac{T_{R,P}}{\Theta_{{\rm clk},P}}
 =\frac{\ln3}{2\pi}\frac{L_\alpha}{R}.
 \label{md:r27:horizonte}
\end{equation}
\end{corollary}
\begin{proof}
Appliquer \eqref{md:r27:transductor} à l'horloge donne la première égalité. Substituer $q_A$ dans $\mathcal T_{I/A}q_A=\epsilon_a$ donne la seconde. Multiplier par $G_a$ et utiliser \eqref{md:r27:secciones} donne $G_a\epsilon_a=c_{\rm int}^4L_\alpha$. Pour l'horizon, on applique le même théorème avec $\beta_R=1/\tau_R$ et l'on divise par l'égalité de l'horloge.
\end{proof}
Les conversions entre sections $i,j$ conservent $\Theta_{i,P}/\Theta_{j,P}=\tau_i/\tau_j$. Tout cycle de conversions se ferme par annulation télescopique. La température de l'horizon et celle de l'horloge ne coïncident que lorsque leurs sections énergétiques coïncident.

Dans une coordinatisation interne, $b_*$ est le coefficient du transducteur défini par \eqref{md:r27:par-lector}. L'unité thermique est spécifiée par la conjugaison entre ces lectures et la section énergétique choisie. La comparaison avec l'application appelée Boltzmann dans un autre protocole transporte également ses lecteurs, ses droites et ses bases ; le chiffre seul conserve la portée scalaire. En particulier, cette extension introduit une réalisation thermométrique explicite et ses preuves. Sa règle est incorporée comme telle, sans l'attribuer rétrospectivement à tout opérateur thermique antérieur ni identifier par leur nom son unité interne et une unité métrologique externe.

% Fragmento ES para inserción, no documento autónomo.
% Destino: VII/manuscrito/main.tex, después de relaciones_termica.tex.
% La sección recibe las salidas y pruebas de normalizacion_conjunta.tex,
% composicion.tex y relaciones_termica.tex; no altera sus generadores.
\section{Composition d’horizon et restitution thermo-géométrique}
\label{ins29:vii:horizonte}

La composition des constantes permet de distinguer ce que conserve chaque
lecture physique du même état engendré. Nous recevons les sections
d’action et de longueur, la paire angulaire et la vitesse constitutive déjà
construites à partir d’APP--TRIT--TPK. Nous abrégeons \(\eta=\eta_{\rm ret}\)
et notons \(s_0\) la section d’action dont l’expression dans la mise en coordonnées
SI est \(10^{-34}\,\mathrm{J\,s}\). Ainsi,
\[
 \hbar=\eta s_0,\qquad
 L_\alpha=r_N\alpha^{16}L_*,\quad r_N=\frac{5759}{23040},\qquad
 G=\frac{c^3L_\alpha^2}{\eta s_0}.
\]
Le retour inscrit détermine \(r_N\) ; le normand local détermine la
puissance \(16\) ; la section \(L_*\) conserve la mise en coordonnées longitudinale.
La fonction \(\eta=H_5-90\mathcal C_\pi/\pi\), avec ses coefficients et ses
dépendances régionales, est celle obtenue dans
\eqref{md:compos:seccion-accion}. Les substitutions suivantes reçoivent ces
sorties et conservent les unités et les domaines de leurs lecteurs.

\subsection{Aire fixe, masse fixe et conservation du produit thermique}

Pour la confrontation d’horizon d’Einstein, l’expression de
Bekenstein--Hawking est
\(S_{\rm BH}=k_Bc^3\mathcal A/(4G\hbar)\).
Le coefficient \(1/4\) appartient ici à cette loi de confrontation.
Dans un horizon de Schwarzschild, stationnaire, sans charge ni rotation,
\(R_S=2GM/c^2\), \(\mathcal A=4\pi R_S^2\) et
\(k_BT_H=\hbar c/(4\pi R_S)\).

\begin{proposition}[Deux contraintes de la composition d’horizon]
\label{ins29:vii:area_masa}
Dans la réalisation indiquée, les sorties structurelles donnent
\begin{align*}
 S_{\mathcal A}&=\frac{k_B\mathcal A}{4L_\alpha^2},&
 T_{\mathcal A}&=\frac{\eta s_0c}{2k_B\sqrt{\pi\mathcal A}},\\
 S_M&=\frac{4\pi k_Bc^2L_\alpha^2M^2}{\eta^2s_0^2},&
 T_M&=\frac{\eta^2s_0^2}{8\pi k_BL_\alpha^2M},
 \qquad S_MT_M=\frac{Mc^2}{2}.
\end{align*}
Dans la comparaison entre sections qui conserve \(c,L_\alpha,k_B\),
la lecture à aire fixe maintient \(S_{\mathcal A}\) et rend
\(T_{\mathcal A}\) proportionnel à \(\eta\) ; la lecture à masse fixe
donne \(S_M\propto\eta^{-2}\) et \(T_M\propto\eta^2\).
\end{proposition}
\begin{proof}
L’identité \(G\hbar=c^3L_\alpha^2\) donne immédiatement la première
entropie. Substituer \(R_S=\sqrt{\mathcal A/(4\pi)}\) dans la température
donne la première température. À masse fixe,
\(R_S=2cL_\alpha^2M/(\eta s_0)\) ; son aire et son inverse produisent les
deux autres formules. Les multiplier annule les facteurs communs et laisse
\(Mc^2/2\).
\end{proof}

L’annulation de l’action dans l’entropie à aire fixe exprime une
compatibilité entre ses lectures gravitationnelle et longitudinale. À masse
fixe, la même section d’action intervient dans la géométrie de l’horizon,
et sa dépendance demeure dans l’entropie. Le bilan doit conserver la
contrainte choisie tout au long de la comparaison.
Le paramètre radial \(R\) de la réalisation précédente
\(\hbar c/(2\pi R)\) coïncide dans cette confrontation avec \(2R_S\) :
les deux expressions représentent alors la même température. Cette
identification conserve le facteur deux entre rayon géométrique et échelle
d’accélération.

\subsection{Séparation angulaire et restitution de la section d’action}

La réalisation modulaire de frontière conserve dix-huit occupations
tritiques dans chaque secteur d’incidence. Nous écrivons \(X=N_{90}\),
\(Y=N_{120}\), \(N=X+Y\) et \(M=X-Y\). Leur séparation déclarée est
\[
 \Delta_{\rm mod}=
 \frac{(A_{\deg}-C^*_{\deg})\pi/180}{9\cdot90}
 \left(1-\frac7{12}\frac\pi{729}\right),\qquad
 d=30\Delta_{\rm mod},\quad f_\pi=1-\frac{7\pi}{8748}.
\]
La coordonnée \(d\) conserve tant la conversion angulaire \(\pi/180\)
que la différence d’incidence \(120-90=30\).
Dans cette réalisation, la distribution sur les configurations microscopiques
est proportionnelle à \(\exp[-d(3X+4Y)]\).
Nous recevons \(A_{\deg}=1000\alpha\),
\(C^*_{\deg}=2(\eta+\alpha)\) et posons
\(\xi=C^*_{\deg}/A_{\deg}=\tanh s\), dans la branche
\(1/500<\xi<1\) d’action positive. La rapidité de forme \(s\) est
adimensionnelle ; les deux coordonnées du quotient angulaire sont exprimées
en degrés.

\begin{proposition}[Courbe de compatibilité et lecteur de restitution]
\label{ins29:vii:curva}
Avec \(d_*=50\pi f_\pi\alpha/243\), on a
\[
 d=d_*(1-\xi)=d_*(1-\tanh s),\qquad
 0<d<\frac{499}{500}d_*.
\]
La distribution précédente satisfait, pour \(m=1,\ldots,36\),
\[
 \frac{\Pr(M=m)}{\Pr(M=-m)}=e^{dm}.
\]
Par conséquent, chaque rapport de probabilités détermine le même \(d\)
et restitue les coordonnées
\[
 \xi=1-\frac d{d_*},\qquad
 \eta=499\alpha-\frac{2430d}{\pi f_\pi},\qquad
 C^*_{\deg}=1000\alpha\left(1-\frac d{d_*}\right).
\]
\end{proposition}
\begin{proof}
La substitution de la paire angulaire donne
\(d=\pi f_\pi(499\alpha-\eta)/2430\).
Comme \(\eta=500\alpha\xi-\alpha\), il en résulte la première égalité et,
par la branche déclarée, l’intervalle. Pour la seconde, chaque configuration
d’occupations \((X,Y)\) a une compagne \((Y,X)\) avec la même
multiplicité. Le quotient de leurs poids est
\(e^{-d(3X+4Y)+d(3Y+4X)}=e^{d(X-Y)}\).
Sommer sur \(X-Y=m\) préserve ce facteur. Appliquer le logarithme et
isoler \(\xi\) et \(\eta\) fournit les trois restitutions.
\end{proof}

Cette composition établit une vérification conjointe : les trente-six
rapports, divisés logarithmiquement par leur \(m\) respectif, doivent donner
la même coordonnée. La restitution agit sur des sorties déjà engendrées et
maintient leur ordre causal. Son résultat relie le biais des occupations
à la forme elliptique et à la section d’action, sans fixer un rayon ni une
unité énergétique supplémentaire.

Les canaux de la fonctionnelle de Barbero sont ensuite restitués :
\[
 q_-=e^{-27d/f_\pi},\qquad
 q_+=D_A/q_-,\qquad D_A=e^{-100\pi\alpha/9}.
\]
En effet, \(27d/f_\pi=\pi(A_{\deg}-C^*_{\deg})/180\), tandis que
\(q_+q_-=D_A\). Substituer ces canaux et \(C^*_{\deg}(d)\) dans
\eqref{md:compos:funcional-barbero} détermine \(\gamma(d)\).
Le coefficient de l’aire conserve alors sa composition
\(\gamma(d)a_0(d)L_\alpha^2\). Lorsqu’il reçoit une valeur propre d’aire
\(\mathcal A_n=\gamma a_0L_\alpha^2n\), la confrontation d’horizon
donne \(S_{\rm BH}(\mathcal A_n)/k_B=\gamma a_0n/4\).
L’entropie de l’état réduit conserve, en outre, ses probabilités et
la bipartition. Les deux évaluations restent disponibles avec leur objet
respectif : géométrie de l’horizon et distribution d’intrication.

% Ampliación para insertar después de hibridacion.tex y del módulo térmico.
% La ecuación electrónica completa y las torres permanecen en sus secciones.
\section{Historique électronique et conservation de la masse sous transport}
\label{md:r27:electron}

L'équation \eqref{md:compos:electron-completo} évalue une section centrale persistante et conserve son retour d'action. La trajectoire, sa mémoire et sa masse sont des lectures différenciées du même état APP--TRIT--TPK. Le transport d'un parcours peut modifier les événements et la mémoire en conservant la valeur du lecteur massique. Pour le déterminer, on utilise le caractère complet de \eqref{md:hib:caracter}, ainsi que l'exposant de retour.

\subsection{Séparation et restitution des deux historiques centraux}
Les registres exprimés des deux orientations sont
\begin{equation}
 e_{++}=(2,2,5,-4,-3,-2)^2,\qquad
 e_{--}=(4,3,2,-2,-2,-5)^2,
 \label{md:r27:historias}
\end{equation}
où l'exposant $2$ indique la concaténation de deux blocs. Ils partagent un total nul et le même histogramme absolu. Pour un registre $e\in\mathbb Z^{12}$, soient
\begin{equation}
 \begin{gathered}
 p_i=e_i+e_{i+6},\quad \mathfrak a_i=e_i-e_{i+6},\\
 s_C=\sum_{i=0}^{5}p_i,\quad M_i=p_i-p_5\quad(0\le i<5).
 \end{gathered}
 \label{md:r27:registro-electron}
\end{equation}
\begin{proposition}[Restitution intégrale du registre]
La donnée $(s_C,M_0,\ldots,M_4,\mathfrak a_0,\ldots,\mathfrak a_5)$ détermine le registre par
\begin{equation}
 \begin{gathered}
 p_5=\frac{s_C-\sum_{i=0}^{4}M_i}{6},\quad p_i=p_5+M_i,\\
 e_i=\frac{p_i+\mathfrak a_i}{2},\quad e_{i+6}=\frac{p_i-\mathfrak a_i}{2}.
 \end{gathered}
 \label{md:r27:inversa-electron}
\end{equation}
Son image entière est caractérisée par la divisibilité par six du premier numérateur et la congruence $p_i\equiv \mathfrak a_i\pmod2$ pour chaque $i$.
\end{proposition}
\begin{proof}
Additionner $p_i=p_5+M_i$ pour les cinq premiers indices et ajouter $p_5$ donne $s_C=6p_5+\sum M_i$. Additionner et soustraire les deux définitions de $p_i,\mathfrak a_i$ donne les formules restantes. Les conditions indiquées sont nécessaires pour obtenir des entiers et suffisantes pour les reconstruire ; la substitution restitue toutes les données originales.
\end{proof}
Pour \eqref{md:r27:historias}, $\mathfrak a=0$ et $M_{++}=(8,8,14,-4,-2)$, tandis que $M_{--}=(18,16,14,6,6)$. La restitution distingue les deux historiques bien que certains lecteurs scalaires partagent leur valeur. Les symboles $\mathfrak a_i$ de cette décomposition sont des différences d'événements et se distinguent du caractère aréal $a_0$ par leur domaine.

\subsection{Forme spectrale du lecteur de retour}
Pour $\tau\in\mathbb R^{12}$, avec des indices cycliques, nous définissons
\begin{equation}
 C_s=\sum_{j=0}^{11}\tau_j\tau_{j+s},\quad
 T_k=\sum_{j=0}^{11}\tau_j e^{-2\pi i kj/12},\quad
 \theta_k=2\pi k/12.
 \label{md:r27:fourier}
\end{equation}
La représentation par caractères s'applique après la production du registre. Le lecteur de retour conserve $\Omega=-2C_1-4C_2-6C_3$ et $P_6=C_6/2$.
\begin{proposition}[Décomposition harmonique exacte]
\begin{equation}
 \begin{split}
 C_s&=\frac1{12}\sum_{k=0}^{11}|T_k|^2\cos(s\theta_k),\\
 \Omega&=\frac1{12}\sum_{k=0}^{11}|T_k|^2
 [-2\cos\theta_k-4\cos2\theta_k-6\cos3\theta_k],\\
 P_6&=\frac1{24}\sum_{k=0}^{11}|T_k|^2(-1)^k.
 \end{split}
 \label{md:r27:espectro}
\end{equation}
\end{proposition}
\begin{proof}
Développer $|T_k|^2$ produit $\sum_{j,l}\tau_j\tau_l e^{-i\theta_k(j-l)}$. La somme $\sum_{k=0}^{11}e^{i\theta_km}$ vaut douze si $12$ divise $m$ et zéro sinon. Appliquer cette identité au déplacement $s$ restitue $C_s$. Comme $\tau$ est réel, les parties imaginaires s'annulent par paires $k,12-k$ et il reste le cosinus. Les deux autres formules résultent d'une combinaison linéaire et de $\cos(6\theta_k)=(-1)^k$.
\end{proof}
Pour $\tau_e=(+++---)^2$, seules subsistent les puissances $|T_2|^2=64$, $|T_6|^2=16$, $|T_{10}|^2=64$. Leurs poids dans $\Omega$ sont $7,4,7$, de sorte que $\Omega=80$ et $P_6=6$. On restitue ainsi la partie harmonique de $\kappa_e=80-54\alpha+6\Delta_4$ dans l'équation électronique intégrale. Le registre alterné $(+-)^6$, de même histogramme, possède $|T_6|^2=144$ et $\Omega=48$ : l'ordre est une dépendance effective du lecteur. Les translations cycliques conservent $|T_k|^2$, même si elles modifient la phase de $T_k$ ; cette lecture spectrale préserve les corrélations, et non l'historique entier.

\subsection{Critère de masse transportée et lectures associées}
Soit $\mathcal H_{\mathcal F}$ la fibre d'un domaine fini de parcours et $\{|\Gamma\rangle\}$ sa base. On fixe l'un des lecteurs de retour $r=D$ ou $r=\Omega$ et une coordinatisation commune avec $R_{\rm act}>0$. La loi des masses de \eqref{md:hib:factor-ruta} détermine
\begin{equation}
 \begin{gathered}
 \frac{m_\Gamma^\beta}{m_e^\beta}
 =\chi_{\rm ref}(\sigma_\Gamma-\sigma_e,\eta_\Gamma-\eta_e)
 R_{\rm act}^{\kappa_\Gamma-\kappa_e},\\
 \Lambda_\Gamma=\ln(m_\Gamma^\beta/m_e^\beta).
 \end{gathered}
 \label{md:r27:masa-relativa}
\end{equation}
Le logarithme agit sur un rapport sans dimension. Dans une fibre non scalaire, on conserve l'opérateur et sa réalisation diagonale déclarée ; le domaine fini peut être une orbite ou une multisection, sans sélectionner une cellule arbitraire.

\begin{theorem}[Conservation du lecteur massique dans un domaine stable]
Soit $F$ un transport TPK agissant bijectivement dans le domaine fixé et $V|\Gamma\rangle=|F\Gamma\rangle$. Pour $M|\Gamma\rangle=m_e^\beta e^{\Lambda_\Gamma}|\Gamma\rangle$, on a
\begin{equation}
 [M,V]|\Gamma\rangle=m_e^\beta
 (e^{\Lambda_{F\Gamma}}-e^{\Lambda_\Gamma})|F\Gamma\rangle.
 \label{md:r27:conmutador-masa}
\end{equation}
Ainsi, $[M,V]=0$ équivaut à $\Lambda_{F\Gamma}=\Lambda_\Gamma$ pour tous les parcours. Les paramètres de la coordinatisation étant fixés, la condition est
\begin{equation}
 \begin{aligned}
 &\delta n_A x+\frac{\delta s_C}{6}y+
 \delta k_{\Delta_4}\Delta_4+\delta b_\zeta\zeta_{270}\\
 &\qquad+\sum_h\delta N_h(q_-^h-q_+^h)
 +\delta\kappa\ln R_{\rm act}=0,
 \end{aligned}
 \label{md:r27:criterio-masa}
\end{equation}
où $\delta$ compare le parcours transporté au parcours initial.
\end{theorem}
\begin{proof}
Appliquer $MV$ et $VM$ à chaque vecteur de la base donne \eqref{md:r27:conmutador-masa}. La positivité de $m_e^\beta$ et l'injectivité de l'exponentielle réelle donnent l'équivalence. Prendre le logarithme de \eqref{md:r27:masa-relativa} et soustraire les caractères des deux parcours donne \eqref{md:r27:criterio-masa}. Dans des tours infinies, on exige pour chaque parcours $\sum_h|N_h(q_-^h-q_+^h)|<\infty$ ; cette convergence absolue légitime la soustraction. La finitude du domaine des parcours conserve les opérateurs bornés utilisés dans cette preuve.
\end{proof}

Le transport est d'abord produit par TPK ; l'égalité détermine ensuite ce qu'il conserve. Les événements individuels et la mémoire peuvent changer et se compenser dans \eqref{md:r27:criterio-masa}. En maintenant la section d'action, l'horloge et la réalisation circulaire de \eqref{md:hib:reciprocidad-radial}, la commutation avec $M$ conserve également les lectures de fréquence propre, de rayon gravitationnel, de longueur de Compton réduite et de poids thermique, qui sont des fonctions du même opérateur positif. Pour l'inverse de la masse, on travaille hors du secteur nul. Dans la réalisation marquée de \eqref{md:r27:area-gravedad}, la quatrième lecture prend la forme
\begin{equation}
 \Xi_\Gamma=\frac{\Omega_\Gamma}{\omega_0}
 =\frac{r_{g,\Gamma}}{L_\alpha}
 =\frac{L_\alpha}{\bar\lambda_\Gamma}
 =\frac{E_\Gamma}{b_*\Theta_{{\rm clk},P}\ln3}.
 \label{md:r27:lecturas-conservadas}
\end{equation}
La preuve consiste à substituer les égalités déjà démontrées dans une même coordinatisation. Ce critère décrit la conservation pendant un transport déclaré ; les sélecteurs physiques d'espèce, de représentation et d'interaction conservent leurs propres conditions. La trajectoire et sa lecture harmonique ajoutent de l'information à la masse scalaire sans remplacer ces conditions.

\subsection{Transport de la section électronique vers le secteur électrofaible}
\label{amp-ew-seccion}

La section centrale de \eqref{md:hib:electron} conserve conjointement le
préfacteur, le correcteur et la mémoire de retour. Ses dépendances proviennent
des feuillets APP, de l’orientation TRIT et de la route enregistrée par TPK ;
le caractère angulaire et les tours agissent ensuite sur les multisections
de la même structure discrète du continu. Nous écrivons
\begin{equation}
 r_e=R_{\rm act}^{\kappa_e},\qquad
 E_e^\beta=\mathcal U_E\varepsilon_e^{(0)}r_e,\qquad
 \kappa_e=80-54\alpha+6\Delta_4.
 \label{amp-ew-electron}
\end{equation}
Ici $\mathcal U_E$ est la section de la droite énergétique et $r_e$ abrège
le retour électronique. La décade d’action reste dans sa droite propre ;
la coordonnée énergétique réalisée est transportée une seule fois.

Les signatures angulaires de \eqref{amp-higgs-firmas} produisent, dans la même
coordinatisation, $E_s^\angle=E_e^\beta\exp(n_sx+s_sy/6)$. En particulier,
\begin{equation}
 E_W^\angle=\mathcal U_E\varepsilon_e^{(0)}r_e
                  e^{94x-y/6},\qquad D=e^{-2x}.
 \label{amp-ew-w}
\end{equation}
Cette réalisation angulaire de W est utilisée. Les raffinements d’une autre tour
conservent leurs facteurs propres lorsqu’on compose avec eux.

\begin{proposition}[Échelle électronique et couplage de Fermi au niveau arbre]
Dans la réalisation électrofaible au niveau arbre, avec les coordonnées énergétiques
exprimées dans une même unité et $g^2=4\pi\alpha/(1-D)$, les relations
$v=2E_W^\angle/g$ et $G_F^{(0)}=1/(\sqrt2v^2)$ donnent
\begin{equation}
 G_F^{(0)}=
 \frac{\pi\alpha\,e^{-188x+y/3}}
 {\sqrt2(1-D)\mathcal U_E^2(\varepsilon_e^{(0)})^2r_e^2}.
 \label{amp-ew-fermi}
\end{equation}
Les signatures, les canaux et la coordinatisation énergétique étant fixes, le passage de
l’étape basale à l’étape après retour satisfait
\begin{equation}
 \frac{(G_F^{(0)})_\beta}{(G_F^{(0)})_0}=r_e^{-2},\qquad
 \frac{v_\beta}{v_0}=r_e.
 \label{amp-ew-escala}
\end{equation}
Les quotients $E_W^\angle/E_Z^\angle$, $E_H^\angle/E_W^\angle$ et
$\lambda_H=(g^2/8)(E_H^\angle/E_W^\angle)^2$ sont indépendants de
ce facteur commun $r_e$.
\end{proposition}
\begin{proof}
La substitution de $v$ donne $G_F^{(0)}=g^2/[4\sqrt2(E_W^\angle)^2]$.
Appliquer l’expression de $g^2$ et élever \eqref{amp-ew-w} au carré
produit \eqref{amp-ew-fermi} : les coefficients $188$ et $1/3$ résultent
respectivement de $2\cdot94$ et $2/6$. Remplacer $E_e^0$ par
$E_e^\beta=r_eE_e^0$ multiplie toutes les énergies angulaires par $r_e$.
La définition de $v$ conserve la première puissance ; celle de $G_F^{(0)}$
conserve l’inverse de son carré. Dans chaque quotient d’énergies, le
même facteur positif se simplifie, et la formule de $\lambda_H$ conserve
cette simplification. Les couplages $g,g'$ dépendent de $\alpha,x$ et
restent fixes dans cette comparaison.
\end{proof}

La grandeur $G_F^{(0)}$ a la dimension de l’inverse du carré d’une énergie.
La réalisation au niveau arbre déclarée spécifie la portée physique de cette
composition ; les corrections radiatives ont leurs opérateurs et leur ordre
postérieurs. Comparer les deux étapes du lecteur conserve leurs coordinatisations et
leurs signatures, au lieu de postuler une variation temporelle indépendante des
constantes corrélées.

\subsection{Conservation des raffinements de route}
\label{amp-ew-refinamiento}
La composition avec la loi complète préserve le caractère et la tour qui
correspondent à chaque route. Pour deux sections propres positives comparables,
dans une même coordinatisation et avec un même lecteur $r\in\{D,\Omega\}$, soit
\begin{equation}
 \rho_{W/e}^{\,r}
 =\frac{\chi_{\rm full}^{\,r}(\sigma_W,\eta_W)}
        {\chi_{\rm full}^{\,r}(\sigma_e,\eta_e)}
 R_{\rm act}^{\kappa_r(\Gamma_W)-\kappa_e}.
 \label{amp-ew-razon-completa}
\end{equation}
L’énergie correspondante est $E_W=E_e^\beta\rho_{W/e}^{\,r}$.
Dans la réalisation de Fermi compatible avec cette section, on obtient
\begin{equation}
 G_F^{(0)}=
 \frac{\pi\alpha}
 {\sqrt2(1-D)(E_e^\beta)^2(\rho_{W/e}^{\,r})^2}.
 \label{amp-ew-fermi-completo}
\end{equation}
La preuve substitue le rapport complet dans l’identité quadratique
précédente. Dans la coordinatisation angulaire explicite, ce rapport prend la forme de
\eqref{amp-ew-w} ; dans une réalisation raffinée, il retient également les
différences de tours et de lecteurs de \eqref{md:hib:cocientes}.
Si $E_s=E_e^\beta\chi_sf_s$, alors
$E_s/E_t=(\chi_s/\chi_t)(f_s/f_t)$. Cette égalité montre exactement
quel facteur se simplifie et quelle information spécifique demeure.

\subsection{Congruence massique et inversion spectrale sans commutation}
\label{amp-ew-congruencia}
Soit $\widehat M^0>0$ l’opérateur basal sur une fibre finie positive,
avec son caractère et ses tours incorporés, et soit $\widehat B^r$
le lecteur autoadjoint de retour. Le transport d’action produit
$C=R_{\rm act}^{\widehat B^r/2}>0$ et la loi
$\widehat M^\beta=C\widehat M^0C$.

\begin{proposition}[Réciprocité opératoire de la longueur d’action]
Avec $\hbar_a,c>0$ dans une même coordinatisation dimensionnelle,
\begin{equation}
 \begin{aligned}
 (\widehat M^\beta)^{-1}
    &=C^{-1}(\widehat M^0)^{-1}C^{-1},\\
 \widehat L_C^\beta:=\frac{\hbar_a}{c}(\widehat M^\beta)^{-1}
    &=C^{-1}\widehat L_C^0C^{-1},\qquad
 \widehat L_C^0=\frac{\hbar_a}{c}(\widehat M^0)^{-1}.
 \end{aligned}
 \label{amp-ew-inversa-operador}
\end{equation}
Ces identités admettent $[C,\widehat M^0]\ne0$.
\end{proposition}
\begin{proof}
La positivité de $C$ implique l’inversibilité. Pour $u\ne0$,
$\langle u,C\widehat M^0Cu\rangle
=\langle Cu,\widehat M^0Cu\rangle>0$ ; par conséquent, l’opérateur final
est également inversible. La multiplication directe dans les deux ordres donne
\[
 (C\widehat M^0C)\bigl(C^{-1}(\widehat M^0)^{-1}C^{-1}\bigr)=I,
 \qquad
 \bigl(C^{-1}(\widehat M^0)^{-1}C^{-1}\bigr)(C\widehat M^0C)=I.
\]
Multiplier par $\hbar_a/c$ prouve la seconde identité. Dans la
décomposition spectrale de l’opérateur final,
$\widehat M^\beta=\sum_jm_jP_j$ avec $m_j>0$, on obtient
$\widehat L_C^\beta=\sum_j\hbar_a/(cm_j)P_j$. La fonction inverse
est injective sur les valeurs propres positives : elle conserve leurs projecteurs
et leurs multiplicités finaux.
\end{proof}

La congruence peut changer les projecteurs de l’opérateur basal ;
l’inversion spectrale conserve ceux de l’opérateur final. Le secteur nul
reste en dehors de cette inversion. Une préparation qui occupe plusieurs
espaces propres conserve sa mesure spectrale et ses multiplicités.
En rang un, l’identité retrouve la longueur électronique
$\bar\lambda_{C,e}=\hbar_ac/E_e^\beta$ et, ensuite,
$a_{{\rm B},e}=\bar\lambda_{C,e}/\alpha$. Le résultat relie
l’échelle électronique, la réponse électrofaible et la longueur d’action
par des opérations explicites sur leurs structures partagées.

% Procedencia: COMPATIBILIDAD_HOLONOMIA_E_INFORMACION.md, secciones 2--9.
% Desarrollo científico conservado; gestión y correspondencia en archivo separado.
\section{Compatibilité constitutive, holonomie relative et information opérationnelle}
\label{md:comp:capitulo}

\subsection{Réciprocité différentielle de la vitesse et de l’impédance}
\label{md:comp:reciprocidad}

L’antécédent constitutif de la section~\ref{iii:sec:constitutiva} construit, sur les feuillets
conservés de sa représentation,
\begin{equation}
 \begin{gathered}
 T=e^{-xI-y\mathcal R},\qquad
 r_\pm=f\bigl(e^{-30(x\pm y)}\bigr),\\
 V=(I-T^{90})(I-T^{120})^{-1}=r_+P_++r_-P_-.
 \end{gathered}
 \label{md:comp:constitutiva}
\end{equation}
où
\[
 f(s)=\frac{1-s^3}{1-s^4},\qquad x>|y|.
\]
Les projecteurs de feuillet restent fixés. La chambre physique marquée
\(x>y>0\) et sa conjuguée sont contenues dans cette collection analytique.
Nous définissons les coordonnées logarithmiques normalisées
\begin{equation}
 c_\ell=\log\widehat c=-\log r_+-\log r_-,\qquad
 z_\ell=\log\widehat Z=\log r_--\log r_+.
 \label{md:comp:logaritmos}
\end{equation}
Ces lettres ne désignent ni de nouvelles constantes ni le contre-angle \(C^*\).
Soient
\[
 u=30(x+y),\qquad v=30(x-y),\qquad
 g(w)=\log f(e^{-w}),\qquad a=g'(u),\quad b=g'(v).
\]

\begin{proposition}[Réciprocité constitutive différentielle]
\label{md:comp:prop-reciprocidad}
Dans le domaine \(x>|y|\), on a
\begin{equation}
 g'(w)=\frac{3}{e^{3w}-1}-\frac{4}{e^{4w}-1}>0,
 \qquad
 D(c_\ell,z_\ell)=-30
 \begin{pmatrix}a+b&a-b\\a-b&a+b\end{pmatrix}.
 \label{md:comp:jacobiano}
\end{equation}
Les valeurs propres du jacobien sont \(-60a\) et \(-60b\). Il est donc
symétrique et défini négatif, son déterminant est \(3600ab>0\), et
\begin{equation}
 \boxed{\partial_y\log\widehat c=\partial_x\log\widehat Z.}
 \label{md:comp:reciprocidad-cruzada}
\end{equation}
\end{proposition}
\begin{proof}
On dérive les identités
\(c_\ell=-g(u)-g(v)\) et \(z_\ell=g(v)-g(u)\).
La positivité de \(g'\) provient de la décroissance stricte, sur \(k>0\),
de \(k/(e^{kw}-1)\). Elle résulte également de l’expression rationnelle positive
de la section~\ref{md:comp:restitucion}. La matrice obtenue possède les
valeurs propres et le déterminant indiqués ; l’égalité de ses entrées
non diagonales donne \eqref{md:comp:reciprocidad-cruzada}.
\end{proof}

L’égalité relie la réponse de la vitesse à la séparation
orientée et la réponse de l’impédance à la coordonnée commune.
Les deux sensibilités proviennent d’une même fonction spectrale. C’est une
réciprocité de cette collection constitutive ; elle n’est pas identifiée sans preuve
à une loi thermodynamique ou d’électromagnétisme macroscopique.

\subsection{Potentiel spectral et connexion de profondeur deux avec Catalan}
\label{md:comp:potencial}

\subsubsection{Construction du potentiel}
\label{md:comp:construccion-potencial}

À partir des canaux déjà construits, nous définissons
\begin{equation}
 \mathscr F(x,y)=\sum_{\sigma=\pm1}\sum_{n\geq1}
 \left[
 \frac{e^{-120n(x+\sigma y)}}{120n^2}
 -\frac{e^{-90n(x+\sigma y)}}{90n^2}
 \right].
 \label{md:comp:serie-potencial}
\end{equation}

\begin{proposition}[Potentiel des deux réponses logarithmiques]
\label{md:comp:prop-potencial}
La série \eqref{md:comp:serie-potencial} définit sur \(x>|y|\) un potentiel
qui satisfait
\begin{equation}
 \boxed{\partial_x\mathscr F=c_\ell,\qquad
        \partial_y\mathscr F=z_\ell.}
 \label{md:comp:gradiente-potencial}
\end{equation}
Sa hessienne est le jacobien de
\eqref{md:comp:jacobiano} ; par conséquent, \(\mathscr F\) est
strictement concave. La constante additive est fixée par
\(\mathscr F\longrightarrow0\) lorsque \(x-|y|\longrightarrow\infty\).
\end{proposition}
\begin{proof}
Sur tout compact du domaine \(x>|y|\), les quantités \(x\pm y\)
ont un minimum positif. Une puissance de \(n\) multipliée par une
exponentielle décroissante est sommable ; ainsi, la série et ses dérivées
de tout ordre fixé convergent uniformément sur ce compact.
En dérivant terme à terme et en utilisant
\(-\log(1-z)=\sum_{n\geq1}z^n/n\), on obtient les deux identités
de \eqref{md:comp:gradiente-potencial}. La hessienne coïncide alors
avec \eqref{md:comp:jacobiano}, qui est défini négatif.
La normalisation à l’infini est celle de la série présentée.
\end{proof}

\(\mathscr F\) est un potentiel mathématique de la réponse. On ne lui
attribue pas d’unités d’énergie et on ne le déclare pas potentiel thermodynamique.

\subsubsection{Réalisation par un même moment spectral}
\label{md:comp:momento-espectral}

Soit
\[
 D_2(z)=\sum_{n\geq1}\frac{z^n}{n^2}.
\]
Sa reconnaissance comme dilogarithme est postérieure à cette définition par
série. Pour \(|z|<1\), il satisfait
\((z\partial_z)^2D_2(z)=z/(1-z)\). Dans cette notation,
\begin{equation}
 \mathscr F=
 \sum_{\sigma\in\{+,-\}}
 \left[\frac{D_2(q_\sigma^{120})}{120}
       -\frac{D_2(q_\sigma^{90})}{90}\right].
 \label{md:comp:potencial-dos-momentos}
\end{equation}
Le moment orienté de Catalan construit dans la section~\ref{iii:sec:catalan} satisfait
\begin{equation}
 \mathcal C_2(s)=\operatorname{Im}D_2(is)
 =\sum_{k\geq0}\frac{(-1)^ks^{2k+1}}{(2k+1)^2},\qquad
 \mathcal C_2(1)=G_{\rm Cat}.
 \label{md:comp:catalan-momento}
\end{equation}

\begin{proof}
Les puissances paires de \(i\) sont réelles et les impaires ont une partie
imaginaire \((-1)^k\). Cela donne la série de
\eqref{md:comp:catalan-momento}, qui converge absolument même
en \(s=1\). Pour \(0\leq s<1\), appliquer terme à terme deux fois
\(s\partial_s\) redonne \(s/(1+s^2)\), le coefficient orienté de la
résolvante du quart de tour déjà présente dans le corpus.
L’identité dérivée en \(s=1\), si elle est utilisée, s’interprète
par prolongement analytique ou limite d’Abel : la série deux fois
dérivée ne converge pas au sens ordinaire à cette extrémité.
L’identification \(\mathcal C_2(1)=G_{\rm Cat}\) provient directement
de la série originale absolument convergente.
\end{proof}

Par conséquent, le moment de Catalan et le potentiel constitutif sont réunis
par un même moment spectral de profondeur deux, évalué sur des
supports différents : le quart de tour orienté et les contractions
des secteurs \(90/120\). Cette composition n’identifie pas Catalan
à \(\mathscr F\), ni n’affirme que sa seule valeur extrême détermine
tous les canaux. La collection analytique et ses arguments conservent
une information qu’une valeur isolée ne contient pas.

\subsubsection{Compatibilité avec l’amplification}
\label{md:comp:amplificacion-potencial}

Pour \(T_m=T\otimes I_{9^m}\) et
\(\tau_m=\operatorname{Tr}/(2\cdot9^m)\), on a exactement
\begin{equation}
 \mathscr F_m=2\tau_m\left[
 \frac{D_2(T_m^{120})}{120}
 -\frac{D_2(T_m^{90})}{90}\right]=\mathscr F.
 \label{md:comp:potencial-amplificado}
\end{equation}
La multiplicité \(9^m\) annule la normalisation. Le même argument
conserve les dérivées. Cela atteste l’amplification de la
représentation, non tout raffinement imaginable du TPK.
L’opérateur temporel n’est pas remplacé par une liste de valeurs propres
sans marques de feuillet.

Pour tout lacet lisse fermé dans cette collection analytique,
\begin{equation}
 \oint(c_\ell\,dx+z_\ell\,dy)=0.
 \label{md:comp:integral-cerrada}
\end{equation}
Cette égalité se rapporte à la projection constitutive à deux
coordonnées. Elle n’implique pas que l’histoire enrichie ou son holonomie
opératoire soient triviales.

\subsection{Restitution exacte et sensibilité finie}
\label{md:comp:restitucion}

\begin{proposition}[Restitution analytique des coordonnées angulaires]
\label{md:comp:prop-restitucion}
Avec \(L=\log(16/9)\), l’application
\((x,y)\longmapsto(c_\ell,z_\ell)\) est un difféomorphisme réel analytique
du domaine \(x>|y|\) sur
\begin{equation}
 0<c_\ell+z_\ell<L,\qquad 0<c_\ell-z_\ell<L.
 \label{md:comp:imagen-logaritmica}
\end{equation}
\end{proposition}
\begin{proof}
La fonction \(g\) croît strictement de \(\log(3/4)\) à \(0\).
Les combinaisons
\[
 c_\ell+z_\ell=-2g(u),\qquad c_\ell-z_\ell=-2g(v)
\]
récupèrent \(u,v\) au moyen de \(g^{-1}\). Alors
\(x=(u+v)/60\) et \(y=(u-v)/60\).
L’injectivité est maintenue lorsqu’on restreint à l’image effectivement
produite par HMT ; on ne déclare pas que HMT produise tout le cône.
\end{proof}

Avec \(s=e^{-w}\), la dérivée a la forme rationnelle
\begin{equation}
 a(w)=g'(w)=
 \frac{s^3(s^2+2s+3)}{(1+s+s^2)(1+s+s^2+s^3)},
 \qquad 0<a(w)<\frac12.
 \label{md:comp:derivada-racional}
\end{equation}
La fonction est strictement décroissante, avec \(a(0+)=1/2\), et
\(a(w)\sim3e^{-3w}\) à l’infini. Pour prouver la décroissance,
on utilise
\begin{equation}
 a'(w)=-\frac{9}{4\sinh^2(3w/2)}
        +\frac{16}{4\sinh^2(2w)}<0,
 \label{md:comp:decrecimiento}
\end{equation}
car \(k/\sinh(kw/2)\) décroît strictement sur \(k>0\).
En outre,
\begin{equation}
 \frac12e^{-3w}<a(w)<3e^{-3w}.
 \label{md:comp:cotas-exponenciales}
\end{equation}
La borne inférieure résulte de la multiplication des dénominateurs et de l’utilisation de
\begin{equation}
 \begin{split}
 &2(s^2+2s+3)-(1+s+s^2)(1+s+s^2+s^3)\\
 &\hspace{1cm}=(1-s)(5+7s+6s^2+3s^3+s^4)>0.
 \end{split}
 \label{md:comp:polinomio-cota}
\end{equation}
La borne supérieure résulte de
\[
 3(1+s+s^2)(1+s+s^2+s^3)-(s^2+2s+3)>0.
\]

Si \(m=30(x-|y|)\) et \(M=30(x+|y|)\), la norme de l’inverse du
jacobien est \(1/[60a(M)]\). Pour ne pas confondre l’application vectorielle avec
le potentiel scalaire, nous écrivons
\(\mathbf F=(c_\ell,z_\ell)\).

\begin{proposition}[Bornes finies de restitution]
\label{md:comp:prop-sensibilidad}
Sur une région convexe où les deux arguments \(u,v\) appartiennent
à \([m_0,M_0]\subset(0,\infty)\), on vérifie
\begin{equation}
 60a(M_0)\|p-q\|
 \leq\|\mathbf F(p)-\mathbf F(q)\|
 \leq60a(m_0)\|p-q\|.
 \label{md:comp:cotas-finitas}
\end{equation}
\end{proposition}
\begin{proof}
On intègre la hessienne de \(\mathscr F\), c’est-à-dire le jacobien
de \(\mathbf F\), sur le segment. Pour la borne inférieure, on utilise
\[
 -\langle\mathbf F(p)-\mathbf F(q),p-q\rangle
 \geq60a(M_0)\|p-q\|^2
\]
et l’inégalité de Cauchy--Schwarz. La borne supérieure utilise la
norme maximale du jacobien. Ce sont des bornes d’erreurs finies dans le
domaine spécifié.
\end{proof}

En \(y=0\) et lorsque \(x\longrightarrow\infty\), les deux directions
se contractent également : le conditionnement matriciel vaut un,
mais la norme de l’inverse croît comme \(e^{90x}/180\).
L’inversibilité exacte et la récupération robuste face à l’erreur absolue
sont donc des propriétés différentes. Lorsque \(|y|\longrightarrow x\)
avec \(x>0\) fixé, il n’y a pas de singularité différentielle : une dérivée tend
vers \(1/2\) et l’autre vers \(a(60x)>0\).

\subsection{Le quart de tour provient du cycle dodécaphasique}
\label{md:comp:cuarto-giro}

On retrouve la construction de \eqref{iii:eq:quarter-intertwiner}. Dans \(\mathbb R^{12}\),
soient
\[
 C_{12}e_j=e_{j+1\bmod12},\qquad
 \Pi_4=P_4^{\rm cyc}\frac{I-C_{12}^6}{2},
\]
et
\begin{equation}
 \begin{gathered}
 u_0=(1,0,-1,0,1,0,-1,0,1,0,-1,0)^{\mathsf T},\qquad
 v_0=C_{12}u_0,\\
 W=(u_0,v_0)/\sqrt6.
 \end{gathered}
 \label{md:comp:isometria-ciclica}
\end{equation}
Le symbole \(\Pi_4\) distingue le projecteur des coordonnées
d’action \(Q,P\). On vérifie
\begin{equation}
 \begin{gathered}
 W^{\mathsf T}W=I,\qquad WW^{\mathsf T}=\Pi_4,\qquad
 C_{12}W=WJ_0,\\
 J_0=\begin{pmatrix}0&-1\\1&0\end{pmatrix},\qquad J_0^2=-I.
 \end{gathered}
 \label{md:comp:entrelazador-ciclico}
\end{equation}
En particulier, \(C_{12}^3W=-WJ_0\). Le plan possède ainsi un opérateur d’entrelacement
explicite avec le cycle ; aucune rotation externe n’a été introduite pour
ajuster une réponse.

L’antécédent radial de \eqref{rec:eq:entrelazamiento} réalise
\[
 S_\zeta=\operatorname{diag}(e^{\zeta/2},e^{-\zeta/2}),
 \qquad J_\zeta=S_\zeta J_0S_\zeta^{-1}.
\]
Sur ce plan, nous définissons
\begin{equation}
 U_\zeta(\theta)=S_\zeta e^{-\theta J_0}S_\zeta^{-1}
 =\begin{pmatrix}
 \cos\theta&e^\zeta\sin\theta\\
 -e^{-\zeta}\sin\theta&\cos\theta
 \end{pmatrix}.
 \label{md:comp:transporte-plano}
\end{equation}
Le signe de \(\theta\) fixe l’orientation choisie. Le rapport engendré
\(\xi=C^*/A\), avec \(|\xi|<1\) et les deux angles exprimés dans la même
unité, satisfait \(\zeta=\operatorname{artanh}\xi\).
Dans l’orientation conjuguée,
\((\varepsilon,\xi)\mapsto(-\varepsilon,-\xi)\) conserve la section
positive d’action au moyen de l’opérateur de récupération orienté du corpus ;
\(\zeta\) change de signe.

\subsection{Lacet relatif de deux transports}
\label{md:comp:lazo}

Soient
\begin{equation}
 U=U_{\zeta_s}(\theta),\qquad
 V=U_{\zeta_r}(\vartheta),\qquad
 H_{\rm lazo}=UVU^{-1}V^{-1}.
 \label{md:comp:def-lazo}
\end{equation}
Les produits se lisent avec une action de droite à gauche.
La lettre \(\vartheta\) désigne ici la seconde phase, pour la distinguer
de la coordonnée régionale \(\varphi\). Posons
\[
 \delta=\zeta_s-\zeta_r,\qquad
 \chi=\sinh\delta\,\sin\theta\,\sin\vartheta.
\]

\begin{proposition}[Trace et spectre du lacet relatif]
\label{md:comp:prop-lazo}
La composition \eqref{md:comp:def-lazo} satisfait
\begin{equation}
 \boxed{\begin{gathered}
 \det H_{\rm lazo}=1,\qquad
 \operatorname{tr}H_{\rm lazo}=2+4\chi^2,\\
 H_{\rm lazo}=I\ \Longleftrightarrow\ \chi=0.
 \end{gathered}}
 \label{md:comp:traza-lazo}
\end{equation}
Ses valeurs propres sont
\begin{equation}
 \lambda_\pm=(\sqrt{1+\chi^2}\pm|\chi|)^2
 =\exp\bigl(\pm2\operatorname{arsinh}|\chi|\bigr).
 \label{md:comp:espectro-lazo}
\end{equation}
\end{proposition}
\begin{proof}
En conjuguant simultanément par \(S_{\zeta_r}^{-1}\), l’opérateur \(V\)
est \(R(-\vartheta)\), et \(U\) possède les termes
\(e^{\pm\delta}\sin\theta\). La multiplication directe donne, dans cette
référence,
\[
 UV-VU=-2\chi\operatorname{diag}(1,-1).
\]
Comme
\[
 H_{\rm lazo}-I=(UV-VU)U^{-1}V^{-1},
\]
on obtient l’équivalence avec \(\chi=0\). La multiplication, ou
l’identité de traces pour les matrices unimodulaires, donne
\(\operatorname{tr}H_{\rm lazo}=2+4\chi^2\).
Le polynôme caractéristique est
\(\lambda^2-(2+4\chi^2)\lambda+1\), et ses racines sont celles de
\eqref{md:comp:espectro-lazo}.
\end{proof}

C’est une holonomie relative de la composition déclarée.
Elle n’est identifiée ni à \(\Gamma_9\) ni à une mise à jour complète
du TPK par similitude de noms.

\subsubsection{Spécialisation entièrement déterminée par la paire conjuguée}
\label{md:comp:lazo-conjugado}

Soient \(U_+=U_\zeta(\pi/2)\) et
\(U_-=U_{-\zeta}(\pi/2)\). Les deux satisfont
\(U_+^2=U_-^2=-I\), mais
\begin{equation}
 \begin{aligned}
 U_+U_-&=-\operatorname{diag}(e^{2\zeta},e^{-2\zeta}),\\
 H_{\rm lazo}=(U_+U_-)^2
 &=\operatorname{diag}(e^{4\zeta},e^{-4\zeta}).
 \end{aligned}
 \label{md:comp:lazo-diagonal}
\end{equation}
En écrivant \(\xi=\tanh\zeta=C^*/A\), avec les deux angles exprimés
dans la même unité, on obtient
\begin{equation}
 \begin{aligned}
 H_{\rm lazo}
 &=\operatorname{diag}\left(
 \left[\frac{1+\xi}{1-\xi}\right]^2,
 \left[\frac{1-\xi}{1+\xi}\right]^2\right),\\
 \operatorname{tr}H_{\rm lazo}-2
 &=\frac{16\xi^2}{(1-\xi^2)^2}.
 \end{aligned}
 \label{md:comp:lazo-razon-angular}
\end{equation}
La formule de \(H_{\rm lazo}\) compose le quart de tour et les formes
conjuguées de la même paire. Son expression est exacte dans le rapport angulaire
engendré \(\xi\).

Sur la branche positive d’action du corpus,
\(A_{\deg}=1000\alpha\) et
\(C^*_{\deg}=2(\eta_{\rm ret}+\alpha)\), avec
\(\hbar_{\rm ret}=s_0\eta_{\rm ret}\). Par conséquent,
\begin{equation}
 \begin{aligned}
 \xi&=\frac{\eta_{\rm ret}+\alpha}{500\alpha},\\
 \rho_{\rm lazo}=e^{4\zeta}
 &=\left(\frac{501\alpha+\eta_{\rm ret}}
                 {499\alpha-\eta_{\rm ret}}\right)^2,
 \qquad 0<\eta_{\rm ret}<499\alpha.
 \end{aligned}
 \label{md:comp:multiplicador-lazo}
\end{equation}
Les coefficients \(501\) et \(499\) sont \(500\pm1\) ;
\(500\) provient du facteur \(1000/2\) de la paire angulaire et n’est pas un
nouveau paramètre. Les axes orientés étant conservés, la restitution est
\begin{equation}
 \boxed{
 \eta_{\rm ret}=\alpha\left[
 500\frac{\sqrt{\rho_{\rm lazo}}-1}
          {\sqrt{\rho_{\rm lazo}}+1}-1\right].}
 \label{md:comp:restitucion-accion}
\end{equation}
La dérivée
\[
 \frac{d\eta_{\rm ret}}{d\rho_{\rm lazo}}
 =\frac{500\alpha}
 {\sqrt{\rho_{\rm lazo}}(\sqrt{\rho_{\rm lazo}}+1)^2}
\]
est positive. La chambre \(\eta_{\rm ret}>0\) correspond à
\(\rho_{\rm lazo}>(501/499)^2\) ; le bord
\(\eta_{\rm ret}\longrightarrow499\alpha\) correspond à
\(\rho_{\rm lazo}\longrightarrow\infty\).
C’est une lecture inverse du coefficient adimensionnel d’action à partir
du transport relatif et de la coordonnée \(\alpha\) déjà construite.
Elle conserve l’échelle \(s_0\) comme antécédent : une transformation
unimodulaire de forme ne détermine pas à elle seule cette échelle dimensionnelle.
L’orientation conjuguée inverse \(\rho_{\rm lazo}\) et transporte
l’étiquette d’action ; elle ne transforme pas \(\eta_{\rm ret}\) en action négative.

\subsubsection{Itération du lacet et accumulation orientée}
\label{md:comp:iteracion}

Soit
\[
 \rho_{\rm lazo}=e^{4\zeta}
 =\left(\frac{1+\xi}{1-\xi}\right)^2.
\]
Pour tout entier \(n\),
\begin{equation}
 \begin{gathered}
 H_{\rm lazo}^n
 =\operatorname{diag}(\rho_{\rm lazo}^{n},\rho_{\rm lazo}^{-n}),
 \\
 Q_n=\rho_{\rm lazo}^{n}Q_0,\quad
 P_n=\rho_{\rm lazo}^{-n}P_0.
 \end{gathered}
 \label{md:comp:iteracion-diagonal}
\end{equation}
La preuve est la multiplication de matrices diagonales, étendue à
\(n<0\) au moyen de l’inverse. Le produit
\(Q_nP_n=Q_0P_0\) et l’aire symplectique sont conservés, tandis que l’état change
pour \(\zeta\ne0\) et \(X_0\ne0\). Dans le programme de quatre
opérations, les incréments de phase déclarés ont une somme nulle à chaque
répétition ; cette comptabilité des phases ne détermine pas le transport
total de l’état.

Pour \(\zeta\ne0\) et \(Q_0\ne0\), on récupère exactement
\[
 n=\frac{\log|Q_n/Q_0|}{4\zeta}.
\]
Si \(Q_0=0\), la condition \(P_0\ne0\) permet d’utiliser
\[
 n=-\frac{\log|P_n/P_0|}{4\zeta}.
\]
La lecture nécessite de connaître la préparation et de conserver une référence.
L’extension multiplicative
\(H_{\rm lazo}^{n+m}=H_{\rm lazo}^{n}H_{\rm lazo}^{m}\)
correspond à un incrément logarithmique additif \(4(n+m)\zeta\).
Inverser le lacet change le signe de cet incrément.

Ainsi, un programme fixe et répété admet une réponse qui enregistre
son itération. C’est une réalisation mathématique d’accumulation
opérationnelle dans la composition étudiée. On n’identifie pas \(n\)
à toute la retenue ou la mémoire de \(\Gamma_9\), et l’on ne déduit pas une
entropie nulle pour la dynamique complète à partir de la périodicité
de son programme de phases. Si l’amplitude initiale est inconnue et
si l’on ne dispose que de l’état final, \(n\) n’est pas déterminé par
cette lecture isolée.

\subsubsection{Relèvement explicite aux douze phases}
\label{md:comp:levantamiento}

Soient
\[
 \widetilde S_\zeta=I-\Pi_4+WS_\zeta W^{\mathsf T},\qquad
 L_\pm=\widetilde S_{\pm\zeta}C_{12}^3
       \widetilde S_{\pm\zeta}^{-1}.
\]
Sur \(\operatorname{Ran}\Pi_4\) agissent \(U_+,U_-\), respectivement ;
sur le complément, tous deux agissent comme \(C_{12}^3\). Par conséquent,
\begin{equation}
 L_+L_-L_+^{-1}L_-^{-1}
 =I-\Pi_4+WH_{\rm lazo}W^{\mathsf T}.
 \label{md:comp:levantamiento-lazo}
\end{equation}
L’action complémentaire s’annule. La trace de l’opérateur à douze
phases est \(10+\operatorname{tr}H_{\rm lazo}\).
L’entrelacement et la composition sont exacts et sont vérifiés
rationnellement à l’aide des vecteurs entiers \(u_0,v_0\).
Ce relèvement réunit l’antécédent dodécaphasique et le transport
elliptique ; il n’affirme pas que sa séquence soit à elle seule la mise à jour
autonome originale du TPK.

\subsubsection{Raffinement et limite}
\label{md:comp:limite}

L’amplification
\(H_{{\rm lazo},m}=H_{\rm lazo}\otimes I_{9^m}\) conserve
\[
 \frac{1}{9^m}\operatorname{Tr}H_{{\rm lazo},m}
 =\operatorname{Tr}H_{\rm lazo},
\]
ainsi que la norme et le spectre, hormis les multiplicités.
Avec les inclusions \(v\mapsto v\otimes e_0\), on a
\[
 H_{{\rm lazo},m+1}\iota_m=\iota_m H_{{\rm lazo},m}.
\]
L’opérateur uniforme définit dans la limite inductive hilbertienne une
extension bornée, de même norme. C’est la limite de cette
amplification ; il ne remplace pas un autre système de raffinements
du corpus.

\subsection{Ce qui distingue un changement de coordinatisation d’un effet relatif}
\label{md:comp:carta-efecto}

Le transport de coordonnées
\[
 T_n=S_{\zeta_{n+1}}R(-\theta_n)S_{\zeta_n}^{-1}
\]
se télescope :
\begin{equation}
 T_{N-1}\cdots T_0
 =S_{\zeta_N}R\left(-\sum_{n=0}^{N-1}\theta_n\right)
 S_{\zeta_0}^{-1}.
 \label{md:comp:telescopaje}
\end{equation}
Si la coordinatisation et la phase reviennent, on obtient \(I\). Le commutateur de la
section~\ref{md:comp:lazo} compose des évolutions de formes distinctes
dans une référence commune ; ce n’est pas ce changement passif.

Sous un changement commun de représentation \(P\), les opérateurs
\(U,V,H_{\rm lazo}\) sont conjugués par \(P\).
La trace et le spectre de \(H_{\rm lazo}\) restent invariants.
Si, en outre, \(X\mapsto PX\) et
\(g\mapsto P^{-\mathsf T}gP^{-1}\), on conserve
\(X^{\mathsf T}gX\) et toute lecture énergétique correctement
transportée. Ainsi, un \(H_{\rm lazo}\ne I\) ne disparaît pas par
renommage des coordonnées.

Cela constitue un protocole mathématique relationnel fermé.
Sa mise en œuvre expérimentale exige que les évolutions et leurs
inverses soient des opérations physiques disponibles et que le système,
la référence et le contrôleur possèdent une réalisation commune.
On n’affirme ici ni force nouvelle ni source d’énergie.

\subsubsection{Conservation de l’aire et conservation d’une même énergie}
\label{md:comp:area-energia}

\begin{proposition}[Absence d’une métrique positive commune]
\label{md:comp:prop-metrica}
Chaque quart de tour conserve sa propre métrique :
\(g_\zeta=\operatorname{diag}(e^{-\zeta},e^\zeta)\) pour \(U_+\),
et \(g_{-\zeta}\) pour \(U_-\). Lorsque \(\zeta\ne0\), il n’existe pas
de métrique positive commune.
\end{proposition}
\begin{proof}
Les conditions
\[
 G=\begin{pmatrix}a&b\\b&d\end{pmatrix}>0,
 \qquad U_+^{\mathsf T}GU_+=G
\]
imposent \(b=0\) et \(d=ae^{2\zeta}\).
La condition \(U_-^{\mathsf T}GU_-=G\) exige
\(d=ae^{-2\zeta}\). Toutes deux ne sont compatibles qu’avec \(\zeta=0\).
\end{proof}

Par conséquent, conserver l’aire symplectique n’équivaut pas à conserver
une même forme énergétique pendant toute composition de ces
transports. Dans la référence
\(g_r=S_r^{-\mathsf T}S_r^{-1}\), avec \(\omega_r>0\),
\(X\ne0\) et \(E_r(X)=\omega_r X^{\mathsf T}g_rX/2\), on a
\begin{equation}
 \begin{gathered}
 \frac{E_r(H_{\rm lazo}X)}{E_r(X)}
 =\frac{x^{\mathsf T}\overline{H}_{\rm lazo}^{\mathsf T}
                   \overline{H}_{\rm lazo}x}{x^{\mathsf T}x},\\
 \overline{H}_{\rm lazo}=S_r^{-1}H_{\rm lazo}S_r,\qquad
 x=S_r^{-1}X.
 \end{gathered}
 \label{md:comp:cociente-energia}
\end{equation}
Les extrêmes sont les carrés des valeurs singulières.
Dans le cas conjugué diagonal, ils sont \(e^{\pm8|\zeta|}\), tandis que
les dilatations de l’état sont \(e^{\pm4\zeta}\).
Cette différence d’exposants est nécessaire. Changer physiquement
d’opération exige de prendre en compte le contrôleur ; le gain d’une
lecture n’atteste pas une création d’énergie.

\subsubsection{L’ordre nécessite une lecture orientée}
\label{md:comp:orden-orientado}

Échanger \(U_+\) et \(U_-\) inverse \(H_{\rm lazo}\).
La trace et l’ensemble des gains extrêmes sont égaux pour
\(H_{\rm lazo}\) et \(H_{\rm lazo}^{-1}\). Ces scalaires détectent
l’effet du lacet, mais non son orientation. Deux entrées
indépendantes avec leurs réponses vectorielles complètes reconstruisent
\(H_{\rm lazo}\) ; une mesure scalaire ne suffit pas en général.

Dans le cas conjugué, avec des axes de référence marqués, soient
\[
 G_Q=\frac{E_r(H_{\rm lazo}e_Q)}{E_r(e_Q)},\qquad
 G_P=\frac{E_r(H_{\rm lazo}e_P)}{E_r(e_P)}.
\]
La lecture
\begin{equation}
 \mathcal O=\frac14\log(G_Q/G_P)=4\zeta
 \label{md:comp:lector-orientado}
\end{equation}
récupère l’orientation et change de signe lorsqu’on inverse le lacet.
Les axes et la métrique doivent être transportés ensemble sous un changement
de représentation. On distingue ainsi l’orientation sans l’attribuer
à une trace qui l’élimine.

\subsection{Information d’histoire et statistique d’occupation}
\label{md:comp:informacion}

Les produits
\[
 U_+U_-U_+^{-1}U_-^{-1}
 \qquad\text{et}\qquad
 U_+U_+^{-1}U_-U_-^{-1}
\]
contiennent les mêmes quatre opérations avec les mêmes multiplicités.
Le second vaut \(I\) ; le premier est \(H_{\rm lazo}\ne I\) pour
\(\zeta\ne0\). Par conséquent, toute statistique qui dépend
uniquement des multiplicités attribue le même résultat à deux
protocoles ayant des actions différentes sur l’état.

La conclusion exacte est que cette statistique ne constitue pas une
description suffisante pour prédire la réponse de la composition.
Shannon n’est pas réfuté : une distribution sur des séquences ordonnées
peut représenter l’ordre. On distingue l’information qui a été
conservée dans l’observable choisie.

\begin{definition}[Information opérationnelle d’une histoire]
\label{md:comp:def-informacion}
Dans cette représentation, nous proposons d’appeler information opérationnelle d’une
histoire la classe de son transport ordonné discernable par les
lectures admises. Soit \(\mathcal D\) une collection de détecteurs
définis sur des opérateurs. Elle peut inclure
\[
 D_{X,L}(H_{\rm lazo})=L(H_{\rm lazo}X),
 \qquad D_{\rm tr}(H_{\rm lazo})=\operatorname{tr}H_{\rm lazo},
\]
avec une préparation \(X\) et un lecteur linéaire \(L\).
Deux histoires sont équivalentes si
\[
 D(H_{{\rm lazo},1})=D(H_{{\rm lazo},2})
 \qquad\text{pour tout }D\in\mathcal D.
\]
\end{definition}

Élargir la collection peut séparer des classes auparavant indiscernables.
Avec des réponses vectorielles complètes sur une base, on reconstruit tout
opérateur linéaire ; avec la seule trace, il peut subsister
\(H_{\rm lazo}\sim H_{\rm lazo}^{-1}\).
On distingue ainsi le détecteur d’une réponse de l’état du détecteur
spectral de l’opérateur.

Cette définition n’identifie pas le transport à toute la généalogie
TPK : des histoires distinctes peuvent produire le même opérateur dans une
représentation. L’information latente dans d’autres feuillets, incidences et
prolongements est conservée dans leurs propres lecteurs.

La dimension statistique et la dimension opérationnelle sont compatibles.
Le taux de croissance combinatoire, ou entropie topologique, compte
la croissance asymptotique des configurations admissibles ; un taux
d’entropie de Shannon nécessite en outre une distribution.
Une application d’observation détermine quelles transformations peuvent
être reconstruites. Aucune des preuves précédentes ne déclare une entropie
thermodynamique universelle nulle. L’apport consiste à préciser quelle description
est suffisante pour une tâche de prédiction ou de restitution.

\subsection{Conséquences et nouvelles définitions proposées}
\label{md:comp:consecuencias}

Les résultats précédents réunissent cinq notions, avec leurs domaines
et leurs lecteurs conservés :
\begin{enumerate}
 \item \textbf{Potentiel spectral constitutif.}
 \(\mathscr F\), défini par les moments \(90/120\), a pour
 ses deux dérivées le logarithme de la vitesse normalisée et le
 logarithme de l’impédance normalisée. C’est une définition mathématique
 postérieure aux canaux, sans interprétation énergétique ajoutée.

 \item \textbf{Réciprocité constitutive différentielle.}
 L’égalité des sensibilités croisées démontrée dans la
 section~\ref{md:comp:reciprocidad} fournit une contrainte
 conjointe vérifiable, outre les valeurs.

 \item \textbf{Défaut de fermeture relatif.}
 \(\operatorname{tr}H_{\rm lazo}-2\) détecte le lacet non trivial.
 Il s’accompagne d’une lecture orientée, comme \(\mathcal O\), lorsqu’il
 faut distinguer l’ordre. La trace seule ne définit pas toute la mémoire.

 \item \textbf{Restitution opérationnelle.}
 C’est la récupération de l’action de l’histoire à partir d’une collection
 spécifiée d’observables. Elle se distingue de la restitution de deux
 canaux et de la reconstruction intégrale de l’état.

 \item \textbf{Marge de restitution.}
 La borne \(60a(M_0)\) de la section~\ref{md:comp:restitucion}
 quantifie la stabilité absolue sur un domaine, tandis que
 l’injectivité affirme seulement l’unicité exacte.
\end{enumerate}

Les coordonnées de masse, de fréquence, de rayons et de température du
développement précédent restent intactes. Ce développement ajoute
deux niveaux : la compatibilité différentielle entre lecteurs scalaires
et la mémoire d’ordre dans les compositions opératoires.
On n’identifie pas le facteur de dilatation de \(H_{\rm lazo}\) à
une masse nouvelle sans passer par le lecteur de masses correspondant.

\clearpage
\begingroup
\fontsize{10.2}{12.2}\selectfont
\setlength{\parskip}{1pt plus .5pt minus .5pt}
\titlespacing*{\paragraph}{0pt}{6pt}{.75em}
\section{Conclusions}
\label{md:conclusiones}
Le résultat transversal de ce travail est la détermination conjointe de
structures que leurs évaluations scalaires présentent séparément. Des
applications explicites de compatibilité et de restitution ont été construites, les invariants
de leurs transformations identifiés, et les observations qui
récupèrent l’information qu’elles conservent démontrées. La composition réunit les constantes
par les opérations qu’elles partagent.
L’action participe à une échelle dimensionnelle, au contre-angle et au
quotient de retour électronique. La paire angulaire détermine des canaux
orientés ; les incidences hexade--octade produisent des lecteurs constitutifs
et la fonctionnelle de Barbero--Immirzi. La composition conserve cette généalogie
et permet de formuler des propriétés conjointes qui dépassent l’évaluation
individuelle de chaque lecteur. Le point de départ demeure dans la composition
APP--TRIT--TPK : les relations inverses agissent sur ses sorties construites
et préservent les orientations, les domaines et les données dimensionnelles
qu’elles utilisent.

\paragraph{Mémoire chronologique et réponse opératoire.}

L’analyse chronologique précise quelle information conserve un prolongement.
Dans le facteur de capacité et avec la mesure de phase déclarée, l’entropie
des blocs croît sans borne tandis que son taux par position tend vers zéro.
Le bilan, qui admet deux valeurs, perd une quantité d’information
chronologique également croissante. L’activité quadratique du lecteur de
courant reste positive bien que sa moyenne soit nulle. Ces propriétés
se rapportent à des observations différentes du même processus et peuvent
coexister sans identifier le taux symbolique à une température.

La restitution opératoire ajoute une distinction dynamique. Dans la
représentation nonadique de lecture et de mémoire, la sensibilité à l’ordre
se répartit entre les contributions \(1/81\) et \(8/81\) du
commutateur interne. Réintroduire la mémoire réunit les deux et produit
\(1/9\). L’histoire retenue participe donc à la composition
de la réponse : connaître le bilan présent et connaître la capacité
de répondre à une continuation sont des informations différentes.

\paragraph{Détermination spectrale et stabilité de la reconstruction.}
La restitution spectrale établit un résultat de détermination conjointe. Dans la collection
constitutive normalisée, la vitesse fixe le produit des deux
réponses et la fonctionnelle de Barbero distingue de manière monotone leur répartition
orientée. Les deux données récupèrent les valeurs propres étiquetées. Sur
l’image engendrée, l’inversion se compose avec la restitution angulaire et celle du
coefficient d’action. Les marques de feuillet et les bases dimensionnelles sont
conservées : ce sont des données de la représentation utilisées par la récupération.

La caractérisation différentielle précise cette reconstruction. Les logarithmes de vitesse et
d’impédance sont les dérivées d’un même potentiel spectral, défini
au moyen des canaux d’incidence. Sa hessienne définie négative produit
une réciprocité exacte des sensibilités et permet de borner l’inversion
sur des domaines contrôlés. L’unicité avec des données exactes et la sensibilité
face à une erreur finie sont ainsi caractérisées par des résultats distincts.
Cette précision est particulièrement pertinente lorsqu’une contraction
intense conserve l’injectivité tout en réduisant la séparation
observable entre états.

\paragraph{Composition constitutive et reconstruction électrofaible.}
L’inverse cubique transporte la multiplication des contractions vers une
loi associative sur les réponses. Sa coordonnée logarithmique est
additive et conserve le sens de la composition dans chaque feuillet. La
reconstruction électrofaible précise une autre complémentarité : le quotient
des couplages de jauge récupère la contraction commune, tandis que
Higgs détermine la coordonnée orientée dans le domaine spécifié.
Les projecteurs marqués transforment cette récupération scalaire en
restitution de l’opérateur constitutif. On maintient ainsi une connexion
démontrée entre des lecteurs distincts des mêmes canaux.

\paragraph{La constante de Catalan et le potentiel constitutif.}
Le moment de profondeur deux établit une connexion entre deux évaluations
d’une même structure fonctionnelle. Sur le
quart de tour orienté, sa partie imaginaire produit la collection dont la
valeur extrême est Catalan ; sur les canaux contractifs, le même
moment intervient dans le potentiel constitutif. Les arguments,
l’orientation et la normalisation de chaque évaluation restent spécifiés.
La structure fonctionnelle transmet plus d’information que n’importe laquelle de ses
valeurs extrêmes considérée isolément.
La formule intégrale de la proposition~\ref{md:compos:inversa-catalan}
récupère en outre le moment à partir de la réponse constitutive. La condition
de limite nulle à l’origine fixe de manière unique les constantes d’intégration.
Production et restitution sont ainsi articulées sur une même collection
fonctionnelle, avec son domaine et sa valeur extrême.

\newpage
\paragraph{Reconstruction de l’action et mémoire de l’ordre.}
La composition relative démontre que l’ordre laisse
une empreinte opératoire récupérable. Les quarts de tour de formes
conjuguées et leurs inverses produisent un lacet unimodulaire dont le facteur est
\[
 \rho_{\rm lazo}=
 \left(\frac{501\alpha+\eta_{\rm ret}}
 {499\alpha-\eta_{\rm ret}}\right)^2,
 \qquad 0<\eta_{\rm ret}<499\alpha.
\]
La branche et les références déclarées permettent de récupérer
\(\eta_{\rm ret}\). Son échelle dimensionnelle reste l’antécédent
nécessaire pour reconstruire la section de Planck. Le relèvement au
cycle de douze phases établit où agit cette transformation, et
l’itération exprime son accumulation par des puissances réciproques.
L’holonomie nonadique conserve son propre domaine et sa loi de mise à jour ;
la composition étudiée constitue un transport ultérieur bien défini.

L’information opérationnelle acquiert ainsi une définition précise : deux
histoires sont discernables par rapport aux préparations et aux
détecteurs admis lorsque leurs transports produisent des réponses
différentes. Les multiplicités d’opérations ne déterminent généralement pas
cette réponse ; l’ordre peut être indispensable. La trace et les réponses
vectorielles ne conservent pas non plus les mêmes données. Une collection suffisante
de lectures reconstruit l’opérateur, tandis que d’autres identifient des
transformations d’orientation opposée. Cette définition explicite
la relation entre observation, mémoire et restitution sans confondre une
représentation avec la généalogie complète de l’état.

\paragraph{Réciprocité massique et articulation des grandeurs physiques.}
La normalisation conjointe possède une preuve opératoire préalable :
le retour inscrit produit $L_\alpha$, et le transport polaire du même
caractère détermine les lecteurs $R$ et $\Lambda$.
Les identités $H\Lambda=\hbar cI$ et $R\Lambda=L_\alpha^2I$ fixent
$G=c^3L_\alpha^2/\hbar$ comme coefficient unique de la réalisation radiale.
La preuve conserve les variations non commutatives et l’élimination compatible
de mémoire, y compris le transport dual de la longueur d’action.
Son évaluation après retour dans la section SI est
$6.674299659414022706\ldots\,10^{-11}\,\mathrm{m^3\,kg^{-1}\,s^{-2}}$ ;
la confrontation CODATA s’effectue après cette évaluation.

La composition réunit quatre lectures d’une même coordonnée
structurelle positive de masse. Sous les conditions circulaire et thermique
déclarées, la fréquence propre, le rayon gravitationnel réduit, la
longueur de Compton réduite et l’énergie normalisée par la décision
tritique satisfont
\[
 \Xi_\Gamma=
 \frac{\Omega_\Gamma}{\omega_0}
 =\frac{r_{g,\Gamma}}{R_{108}}
 =\frac{R_{108}}{\bar\lambda_\Gamma}
 =\frac{E_\Gamma}{k_B\Theta_{108}\log3}.
\]
Le préfacteur électronique, le caractère et les tours appartiennent au
constructeur de \(\Xi_\Gamma\). L’identité conserve leurs fonctions et
leurs normalisations. Le facteur massique est caractérisé, dans cette réalisation,
comme paramètre de dilatation réciproque des deux longueurs ; leur produit
maintient l’aire \(R_{108}^2\). La pondération thermique du cycle
s’exprime par \(3^{-\Xi_\Gamma}\), avec les conditions de
normalisation et de convergence correspondantes.
Ici $R_{108}=L_\alpha$ ; la forme circulaire conserve la longueur déjà
construite. Lorsqu’on fait varier l’horloge, on transporte également le quotient électronique
$b_e$, en préservant la même lecture de masse. Les fenêtres chronologiques,
les phases réduites et les coûts terminaux décrivent des aspects de l’état ;
leur conservation isolée permet des changements que la composition complète distingue.

\paragraph{Échelle électronique, Fermi et longueur d’action.}
La substitution de la section électronique complète conserve son préfacteur
et son retour dans la réalisation électrofaible. L’échelle commune
apparaît avec une puissance inverse double dans Fermi au niveau arbre et une puissance
directe dans la valeur d’espérance de Higgs ; elle se simplifie dans les
quotients angulaires et dans l’autocouplage. Les facteurs spécifiques
de chaque espèce restent dans leurs rapports. La congruence positive
transporte en outre la longueur de Compton par la congruence inverse,
sans exiger de commutation entre masse basale et retour. Ces preuves
distinguent l’échelle absolue, l’information angulaire et la structure
spectrale conservée par chaque opération.

\paragraph{Restitution électronique et conservation sous transport.}
La lecture électronique conserve une dépendance vérifiable à l’ordre.
La somme et les différences entre positions opposées reconstruisent le
registre, y compris ses conditions de caractère entier ; les autocorrélations
et le spectre séparent les registres de mêmes multiplicités. Le
coefficient $80$ correspond à l’évaluation du registre central
spécifié et change lorsqu’on modifie son ordre. Pour une collection finie
de multisections, l’opérateur de masse est positif et son commutateur
avec un transport détermine exactement quand celui-ci conserve la masse.
Ce critère permet de transférer la même conservation aux lectures
temporelle, spatiale, gravitationnelle et thermique, avec leurs références fixes.

\newpage
\paragraph{Trace marquée et réponse thermométrique.}
La composition des capacités et des occupations produit
\[
 b_*=10^{-23}+380\,10^{-26}+649\,10^{-29}
     =1{,}380649\times10^{-23}.
\]
La graduation des $649$ états conserve les secteurs $1$, $24$
et $624$ et leurs indices de mémoire respectifs. Le spectre de rapport
$1/10$ provient du transport de mémoire déjà démontré. Sur cette
référence est construite une réalisation thermométrique explicite :
l’information marquée est $b_*s(\rho)$ et l’énergie est la lecture
conditionnée au secteur marqué. Pour des variations de Gibbs avec
hamiltonien et référence fixes, la température satisfait
$b_*\Theta=\beta^{-1}$. Le coefficient est ainsi caractérisé par
une opération thermique concrète. Conditionner les deux grandeurs de la
même manière produit une autre normalisation ; faire varier la référence ajoute
le terme différentiel correspondant. La preuve identifie ces
dépendances et conserve la coordinatisation dimensionnelle nécessaire à la
lecture physique de Boltzmann. En composant avec l’horloge tritique,
l’aire d’intrication et la réciprocité gravitationnelle, on obtient
la relation conjointe entre énergie de décision, action et aire.

\paragraph{Contenu physique de la transformation.}
L’analyse distingue de manière constructive le transport
passif et la réponse relative. Le premier se télescope lorsque la coordinatisation
et la phase reviennent. La seconde conserve un opérateur final distinct de l’identité
lorsque le commutateur est non trivial. La conservation symplectique de l’aire
coexiste avec l’absence d’une métrique énergétique positive commune aux
deux formes conjuguées. Interpréter un gain énergétique exige
d’inclure la réalisation du contrôleur et la référence. Le résultat
mathématique détermine ce qui doit être confronté et quelles données doivent être conservées
pour attribuer une réponse au système physique.

\paragraph{Définition structurelle et unité du résultat.}
Les développements permettent de préciser les définitions utilisées. Le
coefficient d’action caractérise une section régionale qui participe tant
à la construction angulaire qu’à sa restitution par transport. La
fonctionnelle de Barbero--Immirzi distingue la répartition orientée des canaux
constitutifs ; associée à la vitesse normalisée, elle détermine leur spectre.
La constante de Catalan apparaît comme une évaluation extrême du moment
orienté dont la collection intervient également dans le potentiel de réponse.
La coordonnée de masse paramètre, dans la réalisation spécifiée, une
dilatation réciproque qui conserve le produit des longueurs. La
constante de Boltzmann articule les droites de température et d’énergie de
décision au moyen du transducteur positif et de la normalisation tritique
déclarés. Chaque définition ajoute à la valeur une opération, un domaine et une
fonction dans la composition.

\paragraph{Séparation critique et restitution des lecteurs.}
La loi d’échelle paramétrique porte sur la suite sélectionnée par la construction elle-même : les premières racines de retour du lecteur dodécaphasique uniforme. APP--TRIT--TPK, le continu conjoint, l’orientation et la mémoire demeurent son antécédent. L’identification des centres et le contrôle du reste de \ref{hmtfg:escalado:15}--\ref{hmtfg:escalado:17} aboutissent à \eqref{hmtfg:articulacion:limite} : la constante de Feigenbaum est l’invariant asymptotique de séparation de ces paramètres. Le même développement restitue le profil à partir de sa mesure à douze nœuds par Jacobi et, par inversion constitutive, restitue les canaux orientés et le moment dont l’extrémité est la constante de Catalan. Le bilan réunit ainsi génération, sélection et restitution, en conservant le domaine et l’information qui distinguent chaque lecteur.

% Adición a las conclusiones de VII; las conclusiones anteriores se conservan.
\paragraph{Intégration des quatre interactions et clôture canonique totale.}
La gravitation et les interactions forte, faible et électromagnétique agissent sur une même réalisation matérielle de la structure APP--TRIT--TPK. Le transport conserve les feuillets, l’orientation, la mémoire et l’incidence de l’état enrichi ; ses réalisations de spin et de jauge se composent sur les mêmes fibres matérielles. La composante électromagnétique appartient à la réalisation électrofaible : elle n’est pas introduite comme une interaction supplémentaire qui en serait déconnectée. L’action totale utilise les coefficients et les lecteurs déjà construits, dont la section de Planck, la longueur radiale et la fonctionnelle de Barbero--Immirzi.

Le contenu de cette intégration ne se limite pas à réunir des termes dans une expression. La variation de la même action produit le courant total de spin et la source métrique conjointe. L’élimination de la contorsion conserve les termes croisés entre espèces ; la transformation de Legendre conserve les moments matériels et produit les contraintes totales. Les preuves des sections~\ref{hmtcuatro:gea:seccion} et~\ref{hmtcuatro:accion-retroaccion-restricciones} établissent la rétroaction, le caractère de première classe et la propagation de ces contraintes dans la réalisation régulière spécifiée.

En particulier, soit $\Theta_{\rm tot}$ le potentiel canonique total obtenu dans cette transformation, $\mathcal Z=\{C_A=0\}$ la surface régulière des contraintes et $X_N=N^AX_{C_A}$, $X_M=M^BX_{C_B}$ deux déformations caractéristiques admissibles. La connexion induite par l’action possède $H_N^{\rm tot,can}=-\Theta_{\rm tot}(X_N)$. Sa courbure, calculée sans omettre les dérivées des coefficients de déformation, est
\begin{equation}
\begin{aligned}
\mathcal F^{\rm tot,can}_{N,M}
&=\delta_NH_M^{\rm tot,can}-\delta_MH_N^{\rm tot,can}
 -H_{[X_N,X_M]}^{\rm tot,can}
 +\frac{i}{\hbar}[H_N^{\rm tot,can},H_M^{\rm tot,can}]\\
&=-\mathrm d\Theta_{\rm tot}(X_N,X_M)
 =N^AM^B f_{AB}{}^{D}C_D.
\end{aligned}
\label{hmtcuatro:conclusion-curvatura-total}
\end{equation}
Par conséquent,
\[
 \left.\mathcal F^{\rm tot,can}_{N,M}\right|_{\mathcal Z}=0.
\]
La formule inclut la déformation tangentielle et les actions de jauge du crochet total ; si l’on travaille dans le quotient, elle exprime la même égalité modulo ces directions verticales. La représentation préquantique conserve cette clôture au moyen de l’identité de commutateurs démontrée dans \eqref{hmtcuatro:representacion-precuantica}.

Ainsi, l’incorporation gravitationnelle fait partie de la même action, de ses sources et de son algèbre de contraintes ; ce n’est pas une équation juxtaposée après la construction des autres interactions. L’annulation démontrée est celle de la courbure caractéristique canonique ; elle n’exige pas l’annulation de la courbure de l’espace-temps ni des intensités des champs de jauge. Son domaine conserve le corepère non dégénéré, la carte régulière et les conditions de bord utilisées dans la preuve. Elle ne supprime pas non plus la mémoire ni la monodromie globale des histoires qui ont donné lieu à la réalisation.


La composition d’horizon distingue la dépendance de l’action à aire
fixe et à masse fixe. La section~\ref{ins29:vii:horizonte} démontre en outre
une courbe de compatibilité entre le biais thermique, la forme elliptique et les
canaux de Barbero. Les rapports de probabilités restituent la même
section d’action et fournissent une vérification conjointe de ces lectures.


Production, compatibilité et restitution constituent donc trois
niveaux complémentaires d’une même description. La généalogie fixe
les objets, les relations conjointes contraignent leurs lectures et les
détecteurs déterminent la structure récupérable. Les nouvelles
caractérisations s’appuient sur les opérations exposées et conservent
leurs conditions, ce qui permet de les incorporer à des développements sectoriels
sans dissoudre l’unité de l’argument.

\endgroup
\clearpage\appendix
% Correspondencia editorial: gestion/mapa_transportes_auxiliares.md.
% Realizaciones auxiliares: no sustituyen el eje de compatibilidad y restitución.
\section{Transport compensé, mémoire géométrique et confrontation énergétique}
\label{md:aux:transporte}

\subsection{Origine opératorielle et domaine de la réalisation}
\label{md:aux:origen}

La construction conserve APP--TRIT--TPK et l’état enrichi. APP
réunit les feuillets arithmétiques avec leurs résidus et quotients ; TRIT conserve
le régime et l’orientation ; TPK compose sélection, transport et
mise à jour de la mémoire. Les représentations matricielles utilisées
ci-dessous sont des réalisations ultérieures. L’holonomie nonadique
$\Gamma_9$ et sa mémoire maintiennent leurs propres opérateurs.

Les antécédents employés sont les représentations quadratiques
elliptique, hyperbolique et parabolique du TRIT, la différentielle des
transports inversibles, la fonctionnelle énergétique
$E=\langle Df,WDf\rangle$ avec $W>0$, le retour avec mémoire
translationnelle, l’ellipse d’action avec produit des demi-axes $2\hbar$
et sa réalisation LC, la polarisation des vacances et le couplage
réversible de deux modes. Le retour translationnel et le commutateur
qui suit sont des objets distincts. Le lacet elliptique--parabolique et
son bilan d’aire orientée seront ensuite utilisés comme collection
de confrontation.

La composition compensée et le protocole de confrontation réunissent ces
réalisations. La théorie algébrique des traces de commutateurs dans
$\operatorname{SL}(2)$ constitue un antécédent mathématique
établi.\footnote{Goldman, \emph{Trace Coordinates on
Fricke spaces}, \href{https://arxiv.org/abs/0901.1404}{arXiv:0901.1404}.}
L’objet de cet appendice est leur composition concrète avec les
lecteurs déclarés. Les paramètres $a,b$ sont des paramètres des
réalisations, non des constantes physiques ajustées.

La portée s’arrête à cette réalisation locale. L’identification de
chemins TPK admettant un dispositif matériel avec ces paramètres
exige de composer leurs lecteurs effectifs. Une matrice inversible
arbitraire n’atteste pas à elle seule un chemin TPK, et l’identification
du défaut de commutation à la torsion de Cartan nécessite sa
connexion et sa soudure.

\subsection{Compensation des opérations et conservation de l’aire}
\label{md:aux:conmutador}

Les matrices de régime sont
\begin{equation}
 J_{\mathrm e}=\begin{pmatrix}0&1\\-1&0\end{pmatrix},\qquad
 J_{\mathrm h}=\begin{pmatrix}0&1\\1&0\end{pmatrix}.
 \label{md:aux:regimenes}
\end{equation}
Elles satisfont $J_{\mathrm e}^2=-I$, $J_{\mathrm h}^2=I$ et
$[J_{\mathrm e},J_{\mathrm h}]=2\operatorname{diag}(1,-1)$.
On définit
\begin{equation}
 R(a)=e^{aJ_{\mathrm e}},\qquad B(b)=e^{bJ_{\mathrm h}},\qquad
 C_{\mathrm{aux}}(a,b)=R(a)B(b)R(a)^{-1}B(b)^{-1}.
 \label{md:aux:def-compensado}
\end{equation}
L’action sur les colonnes se lit de droite à gauche. Le protocole
contient chaque transformation et son inverse, mais la suite n’est pas
le retour en arrière complet du parcours. Un retour en arrière complet aurait,
par exemple, la forme $RBB^{-1}R^{-1}=I$. Le symbole
$C_{\mathrm{aux}}$ distingue ce commutateur elliptique--hyperbolique
du contre-angle $C^*$ et du lacet de formes conjuguées
$H_{\mathrm{lazo}}$ étudié dans le corps principal.

\begin{proposition}
\label{md:aux:prop-traza}
Pour $a,b\in\mathbb R$,
\begin{equation}
 \det C_{\mathrm{aux}}=1,\qquad
 \operatorname{tr}C_{\mathrm{aux}}=2+4\sin^2a\sinh^2b.
 \label{md:aux:traza-compensado}
\end{equation}
\end{proposition}
\begin{proof}
Posons $c_a=\cos a$, $s_a=\sin a$, $u_b=\cosh b$ et $v_b=\sinh b$.
Alors
\[
 R=\begin{pmatrix}c_a&s_a\\-s_a&c_a\end{pmatrix},\qquad
 B=\begin{pmatrix}u_b&v_b\\v_b&u_b\end{pmatrix},
\]
avec $c_a^2+s_a^2=1$ et $u_b^2-v_b^2=1$. Toutes les matrices ont un
déterminant égal à un. La multiplication donne
\begin{equation}
 \begin{aligned}
 (C_{\mathrm{aux}})_{11}&=u_b^2+2c_as_au_bv_b-(c_a^2-s_a^2)v_b^2,\\
 (C_{\mathrm{aux}})_{12}&=-2s_a^2u_bv_b-2c_as_av_b^2,\\
 (C_{\mathrm{aux}})_{21}&=-2s_a^2u_bv_b+2c_as_av_b^2,\\
 (C_{\mathrm{aux}})_{22}&=u_b^2-2c_as_au_bv_b-(c_a^2-s_a^2)v_b^2.
 \end{aligned}
 \label{md:aux:entradas-compensado}
\end{equation}
La somme diagonale est
$2u_b^2-2(c_a^2-s_a^2)v_b^2=2+4s_a^2v_b^2$, comme affirmé.
\end{proof}

Lorsque $s_av_b\ne0$, la trace dépasse deux et le commutateur est
hyperbolique, avec des valeurs propres positives réciproques
\begin{equation}
 e^{\lambda},\ e^{-\lambda},\qquad
 \lambda=2\operatorname{arsinh}|\sin a\sinh b|>0.
 \label{md:aux:espectro-compensado}
\end{equation}
Cela s’obtient à partir de $\cosh\lambda=\operatorname{tr}(C_{\mathrm{aux}})/2$
et $\cosh(2\operatorname{arsinh}x)=1+2x^2$.
Si $s_av_b=0$, alors $b=0$ ou $R=\pm I$, et
$C_{\mathrm{aux}}=I$.
Le développement local est
\begin{equation}
 C_{\mathrm{aux}}=I+2ab\operatorname{diag}(1,-1)
 +O(|a|^2|b|+|a||b|^2).
 \label{md:aux:expansion-compensado}
\end{equation}
Le résidu commence par le produit des deux opérations. Appliquer
l’une quelconque d’entre elles suivie de sa propre inverse produit l’identité.

Une matrice réelle bidimensionnelle de déterminant un conserve la forme
d’aire. Ainsi, $C_{\mathrm{aux}}$ transforme l’ellipse
d’action en une autre ellipse d’aire $h$. Si $\Sigma>0$ est une matrice
de covariance,
\begin{equation}
 \Sigma'=C_{\mathrm{aux}}\Sigma C_{\mathrm{aux}}^{\mathsf T},
 \qquad \det\Sigma'=\det\Sigma.
 \label{md:aux:covarianza-area}
\end{equation}
La distribution des demi-axes et leur orientation peuvent changer.
Les valeurs propres $e^{\pm\lambda}$ décrivent des directions propres et des
itérations ; les valeurs singulières déterminent une autre lecture et, en
général, ne coïncident pas avec elles. La variation des demi-axes euclidiens
en une seule application nécessite la métrique correspondante.

La conservation de l’aire d’action est compatible avec une
accumulation de déformation orientée. L’observable géométrique
$\lambda$ dépend de $a,b$ et ne constitue pas une constante universelle
supplémentaire. La conservation du déterminant ou de l’aire, à elle
seule, n’établit pas la conservation de l’énergie physique.

\subsection{Lecteur positif du résidu et inversion du parcours}
\label{md:aux:residuo}

Dans une fibre de retour, avec un champ $f$ et un poids hermitien $W>0$
déclarés, le lecteur énergétique prend la forme
\begin{equation}
 E_C(f)=\langle(I-C_{\mathrm{aux}})f,
                  W(I-C_{\mathrm{aux}})f\rangle.
 \label{md:aux:energia-residuo}
\end{equation}
Pour $s_av_b\ne0$, le commutateur n’a pas de valeur propre égale à un et
$E_C(f)>0$ pour tout $f\ne0$. Le retour en arrière complet produit un résidu
nul. La normalisation dimensionnelle appartient à la fonctionnelle de
départ : ses coordonnées numériques nécessitent leur réalisation avant
d’être interprétées comme des joules, de la chaleur, une masse ou un courant de conduction.

L’inversion complète du parcours inverse $C_{\mathrm{aux}}$.
La trace et $\lambda$ identifient l’opérateur et son inverse. La
restitution de direction nécessite une lecture supplémentaire : le vecteur
résiduel, les composantes du transport ou une sonde avec résolution
de phase. Cette comparaison limite concrètement la suffisance
d’une lecture exclusivement scalaire.

Un exemple rationnel, indépendant de données physiques, est
\begin{equation}
 R=\begin{pmatrix}3/5&4/5\\-4/5&3/5\end{pmatrix},\qquad
 B=\begin{pmatrix}5/4&3/4\\3/4&5/4\end{pmatrix}.
 \label{md:aux:ejemplo-racional}
\end{equation}
On y obtient $\operatorname{tr}C_{\mathrm{aux}}=86/25$,
$\det C_{\mathrm{aux}}=1$ et les valeurs propres
$(43\pm6\sqrt{34})/25$. Les contrôles rationnels du supplément
comprennent cet exemple et $169$ paires, avec des cas nuls, un retour en arrière
complet et la conservation du déterminant de covariance. Ce sont
des vérifications finies des identités démontrées.

\subsection{Couplage réversible et lecture réduite de l’intrication}
\label{md:aux:acoplamiento}

Pour $H\in\operatorname{Sp}(2,\mathbb R)$, la réalisation à deux
modes utilise le couplage
\begin{equation}
 U_H=\frac19\begin{pmatrix}
 8I+H&\sqrt8(H-I)\\
 \sqrt8(H-I)&I+8H
 \end{pmatrix}.
 \label{md:aux:acoplamiento-uh}
\end{equation}
Avec deux gaussiennes pures identiques préparées à partir de l’ellipse
d’action et avec un transport conjoint de la mise en coordonnées initiale, le lecteur
réduit satisfait
\begin{equation}
 \nu(H)^2=1+\frac8{81}
    \bigl\{\operatorname{tr}(HH^{\mathsf T})-2\bigr\},
 \qquad\operatorname{Tr}\rho_{\mathrm{vis}}^2=\nu(H)^{-1}.
 \label{md:aux:mezcla-generica}
\end{equation}
La préparation et les lecteurs restent fixés. La
transformation conjointe est réversible et conserve la pureté globale.

\begin{corollary}
\label{md:aux:cor-mezcla}
Pour $C_{\mathrm{aux}}(a,b)$,
\begin{equation}
 \operatorname{tr}(C_{\mathrm{aux}}C_{\mathrm{aux}}^{\mathsf T})-2
 =16\sin^2a\sinh^2b\cosh^2b,
 \label{md:aux:traza-cuadratica}
\end{equation}
et
\begin{equation}
 \boxed{\nu_C^2=1+\frac{32}{81}\sin^2a\sinh^2(2b).}
 \label{md:aux:mezcla-compensado}
\end{equation}
\end{corollary}
\begin{proof}
Les quatre entrées de \eqref{md:aux:entradas-compensado}
s’écrivent $d+k$, $l-m$, $l+m$, $d-k$, où
\[
 d=1+2s_a^2v_b^2,\quad k=2c_as_au_bv_b,\quad
 l=-2s_a^2u_bv_b,\quad m=2c_as_av_b^2.
\]
La somme de leurs carrés est
$2(d^2+k^2+l^2+m^2)=2+16s_a^2u_b^2v_b^2$.
La substitution dans \eqref{md:aux:mezcla-generica} et
$\sinh(2b)=2\sinh b\cosh b$ prouve le résultat.
\end{proof}
L’exemple \eqref{md:aux:ejemplo-racional} produit
$\nu_C^2=17/9$ et la pureté $3/\sqrt{17}$.

L’annulation de $\hbar$ et de la déformation initiale utilise
la préparation déclarée et le transport conjoint de la mise en coordonnées.
Un protocole avec $s_av_b\ne0$ produit un mélange réduit ;
l’inversion complète du couplage peut le défaire. Mélange
réduit, énergie stockée et chaleur conservent leurs fonctions
comme lectures distinctes.

\subsubsection{Contraste signé dans une collection elliptique--parabolique}
\label{md:aux:familia-hn}

La collection de lacets construite est
\begin{equation}
 H_n=\begin{pmatrix}1+n&n^2\\n&n^2-n+1\end{pmatrix},
 \qquad n\in\mathbb Z.
 \label{md:aux:hn}
\end{equation}
Son bilan alterné de mémoire dépend de $n^2$, tandis que le
quotient des déterminants de covariance est
\begin{equation}
 \mathcal R_n=\frac{\det V_{\mathrm{vis},n}}{\det V_0}
 =1+\frac8{81}n^2(2n^2-2n+5).
 \label{md:aux:covarianza-hn}
\end{equation}
En soustrayant les deux orientations, les termes pairs s’annulent :
\begin{equation}
 \boxed{\mathcal R_{-n}-\mathcal R_n=\frac{32n^3}{81}.}
 \label{md:aux:diferencia-direccional}
\end{equation}
Pour $n=1$ et $n=-1$, les réponses sont $121/81$ et $153/81$,
avec les puretés $9/11$ et $3/\sqrt{17}$, respectivement.
Cette collection satisfait $H_{-n}\ne H_n^{-1}$ en général ;
elle n’est pas non plus l’itération $H_1^n$. Le changement du paramètre
orienté doit être distingué du retour en arrière complet des
opérations. Sa lecture alternée d’aire conserve une
identification distincte de celle de la fonctionnelle de Barbero.

Le critère de mélange emploie le commutateur $[D,J]$, avec
$D=\sqrt8(H-I)/9$. Si $D$ commute avec la structure complexe
compatible avec l’état initial, il peut transporter de la mémoire
sans mélanger ce produit gaussien. S’il ne commute pas, le lecteur
détermine le mélange. L’absence de mélange pour cet état
est une propriété distincte de l’absence de dissipation matérielle.

La différence \eqref{md:aux:diferencia-direccional} est un corollaire
de la collection et de son lecteur, non une constante universelle.
La formule de $C_{\mathrm{aux}}$ compose le même lecteur avec
une autre collection admise par le théorème. Une confrontation matérielle
doit déterminer les transports et les couplages avant de mesurer ;
la programmation arbitraire de ces matrices atteste leur
émulation et conserve une portée différente.

\subsection{Deuxième moment du courant de vacances}
\label{md:aux:vacancias}

Le lecteur de capacité fixe
\[
 \delta=\log_{10}(10/9),\qquad z_n\in\{0,1\},\qquad
 P_N=\sum_{n<N}(z_n-\delta),\qquad |P_b-P_a|<1.
\]
Sur des intervalles égaux $t_0$, le courant centré est
$I_n=A_I(z_n-\delta)$, avec $A_I=u_Q/t_0$.
Le symbole $A_I$ distingue cette amplitude de l’angle principal.
Pour $L=b-a$,
\begin{equation}
 \frac1{A_I^2}\sum_{n=a}^{b-1}I_n^2
 =L\delta(1-\delta)+(1-2\delta)(P_b-P_a).
 \label{md:aux:segundo-momento}
\end{equation}
La preuve consiste à développer $(z_n-\delta)^2$, à utiliser
$z_n^2=z_n$ et à sommer. Le déséquilibre signé demeure
borné, tandis que
$I_{\mathrm{rms}}/|A_I|\longrightarrow\sqrt{\delta(1-\delta)}$.
L’amplitude $A_I$ est le saut entre les deux niveaux de courant.
La confrontation conserve la réalisation temporelle et doit comptabiliser
tout couplage qui filtre la suite.

Un bilan signé nul peut coexister avec une lecture
quadratique positive. Si l’on ajoute une résistance idéale $R_E>0$
et que l’on maintient ce courant par intervalles, le modèle ultérieur
de confrontation donne
\begin{equation}
 Q_{R_E}=R_Et_0\sum_n I_n^2.
 \label{md:aux:contraste-resistivo}
\end{equation}
La résistance appartient à cette réalisation ultérieure. Sa
fonction consiste à confronter une réduction possible de chaleur
à un bilan explicite.

\subsection{Conditions d’une confrontation matérielle}
\label{md:aux:contraste-material}

Le protocole sépare deux questions. Pour la mémoire d’ordre,
on compare des parcours compensés, des échanges d’ordre et des
retours en arrière complets au moyen d’une réponse rémanente vectorielle
ou résolue en phase. Une intensité ou une charge nette isolées
peuvent perdre la différence recherchée. Pour la récupération
d’énergie, on mesure l’entrée, l’énergie utile ou restituée, la variation
du stockage et la chaleur, y compris la préparation et la restauration
du dispositif. Le retour de phase conserve séparée la
variation de l’énergie stockée.

Dans une confrontation scalaire à un port, la relation
$v_2(t)=v_1(T-t)$ conserve la distribution des amplitudes et le
spectre de puissance. En régime périodique, un même système
linéaire stationnaire prédit une puissance moyenne identique. Une loi
instantanée $i=g(v)$ conserve également l’intégrale de $vi$.
Une différence reproductible examinerait une dépendance historique,
mais son attribution nécessite de comparer l’hystérésis, l’échauffement
et des modèles dynamiques alternatifs. Avec plusieurs ports, il faut
en outre conserver les spectres croisés.

Une réduction de la dissipation est testée à transfert
utile, fidélité et rythme égaux. L’atténuation doit distinguer l’énergie,
l’amplitude du signal, la population de porteurs et la charge conservée
avant de lui attribuer une origine gravitationnelle. La distinction entre
stockage, récupération et pertes possède des précédents
expérimentaux dans la récupération électrocalorique.\footnote{Defay
et al., \emph{Enhanced electrocaloric efficiency via energy
recovery} (2018),
\href{https://doi.org/10.1038/s41467-018-04027-9}{doi:10.1038/s41467-018-04027-9}.
Cet antécédent ne constitue pas une preuve expérimentale du
transport HMT composé ici.}

Le résultat de l’appendice est une identité générale au sein
de la collection matricielle spécifiée, un lecteur positif du
résidu, un contrôle d’inversion et une loi exacte du deuxième
moment du courant. Sa portée distingue ces résultats
d’une nouvelle constante universelle, d’un prototype sans échauffement
ou d’une démonstration de désintégration électronique gravitationnelle.
Le lien prospectif avec un dispositif doit être fixé
avant d’obtenir les données de confrontation.

\section{Action, mémoire réduite et travail récupérable}
\label{md:aux:trabajo}

\subsection{Composition de l’action avec une réalisation énergétique}
\label{md:aux:accion}

La composition reçoit les sorties d’action et les transports
de mémoire d’APP--TRIT--TPK, en conservant l’état enrichi
comme antécédent. Son domaine est la réalisation gaussienne et
énergétique déclarée. La passivité et l’ergotropie sont des concepts
établis ; ils sont employés pour étudier le couplage spécifique
et les lecteurs déjà construits.

L’équation d’action maintient ses coefficients explicites :
\begin{equation}
 H_5=\frac12\sqrt{\frac{1000\alpha}{\varphi}}-\alpha
 +\frac9{16}\alpha^2-\frac59\alpha^3
 +\frac7{48}\alpha^4-\frac1{54}\alpha^5,
 \label{md:aux:h5}
\end{equation}
\begin{equation}
 D_A=e^{-100\pi\alpha/9},\qquad
 \mathcal C_\pi=169\alpha^6+\frac{D_A\alpha^7}{1-D_A\alpha},
 \label{md:aux:correccion-accion}
\end{equation}
\begin{equation}
 \begin{aligned}
 \hbar_{\mathrm{pre}}&=s_0H_5,\\
 \hbar_{\mathrm{ret}}&=s_0\left(H_5-\frac{90}{\pi}\mathcal C_\pi\right),
 \qquad h_s=2\pi\hbar_s,\qquad s_0=10^{-34}\mathcal U_S.
 \end{aligned}
 \label{md:aux:secciones-accion}
\end{equation}
La décennie et la mise en coordonnées dimensionnelle conservent des antécédents
distincts. La formule maintient les coordonnées $\pi,\varphi$,
l’orientation et l’échelle du lecteur, en plus de $\alpha$.
L’égalité dimensionnelle conserve son unité d’action.

Une section positive étant fixée, on écrit $\hbar=\hbar_s$.
L’ellipse a pour demi-axes $\sqrt{2\hbar}e^{\zeta/2}$ et
$\sqrt{2\hbar}e^{-\zeta/2}$, où
$\zeta=\operatorname{artanh}(C^*/A)$. Son aire orientée satisfait
\begin{equation}
 \mathcal A(\theta)=\hbar\theta,\qquad
 \mathcal A(2\pi)=h.
 \label{md:aux:area-fase}
\end{equation}
Dans la réalisation LC, avec une horloge $\omega$ et une impédance de
référence déclarées,
\begin{equation}
 \mathcal H_\zeta(Q,P)
 =\frac\omega2\left(e^{-\zeta}Q^2+e^\zeta P^2\right).
 \label{md:aux:hamiltoniano-lc}
\end{equation}
Sur l’orbite, on obtient $\mathcal H_\zeta=\hbar\omega=hf$,
avec $f=\omega/(2\pi)$. La convention hamiltonienne donne
$\dot\theta=-\omega$ et, par conséquent,
$d\mathcal A/dt=-\hbar\omega$. L’énergie positive coïncide
avec la grandeur de cette vitesse aréolaire dans la réalisation et
la section fixe considérées. L’aire accumulée et l’action
lagrangienne d’une trajectoire arbitraire sont des objets distincts.

L’équation interne détermine l’échelle d’action par unité
de phase ; l’horloge fixe son rythme. La déformation de l’ellipse
et l’ordre du transport conservent une information supplémentaire par rapport à
$hf$. La préparation gaussienne employée ensuite est distincte
d’un point de l’orbite classique : son énergie initiale est
$hf/2$, au lieu de $hf$.

\subsection{Préparation et lecteurs du résultat énergétique}
\label{md:aux:hipotesis-gaussianas}

On prépare deux modes gaussiens purs, centrés et identiques,
avec la covariance
\begin{equation}
 V_0=\frac\hbar2MM^{\mathsf T},\qquad
 M=\operatorname{diag}(e^{\zeta/2},e^{-\zeta/2}).
 \label{md:aux:preparacion}
\end{equation}
$U_H$ de \eqref{md:aux:acoplamiento-uh} agit, transporté vers
la mise en coordonnées initiale au moyen de $M\oplus M$. L’état conjoint
demeure pur. Les lecteurs énergétiques finaux ont une
même fréquence $\omega$, conservent la métrique initiale
$G_\zeta=M^{-\mathsf T}M^{-1}$ et ne comprennent pas d’énergie
d’interaction résiduelle entre les modes. Les opérations
d’extraction commencent et se terminent avec ces mêmes lecteurs.
L’accessibilité physique des opérations admises est une
hypothèse de la borne de travail, distincte de la construction
d’un dispositif qui les réalise.

On définit
\begin{equation}
 \begin{gathered}
 \tau=\operatorname{tr}(HH^{\mathsf T})-2\geq0,\\
 W_v=\frac{8I+HH^{\mathsf T}}9,\qquad
 W_m=\frac{I+8HH^{\mathsf T}}9.
 \end{gathered}
 \label{md:aux:covarianzas-reducidas}
\end{equation}
Les covariances sont $V_j=(\hbar/2)MW_jM^{\mathsf T}$.
La pureté visible est notée $\mu=\operatorname{Tr}\rho_v^2$
et sa valeur propre symplectique normalisée $\nu=\mu^{-1}$.
Dans cet appendice, $\nu$ remplit cette fonction ; la fréquence
ordinaire s’écrit $f$ et la fréquence angulaire $\omega$.

\subsection{Identité entre énergie visible et pureté}
\label{md:aux:energia-pureza}

\begin{proposition}
\label{md:aux:prop-energia-pureza}
Avec la préparation et les lecteurs précédents, l’énergie
visible au-dessus de l’état initial satisfait
\begin{equation}
 \frac{\Delta E_v}{hf}=\frac\tau{36},\qquad
 \nu^2=1+\frac{8\tau}{81},\qquad
 \boxed{\mu^{-2}-1=\frac{32}{9}\frac{\Delta E_v}{hf}.}
 \label{md:aux:identidad-energia-pureza}
\end{equation}
\end{proposition}
\begin{proof}
L’état est centré et le lecteur quadratique demeure fixe.
Par conséquent,
\[
 E_v=\frac\omega2\operatorname{tr}(G_\zeta V_v)
 =\frac{\hbar\omega}{4}\operatorname{tr}W_v.
\]
Comme $\operatorname{tr}W_v=2+\tau/9$, en soustrayant
$\hbar\omega/2$, on obtient la première égalité.
La covariance réduite satisfait
$\det W_v=1+8\tau/81$ et $\mu=(\det W_v)^{-1/2}$.
L’élimination de $\tau$ prouve l’identité encadrée.
La démonstration ne nécessite ni la symétrie de $H$ ni la coïncidence
entre ses valeurs propres et ses valeurs singulières.
\end{proof}

La loi élimine les paramètres particuliers du parcours.
Le coefficient $32/9$ provient de la répartition du couplage
et l’échelle $h$ est conservée jusqu’à la normalisation.
Pour vérifier la spécificité de cette répartition, considérons
$W_p=pI+(1-p)HH^{\mathsf T}$, avec $0<p<1$. On obtient
\begin{equation}
 \mu^{-2}-1=4p\frac{\Delta E}{hf}.
 \label{md:aux:reparto-general}
\end{equation}
Le choix déclaré est $p=8/9$. Lorsque $\Delta E>0$,
une confrontation peut estimer
\begin{equation}
 p_{\mathrm{obs}}=
 \frac{\mu^{-2}-1}{4\Delta E/(hf)}
 \label{md:aux:reparto-observado}
\end{equation}
et le comparer à $8/9$ sans ajuster $p$ à ces mêmes données.
Programmer un mélangeur avec cette répartition vérifie une
émulation ; l’attribution de la répartition à un mécanisme naturel
conserve sa preuve indépendante.

L’identité utilise des gaussiennes centrées et la préparation
spécifiée. Un déplacement cohérent ajoute de l’énergie sans
changer la pureté. Le bruit thermique initial et les changements de
métrique modifient la relation. Ces variations délimitent
l’énoncé et fournissent des critères de réfutation d’une extension à des états
arbitraires.

La combinaison convexe $W_v$ est une combinaison de matrices
de covariance produite en réduisant un état gaussien conjoint.
Un mélange probabiliste de deux états purs constitue
un autre objet, généralement non gaussien ; sa pureté n’est pas
donnée par le déterminant de cette covariance.

\subsection{Travail récupérable sous opérations locales}
\label{md:aux:ergotropia}

L’ergotropie est la diminution maximale d’énergie par des
opérations unitaires cycliques avec un lecteur énergétique fixe.
La définition ne comprend pas automatiquement le coût net de
construction, de préparation ou de contrôle de l’appareil.

Le spectre réduit possède les populations
\begin{equation}
 \lambda_j=(1-r)r^j,\qquad r=\frac{\nu-1}{\nu+1}.
 \label{md:aux:espectro-gaussiano}
\end{equation}
Elles sont ordonnées de manière décroissante avec l’énergie de l’oscillateur.
L’état passif a ainsi l’énergie $E_{v,\mathrm{pas}}=hf\nu/2$.
Une opération gaussienne l’atteint :
\begin{equation}
 S=\sqrt\nu\,W_v^{-1/2},\qquad
 \det S=1,\qquad SW_vS^{\mathsf T}=\nu I.
 \label{md:aux:pasivacion}
\end{equation}
L’ordre des populations démontre également la minimalité
parmi toutes les opérations unitaires, au-delà des gaussiennes.

Par conséquent,
\begin{equation}
 \boxed{
 \frac{\mathcal W_v}{hf}
 =\frac{(\nu-1)(9\nu-7)}{32},\qquad
 \frac{E_{v,\mathrm{pas}}-hf/2}{hf}=\frac{\nu-1}{2}.}
 \label{md:aux:trabajo-visible}
\end{equation}
Les deux quantités ont pour somme
$\Delta E_v/(hf)=9(\nu^2-1)/32$. Le terme passif est une
énergie inaccessible sous les opérations locales déclarées ;
sa présence n’équivaut pas à un flux de chaleur réalisé.

\subsubsection{Accès conjoint aux corrélations}
\label{md:aux:acceso-conjunto}

Le deuxième mode satisfait $\det W_m=\det W_v=\nu^2$ et
$\Delta E_m=8\Delta E_v$. Par conséquent,
\begin{equation}
 \Delta E_{\mathrm{tot}}=9\Delta E_v=hf\frac\tau4.
 \label{md:aux:energia-conjunta}
\end{equation}
L’inversion conjointe du couplage restitue le produit
initial. L’état total est pur et ce produit est l’état
fondamental du lecteur conjoint ; l’accès conjoint permet de
récupérer idéalement $\Delta E_{\mathrm{tot}}$.
Les opérations locales indépendantes laissent deux états
passifs, chacun d’énergie $hf\nu/2$. On en déduit
\begin{equation}
 \boxed{\mathcal W_{\mathrm{conjunta}}-\mathcal W_{\mathrm{local}}
 =hf(\nu-1)=hf(\mu^{-1}-1).}
 \label{md:aux:ventaja-conjunta}
\end{equation}

L’identité quantifie le travail supplémentaire accessible
au moyen des corrélations. La préparation initiale a nécessité
de l’énergie et le bilan d’un cycle complet comprend cet apport.
La relation compare des lectures énergétiques et des opérations sur
l’état, sans identifier l’information et l’énergie à une
même grandeur physique.

Dans cet énoncé, ``local'' signifie des opérations unitaires de produit
$U_v\otimes U_m$ ; les mesures,
la communication classique et la rétroaction sont exclues de cette classe. L’accroissement
provient de la restriction opérationnelle sur l’état
corrélé et de ses spectres marginaux, non d’une
énergie d’interaction résiduelle. La répartition $1:8$ concerne
les énergies au-dessus du vide, non les énergies
totales ni les ergotropies.

\subsection{Cas exacts et orientation du transport}
\label{md:aux:casos-energia}

Pour $H_n$ de \eqref{md:aux:hn},
\begin{equation}
 \tau_n=n^2(2n^2-2n+5).
 \label{md:aux:tau-hn}
\end{equation}
Dans $n=1$, on obtient
\begin{equation}
 \begin{gathered}
 \nu=11/9,\qquad\mu=9/11,\qquad
 \Delta E_v=\frac5{36}hf,\\
 \mathcal W_v=\frac1{36}hf,\qquad
 E_{v,\mathrm{pas}}-hf/2=\frac19hf.
 \end{gathered}
 \label{md:aux:caso-positivo}
\end{equation}
L’ensemble possède $\Delta E_{\mathrm{tot}}=5hf/4$, le travail
local $37hf/36$ et l’avantage conjoint $2hf/9$.
Le cas $n=-1$ partage les valeurs propres de $H_1$,
mais possède $\tau=9$ et
\begin{equation}
 \nu=\frac{\sqrt{17}}3,\qquad
 \Delta E_v=\frac{hf}{4},\qquad
 \mathcal W_v=\frac{9-2\sqrt{17}}{12}hf.
 \label{md:aux:caso-negativo}
\end{equation}
L’énergie visible ajoutée diffère dans le rapport $9/5$ à
fréquence et section d’action égales. Le changement $n\mapsto-n$
n’est pas le retour en arrière complet : $H_{-n}\ne H_n^{-1}$ en général.
La véritable inversion matricielle conserve $\tau$ en
dimension deux. L’asymétrie paramétrique de cette collection
conserve ainsi sa distinction par rapport à un renversement complet.

Pour $C_{\mathrm{aux}}$,
$\tau_C=4\sin^2a\sinh^2(2b)$. La composition énergétique donne
\begin{equation}
 \boxed{\Delta E_v(C_{\mathrm{aux}})
 =\frac{hf}{9}\sin^2a\sinh^2(2b),}
 \label{md:aux:energia-compensado}
\end{equation}
\begin{equation}
 \mathcal W_{\mathrm{conjunta}}-\mathcal W_{\mathrm{local}}
 =hf\left(\sqrt{1+\frac{32}{81}\sin^2a\sinh^2(2b)}-1\right).
 \label{md:aux:trabajo-compensado}
\end{equation}
La connexion avec Planck demeure explicite : dans la section
précédente, $hf=2\pi s_0H_5(\alpha,\varphi)f$ ; dans celle du retour,
$H_5$ est remplacé par $\eta_{\mathrm{ret}}$. La composition
maintient fixée la section utilisée ; elle ne suppose pas une
variation expérimentale de $\alpha$ ni une alternance libre
entre sections dimensionnelles.

\subsection{Boltzmann et température équivalente de l’état passif}
\label{md:aux:temperatura}

Pour $\nu>1$, l’état passif a pour température équivalente
\begin{equation}
 k_BT_{\mathrm{pas}}
 =\frac{hf}{\log[(\nu+1)/(\nu-1)]}.
 \label{md:aux:temperatura-pasiva}
\end{equation}
Lorsque l’on utilise l’horloge du transducteur,
$hf=k_B\Theta_{\mathrm{clk}}\log3$, il en résulte
\begin{equation}
 \frac{T_{\mathrm{pas}}}{\Theta_{\mathrm{clk}}}
 =\frac{\log3}{\log[(\nu+1)/(\nu-1)]}.
 \label{md:aux:razon-temperaturas}
\end{equation}
Pour $n=1$, on retrouve $\log3/\log10$. L’entropie est
celle du spectre gaussien, avec l’occupation $(\nu-1)/2$.
La composition énergétique conserve cette lecture thermique
préalablement construite. La température équivalente de
l’état passif n’implique pas que l’état réduit antérieur
ait été à l’équilibre ni que de la chaleur ait été transférée.

\subsection{Spécificité et conditions de confrontation}
\label{md:aux:contraste-trabajo}

La relation générale entre cohérence, corrélations et
extraction de travail appartient à la littérature existante.
La composition examinée spécifie une loi conjointe :
coefficient $32/9$, répartition énergétique $1:8$, collection
orientée $H_n$, échelle d’action et différence entre
accès local et conjoint.

On distingue deux niveaux expérimentaux. L’émulation
réalise $U_H$ et vérifie ses identités ; elle vérifie cette
réalisation du modèle. La prédiction physique exige de
déterminer prospectivement, à partir de la construction, le
dispositif, les degrés de liberté, la fréquence, la
préparation et le couplage qui réalisent ses objets,
sans ajuster $8/9$ aux données de la confrontation. On mesure
séparément l’énergie, la covariance, la pureté, les coûts de
contrôle et l’énergie de restauration.

Une confrontation indépendante obtient la température et la pureté
par des lecteurs indépendants de la loi qu’elle entend
vérifier. Le travail récupéré est comparé à un
bilan comprenant la préparation, les contrôles et la restauration
de la mémoire.

Les comparaisons externes comprennent des études
d’ergotropie gaussienne et de passivation,\footnote{Kua, Serafini
et Genoni, \emph{Daemonic ergotropy of Gaussian quantum
states and the role of measurement-induced purification
via general-dyne detection},
\href{https://arxiv.org/abs/2506.22288}{arXiv:2506.22288}.
Leur rôle est la comparaison avec des formules générales
d’énergie, de pureté et de passivité ; elles ne valident ni le couplage HMT
ni ses constantes.}
l’extraction d’ergotropie cohérente avec des coûts
énergétiques explicites,\footnote{\emph{Experimental
Extraction of Coherent Ergotropy and Its Energetic Cost
in a Superconducting Qubit},
\href{https://arxiv.org/abs/2506.16881}{arXiv:2506.16881}.}
et le contrôle cohérent entre les cycles de machines
quantiques.\footnote{Hou et al., \emph{Combining energy
efficiency and quantum advantage in cyclic machines},
\href{https://doi.org/10.1038/s41467-025-60179-5}{doi:10.1038/s41467-025-60179-5}.
L’expérience citée ne mesure pas le transport nonadique
composé dans cet appendice.}

Le résultat est une composition algébrique avec une
structure de départ explicite. Ses preuves fixent des
identités et des restrictions opérationnelles. Un avantage
expérimental observé, une nouvelle constante universelle,
une violation de la thermodynamique ou un mécanisme matériel
de transport sans chaleur nécessiteraient leurs démonstrations
et confrontations respectives ; ce ne sont pas des conclusions des
identités auxiliaires établies ici.

\section{Énumération exhaustive et détermination finie des opérateurs de transport}
\label{iii:nucleo-finito}

Cet appendice développe deux opérations finies avec des données de départ
différentes. La première énumère les conditions initiales des deux curseurs
et construit leurs émissions au moyen de la récurrence arithmétique déclarée.
La seconde détermine les opérateurs de transport à partir du registre
régional des transitions. La séparation de leurs domaines permet de préciser
ce que produit chaque démonstration : l’énumération construit le catalogue des
émissions ; la reconstruction matricielle détermine les opérateurs compatibles
avec le registre. Le théorème de reconstruction n’est pas employé comme démonstration
de la production antérieure de ce registre.

Les primitives de la première opération sont le support $I_9^2$, les quatre
orientations cardinales, le sélecteur tritique, la lecture antérieure au
déplacement, les inversions d’orientation et la règle d’émission.
Leurs définitions sont réunies ci-dessous. Celles de la seconde sont les
quatre matrices de transitions de~\eqref{iii:nf-registro}.
Leur reconstruction unique est démontrée dans le
théorème~\ref{iii:nf-unicidad} ; le calendrier qui en résulte est exprimé
dans~\eqref{iii:nf-calendario}. Ces déterminations conservent
les données de départ déclarées pour chaque opération.

\subsection{Domaine complet des conditions initiales}

Soit $I_9=\{0,\ldots,8\}$ et soit
\[
 \mathcal D=\{N,E,S,O\},\qquad
 v_N=(-1,0),\quad v_E=(0,1),\quad
 v_S=(1,0),\quad v_O=(0,-1).
\]
Une condition initiale de curseur est $s=(i,j,d)\in
\mathcal S=I_9^2\times\mathcal D$. Les deux curseurs sont indépendants
dans le domaine initial
\begin{equation}
 \Omega=\mathcal S_+\times\mathcal S_\times,
 \qquad |\mathcal S|=9^2\cdot4=324,
 \qquad |\Omega|=324^2=104\,976.
 \label{iii:nf-dominio}
\end{equation}
L’indice $i$ correspond à l’étiquette positive $i+1$ d’APP. Pour
des entiers positifs, la réduction positive est $\rho_9(n)=1+[(n-1)]_9$,
où $[a]_b$ désigne le représentant dans $\{0,\ldots,b-1\}$.
Les évaluations résiduelles des deux feuillets sont
\[
 \sigma(i,j)=\rho_9(i+j+2),\qquad
 \varpi(i,j)=\rho_9((i+1)(j+1)).
\]
Le lecteur multiplicatif fini utilise
\begin{equation}
 \begin{array}{c|rrrrrrrrr}
 a&1&2&3&4&5&6&7&8&9\\ \hline
 \ell_9(a)&0&1&0&2&5&0&4&3&0
 \end{array}.
 \label{iii:nf-log}
\end{equation}
Sur les unités, $\ell_9$ est l’exposant discret en base $2$.
Sur $3,6,9$, la valeur $0$ constitue l’extension déclarée de
l’observable d’émission ; l’appartenance à ce secteur demeure dans le
registre de trajectoire.

Pour $1\le t\le54$, on définit
\[
 r_t=1+[(t-1)]_9,
 \qquad
 \tau_t=
 \begin{cases}
 +1,&r_t\in\{1,4,7\},\\
 -1,&r_t\in\{2,5,8\},\\
 0,&r_t\in\{3,6,9\}.
 \end{cases}
\]
Chaque curseur conserve sa position $z_t$ immédiatement avant le pas
$t$ et son vecteur d’orientation $v_t$. Le curseur additif est actif
lorsque $\tau_t=+1$ ; le multiplicatif, lorsque $\tau_t=-1$. Pour chacun
d’entre eux, avec l’indicateur d’activité $a_t$, la mise à jour est
\begin{equation}
 z_{t+1}=[z_t+a_tv_t]_9,
 \qquad
 v_{t+1}=
 \begin{cases}
 -v_t,&t\in\{27,54\},\\
 v_t,&t\notin\{27,54\}.
 \end{cases}
 \label{iii:nf-actualizacion}
\end{equation}
La lecture du pas $t$ est effectuée en $z_t$, avant le déplacement.
L’événement neutre conserve les deux positions. Cette convention fixe
sans ambiguïté les extrémités des fenêtres et l’effet des
inversions cardinales.

Dans les fenêtres $W_m=\{9m-8,\ldots,9m\}$, $1\le m\le6$, soient
\begin{align}
 S_m(s_+)&=\sum_{\substack{t\in W_m\\\tau_t=+1}}
                 \sigma(z_t^+),&
 E_m(s_\times)&=\sum_{\substack{t\in W_m\\\tau_t=-1}}
                 \ell_9(\varpi(z_t^\times)),
 \label{iii:nf-firmas}\\
 s_m&=[S_m]_{10},&e_m&=[E_m]_{10}.\notag
\end{align}
Dans chaque fenêtre, il y a trois événements de chaque type. Les signatures $s$ et $e$
sont les six composantes de ces résidus, non les entiers $S_m,E_m$
antérieurs à leur réduction.

\subsection{Recensement des signatures et des émissions}

L’émission d’une paire de conditions initiales est déterminée par
\begin{equation}
 U_m=100s_m+10[s_m+e_m]_{10}+[3s_m+5e_m+1]_{10},
 \qquad 1\le m\le6.
 \label{iii:nf-emision}
\end{equation}
Le terme final $1$ est $[7\cdot3]_{10}$, correspondant aux trois
événements neutres. Les coefficients de cette règle font partie de la
récurrence déclarée. Les constantes réelles étudiées ultérieurement
n’interviennent ni dans l’énumération de son domaine ni dans son évaluation.

\begin{proposition}[Factorisation injective du catalogue fini]
\label{iii:nf-inyectividad}
L’émission distingue chaque paire de signatures $(s,e)$. Il y a $18$ signatures
additives et $26$ multiplicatives ; leurs paires produisent $468$ émissions.
Les multiplicités de ces émissions sont $144$, $432$ et $1944$,
avec $432$, $18$ et $18$ émissions respectivement.
\end{proposition}

\begin{proof}
Le chiffre des centaines de $U_m$ restitue $s_m$. Si $b_m$ est son chiffre des
dizaines, alors $e_m=[b_m-s_m]_{10}$. Le chiffre des unités vérifie
la troisième congruence de \eqref{iii:nf-emision}. Par conséquent, le
couplage est injectif.

Le recensement s’obtient en parcourant, dans l’ordre cartésien complet,
$i=0,\ldots,8$, $j=0,\ldots,8$ et $d=N,E,S,O$, et en appliquant les $54$
pas de \eqref{iii:nf-actualizacion}. Les $324$ conditions additives
se répartissent entre les $18$ signatures suivantes, chacune avec $18$
préimages. Dans les tableaux, une suite de six chiffres représente
les six composantes ordonnées d’une signature.
\begin{center}
\begin{tabular}{ccc}
$122564$&$122898$&$212645$\\
$212988$&$221456$&$221889$\\
$456221$&$465898$&$546988$\\
$564122$&$645212$&$654889$\\
$889221$&$889654$&$898122$\\
$898465$&$988212$&$988546$
\end{tabular}
\end{center}
Pour le curseur multiplicatif, les signatures et leurs multiplicités sont
\begin{center}
\begin{tabular}{rr@{\qquad}rr}
Signature&Multiplicité&Signature&Multiplicité\\ \hline
$000000$&$108$&$555555$&$24$\\
$177393$&$8$&$555717$&$8$\\
$177555$&$8$&$555771$&$8$\\
$177771$&$8$&$555933$&$8$\\
$339555$&$8$&$717555$&$8$\\
$339771$&$8$&$717717$&$8$\\
$339933$&$8$&$717933$&$8$\\
$393177$&$8$&$771177$&$8$\\
$393393$&$8$&$771339$&$8$\\
$393555$&$8$&$771555$&$8$\\
$555177$&$8$&$933339$&$8$\\
$555339$&$8$&$933555$&$8$\\
$555393$&$8$&$933717$&$8$
\end{tabular}
\end{center}
Ces tableaux résultent des sommes finies \eqref{iii:nf-firmas}, dont le
domaine, l’ordre de lecture et les inversions sont déjà spécifiés. Leurs sommes
de multiplicités sont $18\cdot18=324$ et
$24\cdot8+24+108=324$. Le domaine cartésien
\eqref{iii:nf-dominio} permet de réaliser n’importe quelle paire de signatures.
L’injectivité fournit $18\cdot26=468$ émissions et les
multiplicités
\[
 (18\cdot24,\ 18\cdot8)=(432,144),\qquad
 (18,\ 18\cdot24)=(18,432),\qquad
 (18,\ 18\cdot108)=(18,1944),
\]
où chaque paire indique le nombre d’émissions et la multiplicité de sa fibre.
Enfin,
\[
 432\cdot144+18\cdot432+18\cdot1944=104\,976.
\]
L’énumération épuise ainsi le domaine complet des conditions initiales.
\end{proof}

La réduction ternaire orientée est
\[
 \Psi(U)=([-U_1]_3,\ldots,[-U_6]_3).
\]
Appliquée aux $18\cdot26$ émissions des tableaux précédents,
elle produit le recensement des fibres suivant. Ici, $r$ compte les émissions
distinctes, non les conditions initiales.
\begin{equation}
 \begin{array}{c|rrrrrr}
 r&1&2&3&4&5&6\\ \hline
 \#\{w:\#\Psi^{-1}(w)=r\}&132&66&18&3&6&18
 \end{array}.
 \label{iii:nf-censo-ternario}
\end{equation}
La somme de la deuxième ligne est $243$, tandis que sa somme pondérée
par $r$ est $468$. On distingue ainsi l’image de $243$ mots
de l’espace ambiant $\mathbb F_3^6$, qui contient $729$ mots. Le certificat
joint conserve toutes les signatures avec leurs préimages, les $468$
émissions et leur réduction ternaire ; les formules précédentes fournissent
une procédure complète pour les reconstruire.

\subsection{Registre des transitions et reconstruction des relèvements}

Dans cette partie, on fixe le registre régional du propriétaire cité.
Soient $b_0^\chi,\ldots,b_4^\chi\in\mathbb F_3^6$, avec
$\chi\in\{\mathsf C,\mathsf P,\mathsf A\}$, ses trois suites.
Le registre organise les transitions $b_0\to b_1\to b_2$ dans
$X_0,Y_0$ et les transitions $b_2\to b_3\to b_4$ dans $X_1,Y_1$,
avec des vecteurs lignes et le même ordre régional :
\begingroup\small
\begin{equation}
\begin{aligned}
X_0&=\begin{pmatrix}
0&1&0&2&1&1\\0&1&2&2&2&2\\2&0&1&1&0&1\\
1&2&1&2&2&1\\1&2&1&2&0&0\\1&1&2&2&0&2
\end{pmatrix},&
Y_0&=\begin{pmatrix}
0&1&2&2&2&2\\0&1&0&2&1&1\\1&2&1&2&2&1\\
1&0&2&0&1&1\\1&1&2&2&0&2\\1&2&1&0&2&0
\end{pmatrix},\\[1ex]
X_1&=\begin{pmatrix}
0&1&0&2&1&1\\0&0&2&1&1&1\\1&0&2&0&1&1\\
0&1&2&2&2&2\\1&2&1&0&2&0\\0&1&0&2&1&0
\end{pmatrix},&
Y_1&=\begin{pmatrix}
0&0&2&1&1&1\\1&1&0&2&2&1\\0&1&2&2&2&2\\
1&0&2&0&1&1\\0&1&0&2&1&0\\0&1&0&2&0&0
\end{pmatrix}.
\end{aligned}
\label{iii:nf-registro}
\end{equation}
\endgroup
La production précédente s’exprime chez son propriétaire au moyen de la
composition de sélection, de transport et de mise à jour :
\[
 x_{j+1}^{\chi}
 =\mathcal U_{9j+9}\circ\cdots\circ\mathcal U_{9j+1}(x_j^{\chi}),
 \qquad b_j^\chi=\Pi_6(x_j^\chi),
 \qquad\mathcal U_t=\mathsf{Upd}_t\circ\mathsf{Tra}_t\circ\mathsf{Sel}_t.
\]
Le résultat suivant part du registre \eqref{iii:nf-registro}.
Sa démonstration permet de reconstruire les opérateurs et de vérifier la
compatibilité des transitions ; elle n’attribue pas à l’inversion matricielle
la production coordonnée des états $x_j^\chi$.

\begin{proposition}[Unicité sur le registre régional]
\label{iii:nf-unicidad}
Les systèmes $X_aL_a=Y_a$, $a=0,1$, ont une solution unique sur
$\mathbb F_3$. Les deux solutions sont
\begingroup\small
\begin{equation}
L_0=\begin{pmatrix}
2&2&2&1&2&1\\2&1&2&2&1&1\\1&0&0&1&0&2\\
0&0&1&1&1&0\\2&2&2&0&2&1\\2&1&2&1&0&0
\end{pmatrix},\qquad
L_1=\begin{pmatrix}
2&1&2&0&2&2\\1&2&1&1&2&0\\1&2&1&1&1&2\\
0&1&0&0&2&1\\2&0&2&1&0&1\\0&2&2&2&1&1
\end{pmatrix}.
\label{iii:nf-lifts}
\end{equation}
\endgroup
\end{proposition}

\begin{proof}
L’élimination dans $\mathbb F_3$ donne $\det X_0=1$ et
$\det X_1=2=-1$. Chaque application $L\mapsto X_aL$ est donc
un automorphisme de l’espace des matrices de dimension $36$ ; dans une
vectorisation par colonnes, elle est représentée par $I_6\otimes X_a$.
On obtient $L_a=X_a^{-1}Y_a$. La multiplication des matrices de
\eqref{iii:nf-lifts} par les matrices $X_a$ respectives reproduit les
deux matrices $Y_a$ de \eqref{iii:nf-registro} ; existence et unicité
sont établies sur ce registre.
\end{proof}

Les trois suites résultantes, avec une séparation explicite en blocs,
sont
\begin{align*}
 B^{\mathsf C}&=010211\mid012222\mid010211\mid002111\mid110221,\\
 B^{\mathsf P}&=201101\mid121221\mid102011\mid012222\mid102011,\\
 B^{\mathsf A}&=121200\mid112202\mid121020\mid010210\mid010200.
\end{align*}
Chaque bloc conserve ses six positions. La concaténation n’est pas
identifiée à un entier en base dix.

\subsection{Exhaustivité du calendrier et stabilité de la reconstruction}

Un calendrier binaire est $c=(c_0,c_1,c_2,c_3)\in\{0,1\}^4$.
On définit le nombre de transitions incompatibles avec le registre par
\begin{equation}
 d(c)=\sum_{\chi\in\{\mathsf C,\mathsf P,\mathsf A\}}
       \sum_{j=0}^{3}
       \mathbf1\{b_j^\chi L_{c_j}\ne b_{j+1}^\chi\}.
 \label{iii:nf-calendario}
\end{equation}
Les multiplications des suites précédentes par
\eqref{iii:nf-lifts} donnent
\[
 d(c)=3\,d_{\rm H}(c,(0,0,1,1)),
\]
où $d_{\rm H}$ est la distance de Hamming : à chacune des
quatre positions, l’opérateur alternatif ne satisfait pas les trois
transitions régionales. Le recensement complet est
\begin{center}
\begin{tabular}{cr@{\qquad}cr@{\qquad}cr@{\qquad}cr}
$c$&$d(c)$&$c$&$d(c)$&$c$&$d(c)$&$c$&$d(c)$\\ \hline
$0000$&$6$&$0100$&$9$&$1000$&$9$&$1100$&$12$\\
$0001$&$3$&$0101$&$6$&$1001$&$6$&$1101$&$9$\\
$0010$&$3$&$0110$&$6$&$1010$&$6$&$1110$&$9$\\
$0011$&$0$&$0111$&$3$&$1011$&$3$&$1111$&$6$
\end{tabular}
\end{center}
Par conséquent, $0011$ est l’unique calendrier qui reproduit les
trois suites complètes.

La sensibilité aux entrées des matrices admet également une
preuve exhaustive sans approximations. Une modification unitaire
remplace $L_a$ par $L_a+\lambda E_{ij}$, avec
$a\in\{0,1\}$, $1\le i,j\le6$ et
$\lambda\in\mathbb F_3\setminus\{0\}$. Il y a
$2\cdot6\cdot6\cdot2=144$ modifications de ce type. Comme
$X_a$ est inversible,
\[
 X_a(L_a+\lambda E_{ij})-Y_a
 =\lambda X_aE_{ij}\ne0.
\]
Toute modification détruit au moins une des transitions du
registre correspondant. Le certificat exécutable énumère en outre
les $144$ modifications et conserve, pour chacune, le nombre de
transitions non satisfaites.

La propriété démontrée est la rigidité des opérateurs par rapport au
registre régional fixé. Pour reproduire une construction antérieure
du registre lui-même, l’action coordonnée de ses opérateurs
de sélection, de transport et de mise à jour est nécessaire. Cette dépendance reste
séparée du théorème de rigidité, du catalogue fini des émissions et
des réalisations analytiques ultérieures.

\subsection{Certificat de reproduction et résidence des données}

Le programme \texttt{verificar\_nucleo\_finito.py} parcourt les $324$
conditions de chaque curseur et forme le produit des classes résultantes.
Le catalogue documentaire est chargé après cette construction pour
comparer les émissions et leurs multiplicités. Dans la seconde partie,
les quatre matrices de transitions sont extraites de la copie littérale
du propriétaire ; les solutions sont calculées par élimination exacte
modulo $3$ avant de les comparer aux matrices imprimées. Le reçu
\texttt{NUCLEO\_FINITO.json} conserve les domaines, les préimages,
les émissions, les $16$ calendriers et les $144$ modifications.

La vérification utilise des entiers et des résidus exacts. Elle ne reçoit ni valeurs
métrologiques ni développements réels des constantes. Sa portée
est celle des opérations finies et des données de départ spécifiées
dans cet appendice. Les démonstrations de compatibilité à profondeur
illimitée et les identifications physiques conservent leurs domaines
et leurs obligations d’exposition respectifs.

\section{Sources, conservation et reproduction}
Le supplément numérique conserve le document cumulatif intégral, ses trente-neuf sources, les programmes de vérification et les reçus correspondants. La narration approuvée et sa continuation sont maintenues littéralement dans cette archive. La présente édition transforme leur énonciation conversationnelle en exposition académique et consigne la correspondance entre sections, sans remplacer les originaux.

Les antécédents proviennent des sources de la série du 21 septembre 2026. Le noyau formel est incorporé à partir de sa clôture documentaire dans l'article constitutif ; les développements de l'action, de l'incidence, de la masse et de la température sont conservés avec leurs démonstrations. Le manifeste indique chaque fichier, intervalle extrait et empreinte. Ces données documentent la provenance et la conservation, tandis que les preuves du corps du texte fixent la portée mathématique des résultats.

Les programmes de composition des constantes, d'hybridation et de compatibilité confrontent des identités exactes et des exemples finis. Les diagnostics numériques sont distingués des contrôles rationnels. Les preuves de monotonie, de convergence, d'inversion et de passage à la limite restent développées dans le texte : les exemples exécutables complètent ces démonstrations.

La compilation utilise LuaLaTeX. Les sources incluses et leurs figures permettent de reproduire le PDF sans dépendre des répertoires originaux de l'auteur. Le fichier d'instructions du paquet spécifie l'ordre de compilation et le mode d'exécution des vérifications.

\clearpage\begin{thebibliography}{99}
\bibitem{mustovicary} B. Musto et J. Vicary, ``Quantum Latin Squares and Unitary Error Bases'', \emph{Quantum Information and Computation} \textbf{16}, 1318--1332 (2016). \href{https://doi.org/10.26421/QIC16.15-16-4}{doi:10.26421/QIC16.15-16-4}.

\bibitem{delascuevas2020} G. De las Cuevas, T. Drescher et T. Netzer, ``Quantum Magic Squares: Dilations and Their Limitations'', \emph{Journal of Mathematical Physics} \textbf{61}, 111704 (2020). \href{https://doi.org/10.1063/5.0022344}{doi:10.1063/5.0022344}.

\bibitem{delascuevas2023} G. De las Cuevas, T. Netzer et I. Valentiner-Branth, ``Magic Squares: Latin, Semiclassical, and Quantum'', \emph{Journal of Mathematical Physics} \textbf{64}, 022201 (2023). \href{https://doi.org/10.1063/5.0127393}{doi:10.1063/5.0127393}.

\bibitem{codata2022} P. J. Mohr, D. B. Newell, B. N. Taylor et E. Tiesinga, ``CODATA Recommended Values of the Fundamental Physical Constants: 2022'', \emph{Reviews of Modern Physics} \textbf{97}, 025002 (2025). \href{https://doi.org/10.1103/RevModPhys.97.025002}{doi:10.1103/RevModPhys.97.025002}. \href{https://physics.nist.gov/cuu/Constants/Table/allascii.txt}{Tableau des constantes}.

\bibitem{flm} I. B. Frenkel, J. Lepowsky et A. Meurman, \emph{Vertex Operator Algebras and the Monster}, Pure and Applied Mathematics, vol. 134, Academic Press, 1988.

\bibitem{borcherds} R. E. Borcherds, ``Monstrous Moonshine and Monstrous Lie Superalgebras'', \emph{Inventiones Mathematicae} \textbf{109}, 405--444 (1992). \href{https://doi.org/10.1007/BF01232032}{doi:10.1007/BF01232032}.

\bibitem{conwaysloane} J. H. Conway et N. J. A. Sloane, \emph{Sphere Packings, Lattices and Groups}, 3.\textsuperscript{e} éd., Grundlehren der mathematischen Wissenschaften 290, Springer, 1999. \href{https://doi.org/10.1007/978-1-4757-6568-7}{doi:10.1007/978-1-4757-6568-7}.

\bibitem{bipm2018} Conférence générale des poids et mesures, \emph{Resolution 1 of the 26th CGPM: On the revision of the International System of Units}, 2018. \href{https://www.bipm.org/en/committees/cg/cgpm/26-2018/resolution-1}{Résolution officielle}. \href{https://doi.org/10.59161/CGPM2018RES1E}{doi:10.59161/CGPM2018RES1E}.

\bibitem{Bose1924} Satyendra~Nath Bose.
\newblock Plancks gesetz und lichtquantenhypothese.
\newblock {\em Zeitschrift für Physik}, 26:178--181, 1924.

\bibitem{Dirac1926} Paul A.~M. Dirac.
\newblock On the theory of quantum mechanics.
\newblock {\em Proceedings of the Royal Society A}, 112:661--677, 1926.

\bibitem{Einstein1924} Albert Einstein.
\newblock Quantentheorie des einatomigen idealen gases.
\newblock {\em Sitzungsberichte der Preussischen Akademie der Wissenschaften},
  pages 261--267, 1924.

\bibitem{Fermi1926} Enrico Fermi.
\newblock Sulla quantizzazione del gas perfetto monoatomico.
\newblock {\em Rendiconti Lincei}, 3:145--149, 1926.

\bibitem{Pauli1925} Wolfgang Pauli.
\newblock Über den zusammenhang des abschlusses der elektronengruppen im atom
  mit der komplexstruktur der spektren.
\newblock {\em Zeitschrift für Physik}, 31:765--783, 1925.

\bibitem{Pathria2011} R.~K. Pathria et Paul~D. Beale.
\newblock {\em Statistical Mechanics}.
\newblock Elsevier, 3e édition, 2011.

\bibitem{bennett2003} C. H. Bennett, ``Notes on Landauer's principle,
reversible computation, and Maxwell's Demon'', \emph{Studies in History
and Philosophy of Modern Physics} \textbf{34}, 501--510 (2003).
\href{https://doi.org/10.1016/S1355-2198(03)00039-X}{doi:10.1016/S1355-2198(03)00039-X}.
\end{thebibliography}

\end{document}
